\documentclass[11pt,a4paper]{article}
\usepackage{fontspec}
\setmainfont{Latin Modern Roman}
\usepackage[margin=23mm]{geometry}
\usepackage{amsmath,amssymb,amsthm,mathtools,mathrsfs}
\usepackage{longtable,booktabs}
\usepackage[unicode,hidelinks]{hyperref}
\providecommand{\C}{\mathbb C}
\providecommand{\R}{\mathbb R}
\newtheorem{theorem}{Theorem}[section]
\newtheorem{lemma}[theorem]{Lemma}
\newtheorem{proposition}[theorem]{Proposition}
\newtheorem{corollary}[theorem]{Corollary}
\newtheorem{definition}[theorem]{Definition}
\newtheorem{remark}[theorem]{Remark}
\allowdisplaybreaks
\emergencystretch=3em
\makeatletter
\renewcommand\@pnumwidth{2.5em}
\renewcommand\@tocrmarg{3.5em}
\makeatother
\hypersetup{pdftitle={Original Mixed Cohomology: Periods, Kernels and Arithmetic Control},pdfauthor={Split-Zero research programme}}
\title{Original Mixed Cohomology\\Periods, Kernels and Arithmetic Control}
\author{Split-Zero research programme}
\date{15 September 2026}
\begin{document}
\maketitle
\paragraph{Complete proof edition.}
The chapters below retain the original periods, amplitude units,
primary multiplicities, cyclic maps, source masses and four endpoint
metrics. Each chapter supplies its full calculation and proof. Earlier
finite bounds remain recorded in their original form; subsequent
chapters give their proved refinements. In particular KAF retains the
factorial denominator of OKV and improves the absolute kernel bound,
while AMT evaluates the source differences left in OMV.
\paragraph{Sources and scope.}
The complete preceding Gamma and observation proofs, including their
historical analytic providers, are supplied in the source closure.
The original all-multiplicity cyclic maps remain in OCS. The evaluated
five-member quadric orbit, exact kernel dimension, and its individual
volume conclusions state their simple-quartet period domain explicitly.
The full all-multiplicity operator estimate is given in KAF.
OCF, OKM and OKA calculate the actual kernel frame, its original
metrics and the boundary action. OMG proves the full all-multiplicity
Segre--Veronese geometry and retains the actual cyclic stabilizer.
\paragraph{Current refinement.}
AKS sharpens the complete kernel bound to
$32\log(25\log3/(4\pi))$ at the $kq$ scale. AKJ evaluates the
original residue inverse and conductor degree flag. AKP proves the
precise theta representative-change and minimum-section maps.
RUI now proves all-rank exterior control of the original finite
conductor, and CEP carries it through the exact source ideals,
four endpoint quotients and original compressed covariances.
The resulting conductor error is $o(kq)$ for each fixed actual period.
CTV now gives the exact outer-root quotient and its controlled
norm tails. ROQ calculates the actual residual observation, all its
positive increments and full graph metric. OSP proves a full-norm
comparison with explicit moments of the original Gamma sources.
OAR and PMR carry these finite errors into the arithmetic return.
OFV now evaluates the outer contribution at the $kq$ scale, using
the full EIQ equilibrium and QLG finite partition estimate supplied
below. FMC constructs fixed original coefficient frames with their
exact Gamma and polynomial-degree maps. OVR propagates the evaluated
outer term into the original kernel, residual observation, arithmetic
boundary and signed mixed return. Its remaining analytic quantity is
the original residual Gamma compression.
OOQ controls every direction of the complete original outer quotient.
GIP gives the full residual source and relation graph, with the exact
kernel--cokernel sequence of its outer-coordinate map. PCL evaluates
the original conductor limit and proves its pivot restriction.
ROG straightens the complete low-degree relation graph with a proved
subleading exterior cost. RDS supplies the exact dual differential
realization and original Gamma coefficient map of that relation quotient.
FEX gives the complete fixed-period exterior transfer; FGC proves its
exact comparison with the original graph, retaining both low-degree
metric corrections and the actual residual observation. ECL evaluates
the original equilibrium coefficient with its explicit certified interval.
EQR carries that calculation into the original earlier centre estimate
while preserving its full error radius. RWC proves strict original
full-window residual growth and every exterior coefficient. OER
propagates the complete fixed-column and window bounds into the
earlier kernel, boundary, signed and same-class control equations.
MCE computes the same arithmetic class's action control, including
every kernel, boundary and mixed exterior term. MCP gives the exact
marked-pairing action and source-derivative identities. RFX proves the
full arithmetic and metric reflection, evaluates the total phase at
the original source, and retains the individual mixed pairings.
OMR carries these calculations into the earlier receiving equations
and the complete arithmetic control-radius budget.
RCF calculates the complete common factor of the actual polynomial
relations and its exact quotient sequence. RCM carries the resulting
quotient through both original Gamma source orders and all four windows,
retaining the actual residual subspace and its metric.
DNE realizes that weighted relation metric by its minimum dual
extension with the complete numerator, source coefficients and masses.
PEC controls its positive exterior expansion and transports the exact
negative spectral deficits into the original kernel, residual and
arithmetic action bounds. Every remaining term is kept in its original
space with the full coordinate maps.
ACR gives finite formulas for the actual common-factor map on the
original kernel, including every coefficient-support case. LBRE proves
the complete boundary exclusions at the original period limit, and ASR
gives an explicit support radius and finite exceptional-period set.
Their zero common-factor conclusion retains the entire complementary
residual space and its metric; all exceptional cases keep their exact maps.
PBR composes these full calculations through the original kernel, residual,
low-image and arithmetic maps, retaining both spectral deficits and the
original centred finite-radius estimate. LJR proves the exact low-jet
quotient Jacobian and the leading term of both low-window deficits, with
the complete finite remainder. LRR installs that result into the same
four-window signed return and arithmetic action calculation. HJR proves
the compatible later-jet graph, both native conditioning denominators and
their exact positive row increments. It transports the original finite
total bound and cancels the common low-window term in the signed deficits.
HCR substitutes the two complete row sums in the original mixed-return
interval and action budget, retaining every finite flag, low-image and
arithmetic error in that same calculation.
All preceding chapters remain in full, with their proved refinements.
\tableofcontents
\clearpage

\clearpage
% BEGIN COMPLETE OCS
% New original-observation calculation; the sealed cut22 is not edited.
\section{The literal cyclic sectors and the original mixed kernel}
\label{sec:OCS}

This calculation evaluates the cyclic projector already present in
the original period observation. Its coefficients retain the original
period integrals, amplitude unit, and all primary nilpotents. The
result includes an exact finite matrix for the actual rank, as well
as unconditional rank bounds obtained from the original symmetric
tensor geometry.

\subsection{Original data and the Fourier basis of the literal matrix}
Use the following data of OPG1--5, including the same contours and
their orientations:
\[
 \begin{gathered}
 \alpha_{\epsilon,\eta}=\tfrac12+\epsilon\delta+i\eta\gamma,
 \quad\epsilon,\eta\in\{+1,-1\},\quad0<\delta<\tfrac12,\quad\gamma>2,\\
 h(z_1)=\prod_{\epsilon,\eta}(z_1-\alpha_{\epsilon,\eta})^m,
 \quad n=4m,\quad N=n+1,\quad E_h=\mathbb C[z_1]/(h),\\
 k=4l+1,\quad l,m\geq1,\quad c=k/2,
 \quad e=1+k(m-1),\quad q=e(k+1)^2,\\
 \beta_{a,b}=c+(2a-k)\delta+i(2b-k)\gamma,
 \quad\chi(S)=\prod_{a,b=0}^k(S-\beta_{a,b})^e,
 \quad E=\mathbb C[S]/(\chi).
 \end{gathered}\tag{OCS1}
\]
Here $A_h$ and $A$ are multiplication by $z_1$ and $S$ on their
respective algebras, and $Y=(A-cI)/i$. The original amplitude
unit is $U=M_{j_hv_h}$, with $v_h=(2\xi)/h$ and every original
$v_h(\alpha)\ne0$. For the specified nonzero $u$ and its branch,
$\Phi_h'=h$, $\Phi_h(0)=0$, and
\[
 \Gamma_j=-\ell_0+\ell_j,\qquad
 \arg\ell_j=\frac{\arg u+\pi+2\pi j}{N},\qquad
 (\Pi[P])_j=\int_{\Gamma_j}e^{\Phi_h(z_1)/u}
                       \operatorname{rem}_hP(z_1)\,dz_1.
 \tag{OCS2}
\]
The representative has degree less than $n$. The convergence and
invertibility of this original period map are the complete period
theorems retained in OPG4. No period value or unit derivative is
chosen in the present calculation. The original injection and
observation are
\[
\begin{gathered}
 \iota[P]=[P(z_1+\cdots+z_k)],\quad
 \mathcal V=\Pi^{\otimes k}U^{\otimes k}\iota,
 \\ P_k=\frac1N\sum_{a=0}^{N-1}(C^{-a})^{\otimes k},
 \\ L=P_k\mathcal V=\jmath\Lambda,
 \quad B=\operatorname{im}L,\quad K=\ker\Lambda.
\end{gathered}\tag{OCS3}
\]
The map $\jmath:B\hookrightarrow(\mathbb C^n)^{\otimes k}$
is the original image inclusion. The injection $\iota$ has precisely
the defining polynomial in OCS1: on each ordered primary tensor
the sum of its $k$ independent nilpotents has index
$e=1+k(m-1)$, its top coefficient being
$(k(m-1))!/((m-1)!)^k\ne0$, and its distinct eigenvalues are
the displayed $\beta_{a,b}$. Thus $\mathcal V$ is injective.

On a period-coordinate column $x\in\mathbb C^{N-1}$ the literal
matrix in OPG4--5 is
\[
 (Cx)_j=x_{j+1}-x_1\quad(1\leq j<N-1),\qquad
 (Cx)_{N-1}=-x_1.
 \tag{OCS4}
\]
Set $\zeta=\exp(2\pi i/N)$ and, for $1\leq t\leq N-1$, let
\[
 (f_t)_j=\zeta^{jt}-1,\qquad
 F=[f_1\ \cdots\ f_{N-1}],\qquad
 (F^{-1}x)_t=\frac1N\sum_{j=1}^{N-1}\zeta^{-tj}x_j.
 \tag{OCS5}
\]
The last expression is the exact inverse. Indeed
$\sum_{j=1}^{N-1}\zeta^{-tj}(\zeta^{uj}-1)=N\mathbf1_{t=u}$
for $t,u\ne0$ modulo $N$, by the finite geometric sum.
The first $N-2$ rows of OCS4 give
$\zeta^{(j+1)t}-\zeta^t=\zeta^t(\zeta^{jt}-1)$, and the
last row gives $1-\zeta^t=\zeta^t(\zeta^{(N-1)t}-1)$.
Consequently $Cf_t=\zeta^tf_t$, proving the orientation of
every inverse power in OCS3, as well as $C^N=I$.
The original Euclidean Gram of this basis is
\[
 F^*F=N(I_{N-1}+\mathbf1\mathbf1^*).
 \tag{OCS6}
\]
To verify it, expand
$\sum_j(\zeta^{-jt}-1)(\zeta^{ju}-1)$ and use the same
geometric sums; its value is $N(\mathbf1_{t=u}+1)$.
Thus the complete coefficient and metric transformations are
both specified.

\subsection{The ordered-to-symmetric map and the exact sector projector}
Let
\[
 \mathcal A_k=\{\nu\in\mathbb Z_{\geq0}^{N-1}:\sum_t\nu_t=k\},
 \quad\operatorname{wt}(\nu)=\sum_{t=1}^{N-1}t\nu_t\pmod N,
 \quad\mathcal I_{k,\tau}=\{\nu\in\mathcal A_k:
                                      \operatorname{wt}(\nu)=\tau\}.
 \tag{OCS7}
\]
For each $\nu$ define the following unscaled symmetric tensor:
\[
 S_\nu=\sum_{\substack{(t_1,\ldots,t_k)\in\{1,\ldots,N-1\}^k\\
                   \#\{j:t_j=t\}=\nu_t\ \text{for each }t}}
                    f_{t_1}\otimes\cdots\otimes f_{t_k}.
 \tag{OCS8}
\]
Every distinct ordered word occurs once. Since the ordered tensors
in the $f_t$ form a basis, their disjoint word supports show that
the $S_\nu$ are linearly independent. A permutation-invariant tensor
has equal coefficients on each word orbit, so they span exactly
$\operatorname{Sym}^k(\mathbb C^{N-1})$, realized as that invariant
subspace of the full ordered tensor product. In particular, for
$x=(x_1,\ldots,x_{N-1})$,
\[
 (Fx)^{\otimes k}=\sum_{\nu\in\mathcal A_k}x^\nu S_\nu,
 \quad
 (C^{-a})^{\otimes k}S_\nu=
                         \zeta^{-a\operatorname{wt}(\nu)}S_\nu,
 \quad
 P_k(Fx)^{\otimes k}=\sum_{\nu\in\mathcal I_{k,0}}x^\nu S_\nu.
 \tag{OCS9}
\]
The finite average in OCS3 is exactly one on weight zero and
zero on every other weight, by its geometric sum. No multinomial
coefficient is omitted: the multiplicity $k!/\prod_t\nu_t!$
is already the number of summands in OCS8.

Write $\mathcal F_k:\mathbb C^{\mathcal A_k}\to
\operatorname{Sym}^k(\mathbb C^{N-1})$ for the basis isomorphism
$e_\nu\mapsto S_\nu$. Let $J_0$ be the coordinate inclusion of
$\mathbb C^{\mathcal I_{k,0}}$ into $\mathbb C^{\mathcal A_k}$,
$\pi_0$ the coordinate restriction in the opposite direction,
and let $\mathcal F_0=\mathcal F_kJ_0$, with its target the original
ordered tensor space. The exact projector factorization is
\[
 P_k\mathcal F_k=\mathcal F_0\pi_0.
 \tag{OCS10}
\]
All images of $\mathcal V$ are symmetric: polynomial evaluation
at a sum is invariant under every factor permutation, and the
same map $\Pi U$ is applied in all factors. Therefore OCS10
applies on the entire original image $\mathcal V(E)$.

For completeness, the full inherited metric is retained in these
coordinates. Put $G=F^*F$ as in OCS6. The Gram of the basis OCS8 is
\[
 \mathcal G_{\nu,\mu}=
 \sum_{\substack{a_{tu}\geq0\\\sum_u a_{tu}=\nu_t\ ,\
                                    \sum_t a_{tu}=\mu_u}}
       \frac{k!}{\prod_{t,u}a_{tu}!}\prod_{t,u}G_{tu}^{a_{tu}}.
 \tag{OCS11}
\]
Indeed the coefficient of $\bar x^\nu y^\mu$ in
$(\bar x^TGy)^k$ is the right side by the multinomial theorem.
Taking the inner product of the two pure tensors in OCS9 gives
the same coefficient as $\langle S_\nu,S_\mu\rangle$.
This proves OCS11, including its off-diagonal entries; it is
positive definite because OCS8 is a basis. On the full symmetric
coordinates the projector matrix is $J_0\pi_0$, and its adjoint
for this exact Gram is
$\mathcal G^{-1}(J_0\pi_0)^*\mathcal G$.
Thus the Fourier-sector computation makes no unproved
orthogonality or norm assertion about the original cyclic average.

\subsection{Every primary period coefficient of the projected curve}
For each original root $\alpha$ let $E_{\alpha,a}$, $0\leq a<m$,
be the degree-less-than-$4m$ primary representatives of OPG10.
Explicitly, if $h_\alpha=h/(z_1-\alpha)^m$, put
\[
 E_{\alpha,a}(z_1)=\operatorname{rem}_h\left[
 h_\alpha(z_1)\sum_{j=0}^{m-1}
       \frac{(h_\alpha^{-1})^{(j)}(\alpha)}{j!}
                  (z_1-\alpha)^{j+a}\right],\qquad
 b_{j;\alpha,a}=\int_{\Gamma_j}e^{\Phi_h(z_1)/u}
                                    E_{\alpha,a}(z_1)\,dz_1.
 \tag{OCS12}
\]
The four primary factors are coprime. The Taylor inverse gives
constant one on its own factor and zero on the other three;
thus the classes $[E_{\alpha,a}]$ form the full primary basis
with all $m$ nilpotent degrees at every root.
Define the exact finite Fourier and unit coefficients
\[
 \widehat b_{t;\alpha,a}=
     \frac1N\sum_{j=1}^{N-1}\zeta^{-tj}b_{j;\alpha,a},
 \qquad
 f_{t;\alpha,d}=\frac{i^{-d}}{d!}
       \sum_{a=d}^{m-1}\widehat b_{t;\alpha,a}
                         \frac{v_h^{(a-d)}(\alpha)}{(a-d)!},
 \quad0\leq d<m.
 \tag{OCS13}
\]
Let $\omega_\alpha=(\alpha-1/2)/i$ and set
\[
 \mathbf v(w)=\Pi U\exp\bigl(w(A_h-\tfrac12I)/i\bigr)[1],
 \quad x(w)=F^{-1}\mathbf v(w),\qquad
 x_t(w)=\sum_{\alpha}e^{\omega_\alpha w}
                          \sum_{d=0}^{m-1}f_{t;\alpha,d}w^d.
 \tag{OCS14}
\]
To verify the final equality, in the primary algebra at $\alpha$
multiply the Taylor series of $v_h$ by the terminating series of
$\exp(w(z_1-\alpha)/i)$. The coefficient on
$[E_{\alpha,a}]$ is
$e^{\omega_\alpha w}\sum_{d=0}^a
v_h^{(a-d)}(\alpha)(w/i)^d/((a-d)!d!)$.
Apply the original polynomial period map and the inverse in OCS5;
this gives precisely OCS13--14. In particular, the unreduced
entire function has passed through its full primary remainder
before any period is integrated.

For $\nu\in\mathcal I_{k,0}$, $0\leq a,b\leq k$ and
$0\leq D<e$, define $c_{\nu;a,b,D}$ by the following finite
multinomial rule. Sum over integers $n_{t,\alpha,d}\geq0$ satisfying
\[
 \begin{gathered}
 \sum_{\alpha,d}n_{t,\alpha,d}=\nu_t\quad\text{for each }t,\qquad
 \sum_{t,\alpha,d}d\,n_{t,\alpha,d}=D,\\
 \sum_{t,d,\alpha:\operatorname{Re}\alpha=1/2+\delta}
                                 n_{t,\alpha,d}=a,\qquad
 \sum_{t,d,\alpha:\operatorname{Im}\alpha=+\gamma}
                                 n_{t,\alpha,d}=b.
 \end{gathered}
\]
For these indices put
\[
 c_{\nu;a,b,D}=
 \sum_{\text{stated indices}}
   \left(\prod_{t=1}^{N-1}
             \frac{\nu_t!}{\prod_{\alpha,d}n_{t,\alpha,d}!}\right)
        \prod_{t,\alpha,d}f_{t;\alpha,d}^{\,n_{t,\alpha,d}}.
 \tag{OCS15}
\]
Empty index sets give zero and zero exponents give factors one.
Every original primary coefficient is present in this finite rule.
The ordinary multinomial expansion of OCS14 now proves
\[
 x(w)^\nu=\sum_{a,b=0}^k e^{\omega_{a,b}w}
                   \sum_{D=0}^{e-1}c_{\nu;a,b,D}w^D,
 \qquad
 \omega_{a,b}=(\beta_{a,b}-c)/i
                 =(2b-k)\gamma-i(2a-k)\delta.
 \tag{OCS16}
\]
Indeed $D\leq k(m-1)=e-1$ and the counts of each original
sign give exactly the displayed exponent. The sums with different
$(a,b)$ are distinct since their real and imaginary parts are
independently distinct.

\subsection{The actual finite observation matrix and its complete kernel}
Define the primary coefficient isomorphism
$\mathcal H:\mathbb C^q\to E$ by sending coordinate $(a,b,D)$
to $\varepsilon_{a,b}(S-\beta_{a,b})^D$, where
$\varepsilon_{a,b}$ is obtained from $\chi$ by the same truncated
Taylor-inverse construction as OCS12, now of order $e$.
The proof by the coprime factors is identical and shows that this
is a basis of all $q$ primary directions. Set
\[
 \mathfrak A_{\nu;(a,b,D)}=i^D D!\,c_{\nu;a,b,D},
 \qquad
 \mathfrak A:\mathbb C^q\longrightarrow
                            \mathbb C^{\mathcal I_{k,0}}.
 \tag{OCS17}
\]
Then, on the original spaces and with the original ordered target,
\[
 \boxed{\ L\mathcal H=\mathcal F_0\mathfrak A,\qquad
 K=\mathcal H(\ker\mathfrak A),\qquad
 \operatorname{rank}\Lambda=\operatorname{rank}\mathfrak A. }
 \tag{OCS18}
\]
Here is a proof giving every factorial and phase. In the primary
sum algebra,
\[
 e^{wY}[1]=\sum_{a,b}e^{\omega_{a,b}w}
          \sum_{D=0}^{e-1}\frac{i^{-D}w^D}{D!}
                  \varepsilon_{a,b}(S-\beta_{a,b})^D.
\]
The exact original tensor identity of OPG8 gives
$\mathcal V e^{wY}[1]=\mathbf v(w)^{\otimes k}$: the centre
$c=k/2$ splits into the $k$ original centres $1/2$, while
$U$ commutes with multiplication by $z_1$.
Applying the literal projector and then OCS9 and OCS16 gives
the right-hand curve with coefficients
$\mathcal F_0c_{a,b,D}$.
Exponential polynomials $w^De^{\omega_{a,b}w}$ with these distinct
exponents are linearly independent. To prove this directly,
apply $\prod_{(a',b')\ne(a,b)}(\partial_w-\omega_{a',b'})^e$
to any zero linear combination. It kills all the other exponents.
On the selected polynomial factor, each remaining operator is
$\partial_w+(\omega_{a,b}-\omega_{a',b'})$, invertible on the
polynomials of degree less than $e$ because its constant term
is nonzero and differentiation is nilpotent there. The surviving
polynomial must therefore vanish. Doing this for every pair proves
the asserted independence. Equating coefficients now gives
$L\mathcal H e_{a,b,D}=\mathcal F_0(i^DD!c_{a,b,D})$.
This proves the first equality in OCS18 on a basis. The other
two follow from injectivity of $\mathcal F_0$ and $\jmath$.

In particular the isomorphism from the actual sector image to
the original boundary is
\[
 \mathcal Q:\operatorname{im}\mathfrak A\xrightarrow{\ \sim\ }B,
 \qquad \mathcal Qx=\jmath^{-1}\mathcal F_0x,
 \qquad
 E/K\xrightarrow{\ \sim\ }\operatorname{im}\mathfrak A,
 \quad[z]\longmapsto\mathfrak A\mathcal H^{-1}z.
 \tag{OCS19}
\]
Both inverses exist on precisely their displayed images by OCS18.
These are full coefficient maps; the inherited Gram of
$\mathcal F_0$ is $J_0^*\mathcal GJ_0$ from OCS11.

For a completely explicit finite rank prescription, define
\[
 \Delta_r(\mathfrak A)=
   \sum_{\substack{I\subseteq\mathcal I_{k,0},\ J\subseteq\{(a,b,D)\}\\
                   |I|=|J|=r}}|\det\mathfrak A_{I,J}|^2,
 \qquad\Delta_0=1.
 \tag{OCS20}
\]
Every entry of every minor is the finite expression OCS12--17
in the actual periods and amplitude derivatives. Gaussian
elimination proves that a matrix has rank at least $r$ exactly
when one of its $r$ by $r$ minors is nonzero: pivot rows and
columns supply such a minor, and a nonzero minor supplies $r$
independent columns. Hence the actual rank is exactly
$\max\{r:\Delta_r(\mathfrak A)>0\}$. This evaluates its
finite deciding data without prescribing generic values or
asserting that any unexamined maximal minor is nonzero.

The Gamma vectors used in the existing endpoint increments have
the equally explicit coordinates
\[
 \jmath\lambda_d=\mathcal F_0\widehat\lambda_d,
 \qquad
 (\widehat\lambda_d)_\nu=
 \frac1{\sqrt{h_d^{(s)}}}\sum_{a,b=0}^k\sum_{D=0}^{e-1}
               c_{\nu;a,b,D}p_d^{(D)}(\omega_{a,b}),
 \quad h_d^{(s)}=(2\pi)^{s/2}d!(s/2)_d,
 \quad s\in\{1,k\}.
 \tag{OCS21}
\]
Taylor expansion of $p_d(Y)[1]$ in the sum primary basis gives
the coefficient $i^{-D}p_d^{(D)}(\omega_{a,b})/D!$.
Multiplication by OCS17 cancels exactly these factorials and
phases, proving OCS21. This also proves that the generating
formula of OPG9 becomes
$(1+z^2)^{-s/4}\mathcal F_0
 (x(\arctan z)^\nu)_{\nu\in\mathcal I_{k,0}}$,
with the original scalar exponent and both existing sources.

\subsection{Exact invariant dimensions and unconditional forced losses}
Put $d_{N,k,\tau}=|\mathcal I_{k,\tau}|$.
By OCS7--10 this is the dimension of the full symmetric sector,
and its exact finite character formula is
\[
 d_{N,k,0}=[T^k]\frac1N
       \sum_{h\mid N}\varphi(h)\frac{1-T}{(1-T^h)^{N/h}}.
 \tag{OCS22}
\]
To prove it, the trace generating series for $C^a$ on symmetric
tensors is
$\prod_{t=1}^{N-1}(1-T\zeta^{at})^{-1}$, by summing the
monomial basis OCS8. If $\zeta^a$ has order $h$, the list
$\zeta^{at}$ with $0\leq t<N$ contains every $h$th root
$N/h$ times. Their product is
$\prod_{t=0}^{N-1}(1-T\zeta^{at})=(1-T^h)^{N/h}$.
Removing the $t=0$ factor gives the series
$(1-T)/(1-T^h)^{N/h}$. There are exactly $\varphi(h)$
indices $a$ of that order, so averaging the traces of the
original finite projector proves OCS22. The inverse powers
give the same average by $a\mapsto-a$.
In ordinary finite binomial terms, for $k\geq1$ this is
\[
 d_{N,k,0}=\frac1N\sum_{h\mid N}\varphi(h)
 \left\{
 \mathbf1_{h\mid k}\binom{k/h+N/h-1}{N/h-1}
 -\mathbf1_{h\mid(k-1)}
                  \binom{(k-1)/h+N/h-1}{N/h-1}
 \right\}.
 \tag{OCS23}
\]
The binomial expansion of $(1-T^h)^{-N/h}$ follows by
counting weak compositions, proving the coefficient extraction
used here. It follows from the actual matrix OCS17 that
\[
 \operatorname{rank}\Lambda\leq\min(q,d_{N,k,0}),
 \qquad \dim K\geq\max(0,q-d_{N,k,0}).
 \tag{OCS24}
\]
These inequalities apply to the original map, since its full
image factorization was proved in OCS18; the ambient dimension
has not been equated with its actual rank.
There is also an unconditional nonzero observation for the
original $k\geq5$. Each entire function $x_t(w)$ in OCS14 is
nonzero. Otherwise its nonzero coordinate covector would
annihilate every derivative at zero of
$\Pi U\exp(w(A_h-1/2)/i)[1]$. Those derivatives span the full
space, since $[1]$ is cyclic for multiplication by $z_1$ and
$\Pi U$ is invertible. A product of finitely many nonzero entire
functions is nonzero. The zero sector is nonempty: take the
weights $1,1,N-2$ and then $(k-3)/2$ pairs $1,N-1$.
Their sum is a multiple of $N$. Its monomial curve in OCS9
is nonzero, so $\operatorname{rank}\Lambda\geq1$.

For $m=1$, $N=5$, and every $k\geq0$, write
\[
\begin{gathered}
 D_k=\binom{k+3}{3},\qquad
 \epsilon_k=\begin{cases}1&k\equiv0\pmod5,\\-1&k\equiv1\pmod5,\\
                         0&\text{otherwise}.
             \end{cases}\\
 d_{5,k,0}=\frac{D_k+4\epsilon_k}{5},\quad
 d_{5,k,\tau}=\frac{D_k-\epsilon_k}{5}\quad(\tau\ne0).
\end{gathered}\tag{OCS25}
\]
For each of the four nonidentity powers of $C$, the trace
series is $(1-T)/(1-T^5)$, whose coefficient is $\epsilon_k$.
The weight-$\tau$ projector introduces the Fourier multiplier;
the sum of its four nontrivial fifth roots is $4$ when
$\tau=0$ and $-1$ otherwise. This proves both formulas.
For the requested original orders the exact counts and resulting
kernel lower bounds are
\[
\begin{array}{c|r|r|r|r}
 k&D_k&d_{5,k,0}&q=(k+1)^2&\max(0,q-d_{5,k,0})\\\hline
 5&56&12&36&24\\
 9&220&44&100&56\\
 13&560&112&196&84\\
 17&1140&228&324&96\\
 21&2024&404&484&80\\
 25&3276&656&676&20\\
 29&4960&992&900&0
\end{array}\tag{OCS26}
\]
These are integer evaluations of OCS25. The strict inequality
$d_{5,k,0}<q$ holds for $k=4l+1$, $l\geq1$, exactly for the
first six orders in the table. They exhaust $k\leq25$ in this
progression. For $k\geq29$,
\[
 D_k-5(k+1)^2=\frac{(k+1)(k^2-25k-24)}6\geq460>4,
\]
where $k^2-25k-24$ is positive and increasing from $k=29$.
Even the most negative correction $4\epsilon_k\geq-4$
therefore gives $d_{5,k,0}>q$. This proves the threshold
without extrapolating the finite table.

\subsection{The actual simple-quartet quadric and a stronger rank bound}
Continue with $m=1$. Order the four roots as
$(++),(+-),(-+),(--)$ and let
$\mathcal E_{\rm prim}:\mathbb C^4\to E_h$ be the primary
basis isomorphism whose columns are the four $[E_{\alpha,0}]$.
Define the actual invertible matrix and the original exponential
quartet by
\[
 \mathsf H=F^{-1}\Pi U\mathcal E_{\rm prim},\qquad
 t_{\epsilon,\eta}(w)=e^{(\eta\gamma-i\epsilon\delta)w},
 \qquad x(w)=\mathsf Ht(w).
 \tag{OCS27}
\]
Every column of $\mathsf H$ is OCS13 with $d=0$, namely the
Fourier transform of its original primary period times its
actual nonzero $v_h(\alpha)$. Each displayed map is invertible,
so $\mathsf H$ is invertible without a genericity assumption.
Put $\mathcal R_k=\mathbb C[x_1,x_2,x_3,x_4]_k$, the homogeneous
polynomials of degree $k$, and define
\[
 Q_0(t)=t_{++}t_{--}-t_{+-}t_{-+},\qquad
 Q(x)=Q_0(\mathsf H^{-1}x).
 \tag{OCS28}
\]
This is the literal transformed quadratic relation of the
original curve; its coefficients are determined by the periods
and unit through OCS27. In particular $Q\ne0$.

The kernel of homogeneous restriction to that curve is exactly
\[
 \ker\{\mathcal R_k\to\operatorname{Hol}(\mathbb C),
                         f\mapsto f(x(w))\}
            =Q\mathcal R_{k-2}\quad(k\geq2).
 \tag{OCS29}
\]
Here is a complete verification. First work in the $t$ coordinates.
A degree-$k$ monomial has exponent
$(2b-k)\gamma-i(2a-k)\delta$, where $a$ is its count of
first sign $+$ and $b$ its count of second sign $+$.
These $(k+1)^2$ complex exponents are distinct: equality of
their real parts forces $b$ equal and equality of their
imaginary parts forces $a$ equal. Equivalently, the two
frequencies $2\gamma$ and $-2i\delta$ admit no nonzero
integer relation, by their real and imaginary parts. Thus
the exponentials have no extra polynomial relation of this
homogeneous type; their linear independence follows, for
example, by differentiating $0,\ldots,(k+1)^2-1$ times and
using the Vandermonde determinant of the distinct exponents.
Every pair $(a,b)$ occurs: the count $r_{++}$ may be any integer
between $\max(0,a+b-k)$ and $\min(a,b)$, a nonempty interval,
and the other three counts are $a-r_{++}$, $b-r_{++}$,
and $k-a-b+r_{++}$.

Within any fixed pair $(a,b)$, two monomial exponent vectors
differ by an integer multiple of $(1,-1,-1,1)$. After their
common monomial is factored out, their difference is therefore
the difference of equal powers of $t_{++}t_{--}$ and
$t_{+-}t_{-+}$. The identity
$A^r-B^r=(A-B)\sum_{j=0}^{r-1}A^{r-1-j}B^j$ shows that
this difference is divisible by $Q_0$. Any polynomial
vanishing on the curve has coefficient sum zero within
each pair $(a,b)$, by the proved exponential independence.
It is thus a sum of these monomial differences and is
divisible by $Q_0$. Its quotient is homogeneous of degree
$k-2$. Conversely $Q_0(t(w))=0$, since each of its two
products is one. This proves the assertion in $t$
coordinates, and the invertible substitution OCS27 proves
OCS29. In particular
$\dim(\mathcal R_k/Q\mathcal R_{k-2})=(k+1)^2=q$;
no information about the original curve has been replaced
by a freely chosen quadric.

Let $\mathcal I_k\subseteq\mathcal R_k$ be the span of the
weight-zero monomials $x^\nu$, using the weights $1,2,3,4$
from the exact Fourier basis. The coordinate functions of
the original projected curve in OCS9 are precisely these
monomials restricted to $x(w)$. The span of the vector
curve $Le^{wY}[1]$ equals $\operatorname{im}L$: its
derivatives at zero contain $LY^d[1]$, $0\leq d<q$, and
$Y^d[1]$ is a basis because $Y$ is an invertible affine
change of multiplication by $S$ in the degree filtration.
The rank of that vector curve equals the dimension of
the span of its coordinate functions, by the equality of
row and column ranks of its finite exponential coefficient
matrix. Hence OCS29 proves the exact equality
\[
 \operatorname{rank}\Lambda
   =d_{5,k,0}-\dim(\mathcal I_k\cap Q\mathcal R_{k-2}).
 \tag{OCS30}
\]

Fix any multiplicative monomial order, for instance
lexicographic order, and write $x^\mu$ for the leading
monomial of the actual $Q$, with $|\mu|=2$. Choose a
row-echelon basis of the finite-dimensional intersection
in OCS30 in descending monomial order. Its leading
monomials are distinct, all have weight zero, and all
are divisible by $x^\mu$: the leading monomial of
$Qf$ is the product of those of $Q$ and $f$ because
the order is multiplicative and the coefficient field
has no zero divisors. Thus each such leading monomial
is $x^\mu x^\eta$ for a distinct degree-$k-2$ monomial
of weight $-\operatorname{wt}(\mu)$. Counting this set
and using OCS30 gives
\[
 \boxed{\quad
 d_{5,k,0}-d_{5,k-2,-\operatorname{wt}(\mu)}
 \ \leq\ \operatorname{rank}\Lambda
 \ \leq\ \min(q,d_{5,k,0}).\quad}
 \tag{OCS31}
\]
The left side depends only on the actual leading monomial
of the actual period quadric, and remains valid whatever
its other coefficients are. A uniform consequence retaining
the strongest count independent of that monomial is
\[
 r_k^-:=d_{5,k,0}-\max_{\tau\in\mathbb Z/5}d_{5,k-2,\tau}
       =\left\lfloor\frac{(k+1)^2}{5}\right\rfloor
                                  +\mathbf1_{5\mid k}
       \leq\operatorname{rank}\Lambda.
 \tag{OCS32}
\]
For the equality, $D_k-D_{k-2}=(k+1)^2$, as direct
subtraction of the two cubic binomials shows. Formula
OCS25 then gives
\[
 r_k^-=\frac{(k+1)^2+4\epsilon_k
                   -\max(4\epsilon_{k-2},-\epsilon_{k-2})}{5}.
\]
For $k$ congruent respectively to $0,1,2,3,4$ modulo five,
its numerator correction is respectively $4,-4,-4,-1,0$.
The corresponding residues of $(k+1)^2$ are $1,4,4,1,0$.
These five cases prove OCS32 for every $k\geq2$.
For the original orders in OCS26 the complete unconditional
intervals are therefore
\[
\begin{array}{c|r|r|r|r}
 k&r_k^-&\min(q,d_{5,k,0})&\max(0,q-d_{5,k,0})&q-r_k^-\\\hline
 5&8&12&24&28\\
 9&20&44&56&80\\
 13&39&112&84&157\\
 17&64&228&96&260\\
 21&96&404&80&388\\
 25&136&656&20&540\\
 29&180&900&0&720
\end{array}\tag{OCS33}
\]
The middle two rank columns bound $\operatorname{rank}\Lambda$;
the last two columns bound $\dim K$ from below and above,
respectively. No endpoint of either interval has been assigned
as the actual value without its original minors.

\subsection{A second exact finite matrix for the simple-quartet kernel}
The quadric calculation also gives the actual kernel a
polynomial quotient model. Write $Q(x)=\sum_{|\mu|=2}q_\mu x^\mu$.
Let $T_Q$ be the full multiplication matrix
$\mathcal R_{k-2}\to\mathcal R_k$ in the monomial bases:
\[
 (T_Q)_{\nu,\eta}=q_{\nu-\eta},
 \quad |\nu|=k,\quad|\eta|=k-2,
 \qquad q_{\nu-\eta}=0\text{ if an entry of }\nu-\eta<0.
\]
Let $T_{\rm out}$ be its rows of nonzero cyclic weight.
Every coefficient $q_\mu$ is explicitly fixed by
$Q_0(\mathsf H^{-1}x)$, and the inverse matrix may be
computed as $\operatorname{adj}\mathsf H/\det\mathsf H$,
with nonzero denominator proved in OCS27. Multiplication
by $Q$ is injective. The condition that $Qf$ belongs
to $\mathcal I_k$ is exactly $T_{\rm out}f=0$.
Thus OCS30 gives the further exact rank identities
\[
 \begin{split}
 \operatorname{rank}\Lambda
   &=d_{5,k,0}-D_{k-2}+\operatorname{rank}T_{\rm out},\\
 \dim K&=D_k-d_{5,k,0}-\operatorname{rank}T_{\rm out}.
 \end{split}\tag{OCS34}
\]
The ranks of this concrete matrix are decided by its finite
minors, by the same elimination proof as OCS20. This is
an alternative exact calculation from the actual periods,
not an assumption about the quadric's coefficients.

To give the typed map behind the kernel equality, let
$E^*$ denote the complex linear dual. Define
$\Theta:\mathcal R_k\to E^*$ as follows: $\Theta(f)$ is
the unique covector such that
\[
 \Theta(f)(e^{wY}[1])=f(x(w))\quad\text{for all }w.
 \tag{OCS35}
\]
For $m=1$ every function on the right is a linear
combination of the $q$ distinct exponentials in OCS16.
The primary expansion of $e^{wY}[1]$ gives exactly those
exponentials as a basis of covector functions. This proves
existence and uniqueness. The proof of OCS29 proves
surjectivity and $\ker\Theta=Q\mathcal R_{k-2}$.
Moreover $\Theta(\mathcal I_k)=\operatorname{im}L^\vee$,
where $L^\vee$ denotes the complex linear dual map:
the monomial coefficient functionals in the basis OCS8
pull back along $L$ to exactly OCS35 for $x^\nu$ of
weight zero. These functionals span the full dual of
the actual image, since every linear functional on a
subspace of a finite-dimensional space extends to the
whole space. Restriction of $E^*$ to $K$ has kernel
$\operatorname{im}L^\vee$: a covector vanishing on $K$
factors uniquely through $E/K\simeq\operatorname{im}L$.
Thus the exact maps give
\[
 \boxed{\quad
 \operatorname{coker}T_{\rm out}
 \simeq\mathcal R_k/(\mathcal I_k+Q\mathcal R_{k-2})
 \xrightarrow{\ \Theta|_K\ }K^*.
 \quad}\tag{OCS36}
\]
For the first isomorphism, use the direct monomial splitting
of $\mathcal R_k$ into its zero and nonzero sectors; after
quotienting by $\mathcal I_k$, multiplication by $Q$
becomes precisely $T_{\rm out}$. For the second, the
surjectivity and kernel were just proved. This identifies
the actual lost directions by a full quotient map.

\subsection{All primary action directions remain in the mixed calculation}
For arbitrary original $m\geq1$, retain the full primary
matrix $\mathfrak A$ in OCS17. In the basis $\mathcal H$
the original sum action is exactly
\[
 J e_{a,b,D}=\omega_{a,b}e_{a,b,D}
       +\begin{cases}i^{-1}e_{a,b,D+1}&D<e-1,\\0&D=e-1,
         \end{cases}
 \qquad Y\mathcal H=\mathcal HJ.
 \tag{OCS37}
\]
This follows by multiplying
$\varepsilon_{a,b}(S-\beta_{a,b})^D$ by
$(S-c)/i=\omega_{a,b}+(S-\beta_{a,b})/i$ in its original
primary algebra. Consequently the actual possible action
leaving the cyclic kernel is the fully specified map
\[
 \ker\mathfrak A\longrightarrow\operatorname{im}\mathfrak A,
 \qquad x\longmapsto\mathfrak A Jx,
 \qquad
 \Lambda Y\mathcal Hx=\mathcal Q\mathfrak A Jx.
 \tag{OCS38}
\]
The codomain assertion follows from $\mathfrak A Jx$ being
a value of the same matrix $\mathfrak A$ on $Jx$.
For the original inclusion $I:K\hookrightarrow E$, this
is the OPG22 defect transported through the exact kernel
isomorphism in OCS18. In particular none of the discarded
sectors has been declared invariant under $Y$.

The part of the kernel undetected by all original iterates
is exactly computed by the stacked matrix
\[
 \mathfrak O=\begin{bmatrix}
          \mathfrak A\\\mathfrak AJ\\\vdots\\\mathfrak AJ^{q-1}
          \end{bmatrix},\qquad
 \mathcal K_\infty=\mathcal H\ker\mathfrak O,
 \quad
 \dim\mathcal K_\infty=q-\operatorname{rank}\mathfrak O,
 \quad
 \dim(K/\mathcal K_\infty)=
        \operatorname{rank}\mathfrak O-\operatorname{rank}\mathfrak A.
 \tag{OCS39}
\]
Indeed $\mathfrak A J^a x$ is the exact sector-coordinate
value of $\Lambda Y^a\mathcal Hx$ by OCS18 and OCS37.
The polynomial $\chi(c+iT)$ annihilates $J$ and has
degree $q$ with leading coefficient $i^q$. It expresses
every higher power in the span of the first $q$, so the
stacked kernel is precisely the intersection over all
iterates and is $J$-invariant. Rank-nullity proves both
dimension formulas. Its quotient identifies the directions
that the original one-step observation loses and subsequent
original iterates detect.

Equations OCS17--21, OCS34--39 supply actual finite maps
and deciding minors for the original observation and its
kernel. The forced kernel losses in OCS26 and the rank
intervals OCS33 concern this unchanged projector and these
unchanged periods. The endpoint volume calculation still
uses the full positive sums of the actual vectors
$\lambda_d$ in OPG16--18 and their exact inherited
coordinates OCS11 and OCS19. No volume growth estimate
is deduced from these dimension bounds alone.

% END COMPLETE OCS

\clearpage
% BEGIN COMPLETE OQC
% Derived from the literal original period quadric and cyclic projector.
\section{The original quadric orbit and an explicit invariant lifting map}
\label{sec:OQC}

This calculation uses the original simple-quartet construction, with
$m=1$, $k=4l+1$, $l\geq1$, $q=(k+1)^2$, $0<\delta<1/2$, and
$\gamma>2$. The period parameter, its branch, the polynomial contours,
the complete amplitude unit, and the observation are those of OPG and
OCS. The degree-eight polynomial below is the product of four actual
orbit equations. Its degree and all of its coefficients are derived
from that orbit; no source order is selected to match an estimate.

\subsection{The actual coefficient ring and the original observation}
Order the original roots as $(++),(+-),(-+),(--)$, with
$\alpha_{\epsilon,\eta}=1/2+\epsilon\delta+i\eta\gamma$.
Let $\mathcal E_{\rm prim}$ have columns the four original primary
idempotents in $E_h=\mathbb C[z_1]/(h)$. For the original matrices
$\Pi,U,F$ in OCS2, OCS5, OCS27 put
\[
\begin{gathered}
 \mathsf H=F^{-1}\Pi U\mathcal E_{\rm prim},\qquad
 t_{\epsilon,\eta}(w)=e^{(\eta\gamma-i\epsilon\delta)w},\\
 x(w)=\mathsf Ht(w),\qquad
 Q(x)=Q_0(\mathsf H^{-1}x),\quad
 Q_0(t)=t_{++}t_{--}-t_{+-}t_{-+}.
\end{gathered}\tag{OQC1}
\]
Every matrix in $\mathsf H$ is the specified original isomorphism;
in particular its nonzero amplitude factors are the actual
$v_h(\alpha)$, and its period entries are the original convergent
polynomial integrals. Thus $\det\mathsf H\ne0$ by the included
period-isomorphism and primary-basis proofs. Define
\[
 \mathscr P_d=\mathbb C[x_1,x_2,x_3,x_4]_d,\qquad
 R_d=\mathscr P_d/(Q\mathscr P_{d-2}),\qquad
 \mathscr P_d=R_d=0\quad(d<0).
 \tag{OQC2}
\]
The matrix of $Q_0(t)=t^TM_0t$ has entries
$(M_0)_{14}=(M_0)_{41}=1/2$ and
$(M_0)_{23}=(M_0)_{32}=-1/2$, all others zero. Hence
\[
 M=\mathsf H^{-T}M_0\mathsf H^{-1},\qquad
 \det M=\frac{1}{16(\det\mathsf H)^2}\ne0.
 \tag{OQC3}
\]
In particular $Q$ is irreducible. Indeed, a reducible homogeneous
quadratic over $\mathbb C$ is a product of two linear forms, whose
symmetric coefficient matrix is a sum of two rank-one matrices and
has rank at most two. This contradicts OQC3. The polynomial ring
over a field is a unique factorization domain, by the repeated
one-variable Gauss lemma, so its irreducible element $Q$ is prime.
Consequently $R=\mathbb C[x_1,x_2,x_3,x_4]/(Q)$ is an integral domain.

For clarity, the complete homogeneous restriction calculation has
the precise kernel $Q\mathscr P_{d-2}$. In the $t$ coordinates a
degree-$d$ monomial restricts to
\[
 \exp\bigl(((2b-d)\gamma-i(2a-d)\delta)w\bigr),
       \qquad 0\leq a,b\leq d.
 \tag{OQC4}
\]
Every pair occurs: the number of $(++)$ factors can range from
$\max(0,a+b-d)$ to $\min(a,b)$. The exponents are distinct by
their real and imaginary parts. Their exponentials are linearly
independent, since differentiating orders zero through $(d+1)^2-1$
gives a Vandermonde matrix with nonzero determinant. Two monomials
giving the same pair have exponent vectors differing by a multiple
of $(1,-1,-1,1)$. After their common monomial is removed, their
difference is divisible by $t_{++}t_{--}-t_{+-}t_{-+}$ through
$A^r-B^r=(A-B)\sum_{j=0}^{r-1}A^{r-1-j}B^j$. A vanishing
polynomial has coefficient sum zero within each pair by the
independence just proved, and is therefore a sum of these differences.
Conversely $Q_0(t(w))=0$. Transport by $\mathsf H$ proves the
stated kernel and
\[
                  \dim R_d=(d+1)^2\quad(d\geq0).
 \tag{OQC5}
\]

Write $\zeta=e^{2\pi i/5}$ and
$D=\operatorname{diag}(\zeta,\zeta^2,\zeta^3,\zeta^4)$.
The literal Fourier calculation gives $CF=FD$. On polynomials
define the original induced action and its average by
\[
 (g f)(x)=f(D^{-1}x),\qquad g^5=1,\qquad
 \mathsf R f=\frac15\sum_{a=0}^4g^af,
 \qquad \mathscr I_d=\mathscr P_d^{\langle g\rangle}.
 \tag{OQC6}
\]
Thus $\mathscr I_d$ consists of the weight-zero monomials, since
$gx^\nu=\zeta^{-\sum j\nu_j}x^\nu$. This is the dual of the
original average of $(C^{-a})^{\otimes d}$; both the inverse powers
and the factor $1/5$ are retained.

For the original algebra $E_d$ of the $d$-fold sum, define
$\Theta_d:R_d\xrightarrow{\sim}E_d^\vee$ by
\[
 \Theta_d([f])(e^{wY_d}[1])=f(x(w)),\qquad
       Y_d=(A_d-dI/2)/i.
 \tag{OQC7}
\]
Here $\vee$ denotes the complex-linear dual. The primary expansion
of the vector on the left gives all distinct exponentials OQC4;
therefore existence and uniqueness of this covector and bijectivity
follow from the preceding complete restriction calculation. The
original observation $\Lambda_d:E_d\to B_d$ obeys
\[
 J_d:=\operatorname{im}(\mathscr I_d\to R_d),\qquad
 \Theta_d(J_d)=\operatorname{im}\Lambda_d^\vee,
 \qquad \dim J_d=\operatorname{rank}\Lambda_d,
 \qquad R_d/J_d\xrightarrow{\sim}(\ker\Lambda_d)^\vee.
 \tag{OQC8}
\]
Indeed the coordinate functions of the projected pure tensor in
OCS9 are precisely $x(w)^\nu$ of weight zero; their pulled-back
covectors span the dual of its actual image. A covector vanishing
on $\ker\Lambda_d$ factors through that image, which proves the
last isomorphism as well. OQC8 specifies the exact original image;
it has not been replaced by the whole invariant tensor space.

\subsection{The complete orbit decision from the actual coefficients}
Put $Q_a=g^aQ$ for $0\leq a<5$. The orbit of the line $\mathbb C Q$
has size one or five because five is prime. Moreover a size-one
orbit has $gQ=Q$, with no other possible scalar. To prove this, write
$gQ=\lambda Q$. Then $\lambda^5=1$, while OQC3 and the determinant
of $D$, which is $\zeta^{10}=1$, imply
\[
 \det(D^{-T}MD^{-1})=\det M=\lambda^4\det M.
 \tag{OQC9}
\]
Thus $\lambda^4=1$ and $\lambda^5=1$, forcing $\lambda=1$.
If any nonidentity power fixed the line, its powers generate the
same cyclic group, so the preceding argument applies. This proves
the entire orbit classification.

There is an exact coefficient test, requiring no genericity premise.
Write the actual $Q=\sum_{|\mu|=2}q_\mu x^\mu$ and put
\[
 \mathfrak e(Q)=\sum_{\substack{|\mu|=2\\
                      \sum j\mu_j\not\equiv0\ (5)}}|q_\mu|^2.
 \tag{OQC10}
\]
All coefficients here are the entries of OQC3, including the full
inverse of the original period matrix. Since the monomials are a
basis and each summand is nonnegative, $\mathfrak e(Q)=0$ exactly
when $Q$ is invariant. The only degree-two weight-zero monomials
are $x_1x_4$ and $x_2x_3$. Hence the size-one case is exactly
$Q=a x_1x_4+b x_2x_3$ with $ab\ne0$; nonvanishing of the two
coefficients follows from OQC3. Positive $\mathfrak e(Q)$ gives
the five distinct irreducible orbit equations. These are exhaustive
outcomes of a test on the given periods, not assumptions on them.

In the size-one outcome, multiplication by $Q$ preserves every
weight. If $Qf$ is invariant, then $Q(gf-f)=0$ in the polynomial
ring, and $gf=f$. Therefore
\[
 \ker(\mathscr I_k\to R_k)=Q\mathscr I_{k-2},\qquad
 \operatorname{rank}\Lambda_k=d_{5,k,0}-d_{5,k-2,0},
 \quad \dim\ker\Lambda_k=q-d_{5,k,0}+d_{5,k-2,0}.
 \tag{OQC11}
\]
This also proves that $J_k=R_k^{\langle g\rangle}$: an invariant
class can be lifted to a polynomial, and averaging that lift gives
an invariant representative. The count itself follows from
\[
 d_{5,k,0}=\frac{\binom{k+3}{3}+4\epsilon_k}{5},\qquad
 \epsilon_k=\mathbf1_{k\equiv0\ (5)}-\mathbf1_{k\equiv1\ (5)}.
 \tag{OQC12}
\]
For a proof, the identity element has symmetric-power generating
series $(1-z)^{-4}$; each of the other four powers has series
$\prod_{j=1}^4(1-z\zeta^j)^{-1}=(1-z)/(1-z^5)$.
Averaging their coefficients gives OQC12. Subtraction and
$\binom{k+3}{3}-\binom{k+1}{3}=(k+1)^2$ evaluate the rank in
OQC11, for every $k\geq2$, as
\[
 \left\lfloor\frac{(k+1)^2}{5}\right\rfloor
                 +\mathbf1_{k\equiv0\text{ or }3\ (5)}.
 \tag{OQC13}
\]
The five numerator corrections before division by five are
$4,-4,-4,4,0$, respectively; the residues of $(k+1)^2$ are
$1,4,4,1,0$. This verifies every case in OQC13.

\subsection{An explicit invariant lift for the five-member orbit}
For the size-five outcome determined by OQC10, retain the full
actual products
\[
 \mathcal H_Q=\prod_{a=0}^4Q_a\quad(\deg\mathcal H_Q=10),
 \qquad \mathcal A_Q=\prod_{a=1}^4Q_a
                                      \quad(\deg\mathcal A_Q=8).
 \tag{OQC14}
\]
The action permutes all five factors cyclically, so
$g\mathcal H_Q=\mathcal H_Q$. None of the four factors of
$\mathcal A_Q$ is a scalar multiple of $Q$. Since each is
irreducible and $Q$ is prime, multiplication by $[\mathcal A_Q]$
is injective on the integral domain $R$.

For every homogeneous $f\in\mathscr P_d$ the exact trace identity is
\[
 \left[5\mathsf R(\mathcal A_Q f)\right]_{(Q)}
                          =[\mathcal A_Q f]_{(Q)}.
 \tag{OQC15}
\]
Indeed $g^a\mathcal A_Q$ is the product of all orbit factors
except $Q_a$. For $a\ne0$ it contains $Q_0=Q$. Those four
terms vanish in the displayed quotient, and the $a=0$ term is
exactly the right side. This proves OQC15 with its original
average and full factors, including every period-dependent
coefficient.

The lift also has a well-defined typed quotient map. Let
$\mathscr U=\mathbb C[x_1,x_2,x_3,x_4]/(\mathcal H_Q)$;
the invariant equation $\mathcal H_Q$ makes the cyclic action
well defined on this quotient. Then
\[
 \begin{gathered}
 \mathfrak L_d:R_d\longrightarrow
                      (\mathscr U_{d+8})^{\langle g\rangle},\qquad
 [f]_{(Q)}\longmapsto[5\mathsf R(\mathcal A_Q f)]_{(\mathcal H_Q)},\\
 \operatorname{res}\mathfrak L_d([f])=[\mathcal A_Q f]_{(Q)},
 \qquad \operatorname{res}:\mathscr U\to R.
 \end{gathered}\tag{OQC16}
\]
If $f$ changes by $Qb$, its proposed image changes by
$5\mathsf R(\mathcal H_Qb)=5\mathcal H_Q\mathsf Rb$, hence
by zero in $\mathscr U$. This proves well-definedness. The
image is invariant because it has an invariant polynomial
representative; OQC15 proves the second line. Injectivity follows
from the injectivity of multiplication by $[\mathcal A_Q]$ on
$R$. All maps in OQC16 are thus established, including their kernels.

The actual original invariant image consequently contains
\[
 [\mathcal A_Q]R_{k-8}\ \subseteq\ J_k\ \subseteq\ R_k.
 \tag{OQC17}
\]
The first inclusion follows by taking the invariant polynomial
representatives OQC15, and uses no dimension inference. Therefore,
for every $k\geq8$ in this outcome,
\[
 \operatorname{rank}\Lambda_k\geq(k-7)^2,\qquad
 \dim\ker\Lambda_k\leq(k+1)^2-(k-7)^2=16k-48.
 \tag{OQC18}
\]
These inequalities follow from the proved injection and OQC5.
They may be combined with the unconditional actual-matrix bounds
in OCS24, OCS31--33 by taking the stronger lower and upper bounds.
For the original degrees $k=9,13,17,21,25,29$, the displayed upper
bounds on the original kernel are respectively
$96,160,224,288,352,416$; at the first two degrees the earlier
unconditional bounds may be stronger. No such upper bound is
assigned to the size-one outcome, whose exact rank is OQC11.
Together OQC10--18 compute the complete orbit alternatives from
the original coefficients. Evaluating $\mathfrak e(Q)$ in the
original period family is a separate calculation on those explicit
coefficients; no outcome is postulated here.

\subsection{The induced original dual map and its complete action defect}
At $k\geq9$ the degree $d=k-8$ is positive, and the original
$d$-fold sum algebra is defined using the same $h,\Pi,U$.
For the size-five outcome, OQC7, OQC8 and OQC17 give the uniquely
typed injective map
\[
 \begin{gathered}
 \mathfrak J_k:E_{k-8}^\vee\hookrightarrow B_k^\vee,\\
 \Lambda_k^\vee\mathfrak J_k
        =\Theta_k\,M_{[\mathcal A_Q]}\,\Theta_{k-8}^{-1},\qquad
 \mathfrak J_k^\vee:B_k\twoheadrightarrow E_{k-8}.
 \end{gathered}\tag{OQC19}
\]
The right side of the second line has image in
$\operatorname{im}\Lambda_k^\vee$ by OQC17, and
$\Lambda_k^\vee$ is injective because $\Lambda_k$ is onto its
actual image. This gives existence and uniqueness. Each of the
three right-hand maps is injective, so $\mathfrak J_k$ is
injective. Finite-dimensional duality then proves the stated
surjection, using the canonical bidual identifications. The new
degree $k-8$ records contraction against the degree-eight
polynomial derived in OQC14. It does not replace either original
Gamma source of order $1$ or $k$.

The surviving action and its defect are fully determined. Put
\[
 \Omega=\operatorname{diag}(\gamma-i\delta,-\gamma-i\delta,
                              \gamma+i\delta,-\gamma+i\delta),
 \quad G=\mathsf H\Omega\mathsf H^{-1},\qquad
 (\mathcal Df)(x)=\sum_{j=1}^4(Gx)_j\partial_jf(x).
 \tag{OQC20}
\]
Here $x'(w)=Gx(w)$ by OQC1. In the original $t$ coordinates,
the weights of each product in $Q_0$ sum to zero, so
$\mathcal DQ=0$. Thus this derivation acts on each $R_d$.
Differentiating OQC7 proves exactly
\[
          \Theta_d\mathcal D=Y_d^\vee\Theta_d.
 \tag{OQC21}
\]
Both sides evaluated at $e^{wY_d}[1]$ are the derivative of
$f(x(w))$, and these vectors span $E_d$ by the distinct
exponentials; hence equality holds as maps, not only on a curve.
Leibniz's rule gives the complete defect
\[
 \begin{split}
 Y_k^\vee\Lambda_k^\vee\mathfrak J_k
       -\Lambda_k^\vee\mathfrak J_kY_{k-8}^\vee
 &=\Theta_k M_{[\mathcal D\mathcal A_Q]}\Theta_{k-8}^{-1},\\
 \mathcal D\mathcal A_Q
 &=\sum_{a=1}^4(\mathcal DQ_a)
                  \prod_{\substack{1\leq b\leq4\\b\ne a}}Q_b.
 \end{split}\tag{OQC22}
\]
The codomain in the first line is $E_k^\vee$. Its image has not
been asserted to lie in $\operatorname{im}\Lambda_k^\vee$.
Restriction to $\ker\Lambda_k$ supplies the exact remaining
map into $(\ker\Lambda_k)^\vee$, through OQC8. In particular
OQC19 is not used as an unproved intertwining of boundary actions.
OQC22 retains every term that can leave the observed dual subspace.

Finally, the failure of the cyclic average to commute with the
original action is explicitly calculable before any quotient:
\[
 ([\mathcal D,g^a]f)(x)
   =(\nabla f)(D^{-a}x)^T
                   (D^{-a}G-GD^{-a})x,
 \qquad
 [\mathcal D,\mathsf R]=\frac15\sum_{a=0}^4[\mathcal D,g^a].
 \tag{OQC23}
\]
The chain rule applied separately to $f(D^{-a}x)$ and to
$(\mathcal Df)(D^{-a}x)$ proves both identities. Thus all
action comparisons in the invariant lift have explicit original
coefficient maps. The rank estimates in OQC18 concern exactly
the kernel used in the original Gamma and arithmetic quotient
Grams, because those constructions use the same $\Lambda_k$.
An estimate for its volume still requires those original metrics;
no volume conclusion has been inferred solely from its dimension.

% END COMPLETE OQC

\clearpage
% BEGIN COMPLETE OQX
% Exact completion of the orbit rank calculation. No period outcome is assumed.
\section{Exact original rank from the complete quadric orbit}
\label{sec:OQX}

Retain every original object and map in OQC1--23. In particular
$m=1$, $k=4l+1$, $l\geq1$, $q=(k+1)^2$,
$\mathscr P=\mathbb C[x_1,x_2,x_3,x_4]$, and
$g f(x)=f(D^{-1}x)$ for the original Fourier matrix
$D=\operatorname{diag}(\zeta,\zeta^2,\zeta^3,\zeta^4)$.
The polynomial $Q=Q_0(\mathsf H^{-1}x)$ has its actual original
period and amplitude coefficients. Write
$\mathscr I=\mathscr P^{\langle g\rangle}$,
$R=\mathscr P/(Q)$ and $J=\operatorname{im}(\mathscr I\to R)$,
with the grading and the actual observation identification of OQC8.
For negative indices put $\mathscr I_d=0$ and $d_{5,d,0}=0$.

The orbit decision is the exhaustive test OQC10. The size-one
case already has the exact rank OQC11--13. The following calculation
completes the other case using the same original coefficients. It
does not assert that a period parameter belongs to that case before
its coefficient test has been evaluated.

\subsection{Every invariant vanishing polynomial contains the full orbit}

In the size-five case, let $Q_a=g^aQ$ and retain exactly
$\mathcal H_Q=\prod_{a=0}^4Q_a$ from OQC14. Each $Q_a$ is
irreducible, because an invertible linear substitution preserves
factorizations and $Q$ is irreducible by OQC3. The five factors are
pairwise nonassociate by the proved orbit classification, and each
is prime in the original polynomial ring. Consequently
\[
 \boxed{\quad
 \mathscr I\cap Q\mathscr P=\mathcal H_Q\mathscr I,
 \qquad
 \ker(\mathscr I_d\longrightarrow R_d)
                       =\mathcal H_Q\mathscr I_{d-10}.
 \quad}\tag{OQX1}
\]
Here is the complete divisibility proof. If $f\in\mathscr I$ and
$f=Qb$, then for every $a=0,\ldots,4$ the exact equality
$g^af=f$ gives $f=Q_a g^ab$. Thus every one of the five original
prime factors divides $f$. Distinctness of the factors implies
that their product divides $f$: after factoring out $Q_0$,
primality of $Q_1$ and the fact that $Q_1$ does not divide $Q_0$
force $Q_1$ to divide the remaining quotient. Repeating this
argument with $Q_2,Q_3,Q_4$ proves
$f=\mathcal H_Q b_0$. The action permutes the five exact factors
cyclically, so $g\mathcal H_Q=\mathcal H_Q$. Applying $g$ to
this last equality and using $gf=f$ gives
$\mathcal H_Q(gb_0-b_0)=0$. The polynomial ring is an integral
domain and $\mathcal H_Q\ne0$, whence $gb_0=b_0$.
This proves the first inclusion in OQX1. Conversely an invariant
$b_0$ gives an invariant $\mathcal H_Qb_0$ divisible by $Q$,
which proves the reverse inclusion. The product is homogeneous
of degree ten. If its product with $b_0$ is homogeneous of degree
$d$, decomposition of $b_0$ into homogeneous parts shows that
only its part of degree $d-10$ can be nonzero: each other part
has zero product with $\mathcal H_Q$ and hence is zero.
This proves the graded kernel formula, including $d<10$.

\subsection{The exact invariant quotient and the restriction map}

Let $\mathscr U=\mathscr P/(\mathcal H_Q)$ as in OQC16. The
following are isomorphisms of graded complex algebras with their
specified maps:
\[
 \boxed{\quad
 \mathscr I/(\mathcal H_Q\mathscr I)
 \xrightarrow{\ [f]\mapsto[f]_{(\mathcal H_Q)}\ }
       \mathscr U^{\langle g\rangle}
 \xrightarrow{\ \operatorname{res}\ } J\subseteq R.
 \quad}\tag{OQX2}
\]
The action on $\mathscr U$ is defined because $\mathcal H_Q$ is
invariant. If $[f]_{(\mathcal H_Q)}$ is invariant, averaging its
representatives gives
$[\mathsf R f]_{(\mathcal H_Q)}=[f]_{(\mathcal H_Q)}$, since
$[g^af]=[f]$ for every $a$. Thus the first map is onto.
Its kernel before passage to the displayed domain is
$\mathscr I\cap\mathcal H_Q\mathscr P
=\mathcal H_Q\mathscr I$: the reverse inclusion is immediate,
and for $f=\mathcal H_Qb\in\mathscr I$, cancellation in
$\mathcal H_Q(gb-b)=0$ proves $b\in\mathscr I$.
This proves that the first map is injective as well.

The full quotient restriction
$\operatorname{res}:\mathscr U\to R$ is well defined because
$(\mathcal H_Q)\subseteq(Q)$; its kernel is exactly
$Q\mathscr P/\mathcal H_Q\mathscr P$.
For its invariant subspace, choose the invariant representative
just constructed. If its restriction vanishes, that representative
belongs to $\mathscr I\cap Q\mathscr P$, which is
$\mathcal H_Q\mathscr I$ by OQX1. Its class in $\mathscr U$ is
therefore zero. The invariant restriction is injective. Its image
is exactly $J$, since an invariant representative restricts to
$J$ by definition and every element of $J$ has such a representative.
Both maps are induced by polynomial quotient maps, so they preserve
products, sums, complex scalars and degrees. No action of $g$ on
$R$ is asserted in the size-five case. In particular the second
map retains precisely the invariant image of the original
observation, with its complete restriction kernel already computed.

\subsection{All original degrees and the exact kernel dimension}

The map $\Theta_k$ and the original observation in OQC8 identify
$\dim J_k$ with $\operatorname{rank}\Lambda_k$ and identify
$R_k/J_k$ with $(\ker\Lambda_k)^\vee$. Multiplication by the
nonzero polynomial $\mathcal H_Q$ is injective in $\mathscr P$.
Taking dimensions in OQX1 thus gives, for every nonnegative degree
$d$, the exact invariant restriction rank
$\dim J_d=d_{5,d,0}-d_{5,d-10,0}$. At the original degrees this
evaluates to
\[
 \boxed{\quad
 \operatorname{rank}\Lambda_k
      =d_{5,k,0}-d_{5,k-10,0}
      =k^2-6k+17,\qquad
 \dim\ker\Lambda_k=8k-16,
 \quad k=4l+1,\ l\geq1,
 \quad}\tag{OQX3}
\]
in the size-five outcome of the actual coefficient test.
To verify every constant and degree range, for $k\geq10$ the
complete character formula OQC12 gives
\[
 d_{5,k,0}-d_{5,k-10,0}
 =\frac15\left\{\binom{k+3}{3}-\binom{k-7}{3}
                         +4(\epsilon_k-\epsilon_{k-10})\right\}.
\]
The two $\epsilon$ values are equal because their indices differ
by ten. Expanding the remaining cubics gives
\[
 \binom{k+3}{3}-\binom{k-7}{3}
 =\frac{(k^3+6k^2+11k+6)
             -(k^3-24k^2+191k-504)}6
 =5(k^2-6k+17).
\]
The only original degrees below ten are five and nine. Their
negative-degree terms in OQX1 vanish, while the same exact
character formula gives
\[
 d_{5,5,0}=(56+4)/5=12=5^2-6\cdot5+17,
 \qquad
 d_{5,9,0}=220/5=44=9^2-6\cdot9+17.
\]
This proves the first equality of OQX3 for every original $l\geq1$.
Since $\dim E_k=\dim R_k=(k+1)^2$, rank-nullity gives
$(k+1)^2-(k^2-6k+17)=8k-16$, proving the second equality.
For the original degrees already tabulated in OCS26, the exact
five-orbit values are
\[
 \begin{array}{c|r|r}
 k&\operatorname{rank}\Lambda_k&\dim\ker\Lambda_k\\\hline
 5&12&24\\
 9&44&56\\
 13&108&88\\
 17&204&120\\
 21&332&152\\
 25&492&184\\
 29&684&216
 \end{array}
\]
These are exact values for this exhaustively classified orbit
outcome, replacing the less precise bounds OQC18 wherever that
outcome occurs. The full invariant lift and its action defect
OQC15--23 remain valid with all original coefficients. The
size-one outcome retains its different exact values OQC11--13.
No choice between the two outcomes has been made here without
evaluating the original period coefficient test, and no estimate
for a metric volume has been inferred from its dimension alone.

% END COMPLETE OQX

\clearpage
% BEGIN COMPLETE PQO
\section{The actual period quadric and its cyclic continuation}
\label{sec:PQO}

This calculation uses the original simple quartet, $m=1$, and the
literal period matrix, amplitude unit and cyclic matrix of BRU and
OPG. Its formulas evaluate the transported quadric from those data.
All source orders elsewhere in the programme remain unchanged.

\subsection{Original roots, entire unit and centred phase}
Order the roots and keep their actual values as
\[
\begin{gathered}
 \alpha_1=\tfrac12+\delta+i\gamma,\quad
 \alpha_2=\tfrac12+\delta-i\gamma,\quad
 \alpha_3=\tfrac12-\delta+i\gamma,\quad
 \alpha_4=\tfrac12-\delta-i\gamma,\\
 0<\delta<\tfrac12,\quad\gamma>2,\quad
 h(s)=\prod_{j=1}^4(s-\alpha_j),\quad g=2\xi,\quad v_h=g/h.
\end{gathered}\tag{PQO1}
\]
The four zero orders of $g$ are complete and equal to one.
Thus $v_h$ is entire and has nonzero value at each displayed root.
Put $w=s-1/2$, and define the following original constants:
\[
\begin{gathered}
 z_+=\delta+i\gamma,\quad z_-=\delta-i\gamma,\quad
 \lambda=z_+^2,\quad\mu=z_-^2,\\
 A=2(\gamma^2-\delta^2),\qquad B=(\delta^2+\gamma^2)^2,\qquad
 \mathfrak a=\Phi_h(1/2)=\frac1{160}+\frac A{24}+\frac B2,\\
 h(1/2+w)=w^4+Aw^2+B,\qquad
 \Phi_h(1/2+w)=\mathfrak a+\frac{w^5}{5}+\frac{Aw^3}{3}+Bw.
\end{gathered}\tag{PQO2}
\]
The product of the two quadratics
$((w-\delta)^2+\gamma^2)((w+\delta)^2+\gamma^2)$ proves the
quartic identity. Integrating it and using $\Phi_h(0)=0$
proves the phase and constant in PQO2. In particular the translation
has not removed the original constant $\mathfrak a$.

The theta integral IPM.2 has real coefficients, real theta kernel,
and an expression unchanged by $s\mapsto1-s$. It therefore gives
$g(\bar s)=\overline{g(s)}$ and $g(1-s)=g(s)$ directly; the same
two identities hold for $h$. Taking their entire quotient gives
\[
 v_h(\alpha_1)=v_h(\alpha_4)=a,\qquad
 v_h(\alpha_2)=v_h(\alpha_3)=\bar a,\qquad a\ne0.
 \tag{PQO3}
\]
These are the full actual unit values. Their magnitude and phase
are not assigned. Set $t_+=a^{-2}$ and $t_-=\bar a^{-2}$.
For $j=0,1,2,3$ retain
\[
 \sigma_j=t_+\lambda^j-t_-\mu^j,\qquad
 \sigma_2=-A\sigma_1-B\sigma_0,\qquad
 \sigma_3=(A^2-B)\sigma_1+AB\sigma_0.
 \tag{PQO4}
\]
The recurrence follows because $\lambda+\mu=-A$ and
$\lambda\mu=B$. Each $\sigma_j$ is purely imaginary.
Furthermore $\sigma_0$ and $\sigma_1$ cannot both vanish:
the first would give $t_+=t_-\ne0$, and the second would then
give $t_+(\lambda-\mu)=0$, whereas
$\lambda-\mu=4i\delta\gamma\ne0$.

\subsection{The original periods as an absolutely convergent series}
Fix $u\ne0$ on its original logarithm branch. Let
$\ell_j(r)=r\exp(i(\arg u+\pi+2\pi j)/5)$,
$\Gamma_j=-\ell_0+\ell_j$, $1\leq j\leq4$.
Let $\Pi$ integrate the degree-less-than-four representative
against $e^{\Phi_h(s)/u}\,ds$ on these original oriented contours.
Write $F_{ja}=\zeta^{ja}-1$, $\zeta=e^{2\pi i/5}$,
$1\leq j,a\leq4$. Direct finite geometric sums prove
\[
 (F^{-1})_{aj}=\frac15\zeta^{-ja},\qquad
 CF=F D,\qquad D=\operatorname{diag}(\zeta,\zeta^2,\zeta^3,\zeta^4),
 \tag{PQO5}
\]
where $(Cx)_j=x_{j+1}-x_1$ for $j<4$ and $(Cx)_4=-x_1$.
For example multiplying the asserted inverse by $F$ gives
$\frac15\sum_{j=1}^4\zeta^{-ja}(\zeta^{jb}-1)=\delta_{ab}$;
substitution into each row gives $CF=FD$.

Let $T$ be the exact increasing-coefficient translation matrix
$T_{rb}=\binom br2^{-(b-r)}$ for $0\leq r\leq b\leq3$, and
zero otherwise. Thus it sends coefficients of $p(s)$ to those
of $p(1/2+w)$, with determinant one. Define the centred
Fourier-period matrix
\[
 P_{ar}(u)=\frac15\sum_{j=1}^4\zeta^{-ja}
 \int_{\Gamma_j}e^{(w^5/5+Aw^3/3+Bw)/u}w^r\,dw,
 \quad 1\leq a\leq4,\quad0\leq r\leq3.
 \tag{PQO6}
\]
The contours in this formula start at $w=0$. They give precisely
the translated original period by
\[
 F^{-1}\Pi=e^{\mathfrak a/u}P(u)T.
 \tag{PQO7}
\]
To justify this exact contour translation, the substitution
$w=s-1/2$ first gives rays starting at $-1/2$. Join each shifted
ray at a large radius to the corresponding unshifted ray. On the
joining segment their arguments differ by $O(1/r)$, and the
leading real phase is $-r^5/(5|u|)+O(r^4)$; all lower phase
terms and polynomial factors are dominated by this negative
leading term. Those outer joining integrals tend to zero.
The two common-origin joining integrals cancel in
$-\ell_0+\ell_j$ with their opposite orientations. Cauchy's
theorem for the entire integrand proves equality in the limit.
Substitution of the polynomial gives exactly the matrix $T$.

For every integer $L\geq1$ define, on the same branch,
\[
 G_L(u)=u^{L/5}5^{L/5-1}e^{\pi iL/5}\Gamma(L/5).
 \tag{PQO8}
\]
Integration of the leading phase on the two rays gives
$\int_{\Gamma_j}e^{w^5/(5u)}w^{L-1}dw
 =G_L(u)(\zeta^{jL}-1)$ by the real substitution
$r^5/(5|u|)$ on each ray, including its phase in $dw$.
Expanding the remaining two exponentials now gives the exact
entries
\[
\boxed{
 P_{ar}(u)=
 \sum_{\substack{p,q\geq0\\r+3p+q+1\equiv a\pmod5}}
 \frac{A^pB^q}{3^p p!q!u^{p+q}}\,
                 G_{r+3p+q+1}(u).
}\tag{PQO9}
\]
Terms with $r+3p+q+1\equiv0$ have zero two-ray integral
in every column and contribute zero before the Fourier map.
To prove every exchange, the sum of the absolute values of the
unintegrated series on either ray is bounded by a constant times
\[
 R^r\exp\left(-\frac{R^5}{5|u|}
             +\frac{|A|R^3}{3|u|}+\frac{|B|R}{|u|}\right).
\]
Here $R$ is the radial coordinate and $r$ is the fixed column
degree. This integrable majorant
proves absolute termwise integration. It is uniform on compact
coefficient and nonzero-$u$ sets on the chosen branch. The finite
Fourier sum in PQO5 then selects exactly the stated residue class.
Thus PQO9 is an actual period formula, with all coefficients and
all terms retained.

For additional finite computation, let $I_d$ denote the vector
of the four centred ray moments of $w^d$. Integration of the
derivative of $e^{(w^5/5+Aw^3/3+Bw)/u}w^d$, with the same
vanishing outer endpoints and cancelling inner endpoints, proves
\[
 I_{d+4}+A I_{d+2}+B I_d=-ud I_{d-1}\quad(d\geq0),
 \qquad u\partial_B I_d=I_{d+1},\quad
 3u\partial_A I_d=I_{d+3}.
 \tag{PQO10}
\]
The first right side is zero at $d=0$. Coefficient derivatives
are justified by the same compact domination. These are the
original polynomial period recurrences, and they compute every
higher moment from the first four and the actual coefficients.

The determinant is especially explicit in these centred
coordinates:
\[
 \det\Pi=-(2\pi)^2u^2e^{4\mathfrak a/u},\qquad
 \det F=-25\sqrt5,\qquad
 \boxed{\det P(u)=\frac{(2\pi)^2}{25\sqrt5}\,u^2.}
 \tag{PQO11}
\]
For the first equality use the full PDE1--13 determinant
calculation at $n=4$: the phase is $e^{15\pi i}=-1$ and
the critical-value sum is $4\mathfrak a$. This critical sum
also follows directly by pairing $w$ with $-w$ in the odd
centred phase. For $F$, subtract row zero from the other rows
in the five-by-five Fourier Vandermonde and expand its first
column. Its phase is $e^{13\pi i}=-1$ and its magnitude
$5^{5/2}$, by the finite Fourier orthogonality identity.
The last equality follows by taking determinants in PQO7 and
using $\det T=1$. The Gamma grid-product proof PGCC1--6 retains
the full evaluated constant in the first equality. In particular
every inverse below exists on the full stated domain.

\subsection{The actual transformed quadric and every deciding coefficient}
Let $V$ be the centred evaluation matrix with rows
$(1,z_j,z_j^2,z_j^3)$, where
$(z_1,z_2,z_3,z_4)=(z_+,z_-,-z_-,-z_+)$. Its inverse sends
the four primary values to their unique degree-less-than-four
polynomial. Thus the actual matrix from primary coordinates
to the Fourier-period space is
\[
 B_{\rm per}=F^{-1}\Pi UJ_{\rm prim}
       =e^{\mathfrak a/u}P(u)V^{-1}
                         \operatorname{diag}(a,\bar a,\bar a,a).
 \tag{PQO12}
\]
This follows either by applying all maps to the four CRT
idempotents or by evaluating the full unit at each root.
It agrees with OCS and OPG in the given original orientation.

For $Q_0(t)=t_1t_4-t_2t_3$, define
$Q_u(x)=Q_0(B_{\rm per}^{-1}x)$. On the unchanged original
period coordinate $\mathbf y\in\mathbb C^4$ the quadratic is
$Q_u^{\rm per}(\mathbf y)=Q_0((\Pi UJ_{\rm prim})^{-1}\mathbf y)
 =Q_u(F^{-1}\mathbf y)$, because $FB_{\rm per}=\Pi UJ_{\rm prim}$.
Thus the exact isomorphism $\mathbf y=Fx$ carries the following
quadric to the original period space. If
$p(w)=p_0+p_1w+p_2w^2+p_3w^3$, multiplication gives
\[
 p(w)p(-w)=p_0^2+(2p_0p_2-p_1^2)w^2
                     +(p_2^2-2p_1p_3)w^4-p_3^2w^6.
\]
Consequently the exact symmetric matrix in the centred
coefficient frame is
\[
 S=\begin{pmatrix}
 \sigma_0&0&\sigma_1&0\\
 0&-\sigma_1&0&-\sigma_2\\
 \sigma_1&0&\sigma_2&0\\
 0&-\sigma_2&0&-\sigma_3
 \end{pmatrix},\qquad
 \boxed{Q_u(x)=x^{\mathsf T}H(u)x,
 \quad H(u)=e^{-2\mathfrak a/u}P(u)^{-\mathsf T}S P(u)^{-1}.}
 \tag{PQO13}
\]
Indeed applying the preceding product formula at $z_+$ and
$z_-$ with multipliers $t_+$ and $-t_-$ gives exactly $S$.
This retains both actual unit values through every $\sigma_j$.
The two parity blocks have determinants
\[
 \sigma_0\sigma_2-\sigma_1^2
          =-t_+t_-(\lambda-\mu)^2,\qquad
 \sigma_1\sigma_3-\sigma_2^2
          =-t_+t_-B(\lambda-\mu)^2.
\]
Thus $S$ and $H(u)$ are nondegenerate, with the exact determinant
\[
 \det H(u)=e^{-8\mathfrak a/u}
       \frac{B(t_+t_-)^2(\lambda-\mu)^4}{(\det P(u))^2}\ne0.
 \tag{PQO14}
\]

Write $N(u)=\operatorname{adj}P(u)$, with entries $N_{ra}$,
$0\leq r\leq3$, $1\leq a\leq4$, and retain
$\Delta(u)=\det P(u)$ from PQO11. Then every coefficient is
given by the explicit minor formula
\[
\begin{aligned}
 H_{ab}(u)=\frac{e^{-2\mathfrak a/u}}{\Delta(u)^2}
 \{&\sigma_0N_{0a}N_{0b}
 +\sigma_1(N_{0a}N_{2b}+N_{2a}N_{0b}-N_{1a}N_{1b})\\
 &+\sigma_2(N_{2a}N_{2b}-N_{1a}N_{3b}-N_{3a}N_{1b})
                         -\sigma_3N_{3a}N_{3b}\}.
\end{aligned}\tag{PQO15}
\]
This is matrix multiplication in PQO13 using
$P^{-1}=N/\Delta$. Each $N_{ra}$ is the signed original
three-by-three period minor. Its entries are given by the
absolutely convergent series PQO9 or by the original integrals
PQO6. The polynomial coefficient of $x_a^2$ is $H_{aa}$;
that of $x_ax_b$, $a<b$, is $2H_{ab}$. Both factors are
retained in this specified quadratic form.

The exact eight deciding entries are
\[
 \mathcal E(u)=
 (H_{11},H_{22},H_{33},H_{44},H_{12},H_{13},H_{24},H_{34}).
 \tag{PQO16}
\]
The cyclic action on quadrics is the pullback
$Q_u\mapsto Q_u(D^{-1}x)$; its symmetric matrix is
$D^{-\mathsf T}H(u)D^{-1}$. Since $D^5=I$ and $\det D=1$,
the line spanned by a nondegenerate form can be stable only
with character one. In fact if its multiplier is $\eta$, then
$\eta^5=1$ by the group action, while taking determinants
gives $\eta^4=1$ because $\det H(u)\ne0$. Thus $\eta=1$.
Entrywise, invariance is precisely
$(\zeta^{-(a+b)}-1)H_{ab}=0$. The only pairs with
$a+b\equiv0$ are $(1,4)$ and $(2,3)$. Therefore
\[
\boxed{\begin{gathered}
 \text{The quadric is invariant exactly on }
       Z=\{u:\mathcal E(u)=0\};\\
 u\in Z\ \Longrightarrow\
 Q_u(x)=2H_{14}(u)x_1x_4+2H_{23}(u)x_2x_3,
                 \quad H_{14}(u)H_{23}(u)\ne0;\\
 \text{at every }u\notin Z\text{ its orbit consists of five distinct quadrics.}
\end{gathered}}\tag{PQO17}
\]
Every assertion follows from the entry test and the rank-four
determinant. The orbit size divides five; it is one exactly
when the line is stable. Thus no two-member orbit occurs.
Moreover $F^{-1}C^{-1}=D^{-1}F^{-1}$ by PQO5, so this is
exactly the orbit of $Q_u^{\rm per}(C^{-j}\mathbf y)$ in the
original period coordinates, with the same five indexed maps.
PQO15--17 determine the exact exceptional locus by original
period minors and original amplitude values, without assigning
generic values to them.

\subsection{The cyclic action is the actual original continuation in \texorpdfstring{$u$}{u}}
Continue the original logarithm of $u$ by $2\pi i$. Its rays
advance from $\ell_j$ to $\ell_{j+1}$, and their oriented
differences advance by the exact original matrix $C$.
For the last row $\ell_5=\ell_0$, giving $-\Gamma_1$.
Equivalently every summand of PQO9 acquires the phase
\[
 e^{2\pi i((r+3p+q+1)/5-p-q)}
       =\zeta^{r+3p+q+1}=\zeta^a.
\]
The original constant $e^{\mathfrak a/u}$ is single-valued
under this continuation. Hence
\[
\begin{gathered}
 P(\log u+2\pi i)=DP(\log u),\quad
 B_{\rm per}(\log u+2\pi i)=DB_{\rm per}(\log u),\\
 H(\log u+2\pi i)=D^{-\mathsf T}H(\log u)D^{-1}.
\end{gathered}\tag{PQO18}
\]
The contour argument and the convergent series prove the same
map. Thus PQO17 describes the actual period continuation of
this original quadric. It introduces no freely chosen cyclic
action on its coefficients.

\subsection{The large-parameter geometry and a positive-ray obstruction}
Let $\varepsilon=u^{-1/5}$ on the same branch, and set
$D_\varepsilon=\operatorname{diag}(\varepsilon^{-1},
\varepsilon^{-2},\varepsilon^{-3},\varepsilon^{-4})$.
The exact series PQO9 gives
\[
 P(u)=\mathscr P(A\varepsilon^2,B\varepsilon^4)D_\varepsilon,
 \qquad
 \mathscr P(0,0)=\operatorname{diag}(g_1,g_2,g_3,g_4),
 \quad g_a=5^{a/5-1}e^{\pi ia/5}\Gamma(a/5).
 \tag{PQO19}
\]
The entries of $\mathscr P$ are entire in the two displayed
arguments, by the same absolute integrable series majorant.
The same determinant proof PQO11 applies to every complex pair
of lower odd-phase coefficients, since the centred phase remains
odd and its critical values pair to zero with full multiplicities.
Thus their determinant is the nonzero constant
$(2\pi)^2/(25\sqrt5)$, by PQO11 and $\det D_\varepsilon=u^2$.
Their inverse is therefore holomorphic everywhere, given by its
adjugate divided by that precise constant. Near $\varepsilon=0$
it is $\operatorname{diag}(g_a^{-1})+O(\varepsilon^2)$.
This bound follows from its convergent power series after
substituting $(A\varepsilon^2,B\varepsilon^4)$; it holds on
a full small complex disc, with the fixed original $A,B$.

If $\sigma_0\ne0$, PQO13 and PQO19 give
\[
 e^{2\mathfrak a/u}H_{11}(u)
              =\frac{\sigma_0}{g_1^2}\varepsilon^2
                                     +O(\varepsilon^4).
\]
If $\sigma_0=0$, then $\sigma_1\ne0$ by PQO4, and they give
\[
 e^{2\mathfrak a/u}H_{22}(u)
              =-\frac{\sigma_1}{g_2^2}\varepsilon^4
                                     +O(\varepsilon^6).
 \tag{PQO20}
\]
For verification of the orders, insert
$P^{-1}=D_\varepsilon^{-1}\mathscr P^{-1}$ into PQO13.
The middle matrix has entries $S_{rs}\varepsilon^{r+s+2}$.
Its first possible degree is two, from $S_{00}$, and its
next possible degree is four, from $S_{02},S_{20},S_{11}$.
The constant term of $\mathscr P^{-1}$ is diagonal, so the
indicated leading entry is exactly the displayed nonzero
coefficient. Every further contribution has the stated higher
degree. This also proves the expansion when some other
$\sigma_j$ vanish.

Consequently for every fixed original quartet and its full
actual amplitude there exists a finite radius outside which
the actual orbit has five distinct quadrics. This conclusion
holds on every original logarithm branch: the estimate is on
the full complex $\varepsilon$ disc, and its leading coefficient
is nonzero. The set $Z$ of PQO17 is a discrete subset of the
logarithmic $u$ surface. Indeed the decisive entry just proved
is holomorphic on that connected surface and is not identically
zero; its zeros are isolated by its local power series, and
$Z$ is contained in that zero set. The explicit finite-radius
bound from the companion period-connection calculation
strengthens this argument while retaining the same actual
coefficients.

Moreover PQO18 multiplies each deciding entry by its nonzero
root-of-unity factor, so the condition $\mathcal E(u)=0$ is
unchanged by every full logarithm turn. Therefore $Z$ is the
inverse image under $\log u\mapsto u$ of a well-defined set
$Z_{\rm abs}\subset\mathbb C^*$. On a sufficiently small disc
about any nonzero $u$, choose one logarithm branch. The preceding
isolated-zero proof shows that $Z_{\rm abs}$ is discrete and
closed on that disc. The explicit radius in the companion
calculation gives $|u|<R_*$ throughout $Z_{\rm abs}$.
Consequently its only possible accumulation point in the finite
$u$-plane is zero. In particular each compact annulus about
zero contains only finitely many such points: an infinite
subset of that compact set would have a nonzero accumulation
point, contradicting the proved local isolated-zero property.

For positive real $u$ retain the actual lift
$\log u=\ln u+2\pi i\ell_{\rm br}$, $\ell_{\rm br}\in\mathbb Z$.
The exact series has an additional real structure. Every
surviving index in PQO9 is $a+5L$, and its phase at the lift
$\ell_{\rm br}=0$ is $e^{\pi ia/5}(-1)^L$. Its remaining
coefficients are real. The actual lift multiplies row $a$ by
$\zeta^{\ell_{\rm br}a}$, by PQO18. Absolute convergence
therefore proves
\[
 P(u)=E_{\ell_{\rm br}} R(u),\qquad
 E_{\ell_{\rm br}}=D^{\ell_{\rm br}}
 \operatorname{diag}(e^{\pi i/5},e^{2\pi i/5},
                       e^{3\pi i/5},e^{4\pi i/5}),
 \qquad R(u)\text{ real and invertible}.
 \tag{PQO21}
\]
Let $M$ be multiplication by $w$ in the increasing centred
polynomial frame, so
\[
 M=\begin{pmatrix}0&0&0&-B\\1&0&0&0\\0&1&0&-A\\0&0&1&0\end{pmatrix},
 \quad c_3=\frac{t_+-t_-}{4i\delta\gamma},\quad
 c_1=\frac{t_++t_-}{2}-(\delta^2-\gamma^2)c_3.
 \tag{PQO22}
\]
Both $c_1,c_3$ are real by PQO3, and the polynomial
$c_1+c_3 w^2$ takes values $t_+,t_-,t_-,t_+$ at the four
original roots. Hence $c_1M+c_3M^3$ is precisely multiplication
by $w(j_hv_h)^{-2}$, with its entire unit retained.

The companion period-connection proof establishes the exact
identity between cyclic invariance of $Q_u$ and
\[
 [P(u)(c_1M+c_3M^3)P(u)^{-1},D]=0.
 \tag{PQO23}
\]
It proves this by the flat residue/reflection form on the
original polynomial periods and retains all its constants.
Using that proved identity, PQO21 gives a further precise
obstruction on the positive ray. Since $D$ has distinct
diagonal entries, a matrix commuting with it is diagonal.
The matrix in PQO23 is $E_{\ell_{\rm br}}$ times a real matrix
times $E_{\ell_{\rm br}}^{-1}$;
if diagonal, its diagonal entries must therefore be real.
Its eigenvalues are exactly
\[
 \frac{\delta+i\gamma}{a^2},\quad
 \frac{\delta-i\gamma}{\bar a^2},\quad
 -\frac{\delta-i\gamma}{\bar a^2},\quad
 -\frac{\delta+i\gamma}{a^2},
 \tag{PQO24}
\]
as follows by evaluation at the four distinct roots of $h$.
Thus every positive-real-$u$ point in $Z$ satisfies the explicit
original-unit equation
\[
 \operatorname{Im}\left(\frac{\delta+i\gamma}{a^2}\right)=0.
 \tag{PQO25}
\]
If that imaginary part is nonzero for the original unit, the
positive ray contains no point of $Z$. When it is zero, the
four displayed eigenvalues are the two nonzero real numbers
$r,-r$, each twice, and the complete period-minor tests
PQO15--17 still determine $Z$. The equation is not assumed
to hold or fail. Together the exact minor tests, the actual
period continuation, the proved large-parameter result and
this real-structure obstruction give control of the original
quadric family with every amplitude value and phase present.

% END COMPLETE PQO

\clearpage
% BEGIN COMPLETE RPC
% Independent mathematical provider. Original periods, amplitude and branch retained.
\section{The original period quadric as a residue commutator}
\label{sec:actual-period-residue-commutator}

Retain $0<\delta<1/2$, $\gamma>2$, the ordered simple quartet
$\alpha_1=1/2+\delta+i\gamma$, $\alpha_2=1/2+\delta-i\gamma$,
$\alpha_3=1/2-\delta+i\gamma$, $\alpha_4=1/2-\delta-i\gamma$,
and $h(s)=\prod_a(s-\alpha_a)$. Put
\[
 w=s-\tfrac12,\quad r_a=\alpha_a-\tfrac12,\quad
 \beta=2(\gamma^2-\delta^2),\quad \eta=(\delta^2+\gamma^2)^2,
 \quad H(w)=w^4+\beta w^2+\eta.
 \tag{RPC1}
\]
Thus $r_4=-r_1$, $r_3=-r_2$, $r_2=\overline{r_1}$ and
$r_1^2-r_2^2=4i\delta\gamma$. No root or original coordinate is
discarded by this translation. Its coefficient isomorphism is
$J_c:\mathbb C[w]_{<4}\to E_h$, $J_cp=[p(s-1/2)]$; its inverse
is substitution $s=w+1/2$ in the unique original representative.
The primitive with its original constant is
\[
 \Phi_h(s)=A+\phi(w),\quad
 \phi(w)=\frac{w^5}{5}+\frac{\beta w^3}{3}+\eta w,\quad
 A=\Phi_h(1/2)=\frac1{160}+\frac\beta{24}+\frac\eta2.
 \tag{RPC2}
\]
The last equality is forced by $\Phi_h(0)=0$ and follows by evaluating
the displayed odd polynomial at $w=-1/2$.

Let $u\ne0$ carry the specified original branch of its logarithm.
Let $\Pi$ be the original period map on the rays and orientations
of PDE2, and let $F_{ja}=\zeta^{ja}-1$, $\zeta=e^{2\pi i/5}$,
$1\leq j,a\leq4$. Its inverse has entries
$(F^{-1})_{aj}=\zeta^{-ja}/5$: multiplying the two matrices and using
$\sum_{j=1}^4\zeta^{jb}=4$ for $5\mid b$ and $-1$ otherwise proves
this entry by entry. The literal centered coefficient-to-Fourier map is
\[
 P=F^{-1}\Pi J_c,\qquad P_0=e^{-A/u}P.
 \tag{RPC3}
\]
These are invertible by the complete original period determinant
PDE3--13. The contours for $P_0$ may be taken to be the rays based at
$w=0$ with the original five angles. Indeed translating the old rays
gives rays based at $-1/2$. Join corresponding finite endpoints and
close each contour at radius $R$. The integrand is entire. On the
bounded-length outer joining curves its exponent is
$-R^5/(5|u|)+O(R^4)$, uniformly along the joining curves, so the
outer integrals tend to zero. The common finite joining path occurs
with opposite coefficients in $-\ell_0+\ell_j$ and cancels exactly.
Cauchy's integral theorem therefore proves this contour identity,
including its signs and the factor $e^{A/u}$.

\subsection{A flat alternating pairing with all constants}

On $\mathbb C[w]/(H)$ define the bilinear, rather than sesquilinear,
pairing
\[
 \mathcal J(p,q)=\sum_{H(r)=0}\frac{p(r)q(-r)}{H'(r)}.
 \tag{RPC4}
\]
The roots are simple at the original packet. The equivalent polynomial
formula is the coefficient of $w^3$ in the remainder of $p(w)q(-w)$
modulo $H$. To prove it, use Lagrange interpolation: the coefficient
of $w^3$ in the idempotent at $r$ is $1/H'(r)$. This polynomial
formula extends the pairing to every complex value of $\beta,\eta$,
including repeated roots, without a division by a repeated-root value.
In the increasing frame $1,w,w^2,w^3$, multiplication by $w$ and
this pairing are the exact matrices
\[
 M=\begin{pmatrix}0&0&0&-\eta\\1&0&0&0\\0&1&0&-\beta\\0&0&1&0\end{pmatrix},
 \qquad
 J=\begin{pmatrix}0&0&0&-1\\0&0&1&0\\0&-1&0&\beta\\1&0&-\beta&0\end{pmatrix}.
 \tag{RPC5}
\]
For the last matrix, the coefficient functional is zero through degree
two, one in degree three, zero in degree four and six, and $-\beta$
in degree five. Inserting the sign $(-1)^j$ from the second argument
gives every displayed entry. Thus $J^{\mathsf T}=-J$, $\det J=1$,
and direct multiplication gives $M^{\mathsf T}J+JM=0$.

Allow the two original odd-phase coefficients to vary complexly,
solely to calculate the existing period map. The ray domination
$e^{-r^5/(5|u|)+O(r^3+r)}$ on compact parameter sets justifies all
coefficient derivatives and integrations by parts below. Ordinary
division of $w^{j+3}$ by $H$, followed by differential reduction,
gives
\[
 \partial_\eta P_0=\frac1uP_0M,\qquad
 \partial_\beta P_0=\frac1{3u}P_0(M^3-uN),\qquad
 N=\begin{pmatrix}0&0&1&0\\0&0&0&2\\0&0&0&0\\0&0&0&0\end{pmatrix}.
 \tag{RPC6}
\]
For clarity, the quotient derivatives for columns $j=0,1,2,3$ are
$0,0,1,2w$, respectively; hence this is the full correction matrix.
The two outer boundary terms vanish, and the two values at the common
origin cancel. No differential remainder was replaced by an arithmetic
remainder in RPC6.

There is also a direct nonvanishing proof on this entire two-parameter
family. The displayed matrices satisfy $\operatorname{Tr}M=0$,
$\operatorname{Tr}M^3=0$ and $\operatorname{Tr}N=0$. The first and
third follow from their diagonals; the second follows by multiplying
RPC5 three times, or from the vanishing third Newton sum of the even
monic polynomial $H$. Multilinearity of the determinant in its columns
then gives $\partial_\eta\det P_0=\partial_\beta\det P_0=0$ from
RPC6, without first assuming an inverse exists. At the origin the
literal Gamma ray integrals give the diagonal matrix $G_u$ defined in
RPC8, whose determinant is nonzero. Hence $P_0$ is invertible at every
point of $\mathbb C^2$, justifying the inverses throughout this family.

The matrix identity $M^{\mathsf T}J+JM=0$ also gives
$(M^3)^{\mathsf T}J+JM^3=0$. Direct multiplication of RPC5--6 gives
\[
 \partial_\beta J+\frac13(N^{\mathsf T}J+JN)=0.
 \tag{RPC7}
\]
In particular $\partial_\beta J$ has entries $(2,3)=1$, $(3,2)=-1$
when indices begin at zero, while $N^{\mathsf T}J+JN$ has the
opposite threefold entries. Differentiate the inverse matrices in
$P_0^{-\mathsf T}JP_0^{-1}$. RPC6--7 show that both its $\beta$ and
its $\eta$ derivatives vanish. This matrix is consequently constant
on the connected coefficient space $\mathbb C^2$.

Define the unchanged positive-Gamma ray constants
\[
 g_a=5^{a/5-1}e^{\pi ia/5}\Gamma(a/5),\qquad
 G=\operatorname{diag}(g_1,g_2,g_3,g_4),\quad
 G_u=\operatorname{diag}(u^{a/5}g_a)_{a=1}^4.
 \tag{RPC8}
\]
At $\beta=\eta=0$, literal substitution in each original ray integral
gives $P_0=G_u$. Every $g_a$ is nonzero, since its Gamma factor is
the integral of a strictly positive function. We have therefore proved
the complete flat-pairing identity
\[
 \begin{split}
 \Omega_0&=P_0^{-\mathsf T}JP_0^{-1}
       =G_u^{-\mathsf T}J\big|_{\beta=0}G_u^{-1},\\
 (\Omega_0)_{14}&=-\frac1{u g_1g_4},\qquad
 (\Omega_0)_{23}=\frac1{u g_2g_3},\qquad
 (\Omega_0)_{ba}=-(\Omega_0)_{ab},\\
 \Omega&=P^{-\mathsf T}JP^{-1}=e^{-2A/u}\Omega_0.
 \end{split}
 \tag{RPC9}
\]
All entries not listed are zero; the indices in RPC9 begin at one.
This formulation retains the Gamma constants directly and requires
no unevaluated normalization of a period or source mass. For
$D=\operatorname{diag}(\zeta,\zeta^2,\zeta^3,\zeta^4)$, the two
paired sums of weights are five. It follows entry by entry that
\[
 D^{\mathsf T}\Omega D=\Omega,
 \qquad D^{\mathsf T}\Omega_0D=\Omega_0.
 \tag{RPC10}
\]

\subsection{The full arithmetic unit and the exact quadratic form}

Retain the original $v_h=2\xi/h$ and its complete simple-root unit
values. Reflection and conjugation of this original function give
\[
 v_h(\alpha_1)=v_h(\alpha_4)=a,\qquad
 v_h(\alpha_2)=v_h(\alpha_3)=\overline a,\qquad a\ne0.
 \tag{RPC11}
\]
Indeed both numerator and denominator are invariant under $s\mapsto1-s$
and conjugate under $s\mapsto\overline s$, first away from the roots
and then by the removable values. The complete zero orders give
nonzero values at these simple roots. Write
\[
 c_3=\frac{a^{-2}-\overline a^{-2}}{4i\delta\gamma},\qquad
 c_1=\frac{a^{-2}+\overline a^{-2}}2-(\delta^2-\gamma^2)c_3,
 \quad \mathfrak r(w)=c_1w+c_3w^3.
 \tag{RPC12}
\]
Both coefficients are real, and $(c_1,c_3)\ne(0,0)$. Directly,
$c_1+c_3r_1^2=a^{-2}$ and $c_1+c_3r_2^2=\overline a^{-2}$;
the same equalities hold at the respective opposite roots. Thus
$\mathfrak r(M)=M\,v_h(M+1/2)^{-2}$ in the original finite algebra, with the
entire function evaluated by its complete original arithmetic unit.
RPC12 is an exact interpolation of this specific multiplier.

For $Q_0(t)=t_1t_4-t_2t_3$, let $J_{\rm prim}$ be the original
Lagrange coordinate isomorphism and set
\[
 B=F^{-1}\Pi UJ_{\rm prim},\qquad Q(x)=Q_0(B^{-1}x),\qquad
 \mathcal H=P\mathfrak r(M)P^{-1}.
 \tag{RPC13}
\]
If $p=P^{-1}x$, the primary components of $B^{-1}x$ are exactly
$p(r_1)/a,p(r_2)/\overline a,p(r_3)/\overline a,p(r_4)/a$.
Since $H'(r)=4r(r^2+\gamma^2-\delta^2)$, its values give
$r/H'(r)=1/(8i\delta\gamma)$ on roots one and four, and
$r/H'(r)=-1/(8i\delta\gamma)$ on roots two and three. Inserting
these four values into RPC4 proves the exact identity
\[
 \boxed{\quad Q(x)=4i\delta\gamma\,
       p^{\mathsf T}\mathfrak r(M)^{\mathsf T}Jp
       =4i\delta\gamma\,x^{\mathsf T}\mathcal H^{\mathsf T}\Omega x.\quad}
 \tag{RPC14}
\]
There is no conjugate transpose in this quadratic identity. Because
$\mathfrak r$ is odd, $\mathfrak r(M)^{\mathsf T}J+J\mathfrak r(M)=0$; hence the matrices in
RPC14 are symmetric. This also proves
$\mathcal H^{\mathsf T}\Omega+\Omega\mathcal H=0$.
The four eigenvalues of $\mathfrak r(M)$ are exactly
$r_1/a^2,\overline{r_1/a^2},-\overline{r_1/a^2},-r_1/a^2$.
They are nonzero, although they need not be distinct. Both $B$ and
$Q_0$ are invertible as linear and quadratic matrices, respectively,
so $Q$ always has rank four, including the possible eigenvalue collisions.

By RPC10, substituting $Dx$ into RPC14 transports $\mathcal H$ to
$D^{-1}\mathcal H D$. The resulting symmetric quadratic matrices
are equal if and only if the corresponding $\mathcal H$ matrices
are equal, because $\Omega$ is invertible. Consequently
\[
 \boxed{\quad Q(Dx)=Q(x)\quad\Longleftrightarrow\quad
 [\mathcal H,D]=0\quad\Longleftrightarrow\quad
 \mathcal H_{ab}=0\text{ for every }a\ne b.\quad}
 \tag{RPC15}
\]
This is an exact test in the actual original period matrix and unit.
The final equivalence uses the four distinct diagonal entries of $D$.
It makes no assumption that their off-diagonal tests are independent.

For completeness, a rank-four quadratic semi-invariant of cyclic
weight $j\ne0$ is impossible: its coefficient in row $j$ could pair
only with the missing weight zero, so that row is zero. Thus any
semi-invariant rank-four quadratic must have weight zero. Since five
is prime, the projective orbit of $Q$ under $D$ has size one or five,
and size one is exactly RPC15. Every failure of any displayed
off-diagonal test proves size five.

\subsection{An explicit original-parameter radius forcing orbit five}

Here we prove that RPC15 fails for all sufficiently large $|u|$,
with a completely specified finite radius for the given original
quartet and unit. The argument evaluates the original periods; it
does not choose or insert a replacement amplitude.

Set $b=\beta u^{-2/5}$, $c=\eta u^{-4/5}$, and let $T(b,c)$
be the Fourier period matrix for
$z^5/5+bz^3/3+cz$ at $u=1$, in columns $1,z,z^2,z^3$ and on
the five rays of angles $(\pi+2\pi j)/5$. Substitution
$w=u^{1/5}z$, with exactly the original branch, gives
\[
 P_0=T(b,c)S_u,\quad S_u=\operatorname{diag}(u^{a/5})_{a=1}^4,
 \quad S_uMS_u^{-1}=u^{1/5}M(b,c),\quad T(0,0)=G.
 \tag{RPC16}
\]
Every column includes its original differential factor. The matrix
$M(b,c)$ is RPC5 with $\beta,\eta$ replaced by $b,c$; in particular
$N_0:=M(0,0)$ has its three subdiagonal entries equal to one.

The following positive constants are explicit finite real integrals
or finite expressions in the original Gamma constants:
\[
 \begin{split}
 E_j&=\frac85\int_0^\infty r^j(r^3/3+r)
                    e^{-r^5/5+r^3/3+r}\,dr\quad(0\le j\le3),\\
 C_T&=2\left(\sum_{j=0}^3E_j^2\right)^{1/2},\qquad
 g_+=\max_a|g_a|,\quad g_-^{-1}=\max_a|g_a|^{-1},\\
 \kappa&=\max_{1\le a<4}|g_{a+1}/g_a|,\qquad
 K_0=GN_0G^{-1},\\
 C_K&=4C_Tg_-^{-1}+2g_+g_-^{-1}
                       +2g_+g_-^{-2}C_T,\\
 C_{K,3}&=C_K(3\kappa^2+3\kappa C_K+C_K^2),\qquad
 \epsilon_0=\min\{1,(2C_Tg_-^{-1})^{-1}\}.
 \end{split}
 \tag{RPC17}
\]
Here $g_-^{-2}$ means $(\max_a|g_a|^{-1})^2$. The negative fifth
degree exponent proves convergence of every integral. All norms in
the following estimates are the Euclidean operator norm, with the
Frobenius norm used explicitly as an upper bound where stated.

Write $\epsilon=|b|+|c|$. When $|b|,|c|\le1$, on each ray
\[
 |e^{bz^3/3+cz}-1|
 \le \epsilon(r^3/3+r)e^{r^3/3+r}.
\]
Each contour has two rays, and each entry of $F^{-1}$ has absolute
value $1/5$ with four summands. The entry in column $j$ of
$T-G$ is therefore at most $E_j\epsilon$. Summing the squares
over all four rows proves
\[
 \|T-G\|\le\|T-G\|_{\rm F}\le C_T\epsilon.
 \tag{RPC18}
\]
For $\epsilon\le\epsilon_0$, the convergent inverse series for
$I+G^{-1}(T-G)$ gives
$\|T^{-1}\|\le2g_-^{-1}$ and
$\|T^{-1}-G^{-1}\|\le2g_-^{-2}C_T\epsilon$.
Also $\|M(b,c)-N_0\|\le\epsilon$, $\|M(b,c)\|\le2$,
and $\|N_0\|=1$. Subtract the products
$K:=T M(b,c)T^{-1}$ and $K_0$ in three terms, changing the left
factor, then the middle factor, then the inverse factor. The preceding
inequalities give, with the complete RPC17 constants,
\[
 \|K-K_0\|\le C_K\epsilon,\quad
 \|K\|\le\kappa+C_K\epsilon,\quad
 \|K^3-K_0^3\|\le C_{K,3}\epsilon.
 \tag{RPC19}
\]
For the last assertion use the exact noncommutative identity
$K^3-K_0^3=(K-K_0)K^2+K_0(K-K_0)K+K_0^2(K-K_0)$.
Its norm is at most $C_K\epsilon$ times
$(\kappa+C_K)^2+\kappa(\kappa+C_K)+\kappa^2$, as claimed.

The original multiplication and full unit in RPC13 now give exactly
\[
 \mathcal H=c_1u^{1/5}K+c_3u^{3/5}K^3,
 \qquad (K_0^3)_{41}=g_4/g_1,\quad (K_0)_{21}=g_2/g_1.
 \tag{RPC20}
\]
The scalar $e^{A/u}$ cancels from this conjugation, while it remains
in the quadratic form through RPC9. Put $B_*=\beta+\eta>0$.
For $|u|\ge1$, $\epsilon\le B_*|u|^{-2/5}$.

If $c_3\ne0$, define the completely specified radius
\[
 \begin{split}
 T_*&=\max\left\{1,\frac{B_*}{\epsilon_0},
 \frac{2\bigl(|c_3|C_{K,3}B_*+|c_1|(\kappa+C_K)\bigr)}
 {|c_3g_4/g_1|}\right\},\qquad R_*=T_*^{5/2}.
 \end{split}
 \tag{RPC21}
\]
For every $u$ on its original branch with $|u|\ge R_*$,
RPC19--20 imply
\[
 \left|u^{-3/5}\mathcal H_{41}-c_3g_4/g_1\right|
 \le |u|^{-2/5}\bigl(|c_3|C_{K,3}B_*
                           +|c_1|(\kappa+C_K)\bigr)
 \le\tfrac12|c_3g_4/g_1|.
 \tag{RPC22}
\]
In particular $\mathcal H_{41}\ne0$. In the quadratic matrix
RPC14 this is the nonzero $(1,1)$ entry
$4i\delta\gamma\,\mathcal H_{41}\Omega_{41}$, whose cyclic
weight is two.

If $c_3=0$, then $c_1\ne0$. In this case define
\[
 T_* =\max\left\{1,\frac{B_*}{\epsilon_0},
                   \frac{2C_KB_*}{|g_2/g_1|}\right\},
 \qquad R_*=T_*^{5/2}.
 \tag{RPC23}
\]
For every $|u|\ge R_*$, the same estimates give
\[
 \left|u^{-1/5}\mathcal H_{21}-c_1g_2/g_1\right|
       \le |c_1|C_KB_*|u|^{-2/5}
       \le\tfrac12|c_1g_2/g_1|.
 \tag{RPC24}
\]
Thus $\mathcal H_{21}\ne0$. The $(1,3)$ entry of the quadratic
matrix is then $4i\delta\gamma\mathcal H_{21}\Omega_{23}$,
which has cyclic weight four; symmetry retains its $(3,1)$ partner.
Both cases prove the unconditional conclusion for the actual original
parameters and the explicit radius belonging to those parameters:
\[
 \boxed{\quad |u|\ge R_*\quad\Longrightarrow\quad
       \#\{[Q(D^jx)]:0\le j<5\}=5.\quad}
 \tag{RPC25}
\]
The bracket denotes the projective class of the quadratic polynomial.
RPC21 or RPC23 supplies its radius for every nonzero original unit;
neither case assumes a condition on an uncomputed period coefficient.

Finally, the original periods are holomorphic on the connected
logarithm cover of $\mathbb C^*$, and on each branch domain: locally the contours may
be held in common decay sectors and differentiated under the dominated
integral, and contour deformation identifies these local definitions.
Their determinant never vanishes. Hence all entries of $\mathcal H$
are holomorphic there. RPC22 or RPC24 proves that its respective
decisive entry is not identically zero on the logarithm cover and,
by the identity theorem on that connected cover, on any open branch domain. The zeros
of that one entry are isolated, by its convergent local power series.
The exact exceptional set
\[
 \mathcal E=\{u:\mathcal H_{ab}(u)=0\text{ for all }a\ne b\}
 \tag{RPC26}
\]
is a subset of those isolated zeros and satisfies $|u|<R_*$.
Thus it has no accumulation inside the branch domain; accumulation
at $u=0$ or a branch boundary has not been excluded. RPC15 is the
full calculation specifying this set. No assertion that it is empty
at arbitrary finite complex $u$ is made or used.

% END COMPLETE RPC

\clearpage
% BEGIN COMPLETE OKV
% Complete finite estimate on the original m=1 packet and its actual kernel.
% Original source and quotient providers: IFB1--11 and GMB1--16.
\section{A finite bound for the original kernel volume from its actual dimension}
\label{sec:OKV}

\paragraph{The original packet, sources, and exact coordinate map.}
Throughout this calculation the original multiplicity is $m=1$ and
\[
 k=4l+1,\quad l\ge1,\quad q=(k+1)^2\ge36,\quad c=k/2,
 \qquad 0<\delta<1/2,\quad\gamma>2.
 \tag{OKV1}
\]
Retain every original root and its two coordinates:
\[
\begin{aligned}
 \chi(S)&=\prod_{a,b=0}^k
  [S-c-(2a-k)\delta-i(2b-k)\gamma],\\
 \omega_{ab}&=(2b-k)\gamma-i(2a-k)\delta,\\
 P(y)&=i^{-q}\chi(c+iy)=\prod_{a,b=0}^k(y-\omega_{ab}).
\end{aligned}
 \tag{OKV2}
\]
The identity follows factor by factor from
$iy-(2a-k)\delta-i(2b-k)\gamma=i(y-\omega_{ab})$;
it retains the full phase $i^q$. The algebra map
\[
 \Phi:\mathbb C[y]/(P)\longrightarrow E_\chi:=\mathbb C[S]/(\chi),
 \qquad [r(y)]\longmapsto[r((S-c)/i)]
 \tag{OKV3}
\]
is well defined because $P((S-c)/i)=i^{-q}\chi(S)$.
Its inverse is $[f(S)]\mapsto[f(c+iy)]$, as verified by substitution.
Both maps preserve degrees and are inverse on each polynomial ring
before taking the quotients. In particular division by $P$ is the
original remainder operation through this exact coordinate map.

Use only the two original source orders $s\in\{1,k\}$ and write
\[
 b_s=s/2,\qquad M_s=(2\pi)^{s/2},\qquad
 dm_{0,s}(y)=\frac{M_s2^{b_s-2}}{\pi\Gamma(b_s)}
     \left|\Gamma\left(\frac{b_s}{2}+\frac{iy}{2}\right)\right|^2dy.
 \tag{OKV4}
\]
The exact source calculation IFB2--5 supplies mass $M_s$, the transform
$\int e^{ty}dm_{0,s}=M_s(\cos t)^{-b_s}$ for $|t|<\pi/2$, and
the original monic orthogonal polynomials and norms
\[
\begin{gathered}
 p_0=1,\quad p_1=y,\qquad
 p_{d+1}=yp_d-d(b_s+d-1)p_{d-1},\\
 h_d^{(s)}=\int|p_d|^2dm_{0,s}=M_s d!(b_s)_d,\qquad
 \varphi_d(y)=p_d(y)/\sqrt{h_d^{(s)}},\qquad
 u_d(S)=\varphi_d((S-c)/i).
\end{gathered}
 \tag{OKV5}
\]
All these integrals are on the same real line, with the original
source mass. For example the stated orthogonality and norms also
follow directly from its transform and the exact generating function
\[
 \sum_{d\ge0}p_d(y)\frac{z^d}{d!}
       =(1+z^2)^{-b_s/2}e^{y\arctan z}.
\]
Indeed the source integral of the product at sufficiently small real
$z,w$ is $M_s(1-zw)^{-b_s}$. The cosine addition identity and the
displayed transform give this equality; the source exponential bound
justifies each finite coefficient derivative under the integral.
The coefficients of $z^dw^e$ then give zero for $d\ne e$ and
$M_s d!(b_s)_d$ for $d=e$. The initial coefficients are $1,y$,
and differentiating the generating function gives precisely the
recurrence in OKV5. Thus the source norms used below retain their
original constants and all finite-degree factors.

Let $C_s:\mathbb C^q\to E_\chi$ have columns $[u_0],\ldots,[u_{q-1}]$
in the original frame $(1,S,\ldots,S^{q-1})$. Its exact leading
coefficients give
\[
 \det C_s=i^{-q(q-1)/2}\prod_{d=0}^{q-1}(h_d^{(s)})^{-1/2}\ne0.
\]
Retain the original GMB vectors and matrices, including the last row:
\[
 v_j=C_s^{-1}[u_{q+j}],\qquad
 K_j=I_q+\sum_{a=0}^{j-1}v_av_a^*,\qquad 0\le j\le q+1.
 \tag{OKV6}
\]
The vector $v_q$ has original polynomial degree $2q$.

\paragraph{Exact remainder coefficients in all original roots.}
Index the full list in OKV2 by $\omega_1,\ldots,\omega_q$, with
each pair $(a,b)$ occurring exactly once. Define the elementary and
complete homogeneous polynomials by their finite sums
\[
 e_h(\omega)=\sum_{1\le a_1<\cdots<a_h\le q}
                      \omega_{a_1}\cdots\omega_{a_h},\qquad
 h_t(\omega)=\sum_{n_1+\cdots+n_q=t}\prod_{a=1}^q\omega_a^{n_a},
 \tag{OKV7}
\]
with $e_0=h_0=1$, $e_h=0$ outside $0\le h\le q$, and
nonnegative integer $n_a$ in the second sum. The notation $h_t(\omega)$
in this paragraph is a polynomial in the roots; the original squared
norm continues to be $h_d^{(s)}$ of OKV5. Multiplication of formal
power series gives the exact identity
\[
 \left(\sum_{h=0}^q(-1)^he_h(\omega)z^h\right)
 \left(\sum_{t\ge0}h_t(\omega)z^t\right)=1,
\]
because each geometric series $(1-\omega_a z)^{-1}$ multiplies its
own factor $1-\omega_a z$. Consequently, for $D=q+j$ with $0\le j\le q$,
the quotient and remainder of $y^D$ upon division by $P$ are exactly
\[
\begin{aligned}
 Q_D(y)&=\sum_{t=0}^jh_t(\omega)y^{j-t},\qquad
 \operatorname{rem}_P(y^D)=y^D-P(y)Q_D(y),\\
 [y^a]\operatorname{rem}_P(y^D)
   &=-\sum_{t=0}^j h_t(\omega)(-1)^{D-t-a}e_{D-t-a}(\omega),
       \qquad 0\le a<q.
\end{aligned}
 \tag{OKV8}
\]
To verify the quotient assertion, the power-series identity makes
the coefficients of $y^D,y^{D-1},\ldots,y^q$ in $P Q_D$ equal
to those of $y^D$. The remaining coefficients are exactly the last
line of OKV8. Their degree is less than $q$, so uniqueness of monic
polynomial division proves both assertions. For $0\le D<q$ the
remainder is $y^D$ itself.

The actual maximal root modulus is
\[
 R=k\sqrt{\delta^2+\gamma^2},\qquad
 D_*:=\max\{3,\sqrt{\delta^2+\gamma^2}\},\qquad T_*:=D_*q.
 \tag{OKV9}
\]
Every $|\omega_{ab}|\le R$, and equality holds at each corner,
so this expression for $R$ is exact. In particular $R\le T_*$.
Counting the summands of OKV7 proves
\[
 |e_h(\omega)|\le\binom qh R^h,\qquad
 |h_t(\omega)|\le\binom{q+t-1}{t}R^t.
\]
The second count follows by placing $q-1$ separators among $t+q-1$
positions to encode $(n_1,\ldots,n_q)$. For fixed $j$, the binomial
$\binom{q+t-1}{t}$ is increasing in $0\le t\le j$, because its
successive ratio is $(q+t)/(t+1)\ge1$. In each nonzero summand
of OKV8 the total root exponent is $(D-t-a)+t=D-a$.
As $t$ varies the indices $D-t-a$ are distinct, so the sum of
their $\binom q{D-t-a}$ is at most $2^q$. It follows that
\[
 \left|[y^a]\operatorname{rem}_P(y^{q+j})\right|
 \le2^q\binom{q+j-1}{j}R^{q+j-a}
 \le2^{3q-1}R^{q+j-a},\qquad 0\le a<q,\quad0\le j\le q.
 \tag{OKV10}
\]
For the last inequality use
$\binom{q+j-1}{j}\le2^{q+j-1}\le2^{2q-1}$, by the binomial
expansion of $(1+1)^{q+j-1}$. These estimates retain the exact
coefficient formula OKV8, including every signed cancellation.

\paragraph{Original Gamma norms of the low monomials.}
The same recurrence and norms in OKV5 give, without changing the measure,
\[
 y\varphi_a=\sqrt{(a+1)(b_s+a)}\,\varphi_{a+1}
          +\sqrt{a(b_s+a-1)}\,\varphi_{a-1},
 \quad \varphi_{-1}=0.
 \tag{OKV11}
\]
For a polynomial of degree at most $q-2$, its orthonormal coefficient
vector is therefore sent by multiplication by $y$ to the sum of
one upward and one downward weighted shift, with output degree at most
$q-1$. Each shift has operator norm at most
$\sqrt{q(b_s+q)}$: the squared norm of its coefficient vector is
bounded by the largest squared weight times the input squared norm.
Since $s\le k<q$ gives $b_s\le q/2$, the triangle inequality yields
\[
 \|yf\|_{L^2(m_{0,s})}
       \le2\sqrt{q(b_s+q)}\,\|f\|_{L^2(m_{0,s})}
       \le3q\,\|f\|_{L^2(m_{0,s})}
       \le T_*\|f\|_{L^2(m_{0,s})}.
\]
Here $2\sqrt{3/2}=\sqrt6<3$. Starting from
$\|1\|_{L^2(m_{0,s})}=\sqrt{M_s}$ and applying this inequality
to $1,y,\ldots,y^{a-1}$ proves
\[
 \boxed{\|y^a\|_{L^2(m_{0,s})}\le\sqrt{M_s}\,T_*^a,
                      \qquad0\le a<q.}
 \tag{OKV12}
\]
All intermediate polynomial degrees satisfy the stated operator domain.

For any polynomial $f=\sum_a f_a y^a$, write
$\mathfrak c_{T_*}(f)=\sum_a|f_a|T_*^a$. This is a finite
coefficient sum, not a change of source or quotient norm.
The exact original recurrence gives
\[
 \mathfrak c_{T_*}(p_{d+1})
 \le T_*\mathfrak c_{T_*}(p_d)
       +d(b_s+d-1)\mathfrak c_{T_*}(p_{d-1}).
\]
For $1\le d<2q$ its actual positive coefficient is bounded by
$d(b_s+d-1)\le2q(b_s+2q)\le5q^2\le T_*^2$.
The initial values are $1,T_*$. Induction now proves
\[
 \boxed{\mathfrak c_{T_*}(p_d)\le(2T_*)^d,
                            \qquad0\le d\le2q.}
 \tag{OKV13}
\]
Indeed substituting the two preceding bounds in the recurrence gives
$T_*(2T_*)^d+T_*^2(2T_*)^{d-1}
 =\tfrac34(2T_*)^{d+1}\le(2T_*)^{d+1}$.

Combining all $q$ coefficients of OKV10 with OKV12, by the triangle
inequality in the original $L^2$ space, gives
\[
 \|\operatorname{rem}_P(y^D)\|_{L^2(m_{0,s})}
           \le\sqrt{M_s}\,q2^{3q}T_*^D\quad(0\le D\le2q).
\]
For $D\ge q$, each summand is bounded by
$\sqrt{M_s}2^{3q-1}R^{D-a}T_*^a
 \le\sqrt{M_s}2^{3q-1}T_*^D$ and there are $q$ summands.
For $D<q$, the assertion follows from OKV12 and $q2^{3q}\ge1$.
Linearity of the exact remainder map, followed by OKV13, proves
\[
 \boxed{\|\operatorname{rem}_P(p_d)\|_{L^2(m_{0,s})}
         \le\sqrt{M_s}\,q2^{3q}(2T_*)^d,
                           \qquad0\le d\le2q.}
 \tag{OKV14}
\]

\paragraph{The original remainder vectors and the full matrix norm.}
The polynomial $\operatorname{rem}_P(p_d)$ has degree less than $q$.
Its expansion in the orthonormal basis
$\varphi_0,\ldots,\varphi_{q-1}$ has squared coefficient norm equal
to its squared original $L^2(m_{0,s})$ norm, by expanding the finite
inner product and using OKV5. The exact quotient map OKV3 and
definition OKV6 therefore give, for $d=q+j$, $0\le j\le q$,
\[
 \|v_j\|_2^2
   =\frac{\|\operatorname{rem}_P(p_{q+j})\|_{L^2(m_{0,s})}^2}
                 {M_s(q+j)!(b_s)_{q+j}}
   \le\frac{q^22^{6q}(2T_*)^{2(q+j)}}{(q+j)!(b_s)_{q+j}}.
 \tag{OKV15}
\]
This equality is the actual map of the original quotient coordinates
to the remainder polynomial. The full source factor $M_s$ occurs in
both the numerator bound and its original norm denominator before
the displayed cancellation. No estimate on the conditioning of
$C_s$ or a companion matrix is used.

Since $b_s\ge1/2$, for each integer $d\ge1$ the first factor of
$(b_s)_d$ is at least $1/2$ and every subsequent factor is at least
one. Thus $d!(b_s)_d\ge1/2$. Also $T_*\ge1$ and $q+j\le2q$.
Consequently
\[
 \|v_j\|_2^2\le2q^22^{6q}(2D_*q)^{4q},\qquad0\le j\le q,
\]
and the original complete matrix obeys
\[
 \boxed{1\le\|K_{q+1}\|_2\le
   1+2(q+1)q^22^{6q}(2D_*q)^{4q}.}
 \tag{OKV16}
\]
To check its norm bound, $\|v_jv_j^*\|_2=\|v_j\|_2^2$
follows directly from Cauchy--Schwarz, with equality on a nonzero
$v_j$; it is zero for a zero vector. Add all $q+1$ summands of
OKV6 and the norm of $I_q$. Positivity also proves its lower bound.

Here is an explicit bound of the requested scale, with constants
depending only on the original two root coordinates:
\[
 \boxed{\log\|K_{q+1}\|_2
       \le C(\delta,\gamma)\,q\log(q+1),\qquad
 C(\delta,\gamma)=12+4\log(2D_*).}
 \tag{OKV17}
\]
For a direct verification, put
$U=2(q+1)q^22^{6q}(2D_*q)^{4q}\ge1$ and use
$\log(1+U)\le\log(2U)$. This upper bound is exactly
\[
 \log4+\log(q+1)+2\log q+6q\log2
       +4q\log(2D_*)+4q\log q.
\]
Write $L=\log(q+1)$. Since $q\ge36$, one has
$\log4\le L$, $\log q\le L$, $6\log2\le2L$, and $L\ge1$.
The last inequality follows already from
$\log4=\int_1^4dt/t>1$, by bounding the integrand below by
$1/2$ on $[1,2]$ and $1/4$ on $[2,4]$.
The preceding display is therefore at most
$4L+6qL+4\log(2D_*)qL$, which is bounded by the right side
of OKV17. The completely finite and sharper bound OKV16 remains
available before this logarithmic estimate.

\paragraph{The four original kernel Grams with their full coefficient frame.}
Retain the original surjection $\Lambda:E_\chi\to B_{\rm obs}$,
the actual subspace $\mathcal K_\Lambda=\ker\Lambda$, and its original
inclusion $I:\mathcal K_\Lambda\hookrightarrow E_\chi$. Put
$r_K=\dim\mathcal K_\Lambda$; no separate formula or upper bound
for this actual dimension is needed here. Use its fixed original
coefficient frame and write
\[
 T_s=C_s^{-1}I,\qquad
 G_{q+j-1}^{(s)}=C_s^{-*}K_j^{-1}C_s^{-1},\qquad
 H_j^{(s)}=I^*G_{q+j-1}^{(s)}I=T_s^*K_j^{-1}T_s,
             \qquad0\le j\le q+1.
 \tag{OKV18}
\]
For completeness, the quotient metric in this identity follows from
the original orthonormal source-coordinate map
$A_j=[I_q,v_0,\ldots,v_{j-1}]$.
It obeys $A_jA_j^*=K_j$. The minimum norm solution of $A_jz=x$
is $z=A_j^*K_j^{-1}x$: it has the required image and is orthogonal
to $\ker A_j$. Every other solution differs by such a kernel vector,
so the Pythagorean identity gives minimum squared norm $x^*K_j^{-1}x$.
The actual remainder is $C_sx$. Substitution proves the original
quotient Gram and then its full restriction in OKV18.

Let $\lambda_s=\|K_{q+1}\|_2$. The original update chain gives
$I_q=K_0\preceq K_j\preceq K_{q+1}\preceq\lambda_s I_q$.
For positive definite matrices, inversion reverses the order:
if $A\preceq B$, conjugate by $A^{-1/2}$ to obtain a matrix
with eigenvalues at least one, invert those eigenvalues and conjugate
back. Thus the full original kernel forms satisfy
\[
 \lambda_s^{-1}H_0^{(s)}\preceq H_j^{(s)}\preceq H_0^{(s)},
 \qquad H_0^{(s)}=T_s^*T_s,
 \qquad H_{j+1}^{(s)}\preceq H_j^{(s)}.
 \tag{OKV19}
\]
For $r_K>0$, the matrix $T_s$ is injective, so $H_0^{(s)}$ is
positive definite. Conjugation in OKV19 by its inverse square root
gives eigenvalues in $[\lambda_s^{-1},1]$. The determinant identity
for that exact congruence is
\[
 \det\bigl[(H_0^{(s)})^{-1/2}H_j^{(s)}(H_0^{(s)})^{-1/2}\bigr]
       =\frac{\det H_j^{(s)}}{\det(T_s^*T_s)}.
\]
Therefore
\[
 -r_K\log\lambda_s
 \le\log\det H_j^{(s)}-\log\det(T_s^*T_s)\le0.
 \tag{OKV20}
\]
In particular the determinant and conditioning of the actual
$T_s^*T_s$ are retained in this identity; no new kernel basis is
substituted for its original frame.

Define the original four-endpoint kernel volume, with its literal
degrees $q-1,q,2q-1,2q$, by
\[
 F_K^{(s)}=
 \log\det H_0^{(s)}+\log\det H_1^{(s)}
       -\log\det H_q^{(s)}-\log\det H_{q+1}^{(s)}.
 \tag{OKV21}
\]
Monotonicity in OKV19 makes each of
$\log\det H_0^{(s)}-\log\det H_q^{(s)}$ and
$\log\det H_1^{(s)}-\log\det H_{q+1}^{(s)}$ nonnegative.
Each is at most $r_K\log\lambda_s$ by OKV20. Consequently the
actual kernel volume satisfies the complete finite estimates
\[
 \boxed{\begin{aligned}
 0\le F_K^{(s)}
 &\le2r_K\log\|K_{q+1}\|_2\\
 &\le2r_K\log\left[1+2(q+1)q^22^{6q}(2D_*q)^{4q}\right]\\
 &\le2r_K C(\delta,\gamma)q\log(q+1),
                    \qquad s\in\{1,k\}.
 \end{aligned}}
 \tag{OKV22}
\]
For $r_K=0$, every kernel determinant has its original empty value
one, and $F_K^{(s)}=0$. Hence OKV22 includes this case.
The four copies of $\log\det(T_s^*T_s)$ in OKV21 have coefficients
$+1,+1,-1,-1$ and cancel exactly only after their presence in
OKV20 has been recorded.

\paragraph{The original boundary and arithmetic transitions.}
At each original degree set
\[
 Q_j^{(s)}=
   \left(\Lambda (G_{q+j-1}^{(s)})^{-1}\Lambda^*\right)^{-1},
 \qquad
 S_j^{(s)}=(G_{q+j-1}^{(s)})^{-1}\Lambda^*Q_j^{(s)}.
 \tag{OKV23}
\]
Their zero dimensional cases have their empty meaning. Surjectivity
of $\Lambda$ and positivity give the displayed inverses. Direct
multiplication proves $\Lambda S_j^{(s)}=I$ and
$I^*G_{q+j-1}^{(s)}S_j^{(s)}=0$.
For any fixed original section $S_*$ of $\Lambda$, the map $[I,S_*]$
is bijective: its inverse records $\Lambda z$ and the unique kernel
coordinate of $z-S_*\Lambda z$. The minimum section differs from
$S_*$ by a kernel column. The two frames therefore differ by a
block triangular matrix with identity diagonal and determinant one.
Taking the determinant of the full metric in the orthogonal frame
$[I,S_j^{(s)}]$ gives
\[
 \det H_j^{(s)}\det Q_j^{(s)}
       =|\det[I,S_*]|^2\det G_{q+j-1}^{(s)}.
 \tag{OKV24}
\]
This proves the exact relation with all kernel and boundary directions
retained. The same frame factor occurs at all four endpoints.
Their signed sum is consequently
\[
 F_K^{(s)}+F_B^{(s)}=\mathcal B_k^{(s)},\qquad
 F_B^{(s)}=\log\det Q_0^{(s)}+\log\det Q_1^{(s)}
               -\log\det Q_q^{(s)}-\log\det Q_{q+1}^{(s)}.
 \tag{OKV25}
\]
The original matrices $G_{q+j-1}^{(s)}$ decrease with $j$, so their
inverse forms increase. Compression by $\Lambda$ and inversion then
show that $Q_j^{(s)}$ decreases as well. Thus $F_B^{(s)}\ge0$,
and the finite companion bound is
\[
 \max\{0,\mathcal B_k^{(s)}-2r_K C(\delta,\gamma)q\log(q+1)\}
       \le F_B^{(s)}\le\mathcal B_k^{(s)}.
 \tag{OKV26}
\]
Here the finer middle bound in OKV22 can be used in place of its
last line, with the same proof.

Finally the original source identifications are
$\mathcal B_k^{(1)}=\mathcal B_k^\sigma$ and
$\mathcal B_k^{(k)}=\mathcal B_k^0$. Their already proved full
source transition is $\mathcal H_k=\mathcal B_k^\sigma-\mathcal B_k^0$,
and the original arithmetic transition is
$\mathcal B_k^{\rm ar}=\mathcal B_k^\sigma+\delta_k^\sigma$,
with its original finite signed allowance for $\delta_k^\sigma$.
Substitution of the proved exact splitting OKV25 gives, with every
source and arithmetic term still present,
\[
 \boxed{\mathcal B_k^{\rm ar}
   =F_K^{(k)}+F_B^{(k)}+\mathcal H_k+\delta_k^\sigma
   =F_K^{(1)}+F_B^{(1)}+\delta_k^\sigma.}
 \tag{OKV27}
\]
Thus OKV22 controls the original Gamma kernel summands in this
same arithmetic identity. It keeps the separate source transition
and arithmetic scalar with their established signs and bounds.
The result is finite for the actual $r_K=\dim\ker\Lambda$ and
uses no additional asserted dimension estimate or replacement source.

% END COMPLETE OKV

\clearpage
% BEGIN COMPLETE OMV
\section{The actual simple-quartet kernel and boundary at the original volume scale}
\label{sec:OMV}

This section combines the complete original period calculation,
the exact cyclic-orbit rank calculation, and the finite original
Gamma kernel estimate. The hypotheses are the actual simple-quartet
data already used in the programme, and the period range below has
an explicit bound proved from those data. No orbit condition or
kernel-rank premise is added to the statement.

\subsection{The original domain and the proved period range}
Fix the original simple quartet $m=1$, its coordinates
$0<\delta<1/2$, $\gamma>2$, and its actual nonzero amplitude value
$a=v_h(1/2+\delta+i\gamma)$, where $v_h=2\xi/h$ and every
zero order in this quartet is complete. Reflection and conjugation
give the other three unit values exactly as in PQO3 and RPC11.
For every integer $l\geq1$ retain
\[
 k=4l+1,\qquad q=(k+1)^2,\qquad c=k/2,
 \qquad s\in\{1,k\}.
 \tag{OMV1}
\]
The polynomial $\chi$, the source measures $m_{0,s}$, their masses
$(2\pi)^{s/2}$, the original observation $\Lambda_k:E_\chi\to B_k$,
its kernel, and all four endpoint degrees are unchanged.

Let $R_*$ be the finite positive real number defined by RPC17 and
RPC21 when the actual $c_3$ of RPC12 is nonzero, and by RPC17 and
RPC23 when it is zero. Those two explicitly evaluated alternatives
exhaust the original nonzero unit, because $c_1,c_3$ cannot both
vanish. In particular $R_*$ depends on the fixed quartet and its
actual unit and is independent of $k$. Its constants are finite
positive Gamma values and the convergent positive integrals RPC17.
The full period proof RPC18--25 establishes
\[
 |u|\geq R_*\quad\Longrightarrow\quad
 \#\{[Q_u(D^j x)]:0\leq j<5\}=5
 \tag{OMV2}
\]
on every original logarithm branch. To recall the actual nonvanishing
calculation, RPC22 bounds the error in
$u^{-3/5}\mathcal H_{41}$ by half the nonzero value
$|c_3g_4/g_1|$ in the first alternative. RPC24 bounds the error
in $u^{-1/5}\mathcal H_{21}$ by half $|c_1g_2/g_1|$ in the
second. Each is a nonzero off-diagonal entry of the exact
$\mathcal H=P(c_1M+c_3M^3)P^{-1}$; the flat residue identity
RPC14--15 makes it a nonzero noninvariant coefficient of the actual
quadric. The degree-five cyclic action and its nondegenerate
quadratic matrix then give OMV2. All powers of $u$ use its original
branch, and the absolute error estimates hold for every such branch.

The exact invariant-ideal proof OQX1--3 now applies at every point
of this proved domain. It gives, for all the original degrees OMV1,
\[
 \boxed{\quad \dim B_k=k^2-6k+17,\qquad
                 r_K:=\dim\ker\Lambda_k=8k-16.\quad}
 \tag{OMV3}
\]
Specifically, an invariant polynomial vanishing on the original
quadric is divisible by all five distinct original orbit factors;
its quotient by their invariant product is invariant. Thus the
actual observed rank is $d_{5,k,0}-d_{5,k-10,0}$, with negative
degrees zero. The full character calculation OQX3 evaluates this
difference to $k^2-6k+17$, including the two original low degrees
$5,9$. Subtraction from $q=(k+1)^2$ proves the second equality.
This is the dimension of the kernel of the literal period observation,
not an assigned ambient dimension or a presumed maximal rank.

\subsection{A uniform finite bound on the actual kernel volume}
Define, from the original root coordinates only,
\[
 D_*=\max\{3,\sqrt{\delta^2+\gamma^2}\},\quad
 C_*=12+4\log(2D_*),\qquad
 \mathfrak M_k=(16k-32)C_*q\log(q+1).
 \tag{OMV4}
\]
Retain the exact original source-coordinate vectors
$v_j=C_s^{-1}[u_{q+j}]$, $K_j=I+\sum_{a<j}v_av_a^*$,
and the full original kernel frame $T_s=C_s^{-1}I$ of OKV6,18.
The original kernel metrics at degrees $q+j-1$ are
$H_j^{(s)}=T_s^*K_j^{-1}T_s$. Their actual four-endpoint volume
satisfies
\[
 \boxed{\begin{aligned}
 0\leq F_K^{(s)}(u)
 &\leq(16k-32)\log\!\left[
       1+2(q+1)q^22^{6q}(2D_*q)^{4q}\right]\\
 &\leq\mathfrak M_k,
 \qquad s\in\{1,k\},\quad |u|\geq R_*.
 \end{aligned}}\tag{OMV5}
\]
Here the source index and the actual period dependence of the kernel
are retained. The proof is direct substitution of the now proved
$r_K=8k-16$ into OKV22. Its finite coefficient proof retains
every original root in the signed remainder formula, all Gamma
recurrence coefficients through degree $2q$, and the full source
mass before its exact norm-ratio cancellation. Its metric proof
retains $\det(T_s^*T_s)$ at each of the four degrees, and cancels
only the four copies with coefficients $+1,+1,-1,-1$.
In particular neither the conditioning nor the changing position
of the original period kernel has been suppressed. The bound is
uniform in the entire period domain OMV2 because its right side
depends only on $k,\delta,\gamma$, and the dimension in OMV3 is
the same at every point of that domain.

Write the four original total quotient baselines as
$\mathcal B_k^{(1)}=\mathcal B_k^\sigma$ and
$\mathcal B_k^{(k)}=\mathcal B_k^0$. The full determinant identity
OKV24--25, with the original fixed coefficient section, gives
\[
 F_B^{(s)}(u)=\mathcal B_k^{(s)}-F_K^{(s)}(u),\qquad
 \max\{0,\mathcal B_k^{(s)}-\mathfrak M_k\}
           \leq F_B^{(s)}(u)\leq\mathcal B_k^{(s)}.
 \tag{OMV6}
\]
Both factors in the exact splitting are nonnegative by monotonicity
of their original minimum metrics. The literal boundary term is
the actual period-vector expression
\[
 F_B^{(s)}(u)=\log\frac{
       \det\mathcal M_{2q}^{(s)}(u)\det\mathcal M_{2q+1}^{(s)}(u)}{
       \det\mathcal M_q^{(s)}(u)\det\mathcal M_{q+1}^{(s)}(u)},
 \qquad
 \mathcal M_D^{(s)}(u)=\sum_{d=0}^{D-1}
       \lambda_d^{(s)}(u)(\lambda_d^{(s)}(u))^*.
 \tag{OMV7}
\]
The vectors are the complete original primary-period coefficients
OPG14 and OCS21 on the actual image of $\Lambda_k$. Formula OPG17
proves that each consecutive determinant ratio is $1+\xi_j$;
the unchanged weights $1,2,\ldots,2,1$ give exactly OMV7.
The companion kernel vectors retain the minimum-section correction
OPR2 and the action defect OPG22/OQC22. No invariance of that kernel
under the original sum action is needed for OMV5--7.

\subsection{The two individual leading coefficients}
The complete original baseline proofs BSL16--19 and HBL establish,
for fixed original $\delta,\gamma$ and $m=1$, both source limits
\[
 \frac{\mathcal B_k^{(s)}}{q^2}\longrightarrow
 C_B=9-8\log2+F(2,\pi)>\frac{769}{17010},
 \qquad s=1\text{ and }s=k.
 \tag{OMV8}
\]
Here $F(2,\pi)$ is exactly the energy functional evaluated in those
complete proofs; its original endpoint convention and coefficients
are retained. No value of it is redefined by the present comparison.
Since $q=(k+1)^2$,
\[
 \frac{\mathfrak M_k}{q^2}
   =C_*\frac{(16k-32)\log((k+1)^2+1)}{(k+1)^2}
               \longrightarrow0.
\]
For an elementary verification, $((k+1)^2+1)<(k+2)^2$, and
the last fraction is at most $32\log(k+2)/k$. The inequality
$\log x=\int_1^x dt/t\leq2(\sqrt x-1)$ for $x\geq1$
proves that this upper bound tends to zero. The finite estimate
OMV5 and the exact identity OMV6 therefore prove
\[
 \boxed{\quad
 \frac{F_K^{(1)}(u)}{q^2},\ \frac{F_K^{(k)}(u)}{q^2}
       \longrightarrow0,\qquad
 \frac{F_B^{(1)}(u)}{q^2},\ \frac{F_B^{(k)}(u)}{q^2}
       \longrightarrow C_B.
 \quad}\tag{OMV9}
\]
These limits are uniform for all $|u|\geq R_*$ and all original
branches. Indeed the kernel error is bounded by
$\mathfrak M_k/q^2$ uniformly, while the boundary error is at
most
$|\mathcal B_k^{(s)}/q^2-C_B|+\mathfrak M_k/q^2$.
The baseline is independent of $u$. This proves the asserted
uniformity, rather than only a conclusion for each separately
fixed observation. The original combined coefficient has now been
resolved into its two individual contributions on the proved
simple-quartet period domain.

\subsection{Exact propagation to the arithmetic identity and the finer scale}
Keep the original signed transitions
$\mathcal H_k=\mathcal B_k^\sigma-\mathcal B_k^0$ and
$\delta_k^\sigma=\mathcal B_k^{\rm ar}-\mathcal B_k^\sigma$.
Their complete finite allowances, including the QGQ simple-
multiplicity refinement, remain those of the full current receiver.
The exact arithmetic receiving equality becomes
\[
 \begin{aligned}
 \mathcal B_k^{\rm ar}-F_B^{(1)}(u)
       &=F_K^{(1)}(u)+\delta_k^\sigma,\\
 \mathcal B_k^{\rm ar}-F_B^{(k)}(u)
       &=F_K^{(k)}(u)+\mathcal H_k+\delta_k^\sigma.
 \end{aligned}\tag{OMV10}
\]
This is substitution in OKV27, so every source and sign is the
original one. In particular the finite upper and lower allowances
are obtained by adding $[0,\mathfrak M_k]$ to the actual interval
for $\delta_k^\sigma$, or to the actual interval for
$\mathcal H_k+\delta_k^\sigma$, respectively. The original proved
relations $\delta_k^\sigma/(kq)\to0$ and
$\mathcal H_k/(kq)\to-\mathcal C_\Gamma/4$, together with
$k/q\to0$ and OMV5, prove that both differences in OMV10,
divided by $q^2$, tend to zero uniformly on the stated period domain.
This is an exact comparison with the full arithmetic baseline.
The literal arithmetic kernel and boundary metrics retain their
own source until their corresponding metric comparison is applied.

The finer original Gamma return is still carried by an exact mixed
identity:
\[
 \begin{aligned}
 F_B^{(1)}(u)-F_B^{(k)}(u)
       &=\mathcal H_k-
                   \bigl(F_K^{(1)}(u)-F_K^{(k)}(u)\bigr),\\
 -\mathfrak M_k\leq F_K^{(1)}(u)-F_K^{(k)}(u)
       &\leq\mathfrak M_k,\\
 \mathcal H_k-\mathfrak M_k
       \leq F_B^{(1)}(u)-F_B^{(k)}(u)
       &\leq\mathcal H_k+\mathfrak M_k.
 \end{aligned}\tag{OMV11}
\]
Subtract OMV6 for the two actual source orders to obtain the first
line. Each kernel term lies in $[0,\mathfrak M_k]$, so their
difference lies in the stated interval, proving the other two lines.
The error bound here is of order $kq\log(q+1)$, whereas the
already proved signed return $\mathcal H_k$ has order $kq$.
No sign or limiting coefficient for either finer mixed difference
is assigned by OMV11. Its explicit remaining term is the difference
of the two original kernel metrics, on the exactly calculated
$8k-16$ dimensional period kernel. The full source multiplier and
its action on this actual kernel remain the input to that next
calculation, rather than a freely selected source order or kernel.

% END COMPLETE OMV

\clearpage
% BEGIN COMPLETE AMT
% Exact original arithmetic mixed transfer. Historical sources remain unchanged.
\section{The original arithmetic kernel and boundary receive the full source comparison}
\label{sec:AMT}

This calculation uses the proved full-polynomial comparison TW12/TWA3,
the original observation OPG1--5, and the complete original Gamma
kernel estimate OKV1--27. It gives the exact maps between their source,
quotient, kernel and boundary metrics, before taking any determinant.
The subsequent simple-quartet conclusion uses the evaluated original
period test RPC1--26 and its exact rank calculation OQX1--3. Every
source order below is one of the already present orders $1,k$.

\subsection{The complete original source and the same observation}
Retain the original full-order quartet, with $0<\delta<1/2$,
$\gamma>2$, and positive integer complete zero order $m$. Thus
\[
\begin{gathered}
 \mathcal R=\{\tfrac12+\delta+i\gamma,\tfrac12+\delta-i\gamma,
              \tfrac12-\delta+i\gamma,\tfrac12-\delta-i\gamma\},\quad
 h(z)=\prod_{\alpha\in\mathcal R}(z-\alpha)^m,\quad v_h=2\xi/h,\\
 k\ge3,\quad c=k/2,\quad e=1+k(m-1),\quad q=e(k+1)^2,\\
 \lambda_{a,b}=c+(2a-k)\delta+i(2b-k)\gamma,\quad
 \chi(S)=\prod_{a,b=0}^k(S-\lambda_{a,b})^e,\quad E=\mathbb C[S]/(\chi),\\
 w_h(y)=|v_h(1/2+iy)|^2/(2\pi),\quad m_{h,k}=w_h^{*k},\quad
 d\sigma(y)=|\Gamma(1/4+iy/2)|^2dy/(2\pi).
\end{gathered}\tag{AMT1}
\]
The original source inner products, conjugate-linear in their first
argument, are
\[
 \langle P,Q\rangle_{\rm ar}
  =\int_{\mathbb R}\overline{P(c+iy)}Q(c+iy)m_{h,k}(y)\,dy,
 \qquad
 \langle P,Q\rangle_\sigma
  =\int_{\mathbb R}\overline{P(c+iy)}Q(c+iy)\,d\sigma(y).
 \tag{AMT2}
\]
No measure mass or complex-coordinate phase is changed. In particular
$\int d\sigma=\sqrt{2\pi}$. Write $\mathcal P_N$ for polynomials
in the original $S$ of degree at most $N$. In its coefficient frame
$(1,S,\ldots,S^N)$ let the two Grams be $H_N^{\rm ar},H_N^\sigma$.
For $N\ge q-1$, the literal remainder map
$J_N:\mathcal P_N\to E$ is onto and has kernel
$\chi\mathcal P_{N-q}$; at $N=q-1$ this is the zero vector space.
The frame on $E$ is $(1,S,\ldots,S^{q-1})$.

Here is the original observation being retained. Put $n=4m$,
$E_h=\mathbb C[z]/(h)$, $U=M_{j_hv_h}$, and
$\iota[P]=[P(z_1+\cdots+z_k)]\in E_h^{\otimes k}$.
The map $j_h$ records every divided Taylor coefficient through order
$m-1$ at every original root. Its unit $j_hv_h$ is invertible because
the stated order is complete. Fix the original $u\ne0$ and its
logarithm branch. With $\Phi_h'=h$ and $\Phi_h(0)=0$, let
$\ell_j$ be the outward ray of angle
$(\arg u+\pi+2\pi j)/(n+1)$ and $\Gamma_j=-\ell_0+\ell_j$.
The period map and projector are exactly
\[
\begin{gathered}
 (\Pi[P])_j=\int_{\Gamma_j}e^{\Phi_h(z)/u}
                           \operatorname{rem}_hP(z)\,dz,\quad1\le j\le n,\\
 C:\Gamma_j\mapsto\Gamma_{j+1}-\Gamma_1\ (j<n),\quad
        \Gamma_n\mapsto-\Gamma_1,\qquad
 P_k=\frac1{n+1}\sum_{a=0}^{n}(C^{-a})^{\otimes k},\\
 L=P_k\Pi^{\otimes k}U^{\otimes k}\iota,\quad
 B=\operatorname{im}L,\quad\jmath:B\hookrightarrow(\mathbb C^n)^{\otimes k},
 \quad\jmath\Lambda=L,\quad\Lambda:E\twoheadrightarrow B,\quad K=\ker\Lambda.
\end{gathered}\tag{AMT3}
\]
These are the convergent original representative integrals; their
existence and full determinant are proved in PDE and OPG. Fix the
original coefficient inclusion $I:K\hookrightarrow E$ and section
$S_*:B\to E$ with $\Lambda S_*=I_B$. Let
$r=\dim K$ and $b=\dim B=q-r$. All matrices below use these same
frames and this same full unit and period observation. Empty Grams
have determinant one, and maps to or from a zero space have their
unique empty values.

\subsection{All constants of the actual full-polynomial inequality}
Retain exactly the constants of TW and TWA:
\[
\begin{gathered}
 \alpha=\pi/2,\quad p=8m-\tfrac92,\quad B_h=42+8m,\quad M=21+4m,\\
 \vartheta_h=\int_{-1}^1w_h(y)\,dy>0,\quad
 M_\alpha=\int_{\mathbb R}e^{\alpha y}w_h(y)\,dy\in(0,\infty),\\
 c_\Gamma=\sqrt2e^{-7/6},\quad C_\Gamma=\sqrt{2\sqrt5}e^{2/3},\\
 a_k=\frac{c_h\vartheta_h^{k-3}e^{-\alpha(k-3)}(k-2)^{-B_h}}
                   {C_\Gamma2^{B_h/2}},\qquad
 A_k=\frac{C_h^{\rm tilt}}{c_\Gamma}k^{p+1}M_\alpha^{k-1},\qquad
 D_N=[1+4(N+M)^2]^M.
\end{gathered}\tag{AMT4}
\]
Here $c_h$ is the proved original convolution lower-envelope constant
TW7 and $C_h^{\rm tilt}$ is the original compact/tail maximum BT12.
Their full constructions occur in the preserved arithmetic endpoint
and BT providers; they depend on this fixed $h$ and not on $k,N,u$.
Their defining inequalities, proved there with the original numerator,
are
\[
\begin{aligned}
 m_{h,k}(y)&\ge c_h\vartheta_h^{k-3}
    e^{-\alpha(|y|+k-3)}(k-2+|y|)^{-B_h},\\
 m_{h,k}(y)&\le k C_h^{\rm tilt}M_\alpha^{k-1}
              e^{-\alpha|y|}(1+|y|/k)^{-p},\\
 c_\Gamma e^{-\alpha|y|}(1+|y|)^{-1/2}
 &\le d\sigma/dy\le
 C_\Gamma e^{-\alpha|y|}(1+|y|)^{-1/2}.
\end{aligned}\tag{AMT5}
\]
For clarity the full polynomial step is given here. The first and
third bounds, $k-2+|y|\le(k-2)(1+|y|)$, and
$(1+|y|)^{B_h}\le2^{B_h/2}(1+y^2)^M$, give
$m_{h,k}\ge a_k(1+y^2)^{-M}d\sigma/dy$. The second and third,
using $1+|y|/k\ge(1+|y|)/k$ and $p>1/2$, give
$m_{h,k}\le A_kd\sigma/dy$.
The fixed Gamma orthonormal polynomials obey
\[
 ye_j=\sqrt{(j+1)(j+1/2)}e_{j+1}
                         +\sqrt{j(j-1/2)}e_{j-1}.
\]
The two shifts on degree at most $N$ have norms at most $N+1$ and
$N$, respectively. Thus $\|yP\|_\sigma\le2(N+1)\|P\|_\sigma$.
Iterating at the successively increased degrees and expanding
$(1+y^2)^M$ proves
$\int|P|^2(1+y^2)^M d\sigma\le D_N\int|P|^2d\sigma$.
Cauchy--Schwarz applied to $|P|(1+y^2)^{-M/2}$ and
$|P|(1+y^2)^{M/2}$ then gives the reciprocal lower bound. The
coordinate map $P(S)\mapsto P(c+iy)$ has inverse
$f(y)\mapsto f((S-c)/i)$, preserves the degree and retains every
coefficient through this explicitly invertible substitution. It
therefore proves on the complete original source
\[
 \boxed{\quad \frac{a_k}{D_N}H_N^\sigma
                  \preceq H_N^{\rm ar}\preceq A_kH_N^\sigma.\quad}
 \tag{AMT6}
\]
All coefficients can be complex in this proof. In particular AMT6
holds on every affine remainder fibre, not only on monic vectors.
The two Grams are positive: the Gamma density is positive and the
lower bound in AMT5 is positive everywhere. Their finite moments
exist by the exponential upper bound.

\subsection{Exact remainder, minimum-section, kernel and boundary maps}
For $a\in\{\sigma,\mathrm{ar}\}$ put
\[
\begin{gathered}
 G_N^a=(J_N(H_N^a)^{-1}J_N^*)^{-1},\qquad
 V_N^a=(H_N^a)^{-1}J_N^*G_N^a:E\longrightarrow\mathcal P_N,\\
 H_{K,N}^a=I^*G_N^aI,\qquad
 Q_N^a=(\Lambda(G_N^a)^{-1}\Lambda^*)^{-1},\qquad
 S_N^a=(G_N^a)^{-1}\Lambda^*Q_N^a:B\longrightarrow E.
\end{gathered}\tag{AMT7}
\]
Inverses are on the indicated nonzero spaces. For an onto matrix
$J$ and positive $H$, direct multiplication gives $JV=I$ for
$V=H^{-1}J^*(JH^{-1}J^*)^{-1}$. Moreover
$Z^*HV=Z^*J^*G=0$ for $JZ=0$. Therefore every lift of $z$ is
$Vz+Z$, with squared norm
$z^*Gz+Z^*HZ$, proving its unique minimum and its Gram. Applying
this proof first to $J_N$, then to $\Lambda$, proves every formula
AMT7 and $\Lambda S_N^a=I_B$, $I^*G_N^aS_N^a=0$.
It also proves that $V_N^aS_N^a$ is the minimum lift for the
composite $\Lambda J_N$.

The exact comparison map on $\mathcal P_N$ is the identity on
original coefficients; it commutes with $J_N$ and $\Lambda J_N$.
It induces the identity on $E,K,B$. Between the two spaces of
minimum lifts the exact bijection and its inverse are
\[
 \operatorname{im}V_N^\sigma\longrightarrow\operatorname{im}V_N^{\rm ar},
 \quad V_N^\sigma z\longmapsto V_N^{\rm ar}z,
 \qquad T_N=V_N^{\rm ar}J_N\big|_{\operatorname{im}V_N^\sigma},
 \quad T_N^{-1}=V_N^\sigma J_N\big|_{\operatorname{im}V_N^{\rm ar}}.
 \tag{AMT8}
\]
Their difference $V_N^{\rm ar}z-V_N^\sigma z$ lies in
$\chi\mathcal P_{N-q}$, with its unique coefficient obtained by
division by the unchanged monic $\chi$. This also gives the
minimum-lift bijection on the preimage of $K$ and the analogous
bijection $V_N^\sigma S_N^\sigma y\mapsto
V_N^{\rm ar}S_N^{\rm ar}y$ for boundary lifts. The latter
differences lie in $\ker(\Lambda J_N)$, exactly as multiplication
by $\Lambda J_N$ verifies.

Minimize both sides of AMT6 over $J_NP=z$. Taking infima preserves
each inequality because these are the same nonempty fibres. Restrict
the resulting inequality to $z=Ix$, and independently minimize it
over $\Lambda z=y$. This proves the three full-vector comparisons
\[
\boxed{\begin{gathered}
 (a_k/D_N)G_N^\sigma\preceq G_N^{\rm ar}\preceq A_kG_N^\sigma,\\
 (a_k/D_N)H_{K,N}^\sigma\preceq H_{K,N}^{\rm ar}
                                      \preceq A_kH_{K,N}^\sigma,\\
 (a_k/D_N)Q_N^\sigma\preceq Q_N^{\rm ar}\preceq A_kQ_N^\sigma.
\end{gathered}}\tag{AMT9}
\]
The source identity is an exact bounded bijection with the two
constants in AMT6. No invariance of $K$ under an additional
multiplier was used to obtain its induced identities.

The matrix $[I,S_*]$ is invertible: its inverse sends $z$ to the
unique kernel coordinate of $z-S_*\Lambda z$ and the boundary
coordinate $\Lambda z$. Each $S_N^a-S_*$ lies in $K$, so the
change from this frame to $[I,S_N^a]$ is triangular with identity
diagonal and determinant one. The latter frame is orthogonal in
$G_N^a$ by AMT7. Consequently
\[
 \boxed{\quad\det H_{K,N}^a\det Q_N^a
       =|\det[I,S_*]|^2\det G_N^a,\qquad a\in\{\sigma,\mathrm{ar}\}.\quad}
 \tag{AMT10}
\]
This includes $r=0$ and $b=0$ and retains the original frame factor.

\subsection{The exact nested deficits and their rank factors}
Define the actual finite source determinant deficit and its width
\[
 \Xi_{k,N}=(N+1)\log A_k-
                  \log\frac{\det H_N^{\rm ar}}{\det H_N^\sigma},
 \qquad L_N=\log(A_kD_N/a_k)\ge0.
 \tag{AMT11}
\]
The last sign follows by applying AMT6 to any nonzero vector.
Put
\[
\begin{aligned}
 d_{E,N}&=q\log A_k-\log\frac{\det G_N^{\rm ar}}{\det G_N^\sigma},\\
 d_{K,N}&=r\log A_k-\log\frac{\det H_{K,N}^{\rm ar}}{\det H_{K,N}^\sigma},\\
 d_{B,N}&=b\log A_k-\log\frac{\det Q_N^{\rm ar}}{\det Q_N^\sigma}.
\end{aligned}\tag{AMT12}
\]
AMT9, by congruence with the inverse square root of each comparison
Gram, places each eigenvalue of its relative Gram in
$[a_k/D_N,A_k]$. Multiplication of all eigenvalues gives
$0\le d_{K,N}\le rL_N$, $0\le d_{B,N}\le bL_N$ and
$0\le d_{E,N}\le qL_N$.

Here is the full source-to-quotient determinant calculation, including
all relation directions. Choose an $H_N^\sigma$-orthonormal
coefficient frame $R_N$ of $\ker J_N$ solely for this computation,
and put $W_N=[R_N,V_N^\sigma(G_N^\sigma)^{-1/2}]$.
AMT7 gives $W_N^*H_N^\sigma W_N=I_{N+1}$. In this frame the
positive contraction
\[
 \mathscr D_N=A_k^{-1}W_N^*H_N^{\rm ar}W_N
       =\begin{pmatrix}F_N&C_N'\\(C_N')^*&A_N'\end{pmatrix}
       \preceq I_{N+1}
\]
has Schur complement
\[
 T_N'=A_N'-(C_N')^*F_N^{-1}C_N'
       =A_k^{-1}(G_N^\sigma)^{-1/2}G_N^{\rm ar}(G_N^\sigma)^{-1/2}.
\]
Indeed minimization in its first block is exactly the original
remainder minimum at $z=(G_N^\sigma)^{-1/2}w$.
Both $F_N$ and $T_N'$ are positive contractions: positivity follows
from positive block elimination, and the upper bound for the Schur
complement follows by testing its minimum at zero first coordinate.
Since $|\det W_N|^2=(\det H_N^\sigma)^{-1}$, elimination gives
\[
 \Xi_{k,N}=d_{R,N}+d_{E,N},\quad d_{R,N}:=-\log\det F_N\ge0,
 \qquad d_{E,N}=-\log\det T_N'\ge0.
 \tag{AMT13}
\]
When $N=q-1$, $F_N$ is empty, $d_{R,N}=0$, and the same equalities
hold without a relation inverse. The determinant identities are
independent of the auxiliary orthonormal frame.
Dividing the two instances of AMT10 proves
$d_{K,N}+d_{B,N}=d_{E,N}$. Thus the complete nested statement is
\[
\boxed{\begin{gathered}
 \Xi_{k,N}=d_{R,N}+d_{K,N}+d_{B,N},\qquad d_{R,N},d_{K,N},d_{B,N}\ge0,\\
 0\le d_{K,N}\le\min\{rL_N,d_{E,N}\}\le\min\{rL_N,\Xi_{k,N}\},\\
 0\le d_{B,N}\le\min\{bL_N,d_{E,N}\}\le\min\{bL_N,\Xi_{k,N}\}.
\end{gathered}}\tag{AMT14}
\]
In particular each metric keeps its own rank factor, and the two
deficits retain their exact common sum. Any proved finite upper
bound for the original $\Xi_{k,N}$, including the degree-specific
HMR44/TWA15 bound, can be substituted in these displayed minima.
No new density or source determinant is introduced by this substitution.

\subsection{All four original endpoints and the correlated mixed return}
Let
\[
 (N_0,N_1,N_2,N_3)=(q-1,q,2q-1,2q),\qquad
 (\epsilon_0,\epsilon_1,\epsilon_2,\epsilon_3)=(1,1,-1,-1),
\]
and define, in their original signs,
\[
\begin{aligned}
 F_K^a&=\sum_{j=0}^3\epsilon_j\log\det H_{K,N_j}^a,\qquad
 F_B^a=\sum_{j=0}^3\epsilon_j\log\det Q_{N_j}^a,\\
 \mathcal B_k^a&=\sum_{j=0}^3\epsilon_j\log\det G_{N_j}^a,
 \quad\Delta_K^\sigma=F_K^{\rm ar}-F_K^\sigma,
 \quad\Delta_B^\sigma=F_B^{\rm ar}-F_B^\sigma,
 \quad\delta_k^\sigma=\mathcal B_k^{\rm ar}-\mathcal B_k^\sigma.
\end{aligned}\tag{AMT15}
\]
AMT10 cancels its four identical frame factors only after these
signs are applied. It gives
$F_K^a+F_B^a=\mathcal B_k^a$ and
$\Delta_K^\sigma+\Delta_B^\sigma=\delta_k^\sigma$ exactly.
The spaces $\mathcal P_N$ increase with the same source norm and
remainder map, so their quotient minima decrease. Restriction to
$K$ and minimization onto $B$ preserve that order. Pairing $N_0$
with $N_2$ and $N_1$ with $N_3$ consequently proves
$F_K^a,F_B^a\ge0$.

For $X=K,B$, put $r_K=r$, $r_B=b$, and retain the exact caps
\[
 c_{X,N}=\min\{r_XL_N,d_{E,N}\},\qquad
 C_X^-=c_{X,q-1}+c_{X,q},\qquad
 C_X^+=c_{X,2q-1}+c_{X,2q}.
 \tag{AMT16}
\]
The entirely source-evaluated caps
$\widehat c_{X,N}=\min\{r_XL_N,\Xi_{k,N}\}$ and the corresponding
$\widehat C_X^\pm$ are at least these values; all conclusions below
also hold with the hats. Substituting AMT12 in AMT15 records each
endpoint and its sign:
\[
\boxed{\begin{aligned}
 \Delta_K^\sigma&=-d_{K,q-1}-d_{K,q}+d_{K,2q-1}+d_{K,2q},\\
 \Delta_B^\sigma&=-d_{B,q-1}-d_{B,q}+d_{B,2q-1}+d_{B,2q},\\
 \delta_k^\sigma&=-d_{E,q-1}-d_{E,q}+d_{E,2q-1}+d_{E,2q},\\
 -C_X^-&\le\Delta_X^\sigma\le C_X^+,\qquad X=K,B.
\end{aligned}}\tag{AMT17}
\]
The four copies of $r_X\log A_k$ cancel here because their signs
sum to zero. Since the original total correction is retained, the
joint interval is stronger than two independent intervals:
\[
\boxed{\begin{aligned}
 \max\{-C_K^-,\delta_k^\sigma-C_B^+\}
   &\le\Delta_K^\sigma\le
          \min\{C_K^+,\delta_k^\sigma+C_B^-\},\\
 \max\{-C_B^-,\delta_k^\sigma-C_K^+\}
   &\le\Delta_B^\sigma\le
          \min\{C_B^+,\delta_k^\sigma+C_K^-\}.
\end{aligned}}\tag{AMT18}
\]
These follow by subtracting the interval for either component from
the exact common sum; they require no independence assertion.

Here are the explicit original four endpoint widths:
\[
\begin{aligned}
 E_k^+&=L_{q-1}+L_q
   =2\log(A_k/a_k)+M\log[1+4(q-1+M)^2]
                         +M\log[1+4(q+M)^2],\\
 E_k^-&=L_{2q-1}+L_{2q}
   =2\log(A_k/a_k)+M\log[1+4(2q-1+M)^2]
                         +M\log[1+4(2q+M)^2].
\end{aligned}\tag{AMT19}
\]
They are the unchanged TWA5 constants, both nonnegative. AMT17
therefore gives in particular
\[
 \boxed{-rE_k^+\le F_K^{\rm ar}-F_K^\sigma\le rE_k^-,\qquad
        -bE_k^+\le F_B^{\rm ar}-F_B^\sigma\le bE_k^-.}
 \tag{AMT20}
\]
The total correction retains TWA15 with all four individual source
deficits:
\[
\begin{aligned}
 \underline\delta_k^\sigma
   &=\max\{-qE_k^+,-\Xi_{k,q-1}-\Xi_{k,q}\},\\
 \overline\delta_k^\sigma
   &=\min\{qE_k^-,\Xi_{k,2q-1}+\Xi_{k,2q}\},\qquad
 \underline\delta_k^\sigma\le\delta_k^\sigma\le\overline\delta_k^\sigma.
\end{aligned}\tag{AMT21}
\]
This also follows immediately from AMT14 and AMT17. It can be
inserted into AMT18 when an interval, rather than the actual finite
total determinant, is wanted. Thus the lower and upper endpoints
are kept in their distinct places throughout the mixed transfer.

For later use, subtraction of the exact logarithms in AMT4 gives
\[
\begin{gathered}
 \mathfrak s_h=\alpha+\log(M_\alpha/\vartheta_h),\qquad
 \mathfrak c_h=-\log M_\alpha+3\log\vartheta_h-3\alpha
       +\log\frac{C_h^{\rm tilt}C_\Gamma2^{B_h/2}}{c_hc_\Gamma},\\
 \log(A_k/a_k)=k\mathfrak s_h+(p+1)\log k
                         +B_h\log(k-2)+\mathfrak c_h,\qquad
 E_k^\pm=O_h(k+\log q).
\end{gathered}\tag{AMT22}
\]
The last statement follows from the displayed exact expressions,
since $M,p,B_h$ and the two integral constants are fixed and
$\log D_N=O_h(1+\log(N+1))$. It is a consequence of the original
constants and contains no change of measure.

\subsection{The full original multiplier and its transported observation}
On the existing tensor class $k=4l+1$, retain
\[
\begin{gathered}
 t_j=2j+\tfrac12,\quad\xi_j=c+t_j,\quad
 f_l(S)=\prod_{j=0}^{l-1}(S-c-t_j),\\
 \beta_k=\frac{(2\pi)^{k/2}}{\sqrt2\Gamma(k/2)},\quad
 dm_{0,k}(y)=\beta_k|f_l(c+iy)|^2d\sigma(y).
\end{gathered}\tag{AMT23}
\]
This is the exact Gamma recurrence identity TWA26 with masses
$(2\pi)^{k/2}$ and $\sqrt{2\pi}$. For $l=0$ both $f_l$ and
$\beta_k$ are one. Write $\mathsf E_\chi:E\to\mathbb C^q$ for
the complete divided Taylor map
$[P]\mapsto(P^{(r)}(\lambda_{a,b})/r!)_{a,b,\,0\le r<e}$.
It is invertible: a vector in its kernel is divisible by the
coprime factors of $\chi$, and dimensions agree. The full
multiplication matrix $U_f$ has, in each primary block, entries
\[
\begin{gathered}
 (U_f)_{r,a}=\begin{cases}
 f_l^{(r-a)}(\lambda)/(r-a)!,&0\le a\le r<e,\\
 0,&a>r.
 \end{cases}\\
 \det U_f=\prod_{a,b=0}^k\prod_{j=0}^{l-1}
       [(2a-k)\delta-t_j+i(2b-k)\gamma]^e\ne0.
\end{gathered}\tag{AMT24}
\]
Indeed $k$ is odd, so every $\lambda_{a,b}$ has nonzero imaginary
part, while all $\xi_j$ are real and distinct. The inverse in each
block is the finite Taylor inverse of its nonzero constant term.
The exact polynomial map and its inverse on its stated range are
\[
 \mathcal P_N\longrightarrow f_l\mathcal P_N\subseteq\mathcal P_{N+l},
 \quad P\longmapsto f_lP,
 \qquad Q\longmapsto Q/f_l.
 \tag{AMT25}
\]
It sends the original fibre $J_NP=z$ bijectively to the product
jet fibre $(\mathsf E_fQ,\mathsf E_\chi Q)=(0,U_f\mathsf E_\chi z)$.
Vanishing of the $l$ distinct $f$-values proves divisibility by
$f_l$, and the inverse of $U_f$ proves the converse fibre statement.
The original relation space maps to $f_l\chi\mathcal P_{N-q}$.
On the free $\chi$ jet coordinate $Y$ of this constrained product
fibre, the exact observed map, its kernel, and a section are
\[
\boxed{\begin{gathered}
 \widetilde\Lambda=\Lambda\mathsf E_\chi^{-1}U_f^{-1},\qquad
 \widetilde\Lambda U_f\mathsf E_\chi=\Lambda,\\
 \ker\widetilde\Lambda=U_f\mathsf E_\chi K,\qquad
 \widetilde I=U_f\mathsf E_\chi I,\qquad
 \widetilde S_*=U_f\mathsf E_\chi S_*,\quad
 \widetilde\Lambda\widetilde S_*=I_B.
\end{gathered}}\tag{AMT26}
\]
Both inclusions in the kernel equality follow by applying the
displayed invertible map and its inverse. This retains the full
original amplitude in $\Lambda$ and every derivative in $U_f$.
The resulting kernel is proved by transport; it has not been
assigned the coordinates of $\mathsf E_\chi K$.

For an explicit metric on this transported space, let $b_n(y)$ be
the fixed monic Gamma polynomials and put
$\pi_n(S)=i^n b_n((S-c)/i)$, whose leading coefficient is one
and squared norm is $\gamma_n=\sqrt{2\pi}(2n)!/4^n$.
The evaluation kernels are the full finite sums
\[
 \mathscr K_{\chi,L}=\sum_{n=0}^L
   \frac{\mathsf E_\chi\pi_n(\mathsf E_\chi\pi_n)^*}{\gamma_n},\quad
 \mathscr K_{f,L}=\sum_{n=0}^L
   \frac{\mathsf E_f\pi_n(\mathsf E_f\pi_n)^*}{\gamma_n},\quad
 \mathscr K_{\chi,f;L}=\sum_{n=0}^L
   \frac{\mathsf E_\chi\pi_n(\mathsf E_f\pi_n)^*}{\gamma_n}.
\]
The product evaluation is onto for $L=N+l\ge q+l-1$, because
a polynomial of degree below $q+l$ vanishing at all its divided
jets is divisible by the degree-$q+l$ product $f_l\chi$ and is
zero. Thus its positive kernel has positive Schur complement
\[
\begin{gathered}
 C_N=\mathscr K_{\chi,N+l}-\mathscr K_{\chi,f;N+l}
            \mathscr K_{f,N+l}^{-1}\mathscr K_{f,\chi;N+l}>0,\\
 U_N=(\mathsf E_\chi\pi_{N+j}/\sqrt{\gamma_{N+j}})_{j=1}^l,
 \quad V_N=\mathscr K_{\chi,f;N+l}\mathscr K_{f,N+l}^{-1/2},\quad
 C_N=\mathscr K_{\chi,N}+U_NU_N^*-V_NV_N^*.
\end{gathered}\tag{AMT27}
\]
Empty $f$ blocks have the identity convention. The inverse of the
joint positive kernel, evaluated on $(0,Y)$, gives the minimum
$Y^*C_N^{-1}Y$: block elimination of the $f$-block proves this
identity with every mixed entry present. Combining it with AMT25
and the coefficient $\beta_k$ proves the exact tensor-source metrics
\[
\boxed{\begin{gathered}
 G_N^0=\beta_k\mathsf E_\chi^*U_f^*C_N^{-1}U_f\mathsf E_\chi,\qquad
 H_{K,N}^0=\beta_k\widetilde I^*C_N^{-1}\widetilde I,\\
 Q_N^0=\beta_k(\widetilde\Lambda C_N\widetilde\Lambda^*)^{-1},\qquad
 S_N^0=\mathsf E_\chi^{-1}U_f^{-1}C_N\widetilde\Lambda^*
                   (\widetilde\Lambda C_N\widetilde\Lambda^*)^{-1}.
\end{gathered}}\tag{AMT28}
\]
Direct multiplication verifies the section and its orthogonality.
The formulas hold on the original $E,K,B$, with the transported
observation exactly as in AMT26. No block or mixed product was
dropped from $C_N$.

The same coefficient identity also supplies a quantitative comparison
directly with this original tensor source. Define the positive exact
finite factors
\[
 \Pi_l=\prod_{j=0}^{l-1}t_j^2,\qquad
 U_{l,N}^{\rm poly}=\prod_{j=0}^{l-1}[4(N+j+1)^2+t_j^2],\qquad
 \widetilde L_N=L_N+\log(U_{l,N}^{\rm poly}/\Pi_l).
 \tag{AMT29}
\]
Pointwise $|f_l(c+iy)|^2\ge\Pi_l$. Also
$\|(iy-t)P\|_\sigma^2=\|yP\|_\sigma^2+t^2\|P\|_\sigma^2$,
so the degree-specific multiplication estimate in AMT6, applied
successively, gives
$\|f_lP\|_\sigma^2\le U_{l,N}^{\rm poly}\|P\|_\sigma^2$.
Writing $H_N^0$ for the unchanged $m_{0,k}$ source Gram yields
\[
 \beta_k\Pi_lH_N^\sigma\preceq H_N^0
                  \preceq\beta_k U_{l,N}^{\rm poly}H_N^\sigma,
 \qquad
 \frac{a_k}{D_N\beta_k U_{l,N}^{\rm poly}}H_N^0
       \preceq H_N^{\rm ar}\preceq\frac{A_k}{\beta_k\Pi_l}H_N^0.
 \tag{AMT30}
\]
The second pair follows by using the first upper bound for its lower
inequality and the first lower bound for its upper inequality.
Minimizing and restricting with the same $J_N,I,\Lambda$ gives
AMT30 for all three metrics in AMT28. Its common upper scalar is
$A_k/(\beta_k\Pi_l)$; it is retained before its four signed copies
cancel. The same determinant proof as AMT12--17 therefore gives
\[
\begin{gathered}
 \widetilde E_k^+=E_k^+
   +\log(U_{l,q-1}^{\rm poly}U_{l,q}^{\rm poly}/\Pi_l^2),\qquad
 \widetilde E_k^-=E_k^-
   +\log(U_{l,2q-1}^{\rm poly}U_{l,2q}^{\rm poly}/\Pi_l^2),\\
 \boxed{-r\widetilde E_k^+\le F_K^{\rm ar}-F_K^0
                 \le r\widetilde E_k^-,\qquad
        -b\widetilde E_k^+\le F_B^{\rm ar}-F_B^0
                 \le b\widetilde E_k^-.}
\end{gathered}\tag{AMT31}
\]
For an explicit bound on the extra widths, $t_j\ge1/2$ gives
\[
 0\le\log(U_{l,N}^{\rm poly}/\Pi_l)
   =\sum_{j=0}^{l-1}\log\left(1+\frac{4(N+j+1)^2}{t_j^2}\right)
   \le l\log[1+16(N+l)^2].
 \tag{AMT32}
\]
Thus $\widetilde E_k^\pm=O_h(k\log(q+k))$ on the stated endpoints.
This is a proved polynomial-source estimate for the exact original
$f_l$, with no inference that its multiplication preserves $K$.

\subsection{The evaluated original simple-quartet arithmetic kernel}
Now take the original simple-quartet class $m=1$, $k=4l+1$,
$l\ge1$, $q=(k+1)^2$. The parameter $u$ is precisely the period
parameter in AMT3, not the real integration coordinate $y$.
RPC21 and RPC23 give an explicit finite radius $R_*$ for the
unchanged quartet and arithmetic unit. For reference its exact
constants and cases are recorded here. Set
\[
\begin{gathered}
 \beta=2(\gamma^2-\delta^2),\quad\eta=(\delta^2+\gamma^2)^2,
 \quad a=v_h(\tfrac12+\delta+i\gamma)\ne0,\\
 c_3=\frac{a^{-2}-\overline a^{-2}}{4i\delta\gamma},\qquad
 c_1=\frac{a^{-2}+\overline a^{-2}}2-(\delta^2-\gamma^2)c_3,\\
 g_j=5^{j/5-1}e^{\pi ij/5}\Gamma(j/5)\ (1\le j\le4),\quad
 g_+=\max_j|g_j|,\quad g_-^{-1}=\max_j|g_j|^{-1},\quad
 \kappa=\max_{1\le j<4}|g_{j+1}/g_j|,\\
 E_j^{\rm ray}=\frac85\int_0^\infty r^j(r^3/3+r)
                  e^{-r^5/5+r^3/3+r}\,dr\ (0\le j\le3),\quad
 C_T=2\left(\sum_{j=0}^3(E_j^{\rm ray})^2\right)^{1/2},\\
 C_K=4C_Tg_-^{-1}+2g_+g_-^{-1}+2g_+g_-^{-2}C_T,\quad
 C_{K,3}=C_K(3\kappa^2+3\kappa C_K+C_K^2),\\
 \epsilon_*=\min\{1,(2C_Tg_-^{-1})^{-1}\},\qquad B_*=\beta+\eta.
\end{gathered}\tag{AMT33}
\]
Here $g_-^{-2}=(g_-^{-1})^2$. All integrals converge by their
negative fifth-degree exponent, and $(c_1,c_3)\ne(0,0)$ by the
original unit values. The exact radius is
\[
 R_*=
 \begin{cases}
 \displaystyle\left[\max\left\{1,\frac{B_*}{\epsilon_*},
 \frac{2(|c_3|C_{K,3}B_*+|c_1|(\kappa+C_K))}{|c_3g_4/g_1|}\right\}\right]^{5/2},
                  &c_3\ne0,\\[6pt]
 \displaystyle\left[\max\left\{1,\frac{B_*}{\epsilon_*},
                     \frac{2C_KB_*}{|g_2/g_1|}\right\}\right]^{5/2},
                  &c_3=0.
 \end{cases}\tag{AMT34}
\]
The complete period calculation RPC16--25 proves, for every original
branch and every $|u|\ge R_*$, a nonzero original commutator entry:
its $(4,1)$ entry in the first case and its $(2,1)$ entry in the
second. RPC14--15 identifies those entries with the exact test for
the original rank-four quadric's cyclic invariance. Thus its actual
projective orbit has five elements throughout this domain. OQX1--3,
using the product of precisely these five original quadrics, proves
\[
 \boxed{r=8k-16,\qquad b=k^2-6k+17,
        \qquad m=1,\ k=4l+1,\ l\ge1,\ |u|\ge R_*.}
 \tag{AMT35}
\]
This is the evaluated original period outcome; no value of the
amplitude or of a period coefficient is chosen. The radius depends
on the fixed one-factor quartet and unit and is independent of $k$.

For the same original packet, OKV22 proves for both existing source
orders $s=1,k$, and for the actual $K$ of any rank,
\[
\begin{gathered}
 D_*=\max\{3,\sqrt{\delta^2+\gamma^2}\},\quad
 C(\delta,\gamma)=12+4\log(2D_*),\\
 Z_q=1+2(q+1)q^2 2^{6q}(2D_*q)^{4q},\\
 0\le F_K^{(s)}\le2r\log Z_q
                \le2r C(\delta,\gamma)q\log(q+1),
 \\ F_K^{(1)}=F_K^\sigma,\quad F_K^{(k)}=F_K^0.
\end{gathered}\tag{AMT36}
\]
Its proof includes the exact source coefficient map and the unchanged
kernel inclusion. Combining it with the already proved AMT17--20
gives the finite arithmetic conclusion
\[
\boxed{\begin{aligned}
 \max\{0,F_K^\sigma-C_K^-\}
       &\le F_K^{\rm ar}\le F_K^\sigma+C_K^+,\\
 0\le F_K^{\rm ar}
       &\le U_K:=2r\log Z_q+\widehat C_K^+\\
       &\le2r C(\delta,\gamma)q\log(q+1)+rE_k^-.
\end{aligned}}\tag{AMT37}
\]
The cap symbols $C_K^\pm,\widehat C_K^\pm$ in this formula mean
the endpoint sums AMT16; the constant $C_K$ without a superscript
in AMT33 belongs only to the explicit period radius. They have
different stated indices and are not identified.
The alternative finite tensor-source bound
$F_K^{\rm ar}\le2r\log Z_q+r\widetilde E_k^-$ also follows from
AMT31. Formula AMT37 retains the sharper fixed-source widths.

On AMT35 the exact ranks give
\[
 \boxed{\quad F_K^{\rm ar}-F_K^\sigma=O_h(k^2),\qquad
 F_K^{\rm ar}=O_h(kq\log(q+1)),\qquad
 \frac{F_K^{\rm ar}}{q^2}\longrightarrow0.\quad}
 \tag{AMT38}
\]
For the first estimate use $r=8k-16$ in AMT20 and AMT22,
with $\log q=2\log(k+1)$. The second follows from AMT37.
For the limit, divide its last line by $q^2$: its first term is
at most $16C(\delta,\gamma)k\log(q+1)/q\to0$, and its second
is $O_h(k^2/q^2)\to0$. Nonnegativity gives the limit by the two
bounds. These estimates are uniform over the entire original
domain $|u|\ge R_*$, since none of their right sides depends on $u$.
The tensor comparison likewise gives
$F_K^{\rm ar}-F_K^0=O_h(k^2\log q)=o(q^2)$ from AMT31--32.

\subsection{The original boundary receives the full baseline coefficient}
Before using a limit, the exact finite arithmetic identity and its
consequence are
\[
\boxed{\begin{gathered}
 F_B^{\rm ar}=\mathcal B_k^\sigma+\delta_k^\sigma-F_K^{\rm ar},\\
 \max\{0,\mathcal B_k^\sigma+\delta_k^\sigma-U_K\}
    \le F_B^{\rm ar}\le\mathcal B_k^\sigma+\delta_k^\sigma,\\
 \max\{0,\mathcal B_k^\sigma+\underline\delta_k^\sigma-U_K\}
    \le F_B^{\rm ar}\le\mathcal B_k^\sigma+\overline\delta_k^\sigma.
\end{gathered}}\tag{AMT39}
\]
The first identity follows from AMT10 and AMT15; AMT37 and
$F_K^{\rm ar},F_B^{\rm ar}\ge0$ prove the first interval, and
AMT21 proves the second. The correlated boundary return also has
the sharper exact interval AMT18, with every original signed
total correction present.

BSL13--19 supplies the complete original finite baseline and its
proved coefficient
\[
 \mathcal B_k^\sigma=V_q+\eta_q+\eta_{q+1}-\eta_1,\qquad
 \frac{\mathcal B_k^{\rm ar}}{q^2}\longrightarrow
 C_B^{\rm base}:=9-8\log2+\mathcal F(2,\pi)
                      >\frac{769}{17010}>0.
 \tag{AMT40}
\]
Here $V_q$, the three distinct original relation errors $\eta_r$,
and the energy $\mathcal F$ are exactly the BSL/HBL/GEL quantities,
whose complete definitions and proofs are in the retained closure.
The superscript on $C_B^{\rm base}$ only distinguishes that existing
baseline coefficient from the endpoint caps $C_B^\pm$ in AMT16.
Its finite BSL14 enclosure can be substituted directly for
$\mathcal B_k^\sigma$ in AMT39, keeping the signs of all three
$\eta$ terms; no new evaluation is needed for that substitution.
Subtract AMT38 from the exact original identity
$F_B^{\rm ar}=\mathcal B_k^{\rm ar}-F_K^{\rm ar}$. It proves
\[
 \boxed{\quad
 \frac{F_B^{\rm ar}}{q^2}\longrightarrow C_B^{\rm base}>0,
 \qquad
 F_K^{\rm ar}+F_B^{\rm ar}=\mathcal B_k^{\rm ar},
 \quad m=1,\ k=4l+1,\ l\to\infty,\ |u|\ge R_*.
 \quad}\tag{AMT41}
\]
The convergence is uniform over that original period domain by
AMT38 and because the full-source baseline is independent of $u$.
Thus the proved original kernel estimate now passes to the actual
arithmetic kernel, and the original arithmetic boundary carries the
entire positive $q^2$ coefficient. All earlier full-sum identities
remain exact, with their individual mixed contributions now receiving
AMT37--41 and their signed finite corrections AMT17--21.

\subsection{The same comparison evaluates the mixed Gamma return on the \texorpdfstring{$kq$}{kq} scale}
The first inequality pair AMT30 directly compares the two existing
Gamma sources on the same original coefficient space. Apply its
proved minimum and restriction maps AMT7--9 to $X=E,K,B$, writing
$H_{E,N}^a=G_N^a$, $H_{B,N}^a=Q_N^a$, and
$r_E=q,r_K=r,r_B=b$. At every original endpoint it gives
\[
 \beta_k\Pi_l H_{X,N}^\sigma\preceq H_{X,N}^0
       \preceq\beta_k U_{l,N}^{\rm poly}H_{X,N}^\sigma.
 \tag{AMT42}
\]
Define the exact relative determinant and the positive degree width
\[
\begin{gathered}
 z_{X,N}=\log\frac{\det H_{X,N}^0}{\det H_{X,N}^\sigma}
                              -r_X\log(\beta_k\Pi_l),\qquad
 \omega_N^\Gamma=\log(U_{l,N}^{\rm poly}/\Pi_l),\\
 0\le z_{X,N}\le r_X\omega_N^\Gamma,\qquad
 z_{K,N}+z_{B,N}=z_{E,N},\\
 \Omega_k^{\rm low}=\omega_{q-1}^\Gamma+\omega_q^\Gamma,\qquad
 \Omega_k^{\rm high}=\omega_{2q-1}^\Gamma+\omega_{2q}^\Gamma.
\end{gathered}\tag{AMT43}
\]
Congruence and multiplication of the eigenvalues in AMT42 prove
the two inequalities. The determinant identity AMT10, applied to
the tensor metric AMT28 as well, proves the common sum. These
facts preserve the complete original kernel and boundary rather
than assigning separate arbitrary errors to their determinants.

Use the precise return orientation
\[
 \Delta_X^\Gamma=F_X^\sigma-F_X^0\quad(X=K,B),\qquad
 \mathcal H_k=\mathcal B_k^\sigma-\mathcal B_k^0=-\mathfrak T_k.
\]
The last equality is the original TWA30/BSL18 sign. Substituting
all four $z$ terms, and canceling $\log(\beta_k\Pi_l)$ only
after the signed sum, proves
\[
\boxed{\begin{aligned}
 \Delta_X^\Gamma
   &=-z_{X,q-1}-z_{X,q}+z_{X,2q-1}+z_{X,2q},\\
 -r_X\Omega_k^{\rm low}
   &\le\Delta_X^\Gamma\le r_X\Omega_k^{\rm high},\\
 \Delta_K^\Gamma+\Delta_B^\Gamma&=\mathcal H_k,\\
 \mathcal H_k-r\Omega_k^{\rm high}
   &\le\Delta_B^\Gamma\le\mathcal H_k+r\Omega_k^{\rm low}.
\end{aligned}}\tag{AMT44}
\]
AMT29 displays every factor of both $\Omega$ values explicitly;
AMT32 gives the finite bounds
\[
\begin{aligned}
 0\le\Omega_k^{\rm low}
 &\le l\log[1+16(q-1+l)^2]+l\log[1+16(q+l)^2],\\
 0\le\Omega_k^{\rm high}
 &\le l\log[1+16(2q-1+l)^2]+l\log[1+16(2q+l)^2].
\end{aligned}\tag{AMT45}
\]
In particular both are $O(k\log(q+1))$ on the original
$m=1$ class. Its actual rank $r=8k-16$ therefore proves directly
\[
 \boxed{\Delta_K^\Gamma=O_h(k^2\log(q+1))=o(kq),\qquad
 \frac{F_K^{\rm ar}-F_K^\sigma}{kq}\longrightarrow0,\qquad
 \frac{F_K^{\rm ar}-F_K^0}{kq}\longrightarrow0.}
 \tag{AMT46}
\]
For the first equality divide the AMT44 bounds by $kq$ and use
$q=(k+1)^2$, so $k\log(q+1)/q\to0$. The second limit uses
AMT38, since $k^2/(kq)=k/q\to0$. The last uses the exact
identity $F_K^{\rm ar}-F_K^0=\Delta_K^\sigma+\Delta_K^\Gamma$.
All bounds are uniform for $|u|\ge R_*$.

For completeness the entire finite arithmetic boundary return is
also evaluated by those same quantities:
\[
\boxed{\begin{aligned}
 F_B^{\rm ar}-F_B^0
   &=\mathcal H_k+\delta_k^\sigma
                         -(F_K^{\rm ar}-F_K^0),\\
 \mathcal H_k+\delta_k^\sigma-r\widetilde E_k^-
   &\le F_B^{\rm ar}-F_B^0
             \le\mathcal H_k+\delta_k^\sigma+r\widetilde E_k^+.
\end{aligned}}\tag{AMT47}
\]
This uses AMT31 with its exact endpoint constants
$\widetilde E_k^+=E_k^++\Omega_k^{\rm low}$ and
$\widetilde E_k^-=E_k^-+\Omega_k^{\rm high}$, rather than a
new source comparison. Both tensor and arithmetic source masses
and the full original resultant have already been retained in
AMT23--30 before these signed cancellations.

The proved GEL18--19 constant is
\[
 \mathcal C_\Gamma=2\int_{a_*^2}^{b_*^2}\log x\,
                       \rho_{2,\pi}(x)\,dx-4(2\log2-1)>0,
 \qquad \frac{\mathfrak T_k}{lq}\longrightarrow\mathcal C_\Gamma.
 \tag{AMT48}
\]
Here $a_*,b_*$ and $\rho_{2,\pi}$ are exactly the proved GEL8/EIQ
equilibrium endpoints and density; their complete definitions and
integral proof remain in that provider. The endpoint letters here
avoid the already fixed original period parameter $u$. GEL18a
gives $4\log(27/(8\pi))\le\mathcal C_\Gamma\le6$ and proves
the strict positivity. TWA15 also proves the actual source return
$\delta_k^\sigma/(kq)\to0$. As $l/k\to1/4$, AMT44, AMT46
and AMT47 now prove
\[
\boxed{\begin{gathered}
 \frac{F_B^\sigma-F_B^0}{kq}\longrightarrow-\frac{\mathcal C_\Gamma}{4},
 \qquad
 \frac{F_B^{\rm ar}-F_B^0}{kq}\longrightarrow-\frac{\mathcal C_\Gamma}{4},
 \qquad
 \frac{F_B^{\rm ar}-F_B^\sigma}{kq}\longrightarrow0,\\
 \frac{F_B^\sigma-F_B^0}{\mathcal H_k}\longrightarrow1,
 \qquad
 \frac{F_B^{\rm ar}-F_B^0}{\mathcal H_k}\longrightarrow1,
 \quad m=1,\ k=4l+1,\ l\to\infty,\ |u|\ge R_*.
\end{gathered}}\tag{AMT49}
\]
Indeed the first two limits subtract an $o(kq)$ kernel return
from $\mathcal H_k=-\mathfrak T_k$, whose quotient by $kq$
tends to $-\mathcal C_\Gamma/4$. The third follows from
$\Delta_B^\sigma=\delta_k^\sigma-\Delta_K^\sigma$.
The nonzero limit of $\mathcal H_k/(kq)$ makes its denominator
nonzero for all sufficiently large original $k$ and proves the
last two ratios by division. The uniformity in $u$ follows from
the uniform kernel bounds and the period-independent total returns.
Thus the full existing source comparison evaluates the mixed
Gamma return also at its $kq$ scale, with its finite signed
remainders AMT44 and AMT47 still present.

% END COMPLETE AMT

\clearpage
% BEGIN COMPLETE KAF
% Complete original-source estimate; preceding editions are unchanged.
\section{Factorial control of the original common kernel at the \texorpdfstring{$kq$}{kq} scale}
\label{sec:KAF}

The exact denominator in OKV15 contains the information needed to improve
its absolute volume estimate. This proof keeps that denominator and also
allows every original primary multiplicity. The original Gamma sources,
root coordinates, coefficient maps and observation are retained. The
simple-quartet conclusion uses the actual period and rank calculations
PQO/RPC and OQC/OQX, with their complete proofs supplied as dependencies.

\subsection{Original coordinates, all primary multiplicities, and source norms}
Fix the original full-order quartet with $0<\delta<1/2$, $\gamma>2$,
and a positive integer primary order $m$. On the original tensor class set
\[
 k=4l+1,\quad l\ge1,\quad e_k=1+k(m-1),\quad
 q=e_k(k+1)^2,\quad c=k/2,\quad s\in\{1,k\},\quad b_s=s/2.
 \tag{KAF1}
\]
The original polynomial and its exact real-coordinate representative are
\[
\begin{split}
 \chi(S)&=\prod_{a,b=0}^k
 [S-c-(2a-k)\delta-i(2b-k)\gamma]^{e_k},\\
 \omega_{ab}&=(2b-k)\gamma-i(2a-k)\delta,\\
 P(y)&=i^{-q}\chi(c+iy)
       =\prod_{a,b=0}^k(y-\omega_{ab})^{e_k}.
\end{split}\tag{KAF2}
\]
In the full list $\omega_1,\ldots,\omega_q$, each root $\omega_{ab}$
occurs $e_k$ times. The exact maximal modulus remains
$R=k\sqrt{\delta^2+\gamma^2}$. Substitution gives inverse algebra maps
\[
 \Phi:\mathbb C[y]/(P)\longrightarrow E=\mathbb C[S]/(\chi),
 \quad[r]\longmapsto[r((S-c)/i)],\qquad
 \Phi^{-1}[f]=[f(c+iy)].\tag{KAF3}
\]
Indeed $P((S-c)/i)=i^{-q}\chi(S)$, and the two substitutions are
inverse before quotienting. Thus every root, multiplicity and factor of
$i$ in the original remainder operation is present.

Retain the complete measures and their monic orthogonal polynomials:
\[
\begin{gathered}
 M_s=(2\pi)^{s/2},\qquad
 dm_{0,s}(y)=\frac{M_s2^{b_s-2}}{\pi\Gamma(b_s)}
       \left|\Gamma\left(\frac{b_s}{2}+\frac{iy}{2}\right)\right|^2dy,\\
 p_0^{(s)}=1,\quad p_1^{(s)}=y,\qquad
 p_{d+1}^{(s)}=yp_d^{(s)}-d(b_s+d-1)p_{d-1}^{(s)},\\
 h_d^{(s)}=\int|p_d^{(s)}|^2dm_{0,s}=M_s d!(b_s)_d,\qquad
 \varphi_d^{(s)}=p_d^{(s)}/\sqrt{h_d^{(s)}},\qquad
 u_d^{(s)}(S)=\varphi_d^{(s)}((S-c)/i).
\end{gathered}\tag{KAF4}
\]
These are precisely the existing orders $1,k$. Their complete source
transform and norm calculation are IFB2--5 and OKV4--5. In particular
$\int e^{ty}dm_{0,s}=M_s(\cos t)^{-b_s}$ for $|t|<\pi/2$.
The generating series
$(1+z^2)^{-b_s/2}\exp(y\arctan z)$ has coefficients
$p_d^{(s)}z^d/d!$. Integrating its product at two sufficiently small
real variables gives $M_s(1-zw)^{-b_s}$, proving the displayed norms
and orthogonality by coefficient comparison. The original exponential
bound justifies those finite coefficient differentiations under the
integral. None of these source statements depends on $m$.

In the original polynomial frame of $E$ let $C_s$ have columns
$[u_0^{(s)}],\ldots,[u_{q-1}^{(s)}]$. Its determinant is exactly
\[
\begin{gathered}
 \det C_s=i^{-q(q-1)/2}\prod_{d=0}^{q-1}(h_d^{(s)})^{-1/2}\ne0,
 \\ v_j^{(s)}=C_s^{-1}[u_{q+j}^{(s)}],\quad
 K_j^{(s)}=I_q+\sum_{a=0}^{j-1}v_a^{(s)}(v_a^{(s)})^*,
 \quad0\le j\le q+1.
\end{gathered}\tag{KAF5}
\]
The last included vector has original degree $2q$.

\subsection{A degree-dependent constant in the unchanged remainder estimate}
Define the actual finite constants
\[
 A_{k,m}=\max\left\{2\sqrt{1+\frac{k}{2q}},\frac{R}{q}\right\},
 \qquad T=qA_{k,m},\qquad
 \mathcal A_{k,m}=6\log2+4+4\log A_{k,m}.
 \tag{KAF6}
\]
The letter $A_{k,m}$ in this section is a coefficient-estimate constant;
it is not the arithmetic upper source constant $A_k$ of TWA/AMT.
In particular $A_{k,m}\ge2$ and $R\le T$. For fixed original quartet
and fixed $m$, $q\ge(k+1)^2$ implies
\[
 A_{k,m}\longrightarrow2,\qquad
 \mathcal A_{k,m}\longrightarrow4+10\log2.
 \tag{KAF7}
\]
Indeed $k/q\to0$ and $R/q=\sqrt{\delta^2+\gamma^2}\,k/q\to0$.
The root-dependent term in the maximum is retained at every finite $k$.

Here is the complete remainder estimate needed below. For the full
root list define elementary and complete homogeneous polynomials
$e_h(\omega)$ and $h_t(\omega)$ by their finite sums, with
$e_0=h_0=1$ and $e_h=0$ outside $0\le h\le q$. The formal identity
\[
 \left(\sum_{h=0}^q(-1)^he_h(\omega)z^h\right)
 \left(\sum_{t\ge0}h_t(\omega)z^t\right)=1
 \tag{KAF8}
\]
follows by multiplying each of the $q$ factors $1-\omega_a z$
by its geometric series. Repeated roots do not affect this identity.
It proves, for $D=q+j$ and $0\le j\le q$, the exact signed formula
\[
 [y^a]\operatorname{rem}_P(y^{q+j})
 =-\sum_{t=0}^j h_t(\omega)(-1)^{q+j-t-a}e_{q+j-t-a}(\omega),
 \qquad0\le a<q.\tag{KAF9}
\]
In detail the quotient is $\sum_{t=0}^jh_t(\omega)y^{j-t}$;
KAF8 cancels all coefficients of degree at least $q$ in its product
with $P$, and the remaining coefficients are KAF9. This proves the
formula by uniqueness of monic division, including all multiplicities.

The defining sums give
$|e_h|\le\binom qh R^h$ and
$|h_t|\le\binom{q+t-1}{t}R^t$. In every nonzero summand of KAF9
the total root exponent is $q+j-a$. The binomial coefficients of
$h_t$ increase with $t$, and the distinct indices in the $e_h$ sum
contribute at most $2^q$. Hence
\[
 |[y^a]\operatorname{rem}_P(y^{q+j})|
 \le2^q\binom{q+j-1}{j}R^{q+j-a}
 \le2^{3q-1}R^{q+j-a}.\tag{KAF10}
\]
The last bound is the binomial expansion at one, since $j\le q$.

The exact orthonormal recurrence from KAF4 has upward and downward
weights $\sqrt{(a+1)(b_s+a)}$ and $\sqrt{a(b_s+a-1)}$.
On degree at most $q-2$ each shift has norm at most
$\sqrt{q(b_s+q)}$. Because $b_s\le k/2$, their sum has norm at
most $2\sqrt{q(k/2+q)}\le T$. Iterating from the constant, with all
intermediate degrees in that range, proves
\[
 \|y^a\|_{L^2(m_{0,s})}\le\sqrt{M_s}\,T^a,
 \qquad0\le a<q.\tag{KAF11}
\]
For the finite coefficient sum $c_T(f)=\sum_a|[y^a]f|T^a$ the
original recurrence gives
$c_T(p_{d+1}^{(s)})\le T c_T(p_d^{(s)})+
d(b_s+d-1)c_T(p_{d-1}^{(s)})$.
For $1\le d<2q$,
\[
 d(b_s+d-1)\le2q(k/2+2q)=4q^2+kq
       \le4q^2+2kq\le T^2.
\]
Starting with $c_T(p_0)=1,c_T(p_1)=T$, induction therefore proves
$c_T(p_d)\le(2T)^d$ for $0\le d\le2q$: the recurrence bound is
$(3/4)(2T)^{d+1}\le(2T)^{d+1}$.
Combining KAF10 with all $q$ monomial norms KAF11, then using the
linearity of the original remainder map, gives
\[
 \|\operatorname{rem}_P(p_d^{(s)})\|_{L^2(m_{0,s})}
       \le\sqrt{M_s}\,q2^{3q}(2T)^d,
       \qquad0\le d\le2q.\tag{KAF12}
\]
For a monomial of degree $D\ge q$, each coefficient term is at
most $\sqrt{M_s}2^{3q-1}R^{D-a}T^a\le
\sqrt{M_s}2^{3q-1}T^D$. For $D<q$ use KAF11 directly.
These two cases and the proved coefficient sum establish KAF12
without deleting a remainder coefficient or changing the source norm.

\subsection{Retaining the factorial denominator}
Expansion of the degree-below-$q$ remainder in the orthonormal
polynomials of KAF4, followed by the exact map KAF3, proves
\[
 \|v_j^{(s)}\|_2^2
 =\frac{\|\operatorname{rem}_P(p_{q+j}^{(s)})\|_{L^2(m_{0,s})}^2}
          {M_s(q+j)!(b_s)_{q+j}}
 \le\frac{q^22^{6q}(2T)^{2(q+j)}}{(q+j)!(b_s)_{q+j}}.
 \tag{KAF13}
\]
The factor $M_s$ occurs in the original numerator and denominator
before this cancellation. For every $d\ge0$,
\[
 d!(b_s)_d\ge d!(1/2)_d=\frac{(2d)!}{4^d},
 \qquad
 \|v_j^{(s)}\|_2^2\le q^22^{6q}\frac{(4T)^{2(q+j)}}{(2q+2j)!}.
 \tag{KAF14}
\]
The inequality is factor by factor because $b_s\ge1/2$.
The equality follows by separating even and odd factors of $(2d)!$.
This is precisely the denominator whose further retention sharpens
OKV16--17.

Define the positive exact quantities
\[
 \eta_{k,m}=\frac{1-1/(4q)}{A_{k,m}^2}<\frac14,\qquad
 Z_{k,m}=1+q^22^{6q}\frac{(4qA_{k,m})^{4q}}{(4q)!}
              \frac{1-\eta_{k,m}^{q+1}}{1-\eta_{k,m}}.
 \tag{KAF15}
\]
For $t_d=(4T)^{2d}/(2d)!$, $q+1\le d\le2q$ gives
$t_{d-1}/t_d=(2d)(2d-1)/(16T^2)\le\eta_{k,m}$.
Consequently $\sum_{d=q}^{2q}t_d\le
t_{2q}\sum_{a=0}^q\eta_{k,m}^a$. Adding the $q+1$ actual rank-one
updates in KAF5 and using KAF14 proves the finite operator bound
\[
 \boxed{1\le\|K_{q+1}^{(s)}\|_2\le Z_{k,m},\qquad s\in\{1,k\}.}
 \tag{KAF16}
\]
Here $\|vv^*\|_2=\|v\|_2^2$ follows by Cauchy--Schwarz with
equality on $v\ne0$; a zero vector contributes zero.

For an explicit linear bound, monotonicity of $\log x$ on each
interval $[j-1,j]$ gives
\[
 \log n!=\sum_{j=1}^n\log j
       \ge\int_1^n\log x\,dx=n\log n-n+1.
\]
Thus $(4qA_{k,m})^{4q}/(4q)!\le
\mathrm e^{-1}(\mathrm e A_{k,m})^{4q}$, where $\mathrm e$ is
the exponential constant and differs from the primary integer $e_k$.
The geometric sum in KAF15 is at most $4/3$. Therefore
$Z_{k,m}\le1+[4/(3\mathrm e)]X$, with
$X=q^2\exp(q\mathcal A_{k,m})$. The positive exponential series
gives $\mathrm e>1+1+1/2+1/6=8/3$, so $4/(3\mathrm e)<1/2$.
Since $A_{k,m}\ge2$ and $q\ge36$, $X>2$, and hence
$1+[4/(3\mathrm e)]X\le X$. We have proved
\[
 \boxed{\log\|K_{q+1}^{(s)}\|_2\le\log Z_{k,m}
       \le q\mathcal A_{k,m}+2\log q.}
 \tag{KAF17}
\]
The full factorial and finite geometric sum in KAF15 remain available
as the sharper finite bound.

\subsection{All four actual kernel and boundary Grams}
Let $\Lambda:E\twoheadrightarrow B$ be the original period and full-unit
observation, let $I:\mathcal K\hookrightarrow E$ be its actual kernel
inclusion, and put $r=\dim\mathcal K$. No rank formula is needed in
this subsection, so it applies to every original $m$. The original
quotient and restricted Grams are
\[
 G_{q+j-1}^{(s)}=C_s^{-*}(K_j^{(s)})^{-1}C_s^{-1},\qquad
 T_s=C_s^{-1}I,\qquad H_j^{(s)}=T_s^*(K_j^{(s)})^{-1}T_s.
 \tag{KAF18}
\]
Indeed the original orthonormal source-coordinate remainder matrix is
$[I_q,v_0,\ldots,v_{j-1}]$. Multiplication by its adjoint gives $K_j$;
its unique minimum lift of $x$ is its adjoint applied to $K_j^{-1}x$,
with squared norm $x^*K_j^{-1}x$. This proves both matrices KAF18
in the original coefficient frame. The positive update chain implies
\[
 Z_{k,m}^{-1}T_s^*T_s\preceq H_j^{(s)}\preceq T_s^*T_s,
 \qquad H_{j+1}^{(s)}\preceq H_j^{(s)}.\tag{KAF19}
\]
Inversion reverses positive order by congruence and inversion of
positive eigenvalues. Restricting this order with the same injective
$T_s$ proves KAF19. After congruence with $(T_s^*T_s)^{-1/2}$,
every eigenvalue lies between $Z_{k,m}^{-1}$ and one. Thus
$-r\log Z_{k,m}\le\log\det H_j^{(s)}-
\log\det(T_s^*T_s)\le0$.

The exact four original degrees are $q-1,q,2q-1,2q$. With their
unchanged signs define
\[
 F_K^{(s)}=\log\det H_0^{(s)}+\log\det H_1^{(s)}
              -\log\det H_q^{(s)}-\log\det H_{q+1}^{(s)}.
\]
Pairing degrees $q-1$ with $2q-1$, and $q$ with $2q$, proves
nonnegativity by KAF19. Each difference is at most $r\log Z_{k,m}$.
The four original frame factors cancel only after these four signs
have been recorded. It follows that
\[
 \boxed{0\le F_K^{(s)}\le2r\log Z_{k,m}
       \le2r\,[q\mathcal A_{k,m}+2\log q],\qquad s\in\{1,k\}.}
 \tag{KAF20}
\]
If $r=0$, every kernel determinant is one and the statement has its
same empty-space meaning.

At each original degree the boundary Gram is
$Q_j^{(s)}=(\Lambda(G_{q+j-1}^{(s)})^{-1}\Lambda^*)^{-1}$.
For any fixed section $S_*$ of $\Lambda$, its minimum section
$S_j^{(s)}=(G_{q+j-1}^{(s)})^{-1}\Lambda^*Q_j^{(s)}$ differs from
$S_*$ by a kernel column and is orthogonal to $I$ in that metric.
The change of frame has determinant one. Taking determinants in this
orthogonal frame proves, with its original coefficient factor,
\[
 \det H_j^{(s)}\det Q_j^{(s)}
       =|\det[I,S_*]|^2\det G_{q+j-1}^{(s)}.
\]
The same four signs now prove the exact splitting
$F_K^{(s)}+F_B^{(s)}=\mathcal B_k^{(s)}$. Boundary minimum norms
decrease with degree, so $F_B^{(s)}\ge0$. In particular
\[
 \boxed{\max\{0,\mathcal B_k^{(s)}-2r\log Z_{k,m}\}
       \le F_B^{(s)}\le\mathcal B_k^{(s)}.}\tag{KAF21}
\]
These are statements about the original kernel and boundary, with
all multiplicities and the same observation throughout.

\subsection{The actual simple quartet and the arithmetic source}
For $m=1$, the complete RPC calculation supplies an explicit finite
$R_*$ from the original unit and quartet. On every logarithm branch
and for $|u|\ge R_*$ its original period quadric has five distinct
cyclic orbit members. The full invariant restriction ideal calculation
OQX then gives $r=8k-16$. Substitution in the finite KAF20 gives
\[
 \boxed{0\le\frac{F_K^{(s)}}{kq}
 \le\left(16-\frac{32}{k}\right)
       \left(\mathcal A_{k,1}+\frac{2\log q}{q}\right),
 \qquad s\in\{1,k\},\quad |u|\ge R_*.}\tag{KAF22}
\]
The original arithmetic-source comparison AMT17--22 proves
$-rE_k^+\le F_K^{\rm ar}-F_K^{(1)}\le rE_k^-$, where the
unchanged full constants are
\[
\begin{gathered}
 E_k^+=2\log(A_k/a_k)+\log D_{q-1}+\log D_q,\\
 E_k^-=2\log(A_k/a_k)+\log D_{2q-1}+\log D_{2q},
 \\ D_N=[1+4(N+21+4m)^2]^{21+4m}.
\end{gathered}
\]
Here $a_k,A_k$ are exactly the original arithmetic constants of AMT4,
whose full source inequality and minimum-section proof are included
there. They retain the actual numerator $2\xi/h$ and all source
masses; AMT22 proves $E_k^\pm=O_h(k+\log q)$.
Thus the actual arithmetic metric satisfies the finite bound
\[
 \boxed{0\le F_K^{\rm ar}
       \le(16k-32)\log Z_{k,1}+(8k-16)E_k^-.}\tag{KAF23}
\]
Combining KAF7, KAF22--23, and $q=(k+1)^2$ proves uniformly on
the entire stated original period domain, for
$a\in\{\sigma,0,\mathrm{ar}\}$,
\[
 \boxed{0\le\liminf_{k\to\infty}\frac{F_K^a}{kq}
       \le\limsup_{k\to\infty}\frac{F_K^a}{kq}
       \le16(4+10\log2).}\tag{KAF24}
\]
For precision, uniformity means the same finite upper bound holds for
every such $u$ and branch at each original $k$; therefore it also
bounds every sequence of those parameters. The arithmetic extra term
divided by $kq$ is at most $(8-16/k)E_k^-/q\to0$.
AMT46 further proves that the three ratios differ by quantities
tending uniformly to zero. Thus their limit points coincide for any
fixed allowed sequence of period parameters; KAF24 does not assert
that this bounded sequence already has a unique limit.

Finally the exact arithmetic identity gives
\[
 \max\{0,\mathcal B_k^{\rm ar}-(16k-32)\log Z_{k,1}
                              -(8k-16)E_k^-\}
       \le F_B^{\rm ar}\le\mathcal B_k^{\rm ar}.\tag{KAF25}
\]
Its full-source value remains
$\mathcal B_k^{\rm ar}=\mathcal B_k^\sigma+\delta_k^\sigma
=\mathcal B_k^0+\mathcal H_k+\delta_k^\sigma$.
All exact Gamma and arithmetic signed returns of AMT44--49 are
retained. The new calculation improves the absolute common kernel
bound to the $kq$ scale, while its exact four-endpoint determinant
and actual period dependence remain the objects for the next leading
term calculation.

% END COMPLETE KAF

\clearpage
% BEGIN COMPLETE OCF
% Original coefficient reduction. All maps are complex-linear until a metric is supplied.
\section{An exact conductor frame and a fixed invariant-generator recursion}
\label{sec:OCF}

Retain the original simple-quartet data, $m=1$, $k=4l+1$,
$l\geq1$, the original root order $(++),(+-),(-+),(--)$, and
the complete original period and amplitude matrix of OQC1:
\[
                  \mathsf H=F^{-1}\Pi U\mathcal E_{\rm prim}.
\]
Thus $\det\mathsf H\ne0$ and
$Q(x)=Q_0(\mathsf H^{-1}x)$, where
$Q_0(t)=t_{++}t_{--}-t_{+-}t_{-+}$.
Use $g f(x)=f(D^{-1}x)$,
$D=\operatorname{diag}(\zeta,\zeta^2,\zeta^3,\zeta^4)$,
$\zeta=e^{2\pi i/5}$, and the original invariant algebra
$\mathscr I=\mathbb C[x_1,x_2,x_3,x_4]^{\langle g\rangle}$.
This calculation applies on the actual five-orbit period locus,
which includes every original logarithm branch with $|u|\geq R_*$
by PQO and RPC. On that locus the five prime polynomials
$Q_a=g^aQ$, $0\leq a<5$, are pairwise nonassociate. Keep their
actual products
\[
 \mathcal A=\prod_{a=1}^4Q_a,\qquad
 \mathcal H_Q=Q\mathcal A=\prod_{a=0}^4Q_a,
 \qquad \deg\mathcal A=8,\quad\deg\mathcal H_Q=10.
 \tag{OCF1}
\]
No coefficient in these polynomials is replaced. The exact orbit
decision and the period domain are those of the cited original
calculations. The construction below treats every possible nonzero
leading coefficient of their actual restriction.

\subsection{The original Segre coefficient isomorphism}

For $d\geq0$ let $\mathcal B_d$ be the space of biforms of
bidegree $(d,d)$ in $(z_0,z_1)$ and $(w_0,w_1)$. Its ordered basis is
\[
 e_{ab}^{(d)}=z_0^a z_1^{d-a}w_0^b w_1^{d-b},
 \qquad 0\leq a,b\leq d.
 \tag{OCF2}
\]
Pairs are ordered lexicographically, first by $a$, then by $b$.
The order is used only for the exact coefficient elimination.
Set $\mathcal B_d=0$ for $d<0$. Multiplication satisfies
$e_{ab}^{(d)}e_{ij}^{(e)}=e_{a+i,b+j}^{(d+e)}$.
The original substitution is
\[
 \begin{gathered}
 t=(z_0w_0,z_0w_1,z_1w_0,z_1w_1)^T,\qquad x=\mathsf Ht,\\
 \mathscr P_d=\mathbb C[x_1,x_2,x_3,x_4]_d,\qquad
 \psi_d:\mathscr P_d\longrightarrow\mathcal B_d,\quad
 f\longmapsto f(\mathsf Ht).
 \end{gathered}\tag{OCF3}
\]
The induced map
$R_d=\mathscr P_d/(Q\mathscr P_{d-2})\longrightarrow\mathcal B_d$
is an isomorphism. Here is a coefficient proof with the given
orientation. For each $(a,b)$, choose an integer
$h$ between $\max(0,a+b-d)$ and $\min(a,b)$. The monomial
$t_{++}^{h}t_{+-}^{a-h}t_{-+}^{b-h}t_{--}^{d-a-b+h}$ maps
to $e_{ab}^{(d)}$, proving surjectivity. Two monomials with
the same pair differ in their exponent vectors by an integer
multiple of $(1,-1,-1,1)$. After removing their common monomial,
their difference is divisible by $t_{++}t_{--}-t_{+-}t_{-+}$,
using $A^n-B^n=(A-B)\sum_{j=0}^{n-1}A^{n-1-j}B^j$.
A polynomial in the kernel has coefficient sum zero within every
pair, so it is a sum of such differences. The reverse inclusion
in the kernel follows by direct substitution. The invertible
substitution $x=\mathsf Ht$ then gives exactly the asserted kernel.
Thus these maps preserve multiplication as well as every grading.

Write the full conductor restriction as
\[
 A=\psi_8(\mathcal A)
     =\sum_{r=0}^8\sum_{s=0}^8 a_{rs}e_{rs}^{(8)}\ne0.
 \tag{OCF4}
\]
To verify its nonvanishing, $Q$ is a rank-four irreducible quadric,
hence prime in the polynomial ring, and none of $Q_1,\ldots,Q_4$
is its associate. Therefore $Q$ does not divide their product.
The kernel calculation OCF3 proves OCF4. The ring
$\mathcal B=\bigoplus_{d\geq0}\mathcal B_d$, with the indicated
diagonal grading, embeds in the four-variable polynomial ring;
it is an integral domain. Multiplication by $A$ is consequently
injective on every $\mathcal B_d$.

\subsection{Every leading-coefficient case gives an exact strip frame}

Let $(r_*,s_*)$ be the lexicographically largest pair for which
$a_{r_*s_*}\ne0$, and put $a_*=a_{r_*s_*}$. This is an
actual nonzero coefficient of OCF4. For $d\geq8$ define
\[
 \begin{split}
 P_d&=\{(r_*+i,s_*+j):0\leq i,j\leq d-8\},\\
 E_d^{\rm strip}&=\{0,\ldots,d\}^2\setminus P_d,\qquad
 W_d=\mathbb C^{E_d^{\rm strip}}.
 \end{split}\tag{OCF5}
\]
For $0\leq d<8$ put $P_d=\varnothing$,
$E_d^{\rm strip}=\{0,\ldots,d\}^2$ and use the same definition
of $W_d$; for $d<0$ put $W_d=0$. Let
$S_d:W_d\hookrightarrow\mathbb C^{(d+1)^2}$ insert zeros in
the positions $P_d$. All coefficient columns below use OCF2.

The following finite elimination defines
$N_d:\mathbb C^{(d+1)^2}\to W_d$ and, for $d\geq8$,
$V_d:\mathbb C^{(d+1)^2}\to\mathbb C^{(d-7)^2}$.
Starting with the complete input biform $f$, traverse all
$(i,j)\in\{0,\ldots,d-8\}^2$ in decreasing lexicographic order.
At the step $(i,j)$, let $c$ be the current coefficient at
$(r_*+i,s_*+j)$, record $b_{ij}=c/a_*$, and replace the current
biform by
\[
       f_{\rm current}-b_{ij}A e_{ij}^{(d-8)}.
 \tag{OCF6}
\]
At completion take the coefficients on $E_d^{\rm strip}$ as
$N_df$ and take $(b_{ij})$ as $V_df$. If $d<8$, $N_d$ is the
identity and $V_d$ is the zero map into the zero space.
The exact multiplication matrix $M_{A,d}$ gives
\[
 S_dN_d+M_{A,d}V_d=I,\qquad
 N_dS_d=I_{W_d},\qquad
 \ker N_d=A\mathcal B_{d-8}.
 \tag{OCF7}
\]
To prove these identities, the largest term of
$A e_{ij}^{(d-8)}$ is $a_*e_{r_*+i,s_*+j}^{(d)}$;
every other term is strictly smaller in the same lexicographic
order. Therefore each step zeros its designated position without
altering an earlier zero. Summing all subtractions gives the
first identity. Input supported on the strip produces zero
coefficients at every elimination step, giving the second.
Finally, the leading term of $Ab$, for a nonzero homogeneous
biform $b$, is the product of the two leading terms and belongs
to $P_d$. Indeed every other product is smaller in the additive
lexicographic order. Thus
$A\mathcal B_{d-8}\cap S_dW_d=0$. The first identity now proves
the kernel formula, the uniqueness of the decomposition, and all
three identities in OCF7. All operations are linear in the input;
their coefficients are explicit rational expressions in the
original $a_{rs}$ with powers of the displayed $a_*$ as denominators.
The polynomial $A$ itself remains the full polynomial OCF4.

In particular the quotient map induced by $N_d$ is the exact
isomorphism
\[
 \mathcal B_d/(A\mathcal B_{d-8})\xrightarrow{\ \sim\ }W_d,
 \qquad
 n_d:=\dim W_d=
 \begin{cases}(d+1)^2,&0\leq d<8,\\16d-48,&d\geq8.
 \end{cases}
 \tag{OCF8}
\]
The second formula is
$(d+1)^2-(d-7)^2=16d-48$. The possible choices
$(r_*,s_*)\in\{0,\ldots,8\}^2$ form an exhaustive finite
partition: the chart for a pair specifies that its coefficient
is nonzero and every larger coefficient is zero. Hence all actual
coefficient cases have been included. At the first original
degree $k=5$ this construction retains all $36$ coordinates;
at every original $k\geq9$ it retains $16k-48$ coordinates.

\subsection{Twenty-five fixed invariant monomials generate all degrees}

Define the following fixed finite sets of exponent vectors:
\[
 G_e=\{\nu\in\mathbb N^4:|\nu|=e,
                   \nu_1+2\nu_2+3\nu_3+4\nu_4\equiv0\pmod5\},
 \quad 1\leq e\leq5,\qquad G=\bigcup_{e=1}^5G_e.
 \tag{OCF9}
\]
Here $G_1$ is empty. The complete remaining lists, in the original
$x_1,x_2,x_3,x_4$ order, are
\[
\begin{array}{c|l}
e&G_e\\\hline
2&(0,1,1,0),(1,0,0,1)\\
3&(0,0,2,1),(0,1,0,2),(1,2,0,0),(2,0,1,0)\\
4&(0,0,1,3),(0,2,2,0),(0,3,0,1),(1,0,3,0),\\
 & (1,1,1,1),(2,0,0,2),(3,1,0,0)\\
5&(0,0,0,5),(0,0,5,0),(0,1,3,1),(0,2,1,2),\\
 & (0,5,0,0),(1,0,2,2),(1,1,0,3),(1,3,1,0),\\
 & (2,1,2,0),(2,2,0,1),(3,0,1,1),(5,0,0,0).
\end{array}\tag{OCF10}
\]
Thus $(|G_1|,\ldots,|G_5|)=(0,2,4,7,12)$ and $|G|=25$.
These lists can also be obtained directly by choosing
$0\leq\nu_1,\nu_2,\nu_3\leq e$ with their sum at most $e$,
putting $\nu_4=e-\nu_1-\nu_2-\nu_3$, and applying the single
displayed congruence, which verifies completeness of each list.

Every positive-degree invariant monomial is a product of members
$x^\nu$, $\nu\in G$. To prove this, list the weights of its
individual factors in any order. If its degree is at most five,
the entire list gives a member of $G$. If its degree exceeds five,
consider the six partial sums of the first five weights, including
the empty partial sum, in $\mathbb Z/5\mathbb Z$. Two coincide;
the nonempty block between them has length at most five and total
weight zero. Its factors give a divisor $x^\nu$ with $\nu\in G$.
The complementary monomial again has total weight zero and smaller
degree. Induction completes the factorization. Because the action
is diagonal, a polynomial is invariant exactly when all its
nonzero monomials have weight zero: comparison of each monomial
coefficient in $gf=f$ proves this statement. Thus
\[
                \mathscr I=\mathbb C[x^\nu:\nu\in G].
 \tag{OCF11}
\]
All intermediate degrees here are degrees in the original
coefficient algebra. This construction introduces no new source
distribution and leaves the original source orders $1$ and $k$ intact.

For each $\nu\in G_e$ the complete generator biform is
\[
 f_\nu=\psi_e(x^\nu)
  =\prod_{j=1}^4
   (H_{j1}z_0w_0+H_{j2}z_0w_1+
                  H_{j3}z_1w_0+H_{j4}z_1w_1)^{\nu_j}
  =\sum_{a,b=0}^e c_{\nu,ab}e_{ab}^{(e)}.
 \tag{OCF12}
\]
For an entirely explicit coefficient formula, sum over all
nonnegative $4$ by $4$ integer arrays $(n_{jt})$ satisfying
$\sum_t n_{jt}=\nu_j$,
$\sum_j(n_{j1}+n_{j2})=a$ and
$\sum_j(n_{j1}+n_{j3})=b$. Then
\[
 c_{\nu,ab}=
 \sum_{(n_{jt})}
     \prod_{j=1}^4\left(
        \frac{\nu_j!}{\prod_{t=1}^4n_{jt}!}
        \prod_{t=1}^4H_{jt}^{n_{jt}}\right).
 \tag{OCF13}
\]
The multinomial theorem proves the formula, including every
original period and amplitude entry. Every factor has degree at
most five and at most $36$ possible biform coefficients.

\subsection{Exact multiplication on the strip and the invariant recursion}

Let $\mathcal L=\mathcal B/(A\mathcal B)$, a graded algebra.
For $\nu\in G_e$ and $d\geq e$ define
\[
 m_{\nu,d}:W_{d-e}\longrightarrow W_d,
 \qquad m_{\nu,d}=N_d M_{f_\nu}S_{d-e}.
 \tag{OCF14}
\]
Here $M_{f_\nu}$ is the exact coefficient convolution:
\[
       (M_{f_\nu}S_{d-e}v)_{ab}
          =\sum_{i,j}c_{\nu,a-i,b-j}(S_{d-e}v)_{ij},
\]
with out-of-range indices giving zero.
Therefore every entry of OCF14 is supplied
by the finite formulas OCF6 and OCF13.
For a representative $h\in\mathcal B_{d-e}$, OCF7 writes
$h-S_{d-e}N_{d-e}h=A b$. Multiplication by $f_\nu$ preserves
this ideal. Applying $N_d$ yields the exact identity
\[
       N_d(f_\nu h)=m_{\nu,d}N_{d-e}h.
 \tag{OCF15}
\]
This proves that OCF14 is multiplication by the specified element
of $\mathcal L$, including its independence of representatives.

The original invariant image is
$J_d=\psi_d(\mathscr I_d)\subseteq\mathcal B_d$.
Set $\overline J_d=N_dJ_d\subseteq W_d$, and set
$\overline J_d=0$ for $d<0$. The exact recursion is
\[
 \overline J_0=\mathbb C\cdot1,\qquad
 \overline J_d=
   \sum_{\substack{\nu\in G_e\\2\leq e\leq5,\ e\leq d}}
                 m_{\nu,d}\overline J_{d-e}
                 \quad(d\geq1).
 \tag{OCF16}
\]
Indeed OCF11 factors every invariant monomial of degree $d>0$
as $x^\nu$ times an invariant monomial of degree $d-e$.
Applying $\psi$, then OCF15, gives the inclusion from left to
right. Conversely every summand is the image of a product of two
invariant polynomials, so it belongs to the left side. Linearity
proves equality for the complete invariant space. In particular
the recursion gives $\overline J_1=0$ and includes all small degrees.

The conductor ideal is contained in this invariant image with
its original trace coefficient:
\[
 A\mathcal B_{d-8}\subseteq J_d,\qquad
 [5\mathsf R(\mathcal A h)]_{(Q)}=[\mathcal A h]_{(Q)},
 \quad\mathsf R=\frac15\sum_{a=0}^4g^a.
 \tag{OCF17}
\]
To prove the identity, $g^a\mathcal A$ is the product of all
five orbit equations except $Q_a$; for $a\ne0$ it contains $Q$.
Thus all four corresponding terms vanish modulo $Q$, and the
$a=0$ term is the displayed product. Every biform in
$\mathcal B_{d-8}$ has a polynomial lift by OCF3, so the identity
proves the inclusion. It also shows exactly why the quotient
used in OCF16 retains the original invariant image.

\subsection{Exact ranks, a finite pivot atlas, and the original quotient map}

For $d\geq0$ put
\[
 d_0(d)=\frac{\binom{d+3}{3}+4\epsilon_d}{5},\qquad
 \epsilon_d=\mathbf1_{d\equiv0\ (5)}-\mathbf1_{d\equiv1\ (5)},
 \qquad d_0(d)=0\quad(d<0).
 \tag{OCF18}
\]
This is the number of degree-$d$ weight-zero monomials: the
identity character has generating function $(1-z)^{-4}$, each
of the four nonidentity characters has generating function
$\prod_{j=1}^4(1-z\zeta^j)^{-1}=(1-z)/(1-z^5)$, and averaging
their coefficients gives OCF18.
Every invariant polynomial divisible by $Q$ is divisible by all
five $Q_a$, hence by $\mathcal H_Q$, since they are pairwise
nonassociate primes. Its quotient by this product is invariant:
apply $g$ and cancel the nonzero invariant product in the integral
domain. Conversely each $\mathcal H_Q$ times an invariant
polynomial is invariant and vanishes modulo $Q$. Consequently
$\ker(\mathscr I_d\to J_d)=\mathcal H_Q\mathscr I_{d-10}$,
and injection of multiplication by the nonzero product gives
$\dim J_d=d_0(d)-d_0(d-10)$.
Injection of multiplication by $A$, OCF7 and OCF17 then prove
\[
 j_d:=\dim\overline J_d
     =d_0(d)-d_0(d-10)-\max(d-7,0)^2.
 \tag{OCF19}
\]
In particular $j_0=1$, $j_1=0$, $j_5=12$, and
$j_d=8d-32$ for $d\geq8$. For the last formula at $d\geq10$,
$\epsilon_d=\epsilon_{d-10}$ and direct expansion gives
$\binom{d+3}{3}-\binom{d-7}{3}=5(d^2-6d+17)$;
subtracting $(d-7)^2$ gives $8d-32$. At $d=8,9$ the ranks
before this subtraction are $33,44$, respectively, which gives
$32,40$ and verifies both remaining cases.

Here is a complete finite construction of all the columns used
by OCF16. Start with $Z_0=(1)$ and the invariant polynomial
certificate $F_{0,1}=1$. Suppose all earlier matrices $Z_b$ have
independent columns spanning $\overline J_b$, and their columns
are $N_b\psi_b(F_{b,h})$ for invariant polynomials $F_{b,h}$.
Form the candidate matrix
\[
 C_d=\bigl[\,m_{\nu,d}Z_{d-e}\,\bigr]_{
                    \nu\in G_e,\ 2\leq e\leq5,\ e\leq d}.
 \tag{OCF20}
\]
Each candidate column has the explicit invariant certificate
$x^\nu F_{d-e,h}$. OCF16 proves that $C_d$ has column image
$\overline J_d$ and rank $j_d$. Select $j_d$ columns and $j_d$
rows giving a nonzero determinant. Let $Z_d$ be those columns,
let $I_d$ be those row positions, and let $T_d$ be the complementary
row positions in $E_d^{\rm strip}$. Retain the corresponding
invariant certificates. This completes the induction. When
$j_d=0$, use the empty matrix with zero columns, empty $I_d$,
and $T_d=E_d^{\rm strip}$.

Every selected determinant and every matrix entry is an explicit
expression obtained from the original entries of $\mathsf H$ and
the full $a_{rs}$. To cover all cases, order the finite choices
of column and row subsets and select the first nonzero determinant;
the corresponding chart asserts its nonvanishing and the vanishing
of all preceding candidates. A nonzero determinant exists because
the rank was proved in OCF19. At rank zero no determinant is needed.
Together with the exhaustive leading-coefficient charts after
OCF8, this gives a finite atlas for each specified $k$. It requires
no generic coefficient premise and supplies no numerical zero test
for an unevaluated analytic coefficient: each chart is specified
by its exact minors.

Define $\rho_d:W_d\to\mathbb C^{T_d}$ and the full original
coefficient map $\pi_d$ by
\[
 \rho_d=\operatorname{row}_{T_d}
            -(Z_d)_{T_d}(Z_d)_{I_d}^{-1}
                                  \operatorname{row}_{I_d},
 \qquad \pi_d=\rho_dN_d.
 \tag{OCF21}
\]
At rank zero the second term is the zero map. If $E_{T_d}$
inserts a column into its $T_d$ positions in $W_d$, then
\[
 \begin{split}
 w&=Z_d(Z_d)_{I_d}^{-1}w_{I_d}+E_{T_d}\rho_dw,\\
 \rho_dE_{T_d}&=I,\qquad\ker\rho_d=\operatorname{im}Z_d
                                  =\overline J_d,\\
 \ker\pi_d&=J_d,\qquad
       \pi_d S_d E_{T_d}=I_{\mathbb C^{T_d}}.
 \end{split}\tag{OCF22}
\]
For the first line, the subtraction on its right leaves zero
$I_d$ rows and exactly $\rho_dw$ on the complementary rows.
This proves the first three claims. If $\pi_df=0$, then
$N_df\in\overline J_d$; choose $j\in J_d$ with $N_dj=N_df$.
Now $f-j\in\ker N_d=A\mathcal B_{d-8}\subseteq J_d$ by
OCF17, so $f\in J_d$. The reverse inclusion follows immediately
from $N_dJ_d=\overline J_d$. The last equality uses $N_dS_d=I$.
This proves every map and kernel in OCF22.

At the original degrees, these are the exact dimensions
\[
 \begin{array}{c|c|c|c}
 k&n_k&j_k&|T_k|=n_k-j_k\\\hline
 5&36&12&24\\
 k\geq9,\ k=4l+1&16k-48&8k-32&8k-16.
 \end{array}\tag{OCF23}
\]
The original map $\Theta_k:R_k\xrightarrow{\sim}E_k^\vee$
in OQC7 sends the coefficient basis OCF2 to the original dual
primary basis in the same order. Indeed its evaluation on
$e^{wY_k}[1]$ is
$\exp(((2b-k)\gamma-i(2a-k)\delta)w)$. The corresponding
eigenvalue of $A_k$ is
$\beta_{ab}=k/2+(2a-k)\delta+i(2b-k)\gamma$; the eigenvalue
of $Y_k=(A_k-kI/2)/i$ is
$\omega_{ab}=(\beta_{ab}-k/2)/i=(2b-k)\gamma-i(2a-k)\delta$.
These eigenvalues are distinct because
$\delta>0$ and $\gamma>0$. Their exponentials are independent
by the Vandermonde determinant, so the evaluation identifies the
covectors exactly. By the original observation identity OQC8,
$\Theta_kJ_k=\operatorname{im}\Lambda_k^\vee$. Therefore OCF22
supplies the specified isomorphism
\[
 E_k^\vee/\operatorname{im}\Lambda_k^\vee
       \xrightarrow{\ \pi_k\ }\mathbb C^{T_k}
       \cong (\ker\Lambda_k)^\vee.
 \tag{OCF24}
\]
The displayed $\pi_k$ uses exactly the original primary coefficient
coordinates under $\Theta_k$. Every dualization in this statement
is complex-linear. No conjugation is inserted in this coefficient
map; the original metric maps can subsequently be composed with it.

Finally, this construction is finite with explicitly bounded
compressed matrices. OCF20 has $n_d$ rows and
$\sum_{e=2}^5 |G_e|j_{d-e}$ columns, with negative subscripts zero.
Both dimensions grow at most linearly in $d$: for example
$n_d\leq16d+1$ and the column number is at most $25(16d+1)$.
For $d\geq13$ the exact column number is $200d-1632$, obtained
from $\sum_e|G_e|=25$, $\sum_e e|G_e|=104$, and
$j_{d-e}=8(d-e)-32$. Only the preceding five graded bases are
needed to form the next matrix. Each column is obtained from
the at-most-$36$-coefficient formula OCF12, one exact convolution,
and OCF6. A full original biform coefficient column, if used as
temporary storage, has $(d+1)^2$ entries; every invariant relation
matrix in OCF20 has the stated linear row and column counts.
The invariant certificates can be retained as their products and
parent-column references, so their expanded high-degree monomials
need never be enumerated. The proof of their completeness is
OCF11--16, and the exact original quotient and its right inverse
are OCF21--24.

\subsection{The same recursion with fourteen proved generators}

The complete fixed set OCF9 can be reduced to the following
fourteen specific members while preserving every map and image
just proved:
\[
\begin{array}{c|l}
e&\mathcal G_e\\\hline
2&x_1x_4,\ x_2x_3\\
3&x_1^2x_3,\ x_1x_2^2,\ x_2x_4^2,\ x_3^2x_4\\
4&x_1^3x_2,\ x_1x_3^3,\ x_2^3x_4,\ x_3x_4^3\\
5&x_1^5,\ x_2^5,\ x_3^5,\ x_4^5 .
\end{array}\tag{OCF25}
\]
Here is a complete proof. Call a positive-degree invariant
monomial indecomposable when it is not the product of two
positive-degree invariant monomials. The divisor proof preceding
OCF11 shows that each indecomposable has degree at most five.
If it contains both $x_1,x_4$, it is their product, since otherwise
division by this invariant quadratic leaves a nonconstant
invariant monomial. The same assertion holds for $x_2,x_3$.
After these cases are removed, at most two variables occur:
every three indices contain one of the two opposite pairs.
One variable gives precisely its fifth power. The remaining
two-element supports are $\{1,2\},\{1,3\},\{2,4\},\{3,4\}$.
For positive exponents $(a,b)$ with $a+b\leq5$, the equation
$a+2b\equiv0\pmod5$ has exactly $(a,b)=(3,1),(1,2)$, and
$a+3b\equiv0\pmod5$ has exactly $(a,b)=(2,1),(1,3)$.
For $\{2,4\}$ divide the weights by $2$ in the field
$\mathbb Z/5\mathbb Z$, giving the first list; for $\{3,4\}$
divide by $3$, giving the second. These are exactly the eight
mixed monomials of degrees three and four in OCF25. Every
listed monomial has no proper invariant divisor, as is also
seen from these lists and the excluded opposite pairs.
This proves the classification. Factoring any invariant
monomial repeatedly into indecomposables terminates because
positive total degree strictly decreases. Hence OCF25 generates
the entire invariant algebra.

Use $\mathcal G=\bigcup_{e=2}^5\mathcal G_e$ in place of
$\{x^\nu:\nu\in G\}$ in OCF12--16 and OCF20. The same
factorization argument proves the same exact invariant image;
all coefficient formulas, invariant certificates, pivot charts,
quotient maps and ranks remain valid. The new recursion has
$2,4,4,4$ generators in degrees $2,3,4,5$, respectively. Thus
for $d\geq13$ its precise matrix dimensions and rank are
\[
 C_d^{(14)}\in
 \operatorname{Mat}_{(16d-48)\,\times\,(112d-864)}(\mathbb C),
 \qquad \operatorname{rank}C_d^{(14)}=8d-32.
 \tag{OCF26}
\]
Indeed the number of generators is $14$ and the sum of their
degrees is $2\cdot2+4\cdot3+4\cdot4+4\cdot5=52$; substituting
the already proved $j_{d-e}=8(d-e)-32$ gives
$14(8d-32)-8\cdot52=112d-864$. At smaller degrees the exact
number of columns is the same sum of the appropriate $j_{d-e}$
from OCF19, with negative subscripts zero. All coefficients
in this reduced recursion are still the original coefficients
of OCF12--14.

% END COMPLETE OCF

\clearpage
% BEGIN COMPLETE OKM
% Complete independent metric/CD provider. Original source orders only.
\section{The exact common kernel frame and its original Gamma metrics}
\label{sec:OKM}

Retain the original simple quartet $m=1$, $k=4l+1$, $l\geq1$,
$q=(k+1)^2$, $c=k/2$, $0<\delta<1/2$ and $\gamma>2$.
The original period parameter and its branch are fixed, on the proved
original OCF rank domain $|u|\geq R_*$, with $R_*$ the retained
explicit period-orbit radius. The kernel
below is the kernel of that one fixed original observation
$\Lambda:E\to B$. It is common to the two source metrics and to
their four degree endpoints. No intersection over different period
parameters or different observations is made.

Order all pairs $(a,b)$, $0\leq a,b\leq k$, lexicographically once
and retain that order. Define
\[
 \begin{gathered}
 \beta_{ab}=c+(2a-k)\delta+i(2b-k)\gamma,\qquad
 \omega_{ab}=\frac{\beta_{ab}-c}{i}
                =(2b-k)\gamma-i(2a-k)\delta,\\
 \chi(S)=\prod_{a,b=0}^k(S-\beta_{ab}),\qquad
 E=\mathbb C[S]/(\chi),\qquad
 \varepsilon_{ab}=
 \left[\prod_{(a',b')\ne(a,b)}
       \frac{S-\beta_{a'b'}}{\beta_{ab}-\beta_{a'b'}}\right],\\
 V_\beta=((\beta_{ab})^d)_{(a,b),\,0\leq d<q}.
 \end{gathered}\tag{OKM1}
\]
The points are distinct: equality of their real and imaginary parts
forces both indices equal. Hence every denominator displayed above
is nonzero. Evaluation at all points gives
$\varepsilon_{ab}(\beta_{a'b'})=\mathbf1_{(a,b)=(a',b')}$.
Polynomial interpolation proves that these are a basis of $E$, and
the evaluation map from increasing polynomial coordinates to this
primary basis is exactly $V_\beta$. Its determinant is the ordered
Vandermonde product, so it is invertible.

\subsection{The exact dual and primal maps of the strip quotient}
Use the coefficient presentation of OQC7--8 in Segre coordinates.
A biform of bidegree $(k,k)$ has the full coefficient expression
\[
 F=\sum_{a,b=0}^k f_{ab}
       X_1^aX_0^{k-a}Y_1^bY_0^{k-b}.
 \tag{OKM2}
\]
The literal variable dictionary to the OCF biform coefficients is
$X_1=z_0$, $X_0=z_1$, $Y_1=w_0$, $Y_0=w_1$.
The Segre substitution has
$t_{++}=X_1Y_1$, $t_{+-}=X_1Y_0$,
$t_{-+}=X_0Y_1$, $t_{--}=X_0Y_0$.
At the original exponential vector take
$X_1=e^{-i\delta w}$, $X_0=e^{i\delta w}$,
$Y_1=e^{\gamma w}$, $Y_0=e^{-\gamma w}$.
Then the monomial with coefficient $f_{ab}$ is $e^{\omega_{ab}w}$.
The original coefficient identification includes the entire invertible
matrix $\mathsf H$ of OQC1 before this Segre presentation; its unit
and period entries remain in the matrices produced by the strip
calculation. OQC7 is characterized by evaluation on $e^{wY}[1]$.
Since
$e^{wY}[1]=\sum_{a,b}e^{\omega_{ab}w}\varepsilon_{ab}$,
the independent exponentials in OQC4 therefore give the precise
complex-linear dual pairing
\[
 \Theta(F)(\varepsilon_{ab})=f_{ab},\qquad
 \Theta(F)(x)=f^{\mathsf T}x_{\rm prim}.
 \tag{OKM3}
\]
There is no binomial coefficient or complex conjugation in this
complex-linear dual pairing. The statement follows by equating all
coefficients of the distinct exponentials, whose independence is
proved by their Vandermonde derivative matrix in OQC4.

Let $J$ be exactly the coefficient subspace of the original observed
dual image, as identified in OQC8. The original strip and invariant
elimination provider OCF constructs its literal surjection
\[
 \pi:\mathbb C^q\twoheadrightarrow\mathbb C^{r_K},
 \qquad\ker\pi=J,\qquad r_K=8k-16.
 \tag{OKM4}
\]
This input is the actual OCF matrix, including its stated period
coefficients and its valid pivot choice. The coefficient pairing
OKM3 and OQC8 identify
$K:=\ker\Lambda=J^\perp$ in primary coordinates, where
$J^\perp=\{x:f^{\mathsf T}x=0\text{ for every }f\in J\}$.
Consequently its exact primal frame and polynomial frame are
\[
 I_{\rm prim}=\pi^{\mathsf T}:\mathbb C^{r_K}\hookrightarrow E_{\rm prim},
 \qquad I_{\rm poly}=V_\beta^{-1}\pi^{\mathsf T}:
                        \mathbb C^{r_K}\hookrightarrow E.
 \tag{OKM5}
\]
Here is a direct proof, including inverse coordinates. Choose any of
the explicit right sections $L_\pi$ of the surjective matrix $\pi$,
so $\pi L_\pi=I_{r_K}$. For $x=\pi^{\mathsf T}v$ and $f\in J$,
$f^{\mathsf T}x=(\pi f)^{\mathsf T}v=0$. The transpose is
injective because $L_\pi^{\mathsf T}\pi^{\mathsf T}=I_{r_K}$.
Conversely, if $x\in J^\perp$, then
$(I-L_\pi\pi)f\in J$ for every $f$, so
$x^{\mathsf T}(I-L_\pi\pi)f=0$ for every $f$. Thus
$x=\pi^{\mathsf T}L_\pi^{\mathsf T}x$.
This proves equality of the image and $K$, and its inverse on that
image is $L_\pi^{\mathsf T}$. In polynomial coordinates the inverse
is $L_\pi^{\mathsf T}V_\beta$. This proves the full typed primal
arrow, rather than selecting a Hermitian transpose in place of it.

\subsection{The complete original source and quotient metric}
For either original source $s\in\{1,k\}$ write $b_s=s/2$ and
$M_s=(2\pi)^{s/2}$. The unchanged real-coordinate measure is
\[
 dm_{0,s}(y)=M_s\frac{2^{b_s}}{4\pi\Gamma(b_s)}
      |\Gamma(b_s/2+iy/2)|^2\,dy,
 \qquad S=c+iy.
 \tag{OKM6}
\]
The complete canonical OGN1--3 proof gives the monic real
orthogonal polynomials and all their squared norms:
\[
 \begin{gathered}
 p_0(y)=1,\quad p_1(y)=y,\quad
 p_{d+1}(y)=yp_d(y)-d(b_s+d-1)p_{d-1}(y),\\
 h_d=M_s d!(b_s)_d,\qquad
 u_d(S)=\frac{p_d((S-c)/i)}{\sqrt{h_d}}.
 \end{gathered}\tag{OKM7}
\]
The square root is positive. Thus $u_d$ has the original leading
coefficient $i^{-d}/\sqrt{h_d}$ in $S$ and is orthonormal on the
line in OKM6. No source mass or phase has been removed.

For a polynomial cutoff $D\geq q-1$ let
\[
 \begin{gathered}
 T_D=\left(\frac{p_d(\omega_{ab})}{\sqrt{h_d}}\right)_
                        {(a,b),\,0\leq d\leq D},
 \qquad R_D=T_DT_D^*,\\
 G_D^{\rm prim}=R_D^{-1},\qquad
 G_D^{\rm poly}=V_\beta^*R_D^{-1}V_\beta,\qquad
 H_D^K=\overline\pi\,R_D^{-1}\pi^{\mathsf T}.
 \end{gathered}\tag{OKM8}
\]
The star denotes conjugate transpose, while the superscript
$\mathsf T$ is ordinary transpose. In particular
$(\pi^{\mathsf T})^*=\overline\pi$, with precisely the indicated
matrix dimensions.

These are the actual minimum quotient metrics, as follows directly
from the original remainder map. In the orthonormal coefficient
frame $u_0,\ldots,u_D$, the remainder map followed by primary
evaluation is exactly $T_D$, column by column. Its first $q$
columns are invertible: the monic polynomials $p_0,\ldots,p_{q-1}$
evaluated at the distinct $\omega_{ab}$ have a nonzero Vandermonde
determinant, and all $\sqrt{h_d}$ are nonzero. Thus $R_D>0$.
For a prescribed primary vector $x$, the vector
$T_D^*R_D^{-1}x$ maps to $x$. It is orthogonal to $\ker T_D$,
because its inner product with a kernel vector contains $T_D$
applied to that vector. Every other lift differs by such a vector,
so its squared norm is
$x^*R_D^{-1}x$ plus the squared norm of that difference.
This proves the first quotient metric in OKM8.
The identity $x_{\rm prim}=V_\beta x_{\rm poly}$ proves the
second by direct pairing. Restriction to the exact frame OKM5
then gives
\[
 I_{\rm poly}^*G_D^{\rm poly}I_{\rm poly}
 = (\pi^{\mathsf T})^*R_D^{-1}\pi^{\mathsf T}
 =H_D^K>0.
 \tag{OKM9}
\]
The last positivity follows from injectivity of $\pi^{\mathsf T}$.
The map $K\hookrightarrow E$ is the fixed original coefficient
map. No invariance of that kernel under multiplication is used.

\subsection{Christoffel--Darboux with the exact source mass}
For arbitrary $z,w\in\mathbb C$ and $D\geq0$, the complete
polynomial identity is
\[
 \sum_{d=0}^D\frac{p_d(z)p_d(w)}{h_d}
 =\begin{cases}
 \displaystyle\frac{p_{D+1}(z)p_D(w)-p_D(z)p_{D+1}(w)}
                     {h_D(z-w)},&z\ne w,\\[6pt]
 \displaystyle\frac{p_{D+1}'(z)p_D(z)-p_D'(z)p_{D+1}(z)}{h_D},
                       &z=w.
 \end{cases}\tag{OKM10}
\]
To prove it, put $a_d=d(b_s+d-1)$, with $a_0=0$, and set
$p_{-1}=0$ for this recurrence expression. The recurrence
in OKM7 gives
\[
 \begin{split}
 (z-w)p_d(z)p_d(w)
 &=p_{d+1}(z)p_d(w)-p_d(z)p_{d+1}(w)\\
 &\quad-a_d\bigl(p_d(z)p_{d-1}(w)-p_{d-1}(z)p_d(w)\bigr).
 \end{split}\tag{OKM11}
\]
For $d\geq1$, the full norm identity is $a_d/h_d=1/h_{d-1}$.
Divide OKM11 by $h_d$ and sum. Every adjacent term cancels,
leaving just the degree-$D$ numerator divided by $h_D$; the
$d=0$ lower term is zero. This proves the first line of OKM10.
Both sides of the resulting equality after multiplication by
$z-w$ are polynomials. Differentiating with respect to $z$ and
then setting $z=w$ proves the second line with the displayed
order and sign. This also proves the value at every confluent
point without an assertion about a limiting singular quotient.
The denominator is the full $h_D=M_sD!(b_s)_D$, including $M_s$.

Since all coefficients of the recurrence are real, the matrix
entries in OKM8 are therefore exactly the values of OKM10 at
\[
 z=\omega_{ab},\qquad w=\overline{\omega_{a'b'}}.
 \tag{OKM12}
\]
Their indexed conjugation and confluent locus are
\[
 \begin{gathered}
 \overline{\omega_{a'b'}}=\omega_{k-a',b'},\qquad
 \omega_{ab}-\overline{\omega_{a'b'}}
       =2(b-b')\gamma-2i(a+a'-k)\delta,\\
 \omega_{ab}=\overline{\omega_{a'b'}}
       \quad\Longleftrightarrow\quad b=b'\text{ and }a+a'=k.
 \end{gathered}\tag{OKM13}
\]
These follow by substituting OKM1; positivity of $\delta,\gamma$
proves both implications in the last line. Thus the conjugation
map on the point indices is precisely $(a,b)\mapsto(k-a,b)$.
The confluent formula of OKM10 applies at these indexed entries.
For the retained odd $k$, its matrix diagonal has $a'=a,b'=b$,
so $2a=k$ cannot occur; the literal Hermitian diagonal is given
by the first line of OKM10. The conjugate-index entries use its
second line. All entries, including the positive diagonal, are
simultaneously defined by the original polynomial sum on the
left of OKM10.

\subsection{Every pivot change and the four original endpoints}
Let $\pi'$ be another valid strip/invariant elimination matrix for
the same original $J$, with the same surjective codomain dimension.
Choose right sections $L_\pi,L_{\pi'}$ and put
\[
 A_{\pi',\pi}=\pi'L_\pi,
 \qquad A_{\pi,\pi'}=\pi L_{\pi'}.
 \tag{OKM14}
\]
Because $\pi'$ vanishes on $J=\ker\pi$, one has
$\pi' L_\pi\pi=\pi'$. Thus
$\pi'=A_{\pi',\pi}\pi$. Reversing the roles gives the
opposite identity, and surjectivity of each matrix then proves
$A_{\pi',\pi}A_{\pi,\pi'}=I$ and
$A_{\pi,\pi'}A_{\pi',\pi}=I$. This proves invertibility and
also independence from the chosen right sections. Writing
$A=A_{\pi',\pi}$ for these identities gives the exact frame
and Gram changes
\[
 \begin{gathered}
 I'_{\rm prim}=I_{\rm prim}A^{\mathsf T},\qquad
 I'_{\rm poly}=I_{\rm poly}A^{\mathsf T},\\
 (H_D^K)'=\overline A H_D^K A^{\mathsf T},\qquad
 \det(H_D^K)'=|\det A|^2\det H_D^K.
 \end{gathered}\tag{OKM15}
\]
The first two formulas follow by transposing $\pi'=A\pi$.
Substitution into OKM9 proves the congruence and determinant.
This retains all actual pivot determinants. They have not been
assigned modulus one.

The kernel is the same fixed $K=\ker\Lambda$ used in MRI1--5.
If its original BRU inclusion is $I_0$ and its fixed-section kernel
coordinate map is $\kappa_0$, let
\[
 \mathsf C_K=\kappa_0I_{\rm poly},\qquad
 \mathsf C_K^{-1}=L_\pi^{\mathsf T}V_\beta I_0.
 \tag{OKM16}
\]
The equalities $I_0\kappa_0x=x$ for $x\in K$ and
$L_\pi^{\mathsf T}V_\beta I_{\rm poly}=I_{r_K}$ prove that
these are inverse matrices and $I_0\mathsf C_K=I_{\rm poly}$.
Consequently
$H_D^K=\mathsf C_K^*(I_0^*G_D^{\rm poly}I_0)\mathsf C_K$.
The matrix $\mathsf C_K$ is independent of $D$ and of the choice of source
order $s$. The source metrics change; the two specified coefficient
frames do not.

It follows that the complete original four-endpoint kernel scalar is
\[
 \boxed{\displaystyle
 F_K^{(s)}=
 \log\det H_{q-1}^K+\log\det H_q^K
       -\log\det H_{2q-1}^K-\log\det H_{2q}^K.}
 \tag{OKM17}
\]
Indeed MRI3 gives the actual consecutive kernel determinant ratio
\[
 \frac{\det(I_0^*G_{q+j-1}^{\rm poly}I_0)}
      {\det(I_0^*G_{q+j}^{\rm poly}I_0)}
 =\frac{1+t_j}{1+\xi_j}.
\]
Each factor $|\det\mathsf C_K|^2$ cancels in that ratio by OKM16.
The weighted sum with endpoint weights $1,2,\ldots,2,1$ therefore
has precisely the four signs in OKM17, as in MRI4. This verifies
the original degree cutoffs, rather than shifting the column count
in $T_D$. Under any valid pivot change OKM15, the added scalar
$2\log|\det A|$ has total coefficient $1+1-1-1=0$, proving
invariance of OKM17 with its original frames.

For emphasis, the full mass in OKM7 appears in every column of
$T_D$ and in the denominator of OKM10. If that common mass is
factored explicitly, write $T_D=M_s^{-1/2}\widehat T_D$, where
$\widehat T_D$ has exactly the same numerator polynomials and
the denominator $\sqrt{d!(b_s)_d}$. Then
$H_D^K=M_s\overline\pi(\widehat T_D\widehat T_D^*)^{-1}
\pi^{\mathsf T}$, so its determinant has factor $M_s^{r_K}$.
It cancels only in the displayed signed four-endpoint scalar,
whose four ranks are equal. The individual metrics in OKM8--9
retain it. No inserted Gamma order, discarded action defect, or
kernel-invariance assertion enters any of these calculations.

% END COMPLETE OKM

\clearpage
% BEGIN COMPLETE OKA
\section{The original kernel graph, boundary basis and complete action}
\label{sec:OKA}

This section uses the complete conductor coefficient construction OCF
and its exact original dual identification in OKM. Fix the original
simple quartet, its complete amplitude unit and periods, the original
parameter on any logarithm branch with $|u|\geq R_*$, and
$k=4l+1$, $l\geq1$. Put $q=(k+1)^2$, $r=8k-16$, and
$\mathcal I_k=\{(a,b):0\leq a,b\leq k\}$. The roots, centre and
centred eigenvalues remain
\[
 c=k/2,\qquad
 \beta_{ab}=c+(2a-k)\delta+i(2b-k)\gamma,\qquad
 \omega_{ab}=(\beta_{ab}-c)/i
             =(2b-k)\gamma-i(2a-k)\delta.
 \tag{OKA1}
\]
The original algebra is $E=\mathbb C[S]/(\chi)$, with
$\chi(S)=\prod_{(a,b)\in\mathcal I_k}(S-\beta_{ab})$.
Let $\varepsilon_{ab}$ be its original primary idempotents and
$\mathcal V:E\to\mathbb C^{\mathcal I_k}$ the evaluation
isomorphism $[p]\mapsto(p(\beta_{ab}))_{ab}$. Its inverse sends
$e_{ab}$ to
\[
 \varepsilon_{ab}(S)
 =\frac{\chi(S)}{(S-\beta_{ab})\chi'(\beta_{ab})}.
 \tag{OKA2}
\]
These denominators are nonzero: equality of two roots in OKA1,
on comparing real and imaginary parts and using $\delta>0$,
$\gamma>2$, implies equality of both indices. Evaluation proves
that the polynomials in OKA2 are the coordinate idempotents,
so their columns give the exact inverse, with no changed root
or polynomial coefficient convention.

\subsection{The coefficient quotient gives a graph in the original primary basis}

Let $\pi:\mathbb C^{\mathcal I_k}\to\mathbb C^r$ be the
surjective coefficient quotient of OCF, whose kernel is precisely
the image $J_k$ of the invariant polynomials under the original
restriction. The domain here consists of the coefficients of the
biforms
$z_0^a z_1^{k-a}w_0^b w_1^{k-b}$, and is the complex-linear
dual of the displayed primary coordinate space. Let
$T\subset\mathcal I_k$ be the final OCF complement of its
invariant-image pivot rows, viewed as a subset of the full
biform coefficient grid. Its ordered coordinate inclusion is
$E_T:\mathbb C^r\to\mathbb C^{\mathcal I_k}$. The explicit
OCF formula gives
\[
 \pi E_T=I_r,\qquad
 N=\pi^T:\mathbb C^r\longrightarrow\mathbb C^{\mathcal I_k},
 \qquad E_T^TN=I_r.
 \tag{OKA3}
\]
Indeed a retained strip monomial in $T$ already has zero
conductor remainder outside its own coordinate, zero coordinates
on the invariant pivot rows, and its indicated coordinate equal
to one. Both reductions in OCF therefore leave that quotient
coordinate equal to its standard unit vector. Transposing gives
the last identity. This uses ordinary transpose, because the
restriction quotient and its dual are complex-linear maps.

Write $J=\mathcal I_k\setminus T$ only for the complementary
index set in this section; the invariant image retains the name
$J_k$. Denote coordinate restrictions by $x_T,x_J$, and put
$N_J=E_J^TN$. The original kernel inclusion is consequently
\[
 I_K=\mathcal V^{-1}N:\mathbb C^r\hookrightarrow E,
 \qquad
 \operatorname{im}I_K=K:=\ker\Lambda,\qquad
 N=\begin{pmatrix}I_r\\N_J\end{pmatrix}_{T,J}.
 \tag{OKA4}
\]
For completeness, a primary vector $v$ pairs with a biform
coefficient vector $f$ by $f^Tv$. The annihilator of
$\ker\pi=J_k$ is $\operatorname{im}\pi^T$: containment
follows by multiplication, while surjectivity of $\pi$ gives
both subspaces dimension $r$. OQC8 identifies precisely this
annihilator with the primary coordinates of the fixed original
$\ker\Lambda$. This proves both the original image assertion
and its direction in OKA4. In polynomial coordinates the same
inclusion is $V_\beta^{-1}N$, where
$(V_\beta)_{(a,b),d}=\beta_{ab}^d$ for $0\leq d<q$;
this is the full coefficient matrix of OKA2.

\subsection{An explicit basis of the actual boundary}

For $j\in J$ define the actual boundary vector
$b_j=\Lambda\varepsilon_j$. These vectors form a basis of
the original image $B=\operatorname{im}\Lambda$. To prove
independence, if $\sum_{j\in J}c_j\varepsilon_j$ lies in
$K$, its primary coordinates equal $N\zeta$. Its $T$
coordinates vanish, so OKA3 gives $\zeta=0$ and all $c_j=0$.
To prove spanning, use the literal decomposition
\[
 x=N x_T+E_J(x_J-N_Jx_T),\qquad
 \Lambda\mathcal V^{-1}x
       =\sum_{j\in J}b_j(x_J-N_Jx_T)_j.
 \tag{OKA5}
\]
The first equality follows separately on its $T$ and $J$
coordinates. Apply $\Lambda$, which kills the first term by
OKA4, to obtain the second. Thus no observation has been
substituted: in this basis the specified original observation
has the exact coefficient matrix $(-N_J,I_{q-r})$.

Here are the complete original period and tensor coordinates of
these boundary vectors. Write the quartet roots as
$\alpha_{\epsilon,\eta}=1/2+\epsilon\delta+i\eta\gamma$,
where $\epsilon,\eta\in\{+1,-1\}$ and the signs $+,-$ in
subscripts denote these same two numbers. In a subscript
$\alpha$ below means that complex root, indexed by its sign
pair. Keep the original one-factor primary idempotents
$\varepsilon^h_\alpha\in E_h$. Set
\[
 v_\alpha=\Pi U\varepsilon^h_\alpha
       =v_h(\alpha)\Pi\varepsilon^h_\alpha,
 \qquad
 Z_+=z_0,\ Z_-=z_1,\ W_+=w_0,\ W_-=w_1,
 \quad
 \mathfrak v(z,w)=\sum_{\epsilon,\eta}
                         v_{\epsilon\eta}Z_\epsilon W_\eta.
 \tag{OKA6}
\]
The vectors belong to the original one-factor period space;
all entries of $\Pi$ and all four actual nonzero amplitude
values remain present. For every $(a,b)\in\mathcal I_k$,
\[
 \boxed{\quad
 \jmath\Lambda\varepsilon_{ab}
  =[z_0^a z_1^{k-a}w_0^b w_1^{k-b}]
       P_k\bigl(\mathfrak v(z,w)^{\otimes k}\bigr),\qquad
 P_k=\frac15\sum_{h=0}^4(C^{-h})^{\otimes k}.
 \quad}\tag{OKA7}
\]
The original sum inclusion $\iota:E\hookrightarrow E_h^{\otimes k}$
sends $\varepsilon_{ab}$ to the sum of the primary tensor
idempotents whose word has $a$ positive real signs and $b$
positive imaginary signs. This assertion follows by evaluation:
on a primary word, the sum variable takes exactly the value
$\beta_{ab}$ of its two counts, and OKA2 evaluates to one
for those counts and zero otherwise. Applying
$(\Pi U)^{\otimes k}$ gives precisely the corresponding sum
of the tensors of the $v_\alpha$ in OKA6. Expanding the tensor
power and taking the indicated biform coefficient selects
exactly the same words, each once. The original definition
$\jmath\Lambda=P_k(\Pi U)^{\otimes k}\iota$ proves OKA7,
including the literal inverse powers and $1/5$. Its specialization
to $j\in J$ therefore supplies the actual basis $b_j$ of
OKA5. Neither a different target nor a numerical period matrix
has been introduced.

The fixed, source-independent kernel coordinate and section are
now explicitly
\[
 \begin{gathered}
 \kappa=E_T^T\mathcal V:E\to\mathbb C^r,
 \qquad S_*:B\to E,\quad S_*b_j=\varepsilon_j,
 \\
 I_K\kappa+S_*\Lambda=I_E,
 \quad\kappa I_K=I_r,\quad\kappa S_*=0,\quad\Lambda S_*=I_B.
 \end{gathered}
 \tag{OKA8}
\]
Every equality follows by applying OKA5 to an arbitrary primary
vector. In particular this is an actual splitting of the fixed
original short exact sequence, used unchanged for both source
orders and for all polynomial degrees.

\subsection{The whole original action and its retained defect}

The original centred sum action is $Y=(S-c)/i$ on $E$.
Under $\mathcal V$ it is the diagonal matrix of the
$\omega_{ab}$ in OKA1. Write its two diagonal submatrices
as $\Omega_T$ and $\Omega_J$. In the exact coordinates
$e=I_K\zeta+S_*b$ of OKA8, it has the complete matrix
\[
 \boxed{\quad
 Y\ \longleftrightarrow\
 \begin{pmatrix}
   \Omega_T&0\\
   \Omega_JN_J-N_J\Omega_T&\Omega_J
 \end{pmatrix},\qquad
 (\Omega_JN_J-N_J\Omega_T)_{jt}
                 =(\omega_j-\omega_t)(N_J)_{jt}.
 \quad}\tag{OKA9}
\]
Indeed the primary vector of $e$ is
$(\zeta,N_J\zeta+b)$. Applying the original diagonal matrix
gives $(\Omega_T\zeta,\Omega_JN_J\zeta+\Omega_Jb)$.
Applying the coordinates in OKA5 subtracts
$N_J\Omega_T\zeta$ from its second component and proves
every block in OKA9. This proves, with its precise codomain,
\[
 YI_K-I_K\Omega_T
     =S_*\bigl(\Omega_JN_J-N_J\Omega_T\bigr)
                :\mathbb C^r\to E,\qquad
 \Lambda YI_K=\Omega_JN_J-N_J\Omega_T:\mathbb C^r\to B,
 \tag{OKA10}
\]
where the matrix on the right uses the basis $b_j$.
The chosen complementary span of the original primary idempotents
is stable under $Y$; no stability of $K$ has been used. Its
entire failure is the displayed coefficient map.

For each original source $s\in\{1,k\}$ retain the monic
Gamma polynomials $p_d^{(s)}$, the complete norms
$h_d^{(s)}=(2\pi)^{s/2}d!(s/2)_d$, and the original
$e_d^{(s)}=[p_d^{(s)}(Y)]/\sqrt{h_d^{(s)}}$. Then the
fixed-section kernel and boundary coordinates of every actual
remainder, including all degrees $d\geq q$, are
\[
 \begin{split}
 (\kappa e_d^{(s)})_t
      &=\frac{p_d^{(s)}(\omega_t)}{\sqrt{h_d^{(s)}}},\\
 (\Lambda e_d^{(s)})_j
      &=\frac{p_d^{(s)}(\omega_j)
                -\sum_{t\in T}(N_J)_{jt}p_d^{(s)}(\omega_t)}
              {\sqrt{h_d^{(s)}}}.
 \end{split}\tag{OKA11}
\]
Polynomial evaluation in the original primary algebra gives the
full vector $p_d^{(s)}(\omega)/\sqrt{h_d^{(s)}}$ even when
its polynomial representative must be reduced modulo $\chi$.
Applying OKA5 and OKA8 proves OKA11. Thus the formula includes
the original finite remainder operation and every phase and source
mass. Composing its second line with the actual columns OKA7
gives exactly the period-vector increments of OPG. Its first line
is the fixed-section kernel vector used by OPG and OPR; the
degree-dependent minimum-section correction in OPR remains
required for the mixed norm. The metric computation in OKM
applies to the common inclusion OKA4 and does not replace that
correction by the fixed-section vector alone.

% END COMPLETE OKA

\clearpage
% BEGIN COMPLETE OMG
% Full original multiplicity geometry. No frozen predecessor is edited.
\section{The complete homogeneous geometry of the original multiplicity curve}
\label{sec:OMG}

\subsection{The actual primary, unit, period and Fourier coefficient map}
Retain the original full-order quartet and every original period
parameter. Write
\[
\begin{gathered}
 \alpha_{\epsilon,\eta}=\tfrac12+\epsilon\delta+i\eta\gamma,
 \quad\epsilon,\eta\in\{-1,1\},\quad0<\delta<\tfrac12,\quad\gamma>2,\\
 m\ge1,\quad d_m=m-1,\quad n=4m,\quad N=4m+1,\quad
 h(z)=\prod_{\epsilon,\eta}(z-\alpha_{\epsilon,\eta})^m,\\
 E_h=\mathbb C[z]/(h),\quad A_h=M_z,\quad
 v_h=2\xi/h,\quad U=M_{j_hv_h},\quad v_h(\alpha_{\epsilon,\eta})\ne0.
\end{gathered}\tag{OMG1}
\]
The last inequalities follow from the stipulated complete original
zero order $m$. The map $j_h$ records all divided Taylor coefficients
of orders $0,\ldots,m-1$ at all four roots. Let
$\mathcal E_{\rm prim}:\mathbb C^{4m}\to E_h$ send its coordinate
$(\alpha,d)$ to the original primary element $\varepsilon_\alpha
(z-\alpha)^d$, $0\le d<m$. Explicitly its representative is
\[
 E_{\alpha,d}(z)=\operatorname{rem}_h\left[
 h_\alpha(z)\sum_{j=0}^{m-1}
       \frac{(h_\alpha^{-1})^{(j)}(\alpha)}{j!}(z-\alpha)^{j+d}
                                      \right],\quad
 h_\alpha=h/(z-\alpha)^m.
\]
The Taylor inverse is defined since $h_\alpha(\alpha)\ne0$.
It gives the idempotent one on its own primary factor and zero
on the others. The four coprime factors therefore prove that
$\mathcal E_{\rm prim}$ is an isomorphism and include every
nilpotent degree. In this basis the full matrix of $U$ is
\[
 U_{\rm prim;\alpha,a;\alpha,d}
 =\begin{cases}v_h^{(a-d)}(\alpha)/(a-d)!,&a\ge d,\\0,&a<d,
   \end{cases}\qquad
 \det U=\prod_\alpha v_h(\alpha)^m\ne0.
 \tag{OMG2}
\]
Other primary blocks are zero. Inverting its constant coefficient
plus its nilpotent part by the finite geometric series proves
invertibility with all derivatives retained.

Keep the original $u\ne0$ and its logarithm branch, the primitive
$\Phi_h'=h$, $\Phi_h(0)=0$, the outward rays $\ell_j$ of angles
$(\arg u+\pi+2\pi j)/N$, and $\Gamma_j=-\ell_0+\ell_j$.
The exact period and Fourier maps are
\[
\begin{gathered}
 (\Pi[P])_j=\int_{\Gamma_j}e^{\Phi_h(z)/u}
                     \operatorname{rem}_hP(z)\,dz,\quad1\le j\le n,\\
 \zeta=e^{2\pi i/N},\quad F_{jt}=\zeta^{jt}-1,\quad
 (F^{-1})_{tj}=\zeta^{-tj}/N\quad(1\le j,t\le n),\quad
 C=F\operatorname{diag}(\zeta,\ldots,\zeta^n)F^{-1}.
\end{gathered}\tag{OMG3}
\]
The full original period theorem PDE/OPG gives the convergence and
invertibility of $\Pi$ on these representatives. The Fourier inverse
follows directly from
$\sum_{j=1}^{N-1}\zeta^{-tj}(\zeta^{sj}-1)=N\mathbf1_{t=s}$.
It is the same literal cyclic matrix as OCS4, including its inverse
powers in the original average. Its inherited Gram is
$F^*F=N(I_n+\mathbf1\mathbf1^*)$, as expansion of its entries gives.

Define the exact diagonal map and coefficient matrix
\[
\begin{gathered}
 D_m^{\rm jet}=\operatorname{diag}\left(i^{-d}/d!\right)_{\alpha,\,0\le d<m},
 \qquad \mathsf H_m=F^{-1}\Pi U\mathcal E_{\rm prim}D_m^{\rm jet},\\
 (\mathsf H_m)_{t;\alpha,d}
  =\frac{i^{-d}}{d!}\sum_{a=d}^{m-1}
     \left[\frac1N\sum_{j=1}^{n}\zeta^{-tj}
              \int_{\Gamma_j}e^{\Phi_h(z)/u}E_{\alpha,a}(z)\,dz\right]
                  \frac{v_h^{(a-d)}(\alpha)}{(a-d)!}.
\end{gathered}\tag{OMG4}
\]
This is exactly OCS13, now regarded as one square matrix. All
factors in its first line are invertible. In particular
\[
 \det\mathsf H_m=(\det F)^{-1}\det\Pi\,
      \left(\prod_\alpha v_h(\alpha)^m\right)
      \det\mathcal E_{\rm prim}
      \prod_\alpha\prod_{d=0}^{m-1}\frac{i^{-d}}{d!}\ne0,
\]
with the determinant of $\mathcal E_{\rm prim}$ taken in the fixed
original coefficient frame. None of these factors is prescribed.

Set $\omega_{\epsilon,\eta}=\eta\gamma-i\epsilon\delta$ and let
\[
 z_{\epsilon,\eta,d}(w)=w^d e^{\omega_{\epsilon,\eta}w},\qquad
 x(w)=F^{-1}\Pi Ue^{w(A_h-1/2)/i}[1].
 \tag{OMG5}
\]
The terminating exponential on each primary factor has coefficient
$i^{-d}w^d/d!$ on $\varepsilon_\alpha(z-\alpha)^d$.
Multiplication by OMG2 and application of the literal period maps
therefore prove the exact identity and inverse coordinate map
\[
 \boxed{x(w)=\mathsf H_m z(w),\qquad z(w)=\mathsf H_m^{-1}x(w).}
 \tag{OMG6}
\]
For the original Euclidean period metric, its Gram on these primary
curve coordinates is $\mathsf H_m^*N(I_n+\mathbf1\mathbf1^*)
\mathsf H_m$. Thus the coefficient change also retains its full
metric and all off-diagonal entries, without claiming orthogonality.

\subsection{The complete homogeneous ideal, with quadratic generators}
Use binary indices $a,b\in\{0,1\}$ for
$\epsilon=2a-1$, $\eta=2b-1$. Let
$\mathscr P_Z=\mathbb C[Z_{a,b,d}:a,b\in\{0,1\},0\le d\le d_m]$
with each variable of degree one. Define the algebra map
\[
 \Phi_m:\mathscr P_Z\longrightarrow
        \mathbb C[X_0,X_1,Y_0,Y_1,T_0,T_1],\qquad
 Z_{a,b,d}\longmapsto X_aY_bT_0^{d_m-d}T_1^d.
 \tag{OMG7}
\]
Its image is the complete diagonal graded algebra
\[
 R_m=\bigoplus_{k\ge0}
  \mathbb C[X_0,X_1]_k\otimes\mathbb C[Y_0,Y_1]_k
                              \otimes\mathbb C[T_0,T_1]_{kd_m}.
 \tag{OMG8}
\]
Indeed a basis monomial of its degree-$k$ component is
\[
 X_0^{k-a}X_1^aY_0^{k-b}Y_1^bT_0^{kd_m-D}T_1^D,
       \quad0\le a,b\le k,\quad0\le D\le kd_m.
 \tag{OMG9}
\]
It is a product of $k$ images in OMG7: choose $k$ binary first
indices with sum $a$, independently $k$ binary second indices with
sum $b$, and $k$ integers in $[0,d_m]$ with sum $D$; the last
choice follows by filling successive slots up to $d_m$. Combining
the three lists into rows gives the required preimage monomial.
Distinct displayed target monomials are linearly independent.

Let $\mathfrak t_m$ be generated by every quadratic binomial
\[
 Z_{a,b,d}Z_{a',b',d'}-Z_{c,e,f}Z_{c',e',f'}
 \quad\text{with}\quad
 a+a'=c+c',\quad b+b'=e+e',\quad d+d'=f+f',
 \tag{OMG10}
\]
where every index belongs to its original range. Repeated variables
are allowed. Then
\[
 \boxed{\ker\Phi_m=\mathfrak t_m,\qquad
           \mathscr P_Z/\mathfrak t_m\simeq R_m.}
 \tag{OMG11}
\]
Here is a complete ideal proof. Each displayed quadratic has zero
image. Two degree-$k$ monomials have the same image exactly when
the sums of their first indices, second indices and third indices
agree. Regard their $k$ factors temporarily as labelled rows. Align
the first indices with the target list by pairwise swaps of first
indices, keeping each row's other indices fixed. Each swap is an
OMG10 relation. Next align the second indices by the same operation,
keeping the aligned first indices fixed. Finally, if the third
indices differ from their target list, choose a row below its target
and a row above its target. Increase the former third index by one
and decrease the latter by one. These entries stay in $[0,d_m]$
because their target entries do; the total sum is unchanged and the
sum of absolute deviations decreases by two. This terminates at
the target list. Each transfer is another OMG10 relation. Thus the
difference between any two monomials with the same image is a sum
of multiples of the quadrics, by telescoping the successive monomial
differences. Labels only specify the sequence of polynomial maps;
the products are the original commutative monomials throughout.
For any polynomial in $\ker\Phi_m$, collect its coefficients by
their target monomial. Each such coefficient sum is zero by target
monomial independence. Subtracting one chosen monomial in each
group expresses that whole group as a linear combination of the
proved monomial differences. This proves the reverse inclusion for
all degrees and hence the entire ideal equality.

Every degree-$k$ original curve monomial has the function
\[
 w^D\exp\{[(2b-k)\gamma-i(2a-k)\delta]w\}
       \quad(0\le a,b\le k,\ 0\le D\le kd_m).
 \tag{OMG12}
\]
The frequencies are distinct: equality of real parts gives equality
of $b$ and equality of imaginary parts gives equality of $a$.
To prove independence with all polynomial powers, suppose
$\sum_{a,b}P_{a,b}(w)e^{\omega_{a,b}w}=0$, with
$\deg P_{a,b}\le kd_m$. Fix a pair and apply
$\prod_{(a',b')\ne(a,b)}(\partial_w-\omega_{a',b'})^{kd_m+1}$.
This kills every other summand. On its surviving polynomial each
factor is $\partial_w+(\omega_{a,b}-\omega_{a',b'})$, invertible
on the space of degree at most $kd_m$ because differentiation is
nilpotent and the constant is nonzero. The surviving polynomial
is therefore zero. Repeating this for every pair proves independence.

Thus the full homogeneous ideal of the original primary curve is
exactly $\mathfrak t_m$: a homogeneous polynomial restricts to zero
if and only if its coefficient sums at every target monomial OMG9
are zero. In particular
\[
 \boxed{\dim R_{m,k}=(k+1)^2[1+k(m-1)]=q_k,
 \quad\ker\{\mathscr P_{Z,k}\to\operatorname{Hol}(\mathbb C),
                      f\mapsto f(z(w))\}=\mathfrak t_{m,k}.}
 \tag{OMG13}
\]
This statement concerns the homogeneous ideal, equivalently the
projective closure of $[z(w)]$. Every one of its degrees has been
computed, including degree zero. No assertion about a different
nonhomogeneous affine ideal is used.

\subsection{The projective model and its complete chart inverses}
For $m\ge2$, the morphism
\[
 (\mathbb P^1)^3\longrightarrow\mathbb P^{4m-1},\qquad
 ([X_0:X_1],[Y_0:Y_1],[T_0:T_1])
      \longmapsto[X_aY_bT_0^{d_m-d}T_1^d]_{a,b,d}
 \tag{OMG14}
\]
has multidegree $(1,1,m-1)$. At least one first, second and third
coordinate is nonzero; their corresponding extreme product is
nonzero, so the map is everywhere defined. Its image and scheme
structure are exactly $\operatorname{Proj}R_m$, with complete
closed ideal OMG10. Here are explicit inverses proving this claim.
The eight charts $Z_{a,b,0}\ne0$ and $Z_{a,b,d_m}\ne0$ cover
$\operatorname{Proj}R_m$. Indeed OMG11 gives
\[
 Z_{a,b,d}^{d_m}=Z_{a,b,0}^{d_m-d}Z_{a,b,d_m}^{d}
             \quad\text{in }R_m.
\]
A homogeneous prime containing every extreme variable contains
every variable by this equality, and hence gives no point of Proj.
On $Z_{a,b,0}\ne0$ the three inverse affine coordinates are
\[
 x=Z_{1-a,b,0}/Z_{a,b,0},\quad
 y=Z_{a,1-b,0}/Z_{a,b,0},\quad
 t=Z_{a,b,1}/Z_{a,b,0}.
\]
For every variable its ratio to $Z_{a,b,0}$ is
$x^{\mathbf1_{a'\ne a}}y^{\mathbf1_{b'\ne b}}t^d$, by OMG7
and its proved injective quotient. Consequently the degree-zero
localized ring is $\mathbb C[x,y,t]$. These three coordinates
are algebraically independent, since their images in the parameter
fraction field are $X_{1-a}/X_a,Y_{1-b}/Y_b,T_1/T_0$.
For the chart at $Z_{a,b,d_m}$ replace the last ratio by
$t=Z_{a,b,d_m-1}/Z_{a,b,d_m}$; the general ratio has exponent
$d_m-d$ in $t$. Its coordinate ring is again $\mathbb C[x,y,t]$.
On chart overlaps these formulas invert the same parameter ratios,
so they glue to the inverse of OMG14. This proves the isomorphism
of schemes and its full defining homogeneous ideal. For $m=1$
the third factor has degree zero and the same argument uses the
four first/second charts: the model is the Segre $(\mathbb P^1)^2$
with ideal $Z_{0,0,0}Z_{1,1,0}-Z_{0,1,0}Z_{1,0,0}$.

In the actual Fourier coordinates let
$\mathscr P_x=\mathbb C[x_1,\ldots,x_n]$ and define
\[
 \mathfrak a=\{f(x):f(\mathsf H_mZ)\in\mathfrak t_m\},\qquad
 R=\mathscr P_x/\mathfrak a.
 \tag{OMG15}
\]
The invertible substitution OMG6 proves that $\mathfrak a$ is the
complete homogeneous ideal of $x(w)$ and that $R\simeq R_m$.
Its explicit generators are all OMG10 quadrics with every $Z$
replaced by its component of $\mathsf H_m^{-1}x$. These are the
actual period and unit coefficients of OMG4, and $\mathfrak a$
is prime because the target ring in OMG7 is a domain. Thus the
original curve's projective closure is the precise linear image
of OMG14 under $\mathsf H_m$, with no prescribed substitute matrix.

\subsection{The exact graded source-dual algebra and every factorial}
For each $k\ge0$ retain the original sum algebra
\[
\begin{gathered}
 e_k=1+kd_m,\quad c_k=k/2,\quad
 \beta_{k,a,b}=c_k+(2a-k)\delta+i(2b-k)\gamma,\\
 \chi_k(S)=\prod_{a,b=0}^k(S-\beta_{k,a,b})^{e_k},\quad
 E_k=\mathbb C[S]/(\chi_k),\quad Y_k=(M_S-c_kI)/i.
\end{gathered}\tag{OMG16}
\]
For $k=0$ this is $E_0=\mathbb C[S]/(S)=\mathbb C$. Let
$H_{k,a,b,D}=\varepsilon_{k,a,b}(S-\beta_{k,a,b})^D$,
$0\le D<e_k$, be the complete original primary basis obtained
by the same Taylor inverse as in OMG2. It satisfies
\[
 e^{wY_k}[1]=\sum_{a,b=0}^k e^{\omega_{k,a,b}w}
             \sum_{D=0}^{e_k-1}\frac{i^{-D}w^D}{D!}H_{k,a,b,D},
 \quad\omega_{k,a,b}=(2b-k)\gamma-i(2a-k)\delta.
 \tag{OMG17}
\]
For $f\in\mathscr P_{x,k}$ expand its original function uniquely as
$f(x(w))=\sum c_{a,b,D}(f)w^De^{\omega_{k,a,b}w}$ using OMG12.
Define the complex linear dual map, without a conjugation,
\[
 \Theta_k(f)(H_{k,a,b,D})=i^DD!c_{a,b,D}(f).
 \tag{OMG18}
\]
Then OMG17 proves exactly
$\Theta_k(f)(e^{wY_k}[1])=f(x(w))$.
The independence and surjectivity in OMG9--13 prove that this
map is onto $E_k^\vee$, with kernel $\mathfrak a_k$.
Consequently it induces a specified vector-space isomorphism
\[
 \boxed{\Theta_k:R_k\xrightarrow{\ \sim\ }E_k^\vee.}
 \tag{OMG19}
\]
All phases and factorials are those of OCS17; for a monomial
$x^\nu$ its coefficients are precisely the finite multinomial
OCS15 rule with the entries OMG4.

The collection of these vector spaces has an exact graded algebra
structure coming from the original sum maps. For $k,j\ge0$ put
\[
 \Delta_{k,j}:E_{k+j}\longrightarrow E_k\otimes E_j,
               \qquad[P(S)]\longmapsto[P(S_1+S_2)].
 \tag{OMG20}
\]
On a primary pair the eigenvalue is the displayed sum
$\beta_{k,a,b}+\beta_{j,a',b'}=\beta_{k+j,a+a',b+b'}$.
The sum of the two nilpotents has index $e_k+e_j-1=e_{k+j}$:
its power $e_k+e_j-2$ has nonzero top coefficient
$\binom{e_k+e_j-2}{e_k-1}$ and its next power is zero. Every
pair of total sign counts occurs. Thus the full minimal polynomial
of $S_1+S_2$ is exactly $\chi_{k+j}$, proving both well-definedness
and injectivity of OMG20. It also gives the complete primary formula
\[
 \Delta_{k,j}H_{k+j,A,B,D}
 =\sum_{\substack{a+a'=A,\ b+b'=B}}
   \sum_{\substack{d+d'=D\\0\le d<e_k,\ 0\le d'<e_j}}
       \binom Dd H_{k,a,b,d}\otimes H_{j,a',b',d'}.
 \tag{OMG21}
\]
Indeed the total primary idempotent is one precisely on the
matching eigenvalue pairs, and the nilpotent power expands by the
ordinary binomial theorem there. All other components vanish.
Substitution of the sum of three coordinates proves coassociativity;
interchange of the two coordinates proves cocommutativity. The
$k=0$ maps give the counit. Hence
\[
 (\ell\star\ell')(v)=(\ell\otimes\ell')(\Delta_{k,j}v)
       \quad(\ell\in E_k^\vee,\ \ell'\in E_j^\vee)
 \tag{OMG22}
\]
is an associative commutative graded product with unit in $E_0^\vee$.
For the exact dual basis $\lambda_{k,a,b,D}=i^DD!H_{k,a,b,D}^\vee$,
OMG21 gives
\[
 \lambda_{k,a,b,D}\star\lambda_{j,a',b',D'}
       =\lambda_{k+j,a+a',b+b',D+D'}.
\]
The coefficient is exactly
$\binom{D+D'}D i^DD!\,i^{D'}D'!=i^{D+D'}(D+D')!$;
this records the cancellation before stating the basis product.
Equivalently, OMG20 sends $e^{wY_{k+j}}[1]$ to
$e^{wY_k}[1]\otimes e^{wY_j}[1]$, with
$c_{k+j}=c_k+c_j$. Evaluating OMG18 on this identity proves
\[
 \boxed{\Theta:\ R\xrightarrow{\ \sim\ }
                    \bigoplus_{k\ge0}(E_k^\vee,\star)
                  \text{ is an isomorphism of graded algebras}.}
 \tag{OMG23}
\]
Its source multiplication is the actual polynomial multiplication;
its target multiplication has been constructed from the original sum
coalgebra rather than assigning a multiplication between unrelated
finite-dual vector spaces.

\subsection{The full semisimple and nilpotent action in this geometry}
On $\mathscr P_Z$ define the degree-preserving derivation
\[
 \partial_Z Z_{a,b,d}=\omega_{2a-1,2b-1}Z_{a,b,d}
                   +\begin{cases}d Z_{a,b,d-1},&d>0,\\0,&d=0.
                     \end{cases}
 \tag{OMG24}
\]
On the parameter ring OMG7 it is induced by
$\partial X_a=-i(2a-1)\delta X_a$,
$\partial Y_b=(2b-1)\gamma Y_b$,
$\partial T_0=0$, $\partial T_1=T_0$.
The product rule verifies OMG24 and commutation with $\Phi_m$;
therefore the complete ideal $\mathfrak t_m$ is stable.
The corresponding exact flow on the parameters is
\[
 (X_a,Y_b,T_0,T_1)\longmapsto
 (e^{-i(2a-1)\delta w}X_a,e^{(2b-1)\gamma w}Y_b,T_0,T_1+wT_0).
\]
At $(X_0,X_1,Y_0,Y_1,T_0,T_1)=(1,1,1,1,1,0)$ it gives precisely
the original functions OMG5. This is the entire stated analytic
flow on the explicit algebraic variety, with all polynomial degrees.

If $B_Z$ is the matrix in OMG24 on the column of variables, its
rows have diagonal $\omega_\alpha$ and entry $d$ in column
$(\alpha,d-1)$ of row $(\alpha,d)$. In actual coordinates the
derivation is
\[
 \partial_x x=\mathcal A x,\qquad
 \mathcal A=\mathsf H_mB_Z\mathsf H_m^{-1}.
 \tag{OMG25}
\]
The inverse substitution proves $\partial_x\mathfrak a\subseteq
\mathfrak a$, and its flow differentiates the original $x(w)$.
On the exact degree-$k$ basis of OMG9/OMG18 the induced action is
\[
 \partial_x\lambda_{k,a,b,D}
   =\omega_{k,a,b}\lambda_{k,a,b,D}
             +\begin{cases}D\lambda_{k,a,b,D-1},&D>0,\\0,&D=0.
              \end{cases}\qquad
 \Theta_k\partial_x=Y_k^\vee\Theta_k.
 \tag{OMG26}
\]
The displayed action on the basis uses its identification OMG23.
For a direct verification, on the original primal primary basis
$Y_kH_D=\omega H_D+i^{-1}H_{D+1}$ with the last term zero at
$D=e_k-1$. Applying $i^DD!H_D^\vee$ proves the exact lowering
coefficient $D$ in its dual action. Each frequency block has all
$e_k$ terms, with nonzero top iterated derivative
$(e_k-1)!$. Thus every original primary nilpotent has exactly the
same index $e_k$ in this model. The full annihilating polynomial
is retained with its leading phase:
\[
 \chi_k(c_k+i\partial_x)=0\quad\text{on }R_k,\qquad
 \chi_k(c_k+iT)=i^{q_k}\prod_{a,b=0}^k(T-\omega_{k,a,b})^{e_k}.
 \tag{OMG27}
\]
Distinct frequencies and the nonzero top chain prove that this
degree $q_k$ is minimal. The original $S$ action corresponds to
$c_k+i\partial_x$ on degree $k$, or to
$\tfrac12\mathsf E+i\partial_x$ on the graded algebra, where
$\mathsf E$ is the degree derivation. These equalities specify
every constant, sign and nilpotent action under OMG23.

\subsection{The exact original cyclic observation as invariant restriction}
Let $D=\operatorname{diag}(\zeta,\ldots,\zeta^n)$ and let
$G=\mathbb Z/N$ act on $\mathscr P_x$ by
$(g f)(x)=f(D^{-1}x)$. Thus a degree-$k$ monomial $x^\nu$
has character $\zeta^{-\sum t\nu_t}$. Write
\[
 \mathscr I=\mathscr P_x^G,\qquad
 \mathscr J=\operatorname{im}(\mathscr I\longrightarrow R).
 \tag{OMG28}
\]
These are graded algebras. The coefficient image of the original
observation is exactly their restriction at every degree, as follows.
Let $f_t$ be column $t$ of $F$ and let $S_\nu$ be the sum of
all distinct ordered words in the $f_t$ having multiplicity vector
$\nu$. Disjoint word supports make the $S_\nu$ a basis of the
permutation-invariant tensor space. The expansion and projector are
\[
 (Fx)^{\otimes k}=\sum_{|\nu|=k}x^\nu S_\nu,\qquad
 P_k(Fx)^{\otimes k}=\sum_{\substack{|\nu|=k\\\sum t\nu_t=0\pmod N}}
                    x^\nu S_\nu,\quad
 P_k=\frac1N\sum_{a=0}^{N-1}(C^{-a})^{\otimes k}.
\]
Every ordered word occurs once; the factor $k!/\prod\nu_t!$ is
the number of terms in $S_\nu$, not a missing coefficient.
The original sum injection is the iterated map OMG20 into
$E_h^{\otimes k}$. Hence the original observation
$L_k=P_k\Pi^{\otimes k}U^{\otimes k}\iota_k=\jmath\Lambda_k$
sends $e^{wY_k}[1]$ to the displayed projected curve with $x=x(w)$.
The coefficient covectors of its surviving $S_\nu$ pull back to
exactly $\Theta_k(x^\nu)$. They span the full dual of the image,
since every covector on a subspace of a finite dimensional space
extends to its ambient space. Consequently
\[
\boxed{\begin{gathered}
 \Theta_k(\mathscr J_k)=\operatorname{im}\Lambda_k^\vee
                          =(\ker\Lambda_k)^\perp,\qquad
 \operatorname{rank}\Lambda_k=\dim\mathscr J_k,\\
 R_k/\mathscr J_k\xrightarrow{\ \Theta_k|_{K_k}\ }K_k^\vee
       \text{ is an isomorphism of vector spaces},\qquad
 K_k=\ker\Lambda_k.
\end{gathered}}\tag{OMG29}
\]
For the last statement, a covector on $E_k$ vanishes on $K_k$
exactly when it factors through $E_k/K_k\simeq\operatorname{im}
\Lambda_k$. Restriction to $K_k$ is onto by extension of a basis.
This proves the claimed quotient and its precise type: the
degree-$k$ quotient is a vector-space quotient by $\mathscr J_k$.

\subsection{The actual orbit ideal, including every stabilizer}
The whole ideal $\mathfrak a$ is generated by the explicit quadratic
space in OMG15. Let $B_2$ be its coefficient matrix in the ordered
degree-two monomial frame of $\mathscr P_x$. Repeated or dependent
generator rows are harmless. For $a\in\mathbb Z/N$ let
$D_{2,a}$ be diagonal with entry $\zeta^{-a\sum t\nu_t}$ at
monomial $x^\nu$. The actual stabilizer and its exact finite test are
\[
 H=\{a:g^a\mathfrak a=\mathfrak a\}
   =\left\{a:\operatorname{rank}\begin{bmatrix}B_2\\B_2D_{2,a}\end{bmatrix}
                      =\operatorname{rank}B_2\right\}.
 \tag{OMG30}
\]
Equality of these row spaces is equivalent to equality of their
quadratically generated ideals; their dimensions agree under the
invertible diagonal map. This proves the finite test using the
actual entries OMG4. It prescribes no generic coefficient value.
The set $H$ is a subgroup by composition and inverse of ideal
equalities. Put $s=|H|$ and $d=N/s$. Since $G$ is cyclic, its
cosets have representatives $0,\ldots,d-1$ and
$H=\{0,d,2d,\ldots,(s-1)d\}$. Define all distinct component ideals
and the union algebra by
\[
 \mathfrak a_a=g^a\mathfrak a\ (0\le a<d),\qquad
 \mathfrak u=\bigcap_{a=0}^{d-1}\mathfrak a_a,
 \qquad A=\mathscr P_x/\mathfrak u.
 \tag{OMG31}
\]
These are distinct homogeneous prime ideals with the same Hilbert
function OMG13. Their intersection is a homogeneous radical ideal
stable under the original cyclic action. It is the complete ideal
of their scheme-theoretic union by the definition of the union's
ideal. No product formula for this intersection is substituted.

For any invariant polynomial, membership in one component ideal
implies membership in all its translates, by applying every $g^a$.
Thus the exact invariant restriction and its inverse description are
\[
\boxed{\begin{gathered}
 \mathscr I\cap\mathfrak a=\mathscr I\cap\mathfrak u,\\
 \mathscr I/(\mathscr I\cap\mathfrak u)
      \xrightarrow{\ [f]\mapsto[f]_{\mathfrak u}\ } A^G
      \xrightarrow{\operatorname{res}_0} \mathscr J\subseteq R
      \quad\text{are isomorphisms of graded algebras}.
\end{gathered}}\tag{OMG32}
\]
To prove surjectivity of the first map, average any representative
of an invariant class by $N^{-1}\sum g^a$; its class is unchanged.
Its kernel is exactly the displayed intersection. For the second,
use an invariant representative: a zero restriction lies in
$\mathscr I\cap\mathfrak a=\mathscr I\cap\mathfrak u$ and hence
is already zero in $A^G$. Its image is $\mathscr J$ by definition.
This gives the invariant image for every actual $H$, including
nontrivial stabilizers of a composite group order $4m+1$.

\subsection{The complete component conductor and an explicit finite witness}
Set $R_a=\mathscr P_x/\mathfrak a_a$. The simultaneous restriction
\[
 A\hookrightarrow\widetilde A:=\prod_{a=0}^{d-1}R_a,
 \qquad[f]_{\mathfrak u}\longmapsto([f]_{\mathfrak a_a})_a
 \tag{OMG33}
\]
is injective precisely because its kernel before quotienting is the
intersection OMG31. Its component conductor ideals are
\[
 \mathfrak c_a=
 \frac{(\bigcap_{b\ne a}\mathfrak a_b)+\mathfrak a_a}{\mathfrak a_a}
       \subseteq R_a,
 \qquad\bigcap_{b\ne a}\mathfrak a_b=\mathscr P_x\quad(d=1).
 \tag{OMG34}
\]
The conductor of this specified inclusion is exactly
\[
 \boxed{\{t\in\widetilde A:t\widetilde A\subseteq A\}
                   =\prod_{a=0}^{d-1}\mathfrak c_a\subseteq A.}
 \tag{OMG35}
\]
Indeed a tuple supported solely on component $a$ lies in the image
of $A$ if and only if its value belongs to $\mathfrak c_a$, by
lifting its other zero values to the intersection of the other
ideals. If $t$ is in the conductor, multiplication by each component
idempotent of $\widetilde A$ gives such a supported tuple; thus each
component of $t$ lies in $\mathfrak c_a$. Conversely these are
ideals of $R_a$. Multiplying a tuple with all its components in
the indicated ideals by any element of $\widetilde A$ gives a sum
of such supported tuples, each lying in $A$. This proves both
inclusions and the full formula. In particular this conductor is
also the annihilator of $\widetilde A/A$ in $A$.

The action of $G$ permutes these components with its original
coefficient substitutions. Projection of an invariant tuple to
component zero gives an isomorphism
\[
 \widetilde A^G\xrightarrow{\ \sim\ }R^H,\qquad
 r\in R^H\longmapsto(g^a r)_{0\le a<d}
                       \text{ for its inverse}.
 \tag{OMG36}
\]
The inverse is independent of coset representatives because $r$
is fixed by $H$, and invariance follows by reindexing the components.
The same argument applies to the conductor, whose ideals are
permuted. Equations OMG32--36 consequently prove the exact
inclusions of graded subalgebras and conductor ideals
\[
 \boxed{\mathfrak c_0^H\subseteq\mathscr J\subseteq R^H,
            \qquad\mathfrak c_0^H\,R^H\subseteq\mathscr J.}
 \tag{OMG37}
\]
Here $\mathfrak c_0^H$ denotes invariant elements of the specified
component ideal and is an ideal of $R^H$. An invariant tuple in
the conductor already lies in $A$, so its projection belongs to
$\mathscr J$ by OMG32. This proves the first inclusion; the
ideal property gives the second.

There is a nonzero homogeneous element of this conductor with a
completely bounded degree. For every $a\ne0$ choose from the
explicit transformed quadratic list of $\mathfrak a_a$ a quadratic
$f_a\notin\mathfrak a_0$. Such a choice exists: containment of
the entire quadratic space in that of $\mathfrak a_0$, which has
the same dimension, would make the spaces equal and their
generated ideals equal, contrary to distinctness. Membership is
the finite row-space test in the same degree-two monomial frame.
The product $f=\prod_{a=1}^{d-1}f_a$ lies in all other component
ideals. It is nonzero modulo $\mathfrak a_0$ because that ideal is
prime and none of its factors belongs to it. Thus
$\bar f\in\mathfrak c_0$ is nonzero of degree $2(d-1)$.
The group $H$ preserves $\mathfrak c_0$ and acts on the domain
$R$. Its exact norm
\[
 c_H=\prod_{h\in H}g^h\bar f\in\mathfrak c_0^H,\qquad
 \deg c_H=D_H=2s(d-1)=2(N-s),\qquad c_H\ne0
 \tag{OMG38}
\]
is invariant by permutation of its factors and nonzero in the
domain. For $d=1$ take the empty product $c_H=1$, $D_H=0$;
all statements remain exact. Multiplication by this actual
conductor element is injective and gives the finite dimension
bounds, with negative graded pieces set to zero,
\[
\boxed{\begin{gathered}
 (R^H)_{k-D_H}\xrightarrow{\ \times c_H\ }\mathscr J_k
                    \hookrightarrow(R^H)_k,\\
 \dim(R^H)_{k-D_H}\le\operatorname{rank}\Lambda_k
                                  \le\dim(R^H)_k,\\
 q_k-\dim(R^H)_k\le\dim K_k
                                  \le q_k-\dim(R^H)_{k-D_H}.
\end{gathered}}\tag{OMG39}
\]
Every ideal, group and matrix in these expressions is defined by
the actual original coefficient map. In particular the stabilizer
has its finite equality test OMG30 throughout; no particular value
of it has been inserted into the dimensions. To compute the displayed
invariant dimensions directly, let the actual $H$ substitutions act
on the basis OMG9 through $Z\mapsto\mathsf H_m^{-1}D^{-h}\mathsf H_mZ$
and reduce with OMG10. This is an explicit finite matrix on $q_k$
coordinates. The idempotent average $s^{-1}\sum_{h\in H}g^h$ has
trace equal to its rank, proving
$\dim(R^H)_k=s^{-1}\sum_{h\in H}\operatorname{Tr}(g^h|R_k)$.
Thus the dimensions in OMG39 also have complete original finite
coefficient formulas for every actual stabilizer.

\subsection{The exact action defect and all subsequently observed directions}
The derivation $\partial_x$ of OMG25 acts on $R$ with all primary
directions OMG26. The possible failure of its restriction to the
observed image is the specified finite linear map
\[
 \mathscr J_k\longrightarrow R_k/\mathscr J_k,
                       \qquad j\longmapsto[\partial_xj].
 \tag{OMG40}
\]
Via OMG29 it is exactly the complex linear dual of the original
map $\Lambda_kY_kI_K:K_k\to B_k$. Indeed a $j$ corresponds to
the covector $\lambda\Lambda_k$ on $E_k$ for a unique boundary
covector $\lambda$; evaluating OMG26 on $I_Kv$ gives
$\lambda\Lambda_kY_kI_Kv$, proving the identification and all
domains and codomains. The action on ambient polynomials also has
the exact conjugation formula
\[
 g^a\partial_xg^{-a}=\partial_{D^a\mathcal A D^{-a}},\qquad
 g^a\partial_xg^{-a}-\partial_x
                   =\partial_{D^a\mathcal A D^{-a}-\mathcal A}.
 \tag{OMG41}
\]
For a linear vector field $\partial_M f=(Mx)\cdot\nabla f$,
the chain rule proves this by substituting $D^{-a}x$ after
differentiating $f(D^a x)$. This formula retains every original
period/unit coefficient in $\mathcal A$ and does not require the
original ideal or the observed image to be preserved by extra
ambient actions.

Finally the subspace invisible to all original iterates has the
complete model
\[
\boxed{\begin{gathered}
 \mathcal K_{\infty,k}=\bigcap_{r=0}^{q_k-1}\ker(\Lambda_kY_k^r),\\
 R_k\Big/\sum_{r=0}^{q_k-1}\partial_x^r\mathscr J_k
       \xrightarrow{\ \Theta_k|_{\mathcal K_{\infty,k}}\ }
                          \mathcal K_{\infty,k}^\vee
                          \quad\text{is a vector-space isomorphism}.
\end{gathered}}\tag{OMG42}
\]
In finite dimensional duality the annihilator of an intersection
of kernels is the sum of the images of their dual maps: this
follows by taking the dual of the single stacked map with all
those components, whose kernel is precisely the intersection.
Equations OMG26 and OMG29 identify these images with
$\partial_x^r\mathscr J_k$, proving OMG42. The exact original
polynomial OMG27 expresses every higher iterate as a linear
combination of the first $q_k$, so the same intersection is taken
over all nonnegative iterates and is $Y_k$-invariant. This proves
the full nilpotent and cyclic-action receiver in the new homogeneous
model for every original multiplicity, with its actual observation,
periods and amplitude jets throughout.

% END COMPLETE OMG


\clearpage
% BEGIN CURRENT AKS
% AKS: new continuation, 14 September 2026.
% This body is self-contained at the level of its explicitly stated inherited inputs.
\section{The absolute common kernel: the new finite result}
\label{sec:AKS}

This complete incoming calculation sharpens the current absolute common
kernel estimate. Its original comparison was against the preceding
logarithmic bound; the intervening KAF chapter already established
the $O(kq)$ scale. The new finite envelope reduces its uniform limsup
constant from $16(4+10\log2)$ to $32\log(25\log3/(4\pi))$.
AKJ below takes the pointwise minimum of both finite bounds. No
limiting kernel coefficient is inferred from these upper bounds.

Retain the actual hypothetical complete quartet and its entire quotient
\[
 \rho=\tfrac12+\delta+i\gamma,\qquad 0<\delta<\tfrac12,\quad \gamma>2,
 \qquad h(z)=\prod_{\alpha\in\{\rho,\bar\rho,1-\rho,1-\bar\rho\}}
 (z-\alpha)^m,\quad v_h=2\xi/h.
\]
The full primary unit is $U=M_{j_hv_h}$. The original tensor order, centre,
polynomial, and coefficient quotient are
\[
\begin{gathered}
 k=4l+1,\quad l\ge1,\quad e=1+k(m-1),\quad q=e(k+1)^2,\quad c=k/2,\\
 \lambda_{a,b}=c+(2a-k)\delta+i(2b-k)\gamma,\qquad
 \chi(S)=\prod_{a,b=0}^k(S-\lambda_{a,b})^e,
 \quad E=\mathbb C[S]/(\chi).
\end{gathered}\tag{AKS1}
\]
For the principal conclusion $m=1$. The source-operator estimates below
also hold with the stated repeated-root list for $m>1$, but no simple-quartet
rank is assigned to those multiplicities.

Let $\Lambda_k(u):E\twoheadrightarrow B_k(u)$ be the original observation,
$K(u)=\ker\Lambda_k(u)$, and $I:K(u)\hookrightarrow E$ its original
coefficient inclusion. For a source $s\in\{1,k\}$, let $G_N^{(s)}$ be the
original minimum-remainder Gram, in the frame $(1,S,\ldots,S^{q-1})$.
Define $H_{K,N}^{(s)}=I^*G_N^{(s)}I$ and retain exactly
\[
 F_K^{(s)}=
 \log\det H_{K,q-1}^{(s)}+\log\det H_{K,q}^{(s)}
 -\log\det H_{K,2q-1}^{(s)}-\log\det H_{K,2q}^{(s)}.
 \tag{AKS2}
\]
Every empty determinant is one. All conjugate transposes refer to the
specified coordinate matrices, not to a bilinear residue pairing.

Here are entirely explicit positive numbers used in the bound:
\[
\begin{gathered}
 R=k\sqrt{\delta^2+\gamma^2},\quad b_s=s/2,\quad
 t_q=\frac\pi2-\frac1q,\quad \mathfrak a_q=\frac{2\log3}{t_q},\quad
 \mathfrak s_q=\frac{\mathfrak a_q^q-1}{\mathfrak a_q-1},\\
 \mathfrak b_{s,q}=\left[\frac9{25}\sin\frac1q\right]^{-b_s}
       \exp\left(2R\left(\log3+\frac{t_q}{2}\right)\right),\\
 \mathcal Z_{s,j}=1+\mathfrak b_{s,q}\mathfrak s_q^2
        \sum_{d=q}^{q+j-1}\frac{d!}{(b_s)_d}
                          \left(\frac{25}{16}\right)^d,
 \quad \mathcal L_{s,j}=\log \mathcal Z_{s,j}\quad (0\le j\le q+1).
\end{gathered}\tag{AKS3}
\]
The empty sum gives $\mathcal Z_{s,0}=1$. In the admitted family $q\ge36$, so
$t_q\in(0,\pi/2)$ and $\mathfrak a_q>1$. These parameters are radii in the existing
Laplace and generating functions; neither is a new source order.

\paragraph{Main theorem (absolute finite bound).}
For the original simple quartet, every original logarithm branch, and
every $|u|\ge R_*$, where $R_*$ is precisely RPC21 or RPC23 (recorded in
Section~\ref{sec:AKS-radius}), put $r=8k-16$. Then
\[
 \boxed{\quad
 0\le F_K^{(s)}(k,u)\le r\bigl(\mathcal L_{s,q}+\mathcal L_{s,q+1}\bigr),
 \qquad s\in\{1,k\}.\quad}
 \tag{AKS4}
\]
There is a stronger, fully period-dependent interval in AKS28 below.
In particular, with
\[
 \kappa_* =32\log\left(\frac{25\log3}{4\pi}\right),
 \qquad 25.02078097649<\kappa_*<25.02078097650,
 \tag{AKS5}
\]
one has the uniform statement
\[
 \boxed{\quad
 0\le\limsup_{k\to\infty\atop k\equiv1\ (4)}
       \ \sup_{\substack{|u|\ge R_*\\\text{original branches}}}
         \frac{F_K^{(s)}(k,u)}{kq}\le\kappa_*.
 \quad}\tag{AKS6}
\]
Here the source may be $s=1$ for every $k$, or $s=k$. The same bound holds
for the actual arithmetic kernel. Its exact finite allowances are in
AKS34--36. Thus the absolute kernel is $O_h(kq)$ on this domain, not only
$O_h(kq\log q)$. Neither AKS5 nor AKS6 calls $\kappa_*$ the actual leading
coefficient.

\section{An exact coefficient realization of the actual lost directions}
\label{sec:AKS-frame}

\subsection{The period matrix and invariant restriction, with their full coefficients}
For $m=1$, order the one-factor roots as $(++),(+-),(-+),(--)$, and let
$\mathcal E_{\rm prim}$ have their original Lagrange idempotents as columns.
Set $\zeta=\exp(2\pi i/5)$ and $F_{jt}=\zeta^{jt}-1$, $1\le j,t\le4$.
The exact inverse is $(F^{-1})_{tj}=\zeta^{-tj}/5$. Retain
\[
\begin{gathered}
 \Pi[P]_j=\int_{\Gamma_j}e^{\Phi_h(z)/u}\operatorname{rem}_hP(z)\,dz,
 \qquad\Phi_h'=h,\quad\Phi_h(0)=0,\\
 \Gamma_j=-\ell_0+\ell_j,\quad
 \arg\ell_j=(\arg u+\pi+2\pi j)/5,\\
 \mathsf H=F^{-1}\Pi U\mathcal E_{\rm prim},\quad
 D=\operatorname{diag}(\zeta,\zeta^2,\zeta^3,\zeta^4),\\
 Q(x)=Q_0(\mathsf H^{-1}x),\qquad
 Q_0(t)=t_1t_4-t_2t_3,\quad Q_a(x)=Q(D^{-a}x).
\end{gathered}\tag{AKS7}
\]
The integral always uses the degree-less-than-four representative.
Invertibility of $\Pi$, all contour phases, and its full determinant are
the PDE/PQO results in the supplied proofs. The four unit values are
$(a,\bar a,\bar a,a)$, with the actual $a=v_h(\rho)\ne0$; no value of $a$
is selected. Hence $\mathsf H$ is the original invertible matrix.

The following formulas make the coefficient realization directly
computable. Write
\[
 \mathcal I_{k,0}=\{\nu\in\mathbb Z_{\ge0}^4:|\nu|=k,
                         \ \nu_1+2\nu_2+3\nu_3+4\nu_4\equiv0\pmod5\}.
\]
For $\nu\in\mathcal I_{k,0}$ define $B_{\nu;(a,b)}$ by summing over
$4\times4$ arrays $n_{t,\alpha}\ge0$ with row sums $\nu_t$, total plus-real
count $a$, and total plus-imaginary count $b$:
\[
 B_{\nu;(a,b)}=
 \sum_{\substack{\sum_\alpha n_{t,\alpha}=\nu_t\\
       \sum_{t,\alpha:\epsilon_\alpha=+}n_{t,\alpha}=a\\
       \sum_{t,\alpha:\eta_\alpha=+}n_{t,\alpha}=b}}
 \left(\prod_{t=1}^4\frac{\nu_t!}{\prod_\alpha n_{t,\alpha}!}\right)
                  \prod_{t,\alpha}\mathsf H_{t,\alpha}^{n_{t,\alpha}}.
 \tag{AKS8}
\]
Let $\mathsf V_{(a,b),d}=\lambda_{a,b}^d$ for $0\le d<q$ and set
$\mathsf A=B\mathsf V$. All entries of $\mathsf A$ are now finite
expressions in the actual periods and unit. There is no generic period
matrix in this definition.

For completeness, put $t_\alpha(w)=\exp((\eta_\alpha\gamma-
 i\epsilon_\alpha\delta)w)$ and $x(w)=\mathsf Ht(w)$. Expansion gives
\[
 x(w)^\nu=\sum_{a,b}B_{\nu;(a,b)}e^{\omega_{a,b}w},\qquad
 \omega_{a,b}=(2b-k)\gamma-i(2a-k)\delta.
\]
In the original primary algebra,
$e^{wY}[1]=\sum_{a,b}e^{\omega_{a,b}w}\varepsilon_{a,b}$, where
$Y=(M_S-cI)/i$. The exponentials are independent by the Vandermonde of
the distinct $\omega_{a,b}$. OCS9's unscaled symmetric word basis
therefore identifies the original observation with
\[
       \jmath\Lambda_k=\mathcal F_0\mathsf A,
       \qquad\ker\mathsf A=K(u)
       \quad\text{in the original coefficient frame.}
       \tag{AKS9}
\]
Every multinomial in AKS8 belongs to its actual ordered words; the
word multiplicity is not inserted a second time. The exact inherited
row Gram of $\mathcal F_0$ is the restriction to $\mathcal I_{k,0}$ of
\[
 \mathcal G_{\nu,\mu}=
 \sum_{\substack{n_{tu}\ge0\\\sum_u n_{tu}=\nu_t\\
                              \sum_t n_{tu}=\mu_u}}
 \frac{k!}{\prod_{t,u}n_{tu}!}\prod_{t,u}
          [5(\delta_{tu}+1)]^{n_{tu}}.
 \tag{AKS10}
\]
This follows by taking the coefficient of $\bar x^\nu y^\mu$ in
$(\bar x^TF^*Fy)^k$, since $F^*F=5(I+\mathbf1\mathbf1^*)$. It is positive
definite and includes every row cross term.

Here is the precise role of the orbit ideal. The map sending a homogeneous
polynomial $f$ to the covector specified by
$\Theta_k(f)(e^{wY}[1])=f(x(w))$ has kernel
$Q\mathbb C[x]_{k-2}$ and is onto $E^\vee$. Indeed equal exponential
indices correspond exactly to monomial differences divisible by
$t_1t_4-t_2t_3$; the displayed Vandermonde proves that there are no other
relations. Its invariant image is $\operatorname{im}\Lambda_k^\vee$.
On the five-member domain, the five $Q_a$ are distinct prime quadrics.
An invariant polynomial divisible by $Q$ is divisible by each $Q_a$,
and therefore by their full product; its quotient is invariant by
cancellation in the polynomial ring. Consequently
\[
\begin{gathered}
 \mathbb C[x]^{\langle D\rangle}\cap (Q)
       =\left(\prod_{a=0}^4Q_a\right)\mathbb C[x]^{\langle D\rangle},\\
 K(u)^\vee\simeq
 \mathbb C[x]_k\big/(\mathbb C[x]_k^{\langle D\rangle}
                              +Q\mathbb C[x]_{k-2}),\\
 b:=\operatorname{rank}\mathsf A
     =d_{5,k,0}-d_{5,k-10,0}=k^2-6k+17,\qquad r=q-b=8k-16.
\end{gathered}\tag{AKS11}
\]
Negative degrees are zero. The count is
$d_{5,d,0}=(\binom{d+3}{3}+4\epsilon_d)/5$,
$\epsilon_d=\mathbf1_{d\equiv0(5)}-\mathbf1_{d\equiv1(5)}$;
subtracting the cubics gives the stated expression for $k\ge10$, and
$k=5,9$ give $12,44$ directly. This is the inherited OQX calculation,
repeated to specify which kernel AKS9 realizes. The source metric is
still to be applied to this kernel, not to the invariant ring itself.

\subsection{A determinant formula requiring no assumed nonzero pivot}
Let $\mathcal G_0$ be the full positive row Gram in AKS10 and set
\[
 W=\mathsf A^*\mathcal G_0\mathsf A,\qquad
 \sigma_j=[z^j]\det(I_q+zW),\qquad
 P_K=\frac1{\sigma_b}\sum_{j=0}^b(-1)^j\sigma_{b-j}W^j.
 \tag{AKS12}
\]
Since $W$ has exactly $b$ positive eigenvalues, $\sigma_b>0$. The scalar
polynomial in AKS12 is one at zero and zero at every nonzero eigenvalue:
it is the normalized product of the factors $\lambda_i-z$ for the $b$
positive eigenvalues. Hence $P_K=P_K^*=P_K^2$ and
$\operatorname{im}P_K=K(u)$. This is the projector derived from the
literal observation, not a freely chosen projection. Replacing
$\mathcal G_0$ by any other positive coordinate Gram leaves this same
coefficient-orthogonal kernel projector, since its image and orthogonal
complement are unchanged. No such replacement is needed here.

At every original endpoint define
\[
 D_{K,N}^{(s)}=[z^r]\det(I_q+zP_KG_N^{(s)}P_K).
 \tag{AKS13}
\]
These numbers are strictly positive for $r>0$ and are one for $r=0$.
To prove their exact metric meaning, choose an orthonormal coefficient
basis $U_K$ of $\operatorname{im}P_K$ just for the proof. In a unitary
extension, $P_KG_NP_K$ has its positive block $U_K^*G_NU_K$ and a zero
block. Thus $D_{K,N}=\det(U_K^*G_NU_K)$. Write the original inclusion
as $I=U_KT$; then $|\det T|^2=\det(I^*I)$ and
\[
 \det H_{K,N}^{(s)}=\det(I^*I)D_{K,N}^{(s)},\qquad
 \boxed{\ F_K^{(s)}=
 \log\frac{D_{K,q-1}^{(s)}D_{K,q}^{(s)}}
               {D_{K,2q-1}^{(s)}D_{K,2q}^{(s)}}.\ }
 \tag{AKS14}
\]
This supplies an equivalent exact coefficient determinant formula for
all the $8k-16$ lost directions, including the entire period dependence.
The original frame factor is present at each endpoint and cancels only
with its four actual signs. It is unnecessary to suppose any particular
minor nonzero. Alternatively, selecting the first independent $r$ columns
of $P_K$ gives an explicit coefficient frame; the determinant formula
avoids even that chart choice.

For $m>1$ replace AKS8--9 by OCS12--18's full confluent matrix
$\mathfrak A_{\nu;(a,b,D)}=i^DD!c_{\nu;a,b,D}$ and its primary isomorphism.
AKS12--14 are then identical with the actual rank, not with the rank in
AKS11. All primary nilpotents and all unit derivatives remain in that
matrix.

\section{A bounded exact remainder operator on the original source}
\label{sec:AKS-operator}

\subsection{The complete source norms and coefficient map}
For either existing source set $M_s=(2\pi)^{s/2}$ and retain
\[
 dm_{0,s}(y)=\frac{M_s2^{b_s}}{4\pi\Gamma(b_s)}
           |\Gamma(b_s/2+iy/2)|^2\,dy,\qquad
 \int e^{ty}dm_{0,s}=M_s(\cos t)^{-b_s},\quad |t|<\pi/2.
 \tag{AKS15}
\]
The supplied IGO1--3 proof derives this transform by the beta integral,
Fourier inversion and a contour shift within $|\operatorname{Im}x|<\pi/2$.
In particular it includes the exact masses $M_1=\sqrt{2\pi}$ and
$M_k=(2\pi)^{k/2}$ and exponential domination on every smaller strip.
The original real monic polynomials satisfy
\[
\begin{gathered}
 p_0=1,\quad p_1=y,\quad
 p_{d+1}=yp_d-d(b_s+d-1)p_{d-1},\\
 \sum_{d\ge0}\frac{p_d(y)}{d!}z^d
      =(1+z^2)^{-b_s/2}e^{y\arctan z},\qquad
 h_d^{(s)}=\|p_d\|^2=M_s d!(b_s)_d.
\end{gathered}\tag{AKS16}
\]
The generating formula follows by differentiating in $z$. To check the
normalization directly, integrate the product of two generating
functions at small real $z,w$. AKS15 and the cosine addition identity
give $M_s(1-zw)^{-b_s}$. Taking any finite pair of derivatives at zero,
justified by the stated exponential domination, gives orthogonality and
the norms in AKS16. Thus the leading phase in
$u_d(S)=p_d((S-c)/i)/\sqrt{h_d^{(s)}}$ is exactly $i^{-d}$.
The same generating function agrees with NIST DLMF 18.23.7 after the
original parameter rescaling; it is derived here from the supplied
recurrence, not substituted for it.

List the original roots
$\omega_\nu=(\lambda_{a,b}-c)/i$ in lexicographic order, repeating each
exactly $e$ times, and put
\[
 P(y)=i^{-q}\chi(c+iy)=\prod_{\nu=1}^q(y-\omega_\nu),\qquad
 \Phi:[f(y)]\mapsto[f((S-c)/i)].
 \tag{AKS17}
\]
The inverse is substitution $S=c+iy$, and
$P((S-c)/i)=i^{-q}\chi(S)$. This is an exact algebra isomorphism, including
the complete confluent orders; there is no monomial comparison divisor.
Let $C_s$ have columns $[u_0],\ldots,[u_{q-1}]$. Its nonzero determinant
is $i^{-q(q-1)/2}\prod_{d<q}(h_d^{(s)})^{-1/2}$. Define
\[
 v_j=C_s^{-1}[u_{q+j}],\qquad
 \mathsf K_j=I_q+\sum_{a=0}^{j-1}v_av_a^*,\qquad 0\le j\le q+1.
 \tag{AKS18}
\]
Parseval in the low-degree original orthonormal polynomial basis proves
\[
 \|v_j\|_2^2=
 \frac{\|\operatorname{rem}_P p_{q+j}\|_{m_{0,s}}^2}
      {M_s(q+j)!(b_s)_{q+j}}.
 \tag{AKS19}
\]
This is the full norm of the actual remainder, not a diagonal entry of
an approximated quotient Gram.

\subsection{Newton interpolation with all roots and their orders}
For $0\le a<q$ define $N_a(y)=\prod_{\nu=1}^a(y-\omega_\nu)$, with
$N_0=1$. For every complex $v$, the unique entire-function remainder is
\[
 \operatorname{rem}_P(e^{vy})
   =\sum_{a=0}^{q-1}[\omega_1,\ldots,\omega_{a+1}]e^{v\cdot}\ N_a(y),
 \qquad
 \left|[\omega_1,\ldots,\omega_{a+1}]e^{v\cdot}\right|
       \le\frac{|v|^a}{a!}e^{R|v|}.
 \tag{AKS20}
\]
Here divided differences at repeated nodes mean Hermite divided
differences, not evaluation with a zero denominator.

A proof covering the confluent case is useful. On the standard
$a$-simplex with barycentric coordinates $\theta_1,\ldots,\theta_{a+1}$,
\[
 [\omega_1,\ldots,\omega_{a+1}]e^{v\cdot}
       =v^a\int_{\Delta_a}e^{v\sum_\nu\theta_\nu\omega_\nu}\,d\theta.
\]
The integral of $\prod\theta_\nu^{n_\nu}$ is
$\prod n_\nu!/(a+\sum n_\nu)!$, as follows by iterated integration of
integer beta integrals. Expanding the exponential therefore gives
$\sum_{n\ge0}v^{a+n}h_n(\omega_1,\ldots,\omega_{a+1})/(a+n)!$.
For monomials, the divided difference of $y^{a+n}$ is exactly $h_n$:
this follows by multiplying the geometric series
$\prod(1-\omega_\nu z)^{-1}$ and the Newton division recurrence.
It is a polynomial identity in the nodes and so also holds at repeated
nodes. This proves the integral formula. Iterated Newton division for
polynomials, followed by the convergent exponential series, proves the
first equality in AKS20, including every required derivative.
Finally the simplex has volume $1/a!$ and
$|\sum\theta_\nu\omega_\nu|\le R$, proving its stated bound.

For any $0<t<\pi/2$, put $\tau=2/t$. Positivity of every term of the
cosh series in AKS15 gives
\[
 \|y^a\|_{m_{0,s}}
 \le\sqrt{M_s}(\cos t)^{-b_s/2}\tau^a a!.
\]
Indeed its square is at most
$M_s(\cos t)^{-b_s}(2a)!t^{-2a}$ and
$(2a)!\le4^a(a!)^2$. Expanding the \emph{whole} Newton polynomial, then
using $|e_j(\omega_1,\ldots,\omega_a)|\le\binom ajR^j$, proves
\[
\begin{split}
 \|N_a\|_{m_{0,s}}
 &\le\sqrt{M_s}(\cos t)^{-b_s/2}
        \sum_{j=0}^a\binom ajR^{a-j}\tau^j j!\\
 &=\sqrt{M_s}(\cos t)^{-b_s/2}\tau^a a!
           \sum_{n=0}^a\frac{(R/\tau)^n}{n!}\\
 &\le\sqrt{M_s}(\cos t)^{-b_s/2}\tau^a a!e^{R/\tau}.
\end{split}\tag{AKS21}
\]
There are no omitted polynomial coefficients. Combining AKS20--21 gives
for every complex $v$
\[
 \|\operatorname{rem}_P(e^{vy})\|_{m_{0,s}}
 \le\sqrt{M_s}(\cos t)^{-b_s/2}e^{R(|v|+t/2)}
                       \sum_{a=0}^{q-1}(2|v|/t)^a.
 \tag{AKS22}
\]

\subsection{Cauchy extraction in the actual Gamma generating function}
Fix $0<z<1$ and write $a_z=\operatorname{arctanh}z$.
For $|w|=z$ the defining power series gives
$|\arctan w|\le a_z$ and $|1+w^2|\ge1-z^2$.
Apply the literal finite-dimensional remainder map to AKS16 and use
Cauchy's formula for its $d$th coefficient. The norm triangle inequality
on that circle and AKS22 give
\[
\begin{split}
 \frac{\|\operatorname{rem}_P p_d\|_{m_{0,s}}}{d!}
 \le{}&z^{-d}(1-z^2)^{-b_s/2}\sqrt{M_s}(\cos t)^{-b_s/2}\\
 &\mathrel{}\times e^{R(a_z+t/2)}
                    \sum_{a=0}^{q-1}(2a_z/t)^a.
\end{split}\tag{AKS23}
\]
Cauchy's formula is legitimate for the finite remainder vector: its
entries are holomorphic for $|w|<1$ by the actual matrix generating
identity OPG7. The source norm is taken only after this finite-vector
identity; there is no unjustified integral of a nonintegrable generating
function on the large circle.

Substitution of AKS23 into AKS19 retains the factorial denominator and
cancels the two occurrences of the original mass, giving
\[
 \boxed{\ 
 \|v_j\|_2^2\le\frac{d!}{(b_s)_d}z^{-2d}
 [\cos t(1-z^2)]^{-b_s}e^{2R(a_z+t/2)}
 \left[\sum_{a=0}^{q-1}(2a_z/t)^a\right]^2,
 \quad d=q+j.\ }
 \tag{AKS24}
\]
Adding the \emph{complete} rank-one positive matrices in AKS18 now proves
$\mathsf K_j\preceq \mathcal Z_{s,j}I_q$ after setting $z=4/5$ and $t=t_q$.
Indeed $a_{4/5}=\log3$, $1-z^2=9/25$, and $\cos t_q=\sin(1/q)$, so
the exact scalar obtained is AKS3. The high index $j=q+1$ includes the
original polynomial of degree $2q$.

There is also a root-by-root finite refinement. Put
$\mu_{s,d}=\int y^{2d}dm_{0,s}$,
$R_a=\max_{\nu\le a+1}|\omega_\nu|$, and
\[
\begin{gathered}
 B_{s,a}^{\rm root}=\sum_{d=0}^a
     |e_{a-d}(\omega_1,\ldots,\omega_a)|\sqrt{\mu_{s,d}},\\
 J_s^{\rm root}(z)=\frac{(1-z^2)^{-b_s}}{M_s}
 \left[\sum_{a=0}^{q-1}\frac{a_z^a}{a!}e^{a_zR_a}
                      B_{s,a}^{\rm root}\right]^2.
\end{gathered}\tag{AKS25}
\]
Replace the scalar multiplying
$d!z^{-2d}/(b_s)_d$ in AKS24 by $J_s^{\rm root}(z)$ to get a valid
sharper bound, by the same proof before AKS21 is relaxed. Every moment
is exactly $M_s(2d)![w^{2d}](\cos w)^{-b_s}$. Thus this refinement keeps
every actual root coefficient and a finite original moment formula.
AKS21 bounds it by the explicit scalar in AKS24. The proof never replaces
the source determinant by a product of diagonal moments: AKS24 is a
full-vector operator bound, and the next step restricts the inverse of
the full covariance.

\section{The four original kernel Grams and the evaluated scale}
\label{sec:AKS-volume}

The original orthonormal source map at degree $q+j-1$ is
$C_s[I_q,v_0,\ldots,v_{j-1}]$. Its minimum lift of $C_sx$ is
$[I_q,v_0,\ldots,v_{j-1}]^*\mathsf K_j^{-1}x$. It has that image and is
orthogonal to the full relation kernel, so Pythagoras proves
\[
 G_{q+j-1}^{(s)}=C_s^{-*}\mathsf K_j^{-1}C_s^{-1},\qquad
 H_j:=H_{K,q+j-1}^{(s)}=T_s^*\mathsf K_j^{-1}T_s,\quad T_s=C_s^{-1}I.
 \tag{AKS26}
\]
This verifies the source-minimum step in the coefficient determinant
AKS13--14. Neither a cross block nor any original relation is removed.

Since $I_q\preceq\mathsf K_j\preceq \mathcal Z_{s,j}I_q$, inversion and compression
give $\mathcal Z_{s,j}^{-1}H_0\preceq H_j\preceq H_0$, where $H_0=T_s^*T_s$.
Define $d_j=-\log\det(H_0^{-1/2}H_jH_0^{-1/2})$. Congruence and all
$r$ eigenvalues give
\[
 0=d_0\le d_1\le\cdots\le d_{q+1},\qquad
 0\le d_j\le r\mathcal L_{s,j},\qquad
 F_K^{(s)}=d_q+d_{q+1}-d_1.
 \tag{AKS27}
\]
The identity uses precisely the four endpoints in AKS2.
In particular it proves AKS4, as well as the sharper interval
\[
 \boxed{\quad d_1\le F_K^{(s)}\le
             r(\mathcal L_{s,q}+\mathcal L_{s,q+1})-d_1.\quad}
 \tag{AKS28}
\]
For a directly calculable expression for its lower endpoint set
$V=\|v_0\|^2$ and
$\theta=v_0^*T_s(T_s^*T_s)^{-1}T_s^*v_0$. Then
\[
 0\le\theta\le V,\qquad
 d_1=\log\frac{1+V}{1+V-\theta}.
 \tag{AKS29}
\]
Indeed the inverse of $I+v_0v_0^*$ is
$I-v_0v_0^*/(1+V)$; take its full compressed determinant. For $r=0$
set $\theta=d_1=0$. In GMB/OPR notation $V=t_0$ and $V-\theta=\xi_0$,
so this is exactly the first original section-corrected kernel increment,
not an independently chosen kernel observation.

Here is an evaluated finite remainder for AKS6. Put
\[
\begin{gathered}
 a_\infty=\frac{4\log3}{\pi},\qquad
 c_0=\log\frac{25\log3}{4\pi}>0,\\
 E_{s,q}(R)=\log2+\frac4{\pi-2/q}
  +\frac{s}{2}\log\frac{25\pi q}{18}
  +2R\left(\log3+\frac\pi4\right)\\
 \hspace{23mm}+\log((q+1)(4q+1))-2\log(a_\infty-1).
\end{gathered}\tag{AKS30}
\]
Then for both source orders, every $q\ge36$, and the same actual kernel,
\[
 \boxed{\quad
 \mathcal L_{s,q},\mathcal L_{s,q+1}\le2c_0q+E_{s,q}(R),\qquad
 0\le F_K^{(s)}\le4rc_0q+2rE_{s,q}(R).
 \quad}\tag{AKS31}
\]
To verify every factor, $b_s\ge1/2$ gives
\[
 \frac{d!}{(b_s)_d}\le\frac{4^d}{\binom{2d}{d}}\le2d+1.
\]
The last inequality uses that the central binomial is the largest of the
$2d+1$ positive coefficients whose sum is $4^d$. Thus the sum in AKS3
is at most $(q+1)(4q+1)(5/4)^{4q}$. Also
$\mathfrak s_q\le \mathfrak a_q^q/(a_\infty-1)$ and
\[
 2q\log \mathfrak a_q\le2q\log a_\infty+\frac4{\pi-2/q},\qquad
 \sin(1/q)\ge\frac2{\pi q}.
\]
The first follows from $-\log(1-x)\le x/(1-x)$ at $x=2/(\pi q)$;
the second is the chord bound for sine on $[0,\pi/2]$.
Use $t_q\le\pi/2$ in the remaining exponential. The order-$q$ terms
combine to
$4q\log(5/4)+2q\log a_\infty=2c_0q$.
The upper bound for the logarithm of the added positive scalar is
positive, so $\log(1+e^x)\le\log2+\max(0,x)$ proves the first bound in
AKS31. AKS28 proves the second.

On the original simple family $q=(k+1)^2$, $R=O_h(k)$ and
$r=8k-16$. Consequently $E_{1,q}(R)=O_h(k+\log q)$ and
$E_{k,q}(R)=O_h(k\log q)$. Dividing AKS31 by $kq$ proves AKS6, with
$32c_0=\kappa_*$. In fact the fixed-source upper bound is
\[
 F_K^{(1)}\le(32k-64)c_0q+(16k-32)E_{1,q}(R),
 \tag{AKS32}
\]
whose explicit second term is $O_h(k^2)$. The tensor-source version has
an $O_h(k^2\log q)$ upper remainder. These are remainders of upper
bounds, not assertions about the actual signed error or leading value.
The proof is uniform in all the period parameters in the stated domain:
no norm or condition number of a period matrix occurs in AKS30.

\paragraph{The precise earlier loss.}
OKV15 had already retained
$\|v_j\|^2\le q^22^{6q}(2D_*q)^{2d}/(d!(b_s)_d)$, $d=q+j$.
Replacing its denominator by $1/2$ caused the old $q\log q$ exponent.
Even without AKS20--24 one obtains the valid repair
\[
 \|\mathsf K_{q+1}\|\le
 1+(q+1)q^22^{6q}\frac{(4D_*q)^{4q}}{(4q)!},\qquad D_*\ge3.
 \tag{AKS33}
\]
Here $d!(b_s)_d\ge(2d)!/4^d$ and the remaining factorial ratio increases
up to $d=2q$ for $D_*\ge3$. The elementary inequality
$\log n!\ge n\log n-n+1$ already makes the logarithm of AKS33 $O_h(q)$.
The Newton calculation strengthens this repair to the packet-independent
limsup constant AKS5, while retaining $R$ in its finite remainder.


\paragraph{Exact enclosure of the displayed numerical ceiling.}
For $x>1$ and $z=(x-1)/(x+1)$ the integrated geometric series gives
\[
 0\le \log x-2\sum_{j=0}^{39}\frac{z^{2j+1}}{2j+1}
      \le\frac{2z^{81}}{81(1-z^2)}.
\]
For $0<x<1$, the alternating sum through $j=39$ for $\arctan x$
has error between zero and $x^{81}/81$. The tangent addition formula,
with the angle in $(0,\pi/2)$, proves
$\pi=16\arctan(1/5)-4\arctan(1/239)$: the tangent of
$4\arctan(1/5)-\arctan(1/239)$ is exactly one. Use the resulting
rational lower and upper bounds for $\pi$ and for $\log3$, then apply
the first displayed bound to the two endpoints of
$25\log3/(4\pi)$. Multiplication by $32$ gives, with outward decimal
rounding, the interval
\[
 \kappa_*\in[25.0207809764994871,\ 25.0207809764994872].
\]
The accompanying rational-arithmetic implementation checks these
inequalities by integer cross multiplication, rather than a floating-point
logarithm. This proves the less finely rounded strict interval AKS5.

\section{The complete arithmetic and mixed receivers}
\label{sec:AKS-transfer}

\subsection{The unchanged finite arithmetic allowances}
Use exactly AMT4's constants. To prevent a loss hidden in order notation,
write them here:
\[
\begin{gathered}
 \alpha=\pi/2,\ p=8m-9/2,\ B_h=42+8m,\ M=21+4m,\\
 \vartheta_h=\int_{-1}^1w_h(y)dy,\quad
 M_\alpha=\int e^{\alpha y}w_h(y)dy,\quad
 w_h(y)=|v_h(1/2+iy)|^2/(2\pi),\\
 c_\Gamma=\sqrt2e^{-7/6},\quad C_\Gamma=\sqrt{2\sqrt5}e^{2/3},\\
 a_k=\frac{c_h\vartheta_h^{k-3}e^{-\alpha(k-3)}(k-2)^{-B_h}}
                    {C_\Gamma2^{B_h/2}},\quad
 A_k=\frac{C_h^{\rm tilt}}{c_\Gamma}k^{p+1}M_\alpha^{k-1},\\
 D_N=[1+4(N+M)^2]^M,\quad L_N^{\rm ar}=\log(A_kD_N/a_k),\\
 E_k^+=L_{q-1}^{\rm ar}+L_q^{\rm ar},\qquad
 E_k^-=L_{2q-1}^{\rm ar}+L_{2q}^{\rm ar}.
\end{gathered}\tag{AKS34}
\]
The constants $c_h,C_h^{\rm tilt}$ are the original TW/BT envelope
constants, not new choices. Their actual inequalities and complete
source construction are preserved in AMT4--6 and the current source
closure. Minimize AMT6 on the same remainder fibre and then restrict to
$I(K)$. AMT9--20 gives, with the original arithmetic source $w_h^{*k}$,
\[
 -rE_k^+\le F_K^{\rm ar}-F_K^{(1)}\le rE_k^-.
\]
Let $\underline U_s=d_1^{(s)}$ and
$\overline U_s=r(\mathcal L_{s,q}+\mathcal L_{s,q+1})-d_1^{(s)}$. The new finite conclusion is
\[
 \boxed{\quad
 \max\{0,\underline U_1-rE_k^+\}\le F_K^{\rm ar}
                  \le\overline U_1+rE_k^-.
 \quad}\tag{AKS35}
\]
This improves AMT37's absolute upper bound. For the sharper correlated
version retain $d_{E,N}=q\log A_k-
\log(\det G_N^{\rm ar}/\det G_N^{(1)})$ and
$c_{X,N}=\min\{r_XL_N^{\rm ar},d_{E,N}\}$, where $r_K=r$, $r_B=q-r$.
Put $C_X^-=c_{X,q-1}+c_{X,q}$,
$C_X^+=c_{X,2q-1}+c_{X,2q}$. AMT18 then gives
\[
\begin{split}
 \max\{0,\underline U_1-C_K^-,
                  \underline U_1+\delta_k^\sigma-C_B^+\}
 &\le F_K^{\rm ar}\\
 &\le\overline U_1+
            \min\{C_K^+,\delta_k^\sigma+C_B^-\}.
\end{split}\tag{AKS36}
\]
One may replace $d_{E,N}$ in the caps by the original full-source
$\Xi_{k,N}$ or its proved finite majorant. The scalar
$\delta_k^\sigma=\mathcal B_k^{\rm ar}-\mathcal B_k^\sigma$ retains
exactly AMT21:
\[
\begin{gathered}
 \max\{-qE_k^+,-\Xi_{k,q-1}-\Xi_{k,q}\}
       \le\delta_k^\sigma\\
       \le\min\{qE_k^-,\Xi_{k,2q-1}+\Xi_{k,2q}\}.
\end{gathered}\tag{AKS37}
\]
Since $E_k^\pm=O_h(k+\log q)$ by the exact AMT22 expression, AKS35 and
AKS32 prove the arithmetic version of AKS6, with an explicit
$O_h(k^2)$ upper remainder. Nonnegativity follows from the actual
nested source minima, as in AMT15, not from comparison with a possibly
negative lower endpoint.

\subsection{The original Gamma multiplier, not a replacement higher source}
Retain $t_j=2j+1/2$, $f_l(S)=\prod_{j=0}^{l-1}(S-c-t_j)$ and
\[
 dm_{0,k}=\beta_k|f_l(c+iy)|^2d\sigma,\quad
 \beta_k=\frac{(2\pi)^{k/2}}{\sqrt2\Gamma(k/2)},\quad
 T_l=\prod_{j=0}^{l-1}t_j^2,\quad
 U_{l,N}=\prod_{j=0}^{l-1}[4(N+j+1)^2+t_j^2].
 \tag{AKS38}
\]
The exact source map is $\mathcal P_N\to f_l\mathcal P_N\subset
\mathcal P_{N+l}$, $P\mapsto f_lP$, with inverse division on this
specified subspace. Let $\mathsf E_\chi$ be the complete divided-jet
map and $U_f$ its Leibniz multiplication matrix,
$(U_f)_{r,a}=f_l^{(r-a)}(\lambda)/(r-a)!$ for $r\ge a$.
AMT25--28 gives the precise transported observation and kernel
\[
\begin{gathered}
 \widetilde\Lambda=\Lambda\mathsf E_\chi^{-1}U_f^{-1},\quad
 \widetilde I=U_f\mathsf E_\chi I,\quad
 \ker\widetilde\Lambda=U_f\mathsf E_\chi K,\\
 C_N=\mathscr K_{\chi,N+l}
   -\mathscr K_{\chi,f;N+l}\mathscr K_{f,N+l}^{-1}
                               \mathscr K_{f,\chi;N+l},\\
 H_{K,N}^{(k)}=\beta_k\widetilde I^*C_N^{-1}\widetilde I.
\end{gathered}\tag{AKS39}
\]
Each covariance is the original finite sum of full polynomial jet outer
products divided by its original monic norm. These formulas retain the
extra $f$ constraints and their full cross covariance. In particular
$M_fK=K$ is not assumed. AKS26 calculates the same tensor metric directly
in its own orthonormal source; the two representations agree by their
identical source norm and identical constrained minimum.

Define
\[
 w_N=\log(U_{l,N}/T_l),\quad
 \Omega^{\rm low}=w_{q-1}+w_q,\quad
 \Omega^{\rm high}=w_{2q-1}+w_{2q}.
\]
The already proved same-coefficient comparison AMT42--44 retains all
mass factors before cancellation and gives
\[
 -r\Omega^{\rm low}\le\Delta_K^\Gamma:=F_K^{(1)}-F_K^{(k)}
                         \le r\Omega^{\rm high}.
 \tag{AKS40}
\]
It therefore sharpens the two independent absolute bounds by
\[
\begin{split}
 \max\{\underline U_k,\underline U_1-r\Omega^{\rm high},0\}
       &\le F_K^{(k)}\\
       &\le\min\{\overline U_k,\overline U_1+r\Omega^{\rm low}\}.
\end{split}\tag{AKS41}
\]
The original finite estimate
$0\le w_N\le l\log[1+16(N+l)^2]$ proves that AKS40 is $o(kq)$ on the
actual rank-$r$ domain. Along with AKS35 this preserves the completed
fact that all three source kernels differ by $o(kq)$. AKS4, not AKS40,
is the new absolute-scale calculation.

\subsection{Boundary, QGQ and arithmetic receiving identities}
At each endpoint retain the original minimum section
$S_N^{(s)}=(G_N^{(s)})^{-1}\Lambda^*
(\Lambda(G_N^{(s)})^{-1}\Lambda^*)^{-1}$. The full Schur identity is
\[
 \det(I^*G_N^{(s)}I)\det Q_N^{(s)}
       =|\det[I,S_*]|^2\det G_N^{(s)},\qquad
 Q_N^{(s)}=(\Lambda(G_N^{(s)})^{-1}\Lambda^*)^{-1}.
 \tag{AKS42}
\]
Thus the new finite absolute boundary interval is
\[
 \max\{0,\mathcal B_k^{(s)}-\overline U_s\}
       \le F_B^{(s)}\le\mathcal B_k^{(s)}-\underline U_s.
 \tag{AKS43}
\]
The baselines in this formula are the \emph{complete} original BSL/HBL
objects, with their full finite relation errors. No $q^2$ coefficient is
assigned to $K$; its proved coefficient there remains zero.

Write $\mathcal H_k=\mathcal B_k^\sigma-\mathcal B_k^0=-\mathfrak T_k$.
The exact full identities still read
\[
\begin{gathered}
 F_B^{(1)}-F_B^{(k)}=\mathcal H_k-\Delta_K^\Gamma,\\
 \boxed{\quad\mathcal B_k^{\rm ar}
   =F_K^{(k)}+F_B^{(k)}+\mathcal H_k+\delta_k^\sigma
   =F_K^{(1)}+F_B^{(1)}+\delta_k^\sigma.\quad}
\end{gathered}\tag{AKS44}
\]
The Gamma interval is not re-estimated in this argument. For either
complete QRI1 or QRI2 interval $[\omega^-,\omega^+]$ for its already
calculated centre $W_k$, retain its exact source errors
$e_0\in[-a_0,b_0]$ and $e_q+e_{q+1}\in[-b_k^+,b_k^-]$. Then
\[
 \underline{\mathcal H}_k=-\omega^+-b_k^--a_0,
 \qquad\overline{\mathcal H}_k=-\omega^-+b_k^++b_0,
 \qquad \mathcal H_k\in
       [\underline{\mathcal H}_k,\overline{\mathcal H}_k].
 \tag{AKS45}
\]
These are exactly QRI3, not new unspecified error choices. The complete
QRI1--9, QGT and QGQ definitions remain verbatim in each updated receiver;
the new interval is intersected with all their retained tighter finite
intervals. In particular the simple-quartet nonzero error radius stays
\[
 \limsup\left|
 \frac{\mathcal H_k+lq\mathcal C_\Gamma}{q}
       +\frac{\partial_\alpha\mathcal L(2,\pi)}8
       -\frac{\log2}{4}\right|\le\frac{A_{\rm QGT}}{\sqrt2}.
 \tag{AKS46}
\]
No zero-radius or $o(q)$ arithmetic assertion is introduced.

For example the complete finite absolute defect relative to the tensor
boundary, with the actual Gamma and arithmetic scalars retained, is
\[
 \underline U_k+\mathcal H_k+\delta_k^\sigma
 \le\mathcal B_k^{\rm ar}-F_B^{(k)}
 \le\overline U_k+\mathcal H_k+\delta_k^\sigma.
 \tag{AKS47}
\]
Inserting AKS41 makes this sharper. Replacing $\mathcal H_k$ by the
appropriate endpoints of AKS45 and $\delta_k^\sigma$ by AKS37 gives its
fully inherited finite scalar enclosure with the same signs.
Also $F_B^{\rm ar}=\mathcal B_k^\sigma+\delta_k^\sigma-F_K^{\rm ar}$;
subtract the interval AKS35 or AKS36 to get its new absolute boundary
interval. This retains the exact correlations rather than taking
independent hypothetical source errors.

The existing limit is
$\mathcal H_k/(kq)\to-\mathcal C_\Gamma/4$ and
$\delta_k^\sigma/(kq)\to0$. Therefore the new interval implies
\[
\begin{gathered}
 \liminf_k\inf_{|u|\ge R_*}
 \frac{\mathcal B_k^{\rm ar}-F_B^{(k)}}{kq}
          \ge-\mathcal C_\Gamma/4,\\
 \limsup_k\sup_{|u|\ge R_*}
 \frac{\mathcal B_k^{\rm ar}-F_B^{(k)}}{kq}
          \le\kappa_*-\mathcal C_\Gamma/4.
\end{gathered}\tag{AKS48}
\]
All original branches are included in the infima and suprema. This does
not decide its sign. The already proved boundary \emph{difference}
limits of AMT49 are unchanged. Every subsequential limit of the absolute
normalized kernel lies in $[0,\kappa_*]$ and is shared by its three source
versions when they are evaluated on the same sequence of period
parameters. Existence of a single limit is not proved by boundedness.

Finally the HC residual remains exactly
$\mathcal R_k=\mathcal B_k^{\rm ar}-4q\log(D_hk)$, with its original
$D_h$ and original nonnegative contraction, phase and trace terms.
AKS44 changes none of them. The supplied positive $q^2$ baseline and
its complete Gamma correction are not re-presented as new results here.

\section{Actions, support faces, and the actual arithmetic class}
\label{sec:AKS-actions}

The coefficient formula AKS12 supplies all four blocks of the same sum
action: $P_KYP_K$, $P_KY(I-P_K)$, $(I-P_K)YP_K$ and
$(I-P_K)Y(I-P_K)$. There is no claim that either off-diagonal block
vanishes. In the actual minimum-section coordinates the corresponding
formula remains
\[
 [I,S_N]^{-1}Y[I,S_N]=
 \begin{pmatrix}
 H_{K,N}^{-1}I^*G_NYI&H_{K,N}^{-1}I^*G_NYS_N\\
 \Lambda YI&\Lambda YS_N
 \end{pmatrix}.
 \tag{AKS49}
\]
Its proof is multiplication by the inverse coordinate maps
$H_{K,N}^{-1}I^*G_N$ and $\Lambda$. In the invariant-ring model, the same
action is the derivation
$\mathcal Df=(\mathsf H\Omega\mathsf H^{-1}x)\cdot\nabla f$,
$\Omega=\operatorname{diag}(\gamma-i\delta,-\gamma-i\delta,
\gamma+i\delta,-\gamma+i\delta)$, with $\mathcal DQ=0$.
For the actual conductor $\mathcal A_Q=\prod_{a=1}^4Q_a$, its full defect
is $\mathcal D\mathcal A_Q=\sum_{a=1}^4(\mathcal DQ_a)
\prod_{b\ne a,\,1\le b\le4}Q_b$, exactly OQC22. Thus this coefficient
realization preserves the original action leaving the observed dual.

In particular the section-corrected vector is still OPR2's
\[
 C_s r_j^{\rm ker}=e_{q+j}-S_{q+j-1}^{(s)}\Lambda e_{q+j}
   =I\bigl(\zeta_{q+j}-\kappa S_{q+j-1}^{(s)}\Lambda e_{q+j}\bigr),
\]
with its full denominator $1+\xi_j$ in the kernel increment. The operator
proof estimates the complete covariance and its actual restriction, so
none of these coupled terms is replaced by a raw fixed-section vector.
The all-iterate kernel stays
$\bigcap_{a=0}^{q-1}\ker(\Lambda Y^a)$; nothing in AKS4 declares it zero.

On the existing arithmetic coefficient realization retain
$\mathcal V_{h,k}P=P(D_1+\cdots+D_k)F_h^{\otimes k}$ and
$J^{(k)}\mathcal V_{h,k}=\eta\operatorname{rem}_\chi$, where
$\eta=U^{\otimes k}\iota$ is the original injection. AKS17 and the
Newton coefficients compute precisely this same remainder. Any two
polynomial lifts differ by $\chi Q$ in the original relation space,
so their original arithmetic cohomology classes agree. The specified
primitive of the difference is EP.48, with the full differential and
cochain signs proved in AKP3--5 below; it depends on the lift difference.
The source comparison to $w_h^{*k}$ is the identity on original
polynomial coefficients, followed by the actual minima as in AMT8;
it is not a replacement arithmetic source.

The scalar object stays $G(R)=\{\tau\}\sqcup R^\bullet$, with supported
zero $0_R^\bullet$, injective amplitude/support pair, and infinite
integer target. Every coefficient map above is applied on each original
present support label and outer leg. It commutes with the specified
coefficient inclusions because both compositions apply the same linear
map to the same coefficient; a present zero is not converted into
absence. No new base object, purity condition, or identification of the
sum action with a Frobenius action is used. The actual distinguished
class and its higher-jet directions remain in the full fibre, even
when the observation kills them.

\section{Exact period domain and the interior alternative}
\label{sec:AKS-radius}

For reference the inherited radius used above is completely specified
as follows; these are AMT33--34/RPC17--23, with no new constants. Put
\[
\begin{gathered}
 \beta=2(\gamma^2-\delta^2),\quad\eta=(\delta^2+\gamma^2)^2,\quad
 a=v_h(\rho)\ne0,\\
 c_3=(a^{-2}-\bar a^{-2})/(4i\delta\gamma),\quad
 c_1=(a^{-2}+\bar a^{-2})/2-(\delta^2-\gamma^2)c_3,\\
 g_j=5^{j/5-1}e^{\pi ij/5}\Gamma(j/5),\quad
 g_+=\max_j|g_j|,\quad g_-^{-1}=\max_j|g_j|^{-1},\quad
 \kappa=\max_{j<4}|g_{j+1}/g_j|,\\
 E_j^{\rm ray}=\frac85\int_0^\infty r^j(r^3/3+r)
       e^{-r^5/5+r^3/3+r}dr,\quad
 C_T=2\left(\sum_{j=0}^3(E_j^{\rm ray})^2\right)^{1/2},\\
 C_K=4C_Tg_-^{-1}+2g_+g_-^{-1}+2g_+g_-^{-2}C_T,\\
 C_{K,3}=C_K(3\kappa^2+3\kappa C_K+C_K^2),\quad
 \epsilon_* =\min\{1,(2C_Tg_-^{-1})^{-1}\},\quad B_* =\beta+\eta.
\end{gathered}\tag{AKS50}
\]
Then
\[
R_*=
\begin{cases}
\displaystyle\left[\max\left\{1,\frac{B_*}{\epsilon_*},
\frac{2(|c_3|C_{K,3}B_*+|c_1|(\kappa+C_K))}{|c_3g_4/g_1|}
\right\}\right]^{5/2},&c_3\ne0,\\[4pt]
\displaystyle\left[\max\left\{1,\frac{B_*}{\epsilon_*},
\frac{2C_KB_*}{|g_2/g_1|}\right\}\right]^{5/2},&c_3=0.
\end{cases}\tag{AKS51}
\]
The actual unit implies $(c_1,c_3)\ne(0,0)$ and all integrals converge.
RPC's full period comparison proves a nonzero $(4,1)$ or $(2,1)$
commutator entry, respectively, for $|u|\ge R_*$ on every branch. It
therefore proves the five-member orbit needed in AKS11.

Inside this radius the same coefficient determinant AKS12--14 and the
operator bound AKS24 remain valid. The orbit is decided by the original
noninvariant quadratic coefficients, not by the radius alone. Since
$Q$ is nondegenerate, its projective orbit has either one or five
members. A one-member orbit must have character one: its character has
both fourth and fifth powers equal to one, by the quadratic determinant
and $D^5=I$. In that case $Q=a'x_1x_4+b'x_2x_3$ with $a'b'\ne0$, and
\[
 b=d_{5,k,0}-d_{5,k-2,0}
   =\left\lfloor\frac{(k+1)^2}{5}\right\rfloor
           +\mathbf1_{k\equiv0\text{ or }3\ (5)},\qquad r=q-b.
 \tag{AKS52}
\]
If a noninvariant coefficient is nonzero, AKS11 applies instead. These
are the complete OQC alternatives. The bound AKS31 always uses that
actual $r$; it does not extend the $8k-16$ specialization to an
unevaluated exceptional point. This calculation does not assert that
the interior exceptional set is empty.

\section{Verification scope and the remaining absolute coefficient}
\label{sec:AKS-scope}

The new mathematical content is AKS20--32 and its application to the
actual coefficient determinant AKS12--14, followed by the finite
receiving improvements AKS35--48. The complete source and orbit
identifications used in that application are the supplied OCS/OQC/OQX,
RPC, OPG/OPR and AMT results; their exact predecessors are preserved.
AKS33 pinpoints a repair obtainable immediately from the earlier proof,
whereas the Newton bound supplies the sharper explicit constant.

The bound \emph{at} scale $kq$ is now finite and uniform. What is not
proved is a matching lower coefficient, a unique limit for
$F_K^{(1)}/(kq)$, or a sign for the interval AKS48. Those questions involve
the orientation of the original period kernel within the inverse
remainder covariance; replacing that compression by a scalar bound
loses precisely that information. This is the finite-bounds outcome
allowed by the current request, not a claim to have evaluated an
asymptotic coefficient. The source differences are already completed
inputs and are not renamed as the remaining problem.

The accompanying exact checks distinguish finite algebraic diagnostics
from the analytic proof for the full family. They do not assert that a
diagnostic root is a zero of $\xi$, that a diagnostic matrix equals an
actual period matrix, or that a new Lean theorem has been checked.
No RH conclusion is drawn from this kernel bound.

% END CURRENT AKS

\clearpage
% BEGIN CURRENT AKJ
% Root continuation: full original residue, conductor and metric maps.
\section{The new ceiling in the original coefficient frame}
\label{sec:AKJ}
The complete incoming AKS proof sharpens the previously proved KAF bound.
Both estimates apply to the same original remainder and the same kernel.
We retain all their finite constants. In particular, for $m=1$,
$k=4l+1$, $l\ge1$, $q=(k+1)^2$, $c=k/2$, $0<\delta<1/2$,
$\gamma>2$, and every original branch with $|u|\ge R_*$, put
$r=8k-16$. The AKS constant is
\[
 \kappa_*=32\log\frac{25\log3}{4\pi}
 \in(25.02078097649,25.02078097650).
 \tag{AKJ1}
\]
AKS20--32 proves the full finite operator inequality from confluent
Newton division and coefficient extraction. Its parameters $t$ and $z$
are extraction radii in the existing source, whose orders remain $1,k$.
For $s\in\{1,k\}$ define, using the exact KAF and AKS quantities,
\[
 \begin{split}
 L_s&=d_1^{(s)},\\
 U_s&=\min\{2r\log Z_{k,1},\,
 r(\mathcal L_{s,q}+\mathcal L_{s,q+1})-d_1^{(s)}\}.
 \end{split}\tag{AKJ2}
\]
The first entry of the minimum is KAF's factorial estimate; the second
is AKS28, with its original first compressed increment retained. Thus
$L_s\le F_K^{(s)}\le U_s$. This pointwise minimum retains whichever
finite estimate is sharper at the actual original parameters. It does
not assume that the asymptotically smaller ceiling wins at every degree.

The two exact original coordinate constructions are related as follows.
Use OKM's invertible ordered Vandermonde $V_\beta$ and its actual
quotient matrix $\pi$. Put $I_\pi=V_\beta^{-1}\pi^{\mathsf T}$.
AKS12 constructs the coefficient-orthogonal projection $P_K$ from the
actual observation and its complete ordered-word Gram. Directly,
\[
 P_K=I_\pi(I_\pi^*I_\pi)^{-1}I_\pi^*,\qquad
 D_{K,N}^{(s)}
 =\frac{\det(\overline\pi(R_N^{(s)})^{-1}\pi^{\mathsf T})}
        {\det(I_\pi^*I_\pi)}.
 \tag{AKJ3}
\]
Indeed $I_\pi$ is injective with image $K$ by OKM5. The first displayed
matrix is Hermitian and idempotent, fixes this image, and kills its
coefficient-orthogonal complement, so it is exactly AKS12's projection.
Apply AKS14 with inclusion $I_\pi$ and then OKM8--9 to obtain the second
identity. The positive denominator is independent of endpoint and
source. Its exponent is $1+1-1-1=0$ in the four-endpoint quotient.
This proves the exact morphism between the incoming projector
determinant and the existing conductor-frame determinant, retaining the
original coefficient metric and every frame factor.

With the original AMT correlated caps $C_X^\pm$ and scalar
$\delta_k^\sigma$, the following joint intervals hold:
\[
\begin{split}
 \max\{0,L_1-C_K^-,L_1+\delta_k^\sigma-C_B^+\}
 &\le F_K^{\rm ar}\\
 &\le U_1+\min\{C_K^+,\delta_k^\sigma+C_B^-\},\\
 \max\{L_k,L_1-r\Omega^{\rm high},0\}
 &\le F_K^{(k)}
 \le\min\{U_k,U_1+r\Omega^{\rm low}\}.
\end{split}\tag{AKJ4}
\]
For proof, add each retained AMT lower or upper source-difference bound
to the corresponding endpoint of AKJ2, and intersect the resulting
intervals. In the first line the upper and lower source differences
are the same AMT18 correlated caps used in AKS36; no independent errors
have replaced their actual values. Every boundary interval is obtained
by subtracting these kernel intervals from the exact original total
$\mathcal B_k^a=F_K^a+F_B^a$. Consequently all three original kernels
have uniform limsup at most $\kappa_*$ after division by $kq$.
The source differences remain $o(kq)$. A matching leading value is
still to be calculated from these same matrices.

\section{Residue detection expressed in every original primary coefficient}
Retain now every original multiplicity $m\ge1$, $e=1+k(m-1)$,
$q=e(k+1)^2$, and the full polynomial
\[
 \chi(S)=\prod_{a,b=0}^k(S-\beta_{ab})^e,
 \quad \beta_{ab}=c+(2a-k)\delta+i(2b-k)\gamma,
 \quad E=\mathbb C[S]/(\chi).
 \tag{AKJ5}
\]
All distinct pairs give distinct roots. For a root $\beta$ write
$z=S-\beta$, $h_\beta(S)=\chi(S)/(S-\beta)^e$, and let
$\varepsilon_\beta$ be its original CRT idempotent. Define exactly
$\ell([f])=[S^{q-1}]\operatorname{rem}_\chi f$.
For every polynomial class the primary expansion gives
\[
 \ell(f)=\sum_\beta[z^{e-1}]
       \frac{f(\beta+z)}{h_\beta(\beta+z)}.
 \tag{AKJ6}
\]
The quotient on the right is a truncated power series through degree
$e-1$. To prove the identity, partial-fraction division of $f/\chi$
gives at each finite pole the coefficient of $(S-\beta)^{-1}$ shown
on the right. Expanding those principal parts at infinity, only their
first-order terms contribute to $S^{-1}$. Polynomial terms contribute
nothing. For the unique representative of degree below $q$, the same
$S^{-1}$ coefficient is its coefficient of $S^{q-1}$. This proves AKJ6
without deleting any repeated-root term. Since $h_\beta(\beta)\ne0$,
the bilinear pairing $(f,g)\mapsto\ell(fg)$ is nondegenerate: if the
first nonzero local coefficient of $f$ is $f_tz^t$, take
$g=\varepsilon_\beta z^{e-1-t}h_\beta$ to obtain $\ell(fg)=f_t\ne0$.

Let $\lambda_\nu$ be the specified zero-charge coefficient rows of
the original OCS observation, so the injection $\mathcal F_0$ gives
$\jmath\Lambda=\mathcal F_0\mathfrak A\mathcal H^{-1}$.
Retain the exact entries
\[
 a_{\nu,\beta,D}:=\lambda_\nu(\varepsilon_\beta z^D)
       =i^DD!c_{\nu;abD},\qquad 0\le D<e.
 \tag{AKJ7}
\]
Here $c_{\nu;abD}$ is the full OCS multinomial in the original period
and unit jets; $\beta=\beta_{ab}$. It is not an assigned coefficient.
The residue representative of each row has the explicit formula
\[
 A_\nu=\left[\sum_\beta\varepsilon_\beta h_\beta(S)
           \sum_{D=0}^{e-1}a_{\nu,\beta,D}(S-\beta)^{e-1-D}\right],
 \qquad \lambda_\nu(x)=\ell(xA_\nu).
 \tag{AKJ8}
\]
Indeed, restrict to one primary factor and substitute AKJ8 into AKJ6.
For $x=\varepsilon_\beta z^D$, multiplication cancels $h_\beta$ and
extracts exactly $a_{\nu,\beta,D}$. These vectors are a basis, proving
the asserted equality for every $x$. This gives the concrete inverse
of the residue equivalence used in the incoming formalization, on
the actual OCS rows, with their factorials and phases.

Define the integer $t_\beta$ from these actual coefficients by
\[
 t_\beta=\begin{cases}
 -1,&a_{\nu,\beta,D}=0\text{ for every }\nu,D,\\
 \max\{D:a_{\nu,\beta,D}\ne0\text{ for some }\nu\},&\text{otherwise},
 \end{cases}\qquad n_\beta=e-1-t_\beta.
 \tag{AKJ9}
\]
Define $\mathcal O_q:E\longrightarrow B^q$ by
$\mathcal O_q(x)=(\Lambda M_S^jx)_{0\le j<q}$, where $B=\operatorname{im}\Lambda$
is the original observation image. Thus $0\le n_\beta\le e$.
The complete original permanently invisible
submodule and its exact dimension are
\[
 \begin{gathered}
 K_\infty=\bigcap_{j=0}^{q-1}\ker(\Lambda M_S^j)
 =\bigoplus_\beta
   \varepsilon_\beta(z^{e-n_\beta})\subset E,\\
 \dim K_\infty=\sum_\beta n_\beta,
 \qquad \operatorname{rank}\mathcal O_q=q-\sum_\beta n_\beta.
 \end{gathered}\tag{AKJ10}
\]
Here $(z^e)=0$ and $(z^0)$ is the entire local algebra. To prove this,
vanishing of the first $q$ observations says
$\ell(S^jxA_\nu)=0$ for all $j<q,\nu$. Nondegeneracy just proved and
the power basis imply $xA_\nu=0$ for every row. Conversely those
products being zero kills every iterate. At the factor $\beta$,
AKJ8 is a unit $h_\beta$ times its displayed reverse coefficient
polynomial. The minimum valuation of all $A_\nu$ is precisely
$n_\beta$. Its common annihilator is $(z^{e-n_\beta})$: multiplication
by a coefficient of minimal valuation is a unit times $z^{n_\beta}$,
and all other coefficients have at least that valuation. This proves
the module formula and its dimension; rank-nullity proves the last
equality. Affine substitution $M_S=c+iY$ gives an invertible triangular
change between the first $q$ powers of the two operators, so the same
kernel is obtained using the original $Y$-stack of OCS.

In particular all iterates are jointly injective exactly when every
original highest-jet column $a_{\nu,\beta,e-1}$ has a nonzero entry.
For $m=1$, AKJ8 reduces to
\[
 A_\nu(\beta)=\chi'(\beta)B_{\nu,\beta},\qquad
 K_\infty=\operatorname{span}\{\varepsilon_\beta:
                         B_{\nu,\beta}=0\text{ for all }\nu\}.
 \tag{AKJ11}
\]
This is an exact test on the AKS8 matrix, not an assumption that it
passes. It constructs the inclusion $K_\infty\hookrightarrow K$:
each listed primary ideal is killed by all rows, hence in particular
by their original one-step observation. The source-minimum metric on
this subspace is its restriction from OKM8, retaining the same original
Gram. Full-iterate injectivity alone therefore supplies no estimate of
the one-step kernel determinant AKJ3.

\section{The actual conductor confines the kernel's polynomial degrees}
Return to $m=1$, $k=4l+1\ge9$ and the original five-member period domain.
Let $\mathcal A_Q=\prod_{j=1}^4Q_j$ and write its full Segre restriction
from OCF as
\[
 A=\sum_{r,s=0}^8a_{rs}
 X_1^rX_0^{8-r}Y_1^sY_0^{8-s}\ne0.
 \tag{AKJ12}
\]
OCF proves nonvanishing and retains all actual $\mathsf H$ coefficients.
No particular $a_{rs}$ is presumed nonzero. Write $k'=k-8$,
$q'=(k-7)^2$ and
\[
 b_{rs}=4+(2r-8)\delta+i(2s-8)\gamma,\quad
 (\mathcal T_Ap)(S')=\sum_{r,s=0}^8a_{rs}p(S'+b_{rs}),\quad
 \mu_j=\sum_{r,s}a_{rs}b_{rs}^j.
 \tag{AKJ13}
\]
All $81$ numbers $b_{rs}$ are distinct. Their Vandermonde is invertible,
so some $\mu_j$ with $0\le j\le80$ is nonzero. Define $v$ to be the first
such index. The exact finite differential expression on degree-below-$q$
polynomials is
\[
 \mathcal T_A=\sum_{j=0}^{q-1}\frac{\mu_j}{j!}\partial_S^j
 =\frac{\mu_v}{v!}\partial_S^v B(\partial_S),\quad
 B(X)=1+\sum_{j=v+1}^{q-1}\frac{v!\mu_j}{j!\mu_v}X^{j-v}.
 \tag{AKJ14}
\]
This is Taylor's finite formula. Hence it kills polynomials of degree
below $v$ and lowers every higher degree by exactly $v$. The operator
$B(\partial_S)$ is invertible: with $U=B(\partial_S)-I$, its inverse
is the finite sum $\sum_{a=0}^{q-1}(-U)^a$. Terms beyond nilpotency are
zero. Thus $\mathcal T_A$ maps polynomials of degree at most $d$
onto those of degree at most $d-v$ when $d\ge v$.

Let $\chi_{k'}$ be the original polynomial for tensor degree $k'$.
The conductor's exact annihilator in the original coefficient algebra is
\[
 V_A=\{p\in\mathbb C[S]_{<q}:\chi_{k'}\mid\mathcal T_Ap\},
 \qquad K\subset V_A.
 \tag{AKJ15}
\]
To prove this map, pair the biform $Ag$, of bidegree $(k,k)$, with
the primary values of $p$. If $g$ has coefficient $g_{ab}$ at bidegree
$(k',k')$, convolution of all original coefficients gives the pairing
\[
 \sum_{a,b=0}^{k'}g_{ab}
       \sum_{r,s=0}^8a_{rs}
       p(\beta^{(k')}_{ab}+b_{rs})
 =\sum_{a,b=0}^{k'}g_{ab}
                      (\mathcal T_Ap)(\beta^{(k')}_{ab}).
 \tag{AKJ16}
\]
The addition identity for the roots follows directly from their retained
centres and both signs. Vanishing for every $g$ is exactly vanishing at
all $q'$ distinct roots, hence the divisibility in AKJ15. OCF's proved
inclusion $A\mathbb C[X,Y]_{(k',k')}\subset J_k$ then gives $K\subset V_A$
by its original complex-linear coefficient pairing.

The following is a full degree-adapted basis of $V_A$. Define
$\mathcal J_v S^d=d!S^{d+v}/(d+v)!$, so
$\partial_S^v\mathcal J_v=I$, and set
\[
 \begin{gathered}
 1,S,\ldots,S^{v-1},\qquad
 P_h=\frac{(q'+h+v)!}{(q'+h)!}
       B(\partial_S)^{-1}\mathcal J_v(\chi_{k'}S^h),\\
 0\le h\le q-q'-v-1.
 \end{gathered}\tag{AKJ17}
\]
The first list is empty when $v=0$. Since $q-q'=16k-48\ge96>v$,
the range of the second list is nonnegative. Each $P_h$ is monic of
degree $q'+h+v$ and maps under $\mathcal T_A$ to
$\mu_v(q'+h+v)!/(v!(q'+h)!)$ times $\chi_{k'}S^h$. These statements
follow from AKJ14, the inverse's constant term one, and the leading
coefficient of $\mathcal J_v$. The images span every multiple of
$\chi_{k'}$ of degree at most $q-v-1$, and the low list spans the exact
kernel of $\mathcal T_A$. This proves spanning and independence. In
particular $\dim V_A=q-q'=16k-48$.

Choose the original OCF basis of $J_k/(A\mathbb C[X,Y]_{(k',k')})$,
whose exact dimension is $j_k=8k-32$. Evaluate each chosen biform row
on the primary values of every polynomial in AKJ17. This gives an
explicit $j_k$ by $(16k-48)$ matrix $\Xi$, with the full original
period coefficients. Its kernel maps isomorphically onto $K$ by the
displayed basis. Its rank is $j_k$: the original $J_k$ and conductor
annihilator duality, or dimension $\dim V_A-\dim K$, gives this rank.
For the submatrix $\Xi_{\le d}$ containing only basis columns of degree
at most $d$, the exact formula is
\[
 \dim(K\cap\mathbb C[S]_{\le d})
 =\#\{\text{degrees in AKJ17 at most }d\}
                      -\operatorname{rank}\Xi_{\le d}.
 \tag{AKJ18}
\]
Distinct monic degrees ensure that no omitted higher column can cancel
to form a degree-at-most-$d$ polynomial. Restricting the coefficient
isomorphism to this prefix and applying rank-nullity proves AKJ18.
It uses actual minors and makes no assertion of a generic prefix rank.

Let $d_0<\cdots<d_{r-1}$ be the degree jumps of the actual kernel.
The basis degrees of $V_A$, in increasing order, are $i$ for $i<v$,
and $q'+i$ for $i\ge v$. Inclusion and dimension of each degree prefix,
together with $K\subset\mathbb C[S]_{<q}$, imply
\[
 i+q'\mathbf1_{i\ge v}\le d_i\le q-r+i.
 \tag{AKJ19}
\]
For the lower inequality the $i$th independent kernel vector cannot
appear earlier than the $i$th independent vector of $V_A$. For the
upper inequality, intersection of an $r$-dimensional subspace of
$\mathbb C^q$ with its $(q-r+i+1)$-dimensional degree prefix has dimension
at least $i+1$. Summing yields the exact exterior-degree bounds
\[
 \begin{split}
 \frac{r(r-1)}2+\max\{r-v,0\}q'
 &\le\sum_{i=0}^{r-1}d_i
 \le r(q-r)+\frac{r(r-1)}2,\\
 0\le rq-\sum_i d_i
 &\le\frac{r(r+1)}2+rj_k+\min\{r,v\}q'.
 \end{split}\tag{AKJ20}
\]
The difference between the first line's upper and lower endpoints is
$rj_k+\min\{r,v\}q'$, since $q-q'=r+j_k$. Therefore the original
kernel's exterior degree is $rq-O(k^2)$, with the exact bound above
and $v\le80$, uniformly in the proved period domain. The affine map
$S=c+iy$ preserves all degree prefixes and multiplies a degree-$d$
leading coefficient by $i^d$. Thus these are also the degrees in the
original Gamma coordinate, with all phases known. They specify the
kernel's position without assigning a condition number to its frame.

\section{The exact conductor translations in the original source norm}
Keep either source $s\in\{1,k\}$, $b=s/2$, $M=(2\pi)^{s/2}$ and
the original $p_n,h_n=M n!(b)_n$. On polynomials of degree at most $D$,
let $\varphi_n=p_n/\sqrt{h_n}$ and $w_n=n!/(b)_n$.
Differentiating the original generating identity AKS16 in $y$ gives
\[
 \partial_y\varphi_n=
 \sum_{\substack{1\le a\le n\\a\text{ odd}}}
 \frac{(-1)^{(a-1)/2}}{a}
       \sqrt{\frac{w_n}{w_{n-a}}}\,\varphi_{n-a}.
 \tag{AKJ21}
\]
Indeed differentiation multiplies the generating series by
$\arctan z=\sum_{a\text{ odd}}(-1)^{(a-1)/2}z^a/a$.
Equating its finite coefficient and substituting the full norms gives
exactly AKJ21; the two copies of $M$ are cancelled only at this point.

Write $H_D=\sum_{j=1}^D1/j$ and $H_0=0$. The original source norm obeys
\[
 \|\partial_y\|\le\sqrt{H_D(H_D+4)}\le H_D+2,
 \qquad
 \|f(y+z)\|\le e^{|z|(H_D+2)}\|f\|.
 \tag{AKJ22}
\]
Here is a proof retaining both source orders. For $n>m\ge1$ and
$b\ge1/2$,
\[
 \frac{w_n}{w_m}\le\prod_{j=m+1}^n\frac{j}{j-1/2}
 \le\sqrt{\frac nm},\qquad w_n\le2\sqrt n\quad(n\ge1).
 \tag{AKJ23}
\]
The factor inequality follows by squaring and using
$j(j-1)\le(j-1/2)^2$; the product telescopes. For $m=0$, separate
the first factor, at most two. Apply the weighted Schur estimate to
the absolute matrix in AKJ21 with positive weights $\vartheta_n=1/\sqrt{w_n}$.
Its weighted row sum is at most $H_D$. Its weighted column sum at $n$
is at most $H_n+4$: for $1\le m<n$ the excess over $1/(n-m)$ is at most
\[
 \frac{\sqrt{n/m}-1}{n-m}
 =\frac1{\sqrt m(\sqrt n+\sqrt m)}\le\frac1{\sqrt{nm}},
\]
whose sum is at most two by comparison with $\int_0^n t^{-1/2}dt$.
The $m=0$ term is at most $2/\sqrt n\le2$. Enlarging the odd-index
sums to all indices preserves these bounds. For completeness, if a
nonnegative matrix has weighted row bound $A$ and column bound $B$,
Cauchy--Schwarz gives
$|(Tx)_m|^2\le A\vartheta_m\sum_nT_{mn}|x_n|^2/\vartheta_n$;
summing and using the column bound gives $\|Tx\|^2\le AB\|x\|^2$.
This proves the first bound in AKJ22, including $D=0$ where the
derivative is zero. The exact finite Taylor identity
$f(y+z)=\exp(z\partial_y)f(y)$ and its norm series prove the second.
The inverse translation is obtained by $-z$ with the same bound.

There is consequently an exact original-source upper control for AKJ13.
For $f(y)=p(c_k+iy)$, evaluate its output at $c_{k-8}+iy$, retaining
$c_k=k/2$ and $c_{k-8}=k/2-4$. Then
\[
 (\mathcal T_Ap)(c_{k-8}+iy)
 =\sum_{r,s=0}^8a_{rs}
   f\left(y+\frac{b_{rs}-4}{i}\right),
 \tag{AKJ24}
\]
and therefore
\[
 \|\mathcal T_Ap\|_{s,c_{k-8}}
 \le\left(\sum_{r,s}|a_{rs}|
          e^{|b_{rs}-4|(H_D+2)}\right)\|p\|_{s,c_k}.
 \tag{AKJ25}
\]
The notation $\|p\|_{s,c}=\|p(c+iy)\|_{m_{0,s}}$ changes only the
explicit centre; the measure, mass and source order $s$ are unchanged.
AKJ24 proves the centre correction rather than discarding it. For a
fixed original period and quartet the coefficient in AKJ25 grows at
most polynomially in $D$, since $H_D\le1+\log D$ for $D\ge1$ and
$|b_{rs}-4|\le8\sqrt{\delta^2+\gamma^2}$. The entire actual coefficient
sum remains visible. The finite inverse in AKJ17 has the specified
moment denominators $\mu_v$; AKJ25 asserts no bound on that inverse.
Together AKJ18--25 provide exact degree information and a source-norm
estimate for the actual conductor map, to be used in the remaining
four-endpoint compression calculation. No leading kernel value or
kernel-invariance assumption has replaced that calculation.

% END CURRENT AKJ

\clearpage
% BEGIN CURRENT AKP
\section{The original theta class and the checked metric interface}
\label{sec:AKP}
The absolute base and original coefficient resolution are retained:
\[
 \mathbf F_{1,\tau}=\{\tau,1\},\qquad
 \mathsf A_\tau(\mathcal T)=
 [V\oplus V\xrightarrow{\Theta(\phi-\widehat\psi)}\mathscr B],
 \qquad Q=\mathscr B/\Theta V.
 \tag{AKP1}
\]
The four-point chart has $+,-<\eta<\sigma$ and coefficients
$(V,V,\mathscr B,0)$, with the original restrictions $\Theta$ and
$J\Theta=\Theta\mathcal F$. Its tensor resolution identifies the top
cohomology with $Q^{\otimes k}$. AT1--2a supplies the complete original
construction. In its notation, retaining the full unit
$\upsilon_h=j_h(2\xi/h)$, the actual cochain and arithmetic class maps are
\[
 \begin{gathered}
 \mathcal V_{h,k}(P)=P(D_1+\cdots+D_k)F_h^{\otimes k},\qquad
 E=\mathbb C[S]/(\chi),\\
 \eta_{h,k}[P]=\upsilon_h^{\otimes k}P(A_h^{(k)})1,\qquad
 q^{(k)}\mathcal V_{h,k}(P)
   =\sigma_h^{\otimes k}\eta_{h,k}[P].
 \end{gathered}\tag{AKP2}
\]
Here $\sigma_h:E_h\hookrightarrow Q$ is the original section and
$\eta_{h,k}:E\hookrightarrow E_h^{\otimes k}$ is the full-jet injection.
Both injections are part of the supplied source proof; the cochain map
has not been replaced by an arbitrary observation.

For completeness, the exact representative-change calculation is as
follows. If $P'=P+\chi\kappa$, perform successive monic divisions of
$\chi(s_1+\cdots+s_k)$ by $h(s_a)$. Its final remainder is zero because
$\chi$ is the original annihilator of the sum in
$\bigotimes_a\mathbb C[s_a]/(h)$. Thus the actual divisions produce
$Q_a(\mathbf s)$ satisfying
\[
 \chi(s_1+\cdots+s_k)=\sum_{a=1}^kh(s_a)Q_a(\mathbf s).
 \tag{AKP3}
\]
The chosen ordered division determines these polynomials; no uniqueness
of all such decompositions is asserted. On the original tensor complex set
\[
 \Psi_\kappa=\sum_{a=1}^k(-1)^{a-1}
 Q_a(\mathbf D)\kappa\left(\sum_bD_b\right)
 \bigl(F_h^{\otimes(a-1)}\otimes\phi_0
                  \otimes F_h^{\otimes(k-a)}\bigr).
 \tag{AKP4}
\]
Its differential is
\[
 d\Psi_\kappa=\mathcal V_{h,k}(\chi\kappa)
             =\mathcal V_{h,k}(P')-\mathcal V_{h,k}(P).
 \tag{AKP5}
\]
Indeed the only nonzero differential in each summand acts at its
$\phi_0$ position. The tensor sign is $(-1)^{a-1}$, cancelling the
displayed sign. The identity $\Theta\phi_0=h(D)F_h$ gives the $a$th
term of AKP3; all operators commute with this differential. Summing
proves AKP5 with the original signs. This is EP.48, expressed at the
current receiving map. Consequently the quotient classes in AKP2 agree.
The chosen primitive for their difference depends on $\kappa$.
At an existing support label $\ell$, this quotient map sends the
relation to $(\ell,0)$ and sends external absence $\tau$ to $\tau$.
The coordinate change has therefore preserved the arithmetic class
and its support through an explicit cochain homotopy.

This calculation specifies the correction to the sentence following
AKS49: equality of cohomology classes is retained, with AKP4 the
primitive of the change. Membership in an additional period-observation
kernel alone supplies no such primitive. Its proved map is the inclusion
$K_\infty\hookrightarrow K\hookrightarrow E$ of AKJ10--11, followed by
the original arithmetic injection in AKP2. In particular, a nonzero
coefficient in this additional kernel remains a nonzero original
arithmetic class under that injection. Its zero observed value and its
arithmetic class are related by these exact maps.

\subsection{The complete minimum-section calculation at each endpoint}
Let $G$ be one of the original positive Hermitian coefficient Grams,
write $K_G=G^{-1}$, and let $\Lambda:E\twoheadrightarrow B$ be the
actual observation onto its image. Define
\[
 Q_G=(\Lambda K_G\Lambda^*)^{-1},\qquad
 S_G=K_G\Lambda^*Q_G.
 \tag{AKP6}
\]
The inverse exists because $\Lambda^*$ is injective and $K_G$ is
positive definite. Direct multiplication, using Hermitian symmetry,
gives
\[
 \Lambda S_G=I_B,\quad S_G^*G=Q_G\Lambda,\quad
 S_G^*GS_G=Q_G.
 \tag{AKP7}
\]
For every rectangular column family $X$ put $Z=\Lambda X$ and
$R=X-S_GZ$. Then
\[
 \Lambda R=0,\qquad R^*GS_G=0,\qquad
 R^*GR=X^*GX-Z^*Q_GZ\succeq0.
 \tag{AKP8}
\]
The first two equalities follow from AKP7. In the expansion of the
last Gram, each cross term is $Z^*Q_GZ$ and the section Gram is the
same matrix. The resulting single subtraction proves the identity,
including every off-diagonal pairing. Any other lift is $S_GZ+R'$
with $\Lambda R'=0$. Its Gram is $Z^*Q_GZ+R'^*GR'$, proving the
minimum property by the same orthogonality calculation.

Retain a fixed section $S_*$, the actual inclusion $I_K$, and original
coordinate map $\kappa_K$ satisfying
$I_K\kappa_K=I_E-S_*\Lambda$. Applying this identity to $R$ gives
\[
 R=I_K\bigl(\kappa_KX-\kappa_KS_G\Lambda X\bigr).
 \tag{AKP9}
\]
Thus the fixed-section coordinates include their specified minimum
correction. For one original new column $e_D$, define
$t_D=e_D^*Ge_D$, $\lambda_D=\Lambda e_D$ and
$\xi_D=\lambda_D^*Q_G\lambda_D$. AKP8 yields
$t_D-\xi_D\ge0$, so the exact ordered factors satisfy
\[
 (1+\xi_D)\left(1+\frac{t_D-\xi_D}{1+\xi_D}\right)=1+t_D.
 \tag{AKP10}
\]
The denominator is positive and is retained. These are the full OPR
minimum-section identities, now proved at the original Grams of
AKS26 and OKM8. For an empty observation image, the section is the
zero map, the Gram on that image has determinant one, and the same
identities reduce to $R=X$. No empty present support is deleted.

The new formalization in draft PR32 checks these finite identities and
the original observation-iterate and residue-annihilator maps. Its
checked implementation is
\[\texttt{8e78bc7c240b04d297ade6afdadfd863e0c6db7b}.\]
The later source and status head is
\[\texttt{9cec8482f2ecf978cf8bdb67b4c080b9dd74f5d6}.\]
AKJ6--11 evaluates the residue inverse in every actual primary-jet
coefficient, giving the exact map to those checked declarations.
The analytic AKS estimate, the conductor metric estimates, and the
polynomial-gcd dimension corollary have the full written proofs stated
here or in their supplied sources; they are outside that Lean run.

% END CURRENT AKP

\clearpage
% BEGIN CURRENT AIE
% Full original conductor inverse estimate. Child calculation, 2026-09-14.
\section{The finite conductor inverse in the two original Gamma norms}
\label{sec:AIE}
\paragraph{Current refinement.}
The finite estimates in this chapter remain valid. RUI5--15 and
CEP29--49 below sharpen their exterior consequence: complementary
singular values and the exact source quotient give an $o(kq)$
conductor error at all four original endpoints, for a fixed actual
period. The full original constants, moments and low-degree graph
are retained in those proofs.


Retain AKJ12--25 and the complete OCF conductor coefficient polynomial.
Thus $m=1$, $k=4l+1\ge9$, $q=(k+1)^2$, $0<\delta<1/2$,
$\gamma>2$, and the original period belongs to the proved five-orbit
locus. Its fixed $4$ by $4$ matrix $\mathsf H$ determines the four
quadrics and all $81$ coefficients $a_{rs}$ of their full product.
They remain the same coefficients as the tensor degree $k$ varies.
Put
\[
 b_{rs}=4+(2r-8)\delta+i(2s-8)\gamma,\quad
 \mu_j=\sum_{r,s=0}^8a_{rs}b_{rs}^j,\quad
 v=\min\{j:\mu_j\ne0\}\le80.
 \tag{AIE1}
\]
The existence and bound for $v$ are the full distinct-node Vandermonde
argument of AKJ13. Throughout, $0\le D\le q-1$ denotes the polynomial
degree cutoff. Every norm below is the original Gamma norm of source
order $s\in\{1,k\}$, with
\[
 \begin{gathered}
 b=s/2,\quad M=(2\pi)^{s/2},\quad h_n=M n!(b)_n,\\
 \varphi_n=p_n/\sqrt{h_n},\quad w_n=n!/(b)_n,
 \quad \|p\|_{s,c}=\|p(c+iy)\|_{m_{0,s}}.
 \end{gathered}\tag{AIE2}
\]
The centre $c$ is displayed wherever it changes. No source order,
source mass, coefficient, or period is assigned a replacement value.

\subsection{A coefficient-controlled disk for the literal finite symbol}

The exact symbol used in AKJ14 and AKJ17 is
\[
 B_q(z)=1+\sum_{j=v+1}^{q-1}
                \frac{v!\mu_j}{j!\mu_v}z^{j-v}.
 \tag{AIE3}
\]
The sum is empty if its lower index exceeds its upper index. Define
the positive, original-data constants
\[
 \begin{gathered}
 A_*^{\rm abs}=\sum_{r,s}|a_{rs}|,\qquad
 B_*^{\rm sh}=\max_{r,s}|b_{rs}|,\qquad
 C_v=\frac{A_*^{\rm abs}(B_*^{\rm sh})^v}{|\mu_v|},\\
 R_0=\frac{1}{B_*^{\rm sh}}
       \min\left\{1,\frac{v+1}{2\mathrm e C_v}\right\},
 \quad \rho_0=\tanh R_0.
 \end{gathered}\tag{AIE4}
\]
Here $A_*^{\rm abs}>0$, $B_*^{\rm sh}\ge4$, $C_v\ge1$,
$R_0>0$, and $0<\rho_0<1$. The superscripts distinguish these
constants from the original conductor polynomial and from every
earlier radius or pivot coefficient. The exact moment $\mu_v$,
including any cancellation in it, remains in the denominator.

For every $|z|\le R_0$ one has
\[
 |B_q(z)-1|
 \le \frac{v!}{|\mu_v|}
       \sum_{j=v+1}^{q-1}\frac{|\mu_j|}{j!}|z|^{j-v}
 \le C_v\frac{B_*^{\rm sh}|z|}{v+1}
                    \exp(B_*^{\rm sh}|z|)
 \le\frac12.
 \tag{AIE5}
\]
To prove the middle inequality, use
$|\mu_j|\le A_*^{\rm abs}(B_*^{\rm sh})^j$, write
$j=v+1+h$, and apply
$v!/(v+1+h)!\le1/((v+1)h!)$. Summation of the exponential
series gives the displayed bound. Definition AIE4 gives
$B_*^{\rm sh}R_0\le1$ and
$C_v B_*^{\rm sh}R_0/(v+1)\le1/(2\mathrm e)$, proving
the last inequality. The first upper bound in AIE5 retains every
individual finite moment and may be used directly for a larger
verified radius. The explicit $R_0$ proves a positive disk for
every actual coefficient chart; no uniform lower bound for
$|\mu_v|$ over varying periods is asserted.

The analytic function
\[
 F_q(z)=\frac1{B_q(-i\arctan z)}
          =\sum_{a=0}^\infty f_{q,a}z^a
 \tag{AIE6}
\]
is well defined on a neighbourhood of the closed disk
$|z|\le\rho_0$, and
$|F_q(z)|\le2$ there. Indeed the power series of $\arctan$
gives $|\arctan z|\le\operatorname{arctanh}|z|\le R_0$;
AIE5 gives $|B_q(-i\arctan z)|\ge1/2$. Compactness gives
an open neighbourhood with nonzero denominator. Cauchy's coefficient
formula therefore proves
\[
 f_{q,0}=1,\qquad |f_{q,a}|\le2\rho_0^{-a}
                    \quad(a\ge1).
 \tag{AIE7}
\]
Although the literal symbol $B_q$ changes with $q$, the disk and its
bound depend only on the original fixed conductor and shifts.

\subsection{The exact coefficient map and its norm}

Let $\mathcal P_D$ denote the polynomials in $y$ of degree at most
$D$, with the orthonormal basis in AIE2. Define the lowering map
\[
 L_s\varphi_0=0,\qquad
 L_s\varphi_n=\sqrt{\frac{w_n}{w_{n-1}}}\,\varphi_{n-1}
                         \quad(1\le n\le D).
 \tag{AIE8}
\]
Its powers have exactly the entries
$L_s^a\varphi_n=\sqrt{w_n/w_{n-a}}\varphi_{n-a}$ for
$a\le n$ and vanish for $a>n$. Thus $L_s^{D+1}=0$.
AKJ21 gives the polynomial identity
$\partial_y=\arctan L_s$, with every power above degree $D$
zero. For the coordinate pullback $\mathcal C_c p(y)=p(c+iy)$,
the chain rule is
$\mathcal C_c\partial_S\mathcal C_c^{-1}=-i\partial_y$.
Consequently the inverse of the literal finite operator is
\[
 \begin{split}
 \mathcal C_c B_q(\partial_S)^{-1}\mathcal C_c^{-1}
   &=F_q(L_s),\\
 F_q(L_s)\varphi_n
   &=\sum_{a=0}^n f_{q,a}
                     \sqrt{\frac{w_n}{w_{n-a}}}\,\varphi_{n-a}.
 \end{split}\tag{AIE9}
\]
For proof, evaluate the analytic power-series identity
$B_q(-i\arctan z)F_q(z)=1$ in the nilpotent matrix $L_s$.
All relevant coefficients are finite, so multiplication is an exact
finite polynomial identity. This also proves that AIE9 is the same
inverse as the finite nilpotent Neumann sum in AKJ14. The factor $-i$
is forced by the original coordinate map and has been retained.

Define the exact weight ratio and the resulting bound
\[
 \begin{split}
 \alpha_{D,s}&=\max_{0\le m\le n\le D}
                         \sqrt{\frac{w_n}{w_m}},\\
 N_{D,s}&=\alpha_{D,s}
            \left(1+2\sum_{a=1}^{D}\rho_0^{-a}\right)
 \le \frac{2\alpha_{D,s}}{1-\rho_0}\rho_0^{-D}.
 \end{split}\tag{AIE10}
\]
Every absolute row sum and every absolute column sum of AIE9 is at
most $\alpha_{D,s}\sum_{a=0}^D|f_{q,a}|$. Cauchy--Schwarz in each
row followed by the column bound proves that the operator norm is
at most this number; AIE7 gives $N_{D,s}$. It follows that, at every
unchanged centre $c$,
\[
 \|B_q(\partial_S)^{-1}p\|_{s,c}
                \le N_{D,s}\|p\|_{s,c}.
 \tag{AIE11}
\]
When $s=k$, $b=k/2\ge1$ and
$w_n/w_{n-1}=n/(b+n-1)\le1$, so $\alpha_{D,k}=1$.
When $s=1$, the same ratio is greater than one, so
$\alpha_{D,1}=\sqrt{D!/(1/2)_D}$; for $D\ge1$, the exact
factor estimate of AKJ23 yields
$\alpha_{D,1}\le\sqrt2 D^{1/4}$. At $D=0$, both values are
one. In particular
\[
 \log\|B_q(\partial_S)^{-1}\|_{s,c}
 \le D\log(\coth R_0)
       +\log\frac{2\alpha_{D,s}}{1-\tanh R_0}
             =O_{\mathsf H,\delta,\gamma}(D+\log(D+1)).
 \tag{AIE12}
\]
For the fixed original quartet and period, $q=(k+1)^2$ and
$D\le q-1$ imply that this logarithm divided by $kq$ tends to
zero. This estimate holds for both original source orders at once;
its constants explicitly retain the actual $\mu_v$.

\subsection{The infinite symbol differs by an exact low-degree map}

To record precisely the truncation in AIE3, define the entire function
\[
 B_\infty(z)=\frac{v!}{\mu_v z^v}
                 \sum_{r,s}a_{rs}e^{b_{rs}z},\qquad B_\infty(0)=1.
 \tag{AIE13}
\]
Its apparent singularity is removable by the definition of $v$.
On polynomials of degree at most $D\le q-1$, set
$E_q=B_\infty(\partial_S)-B_q(\partial_S)$; its exact expression is
\[
 E_q=\sum_{n=q-v}^{D}
       \frac{v!\mu_{n+v}}{(n+v)!\mu_v}\partial_S^n.
 \tag{AIE14}
\]
An empty sum is zero. Its image has degree at most
$D-q+v\le v-1$. Both symbols have constant term one, hence both
are invertible in this finite polynomial algebra. Their inverses obey
\[
 B_q(\partial_S)^{-1}-B_\infty(\partial_S)^{-1}
 =B_q(\partial_S)^{-1}E_qB_\infty(\partial_S)^{-1},
 \tag{AIE15}
\]
whose image is contained in $\mathbb C[S]_{<v}$ and whose rank is
at most $v$. Equation AIE15 follows by multiplying out its right
side. The degree statement holds because every factor is a power
series in the degree-lowering derivative. Thus the infinite-symbol
formula can be related to the literal AKJ inverse by the displayed
exact correction; it cannot replace that inverse without this map.

\subsection{An explicit right inverse of the original conductor}

Assume now $D\ge v$. Define the raising map, on the degree ranges
where its image lies in $\mathcal P_D$, by
\[
 J_s\varphi_n=\sqrt{\frac{w_n}{w_{n+1}}}\varphi_{n+1}
       =\sqrt{\frac{b+n}{n+1}}\varphi_{n+1}.
 \tag{AIE16}
\]
Then $L_s^vJ_s^v=I$ on $\mathcal P_{D-v}$ and
$\|J_s^v\|\le\max\{1,b^{v/2}\}$. The latter follows from
the displayed coefficient, at each raising step, by
$(b+n)/(n+1)\le\max\{1,b\}$.

The analytic function $G(z)=z/\arctan z$, with $G(0)=1$,
satisfies $|G(z)|\le2$ on $|z|<1$. Indeed
\[
 \frac{\arctan z}{z}=\int_0^1\frac{dt}{1+t^2z^2},\qquad
 \operatorname{Re}\frac1{1+u}>\frac12\quad(|u|<1).
 \tag{AIE17}
\]
The real-part inequality follows from
$2\operatorname{Re}(1+\overline u)>|1+u|^2$, which reduces
to $1>|u|^2$. It proves that the integral is nonzero with
real part greater than $1/2$, including the removable value at zero.
Every Taylor coefficient of $G(z)^v$ is therefore bounded by $2^v$:
apply Cauchy's formula on radius $t<1$ and let $t$ increase to one
for each fixed coefficient.

Define a map between the specified degree spaces by
\[
 R_{v,s}=G(L_s)^v J_s^v:\mathcal P_{D-v}\longrightarrow\mathcal P_D.
 \tag{AIE18}
\]
Since all functions of $L_s$ commute, AIE8 and the power-series
identity $(\arctan z)^vG(z)^v=z^v$ give
$\partial_y^vR_{v,s}=L_s^vJ_s^v=I$. This proves the right-inverse
identity even at the top allowed degree. The same row and column
argument as for AIE10, now with every coefficient bounded by $2^v$,
gives
\[
 \|R_{v,s}\|\le
   2^v(D+1)\alpha_{D,s}\max\{1,b^{v/2}\}.
 \tag{AIE19}
\]
When $v=0$ the map is exactly the identity; AIE19 remains a valid
upper bound. Define the corresponding map on polynomials in the
original $S$ coordinate by
\[
 \mathcal R_{v,s,c}
       =\mathcal C_c^{-1}i^vR_{v,s}\mathcal C_c.
 \tag{AIE20}
\]
The chain rule in AIE9 proves
$\partial_S^v\mathcal R_{v,s,c}=I$, with exactly the factor
$(-i)^vi^v=1$. This is a map in the original polynomial algebra;
the source norm only specifies its coefficients through the explicitly
given original orthonormal basis.

Let $c_k=k/2$ and $c'=c_{k-8}=k/2-4$. The full conductor of AKJ13
has the following right inverse on $\mathbb C[S']_{\le D-v}$:
\[
 \mathcal R_{A,s}g
     =\frac{v!}{\mu_v}
             B_q(\partial_S)^{-1}\mathcal R_{v,s,c'}g,
 \qquad \mathcal T_A\mathcal R_{A,s}g=g.
 \tag{AIE21}
\]
To check it, use the exact finite factorization
$\mathcal T_A=(\mu_v/v!)\partial_S^vB_q(\partial_S)$ from
AKJ14 on degree at most $D$. The two inverse factors cancel and
AIE20 gives the stated identity. In the identity the output variable
is $S'$; the differential expression uses the same numerical
polynomial coefficients, so the domain and codomain are as displayed.

Set $H_D=\sum_{j=1}^D1/j$, $H_0=0$, and
\[
 C_{A,D,s}=\frac{v!}{|\mu_v|}\,
  e^{4(H_D+2)}N_{D,s}\,
  2^v(D+1)\alpha_{D,s}\max\{1,b^{v/2}\}.
 \tag{AIE22}
\]
Then the original-centre norm estimate is
\[
 \|\mathcal R_{A,s}g\|_{s,c_k}
                    \le C_{A,D,s}\|g\|_{s,c'}.
 \tag{AIE23}
\]
Indeed AIE11 and AIE19 first prove the bound without the exponential
factor, with both norms at $c'$. For any output polynomial $p$,
$p(c_k+iy)=p(c'+i(y-4i))$, so AKJ22 applied to the exact shift
$-4i$ supplies the remaining factor. Thus the four-unit centre
difference is retained in the typed norm map. For the fixed actual
quartet and period, $v\le80$ and $b\in\{1/2,k/2\}$ give
\[
 \log C_{A,D,s}
 \le D\log(\coth R_0)+O_{\mathsf H,\delta,\gamma}
                         (\log(D+1)+\log(k+1)),
 \qquad \frac{\log C_{A,D,s}}{kq}\longrightarrow0.
 \tag{AIE24}
\]
Every factor in the first assertion is written in AIE22; the
asymptotic statement follows from $D\le(k+1)^2-1$ and the fixed
finite $v$. No varying-period bound is inferred from this statement.

For comparison with the exact monomial lift $\mathcal J_v$ of AKJ17,
put
$L_g=\mathcal R_{v,s,c'}g-\mathcal J_vg$. Both terms have
$v$th derivative $g$, so $L_g\in\mathbb C[S]_{<v}$. More explicitly,
\[
 L_g=\sum_{a=0}^{v-1}
       [S^a](\mathcal R_{v,s,c'}g)S^a,
 \quad
 \mathcal R_{A,s}g-\frac{v!}{\mu_v}B_q(\partial_S)^{-1}\mathcal J_vg
        =\frac{v!}{\mu_v}B_q(\partial_S)^{-1}L_g.
 \tag{AIE25}
\]
The first identity uses the exact zero coefficients of $\mathcal J_vg$
in degrees below $v$. This proves the complete degree-below-$v$
correction between the two lifts, rather than discarding their
different choices of integration constants.

\subsection{The exact quotient and exterior transfer supplied by this bound}

The kernel of $\mathcal T_A$ on degree at most $D\ge v$ is exactly
$\mathbb C[S]_{<v}$ by AKJ14. Give the quotient
$\mathbb C[S]_{\le D}/\mathbb C[S]_{<v}$ its Hilbert quotient norm
from the original source at $c_k$, and give its target
$\mathbb C[S']_{\le D-v}$ the original source at $c'$.
The induced conductor map $\overline{\mathcal T}_A$ is an
isomorphism, because AIE21 proves surjectivity and the displayed
kernel proves injectivity. Define the full upper norm from AKJ25 by
\[
 N_{A,D}=\sum_{r,s}|a_{rs}|
            e^{|b_{rs}-4|(H_D+2)}.
 \tag{AIE26}
\]
For every target $g$, its minimum-norm preimage class obeys
\[
 \frac{\|g\|_{s,c'}}{N_{A,D}}
 \le\|\overline{\mathcal T}_A^{-1}g\|_{\rm quot}
 \le C_{A,D,s}\|g\|_{s,c'}.
 \tag{AIE27}
\]
For the left inequality, apply AKJ25 to every preimage and take the
infimum. For the right inequality, use the actual preimage AIE21.
The infimum is attained since the omitted subspace is a closed
finite-dimensional subspace. These arguments prove the norm statement
with the specified quotient map, rather than assigning the coefficient
norm to the quotient.

In particular AIE21 gives the exact decomposition
\[
 V_A=\mathbb C[S]_{<v}\ \mathbin{\oplus}\
   \mathcal R_{A,s}\bigl(
         \chi_{k-8}\mathbb C[S']_{\le D-v-q'}\bigr),
 \quad q'=(k-7)^2,
 \tag{AIE28}
\]
where a negative degree means the zero space and $V_A$ here denotes
the AKJ15 space restricted to degree at most $D$. Indeed the image
under $\mathcal T_A$ of a member of $V_A$ is precisely a multiple
of $\chi_{k-8}$ in the displayed degree range. Subtracting its lift
leaves the exact kernel. The directness follows by applying
$\mathcal T_A$. Equation AIE25 identifies this decomposition with
the earlier monic basis, including its low-degree correction.

For any $t$-dimensional target subspace and any basis column matrix
$X$ in its original orthonormal coordinates, AIE27 implies the
full Gram determinant inequalities
\[
 N_{A,D}^{-2t}\det(X^*X)
 \le\det\bigl(X^*(\overline{\mathcal T}_A^{-1})^*
                   \overline{\mathcal T}_A^{-1}X\bigr)
 \le C_{A,D,s}^{2t}\det(X^*X).
 \tag{AIE29}
\]
For proof, AIE27 is a two-sided positive Hermitian matrix inequality
after congruence by $X$. Conjugating by $(X^*X)^{-1/2}$ gives
eigenvalues between the two scalar squared bounds, whose product is
the displayed determinant ratio. This proves the exact exterior
transfer and all powers of the norm constants.

Similarly, for any $t$-plane $U\subset\mathbb C[S]_{\le D}$,
AIE11 gives
$\|\bigwedge^t(B_q(\partial_S)^{-1}|_U)\|\le N_{D,s}^t$.
Although $B_q(\partial_S)$ is triangular with every diagonal entry
one and hence determinant one on the full polynomial space, that
full-space determinant does not cancel its restriction to an
arbitrary subspace. The exact restricted determinant is
$\det(X^*(B_q^{-1})^*B_q^{-1}X)/\det(X^*X)$ for a basis matrix
$X$ of that subspace; this is the map to which AIE11 applies.

For the original kernel rank $r=8k-16$ and $D\asymp q$, the upper
logarithm in these exterior bounds is $O_{\mathsf H,\delta,\gamma}(rq)
=O_{\mathsf H,\delta,\gamma}(kq)$. Therefore AIE12 and AIE24 prove
a sub-$kq$ logarithmic bound for each operator, while AIE29 retains
the rank factor required for volumes. They provide the original
conductor inverse, its full low-degree comparison and its quotient
metric control. A sub-$kq$ error for the actual rank-$r$ four-endpoint
determinant is not a consequence of these scalar bounds; that
remaining calculation must use the actual restricted matrices and
the existing endpoint correlation.

% END CURRENT AIE


\clearpage
% BEGIN CURRENT RUI
\section{Subleading exterior control from the actual finite conductor}
\label{sec:RUI}
Retain the original conductor of AKJ12--25 and AIE1--29. In particular
the quartet, full unit, original period and branch are fixed in the
proved five-orbit domain, $m=1$, $k=4l+1\ge9$, and $q=(k+1)^2$.
All $81$ coefficients $a_{rs}$ are those of the actual product of four
period quadrics. They are fixed as $k$ varies. Retain
\[
 \begin{gathered}
 b_{rs}=4+(2r-8)\delta+i(2s-8)\gamma,\qquad
 \mu_j=\sum_{r,s=0}^8a_{rs}b_{rs}^j,\qquad
 v=\min\{j:\mu_j\ne0\}\le80,\\
 A_{\rm abs}=\sum_{r,s}|a_{rs}|,\quad
 B_{\rm sh}=\max_{r,s}|b_{rs}|,\quad
 C_v=\frac{A_{\rm abs}B_{\rm sh}^v}{|\mu_v|}\ge1.
 \end{gathered}\tag{RUI1}
\]
No coefficient or moment denominator is replaced by a generic value.
The finite bound for $v$ is AKJ13's distinct-node Vandermonde proof.
The inequality $C_v\ge1$ follows from the triangle inequality applied
to the actual $v$th moment. The source measure at order $s\in\{1,k\}$
is exactly
\[
 dm_{0,s}(y)=(2\pi)^{s/2}
 \frac{2^{s/2}}{4\pi\Gamma(s/2)}
 |\Gamma(s/4+iy/2)|^2\,dy,\qquad
 \|p\|_{s,c}=\|p(c+iy)\|_{L^2(m_{0,s})}.
 \tag{RUI2}
\]
Every centre appearing below is stated, and the whole mass remains in
the source norm. Put $H_D=\sum_{j=1}^D1/j$, $H_0=0$.

\subsection{The forward polynomial bound and exact determinant}
For an integer cutoff $L\ge v+1$, define the literal finite symbol
\[
 B_L(z)=1+\sum_{a=1}^{L-v-1}
       \frac{v!\mu_{v+a}}{(v+a)!\mu_v}z^a.
 \tag{RUI3}
\]
The empty sum is zero. This includes the actual coefficient symbol
$B_q$ of AKJ14 and the source symbol $B_{N+1}$ needed at a larger
source cutoff $N$. For every displayed coefficient, including the
constant coefficient, one has
\[
 \left|\frac{v!\mu_{v+a}}{(v+a)!\mu_v}\right|
 \le C_v\frac{v!B_{\rm sh}^a}{(v+a)!}
 \le C_v\frac{B_{\rm sh}^a}{a!}.
 \tag{RUI4}
\]
The first inequality is the full moment bound
$|\mu_{v+a}|\le A_{\rm abs}B_{\rm sh}^{v+a}$.
The second follows from $(v+a)!\ge v!a!$.

On $\mathbb C[S]_{\le D}$ with the norm RUI2, the original
coordinate pullback sends $\partial_S$ to $-i\partial_y$.
AKJ21--23 proves $\|\partial_S\|\le H_D+2$, uniformly in both
original source orders. Applying RUI4 term by term to the finite
operator polynomial gives
\[
 \|B_L(\partial_S)\|_{s,c}\le M_D,
 \qquad M_D=C_v\exp\bigl(B_{\rm sh}(H_D+2)\bigr)\ge1.
 \tag{RUI5}
\]
Indeed each power has norm at most $(H_D+2)^a$; extending this
nonnegative finite sum to the exponential series proves the displayed
bound. The constant term is one and every other term strictly lowers
degree. Thus, on this exact $(D+1)$-dimensional polynomial space,
\[
 \det B_L(\partial_S)=1.
 \tag{RUI6}
\]
This is its determinant in the original increasing-power basis, and
therefore in every basis of the same source space, including the
original orthonormal Gamma basis. The assertion preserves the source
measure; it is not a rescaling of this determinant.

\subsection{Complementary singular values retain the full exterior norm}
Let $n=D+1$ and $T=B_L(\partial_S)$ acting on this same Hilbert space.
Write its singular values as $\sigma_1\ge\cdots\ge\sigma_n>0$.
RUI6 gives $\prod_i\sigma_i=1$. For every $0\le t\le n$,
\[
 \begin{split}
 \|\bigwedge^tT^{-1}\|
 &=\prod_{j=n-t+1}^n\sigma_j^{-1}
   =\prod_{j=1}^{n-t}\sigma_j
   =\|\bigwedge^{n-t}T\|
   \le M_D^{n-t},\\
 \|\bigwedge^tT\|&\le M_D^t.
 \end{split}\tag{RUI7}
\]
For a proof of the exterior norm formula, the singular-value
decomposition expresses $T$ as two unitary maps and a diagonal map
with these positive entries. Exterior powers of the unitary maps
preserve the induced norm; the diagonal exterior map has entries
equal to the products over all $t$-subsets. Its largest entry is the
product of the largest $t$ singular values. Applying this to $T^{-1}$
and using the full determinant proves every equality in RUI7. Empty
products cover $t=0,n$.

Consequently, for any actual $t$-dimensional subspace and any full-rank
basis column matrix $X$ in the original orthonormal source coordinates,
\[
 M_D^{-2t}\det(X^*X)
 \le\det\bigl(X^*(T^{-1})^*T^{-1}X\bigr)
 \le M_D^{2(n-t)}\det(X^*X).
 \tag{RUI8}
\]
For the upper inequality apply the first exterior norm of RUI7 to
the decomposable exterior vector of the columns of $X$. Its squared
norm is their Gram determinant. For the lower inequality, put
$Y=T^{-1}X$ and apply the second bound of RUI7 to $TY=X$.
This proves the lower bound on this same subspace without requiring
that it be invariant. Every original off-diagonal Gram entry is
retained by the displayed matrices.

In particular, for fixed original period and quartet, uniformly in
$s\in\{1,k\}$, every $0\le D\le2q$ and every allowed $t$ and $X$,
\[
 \left|\log\frac{
 \det(X^*(T^{-1})^*T^{-1}X)}{\det(X^*X)}\right|
 \le2(D+1)\bigl(\log C_v+B_{\rm sh}(H_D+2)\bigr)
 =O_{\rm actual}(q\log(q+1))=o(kq).
 \tag{RUI9}
\]
The harmonic estimate $H_D\le1+\log D$ for $D\ge1$ follows by
integrating $1/x$ on consecutive intervals; $D=0$ is immediate.
Since $q=(k+1)^2$, the displayed upper bound divided by $kq$ tends
to zero, independently of $D,t,X$. This is a uniform estimate over
the actual subspaces, not a presumed favourable position of the
kernel. Its constants are fixed only after fixing the original period
and quartet; no uniform bound on $1/|\mu_v|$ over moving periods is
asserted. The same inequalities for $T$ follow by exchanging the
image subspace and its inverse image.

RUI9 improves the rank-multiplied estimate at the end of AIE29.
The earlier estimate is valid, but its factor $t$ on a single inverse
operator bound is unnecessary. The complementary exterior identity
uses the exact full determinant and supplies the sharper bound on
the same original matrices.

\subsection{The actual conductor quotient inherits the sharper volume control}
Let $D\ge v$ and use the literal source truncation $L=D+1$.
Taylor's finite formula gives, on degree at most $D$,
\[
 \mathcal T_A=\sum_{j=v}^{D}\frac{\mu_j}{j!}\partial_S^j
             =\frac{\mu_v}{v!}\partial_S^vB_{D+1}(\partial_S),
 \qquad\ker\mathcal T_A=\mathbb C[S]_{<v}.
 \tag{RUI10}
\]
The first equality is exactly the translation operator in AKJ13;
all lower moments vanish by their proved definition. The factor
$B_{D+1}(\partial_S)$ is invertible and preserves each degree
filtration, so the last equality follows from the derivative kernel.

Use the original source centre $c_k=k/2$ and target centre
$c'=c_{k-8}=k/2-4$, both with source order $s$. AIE16--20 gives
the full right inverse $\mathcal R_{v,s,c'}$ of $\partial_S^v$
with norm at most
\[
 P_{D,s}=2^v(D+1)\alpha_{D,s}\max\{1,(s/2)^{v/2}\},
 \quad \alpha_{D,s}=\max_{0\le m\le n\le D}
          \sqrt{\frac{n!/(s/2)_n}{m!/(s/2)_m}}.
 \tag{RUI11}
\]
In particular $\alpha_{D,k}=1$, while
$\alpha_{D,1}\le\sqrt2\max\{1,D^{1/4}\}$; the full derivation
and the phase $i^v$ are the original AIE16--20 construction.
The exact conductor right inverse is
\[
 \mathcal R_{A,s,D}=
 \frac{v!}{\mu_v}B_{D+1}(\partial_S)^{-1}\mathcal R_{v,s,c'}.
 \tag{RUI12}
\]
Multiplying by RUI10 proves $\mathcal T_A\mathcal R_{A,s,D}=I$.
The numerical polynomial coefficients use the original source and
target variables, as in AIE21. The centre change from $c'$ to $c_k$
is the original shift $y\mapsto y-4i$, of norm at most
$\exp(4(H_D+2))$ by AKJ22.

Give the quotient $\mathbb C[S]_{\le D}/\mathbb C[S]_{<v}$ its
minimum norm from the source at $c_k$. The induced conductor
$\overline{\mathcal T}_A$ is an isomorphism onto
$\mathbb C[S']_{\le D-v}$ with its source norm at $c'$. Define
\[
 L_{D,s}=\frac{v!}{|\mu_v|}
          e^{4(H_D+2)}P_{D,s},\qquad
 N_{A,D}=\sum_{r,s}|a_{rs}|
             e^{|b_{rs}-4|(H_D+2)}.
 \tag{RUI13}
\]
For any $t$-dimensional target subspace, with original basis matrix
$X$, its exact quotient Gram ratio satisfies
\[
 N_{A,D}^{-2t}\det(X^*X)
 \le\det\bigl(X^*(\overline{\mathcal T}_A^{-1})^*
                       \overline{\mathcal T}_A^{-1}X\bigr)
 \le L_{D,s}^{2t}M_D^{2(D+1-t)}\det(X^*X).
 \tag{RUI14}
\]
To prove the upper bound, use the actual lift RUI12. Apply
$\mathcal R_{v,s,c'}$ to the $t$ basis columns, retaining its
specified image subspace. RUI7 bounds the exterior power of the
inverse $B_{D+1}$ on that image by $M_D^{D+1-t}$. The scalar,
right-inverse norm, and centre translation give $L_{D,s}^t$.
Projection to the minimum lift in the quotient is an orthogonal
projection on the same source Hilbert space, so its exterior norm
is at most one. This proves the complete upper Gram inequality.
For the lower bound, apply the forward original conductor to the
minimum lifts. AKJ24--25 gives its norm $N_{A,D}$; taking its
$t$th exterior power proves the result. This proof retains the
coupled Gram and the full quotient metric.

Every factor in RUI13 is explicit in the fixed original data.
For the asymptotic assertion restrict now to $v\le D\le2q$.
Since $v\le80$, $s\in\{1,k\}$ and $t\le D+1-v\le2q+1$,
both logarithmic endpoints of RUI14 have absolute value at most
\[
 2t\bigl(|\log L_{D,s}|+|\log N_{A,D}|\bigr)
 +2(D+1)\log M_D
 =O_{\rm actual}(q\log(q+1)+q\log(k+1))=o(kq).
 \tag{RUI15}
\]
For the logarithm of $N_{A,D}$, use
$A_{\rm abs}\le N_{A,D}\le
A_{\rm abs}\exp(\max|b_{rs}-4|(H_D+2))$.
For $L_{D,s}$, its displayed factors and the positive lower bound
$P_{D,s}\ge1$ give the same logarithmic estimate, retaining
$v!/|\mu_v|$. This proves the absolute bounds even if either
constant is below one. The harmonic bound and $q=(k+1)^2$
then prove the final limit. No dimension-$q$ scalar comparison has
silently acquired an extra factor of $q$ in this argument.

At the original coefficient cutoff $D=q-1$, RUI14 is the sharpened
version of AIE29, on precisely the same conductor quotient and
original coefficient subspaces. At a larger source cutoff $D=N$,
its literal symbol is $B_{N+1}$, and RUI10 proves the exact extension.
The old symbol $B_q$ cannot be substituted without recording the
finite difference. Passing this estimate through the original
remainder square requires the actual image of the relation ideal;
that square, its image graph, and the endpoint quotient metrics
are retained in the succeeding calculation.

% END CURRENT RUI

\clearpage
% BEGIN CURRENT CEP
% Original conductor endpoint square, graph quotient and determinant transfer.
\section{The original conductor at every source endpoint}
\label{sec:CEP}

Retain the complete simple-quartet conductor of OCF and AKJ, its
actual original period matrix, and $k=4l+1\ge9$ on the proved
five-orbit locus. Write
\[
 \begin{gathered}
 q=(k+1)^2,\quad k'=k-8,\quad q'=(k-7)^2,\quad
 \Delta=q-q'=16k-48,\\
 c=k/2,\quad c'=k/2-4,\quad
 \beta_{ij}=c+(2i-k)\delta+i(2j-k)\gamma,\\
 \beta'_{ab}=c'+(2a-k')\delta+i(2b-k')\gamma,\\
 A=\sum_{r,s=0}^8a_{rs}e_{rs}^{(8)}\ne0,\quad
 b_{rs}=4+(2r-8)\delta+i(2s-8)\gamma,\\
 \mu_j=\sum_{r,s}a_{rs}b_{rs}^j,\qquad
 v=\min\{j:\mu_j\ne0\}\le80.
 \end{gathered}\tag{CEP1}
\]
The factor $i$ before an imaginary term is the imaginary unit;
subscript indices are integers in their displayed ranges. The
original polynomials are
$\chi_k(S)=\prod_{i,j=0}^k(S-\beta_{ij})$ and
$\chi_{k'}(S')=\prod_{a,b=0}^{k'}(S'-\beta'_{ab})$.
Both have distinct roots. Set
$E_k=\mathbb C[S]/(\chi_k)$ and
$E_{k'}=\mathbb C[S']/(\chi_{k'})$.
Every source below has the original order $s\in\{1,k\}$ and
full mass $M_s=(2\pi)^{s/2}$, with $b=s/2$ and
$h_n=M_s n!(b)_n$. In particular the lower grid does not change
source order $k$ to $k-8$.

\subsection{The full source, remainder and primary-coordinate square}

For every polynomial define the original finite-translation map
\[
 (\mathcal T_Ap)(S')=\sum_{r,s=0}^8a_{rs}p(S'+b_{rs}).
 \tag{CEP2}
\]
For $0\le a,b\le k'$ and $0\le r,s\le8$, direct substitution gives
$\beta'_{ab}+b_{rs}=\beta_{a+r,b+s}$. All indices on the right lie
between zero and $k$. Consequently
\[
 (\mathcal T_A(\chi_k f))(\beta'_{ab})=0
 \quad\hbox{for every polynomial }f.
 \tag{CEP3}
\]
Division by the monic $\chi_{k'}$ and evaluation at its distinct roots
then show $\chi_{k'}\mid\mathcal T_A(\chi_k f)$. Thus CEP2 induces
the specified complex-linear map
$\overline{\mathcal T}_A:E_k\longrightarrow E_{k'}$.

In the original ordered primary coordinates its matrix is
\[
 C_{(a,b),(i,j)}
 =\begin{cases}
 a_{i-a,j-b},&0\le i-a,j-b\le8,\\
 0,&\text{otherwise}.
 \end{cases}
 \qquad C\in\operatorname{Mat}_{q'\times q}(\mathbb C).
 \tag{CEP4}
\]
Evaluation of CEP2 proves this entry formula, including its absence
of conjugation or binomial coefficients. Its ordinary transpose
$C^{\mathsf T}$ is the exact biform coefficient multiplication
by the nonzero polynomial $A$. That multiplication is injective
because the biform ring embeds in a polynomial integral domain.
Therefore $C$ has rank $q'$ and is surjective. With the original
ordered Vandermonde matrices $V_\beta,V_{\beta'}$, the polynomial
coordinate matrix is exactly $V_{\beta'}^{-1}CV_\beta$.
The kernel is the AKJ15 space
\[
 V_A=\{[p]\in E_k:\chi_{k'}\mid\mathcal T_Ap\},
 \qquad \dim V_A=\Delta.
 \tag{CEP5}
\]
CEP3 proves independence of the representative, and CEP4 proves
the dimension. The original common kernel $K=\ker\Lambda$ is
contained in $V_A$ by the full conductor inclusion OCF17.

For every endpoint
$N\in\{q-1,q,2q-1,2q\}$, put $d=N-v$ and define
\[
 \begin{gathered}
 \mathcal P_N=\mathbb C[S]_{\le N},\qquad
 \mathcal P'_d=\mathbb C[S']_{\le d},\\
 U_N:\mathcal P_N\longrightarrow E_k,\quad p\longmapsto[p],
 \qquad
 V_d:\mathcal P'_d\longrightarrow E_{k'},\quad g\longmapsto[g].
 \end{gathered}
\]
The top horizontal map is precisely
$\mathcal T_{A,N}:\mathcal P_N\longrightarrow\mathcal P'_{N-v}$.
Taylor's finite identity and the first nonzero moment show that
CEP2 lowers every degree at least $v$ by exactly $v$ and kills
exactly $\mathbb C[S]_{<v}$. The exact commutative square is
\[
 \begin{array}{ccc}
 \mathcal P_N &\xrightarrow{\ \mathcal T_{A,N}\ }&\mathcal P'_{N-v}\\
 {\scriptstyle U_N}\downarrow&&\downarrow{\scriptstyle V_{N-v}}\\
 E_k&\xrightarrow{\ \overline{\mathcal T}_A\ }&E_{k'} .
 \end{array}\tag{CEP6}
\]
The lower remainder is surjective because
$N-v\ge q-1-v\ge q'-1$; the latter inequality is
$\Delta\ge v$, and $\Delta\ge96>80\ge v$.

The map is complex-linear with its following full action defect.
If $M_S,M_{S'}$ denote the original multiplication operators, then
\[
 CM_S-M_{S'}C=C_b,\qquad
 (C_b)_{(a,b),(i,j)}=
 \begin{cases}a_{i-a,j-b}b_{i-a,j-b},&0\le i-a,j-b\le8,\\
 0,&\text{otherwise}.
 \end{cases}
 \tag{CEP7}
\]
Indeed each entry is multiplied by
$\beta_{ij}-\beta'_{ab}=b_{i-a,j-b}$. In the original centred
coordinates $Y=(M_S-cI)/i$, $Y'=(M_{S'}-c'I)/i$, this becomes
$CY-Y'C=(C_b-4C)/i$. Thus the exact relation to multiplication
retains its coefficient defect.

\subsection{The literal finite inverse at the larger cutoffs}

On $\mathcal P_N$, define the literal finite symbol
\[
 B_{N+1}(z)=1+\sum_{j=v+1}^{N}
              \frac{v!\mu_j}{j!\mu_v}z^{j-v},\qquad
 \mathcal T_{A,N}=\frac{\mu_v}{v!}
                   \partial_S^vB_{N+1}(\partial_S).
 \tag{CEP8}
\]
The expression on the right is evaluated in the output variable
$S'$. Taylor expansion of each term in CEP2 through degree $N$
proves this identity coefficient by coefficient. Its constant term
one makes $B_{N+1}(\partial_S)$ invertible: its difference from
the identity lowers degree, so the inverse is its finite nilpotent
Neumann sum through power $N$. A right inverse of $\partial_S^v$
on the polynomial spaces exists by integration of each monomial;
composition proves surjectivity of the top map in CEP6.

The relation to the earlier literal $B_q$ of AKJ14 is
\[
 \begin{split}
 D_{N,q}&=B_{N+1}(\partial_S)-B_q(\partial_S)
       =\sum_{j=q}^{N}\frac{v!\mu_j}{j!\mu_v}\partial_S^{j-v},\\
 B_{N+1}(\partial_S)^{-1}-B_q(\partial_S)^{-1}
       &=-B_{N+1}(\partial_S)^{-1}D_{N,q}B_q(\partial_S)^{-1},\\
 \mathcal T_{A,N}-\frac{\mu_v}{v!}\partial_S^vB_q(\partial_S)
       &=\sum_{j=q}^{N}\frac{\mu_j}{j!}\partial_S^j.
 \end{split}\tag{CEP9}
\]
Empty sums are zero. The first two differences have image contained
in $\mathbb C[S]_{\le N-q+v}$; the last has image contained in
$\mathbb C[S']_{\le N-q}$. The inverse formula follows by
multiplication, and its degree assertion follows since both inverse
factors preserve every degree prefix. At $N=q-1$ all differences
vanish. At the other three endpoints these are the full retained
defects; their ranks are not assigned the cutoff-independent bound
$v$.

For comparison, the entire symbol
$B_\infty(z)=v!\sum_{r,s}a_{rs}e^{b_{rs}z}/(\mu_vz^v)$,
with removable value one at zero, obeys at each cutoff
\[
 B_\infty(\partial_S)-B_{N+1}(\partial_S)
 =\sum_{n=N+1-v}^{N}
       \frac{v!\mu_{n+v}}{(n+v)!\mu_v}\partial_S^n.
 \tag{CEP10}
\]
Its image and the image of the corresponding inverse difference
lie in $\mathbb C[S]_{<v}$. This follows from the first derivative
degree $N+1-v$, and from the inverse identity with the same order
as AIE15. Thus the rank-at-most-$v$ comparison is with the
correctly growing finite symbol in CEP8.

For completeness the original-source inverse estimate extends
directly to these literal symbols. Retain the exact constants
$A_*^{\rm abs},B_*^{\rm sh},C_v,R_0,\rho_0$ of AIE4.
For every $|z|\le R_0$ the finite coefficient sum in CEP8 is at most
the full absolute moment sum:
\[
 |B_{N+1}(z)-1|
 \le \frac{v!}{|\mu_v|}
          \sum_{j=v+1}^N\frac{|\mu_j|}{j!}|z|^{j-v}
 \le C_v\frac{B_*^{\rm sh}|z|}{v+1}e^{B_*^{\rm sh}|z|}
 \le\frac12.
 \tag{CEP11}
\]
The middle inequality follows by writing $j=v+1+h$ and using
$v!/(v+1+h)!\le1/((v+1)h!)$. Hence
$F_{N+1}(z)=1/B_{N+1}(-i\arctan z)$ has coefficient absolute
values at most $2\rho_0^{-a}$ on degree $a\ge1$, and constant
coefficient one. This is Cauchy's formula on $|z|=\rho_0$, since
$|\arctan z|\le\operatorname{arctanh}\rho_0=R_0$.

In the original orthonormal Gamma basis, the lowering matrix
$L_s\varphi_n=\sqrt{w_n/w_{n-1}}\varphi_{n-1}$,
$w_n=n!/(b)_n$, satisfies
$\partial_y=\arctan L_s$ and $L_s^{N+1}=0$.
Evaluation of the power-series inverse in this nilpotent matrix
therefore gives the exact finite operator
$F_{N+1}(L_s)$, with entries
$[z^{n-m}]F_{N+1}(z)\sqrt{w_n/w_m}$.
The chain rule $S=c+iy$ supplies the retained $-i$.
Bounding the absolute row and column sums gives
\[
 \|B_{N+1}(\partial_S)^{-1}\|_{s,c}
 \le \alpha_{N,s}\left(1+2\sum_{a=1}^N\rho_0^{-a}\right),
 \qquad
 \alpha_{N,s}=\max_{0\le m\le n\le N}\sqrt{w_n/w_m}.
 \tag{CEP12}
\]
This proves the larger-cutoff assertion for the actual finite
symbol rather than substituting $N$ into an identity for $B_q$.
Likewise the explicit AIE18 raising right inverse on these degree
spaces, followed by $i^v$, $B_{N+1}(\partial_S)^{-1}$ and
$v!/\mu_v$, is a right inverse of CEP8. Its comparison with the
monomial integral remains exactly the degree-below-$v$ correction
in AIE25, now with this literal inverse.

\subsection{The exact source operator matrix and its full determinant}

Use the unchanged original orthonormal polynomials
$u_n(S)=p_n((S-c)/i)/\sqrt{h_n}$ in $\mathcal P_N$ and
$u'_m(S')=p_m((S'-c')/i)/\sqrt{h_m}$ in $\mathcal P'_{N-v}$.
Let $A_N$ denote the matrix of $\mathcal T_{A,N}$ in these two
frames. Define the analytic coefficient series at zero
\[
 t(z)=\sum_{r,s}a_{rs}
       \exp\left(\frac{b_{rs}-4}{i}\arctan z\right)
       =\sum_{\ell\ge0}t_\ell z^\ell .
 \tag{CEP13}
\]
Then its exact entries are
\[
 (A_N)_{m,n}=
 \begin{cases}
 t_{n-m}\sqrt{w_n/w_m},&n\ge m,\\
 0,&n<m,
 \end{cases}
 \quad
 0\le m\le N-v,\quad0\le n\le N.
 \tag{CEP14}
\]
To prove this, the original generating identity
$\sum p_n(y)z^n/n!=(1+z^2)^{-b/2}e^{y\arctan z}$
shows that translation by $\eta$ multiplies the series by
$e^{\eta\arctan z}$. Extracting its finite coefficients and
dividing by the complete $\sqrt{h_n}$ gives
$\sqrt{w_n/w_m}$ after the two copies of $M_s$ cancel at this
specified step. Evaluation of CEP2 at $c'+iy$ uses
$p(c+i(y+(b_{rs}-4)/i))$, proving the exact centre shift in CEP13.

For $\ell<v$, the shifted moment
$\sum a_{rs}(b_{rs}-4)^\ell$ is zero by the binomial theorem and
the definition of $v$; at $\ell=v$ it equals $\mu_v$.
Since $\arctan z=z+O(z^3)$, this gives
\[
 t_0=\cdots=t_{v-1}=0,\qquad
 t_v=\frac{(-i)^v\mu_v}{v!}\ne0.
 \tag{CEP15}
\]
The first $v$ columns of $A_N$ vanish. Its remaining square block
$\widehat A_N$, with columns $v,\ldots,N$, is upper triangular
with diagonal $t_v\sqrt{w_{m+v}/w_m}$.
Thus its rank is $N-v+1$, its kernel in the source is exactly
the first $v$ orthonormal coordinates, and
\[
 \begin{split}
 \det\widehat A_N
   &=\left(\frac{(-i)^v\mu_v}{v!}\right)^{N-v+1}
                   \prod_{m=0}^{N-v}\sqrt{\frac{w_{m+v}}{w_m}},\\
 \det(A_NA_N^*)
   &=\left|\frac{\mu_v}{v!}\right|^{2(N-v+1)}
                   \prod_{m=0}^{N-v}\frac{w_{m+v}}{w_m}.
 \end{split}\tag{CEP16}
\]
This retains the complex determinant phase and the original moment.
The operator norm also has the exact original-source upper bound
\[
 \|A_N\|\le
 M_{A,N}:=\sum_{r,s}|a_{rs}|
               e^{|b_{rs}-4|(H_N+2)},\qquad
 H_N=\sum_{j=1}^N1/j.
 \tag{CEP17}
\]
Indeed AKJ22 bounds each exact translation in CEP13 by its
displayed factor, and the triangle inequality gives the sum.

There is a closed four-endpoint expression for CEP16. For the
signs $\sigma=(1,1,-1,-1)$ at the ordered endpoints, write
$\mathcal D_A=\sum_N\sigma_N\log\det(A_NA_N^*)$.
Cancellation of the full products gives
\[
 \begin{split}
 \mathcal D_A
  &=-4q\log\left|\frac{\mu_v}{v!}\right|+\mathcal W_{v,s},\\
 \mathcal W_{v,s}
  &=\sum_{a=0}^{v-1}\left(
      \log w_{q-1-a}+\log w_{q-a}
      -\log w_{2q-1-a}-\log w_{2q-a}\right).
 \end{split}\tag{CEP18}
\]
For $v=0$ the second sum is zero. The product identity used here is
$\prod_{m=0}^{N-v}w_{m+v}/w_m
 =\prod_{a=0}^{v-1}w_{N-a}/\prod_{a=0}^{v-1}w_a$,
obtained by cancellation of consecutive products; it remains
valid when the two index intervals overlap.

Put $c_b^{\log}=\max\{1,b-1\}$. For every $j\ge1$ and $b\ge1/2$,
\[
 \left|\log\frac{j}{b+j-1}\right|\le\frac{c_b^{\log}}j .
 \tag{CEP19}
\]
For $b\ge1$ use $\log(1+(b-1)/j)\le(b-1)/j$.
For $1/2\le b<1$, put $a=1-b\le1/2$ and use
$-\log(1-a/j)\le a/(j-a)\le1/j$; the last inequality is
$aj\le j-a$, which follows from $a\le1/2$ and $j\ge1$.
Summing CEP19 over the two intervals of length $q$ in each
summand of CEP18 proves the completely explicit estimate
\[
 |\mathcal W_{v,s}|
 \le \frac{2v\,c_b^{\log}q}{q-v},\qquad
 |\mathcal D_A|
 \le4q\left|\log\left|\frac{\mu_v}{v!}\right|\right|
             +\frac{2v\,c_b^{\log}q}{q-v}.
 \tag{CEP20}
\]
All denominators are positive since $q\ge100>v$.
For either original source and the fixed actual quartet and period,
the displayed bound is $O_{\rm actual}(q+k)=o(kq)$.

\subsection{The induced original quotient metric and its correction}

Let the matrices of the two remainder maps in the orthonormal source
frames and original primary target frames be denoted by the same
$U_N,V_d$ as their maps. Explicitly
\[
 (U_N)_{(i,j),n}=\frac{p_n((\beta_{ij}-c)/i)}{\sqrt{h_n}},
 \qquad
 (V_d)_{(a,b),m}=\frac{p_m((\beta'_{ab}-c')/i)}{\sqrt{h_m}} .
 \tag{CEP21}
\]
They have sizes $q$ by $N+1$ and $q'$ by $N-v+1$,
respectively. Their first $q$, respectively $q'$, columns have
nonzero evaluation Vandermonde determinant, so both have full row
rank. The square CEP6 is the precise matrix identity
\[
 C U_N=V_dA_N,\qquad
 R_N=U_NU_N^*,\quad R'_d=V_dV_d^*,\quad
 G_N=R_N^{-1},\quad G'_d=(R'_d)^{-1}.
 \tag{CEP22}
\]
Here $G_N$ is the original upper primary quotient metric of OKM8,
and $G'_d$ is the lower-grid metric in the same source order $s$.
The induced metric on $E_k/V_A$, in its specified target coordinates
$E_{k'}$, is
\[
 Q_N=(C G_N^{-1}C^*)^{-1}
       =(V_dA_NA_N^*V_d^*)^{-1}.
 \tag{CEP23}
\]
To prove the first assertion, the minimum of $x^*G_Nx$ under
$Cx=y$ is attained at
$x=G_N^{-1}C^*(CG_N^{-1}C^*)^{-1}y$.
Every other lift differs by an element of $\ker C$ orthogonal
to this lift in the $G_N$ metric, proving the minimum and its
Gram matrix. The second equality is CEP22.

The exact correction to the lower original Gamma quotient metric is
\[
 \begin{gathered}
 W_d=V_d^*(R'_d)^{-1/2},\quad W_d^*W_d=I_{q'},\\
 K_{A,N}=W_d^*A_NA_N^*W_d>0,\qquad
 Q_N=(R'_d)^{-1/2}K_{A,N}^{-1}(R'_d)^{-1/2},\\
 \log\det Q_N=\log\det G'_d-\log\det K_{A,N}.
 \end{gathered}\tag{CEP24}
\]
All square roots are the unique positive Hermitian square roots.
These formulas follow by multiplying out the definition of
$K_{A,N}$ and inverting. In particular the passage between the
induced metric and the lower original source metric retains this
actual compressed covariance.

If $Z_d$ is any orthonormal basis of $\ker V_d$, completing $W_d$
to a unitary matrix, the further exact determinant factorization is
\[
 \det K_{A,N}
   =\det(A_NA_N^*)\,
        \det\bigl(Z_d^*(A_NA_N^*)^{-1}Z_d\bigr).
 \tag{CEP25}
\]
For a proof put $D=A_NA_N^*>0$ and write its block matrix in
the unitary frame $[W_d,Z_d]$. The upper left block is $K_{A,N}$.
Block Gaussian elimination gives
$\det D=\det K_{A,N}\det(D_{22}-D_{21}K_{A,N}^{-1}D_{12})$.
The lower right block of $D^{-1}$ is the inverse of that Schur
complement, proving CEP25. The empty-complement case has determinant
one. The dimension of this retained complement is
$N-v+1-q'$.

Choose any fixed full frame $I_V:\mathbb C^\Delta\hookrightarrow E_k$
of $\ker C$ and a fixed right section $L_C$ of $C$.
Then $F_V=[I_V,L_C]$ is invertible and
\[
 \det(I_V^*G_NI_V)
   =|\det F_V|^2\,\frac{\det G_N}{\det Q_N}
   =|\det F_V|^2\,\frac{\det G_N}{\det G'_d}\det K_{A,N}.
 \tag{CEP26}
\]
Indeed $CF_V=[0,I]$, so the lower right block of
$(F_V^*G_NF_V)^{-1}$ is $CG_N^{-1}C^*=Q_N^{-1}$.
The Schur determinant argument just used gives CEP26.
The frame factor is independent of source cutoff and cancels
only after applying the four signs.

For the actual inclusion $K\subset V_A$, choose its original
coefficient frame, and let $\Xi:\mathbb C^\Delta\twoheadrightarrow
\mathbb C^{\Delta-r}$ have its coordinate kernel, $r=8k-16$,
as in AKJ18. With $H_{V,N}=I_V^*G_NI_V$, the quotient metric is
$Q_{\Xi,N}=(\Xi H_{V,N}^{-1}\Xi^*)^{-1}$.
The same proof with a frame of $\ker\Xi$ and a right section gives,
with all constant frame factors cancelled by the endpoint signs,
\[
 F_K^{(s)}
 =\sum_N\sigma_N\log\det G_N
   -\sum_N\sigma_N\log\det G'_{N-v}
   +\sum_N\sigma_N\log\det K_{A,N}
   -\sum_N\sigma_N\log\det Q_{\Xi,N}.
 \tag{CEP27}
\]
Thus the conductor quotient and the remaining original observation
quotient have both been retained in the determinant relation.
The lower source cutoffs in this equation are exactly
$q-1-v,q-v,2q-1-v,2q-v$. Relative to the lower grid's own
four cutoffs they have the respective shifts
$\Delta-v,\Delta-v,2\Delta-v,2\Delta-v$.

\subsection{An exact graph quotient with controlled exterior metric error}

In original orthonormal coordinates define
\[
 S_N:\mathcal P_N\longrightarrow
         \mathbb C^v\oplus\mathcal P'_{N-v},\qquad
 S_Np=(P_{<v}p,\mathcal T_{A,N}p).
 \tag{CEP28}
\]
Here $P_{<v}$ records the first $v$ original upper orthonormal
coefficients, and the target carries the orthogonal direct-sum norm
with those coordinates and the unchanged lower Gamma source.
The first $v$ source coordinates span the polynomials of degree
below $v$. Thus CEP14--16 give its exact block matrix
$S_N=\operatorname{diag}(I_v,\widehat A_N)$ in the ordered
source split at degree $v$. In particular it is invertible, and
\[
 \begin{gathered}
 \|S_N\|\le M_N:=\max\{1,M_{A,N}\},\\
 |\det S_N|=\left|\frac{\mu_v}{v!}\right|^{N-v+1}
              \prod_{m=0}^{N-v}\sqrt{\frac{w_{m+v}}{w_m}},\\
 |\log|\det S_N||
 \le L_N:=(N-v+1)\left|\log\left|\frac{\mu_v}{v!}\right|\right|
                      +v c_b^{\log}H_N .
 \end{gathered}\tag{CEP29}
\]
The last bound follows by the consecutive-product identity preceding
CEP19 and $|\log w_j|\le c_b^{\log}H_j$: the two lists of $v$
logarithms are each bounded by $v c_b^{\log}H_N$, followed by
the factor $1/2$.

The upper and lower source ideals at this cutoff are
\[
 I_N=\chi_k\mathbb C[S]_{\le N-q},\qquad
 I'_N=\chi_{k'}\mathbb C[S']_{\le N-v-q'} .
 \tag{CEP30}
\]
A negative degree specifies the zero subspace.
Let $\mathcal W_N=U_N^{-1}(V_A)$. The exact image and ideal graph
under CEP28 are
\[
 \begin{gathered}
 S_N\mathcal W_N=\mathbb C^v\oplus I'_N,\\
 S_N I_N=\{(P_{<v}h,\mathcal T_{A,N}h):h\in I_N\}
                    \subset\mathbb C^v\oplus I'_N,\\
 V_A\ \xrightarrow{\ \simeq\ }\
   (\mathbb C^v\oplus I'_N)/(S_N I_N),\qquad
 U_Np\longmapsto[S_Np].
 \end{gathered}\tag{CEP31}
\]
The first equality follows because $\mathcal W_N$ is the preimage
of $I'_N$ under the top map of CEP6, and $S_N$ is invertible.
CEP3 proves the ideal inclusion. If representatives differ by
$I_N$, their images differ by $S_NI_N$, so the last map is
well defined. Conversely every class in the target has a preimage
under $S_N$ in $\mathcal W_N$, and the kernel consists exactly
of $I_N$. This proves the isomorphism. Its norm on the left is
the restriction of $G_N$: all lifts of a member of $V_A$ already
belong to $\mathcal W_N$, so their minimum source norm is the
quotient norm of $\mathcal W_N/I_N$.
The low-degree graph component $P_{<v}h$ is explicitly retained.

For reference the ideal dimensions at the four original endpoints are
\[
 \begin{array}{c|c|c|c}
 N&\dim I_N&\dim I'_N&
       \dim((\mathbb C^v\oplus I'_N)/S_NI_N)\\ \hline
 q-1&0&\Delta-v&\Delta\\
 q&1&\Delta+1-v&\Delta\\
 2q-1&q&q+\Delta-v&\Delta\\
 2q&q+1&q+\Delta+1-v&\Delta .
 \end{array}\tag{CEP32}
\]
These follow by counting monomials in CEP30; $S_N$ preserves
dimensions. Moreover $\mathcal T_A|_{I_N}$ is injective because
a nonzero member has degree at least $q>v$ and its image has
degree exactly $v$ less. Hence the second coordinate of the
retained ideal graph already has dimension $\dim I_N$.

Here is the complete exterior norm estimate for this exact map.
For an invertible complex matrix $S$ of dimension $n$ and
$0\le t\le n$, its singular values
$\sigma_1\ge\cdots\ge\sigma_n>0$ give
\[
 \|\bigwedge^t S^{-1}\|
 =\frac{\|\bigwedge^{n-t}S\|}{|\det S|}.
 \tag{CEP33}
\]
Indeed the numerator is $\sigma_1\cdots\sigma_{n-t}$ and
the denominator is the product of all singular values; their
ratio is the product of the $t$ largest singular values of
$S^{-1}$. This includes $t=0,n$.
For $n=N+1$, CEP29 therefore gives the finite bounds
\[
 \begin{split}
 \log^+\|\bigwedge^t S_N\|&\le t\log M_N,\\
 \log^+\|\bigwedge^t S_N^{-1}\|
        &\le(N+1-t)\log M_N+L_N .
 \end{split}\tag{CEP34}
\]
Here $\log^+x=\max\{0,\log x\}$. The right sides are nonnegative.
Thus for all these ranks both bounds are at most
\[
 E_N=(N+1)\log M_N+L_N.
 \tag{CEP35}
\]
With $R_{\rm sh}=\max_{r,s}|b_{rs}-4|$ and
$A_*^{\rm abs}=\sum|a_{rs}|$, equation CEP17 gives
$\log M_N\le\log^+A_*^{\rm abs}+R_{\rm sh}(H_N+2)$.
For the fixed original quartet and period, $v,R_{\rm sh},\mu_v$
and $A_*^{\rm abs}$ are fixed while $c_b^{\log}=O(k)$.
Since $N\le2q$ and $H_N\le1+\log N$, all four explicit bounds
in CEP35 satisfy
\[
 E_N=O_{\rm actual}(q\log(q+1)+k\log(q+1))=o(kq).
 \tag{CEP36}
\]
This uses only the forward norm and the exact determinant of
the original conductor map.

We now prove the exact relation to the metric of the original
kernel, including the quotient projection. In a finite-dimensional
Hilbert space let $I\subset W$ and let $S$ be invertible.
The induced map $W/I\longrightarrow SW/SI$ sends a minimal lift
$x\in W\cap I^\perp$ to the orthogonal projection of $Sx$
off $SI$. Orthogonal projection is a contraction on every exterior
power: in an orthonormal basis its singular values are zeros and
ones. Therefore, on any $t$-dimensional subspace of $W/I$, the
exterior norm of the induced map is at most
$\|\bigwedge^t S\|$. Applying the identical argument to $S^{-1}$
gives the inverse bound. The minimal-lift identification is
isometric since every other lift differs by a perpendicular
element of $I$, whose squared norm adds. This proves the
quotient assertion with no choice of coefficient norm.

Apply this result with $W=\mathcal W_N$, $I=I_N$, and the
actual rank-$r$ subspace $K\subset V_A$. The target subspace is
the explicit image of $K$ under CEP31; denote its Gram matrix
in the transported original frame by $\widetilde H_{K,N}$.
Let $I_K:\mathbb C^r\hookrightarrow E_k$ be that original
primary frame, and write its original minimum source lift as
$Y_N=U_N^*G_N I_K$.
Let $Z_N$ be the matrix in the original orthonormal source
coordinates of the independent columns
$\chi_k(S)S^a$, $0\le a\le N-q$. If there are no columns,
all following products involving them are zero.
The exact target orthogonal projection and Gram matrix are
\[
 \begin{gathered}
 P_N^{\rm graph}
  =I-S_NZ_N(Z_N^*S_N^*S_NZ_N)^{-1}Z_N^*S_N^*,\\
 \widetilde H_{K,N}
       =Y_N^*S_N^*P_N^{\rm graph}S_NY_N,\qquad
 H_{K,N}=Y_N^*Y_N=I_K^*G_NI_K .
 \end{gathered}\tag{CEP37}
\]
The inverted Gram matrix is positive definite because its
columns are independent and $S_N$ is invertible. Multiplication
shows that $P_N^{\rm graph}$ is Hermitian, idempotent, kills
$S_NI_N$ and fixes its orthogonal complement. Thus CEP37 is
precisely the minimum quotient Gram of CEP31, and its kernel
frame has full rank because CEP31 is injective.
Replacing the lift $Y_N$ by any lift of the same primary frame
adds columns in $I_N$, which the projection kills; hence every
choice of lift gives this same matrix.

Applying the proved exterior quotient estimate and squaring
volumes gives
\[
 \begin{split}
 -2\log\|\bigwedge^r S_N^{-1}\|
 &\le \log\det\widetilde H_{K,N}-\log\det H_{K,N}\\
 &\le2\log\|\bigwedge^r S_N\|,\\
 \left|\log\det\widetilde H_{K,N}-\log\det H_{K,N}\right|
 &\le2E_N .
 \end{split}\tag{CEP38}
\]
In detail, divide the chosen frame by its positive Gram square
root so it is orthonormal in the source quotient. The squared
norm of its exterior image is exactly the determinant ratio.
The forward and inverse exterior estimates prove the two
inequalities, and CEP34--35 prove the last assertion.

Consequently the following reduction of the actual four-endpoint
kernel determinant is proved:
\[
 \begin{gathered}
 \widetilde F_K^{(s)}
       =\sum_N\sigma_N\log\det\widetilde H_{K,N},\\
 |F_K^{(s)}-\widetilde F_K^{(s)}|
       \le2\sum_{N\in\{q-1,q,2q-1,2q\}}E_N=o(kq).
 \end{gathered}\tag{CEP39}
\]
The frame is the original fixed frame of $K$ at every endpoint;
CEP37 gives all four transported matrices explicitly.
Their target spaces retain the original ideal graphs, the
different source cutoffs, all conductor coefficients, and the
full original Gamma masses. CEP39 proves that the leading
$kq$ coefficient is unchanged by this particular graph-quotient
transport. It supplies no value for that coefficient: its
remaining determinant is the displayed, completely specified
four-endpoint graph quotient.

\subsection{All ideal graph entries as original polynomial products}

Each individual shift in CEP2 induces the unital algebra map
$\sigma_{rs}:E_k\to E_{k'}$, $[p]\mapsto[p(S'+b_{rs})]$.
The divisibility in CEP3, applied to the individual shifted
polynomial, proves independence of the representative; substitution
preserves multiplication and one. Thus
$\overline{\mathcal T}_A=\sum a_{rs}\sigma_{rs}$ is its exact
linear combination. For any $x,y\in E_k$, the complete
multiplicative defect of that linear combination is
\[
 \overline{\mathcal T}_A(xy)
   -\overline{\mathcal T}_A(x)\overline{\mathcal T}_A(y)
 =\sum_{r,s}a_{rs}\sigma_{rs}(x)\sigma_{rs}(y)
   -\sum_{r,s,t,u}a_{rs}a_{tu}
                        \sigma_{rs}(x)\sigma_{tu}(y).
 \tag{CEP40}
\]
This follows by applying the proved individual multiplicativity
and expanding the product of the two complete sums.

There is a complete finite product for the map on the source ideal.
For each original shift put
\[
 d_{rs}(S')=
 \prod_{\substack{0\le i,j\le k\\
  (i,j)\notin\{r,\ldots,r+k'\}\times\{s,\ldots,s+k'\}}}
                   (S'+b_{rs}-\beta_{ij}).
 \tag{CEP41}
\]
This is monic of degree $\Delta$, and
$\chi_k(S'+b_{rs})=\chi_{k'}(S')d_{rs}(S')$.
Indeed CEP3's root addition identifies exactly the $q'$ factors
in the indicated rectangle with the factors of $\chi_{k'}$.
All remaining factors are precisely the product in CEP41.
Consequently for every $h\in\mathbb C[S]_{\le L}$,
$L=N-q$, the exact ideal map is
\[
 \begin{split}
 \mathcal T_{A,N}(\chi_kh)&=\chi_{k'}\,\mathcal U_{A,L}h,\\
 (\mathcal U_{A,L}h)(S')&=\sum_{r,s}a_{rs}
                     d_{rs}(S')h(S'+b_{rs})
                  \in\mathbb C[S']_{\le L+\Delta-v}.
 \end{split}\tag{CEP42}
\]
The equality follows by the product identity just proved. The
degree statement follows from CEP8: if $h$ has degree $t$ and
leading coefficient $h_t\ne0$, the leading coefficient of its
displayed image is
$h_t\binom{q+t}{v}\mu_v$ at degree $t+\Delta-v$.
This follows because $\chi_k,\chi_{k'}$ are monic and
$\mathcal T_A$ lowers degree exactly by $v$ with its stated
leading derivative. Thus $\mathcal U_{A,L}$ is injective.
At $L=-1$ the domain and image are zero.

This gives a further exact sequence, with all maps specified:
\[
 0\longrightarrow\mathbb C[S]_{<v}
 \longrightarrow V_A
 \longrightarrow
 \mathbb C[S']_{\le L+\Delta-v}/
                     \mathcal U_{A,L}\mathbb C[S]_{\le L}
 \longrightarrow0 .
 \tag{CEP43}
\]
The first map is the original quotient class. For the second,
lift a member of $V_A$ to $p\in\mathcal W_N$, divide
$\mathcal T_Ap$ by $\chi_{k'}$, and take its class in the last
space. A change of lift by $\chi_kh$ adds exactly
$\mathcal U_{A,L}h$ by CEP42, proving independence. Every target
polynomial times $\chi_{k'}$ has a preimage under the onto
map $\mathcal T_{A,N}$, proving surjectivity. If the image is
zero, subtract the corresponding $\chi_kh$; the remaining
polynomial is killed by $\mathcal T_A$ and hence has degree
below $v$. Its class injects because $v<q$.
This proves exactness. The last space has dimension $\Delta-v$
at every original endpoint, as is also obtained by subtracting
the two polynomial dimensions.

For a literal matrix version of the ideal graph, equip
$\mathbb C[S']_{\le L+\Delta-v}$ with the ordered lower
orthonormal Gamma coefficient basis, and write $M_{\chi',N}$
for multiplication by the full original $\chi_{k'}$ into
$\mathcal P'_{N-v}$. It is injective. The map
\[
 D_{\chi',N}=\operatorname{diag}(I_v,M_{\chi',N})
 :\mathbb C^v\oplus\mathbb C[S']_{\le L+\Delta-v}
       \longrightarrow\mathbb C^v\oplus I'_N
 \tag{CEP44}
\]
is an isomorphism onto the indicated subspace. Its pullback metric
is exactly
$G_{\chi',N}=\operatorname{diag}(I_v,M_{\chi',N}^*M_{\chi',N})$;
this follows by taking the squared norm of each component.
For the domain polynomial coefficient frame of $h$, let
$Z_{\chi,N}h$ be the original upper orthonormal coefficients
of $\chi_kh$. In these coordinates the complete graph relation
matrix is
\[
 \mathcal Z_{A,N}=
 \begin{pmatrix}
 P_{<v}Z_{\chi,N}\\[2pt]
 [\mathcal U_{A,L}]
 \end{pmatrix},\qquad
 D_{\chi',N}\mathcal Z_{A,N}=S_N Z_{\chi,N}.
 \tag{CEP45}
\]
The first component is the original low-degree projection,
and the second consists of the exact finite products and
translations in CEP41--42. The equality is CEP28 and CEP42.
It proves that the graph quotient in CEP31 is precisely the
quotient by $\operatorname{im}\mathcal Z_{A,N}$ in the
displayed metric $G_{\chi',N}$. All matrix inverses and
minimum projections in CEP37 can therefore be computed using
these full original polynomial products. No coefficient of
the ideal, source metric or low-degree graph component is
discarded by this additional coordinate isomorphism.

\subsection{The compressed conductor correction is subleading at every endpoint}

The exact exterior estimate also controls the large compressed
covariance in CEP24 without multiplying a scalar inverse bound
by $q'$. Write $n'=N-v+1$ and retain
$A_N=[0,\widehat A_N]$ with its exact square block in CEP16.
Since $W_d^*W_d=I_{q'}$, the coefficient matrix of its image under
$\widehat A_N^*$ gives
\[
 K_{A,N}
   =(\widehat A_N^*W_d)^*(\widehat A_N^*W_d).
 \tag{CEP46}
\]
The squared norm of the exterior product of the $q'$ columns
on the right is its determinant. The norm of the exterior
product of the columns of $W_d$ is one, because those columns
are orthonormal. The forward estimate therefore gives
$\det K_{A,N}\le\|\widehat A_N\|^{2q'}\le M_N^{2q'}$.
For the reverse inequality apply
$\bigwedge^{q'}\widehat A_N^{-*}$ to the same image columns;
their image is the orthonormal column wedge of $W_d$. Thus
\[
 \det K_{A,N}\ge
      \|\bigwedge^{q'}\widehat A_N^{-1}\|^{-2}.
 \tag{CEP47}
\]
The singular values of an adjoint are the same, so the adjoint
inverse in this argument has exactly the displayed exterior norm.
The singular complement identity CEP33 in dimension $n'$
gives
$\|\bigwedge^{q'}\widehat A_N^{-1}\|
 \le M_N^{n'-q'}/|\det\widehat A_N|$.
CEP29 bounds the absolute logarithm of the last determinant
by $L_N$. Consequently the complete two-sided estimate is
\[
 \begin{gathered}
 -2\bigl((N-v+1-q')\log M_N+L_N\bigr)
       \le\log\det K_{A,N}\le2q'\log M_N,\\
 |\log\det K_{A,N}|\le2E_N=o(kq)
 \quad\hbox{at every original endpoint and source order.}
 \end{gathered}\tag{CEP48}
\]
All ranks in this argument are their actual ranks,
$q'\le N-v+1$ as proved after CEP6. In particular the estimate
applies to the complete original lower quotient, whose rank
grows quadratically.

Define the fully specified total-source difference and the
remaining original observation quotient return by
\[
 \begin{split}
 \mathcal T_k^{(s)}
   &=\sum_N\sigma_N
          \bigl(\log\det G_N-\log\det G'_{N-v}\bigr),\\
 F_{\Xi,k}^{(s)}&=\sum_N\sigma_N\log\det Q_{\Xi,N},\\
 \mathcal E_{A,k}^{(s)}
   &=\sum_N\sigma_N\log\det K_{A,N}.
 \end{split}
\]
Equations CEP26--27 and CEP48 prove the exact identities and bounds
\[
 \begin{gathered}
 F_{V_A}^{(s)}
      =\mathcal T_k^{(s)}+\mathcal E_{A,k}^{(s)},\qquad
 F_K^{(s)}
      =\mathcal T_k^{(s)}-F_{\Xi,k}^{(s)}
                              +\mathcal E_{A,k}^{(s)},\\
 |\mathcal E_{A,k}^{(s)}|\le2\sum_N E_N=o(kq),\\
 \frac{F_K^{(s)}-\mathcal T_k^{(s)}+F_{\Xi,k}^{(s)}}{kq}
        \longrightarrow0 .
 \end{gathered}\tag{CEP49}
\]
Here $F_{V_A}^{(s)}$ uses its fixed original coefficient frame,
and $Q_{\Xi,N}$ is the exact quotient metric of CEP27.
The frame factors cancel with coefficients $1+1-1-1$ only in
these returns, as proved in CEP26. Thus the leading metric
effect of the full conductor source transport has now been
controlled. The remaining rank-$8k-32$ observation quotient
and the shifted lower-grid source determinants are still their
explicit original matrices; CEP49 carries them forward into
the next calculation.

% END CURRENT CEP


\clearpage
% BEGIN CURRENT CTV
% Exact total-source variation in the original root polynomials and Gamma norms.
\section{The total-source difference as an outer-strip quotient}
\label{sec:CTV}

Retain CEP1--49, with $m=1$, $k=4l+1\ge9$,
$q=(k+1)^2$, $k'=k-8$, $q'=(k-7)^2$,
$\Delta=q-q'=16k-48$, and the original fixed
$0<\delta<1/2$, $\gamma>2$. The conductor has its actual
first nonzero moment index $0\le v\le80$. At each of the four
original endpoints $N\in\{q-1,q,2q-1,2q\}$ put
$t=N-q+1$ and $d=N-v$. The source order remains
$s\in\{1,k\}$ at both grids, with
\[
 b=s/2,\quad M_s=(2\pi)^{s/2},\quad
 dm_s(y)=M_s\frac{2^b}{4\pi\Gamma(b)}
             |\Gamma(b/2+iy/2)|^2\,dy,\quad
 h_n=M_s n!(b)_n .
 \tag{CTV1}
\]
These are the complete original measure and monic norms of
OKM6--7. The source centres are $c=k/2$ and $c'=k/2-4$.
The total-source difference in CEP49 is
\[
 \mathcal T_k^{(s)}
   =\sum_N\sigma_N\bigl(\log\det G_N-\log\det G'_{N-v}\bigr),
 \qquad \sigma=(1,1,-1,-1),
 \tag{CTV2}
\]
where both matrices use their original primary frames.
All calculations below retain these four cutoffs and the same
two source orders.

\subsection{The exact nested root polynomials in the original real coordinate}

Define the full monic polynomials
\[
 \begin{gathered}
 P_k(y)=i^{-q}\chi_k(c+iy)
       =\prod_{a,b=0}^k(y-\omega_{ab}),\qquad
 \omega_{ab}=(2b-k)\gamma-i(2a-k)\delta,\\
 P_{k'}(y)=i^{-q'}\chi_{k'}(c'+iy)
       =\prod_{a,b=0}^{k'}(y-\omega'_{ab}),\qquad
 \omega'_{ab}=(2b-k')\gamma-i(2a-k')\delta .
 \end{gathered}\tag{CTV3}
\]
The original phases $i^{-q}$ and $i^{-q'}$ have modulus one;
they occur exactly as written in the coefficient maps.
For every lower pair one has
$\omega'_{ab}=\omega_{a+4,b+4}$. Hence, with the complete
outer index set
$\mathcal O_k=\{0,\ldots,k\}^2\setminus\{4,\ldots,k-4\}^2$,
the exact factorization is
\[
 P_k(y)=P_{k'}(y)P_{\partial,k}(y),\qquad
 P_{\partial,k}(y)=\prod_{(a,b)\in\mathcal O_k}(y-\omega_{ab}),
 \qquad \deg P_{\partial,k}=\Delta .
 \tag{CTV4}
\]
Direct substitution proves the root equality; counting the retained
central rectangle gives $q'$ roots, leaving $\Delta$ distinct
outer roots. The products are monic, so the root equality fixes
the factor without any additional constant.

Since $k'$ is odd, $|2a-k'|\ge1$ for every lower index.
For real $y$ this proves, factor by factor,
\[
 |y-\omega'_{ab}|\ge|\operatorname{Im}\omega'_{ab}|
        \ge\delta,\qquad
 |P_{k'}(y)|\ge\delta^{q'}.
 \tag{CTV5}
\]
Also $|\omega'_{ab}|\le
R'=k'\sqrt{\delta^2+\gamma^2}$. These estimates keep every
original factor; no root has been moved.

\subsection{The full monic source and relation determinant identities}

For a nonzero polynomial $P$ and an integer $a\ge0$ define the relation
moment determinant
\[
 D_a(P;s)=
 \det\left[\int_{\mathbb R}y^{i+j}|P(y)|^2\,dm_s(y)
                                    \right]_{0\le i,j<a},
 \qquad D_0(P;s)=1 .
 \tag{CTV6}
\]
Every nonempty matrix is positive definite: the measure is strictly
positive on the real line, a nonzero polynomial has only finitely
many zeros, and the squared modulus of its product with $P$
is positive on an open interval. All moments exist because the
original source has its finite exponential moments in CTV14 below.

In increasing monomial $y$ coordinates, the quotient by a monic
polynomial $P$ of degree $a$ from source degree $L\ge a-1$
has determinant
\[
 \det G^{y}_{P,L}
   =\frac{\prod_{n=0}^L h_n}{D_{L-a+1}(P;s)} .
 \tag{CTV7}
\]
Here the quotient frame is $1,y,\ldots,y^{a-1}$.
For proof use the source frame
\[(1,y,\ldots,y^{a-1},P,yP,\ldots,y^{L-a}P).\]
Its increasing-degree coefficient matrix has diagonal entries one
and no entries above the degree of a column, hence determinant
one. The last block is exactly the relation Gram in CTV6.
Minimizing over those last coordinates gives its Schur complement,
whose determinant is the full source determinant divided by the
relation determinant. Monic orthogonalization by the original
$p_n$ also has determinant one, so that full determinant is
$\prod_{n=0}^L h_n$. This proves CTV7 including the empty relation
case. It is the full-source analogue of BSL3, with the original
Gamma order retained.

Let $V_\beta,V_{\beta'}$ be the original ordered Vandermonde
matrices in the $S,S'$ variables. The exact coordinate pullback
$S=c+iy$ sends $S^j$ to $(c+iy)^j$ and has diagonal coefficients
$i^j$; its determinant is $i^{a(a-1)/2}$ in degree below $a$.
Its modulus is one, so its Gram congruence has determinant factor
one. The same statement holds with $c'$. Evaluation in primary
coordinates then gives
\[
 \begin{split}
 \det G_N
  &=\frac{\prod_{n=0}^N h_n}
          {D_t(P_k;s)\,|\det V_\beta|^2},\\
 \det G'_{N-v}
  &=\frac{\prod_{n=0}^{N-v}h_n}
     {D_{t+\Delta-v}(P_{k'};s)\,|\det V_{\beta'}|^2}.
 \end{split}\tag{CTV8}
\]
The relation phase $i^q$ or $i^{q'}$ multiplying each source
relation column has squared modulus one, as used in this equality.
The Vandermonde factors are nonzero by the original distinctness
of all roots. Their logarithmic difference is constant across
the four endpoints and has signed coefficient zero. Therefore
CTV8 proves the exact finite expression
\[
 \mathcal T_k^{(s)}
  =\sum_N\sigma_N\left[
      \sum_{n=N-v+1}^{N}\log h_n
      +\log D_{t+\Delta-v}(P_{k'};s)
      -\log D_t(P_k;s)\right].
 \tag{CTV9}
\]
The first sum is empty for $v=0$. In particular both low relation
determinants and both high relation determinants remain present.

\subsection{The exact quotient of the outer-strip polynomial}

Introduce the positive measure
\[
 d\zeta_{k',s}(y)=|P_{k'}(y)|^2\,dm_s(y).
 \tag{CTV10}
\]
This measure is the exact pullback of the original source metric
under the multiplication map
$\mathcal M_{P_{k'}}:f\mapsto P_{k'}f$.
Indeed $\|\mathcal M_{P_{k'}}f\|_{m_s}^2
=\int|f|^2d\zeta_{k',s}$, including the same full factor $M_s$.
On polynomials of degree at most $N-q'$ its image is the specified
subspace $P_{k'}\mathbb C[y]_{\le N-q'}$ of the original
degree-at-most-$N$ source. The map is injective because $P_{k'}$
is nonzero. Thus CTV10 records a proved metric map to the
original source, rather than an assigned new Gamma order.

For $L=N-q'=t+\Delta-1$, let
$Q_{\partial,L}^{(s)}$ be the minimum quotient Gram of
\[
 \mathbb C[y]_{\le L}\longrightarrow
       \mathbb C[y]/(P_{\partial,k})
 \tag{CTV11}
\]
in the source norm CTV10 and in the increasing monomial
remainder frame of dimension $\Delta$.
The source degree is at least $\Delta-1$ at every endpoint.
The same monic-frame Schur proof as in CTV7, now using CTV4,
gives
\[
 \det Q_{\partial,L}^{(s)}
     =\frac{D_{t+\Delta}(P_{k'};s)}{D_t(P_k;s)} .
 \tag{CTV12}
\]
In fact the ideal columns are $P_{\partial,k}y^j$, $0\le j<t$;
their Gram in CTV10 is exactly $D_t(P_k;s)$ by CTV4.
The numerator is the Gram of the complete monomial source
of dimension $L+1=t+\Delta$. This proves every source and
relation factor in CTV12.

Write $\lambda_j^{(k',s)}$ for the squared norm of the monic
orthogonal polynomial of degree $j$ in CTV10. This polynomial
and its norm are completely determined by the positive moment
matrices: subtract from $y^j$ its projection onto lower powers.
If $H_j$ is that lower Gram and $a_j$ its inner-product column
with $y^j$, the lower coefficients are $-H_j^{-1}a_j$.
Completing the square proves minimality among all monic
polynomials of degree $j$; its squared norm is the Schur
complement. Thus
$\lambda_j^{(k',s)}=D_{j+1}(P_{k'};s)/D_j(P_{k'};s)>0$.
Multiplying these identities and substituting CTV12 into CTV9
gives
\[
 \begin{gathered}
 \mathcal F_{\partial,k}^{(s)}
    =\log\det Q_{\partial,\Delta-1}^{(s)}
       +\log\det Q_{\partial,\Delta}^{(s)}
       -\log\det Q_{\partial,q+\Delta-1}^{(s)}
       -\log\det Q_{\partial,q+\Delta}^{(s)},\\
 \mathcal T_k^{(s)}
    =\mathcal F_{\partial,k}^{(s)}
                        +\sum_N\sigma_N\mathcal R_{N,v}^{(s)},\\
 \mathcal R_{N,v}^{(s)}
    =\sum_{a=0}^{v-1}
          \left(\log h_{N-a}
                -\log\lambda_{N-q'-a}^{(k',s)}\right).
 \end{gathered}\tag{CTV13}
\]
Every weighted norm index is nonnegative:
$N-q'-(v-1)\ge\Delta-v\ge16$.
The total difference therefore has an exact rank-$\Delta$
quotient presentation with precisely the displayed norm tails.

\subsection{Finite bounds for all norm tails in both original sources}

The original Gamma generating identity gives the full exponential
moment
$\int e^{t y}\,dm_s(y)=M_s(\cos t)^{-b}$ for
$|t|<\pi/2$. By the evenness of the measure this also equals
$\int\cosh(t y)\,dm_s(y)$.
The nonnegative Taylor terms of $\cosh$ imply
\[
 \begin{gathered}
 \int y^{2j}\,dm_s(y)
  \le M_s(\cos t_0)^{-b}\frac{(2j)!}{t_0^{2j}},\\
 \|y^j\|_{m_s}
  \le\sqrt{M_s}(\cos t_0)^{-b/2}\mathfrak c_{\rm mom}^j j!,\qquad
 t_0=\frac\pi4,\quad \mathfrak c_{\rm mom}=\frac2{t_0}=\frac8\pi .
 \end{gathered}\tag{CTV14}
\]
The second inequality uses
$(2j)!\le4^j(j!)^2$, proved by
$\binom{2j}{j}\le\sum_{a=0}^{2j}\binom{2j}{a}=4^j$.
Thus every constant in this moment estimate includes the
original source mass and order.

For a monic polynomial of degree $n$ whose roots have modulus
at most $R'$, its coefficient of $y^a$ has absolute value at most
$\binom na(R')^{n-a}$: it is the sum of those products of
$n-a$ distinct roots. The triangle inequality and CTV14 therefore
give the full bound
\[
 \|P\|_{m_s}
  \le\sqrt{M_s}(\cos t_0)^{-b/2}
          \mathfrak c_{\rm mom}^n n!\sum_{h=0}^n
                              \frac{(R'/\mathfrak c_{\rm mom})^h}{h!}
  \le\sqrt{M_s}(\cos t_0)^{-b/2}
          \mathfrak c_{\rm mom}^n n!\,e^{R'/\mathfrak c_{\rm mom}}.
 \tag{CTV15}
\]
The monic test polynomial $y^j$ in the minimization defining
$\lambda_j^{(k',s)}$ gives the polynomial
$P_{k'}(y)y^j$ in CTV15, of degree $n=q'+j$; its extra
$j$ roots are exactly zero and also satisfy the stated bound.
For the lower bound, CTV5 applies to every competing monic
polynomial, whose original Gamma squared norm is at least
$h_j$. Taking the minimum proves
\[
 \delta^{2q'}h_j
 \le\lambda_j^{(k',s)}
 \le M_s(\cos t_0)^{-b}\mathfrak c_{\rm mom}^{2(q'+j)}
                ((q'+j)!)^2e^{2R'/\mathfrak c_{\rm mom}}.
 \tag{CTV16}
\]
No lower bound on an actual period coefficient is used here.
The lower root polynomial itself supplies the strictly positive
real-line lower bound.

For each term of CTV13 put $n=N-a$ and $j=n-q'$.
Subtracting the logarithms in CTV16 from the exact
$\log h_n=\log M_s+\log n!+\log(b)_n$ yields
\[
 \begin{split}
 b\log(\cos t_0)+\log\frac{(b)_n}{n!}
          -2n\log\mathfrak c_{\rm mom}
                         -\frac{2R'}{\mathfrak c_{\rm mom}}
 &\le \log h_n-\log\lambda_j^{(k',s)}\\
 &\le
   \sum_{r=j+1}^{n}\bigl(\log r+\log(b+r-1)\bigr)
                          -2q'\log\delta .
 \end{split}\tag{CTV17}
\]
The two full $\log M_s$ contributions cancel in this displayed
comparison because both norms arise from exactly CTV1.
They have not been removed from either individual norm.

Let $c_b^{\log}=\max\{1,b-1\}$. As proved directly in CEP19,
$|\log(j/(b+j-1))|\le c_b^{\log}/j$ for $b\ge1/2$:
use $\log(1+x)\le x$ for $b\ge1$; for $1/2\le b<1$
use $-\log(1-a/j)\le a/(j-a)\le1/j$, $a=1-b$.
Summation gives
$|\log((b)_n/n!)|\le c_b^{\log}H_n$.
With all original endpoints $N\ge99$, define the explicit
nonnegative numbers
\[
 \begin{split}
 L_N^{\rm tail}
    &=c_b^{\log}H_N+\frac b2\log2
                        +2N\log(8/\pi)+\frac{\pi R'}4,\\
 U_N^{\rm tail}
    &=q'\bigl(\log N+\log(b+N-1)+2|\log\delta|\bigr),\\
 C_{N,v}^{\rm tail}
    &=v\max\{L_N^{\rm tail},U_N^{\rm tail}\}.
 \end{split}\tag{CTV18}
\]
Indeed $\cos(\pi/4)=2^{-1/2}$ and
$2R'/\mathfrak c_{\rm mom}=\pi R'/4$.
In the upper sum of CTV17 there are exactly $q'$ terms, and
each is at most $\log N+\log(b+N-1)$.
Thus CTV17 proves
\[
 -vL_N^{\rm tail}\le\mathcal R_{N,v}^{(s)}
               \le vU_N^{\rm tail},\qquad
 |\mathcal R_{N,v}^{(s)}|\le C_{N,v}^{\rm tail}.
 \tag{CTV19}
\]
When $v=0$ the equality $\mathcal R_{N,0}^{(s)}=0$ and the
zero bound follow from the empty sum.
These estimates apply for all actual endpoints and for
both unchanged source orders simultaneously.

\subsection{The proved first-variation reduction and its remaining matrix}

Since $N\le2q$, $q'\le q$, $b\le k/2$ and
$R'\le k\sqrt{\delta^2+\gamma^2}$, equations CTV18--19 give
$C_{N,v}^{\rm tail}=O_{\delta,\gamma}(q\log(q+1)+
k\log(q+1))$ uniformly for every integer $0\le v\le80$.
To see this directly, use $H_N\le1+\log N$ in the first
line of CTV18 and the displayed upper bounds in its second;
$v$ is bounded by the fixed integer 80.
Dividing by $kq$ tends to zero, since
$\log(q+1)/k\to0$ for $q=(k+1)^2$.
Consequently CTV13 and the finite bounds just proved yield
\[
 \left|\mathcal T_k^{(s)}-\mathcal F_{\partial,k}^{(s)}\right|
        \le\sum_N C_{N,v}^{\rm tail}=o(kq).
 \tag{CTV20}
\]
The estimate is uniform over the possible actual moment indices
$v$; it requires no bound on the remaining period coordinates.
Every polynomial and measure defining
$\mathcal F_{\partial,k}^{(s)}$ is independent of that index.

For an explicit matrix target, let
$H_{L}^{\zeta}$ be the moment Gram of
$1,y,\ldots,y^L$ in the full measure CTV10, and let
$\mathsf R_{\partial,L}$ be the literal remainder matrix
modulo the full polynomial $P_{\partial,k}$ in CTV4.
The exact metric is
\[
 Q_{\partial,L}^{(s)}
  =\bigl(\mathsf R_{\partial,L}
          (H_{L}^{\zeta})^{-1}
          \mathsf R_{\partial,L}^*\bigr)^{-1}.
 \tag{CTV21}
\]
The map is onto because its first $\Delta$ columns are the
identity in the remainder frame. The minimum source lift
of a coefficient vector $x$ is
$(H_L^\zeta)^{-1}\mathsf R_{\partial,L}^*
 Q_{\partial,L}^{(s)}x$; it maps to $x$ and is perpendicular
to the kernel. Adding any other lift adds a perpendicular
kernel vector, proving the stated Gram. This also proves
direct equality with the quotient used in CTV11--12.

Combining the complete exact CEP49 and CTV13 identities gives
\[
 \begin{gathered}
 F_K^{(s)}
  =\mathcal F_{\partial,k}^{(s)}
       -F_{\Xi,k}^{(s)}
       +\mathcal E_{A,k}^{(s)}
       +\sum_N\sigma_N\mathcal R_{N,v}^{(s)},\\
 \left|F_K^{(s)}-\mathcal F_{\partial,k}^{(s)}
                         +F_{\Xi,k}^{(s)}\right|
 \le2\sum_NE_N+\sum_NC_{N,v}^{\rm tail}=o(kq).
 \end{gathered}\tag{CTV22}
\]
Here the first error has its fixed-actual-period scope from CEP49;
the new tail bound has the stronger uniform-in-$v$ scope just
proved. Thus the remaining leading calculation now consists
of the explicit outer-strip quotient of rank $16k-48$ and the
original observation quotient of rank $8k-32$. Both have their
complete original metric maps and four cutoffs specified.
No value of either leading volume is assigned by this reduction,
and no derivative of the earlier $q^2$ asymptotic law was used.

% END CURRENT CTV

\clearpage
% BEGIN CURRENT ROQ
% ROQ: original residual observation quotient, with full coefficient maps.
\section{The actual observation remaining inside the conductor kernel}
\label{sec:ROQ}

Retain the original simple quartet, its actual unit and period matrix
$\mathsf H=F^{-1}\Pi U\mathcal E_{\rm prim}$, and its proved five-member
quadric orbit. Put $k=4l+1\ge9$, $q=(k+1)^2$, $q'=(k-7)^2$,
$\Delta=q-q'=16k-48$, $r=8k-16$, and $j=8k-32$.
The period parameter and logarithm branch are fixed at a point on this
proved locus. All constructions below use the original ordered primary
coordinates at
\[
 \beta_{ab}=k/2+(2a-k)\delta+i(2b-k)\gamma,
 \qquad \omega_{ab}=(2b-k)\gamma-i(2a-k)\delta,
 \quad 0\le a,b\le k.
 \tag{ROQ1}
\]
The full conductor is the actual biform
$A=\sum_{a,b=0}^8a_{ab}e_{ab}^{(8)}$ of OCF4. Its matrix
$C:\mathbb C^q\to\mathbb C^{q'}$ is exactly CEP4:
$C_{(a,b),(i,h)}=a_{i-a,h-b}$ when both differences belong to
$\{0,\ldots,8\}$, and zero otherwise. Thus $C^{\mathsf T}$
is multiplication by the entire $A$, and $V_A=\ker C$.
The original observation is $\Lambda$, and $K=\ker\Lambda\subset V_A$.
Ordinary transpose throughout the coefficient duality is denoted by
$\mathsf T$; a star denotes conjugate transpose in a displayed metric.

\subsection{An explicit residual map and its actual boundary basis}

Use the complete OCF6 elimination, with its actual largest nonzero
coefficient $a_{r_*s_*}$ and every remaining $a_{ab}$. To avoid confusing
its degree-$k$ matrices with the source cutoff, write them as
\[
 \mathsf N:\mathbb C^q\twoheadrightarrow\mathbb C^\Delta,
 \quad \mathsf S:\mathbb C^\Delta\hookrightarrow\mathbb C^q,
 \quad \mathsf V:\mathbb C^q\to\mathbb C^{q'},
 \qquad
 \mathsf S\mathsf N+C^{\mathsf T}\mathsf V=I_q,
 \quad \mathsf N\mathsf S=I_\Delta.
 \tag{ROQ2}
\]
These are the exact identities OCF7. Specifically, at each decreasing
lexicographic conductor position the elimination subtracts its current
coefficient divided by $a_{r_*s_*}$ times the entire translated biform
$A e_{ab}^{(k-8)}$. This specifies all entries of $\mathsf N,\mathsf V$;
$\mathsf S$ inserts the retained strip coordinates. It also proves
$\ker\mathsf N=\operatorname{im}C^{\mathsf T}$.
Taking the annihilator for the bilinear coefficient pairing therefore gives
\[
 I_V=\mathsf N^{\mathsf T}:\mathbb C^\Delta\xrightarrow{\ \sim\ }V_A,
 \qquad I_V^{-1}=\mathsf S^{\mathsf T}|_{V_A}.
 \tag{ROQ3}
\]
For a direct verification, $C\mathsf N^{\mathsf T}=0$ by transposing
$\mathsf NC^{\mathsf T}=0$, and $\mathsf S^{\mathsf T}\mathsf N^{\mathsf T}=I$.
If $Cx=0$, transposing the first identity of ROQ2 gives
$x=\mathsf N^{\mathsf T}\mathsf S^{\mathsf T}x$.

Retain the $j$ independent original invariant certificates $F_{k,h}(x)$
selected by OCF20, and their entire biform coefficient columns $f_h$:
\[
 \mathsf F=[f_1,\ldots,f_j],\qquad
 f_h=\operatorname{coeff}\bigl(F_{k,h}(\mathsf Ht)\bigr),\qquad
 Z=\mathsf N\mathsf F\in\operatorname{Mat}_{\Delta\times j}(\mathbb C),
 \qquad D=\mathsf F^{\mathsf T}.
 \tag{ROQ4}
\]
Here $t=(z_0w_0,z_0w_1,z_1w_0,z_1w_1)^{\mathsf T}$ and the coefficient
order is the full original biform order. Every certificate is supplied
by the fourteen-generator recursion OCF25--26, or equivalently the
complete twenty-five-generator recursion OCF9--20. At a recursion step
its certificate is the actual product $x^\nu F_{k-e,h}$, and its column
is $\mathsf N_k M_{f_\nu}\mathsf S_{k-e}Z_{k-e,h}$. The coefficients
of $f_\nu$ are the full multinomials OCF13 in the actual $\mathsf H$.
Thus ROQ4 specifies finite entries without choosing a substitute period
or introducing a freely chosen observation. The proved image and rank are
$\operatorname{im}Z=\mathsf N J_k$ and $\operatorname{rank}Z=j$.

Let $I$ be the $j$ retained pivot rows of $Z$, let $T$ be their
complement among the $\Delta$ strip rows, and put $Z_I=E_I^{\mathsf T}Z$,
$Z_T=E_T^{\mathsf T}Z$. The original OCF selection guarantees that
$Z_I$ is invertible on its exhaustive exact pivot chart. ROQ2 applied
to $\mathsf F$ gives the full coefficient identity
\[
 \mathsf F=\mathsf SZ+C^{\mathsf T}\mathsf V\mathsf F,
 \qquad
 D=Z^{\mathsf T}\mathsf S^{\mathsf T}
                      +(\mathsf V\mathsf F)^{\mathsf T}C.
 \tag{ROQ5}
\]
In particular the unreduced invariant rows and the reduced strip rows
agree on $V_A$, with precisely the displayed conductor multiple as their
ambient difference. The residual observation has the explicit forms
\[
 \Xi:V_A\longrightarrow\mathbb C^j,\quad x\longmapsto Dx,
 \qquad \Xi I_V=Z^{\mathsf T},\qquad
 L=\mathsf N^{\mathsf T}E_I(Z_I^{\mathsf T})^{-1},
 \quad CL=0,\quad DL=I_j.
 \tag{ROQ6}
\]
The last three equalities follow by multiplication and ROQ3--5.
Moreover $\ker\Xi=K$: the rows $C$ and $D$ together span the
coefficient space $J_k$. Indeed conductor multiplication spans
$\ker\mathsf N$, its image is contained in $J_k$ by OCF17, and
the $f_h$ project to the basis $Z$ of $\mathsf NJ_k$.
Thus a primary vector killed by both $C$ and $D$ is exactly the
annihilator of $J_k$, which is the original $K$ by OKM3--5.
This proves rank $j$, surjectivity, and the complete exact sequence
\[
 0\longrightarrow K\longrightarrow V_A
       \xrightarrow{\ \Xi\ }\mathbb C^j\longrightarrow0.
 \tag{ROQ7}
\]

For every $y\in\mathbb C^j$ define
$\mathfrak b(y)=\Lambda Ly$. These vectors give the literal isomorphism
$\mathfrak b:\mathbb C^j\to\Lambda(V_A)$. They span because for
$x\in V_A$, $x-LDx\in K$ by ROQ6--7. They are independent because
$\Lambda Ly=0$ implies $Ly\in K$, hence $y=DLy=0$.
The inverse on the displayed image is $\Lambda x\mapsto Dx$ for
$x\in V_A$: replacing $x$ by another such lift changes it by $K$,
which $D$ kills. In particular the boundary basis vectors are exactly
$\mathfrak b_h=\sum_{a,b}L_{(a,b),h}\Lambda\varepsilon_{ab}$.
Each $\Lambda\varepsilon_{ab}$ is the full original projected tensor
coefficient in OKA7, including its literal $1/5$ cyclic average.

The whole original observation can be recovered from these rows as well.
Take an actual right inverse $L_C$ of $C$, for example the inverse of
the first nonzero square column minor with its coordinate insertion,
and define
\[
 \mathsf L=\bigl[L_C-LDL_C,\ L\bigr],\qquad
 \mathsf R=\begin{pmatrix}C\\D\end{pmatrix},\qquad
 \mathsf R\mathsf L=I_{q'+j},\qquad
 \Lambda=(\Lambda\mathsf L)\mathsf R.
 \tag{ROQ8}
\]
The first identity follows by multiplying its four blocks. For the
last, $\mathsf R(x-\mathsf L\mathsf Rx)=0$, whose kernel is $K$
by ROQ7. Thus the $q'+j=q-r$ columns of $\Lambda\mathsf L$ form an
actual basis of the original boundary: applying its kernel and then
$\mathsf R\mathsf L=I$ proves independence, and ROQ8 proves spanning.
The last $j$ columns are exactly $\mathfrak b_h$.

For completeness the inherited period metric on this residual basis is
also explicit. Let $B$ be the complete matrix AKS8 in primary
coordinates, so $\jmath\Lambda=\mathcal F_0B$. If $\nu,\mu$ range over
the original weight-zero degree-$k$ monomials, set
\[
 (\mathcal G_0)_{\nu\mu}
 =\sum_{\substack{n_{tu}\ge0\\\sum_u n_{tu}=\nu_t\\
                                      \sum_t n_{tu}=\mu_u}}
       \frac{k!}{\prod_{t,u}n_{tu}!}
               \prod_{t,u}[5(\delta_{tu}+1)]^{n_{tu}},
 \qquad P_{\rm per}=L^*B^*\mathcal G_0BL>0.
 \tag{ROQ9}
\]
Expansion of the original word inner product gives
$\mathcal G_0=\mathcal F_0^*\mathcal F_0$ as in AKS10, so the displayed
matrix is exactly the Gram of the independent vectors
$\jmath\mathfrak b_h$. Its strict positivity follows from that
independence and injectivity of $\jmath$. All off-diagonal word terms
and every actual entry of $\mathsf H$ remain in ROQ9.

\subsection{All four induced Gamma metrics and their fixed determinant factor}

For either original source $s\in\{1,k\}$ put
$b_s=s/2$, $M_s=(2\pi)^{s/2}$ and retain
\[
 \begin{gathered}
 p_0=1,\quad p_1=y,\quad
 p_{n+1}=yp_n-n(b_s+n-1)p_{n-1},\\
 h_n=M_s n!(b_s)_n,
 \qquad (u_n)_{ab}=p_n(\omega_{ab})/\sqrt{h_n}.
 \end{gathered}
 \tag{ROQ10}
\]
Thus $U_N=[u_0,\ldots,u_N]$ is precisely CEP21 and
$R_N=U_NU_N^*>0$, $G_N=R_N^{-1}$ for every $N\ge q-1$.
The positivity follows from its first $q$ columns, whose determinant
is the original Vandermonde times their nonzero leading coefficients.
Define
\[
 H_{V,N}=\overline{\mathsf N}G_N\mathsf N^{\mathsf T},\qquad
 Q_{\Xi,N}=
       (Z^{\mathsf T}H_{V,N}^{-1}\overline Z)^{-1}>0.
 \tag{ROQ11}
\]
These are matrices of sizes $\Delta$ and $j$, respectively. The
minimum of $z^*H_{V,N}z$ under $Z^{\mathsf T}z=y$ is attained at
$z=H_{V,N}^{-1}\overline ZQ_{\Xi,N}y$. Its difference from any other
lift lies in $\ker Z^{\mathsf T}$ and is perpendicular in the
$H_{V,N}$ metric. Its squared norm is $y^*Q_{\Xi,N}y$.
Thus ROQ11 calculates exactly the induced metric of ROQ7 in the
actual boundary coordinates ROQ6. This proof also identifies it with
the quotient metric appearing in CEP27 and CEP49.

The actual kernel frame in the strip is
\[
 \rho=E_T^{\mathsf T}-Z_T Z_I^{-1}E_I^{\mathsf T},\qquad
 I_K=\mathsf N^{\mathsf T}\rho^{\mathsf T},\qquad
 F_{\rm str}=[\rho^{\mathsf T},\ E_I(Z_I^{\mathsf T})^{-1}].
 \tag{ROQ12}
\]
The first two matrices are precisely the OCF/OKM frame. In the ordered
rows $(T,I)$ the last one has blocks
$\left(\begin{smallmatrix}I_r&0\\
 -(Z_I^{\mathsf T})^{-1}Z_T^{\mathsf T}&(Z_I^{\mathsf T})^{-1}
 \end{smallmatrix}\right)$.
Consequently $|\det F_{\rm str}|^2=|\det Z_I|^{-2}$, including the
absolute value of the row-order permutation. Block elimination in
$F_{\rm str}^*H_{V,N}F_{\rm str}$ gives the exact identity
\[
 \det(I_K^*G_NI_K)\det Q_{\Xi,N}
                 =\frac{\det H_{V,N}}{|\det Z_I|^2}.
 \tag{ROQ13}
\]
Indeed its upper block is the kernel Gram and its Schur complement
is the minimum quotient metric of the previous paragraph, since
$Z^{\mathsf T}F_{\rm str}=[0,I_j]$. The determinant of a block
matrix is the determinant of its upper block times that Schur
complement, while congruence contributes $|\det F_{\rm str}|^2$.
This proves ROQ13 without assigning modulus one to the actual pivot.

Define $F_{\Xi,k}^{(s)}$ by the four ordered endpoints
$\mathcal N=(q-1,q,2q-1,2q)$ and signs $\sigma=(1,1,-1,-1)$:
\[
 F_{\Xi,k}^{(s)}=\sum_{N\in\mathcal N}\sigma_N\log\det Q_{\Xi,N}.
 \tag{ROQ14}
\]
The fixed factor $|\det Z_I|^{-2}$ in ROQ13 cancels with these signs.
If a different exact certificate basis is selected, its residual rows
on $V_A$ are $D'=MD$ for an invertible $j\times j$ matrix $M$.
This follows because both lists are bases of the dual of $V_A/K$.
For an explicit inverse formula, $M=D'L$, and the old and new right
sections give its inverse; composing the two surjections proves both
inverse identities. The new metric is
$Q'_{\Xi,N}=M^{-*}Q_{\Xi,N}M^{-1}$, directly from their identical
minimum problem. Its determinant changes by $|\det M|^{-2}$, which
also cancels in ROQ14. No pivot or source mass has been removed at
an individual endpoint. In fact ROQ10 gives $Q_{\Xi,N}=M_s\widehat Q_{\Xi,N}$
when the common mass is explicitly factored out of every column;
its determinant carries $M_s^j$. The metric relative to the retained
period Gram ROQ9 has determinant $\det Q_{\Xi,N}/\det P_{\rm per}$,
whose fixed denominator cancels only in the same four-endpoint sum.

\subsection{The exact constrained covariance in the original source}

The residual metric can be computed without inverting the strip frame.
Define, using the complete original covariance,
\[
 \begin{split}
 A_N^C&=CR_NC^*,\qquad B_N^D=DR_NC^*,\qquad T_N^D=DR_ND^*,\\
 W_N&=R_N-R_NC^*(A_N^C)^{-1}CR_N,\\
 \Omega_N&=DW_ND^*
          =T_N^D-B_N^D(A_N^C)^{-1}(B_N^D)^*.
 \end{split}\tag{ROQ15}
\]
Since $C$ has full row rank and $R_N>0$, $A_N^C>0$. There is the exact
identity
\[
 W_N=\mathsf N^{\mathsf T}H_{V,N}^{-1}\overline{\mathsf N},
 \qquad \Omega_N=Z^{\mathsf T}H_{V,N}^{-1}\overline Z>0,
 \qquad Q_{\Xi,N}=\Omega_N^{-1}.
 \tag{ROQ16}
\]
To prove its first equality, $CW_N=0$, $W_N$ is Hermitian, and for
$x\in\ker C$ direct multiplication gives $W_NG_Nx=x$.
The other displayed candidate is Hermitian, has image in $\ker C$
by ROQ3, and has the same property on $x=\mathsf N^{\mathsf T}z$.
Such a matrix is unique: for their Hermitian difference $H$, its image
lies in $V_A$ and $HG_Nx=0$ for $x\in V_A$. For arbitrary $y$ and
$x\in V_A$, this yields $x^*G_NHy=(HG_Nx)^*y=0$.
Putting $x=Hy\in V_A$ gives $(Hy)^*G_N(Hy)=0$, hence $Hy=0$.
This proves equality. Now use $D\mathsf N^{\mathsf T}=Z^{\mathsf T}$
from ROQ6 and take the adjoint to prove the other two identities.

Every entry in ROQ15 is the finite sum of original column products:
$A_N^C=\sum_{n=0}^N(Cu_n)(Cu_n)^*$,
$B_N^D=\sum_{n=0}^N(Du_n)(Cu_n)^*$, and
$T_N^D=\sum_{n=0}^N(Du_n)(Du_n)^*$.
The conductor column is exactly
$Cu_n=V_{N-v}A_Ne_n$ by CEP22 when $n\le N$; its vanishing low
columns and centre difference are precisely CEP13--16.
The residual column is explicitly
$(Du_n)_h=\sum_{a,b}(f_h)_{ab}p_n(\omega_{ab})/\sqrt{h_n}$.
In particular ROQ15 retains all conductor--observation cross terms.
Replacing $D$ by $D+TC$ leaves $\Omega_N$ exactly unchanged because
$CW_N=0=W_NC^*$. Applied to ROQ5, this proves that the full and
strip row presentations give exactly the same metric, rather than
assuming their ambient coefficient matrices coincide.

\subsection{Exact increments and a positive residual four-endpoint return}

For $N\ge q-1$ let $u=u_{N+1}$ be the next original column ROQ10 and
define
\[
 \begin{gathered}
 a_N=1+(Cu)^*(A_N^C)^{-1}(Cu)>0,\\
 w_N=u-R_NC^*(A_N^C)^{-1}Cu\in V_A,\qquad
 z_N=Dw_N=Du-B_N^D(A_N^C)^{-1}Cu,\\
 \eta_N=\frac{z_N^*Q_{\Xi,N}z_N}{a_N}\ge0.
 \end{gathered}\tag{ROQ17}
\]
The complete updates are
\[
 \begin{gathered}
 W_{N+1}=W_N+w_Nw_N^*/a_N,\qquad
 \Omega_{N+1}=\Omega_N+z_Nz_N^*/a_N,\\
 Q_{\Xi,N+1}=Q_{\Xi,N}
       -\frac{Q_{\Xi,N}z_Nz_N^*Q_{\Xi,N}}
                         {a_N+z_N^*Q_{\Xi,N}z_N},\\
 \log\det Q_{\Xi,N}-\log\det Q_{\Xi,N+1}=\log(1+\eta_N).
 \end{gathered}\tag{ROQ18}
\]
Here is a direct proof of the covariance update. Put $a=Cu$ and
$v=R_NC^*(A_N^C)^{-1}a$. Since $R_{N+1}=R_N+uu^*$,
$A_{N+1}^C=A_N^C+aa^*$, whose inverse is
$ (A_N^C)^{-1}-(A_N^C)^{-1}aa^*(A_N^C)^{-1}/a_N$;
multiplying by $A_N^C+aa^*$ verifies this inverse exactly.
For the full covariance calculation write
$\alpha=a^*(A_N^C)^{-1}a$.
The subtracted covariance product expands as
\[
 R_NC^*(A_N^C)^{-1}CR_N
       +vu^*+uv^*+\alpha uu^*
       -\frac{(v+\alpha u)(v+\alpha u)^*}{1+\alpha}.
\]
Subtracting this from $R_N+uu^*$ and the old $W_N$ leaves
$(u-v)(u-v)^*/(1+\alpha)$, proving the first update.
Multiplying by $D,D^*$ proves the second. The same inverse identity
proves the third, and the determinant identity
$\det(H+zz^*/a)=\det H(1+z^*H^{-1}z/a)$ proves the fourth.
For that identity conjugate by $H^{-1/2}$; the rank-one matrix
$H^{-1/2}zz^*H^{-1/2}/a$ has one eigenvalue $z^*H^{-1}z/a$
and all its others zero. This completes the update proof.

Telescoping with the four actual endpoints, rather than changing the
source degree, gives the exact finite sum
\[
 F_{\Xi,k}^{(s)}=\log(1+\eta_{q-1})
       +2\sum_{N=q}^{2q-2}\log(1+\eta_N)
       +\log(1+\eta_{2q-1})\ge0.
 \tag{ROQ19}
\]
Indeed $\log\det Q_{q-1}-\log\det Q_{2q-1}$ sums the differences
for $N=q-1,\ldots,2q-2$, while
$\log\det Q_q-\log\det Q_{2q}$ sums those for $N=q,\ldots,2q-1$.
This proves each endpoint weight. Its first increment
$\lambda_s=\log(1+\eta_{q-1})$ is a completely specified
period-dependent nonnegative lower bound for $F_{\Xi,k}^{(s)}$.
The formulas also characterize equality: $\lambda_s=0$ exactly when
$Du_q=B_{q-1}^D(A_{q-1}^C)^{-1}Cu_q$; the full return vanishes
exactly when this corresponding row residual is zero for every
column $u_q,\ldots,u_{2q}$. Positivity of $Q_{\Xi,N}$ proves both
directions, without asserting a nonzero generic residual.

\subsection{Finite scalar enclosures and propagation into CEP49}

Retain the exact AKS3 scalars $\mathcal Z_{s,t}$ and
$\mathcal L_{s,t}=\log\mathcal Z_{s,t}$, whose full original polynomial
covariance bound is proved in AKS20--27. In primary coordinates it says
\[
 \mathcal Z_{s,t}^{-1}G_{q-1}\preceq G_{q+t-1}\preceq G_{q-1},
                  \qquad 0\le t\le q+1.
 \tag{ROQ20}
\]
The unchanged coefficient map from AKS26 to primary coordinates is
the ordered $V_\beta$, so congruence by that invertible map proves
exactly the displayed statement. Restrict it to $V_A$ and take the
minimum on the fixed fibre $Dx=y$ in $V_A$. Inequalities of squared
norms hold for every such $x$, hence also for their attained minima.
Thus
$\mathcal Z_{s,t}^{-1}Q_{\Xi,q-1}\preceq Q_{\Xi,q+t-1}
\preceq Q_{\Xi,q-1}$.
If $d_t=\log\det Q_{\Xi,q-1}-\log\det Q_{\Xi,q+t-1}$,
the eigenvalues after congruence by $Q_{\Xi,q-1}^{-1/2}$ give
$0\le d_t\le j\mathcal L_{s,t}$, and ROQ18 gives their monotonicity.
Since $d_1=\lambda_s$, the exact four-endpoint identity
$F_{\Xi,k}^{(s)}=d_q+d_{q+1}-d_1$ proves
\[
 \lambda_s\le F_{\Xi,k}^{(s)}
       \le j(\mathcal L_{s,q}+\mathcal L_{s,q+1})-\lambda_s.
 \tag{ROQ21}
\]
The actual rank $j$ enters only through this proved eigenvalue bound;
the determinant and increment formulas ROQ11--19 retain its full
orientation. No leading volume is assigned from a dimension.

Let $\epsilon_{A,k}^{(s)}=2\sum_{N\in\mathcal N}E_N$ with the
complete original constants of CEP29--36. CEP49 gives
$F_K^{(s)}=\mathcal T_k^{(s)}-F_{\Xi,k}^{(s)}+
\mathcal E_{A,k}^{(s)}$ and
$|\mathcal E_{A,k}^{(s)}|\le\epsilon_{A,k}^{(s)}=o(kq)$
for a fixed actual period. Substitution of ROQ21 proves
\[
 \begin{split}
 \mathcal T_k^{(s)}-\epsilon_{A,k}^{(s)}
       -j(\mathcal L_{s,q}+\mathcal L_{s,q+1})+\lambda_s
       &\le F_K^{(s)},\\
 F_K^{(s)}&\le
          \mathcal T_k^{(s)}+\epsilon_{A,k}^{(s)}-\lambda_s.
 \end{split}\tag{ROQ22}
\]
Intersecting this interval with the full AKJ2 or AKJ4 intervals preserves
each of their finite constants and source correlations. More precisely,
one may use the exact value $\mathcal E_{A,k}^{(s)}$ in place of its
upper or lower envelope in the appropriate side. Formula ROQ19 supplies
every additional original positive subtraction, rather than only the
first one. These are finite results on the residual quotient itself;
they do not claim that its leading $kq$ coefficient has been evaluated.

\subsection{The complete residual quotient in the conductor graph}

At a fixed original endpoint $N$, let $\mathcal P_N$, $I_N$,
$\mathcal W_N=U_N^{-1}(V_A)$ and the actual source isomorphism
$S_Np=(P_{<v}p,\mathcal T_{A,N}p)$ be precisely CEP28--31.
The source carries its original orthonormal Gamma coordinates of
order $s=1$ or $s=k$. The induced linear isomorphism on the residual
observation quotient is
\[
 \mathcal W_N/U_N^{-1}(K)
       \xrightarrow{\ [p]\mapsto[S_Np]\ }
 S_N\mathcal W_N/S_NU_N^{-1}(K),
 \qquad [p]\longmapsto DU_Np\in\mathbb C^j.
 \tag{ROQ23}
\]
The last map identifies the left quotient with $V_A/K$ and its fixed
coordinates ROQ6. Its kernel and surjectivity follow from ROQ7 and
surjectivity of $U_N$; applying the invertible $S_N$ proves the other
isomorphism. Thus all relation directions, including the source ideal
and the entire actual $K$, are retained in the transported denominator.

Here are exact matrices of that denominator and the resulting metric.
Let $Z_N^{\chi}$ be the original source coefficient columns
$\chi_k S^a$, $0\le a\le N-q$, with no columns when $N=q-1$.
Take the actual primary frames $I_K$ of ROQ12 and $L$ of ROQ6 and set
\[
 Y_{K,N}=U_N^*G_NI_K,\qquad Y_{L,N}=U_N^*G_NL,
 \qquad J_N=[Z_N^{\chi},Y_{K,N}].
 \tag{ROQ24}
\]
The columns of $Y_{K,N}$ are the minimum lifts of $I_K$ by the same
orthogonality proof as OKM8. In particular they are perpendicular to
$I_N=\ker U_N$, so $J_N$ has independent columns and its image is
exactly $U_N^{-1}(K)$. It has $N+1-q+r$ columns, while
$\dim\mathcal W_N=N+1-q+\Delta$, leaving precisely $j$ quotient
dimensions. The transported complete relation projector and Gram are
\[
 \begin{gathered}
 P_N^{\Xi,\rm graph}
   =I-S_NJ_N(J_N^*S_N^*S_NJ_N)^{-1}J_N^*S_N^*,\\
 \widetilde Q_{\Xi,N}
       =Y_{L,N}^*S_N^*P_N^{\Xi,\rm graph}S_NY_{L,N}>0.
 \end{gathered}\tag{ROQ25}
\]
The inverse exists because $S_N$ is invertible and $J_N$ has independent
columns. Multiplication verifies that the projector is Hermitian,
idempotent and projects onto the orthogonal complement of the entire
$S_NU_N^{-1}(K)$. Thus its Gram is the minimum norm of the transported
quotient ROQ23. Replacing $Y_{L,N}$ by another source lift of $L$ adds
columns in $I_N$, which this projector kills. Replacing the section
$L$ by another section of ROQ6 adds columns in $K$; their lifts lie
in the span of $J_N$ and are killed. Therefore ROQ25 computes
the specified quotient metric independently of either choice.

Before applying $S_N$, the same projection gives
\[
 Q_{\Xi,N}
 =L^*G_NL-L^*G_NI_K(I_K^*G_NI_K)^{-1}I_K^*G_NL.
 \tag{ROQ26}
\]
To prove it, $Y_{L,N}$ and $Y_{K,N}$ are perpendicular to the ideal,
so projection off $J_N$ only subtracts
$Y_{K,N}(Y_{K,N}^*Y_{K,N})^{-1}Y_{K,N}^*Y_{L,N}$.
Using $U_NU_N^*=R_N$ in their Gram products yields exactly ROQ26.
It is also the minimum of the original $G_N$ norm on $Ly+K$, so
it agrees with the already proved ROQ11 and ROQ16.

The all-rank exterior estimate CEP33--35 applies to ROQ23: represent
a quotient class by its minimum lift in the denominator's orthogonal
complement, apply $S_N$, then project off the target denominator.
This projection has every exterior norm at most one. Consequently the
induced $j$th exterior norm is at most $\|\bigwedge^jS_N\|$, and its
inverse is at most $\|\bigwedge^jS_N^{-1}\|$ by the identical argument
for $S_N^{-1}$. CEP34 bounds both logarithms from above by $E_N$.
Applying them to the fixed frame ROQ6 and squaring exterior volumes gives
\[
 \left|\log\det\widetilde Q_{\Xi,N}
                -\log\det Q_{\Xi,N}\right|\le2E_N,\qquad
 \left|\sum_{N\in\mathcal N}\sigma_N\log\det\widetilde Q_{\Xi,N}
                       -F_{\Xi,k}^{(s)}\right|
                       \le\epsilon_{A,k}^{(s)}=o(kq).
 \tag{ROQ27}
\]
All cross terms in $J_N^*S_N^*S_NJ_N$ remain present in ROQ25,
including $P_{<v}(\chi_k S^a)$ and its pairings with the minimum
kernel lifts. Its ideal columns can equally be generated from the full
products $d_{rs}$ and maps $\mathcal U_{A,L}$ in CEP41--45, because
those formulas prove exactly $S_NZ_N^{\chi}$ column by column.
This supplies the typed relation between the explicit residual
observation and the conductor graph at all four original cutoffs.

% END CURRENT ROQ

\clearpage
% BEGIN CURRENT OSP
% Full original-root comparison; no sealed provider is modified.
\section{A quantitative power comparison for the actual outer-strip quotient}
\label{sec:OSP}

Retain the original real parameters and integers
\[
\begin{gathered}
0<\delta<\tfrac12,\quad \gamma>2,\quad k=4l+1\ge9,\\
q=(k+1)^2,\quad k'=k-8,\quad q'=(k-7)^2,\quad
\Delta=q-q'=16k-48.
\end{gathered}\tag{OSP1}
\]
The two source orders in every formula below are precisely $s=1,k$.
Set $b=s/2$, $M_s=(2\pi)^{s/2}$, and retain the full original measure
\[
dm_s(y)=M_s\frac{2^b}{4\pi\Gamma(b)}
       |\Gamma(b/2+iy/2)|^2\,dy.
\tag{OSP2}
\]
In the original coordinate $S=k/2+iy$, and respectively
$S'=k/2-4+iy$, the complete monic root polynomials are
\[
\begin{split}
P_k(y)&=i^{-q}\chi_k(k/2+iy)
 =\prod_{a,b_0=0}^{k}
 [y-(2b_0-k)\gamma+i(2a-k)\delta],\\
P_{k'}(y)&=i^{-q'}\chi_{k'}(k/2-4+iy)
 =\prod_{a,b_0=0}^{k'}
 [y-(2b_0-k')\gamma+i(2a-k')\delta].
\end{split}\tag{OSP3}
\]
The index $b_0$ in these products is distinct from the source parameter
$b=s/2$. Each root occurs with exactly its displayed multiplicity.
The map $(a,b_0)\mapsto(a+4,b_0+4)$ identifies the second root list
with the central subrectangle of the first, so
\[
P_k=P_{k'}P_{\partial,k},\qquad
P_{\partial,k}(y)=
\prod_{(a,b_0)\in\{0,\ldots,k\}^2\setminus\{4,\ldots,k-4\}^2}
[y-(2b_0-k)\gamma+i(2a-k)\delta].
\tag{OSP4}
\]
This proves the factorization and its monic constant; the outer degree
is $\Delta$. Both complete root lists are invariant under
$\omega\mapsto-\omega$. Every root has modulus at most
\[
R=k\sqrt{\delta^2+\gamma^2}.
\tag{OSP5}
\]

\subsection{The two original source masses and a proved interval lower bound}

We first record the source facts needed in the comparison, including
their constants and proofs. For $a>0$ and real $x$, the two Euler
integrals and the change of variables in the beta integral give
\[
\frac{\Gamma(a+ix)\Gamma(a-ix)}{\Gamma(2a)}
 =\int_{\mathbb R}e^{ixu}(2\cosh(u/2))^{-2a}\,du.
\tag{OSP6}
\]
For example the intermediate beta integral is
$\int_0^\infty t^{a+ix-1}(1+t)^{-2a}dt$, and $t=e^u$
gives the displayed expression. Absolute convergence follows from
$a>0$. Rotating Euler's integral through
$\operatorname{sgn}(y)\theta$, $0<\theta<\pi/2$, gives
\[
|\Gamma(a+iy/2)|
 \le e^{-\theta|y|/2}\Gamma(a)(\cos\theta)^{-a}.
\tag{OSP7}
\]
The rotation follows from Cauchy's theorem on the intervening
sector: the large circular arc tends to zero because its real
exponential has positive cosine, and the small one tends to zero
because $a>0$. The remaining radial integral has absolute value
at most $\Gamma(a)(\cos\theta)^{-a}$. In particular the Fourier
transform in OSP6 is integrable. Fourier inversion, followed by
$y=2x$, now gives
\[
\int_{\mathbb R}e^{i\xi y}\,dm_s(y)
  =M_s(\cosh\xi)^{-b},\qquad
\int_{\mathbb R}dm_s=M_s.
\tag{OSP8}
\]
The same identity holds analytically for $|\operatorname{Im}\xi|<\pi/2$:
OSP7 with any larger $\theta<\pi/2$ supplies an integrable dominating
function on each smaller strip. Substituting $\xi=-it$ proves
\[
\int_{\mathbb R}e^{ty}\,dm_s(y)=M_s(\cos t)^{-b},
\qquad |t|<\pi/2.
\tag{OSP9}
\]
This also proves finiteness of all moments used below.

Write
\[
d\sigma(y)=\frac{|\Gamma(1/4+iy/2)|^2}{2\pi}\,dy,\quad
C_L=\frac{\pi}{\Gamma(3/4)^2},\quad
\ell_s=\frac{s-1}{4},\quad
\beta_s=\frac{M_s}{\sqrt2\,\Gamma(s/2)}.
\tag{OSP10}
\]
The two values of $\ell_s$ are the integers $0,l$. The reflection
identity gives a useful explicit lower bound
\[
\frac{d\sigma}{dy}(y)\ge C_Le^{-\pi|y|}.
\tag{OSP11}
\]
Here is its proof on the required domain. For $0<\operatorname{Re}z<1$
the Euler integrals give
$\Gamma(z)\Gamma(1-z)=\int_0^\infty t^{z-1}/(1+t)\,dt$.
On a keyhole contour with argument in $(0,2\pi)$, the circular
terms vanish with bounds $O(r^{\operatorname{Re}z})$ at zero
and $O(R^{\operatorname{Re}z-1})$ at infinity. The residue at $-1$
is $e^{i\pi(z-1)}$; the two radial integrals therefore give
$(1-e^{2\pi iz})\int_0^\infty t^{z-1}/(1+t)\,dt
=2\pi i e^{i\pi(z-1)}$. Their ratio is $\pi/\sin(\pi z)$.
At $z=1/4+iy/2$ this proves
\[
|\Gamma(1/4+iy/2)|^2|\Gamma(3/4-iy/2)|^2
 =\frac{2\pi^2}{\cosh(\pi y)}.
\]
The unrotated Euler integral bounds the second squared factor by
$\Gamma(3/4)^2$, proving
$d\sigma/dy\ge C_L/\cosh(\pi y)\ge C_Le^{-\pi|y|}$.

Gamma recurrence, proved by one integration by parts in Euler's
integral, gives the exact original-source identity
\[
dm_s(y)=\beta_s
 \prod_{j=0}^{\ell_s-1}\bigl(y^2+(2j+\tfrac12)^2\bigr)\,d\sigma(y).
\tag{OSP12}
\]
Indeed $b/2=\ell_s+1/4$, and each recurrence factor before
squaring and clearing denominators is
$j+1/4+iy/2$; its squared modulus is
$[y^2+(2j+1/2)^2]/4$. The factor $4^{-\ell_s}$ together with
the prefactor in OSP2 is exactly $\beta_s/(2\pi)$, because
$b=2\ell_s+1/2$. This proves every factor in OSP12.
For $s=1$ its product is empty and $\beta_1=1$.
In particular, on the full interval $q\le y\le2q$,
\[
\frac{dm_s}{dy}(y)
 \ge \beta_s q^{2\ell_s}C_Le^{-2\pi q}.
\tag{OSP13}
\]
The source mass-to-lower-density factor obeys the uniform finite bound
\[
\frac{M_s}{\beta_s q^{2\ell_s}}
 =\sqrt2\,\Gamma(b)q^{-2\ell_s}
 \le\sqrt{2\pi}\qquad(s=1,k).
\tag{OSP14}
\]
To prove the inequality, use
$\Gamma(b)=\sqrt\pi\prod_{j=0}^{2\ell_s-1}(j+1/2)
\le\sqrt\pi\,b^{2\ell_s}$ and $b\le q$.
All source masses have thus been retained before taking this exact
ratio.

\subsection{A complete coefficient-space estimate for the inner interval}

Let $p\in\mathbb C[y]$ have degree at most $d$, and let
$0<T\le q/2$. Then
\[
\sup_{|y|\le T}|p(y)|^2
 \le \frac{(d+1)^2\,10^{2d}}q
                  \int_q^{2q}|p(u)|^2\,du.
\tag{OSP15}
\]
This controls every coefficient direction; no assumption on zeros,
leading coefficient, or extremizers of $p$ is made.

For proof define the Legendre polynomial
$L_j(x)=(2^j j!)^{-1}(d/dx)^j(x^2-1)^j$.
Its leading coefficient is $(2j)!/[2^j(j!)^2]$. Integrating by
parts $j$ times, with all lower boundary derivatives zero, proves
orthogonality to lower powers and
\[
\int_{-1}^1L_j(x)^2\,dx
=\frac{(2j)!}{2^{2j}(j!)^2}\int_{-1}^1(1-x^2)^j\,dx
=\frac2{2j+1}.
\]
For the middle integral set $x=2t-1$; repeated integration by parts
gives $\int_0^1t^j(1-t)^jdt=(j!)^2/(2j+1)!$.
Consequently $\sqrt{2j+1}L_j(2z-3)$, $0\le j\le d$,
form an orthonormal basis for degree-at-most-$d$ polynomials on
$[1,2]$. Leibniz's rule also gives the explicit identity
\[
L_j(x)=2^{-j}\sum_{h=0}^j\binom jh^2
                         (x-1)^{j-h}(x+1)^h.
\]
For $|x|\le4$ its modulus is at most
$(5/2)^j\sum_h\binom jh^2\le(5/2)^j4^j=10^j$;
the sum inequality follows by comparing with
$(\sum_h\binom jh)^2=4^j$.
Expanding $p(qz)$ in the displayed orthonormal basis and applying
Cauchy--Schwarz gives, for $|z|\le T/q\le1/2$,
\[
|p(qz)|^2
\le\sum_{j=0}^d(2j+1)10^{2j}
                  \int_1^2|p(qv)|^2\,dv
\le(d+1)^2\,10^{2d}\int_1^2|p(qv)|^2\,dv .
\]
Here $|2z-3|\le4$ and $\sum_{j=0}^d(2j+1)=(d+1)^2$.
Changing variables proves OSP15, including $d=0$.

\subsection{The exact original-root map and a finite full-norm comparison}

In this subsection $P$ is either of the complete polynomials
$P_k,P_{k'}$ in OSP3, and $Q=\deg P$ is respectively $q,q'$.
For $d\ge0$ retain the two injections into the same original
polynomial Hilbert space
\[
\begin{aligned}
\mathcal J_P:\mathbb C[y]_{\le d}&\longrightarrow
          \mathbb C[y]_{\le Q+d},&p&\longmapsto Pp,\\
\mathcal J_{0,Q}:\mathbb C[y]_{\le d}&\longrightarrow
          \mathbb C[y]_{\le Q+d},&p&\longmapsto y^Qp .
\end{aligned}\tag{OSP16}
\]
Their codomain norm is precisely $m_s$ in OSP2.
If $P(y)=\sum_{a=0}^QP_a y^a$, set $P_a=0$ outside this range.
In increasing coefficient frames their matrices are
\[
(\mathsf J_P)_{rj}=P_{r-j},\qquad
(\mathsf J_{0,Q})_{rj}=\mathbf1_{\{r=Q+j\}},
\quad 0\le r\le Q+d,\quad 0\le j\le d.
\tag{OSP17}
\]
Thus the exact isomorphism between their images is
$y^Qp\mapsto Pp$, with inverse $Pp\mapsto y^Qp$.
It preserves the original coefficient vector $p$. The first
matrix contains every actual root coefficient; none is replaced
inside that map. Returning to $S$ or $S'$ multiplies the relevant
relation columns by the actual phase $i^Q$ and applies the
triangular coordinate pullback with diagonal entries $i^j$.
Each has the stated exact inverse and determinant of modulus one.

Suppose
\[
R<T\le q/2.
\tag{OSP18}
\]
Define the finite positive constants
\[
\begin{split}
\eta_{Q,T}&=\frac{Q R^2}{T^2-R^2},\\
\kappa_{Q,d,T,s}
 &=\frac{M_s}{\beta_s q^{2\ell_s}C_L}
    \frac{(d+1)^2}{q}\,
    e^{2\pi q}\,10^{2d}
             \left(\frac{T^2+R^2}{q^2}\right)^Q,\\
e_{Q,d,T,s}&=\eta_{Q,T}+\log(1+\kappa_{Q,d,T,s}).
\end{split}\tag{OSP19}
\]
Then every $p\in\mathbb C[y]_{\le d}$ satisfies
\[
e^{-e_{Q,d,T,s}}\int_{\mathbb R}|y|^{2Q}|p(y)|^2\,dm_s(y)
\le\int_{\mathbb R}|P(y)p(y)|^2\,dm_s(y)
\le e^{e_{Q,d,T,s}}
      \int_{\mathbb R}|y|^{2Q}|p(y)|^2\,dm_s(y).
\tag{OSP20}
\]

Here is a proof retaining the entire inner region. Pair the actual
roots as $\omega,-\omega$. On $|y|\ge T$, each paired quotient is
$1-\omega^2/y^2$, so
\[
\left|\log\frac{|P(y)|^2}{|y|^{2Q}}\right|
\le -Q\log(1-R^2/T^2)
\le\frac{Q R^2}{T^2-R^2}=\eta_{Q,T}.
\tag{OSP21}
\]
The first inequality uses $1-u\le|1-z|\le1+u$ for $|z|\le u<1$
on each of $Q/2$ pairs. The second follows by integrating
$1/(1-u)\le1/(1-u_0)$ on $0\le u\le u_0=R^2/T^2$.
For every real $y$, the same pairing gives
$|P(y)|^2\le(y^2+R^2)^Q$. Write
\[
A_B=\int_q^{2q}|y|^{2Q}|p(y)|^2\,dm_s(y).
\]
For nonzero $p$ this number is positive, and OSP13 gives
$A_B\ge q^{2Q}\beta_s q^{2\ell_s}C_Le^{-2\pi q}
 \int_q^{2q}|p(y)|^2dy$.
OSP15 and the exact full mass OSP8 imply
\[
\begin{split}
\int_{|y|<T}|P(y)p(y)|^2\,dm_s(y)
 &\le M_s(T^2+R^2)^Q\sup_{|y|\le T}|p(y)|^2
 \le\kappa_{Q,d,T,s}A_B,\\
\int_{|y|<T}|y|^{2Q}|p(y)|^2\,dm_s(y)
 &\le M_sT^{2Q}\sup_{|y|\le T}|p(y)|^2
 \le\kappa_{Q,d,T,s}A_B .
\end{split}\tag{OSP22}
\]
Let $A_{\rm out}$ and $B_{\rm out}$ be the integrals outside
$(-T,T)$ with factors $|y|^{2Q}$ and $|P|^2$, respectively.
Since $[q,2q]$ is outside that interval, $A_B\le A_{\rm out}$.
OSP21 gives
$e^{-\eta}A_{\rm out}\le B_{\rm out}\le e^\eta A_{\rm out}$.
OSP22 bounds both inner contributions by $\kappa A_{\rm out}$.
Therefore the full reference norm is at most
$(1+\kappa)A_{\rm out}$ and at least $A_{\rm out}$, while the
full actual norm is at least $e^{-\eta}A_{\rm out}$ and at most
$(e^\eta+\kappa)A_{\rm out}\le e^\eta(1+\kappa)A_{\rm out}$.
These four inequalities prove OSP20. The zero polynomial satisfies
it directly.

\subsection{Uniform bounded logarithmic cost at every required dimension}

Choose the fixed number and cutoff
\[
\epsilon=2^{-32},\qquad T=\epsilon q,
\]
and impose the explicit finite conditions
\[
k\ge29,\qquad
k\sqrt{\delta^2+\gamma^2}\le\frac{\epsilon q}{2}.
\tag{OSP23}
\]
In particular $T\ge2R$, as required for a uniform margin in OSP21.
These conditions hold eventually for every fixed original
$\delta,\gamma$; the sufficient bound
$k\ge\max\{29,2\sqrt{\delta^2+\gamma^2}/\epsilon\}$
follows from $q\ge k^2$. The first condition implies $q'\ge q/2$:
$2(k-7)^2-(k+1)^2=k^2-30k+97$ is positive at $k=29$
and is increasing thereafter.

For $Q\in\{q',q\}$ and every $0\le d\le2q$, set
\[
\begin{split}
K_q&=\frac{\sqrt{2\pi}}{C_L}\frac{(2q+1)^2}{q}e^{-2q},\\
\mathcal E_k&=\frac{qR^2}{\epsilon^2q^2-R^2}+\log(1+K_q).
\end{split}\tag{OSP24}
\]
Then OSP19 satisfies $e_{Q,d,T,s}\le\mathcal E_k$ for both
original source orders. To verify this uniformly, OSP14,
$T^2+R^2\le(5/4)\epsilon^2q^2$, $Q\ge q/2$ and $d\le2q$
give
\[
\kappa_{Q,d,T,s}
\le\frac{\sqrt{2\pi}}{C_L}\frac{(2q+1)^2}{q}
\exp\left(q\{2\pi+4\log10+\tfrac12\log(5/4)+\log\epsilon\}\right).
\tag{OSP25}
\]
The expression in braces is less than $-2$. Indeed
$\pi<4$, $\log10<4\log2$, $\log(5/4)<1/4$, and
$\log2>2/3$ give an upper bound
$8-16\log2+1/8<-61/24<-2$.
For completeness, $\pi<4$ follows from
$\pi/4=\int_0^1(1+x^2)^{-1}dx<1$; and the substitution
$x=(1+t)/(1-t)$ in $\log2=\int_1^2dx/x$ gives
$\log2=2\int_0^{1/3}(1-t^2)^{-1}dt>2/3$.
The logarithm bound follows by integrating $1/(1+x)<1$.
This proves OSP25 with $K_q$ as stated. Also
\[
0\le\mathcal E_k
\le \frac{4(\delta^2+\gamma^2)}{3\epsilon^2}+\log(1+K_q).
\tag{OSP26}
\]
Here $R^2\le\epsilon^2q^2/4$ and $R^2/q\le\delta^2+\gamma^2$.
Thus the logarithmic norm cost is bounded at fixed original
parameters, uniformly over all these coefficient degrees and both
unchanged source orders.

\subsection{Every relation determinant and the exact outer quotient}

For a nonzero polynomial $P$ and $a\ge0$ define in the unchanged monomial frame
\[
\begin{split}
D_a(P;s)&=\det\left[\int_{\mathbb R}y^{i+j}|P(y)|^2\,dm_s(y)
                                      \right]_{0\le i,j<a},\\
\mathcal D_a(Q;s)&=\det\left[\int_{\mathbb R}y^{2Q+i+j}\,dm_s(y)
                                      \right]_{0\le i,j<a},
\qquad D_0=\mathcal D_0=1 .
\end{split}\tag{OSP27}
\]
All nonempty Grams are positive definite, because a nonzero
polynomial has positive squared modulus on an open subinterval
of the strictly positive original measure. In the notation of
OSP17 their exact matrices are
$\mathsf J_P^*H^{(s)}_{Q+a-1}\mathsf J_P$ and
$\mathsf J_{0,Q}^*H^{(s)}_{Q+a-1}\mathsf J_{0,Q}$,
where $H^{(s)}_{Q+a-1}$ is the complete original moment Gram.
Consequently OSP20--26 prove for $1\le a\le2q+1$
\[
\left|\log D_a(P;s)-\log\mathcal D_a(Q;s)\right|
\le a\mathcal E_k,\qquad P=P_k\ \hbox{or}\ P_{k'} .
\tag{OSP28}
\]
To see the determinant step exactly, conjugate the actual Gram by
the inverse positive square root of the reference Gram. OSP20
places every eigenvalue of the resulting positive Hermitian
matrix in $[e^{-\mathcal E_k},e^{\mathcal E_k}]$. The logarithm
of its determinant is the sum of its $a$ eigenvalue logarithms,
proving OSP28. This is the full exterior determinant estimate,
including every polynomial coefficient direction. The estimate
also holds at $a=0$ with both sides zero.

For $t\in\{0,1,q,q+1\}$, let $L=t+\Delta-1$. The actual quotient
has source $\mathbb C[y]_{\le L}$ in measure
$|P_{k'}(y)|^2dm_s(y)$, target the increasing remainder frame of
$\mathbb C[y]/(P_{\partial,k})$, and its minimum Gram
$Q_{\partial,L}^{(s)}$. Define a comparison quotient with source
$\mathbb C[y]_{\le L}$ in measure $|y|^{2q'}dm_s(y)$, target the
increasing remainder frame of $\mathbb C[y]/(y^\Delta)$, and
minimum Gram $Q_{0,L}^{(s)}$. Its source measure still has order
$s$ and the full mass factor in OSP2.

For precision, the exact map relating the coefficient presentations
is the degree-preserving isomorphism
\[
\mathcal U_L(y^j)=
\begin{cases}
y^j,&0\le j<\Delta,\\
P_{\partial,k}(y)y^{j-\Delta},&\Delta\le j\le L.
\end{cases}\tag{OSP29}
\]
Its coefficient matrix has diagonal entries one and no terms of
degree greater than that of its input column. It is therefore
invertible with determinant one. If $\mathsf R_{\partial,L}$
is division remainder modulo $P_{\partial,k}$ and
$\mathsf R_{0,L}$ is division remainder modulo $y^\Delta$,
then
$\mathsf R_{\partial,L}\mathcal U_L=\mathsf R_{0,L}$:
the first $\Delta$ basis vectors map to their stated remainders,
and every later vector maps into the actual ideal. This proves
the precise morphism between the quotient coefficient spaces;
no equality of their metrics is assumed.

Using the original monic relation frames on the two sources gives
the exact identities
\[
\det Q_{\partial,t+\Delta-1}^{(s)}
 =\frac{D_{t+\Delta}(P_{k'};s)}{D_t(P_k;s)},\qquad
\det Q_{0,t+\Delta-1}^{(s)}
 =\frac{\mathcal D_{t+\Delta}(q';s)}{\mathcal D_t(q;s)} .
\tag{OSP30}
\]
For proof, append to $1,y,\ldots,y^{\Delta-1}$ the ideal frame
$P_{\partial,k},yP_{\partial,k},\ldots,y^{t-1}P_{\partial,k}$.
Its source coefficient determinant is one by OSP29. The relation
Gram is $D_t(P_k;s)$ because $P_{k'}P_{\partial,k}=P_k$.
Completing the square in the ideal coordinates gives the minimum
quotient Gram as the Schur complement; the determinant of the
full source Gram is the product of that determinant and the
relation determinant. This proves the first identity.
For the second repeat the same calculation with $y^\Delta$:
$y^{q'}y^\Delta=y^q$, with the original source order $s$.
The same proof includes $t=0$ using the empty relation block.

OSP28 applies to every index here, since
$t+\Delta\le q+\Delta+1\le2q+1$. Define the exact endpoint
difference $\varepsilon_t^{(s)}$ by
\[
\log\det Q_{\partial,t+\Delta-1}^{(s)}
=\log\mathcal D_{t+\Delta}(q';s)-\log\mathcal D_t(q;s)
 +\varepsilon_t^{(s)} .
\]
Then, separately at each of the four endpoints,
\[
|\varepsilon_t^{(s)}|\le(2t+\Delta)\mathcal E_k.
\tag{OSP31}
\]
In particular set
\[
\begin{split}
\mathcal F_{0,k}^{(s)}
={}&\log\mathcal D_\Delta(q';s)
    +\log\mathcal D_{\Delta+1}(q';s)
    -\log\mathcal D_{q+\Delta}(q';s)
    -\log\mathcal D_{q+\Delta+1}(q';s)\\
 &-\log\mathcal D_1(q;s)
    +\log\mathcal D_q(q;s)+\log\mathcal D_{q+1}(q;s).
\end{split}\tag{OSP32}
\]
The actual four-endpoint quantity of CTV13 therefore obeys the
proved finite estimate
\[
\boxed{\displaystyle
\left|\mathcal F_{\partial,k}^{(s)}-\mathcal F_{0,k}^{(s)}\right|
\le4(q+\Delta+1)\mathcal E_k
=O_{\delta,\gamma}(q)=o(kq),\qquad s=1,k .}
\tag{OSP33}
\]
Indeed its exact error is
$\varepsilon_0+\varepsilon_1-\varepsilon_q-\varepsilon_{q+1}$;
summing OSP31 gives
$\Delta+(\Delta+2)+(2q+\Delta)+(2q+\Delta+2)
=4(q+\Delta+1)$. OSP26 and $\Delta<q$ prove the last two
estimates. This calculation uses all four source degrees
$\Delta-1,\Delta,q+\Delta-1,q+\Delta$ and every source and
relation determinant in OSP30.

\subsection{The resulting same-source moment target}

The moments in OSP27 are completely specified by the original
analytic identity OSP9:
\[
\mu_j^{(s)}=\int y^j\,dm_s(y)
=M_s\left.\frac{d^j}{dz^j}(\cos z)^{-s/2}\right|_{z=0},
\qquad
\mathcal D_a(Q;s)=
\det[\mu_{2Q+i+j}^{(s)}]_{0\le i,j<a}.
\tag{OSP34}
\]
The differentiated integral is justified by OSP7, or directly by
the finite exponential moments on a slightly larger interval.
Odd moments vanish because the density is even. Reordering rows
and columns together by parity therefore proves the exact further
factorization
\[
\mathcal D_a(Q;s)
=\det[\mu_{2(Q+i+j)}^{(s)}]_{0\le i,j<\lceil a/2\rceil}
 \det[\mu_{2(Q+1+i+j)}^{(s)}]_{0\le i,j<\lfloor a/2\rfloor}.
\tag{OSP35}
\]
There is no permutation sign, since the same permutation is applied
to both rows and columns. Empty factors are one.
In OSP32 the coefficient of $\log M_s$ is
$\Delta+(\Delta+1)-(q+\Delta)-(q+\Delta+1)-1+q+(q+1)=0$.
This is an exact cancellation only after all original moments
and determinant dimensions have been retained.

Equation OSP33 proves an error-controlled comparison with the exact
finite determinant difference OSP32. It does not obtain a first
variation by differentiating a leading $q^2$ law. In particular no
numerical value of the limit of $\mathcal F_{0,k}^{(s)}/(kq)$ is
asserted by this comparison.

% END CURRENT OSP

\clearpage
% BEGIN CURRENT OAR
\section{The outer-root and actual residual bounds in the original arithmetic return}
\label{sec:OAR}
Retain the complete original data of CEP, CTV and ROQ. In particular
$m=1$, $k=4l+1\ge9$, $q=(k+1)^2$, $q'=(k-7)^2$,
$\Delta=16k-48$, $r=8k-16$ and $j=8k-32$. The quartet, full
amplitude unit, period and branch are the specified original ones on
the proved five-orbit locus. The source orders remain $s\in\{1,k\}$.
Every matrix below is the actual matrix in those chapters, with its
original mass and coefficient frame.

\subsection{The two finite error endpoints retain their signs}
Write
\[
 \mathcal N_+=\{q-1,q\},\qquad
 \mathcal N_-=\{2q-1,2q\},\qquad
 \epsilon_{A,k}^{(s)}=2\sum_{N\in\mathcal N_+\cup\mathcal N_-}E_N^{(s)}.
 \tag{OAR1}
\]
Here $E_N^{(s)}$ is the explicit full conductor constant of CEP35.
The nonnegative constants $L_N^{\rm tail},U_N^{\rm tail}$ are
exactly CTV18 in the same source order and with the original
$q'$, $\delta$, $\gamma$. Define
\[
 \begin{split}
 a_{k,s}^-&=v\left(\sum_{N\in\mathcal N_+}L_N^{\rm tail}
                         +\sum_{N\in\mathcal N_-}U_N^{\rm tail}\right),\\
 a_{k,s}^+&=v\left(\sum_{N\in\mathcal N_+}U_N^{\rm tail}
                         +\sum_{N\in\mathcal N_-}L_N^{\rm tail}\right),\\
 e_{k,s}^-&=\epsilon_{A,k}^{(s)}+a_{k,s}^-,\qquad
 e_{k,s}^+=\epsilon_{A,k}^{(s)}+a_{k,s}^+.
 \end{split}\tag{OAR2}
\]
These are finite quantities for every original period in the stated
domain. For each fixed original period, both $e_{k,s}^\pm=o(kq)$.
Indeed CEP36 proves this for $\epsilon_{A,k}^{(s)}$ and CTV18--20
proves it for each of the four tail contributions, uniformly over
the possible indices $0\le v\le80$ and both original source orders.
When $v=0$, both $a_{k,s}^\pm$ are exactly zero.

Set $\Phi_{k,s}=F_{\Xi,k}^{(s)}$ and retain the full outer quotient
$\mathcal F_{\partial,k}^{(s)}$ of CTV11--13. The exact identity is
\[
 \begin{gathered}
 F_K^{(s)}=\mathcal F_{\partial,k}^{(s)}-\Phi_{k,s}+\mathfrak e_{k,s},\\
 \mathfrak e_{k,s}=\mathcal E_{A,k}^{(s)}
     +\sum_{N\in\mathcal N_+}\mathcal R_{N,v}^{(s)}
     -\sum_{N\in\mathcal N_-}\mathcal R_{N,v}^{(s)},\\
 -e_{k,s}^-\le\mathfrak e_{k,s}\le e_{k,s}^+.
 \end{gathered}\tag{OAR3}
\]
The equality is CTV22, keeping its original two errors. CTV19 says
$-vL_N^{\rm tail}\le\mathcal R_{N,v}^{(s)}\le vU_N^{\rm tail}$.
At a positive endpoint this contributes lower bound $-vL_N^{\rm tail}$
and upper bound $vU_N^{\rm tail}$; at a negative endpoint it
contributes lower bound $-vU_N^{\rm tail}$ and upper bound
$vL_N^{\rm tail}$. Adding these four inequalities and the bounds
$-\epsilon_{A,k}^{(s)}\le\mathcal E_{A,k}^{(s)}
\le\epsilon_{A,k}^{(s)}$ proves each sign of OAR3.

\subsection{Original positive increments give explicit kernel caps}
The full value of $\Phi_{k,s}$ is ROQ19's finite sum
\[
 \Phi_{k,s}=\log(1+\eta_{q-1})
       +2\sum_{N=q}^{2q-2}\log(1+\eta_N)
       +\log(1+\eta_{2q-1}),\qquad
 \lambda_s=\log(1+\eta_{q-1})\ge0.
 \tag{OAR4}
\]
Every $\eta_N$ here is ROQ17's original constrained column, including
the complete conductor--observation cross covariance and its
positive denominator. No new observation is selected.
ROQ21 gives, in the original AKS notation,
\[
 \lambda_s\le\Phi_{k,s}\le
 V_s:=j(\mathcal L_{s,q}+\mathcal L_{s,q+1})-\lambda_s.
 \tag{OAR5}
\]
Retain the complete earlier kernel interval $L_s\le F_K^{(s)}\le U_s$
of AKJ2/MRI4e, with its pointwise minimum of the KAF and AKS bounds.
Define two explicit quantities
\[
 \begin{split}
 L_s^{\partial}&=\max\left\{L_s,
       \mathcal F_{\partial,k}^{(s)}-e_{k,s}^--V_s\right\},\\
 U_s^{\partial}&=\min\left\{U_s,
       \mathcal F_{\partial,k}^{(s)}+e_{k,s}^+-\lambda_s\right\}.
 \end{split}\tag{OAR6}
\]
Then
\[
 \begin{gathered}
 L_s^{\partial}\le F_K^{(s)}\le U_s^{\partial},\\
 \mathcal F_{\partial,k}^{(s)}-\Phi_{k,s}-e_{k,s}^-
       \le F_K^{(s)}
       \le\mathcal F_{\partial,k}^{(s)}-\Phi_{k,s}+e_{k,s}^+.
 \end{gathered}\tag{OAR7}
\]
For the first line substitute both ends of OAR5 into OAR3 and
intersect with the retained interval. The second line substitutes
the full finite value OAR4 into OAR3. Intersecting these intervals
preserves every bound. Both contain the actual value, so their
intersection is nonempty. The dimension $j$ is used only in the
proved finite comparison OAR5. It assigns no leading determinant.
Each additional evaluated positive term of OAR4 can be included
in the upper subtraction in OAR6 by the same inequality.

\subsection{The correlated arithmetic kernel and boundary}
Keep the original arithmetic total $\mathcal B_k^{\rm ar}$,
the scalar $\delta_k^\sigma$ and the four source-difference constants
$C_K^\pm,C_B^\pm$ of MRI4f. They satisfy, for the actual
$d=F_K^{\rm ar}-F_K^{(1)}$,
\[
 \max\{-C_K^-,\delta_k^\sigma-C_B^+\}\le d
       \le\min\{C_K^+,\delta_k^\sigma+C_B^-\}.
 \tag{OAR8}
\]
This is the pair of original kernel bounds together with the
boundary bounds, transferred by
$d+(F_B^{\rm ar}-F_B^{(1)})=\delta_k^\sigma$.
It retains their common scalar and signs. Define
\[
 \begin{split}
 L_{\rm ar}^{\partial}
   &=\max\{L_{\rm ar},L_1^{\partial}-C_K^-,
                    L_1^{\partial}+\delta_k^\sigma-C_B^+\},\\
 U_{\rm ar}^{\partial}
   &=\min\{U_{\rm ar},U_1^{\partial}+C_K^+,
                    U_1^{\partial}+\delta_k^\sigma+C_B^-\}.
 \end{split}\tag{OAR9}
\]
Here $[L_{\rm ar},U_{\rm ar}]$ is the complete interval of MRI4f.
Adding OAR8 to the source-$1$ first line of OAR7 proves
\[
 \begin{gathered}
 L_{\rm ar}^{\partial}\le F_K^{\rm ar}\le U_{\rm ar}^{\partial},\\
 \mathcal B_k^{\rm ar}-U_{\rm ar}^{\partial}
       \le F_B^{\rm ar}
       \le\mathcal B_k^{\rm ar}-L_{\rm ar}^{\partial}.
 \end{gathered}\tag{OAR10}
\]
The second line subtracts the first from the exact original identity
$F_K^{\rm ar}+F_B^{\rm ar}=\mathcal B_k^{\rm ar}$.
It uses the same arithmetic total at both endpoints. If the second
line of OAR7 is used instead, exactly the same addition and
subtraction yield the further interval containing the full residual
sum. Thus every exact positive residual subtraction is available to
the original arithmetic kernel and its correlated boundary.

\subsection{The signed mixed return with its complete finite error}
Write
$\mathfrak S_k^{\rm mix}=F_K^{\rm ar}+F_B^{\rm ar}-F_B^{(k)}$.
The original BRI6d identity is
$\mathfrak S_k^{\rm mix}=F_K^{(k)}+\mathcal H_k+\delta_k^\sigma$.
Consequently OAR3 gives the exact current expression
\[
 \begin{gathered}
 \mathfrak S_k^{\rm mix}=\mathcal F_{\partial,k}^{(k)}-\Phi_{k,k}
              +\mathcal H_k+\delta_k^\sigma+\mathfrak e_{k,k},\\
 -e_{k,k}^-\le\mathfrak e_{k,k}\le e_{k,k}^+,
 \qquad e_{k,k}^\pm=o(kq).
 \end{gathered}\tag{OAR11}
\]
This substitution proves every sign, keeping the original Gamma
return and arithmetic scalar. Addition of their unchanged values
to the source-$k$ bounds OAR7 yields the corresponding finite
intervals for this same signed quantity. The full original
four-endpoint metrics, outer-root polynomial, weighted source,
period-dependent certificates and graph denominator remain in
OAR11 through their proved maps CTV and ROQ. The next leading-scale
calculation is the difference of these two explicit quotient returns.

% END CURRENT OAR

\clearpage
% BEGIN CURRENT PMR
\section{The same-source Gamma moment comparison in the arithmetic return}
\label{sec:PMR}
Retain every original parameter, matrix and error in OAR1--11.
The complete OSP comparison supplies an additional bound directly
on the actual outer-root quotient. Set $\epsilon=2^{-32}$ and
retain its explicit range
\[
 k\ge29,\qquad k\sqrt{\delta^2+\gamma^2}
           \le\epsilon(k+1)^2/2.
 \tag{PMR1}
\]
This range contains all sufficiently large original $k=4l+1$ for
each fixed original quartet: $k\ge\max\{29,
2\sqrt{\delta^2+\gamma^2}/\epsilon\}$ is sufficient because
$(k+1)^2\ge k^2$. For the remaining finite values the full earlier
OAR estimates remain valid without using this comparison.

Keep OSP24's nonnegative $\mathcal E_k$ and put
\[
 c_k^\circ=4(q+\Delta+1)\mathcal E_k,
 \qquad e_{k,s}^{0,-}=e_{k,s}^-+c_k^\circ,
 \qquad e_{k,s}^{0,+}=e_{k,s}^++c_k^\circ.
 \tag{PMR2}
\]
OSP26 and OSP33 prove $c_k^\circ=O_{\delta,\gamma}(q)=o(kq)$
on PMR1, uniformly for the unchanged source orders $s=1,k$.
The actual outer quotient differs from OSP32's exact seven
original Gamma moment determinants by at most $c_k^\circ$.
Consequently the full identity and signed bounds are
\[
 \begin{gathered}
 F_K^{(s)}=\mathcal F_{0,k}^{(s)}-\Phi_{k,s}+\mathfrak e^0_{k,s},\\
 \mathfrak e^0_{k,s}=\mathfrak e_{k,s}
       +\mathcal F_{\partial,k}^{(s)}-\mathcal F_{0,k}^{(s)},\\
 -e_{k,s}^{0,-}\le\mathfrak e^0_{k,s}\le e_{k,s}^{0,+},
 \qquad e_{k,s}^{0,\pm}=o(kq).
 \end{gathered}\tag{PMR3}
\]
The first two equalities are OAR3 with the displayed exact
difference added and subtracted. The two bounds follow by adding
the interval $[-c_k^\circ,c_k^\circ]$ of OSP33 to the unequal
OAR3 endpoints. No metric equality is substituted for the
proved OSP coefficient maps and determinant comparison.
The residual sum $\Phi_{k,s}$ retains all actual period data.

Let $V_s$ and $\lambda_s$ be OAR4--5's exact residual upper
bound and first positive increment. Define
\[
 \begin{split}
 L_s^0&=\max\{L_s^\partial,
                \mathcal F_{0,k}^{(s)}-e_{k,s}^{0,-}-V_s\},\\
 U_s^0&=\min\{U_s^\partial,
                \mathcal F_{0,k}^{(s)}+e_{k,s}^{0,+}-\lambda_s\}.
 \end{split}\tag{PMR4}
\]
Then $L_s^0\le F_K^{(s)}\le U_s^0$. To prove this, insert
$\lambda_s\le\Phi_{k,s}\le V_s$ into PMR3 and intersect
with OAR7. Using the full finite sum for $\Phi_{k,s}$ in PMR3
gives the sharper interval centred at
$\mathcal F_{0,k}^{(s)}-\Phi_{k,s}$. Both retain the same
unequal finite errors of PMR2.

The exact arithmetic correlations of OAR8 give
\[
 \begin{split}
 L_{\rm ar}^0&=\max\{L_{\rm ar}^\partial,L_1^0-C_K^-,
                                      L_1^0+\delta_k^\sigma-C_B^+\},\\
 U_{\rm ar}^0&=\min\{U_{\rm ar}^\partial,U_1^0+C_K^+,
                                      U_1^0+\delta_k^\sigma+C_B^-\}.
 \end{split}\tag{PMR5}
\]
Adding those correlations to PMR4 at $s=1$ proves
$L_{\rm ar}^0\le F_K^{\rm ar}\le U_{\rm ar}^0$.
Subtracting from the exact same arithmetic total gives
\[
 \mathcal B_k^{\rm ar}-U_{\rm ar}^0
      \le F_B^{\rm ar}\le\mathcal B_k^{\rm ar}-L_{\rm ar}^0.
 \tag{PMR6}
\]
Thus the actual kernel and boundary retain their common scalar
and both original source comparisons.

Finally the original signed-return identity BRI6d and PMR3 at
source order $k$ give
\[
 \begin{gathered}
 F_K^{\rm ar}+F_B^{\rm ar}-F_B^{(k)}
   =\mathcal F_{0,k}^{(k)}-\Phi_{k,k}
      +\mathcal H_k+\delta_k^\sigma+\mathfrak e^0_{k,k},\\
 -e_{k,k}^{0,-}\le\mathfrak e^0_{k,k}\le e_{k,k}^{0,+}.
 \end{gathered}\tag{PMR7}
\]
This follows by substituting into
$F_K^{\rm ar}+F_B^{\rm ar}-F_B^{(k)}=F_K^{(k)}+
\mathcal H_k+\delta_k^\sigma$; every sign is retained.
The entire original same-source moment expression is OSP32--35.
The next asymptotic calculation must evaluate it with the actual
residual sum and an error small compared with $kq$.

\paragraph{The exact finite endpoints under this comparison.}
The interval transport in PMR4--6 preserves the already retained
outer-root endpoints exactly:
\[
 L_s^0=L_s^\partial,\qquad U_s^0=U_s^\partial,\qquad
 L_{\rm ar}^0=L_{\rm ar}^\partial,\qquad
 U_{\rm ar}^0=U_{\rm ar}^\partial.
 \tag{PMR8}
\]
For proof, OSP33 gives
$\mathcal F_{0,k}^{(s)}-c_k^\circ\le\mathcal F_{\partial,k}^{(s)}
\le\mathcal F_{0,k}^{(s)}+c_k^\circ$.
The new lower candidate in PMR4 is therefore at most
$\mathcal F_{\partial,k}^{(s)}-e_{k,s}^--V_s$, which is at
most $L_s^\partial$ by OAR6. The new upper candidate is at
least $\mathcal F_{\partial,k}^{(s)}+e_{k,s}^+-\lambda_s$,
which is at least $U_s^\partial$. Taking the displayed maximum
and minimum proves the source equalities. Substitution in PMR5
and the defining maxima and minima OAR9 proves the arithmetic
equalities; subtracting from the same total proves equality of
the boundary endpoints as well. Thus this comparison supplies
the exact same-source moment target and controlled error; a
further numerical improvement must use sharper evaluated moment
or residual information. The preceding ROQ/OAR bounds remain
the proved finite intervals throughout this transport.

% END CURRENT PMR


\clearpage
% BEGIN CURRENT EIQ
% Independent proof. No external one-cut or equilibrium theorem is used.
% May be included as a section in the cumulative paper.
\section{Exact equilibrium for the square-root logarithmic potential}
\label{sec:independent-sqrt-log-equilibrium}

Throughout this section $\alpha>0$ and $\beta>0$.  On $(0,\infty)$ set
\begin{equation}\tag{EIQ1}
 V_{\alpha,\beta}(x)=\beta\sqrt{x}-\alpha\log x,
 \qquad
 \mathcal E_{\alpha,\beta}(\mu)
 =\int V_{\alpha,\beta}(x)\,d\mu(x)
       -\iint\log|x-y|\,d\mu(x)d\mu(y).
\end{equation}
The measures considered initially are probability measures for which
$\int V_{\alpha,\beta}\,d\mu$ is finite.  Because this potential is bounded
below and tends to $+\infty$ at both ends, this implies
$\int(\sqrt{x}+|\log x|)\,d\mu<\infty$.
The positive part of $\log|x-y|$ is then integrable, by
$\log^+|x-y|\leq\log(1+x)+\log(1+y)$.
Thus the energy is well defined with value in $\mathbb R\cup\{+\infty\}$.
All comparisons below also apply to the usual extended energy defined by
integrating the kernel
$\tfrac12V_{\alpha,\beta}(x)+\tfrac12V_{\alpha,\beta}(y)-\log|x-y|$:
the kernel is bounded below, and finite energy for this kernel implies the
same moment conditions.  To see the latter assertion without cancellation,
put $g(x)=\tfrac12V_{\alpha,\beta}(x)-\log(1+x)$.
The kernel is at least $g(x)+g(y)$; $g$ is bounded below and controls positive
constant multiples of both $\sqrt{x}$ near infinity and $|\log x|$ near zero.

\subsection{Existence and uniqueness of the two endpoints}
For $0<k<1$ define the complete elliptic integrals by their actual real
integrals
\begin{equation}\tag{EIQ2}
 K(k)=\int_0^{\pi/2}\frac{d\theta}
            {\sqrt{1-k^2\sin^2\theta}},\qquad
 E(k)=\int_0^{\pi/2}\sqrt{1-k^2\sin^2\theta}\,d\theta.
\end{equation}
There is exactly one $k\in(0,1)$ satisfying
\begin{equation}\tag{EIQ3}
 \frac{\sqrt{1-k^2}K(k)}{E(k)}=\frac{\alpha}{\alpha+2}.
\end{equation}
For this $k$ put
\begin{equation}\tag{EIQ4}
 r=\sqrt{1-k^2},\qquad
 v=\frac{(\alpha+2)\pi}{\beta E(k)},\qquad u=rv,
 \qquad a=u^2,\quad b=v^2.
\end{equation}
In particular $0<a<b<\infty$, and the endpoint equations, with all
constants retained, are
\begin{equation}\tag{EIQ5}
 uK(k)=\frac{\alpha\pi}{\beta},\qquad
 vE(k)=\frac{(\alpha+2)\pi}{\beta},\qquad
 k^2=1-\frac{u^2}{v^2}.
\end{equation}

Here is a proof of the assertion of uniqueness. Differentiation of the
integrals in (EIQ2), and integration of the derivative of
$\sin\theta\cos\theta/\sqrt{1-k^2\sin^2\theta}$ between its two zero
endpoint values, give
\[
 E'(k)=\frac{E(k)-K(k)}{k},\qquad
 K'(k)=\frac{E(k)}{k(1-k^2)}-\frac{K(k)}{k}.
\]
For completeness the second identity is equivalent, after differentiation
under the integral sign, to
\[
 \int_0^{\pi/2}
 \left\{\frac{k^2\sin^2\theta}{(1-k^2\sin^2\theta)^{3/2}}
 -\frac{1-k^2\sin^2\theta}{(1-k^2)\sqrt{1-k^2\sin^2\theta}}
 +\frac1{\sqrt{1-k^2\sin^2\theta}}\right\}d\theta=0;
\]
the expression in braces is
$-k^2(1-k^2)^{-1}\,d(\sin\theta\cos\theta/
\sqrt{1-k^2\sin^2\theta})/d\theta$.
With $t=E(k)/K(k)$, the integral definitions show strictly
$r^2<t<1$.  Consequently, if $f(k)=rK(k)/E(k)$, then
\begin{equation}\tag{EIQ6}
 \frac{kf'(k)}{f(k)}
 =\frac{t}{r^2}+\frac1t-1-\frac1{r^2}
 =\frac{(t-1)(t-r^2)}{r^2t}<0.
\end{equation}
At $k=0$ the continuous limiting value is $f(0)=1$.
As $k\uparrow1$, the integrand of $rK(k)$ tends to zero at every
$\theta<\pi/2$ and is bounded by $1$, while $E(k)\to1$ by dominated
convergence.  Hence $f(k)\to0$.  This proves (EIQ3), including existence.
It also proves smooth dependence of $k,u,v,a,b$ on $(\alpha,\beta)$ in
$(0,\infty)^2$, by the nonzero derivative in (EIQ6).

Let $m=(a+b)/2$, $R_s=\sqrt{(a+s)(b+s)}$ for $s\geq0$, and
$R_0=uv$.  Substitution
$t=a\cos^2\theta+b\sin^2\theta$ gives, directly from (EIQ5),
\begin{equation}\tag{EIQ7}
 \int_a^b\frac{V'_{\alpha,\beta}(t)}
                  {\sqrt{(b-t)(t-a)}}\,dt=0,
 \qquad
 \int_a^b\frac{tV'_{\alpha,\beta}(t)}
                  {\sqrt{(b-t)(t-a)}}\,dt=2\pi.
\end{equation}
Indeed the three integrals with numerators $t^{-1/2},t^{-1},t^{1/2}$
are respectively $2K(k)/v$, $\pi/(uv)$, and $2vE(k)$, and
$V'(t)=\beta/(2\sqrt t)-\alpha/t$.

\subsection{A positive density and its exact Cauchy transform}
The elementary Stieltjes integral, obtained by $s=xz^2$, is
\begin{equation}\tag{EIQ8}
 \int_0^\infty\frac{ds}{\sqrt s(x+s)}=\frac\pi{\sqrt x}
 \quad(x>0).
\end{equation}
Using this identity in (EIQ7) and applying Tonelli to its positive
constituents gives
\begin{equation}\tag{EIQ9}
 \frac\alpha{R_0}=
 \frac\beta{2\pi}\int_0^\infty\frac{ds}{\sqrt s R_s},
 \qquad
 \int_0^\infty\frac{1-s/R_s}{\sqrt s}\,ds=2vE(k).
\end{equation}
Both integrals converge; their integrands are $O(s^{-1/2})$ at zero
and $O(s^{-3/2})$ at infinity.  The second formula also follows by
averaging $\sqrt t=\pi^{-1}\int_0^\infty
s^{-1/2}t/(t+s)\,ds$ against the arcsine probability measure on $[a,b]$.

Define for every $x>0$
\begin{align}\tag{EIQ10}
 h(x)
 &=\frac\alpha{xR_0}
   -\frac\beta{2\pi}
       \int_0^\infty\frac{ds}{\sqrt s(x+s)R_s} \\
 &=\frac\beta{2\pi x}
       \int_0^\infty\frac{\sqrt s}{(x+s)R_s}\,ds>0.\notag
\end{align}
The second equality is exactly (EIQ9), with no change of the potential.
All integrals in (EIQ10) converge absolutely and locally uniformly in
$x>0$.  The candidate density on the original positive half-line is
\begin{equation}\tag{EIQ11}
 \rho_{\alpha,\beta}(x)=
 \begin{cases}
 \displaystyle
 \frac{\sqrt{(b-x)(x-a)}}{2\pi}\,h(x),&a<x<b,\\
 0,&x\notin(a,b).
 \end{cases}
\end{equation}
It is strictly positive on $(a,b)$ and vanishes continuously with the
square-root factor at each endpoint.  Positivity here follows from the
last integral in (EIQ10), not from a one-interval ansatz.

The following elementary transform computes its mass and resolvent.
Choose \[R(z)=\sqrt{(z-a)(z-b)}\]
on $\mathbb C\setminus[a,b]$ with
$R(z)/z\to1$ at infinity. Thus $R(x)>0$ for $x>b$,
$R(x)<0$ for $x<a$, and $R(x+i0)=i\sqrt{(b-x)(x-a)}$ on $(a,b)$.
Direct substitution $t=m+(b-a)\cos\theta/2$, followed by
$w=\tan(\theta/2)$, proves
\[
 \frac1\pi\int_a^b\frac{dt}{(z-t)\sqrt{(b-t)(t-a)}}=\frac1{R(z)}.
\]
For real $z>b$, the substitution reduces the integral to
$\pi^{-1}\int_0^\infty2\,dw/[(z-b)+(z-a)w^2]$;
the identity elsewhere follows by analytic continuation.
Polynomial division then gives
\begin{equation}\tag{EIQ12}
 \frac1\pi\int_a^b\frac{\sqrt{(b-t)(t-a)}}{z-t}\,dt
     =z-m-R(z),\qquad
 \frac1\pi\int_a^b\frac{\sqrt{(b-t)(t-a)}}{t+s}\,dt
     =m+s-R_s.
\end{equation}
For example the first identity follows from
$(b-t)(t-a)=-(t-z)^2+2(m-z)(t-z)-(z-a)(z-b)$,
the preceding arcsine transform, and the arcsine first moment $m$.
Partial fractions in these two identities yield
\begin{equation}\tag{EIQ13}
 H_s(z):=\frac1\pi\int_a^b
       \frac{\sqrt{(b-t)(t-a)}}{(t+s)(z-t)}\,dt
       =1-\frac{R_s+R(z)}{z+s}.
\end{equation}
The formula at $z=-s$ is interpreted by its removable limit.

The mass of (EIQ11), computed using (EIQ12), is
\begin{align}\tag{EIQ14}
 \int_a^b\rho_{\alpha,\beta}(t)\,dt
 &=\frac\alpha{2R_0}(m-R_0)
 -\frac\beta{4\pi}\int_0^\infty
       \frac{(m+s)/R_s-1}{\sqrt s}\,ds\\
 &=-\frac\alpha2+
   \frac\beta{4\pi}\int_0^\infty\frac{1-s/R_s}{\sqrt s}\,ds
 =-\frac\alpha2+\frac{\beta vE(k)}{2\pi}=1.\notag
\end{align}
In particular this is a probability density.

Define its Cauchy transform with the explicit sign convention
$G(z)=\int_a^b\rho_{\alpha,\beta}(t)/(z-t)\,dt$.
For $z$ outside both $[a,b]$ and $(-\infty,0]$, formulas (EIQ10) and
(EIQ13) give
\begin{align}\tag{EIQ15}
 G(z)
 &=\frac\alpha{2R_0}H_0(z)
       -\frac\beta{4\pi}
          \int_0^\infty\frac{H_s(z)}{\sqrt sR_s}\,ds\\
 &=-\frac\alpha{2z}+\frac\beta{4\sqrt z}
       -\frac{R(z)h(z)}2
 =\frac{V'_{\alpha,\beta}(z)-R(z)h(z)}2.\notag
\end{align}
Here the positive square root is continued with its negative-real-axis
cut, and $h(z)$ is the first integral formula of (EIQ10).  Every integral
exchange is absolutely convergent, uniformly on compact subsets of this
domain: near zero the auxiliary integrand is $O(s^{-1/2})$, and near
infinity the original integral defining $H_s$ is $O(s^{-1})$.
The constant terms on the second line cancel by (EIQ9), and (EIQ8)
supplies the term $\beta/(4\sqrt z)$.

\subsection{The Euler equality, its strict exterior inequality, and uniqueness}
Set
\begin{equation}\tag{EIQ16}
 U(x)=V_{\alpha,\beta}(x)
             -2\int_a^b\log|x-t|\rho_{\alpha,\beta}(t)\,dt
 \quad(x>0).
\end{equation}
This function is continuous on $(0,\infty)$.  The integral is continuous
because the density is bounded, compactly supported, and the integral
of $|\log|x-t||$ over an interval of length $\varepsilon$ tends to zero
uniformly in its centre.  In the interior of $[a,b]$, differentiation of
the logarithmic potential is the principal-value integral
$\operatorname{PV}\int\rho(t)/(x-t)\,dt$.
This follows by subtracting $\rho(x)$ on a symmetric interval about
$x$; the remaining numerator is $O(|t-x|)$ because the density is
continuously differentiable locally in the interior.  The real part of
(EIQ15) on the upper edge of the interval is precisely this principal
value.  Indeed subtracting $\rho(x)$ in
$\int\rho(t)/(x+i\varepsilon-t)\,dt$ gives the same limit, while the
constant term has real part equal to its elementary principal value.
Since $R(x+i0)$ is purely imaginary and $h(x)$ is real, (EIQ15) proves
$U'(x)=0$ on $(a,b)$.  There is therefore a real constant $\ell$ such that
\begin{equation}\tag{EIQ17}
 U(x)=\ell\quad(a\leq x\leq b).
\end{equation}
For positive $x$ outside this interval, (EIQ15) instead gives the exact
real derivative
\begin{equation}\tag{EIQ18}
 U'(x)=R(x)h(x)
 \begin{cases}<0,&0<x<a,\\>0,&x>b.\end{cases}
\end{equation}
Combining (EIQ17), continuity, and these signs proves
\begin{equation}\tag{EIQ19}
 U(x)>\ell\quad\hbox{for }x\in(0,a)\cup(b,\infty).
\end{equation}
This verifies the entire exterior inequality on the original domain,
including the interval next to zero.

Here is a complete positivity argument for the energy of the difference
of two measures.  If $\nu$ is a finite signed real measure of total mass
zero for which $|\log|x-y||$ is integrable against
$|\nu|\otimes|\nu|$, the scalar identity
\begin{equation}\tag{EIQ20}
 -\log d=\frac12\int_0^\infty
           \frac{e^{-td^2}-e^{-t}}t\,dt\quad(d>0)
\end{equation}
follows by differentiating in $d$ and evaluating at $d=1$.
One half of the integral over any finite subinterval of $(0,\infty)$
has the same sign as $-\log d$ and absolute value at most $|\log d|$.
Consequently dominated convergence and the zero total mass give
\begin{equation}\tag{EIQ21}
 -\iint\log|x-y|\,d\nu(x)d\nu(y)
 =\frac12\int_0^\infty\frac1t
       \iint e^{-t(x-y)^2}\,d\nu(x)d\nu(y)\,dt\geq0.
\end{equation}
The nonnegative sign is proved, rather than assumed, by the Gaussian
Fourier identity
\begin{equation}\tag{EIQ22}
 \iint e^{-t(x-y)^2}\,d\nu(x)d\nu(y)
 =\frac1{\sqrt{4\pi t}}\int_{\mathbb R}
          e^{-\xi^2/(4t)}|\widehat\nu(\xi)|^2\,d\xi,
 \qquad \widehat\nu(\xi)=\int e^{i\xi x}\,d\nu(x).
\end{equation}
The scalar Gaussian Fourier identity follows by completing the square
in the Gaussian integral (or by differentiating its Fourier transform
and its value at zero); its integration against the finite measures is
justified by absolute integrability.  Equality in (EIQ21) forces
$\widehat\nu=0$: the right side of (EIQ22) is strictly positive for every
$t>0$ if this continuous Fourier transform has a nonzero value.
Finally $\widehat\nu=0$ implies $\nu=0$ without an extra uniqueness
assumption. Convolve $\nu$ with a centred Gaussian approximate identity.
Its value at every point is zero by that same Fourier integral.
For every compactly supported continuous test function, its convolution
with these Gaussians converges uniformly to the test function; integration
against $\nu$ then proves that $\nu$ annihilates every such function.

Let $\mu$ be any probability measure of finite energy in (EIQ1), and
write $\mu_*(dx)=\rho_{\alpha,\beta}(x)\,dx$ and $\nu=\mu-\mu_*$.
The integrability needed for (EIQ21) holds.  The $\mu\otimes\mu$ term
has it by finite energy and the already established logarithmic moment;
the $\mu_*\otimes\mu_*$ term has it by bounded compact support; and the
cross term has it because the negative logarithmic singularity integrated
against the bounded density is uniformly bounded, while the positive
part is controlled by $\log(1+x)+\log(1+y)$.
Expanding the energy therefore yields exactly
\begin{equation}\tag{EIQ23}
 \mathcal E_{\alpha,\beta}(\mu)-
       \mathcal E_{\alpha,\beta}(\mu_*)
 =\int(U(x)-\ell)\,d\mu(x)
       -\iint\log|x-y|\,d\nu(x)d\nu(y)\geq0.
\end{equation}
The first term is nonnegative by (EIQ17)--(EIQ19), and the second is
strictly positive unless $\mu=\mu_*$. Infinite energy competitors cannot
improve this finite value. Thus (EIQ11) is the unique global minimizing
probability measure, with exactly the support $[a,b]$.

\subsection{An explicit endpoint integral for the logarithmic moment}
Retain $m,R_s,u,v$ as defined above, and define
\begin{equation}\tag{EIQ24}
 C=\frac{(u+v)^2}4,\qquad
 \eta=\frac{v-u}{v+u},\qquad
 \zeta_s=\frac{b-a}{2(m+s+R_s)}\quad(s\geq0).
\end{equation}
Thus $0<\zeta_s\leq\eta<1$ and $\zeta_0=\eta$.
For every $s\geq0$ put
\begin{equation}\tag{EIQ25}
 L_s=(m+s-R_s)\log C-(m-R_0)
                  -2R_s\log(1-\eta\zeta_s).
\end{equation}
Then the logarithmic moment of the equilibrium measure is exactly
\begin{equation}\tag{EIQ26}
 \mathcal L(\alpha,\beta):=\int_a^b\log x\rho_{\alpha,\beta}(x)\,dx
 =\frac\beta{4\pi}\int_0^\infty
                    \frac{L_0-L_s}{\sqrt s R_s}\,ds.
\end{equation}
This is a convergent one-dimensional endpoint integral of explicitly
given real elementary functions; its only endpoints are those specified
uniquely in (EIQ3)--(EIQ5).

Here is a direct derivation, including the constants in (EIQ25).
Set $h_0=(b-a)/2$ and substitute $t=m+h_0\cos\theta$ in arcsine
integrals.  The identities
\[
 t=C(1+2\eta\cos\theta+\eta^2),\qquad
 t+s=\frac{m+s+R_s}2
                (1+2\zeta_s\cos\theta+\zeta_s^2)
\]
give the absolutely convergent Fourier series
\[
 \log t=\log C+
       2\sum_{j\geq1}\frac{(-1)^{j+1}\eta^j}j\cos(j\theta),
 \qquad
 \frac1{t+s}=\frac1{R_s}
       \left(1+2\sum_{j\geq1}(-\zeta_s)^j\cos(j\theta)\right).
\]
Orthogonality on $[0,\pi]$ gives
\begin{align}\tag{EIQ27}
 \frac1\pi\int_a^b\frac{\log t}{\sqrt{(b-t)(t-a)}}\,dt
 &=\log C,\\
 \frac1\pi\int_a^b\frac{t\log t}{\sqrt{(b-t)(t-a)}}\,dt
 &=m\log C+(m-R_0),\notag\\
 \frac1\pi\int_a^b\frac{\log t}{(t+s)\sqrt{(b-t)(t-a)}}\,dt
 &=\frac{\log C+2\log(1-\eta\zeta_s)}{R_s}.\notag
\end{align}
The second identity uses $h_0\eta=m-R_0$.
Polynomial division
\[
 \frac{(b-t)(t-a)}{t+s}
       =-t+2m+s-\frac{R_s^2}{t+s}
\]
now proves, exactly,
\begin{equation}\tag{EIQ28}
 L_s=\frac1\pi\int_a^b
               \frac{\log t\sqrt{(b-t)(t-a)}}{t+s}\,dt.
\end{equation}
Using the positive integral expression for $h$ in (EIQ10), and the
identity $1/[t(t+s)]=(1/t-1/(t+s))/s$, yields (EIQ26).
The exchanges are absolutely convergent: $\log t$ is bounded on
$[a,b]$, $L_0-L_s=O(s)$ at zero by (EIQ28), and $L_s=O(s^{-1})$
at infinity.  Thus the integrand in (EIQ26) is $O(s^{1/2})$ at zero
and $O(s^{-3/2})$ at infinity.

\subsection{Parameter continuity and the variational derivative}
The density and its logarithmic moment depend smoothly on
$(\alpha,\beta)\in(0,\infty)^2$ in the following precise sense.
On a compact parameter neighbourhood, (EIQ3)--(EIQ6) put $a,b$ in a
fixed compact subset of $(0,\infty)$, with $b-a$ bounded below by a
positive number.  The integral for $h$ and each of its parameter and
$x$ derivatives converge uniformly for $x$ in this fixed compact set:
their absolute integrands are bounded by constants times $s^{1/2}$
near zero and $s^{-3/2}$ near infinity. Parameter derivatives of the
endpoints are bounded on a smaller compact neighbourhood.
Under the fixed-interval coordinate
$x=a+(b-a)y$, $0\leq y\leq1$, the measure is exactly
\begin{equation}\tag{EIQ29}
 \rho_{\alpha,\beta}(x)\,dx
 =\frac{(b-a)^2}{2\pi}\sqrt{y(1-y)}
                  h\bigl(a+(b-a)y\bigr)\,dy.
\end{equation}
All parameter derivatives of its coefficient are bounded by a constant
times $\sqrt{y(1-y)}$. The same is true after multiplication by
$\log(a+(b-a)y)$ or by $\sqrt{a+(b-a)y}$.
Differentiation under this integral proves smoothness, in particular
continuity, of both moments on the entire positive parameter quadrant.

Define $E_{\min}(\alpha,\beta)=\mathcal E_{\alpha,\beta}(\mu_{\alpha,\beta})$,
where the measure is the unique minimizer proved above.
Testing the minimizer for a parameter pair against the neighbouring
potential, and then reversing the two tests, gives for $h>0$
\[
 -\mathcal L(\alpha+h,\beta)
 \leq\frac{E_{\min}(\alpha+h,\beta)-E_{\min}(\alpha,\beta)}h
 \leq-\mathcal L(\alpha,\beta).
\]
The corresponding inequalities for $h<0$ and the proved continuity
therefore give
\begin{equation}\tag{EIQ30}
 \partial_\alpha E_{\min}(\alpha,\beta)=-\mathcal L(\alpha,\beta),\qquad
 \partial_\beta E_{\min}(\alpha,\beta)
        =\int\sqrt{x}\,d\mu_{\alpha,\beta}(x).
\end{equation}
No unproved differentiability of a minimizing selection is being used.
For the opposite convention
$F(\alpha,\beta)=\sup_\mu\{\iint\log|x-y|\,d\mu(x)d\mu(y)
-\int V_{\alpha,\beta}\,d\mu\}=-E_{\min}(\alpha,\beta)$,
the exact derivative is therefore
$\partial_\alpha F(\alpha,\beta)=\mathcal L(\alpha,\beta)$.

The second moment in (EIQ30) has an exact value.  Push the minimizing
measure forward by the explicit dilation $x\mapsto c^2x$, $c>0$.
Writing $M=\int\sqrt{x}\,d\mu_{\alpha,\beta}(x)$, its energy minus
the energy at $c=1$ is precisely
$\beta(c-1)M-2(\alpha+1)\log c$.
Its derivative at the minimizing value $c=1$ is zero. Hence
\begin{equation}\tag{EIQ31}
 \int\sqrt{x}\,d\mu_{\alpha,\beta}(x)
       =\frac{2(\alpha+1)}\beta.
\end{equation}
The same explicit dilation, or (EIQ3)--(EIQ5), proves that the map
$x\mapsto(\pi/\beta)^2x$ pushes $\mu_{\alpha,\pi}$ to
$\mu_{\alpha,\beta}$, including the density and its two endpoints.
Consequently
\begin{equation}\tag{EIQ32}
 \mathcal L(\alpha,\beta)
       =\mathcal L(\alpha,\pi)+2\log(\pi/\beta).
\end{equation}

For the particular original limiting potential
$V(x)=\pi\sqrt{x}-2\log x$, the substitutions are exactly
$\alpha=2$, $\beta=\pi$. Equations (EIQ5), (EIQ10), and (EIQ26) become
\begin{equation}\tag{EIQ33}
 uK(k)=2,\qquad vE(k)=4,\qquad
 h(x)=\frac1{2x}\int_0^\infty\frac{\sqrt s}{(x+s)R_s}\,ds,
 \qquad
 \mathcal L(2,\pi)=\frac14\int_0^\infty
                         \frac{L_0-L_s}{\sqrt sR_s}\,ds.
\end{equation}
In particular the preceding proof includes positivity, unit mass,
the full Euler inequality, the unique global minimizer, and the
parameter continuity needed at these actual parameters.

% END CURRENT EIQ

\clearpage
% BEGIN CURRENT QLG
% Complete finite-n proof, 14 September 2026.
% Provider: equilibrium/independent/EXACT_SQRT_LOG_EQUILIBRIUM.tex, EIQ1--33.
% The Heine factor 1/n! and both reference measures are retained explicitly.
\section{A uniform finite partition estimate from the original equilibrium}
\label{sec:QLG}

\paragraph{QLG1. Exact measures and parameter domain.}
Fix the compact parameter rectangle
\[
 \mathcal K=[3/2,5/2]\times[\pi/2,3\pi/2],\qquad
 V_{\alpha,\beta}(x)=\beta\sqrt{x}-\alpha\log x\quad(x>0).
 \tag{QLG1}
\]
The equilibrium proved in EIQ1--33 has density
\[
 \rho(x)=\frac{\beta\sqrt{(b-x)(x-a)}}{4\pi^2x}
       \int_0^\infty
       \frac{\sqrt{s}\,ds}{(x+s)\sqrt{(a+s)(b+s)}}\quad(a<x<b),
 \qquad \rho(x)=0\quad(x\notin(a,b)),
 \tag{QLG2}
\]
where the exact endpoints are those of EIQ3--5. In particular
\(0<a<b<\infty\), the density is bounded and continuous, positive
in the interior, and has integral one. Write
\[
 d\mu(x)=\rho(x)\,dx,\quad
 P(z)=\int_a^b\log|z-t|\,d\mu(t)\quad(z\in\mathbb C),\quad
 I(\mu)=\iint\log|x-y|\,d\mu(x)d\mu(y),
\]
\[
 F(\alpha,\beta)=I(\mu)-\int V_{\alpha,\beta}\,d\mu.
 \tag{QLG3}
\]
The original Euler equality and full exterior inequality in EIQ17--19
give a real number \(\ell\) for which
\[
 D(x)=V_{\alpha,\beta}(x)-2P(x)-\ell\geq0\quad(x>0),
 \qquad D(x)=0\quad(a\leq x\leq b).
 \tag{QLG4}
\]
Integrating the equality on the support gives the exact sign identities
\[
 \ell=\int V_{\alpha,\beta}\,d\mu-2I(\mu),\qquad
 F(\alpha,\beta)=-\ell-I(\mu).
 \tag{QLG5}
\]
The two integration measures treated below are, separately,
\[
 d\lambda_0(x)=dx,\qquad
 d\lambda_\omega(x)=d\omega(x)=\tfrac12 x^{-1/2}e^{-\sqrt{x}}\,dx.
 \tag{QLG6}
\]
The map \(x=r^2\) gives \(\int d\omega=\int_0^\infty e^{-r}dr=1\).
For either of these measures, with its stated density \(w_\lambda\),
define the literal finite integral
\[
 J_n^\lambda(\alpha,\beta)=\frac1{n!}
 \int_{(0,\infty)^n}
 \exp\!\left(2\sum_{i<j}\log|x_i-x_j|
                 -n\sum_{i=1}^n V_{\alpha,\beta}(x_i)\right)
       \prod_{i=1}^n d\lambda(x_i),\qquad n\geq1.
 \tag{QLG7}
\]
No source measure, scale, factorial, or exponent has been absorbed into
\(F\). The unweighted integral without the Heine factor is exactly
\(n!J_n^\lambda\).

\paragraph{QLG2. Constants proved finite on the whole rectangle.}
Let
\[
 a_* =\min_{(\alpha,\beta)\in\mathcal K} a(\alpha,\beta)>0,
 \qquad B_* =\max_{(\alpha,\beta)\in\mathcal K} b(\alpha,\beta)<\infty,
 \qquad \beta_-=\pi/2,\quad\beta_+=3\pi/2.
 \tag{QLG8}
\]
These are explicit compact-parameter extrema of the uniquely specified
endpoint integrals EIQ3--5. Their asserted strict inequalities follow
from the continuous positive endpoints proved in EIQ6 and compactness
of \(\mathcal K\). Thus their definition requires no existence or
regularity assumption about an unspecified minimizing measure.
For \(x\in[a,b]\), the integral in QLG2 satisfies
\[
 \int_0^\infty\frac{\sqrt{s}\,ds}{(x+s)\sqrt{(a+s)(b+s)}}
 \leq\int_0^\infty\frac{\sqrt{s}\,ds}{(a_*+s)^2}
 =\frac{\pi}{2\sqrt{a_*}}.
\]
For the last equality substitute \(s=a_*r^2\), followed by
\(r=\tan\theta\); the resulting integral is
\[2a_*^{-1/2}\int_0^{\pi/2}\sin^2\theta\,d\theta.\]
Since \(\sqrt{(b-x)(x-a)}\leq(b-a)/2\leq B_*/2\), it follows that
\[
 0\leq\rho(x)\leq M_*:=\frac{\beta_+B_*}{16\pi a_*^{3/2}}
       \quad\hbox{for all real }x\hbox{ and all parameters in }\mathcal K.
 \tag{QLG9}
\]
In particular \(B_*M_*\geq1\), by the mass-one identity and support
length at most \(B_*\).

The entropy relative to either stated measure is
\[
 H_\lambda(\mu)=\int_a^b
                 \rho(x)\log\frac{\rho(x)}{w_\lambda(x)}\,dx,
       \qquad 0\log0=0.
 \tag{QLG10}
\]
It is absolutely integrable: the function \(u|\log u|\) is bounded
on \([0,M_*]\), and \(\log w_\lambda\) is bounded on the fixed
compact interval \([a_*,B_*]\). In the two cases its upper bounds are
\[
 H_{\lambda_0}(\mu)\leq\log M_* ,\qquad
 H_{\lambda_\omega}(\mu)\leq
 \log M_*+\log2+\tfrac12\log B_*+\sqrt{B_*}.
 \tag{QLG11}
\]
All pair logarithmic energies on this support are absolutely
integrable. For the negative part, integration against a density
bounded by \(M_*\) is bounded above by
\(M_*\int_{-1}^1-\log|s|\,ds=2M_*\), uniformly in the centre.
For the positive part the bound \(\log^+ B_*\) suffices.

Put
\[
 T_\lambda=\int_0^\infty e^{-D(x)}\,d\lambda(x).
 \tag{QLG12}
\]
For \(\lambda=\lambda_\omega\), QLG4 and the mass of \(\omega\)
give \(0<T_{\lambda_\omega}\leq1\). For the Lebesgue measure a finite
uniform bound can also be displayed entirely. For every \(x>0\),
\[
 P(x)\leq\log(1+x)+\log(1+B_*),\qquad P(x)\geq-2M_*.
\]
The lower bound uses only the negative part estimate just proved.
Evaluating QLG4 at any point of \([a,b]\) therefore yields
\[
 \ell\leq L_*:=\beta_+\sqrt{B_*}
    +\tfrac52\max\{|\log a_*|,|\log B_*|\}+4M_*.
\]
Consequently
\[
 T_{\lambda_0}\leq e^{L_*}(1+B_*)^2
    \int_0^\infty (1+x)^2(1+x^{5/2})e^{-\beta_-\sqrt{x}}\,dx
    =:T_*<\infty.
 \tag{QLG13}
\]
Here \(x^\alpha\leq1+x^{5/2}\) for all parameters in \(\mathcal K\).
The integral in QLG13 equals, by \(x=r^2\) and elementary repeated
integration by parts,
\[
 2\left(\frac{1!}{\beta_-^2}
       +\frac{2\,3!}{\beta_-^4}
       +\frac{5!}{\beta_-^6}
       +\frac{6!}{\beta_-^7}
       +\frac{2\,8!}{\beta_-^9}
       +\frac{10!}{\beta_-^{11}}\right).
 \tag{QLG14}
\]

\paragraph{QLG3. The planar logarithmic comparison, with its proof.}
We use the plane only to construct an exact comparison of logarithmic
kernels; the original particles and potential remain on \((0,\infty)\).
If \(\eta\) is a finite signed measure of compact support on
\(\mathbb R^2\), has total mass zero, and
\(\iint|\log|z-w||\,d|\eta|(z)d|\eta|(w)<\infty\), then
\[
 I(\eta):=\iint\log|z-w|\,d\eta(z)d\eta(w)\leq0.
 \tag{QLG15}
\]
To prove this, for \(d>0\) integrate the elementary identity
\[
 -\log d=\frac12\int_0^\infty
                    \frac{e^{-td^2}-e^{-t}}{t}\,dt.
\]
It follows by differentiation in \(d\) and evaluation at \(d=1\).
On any finite subinterval of \((0,\infty)\), the integral on the
right has the same sign as \(-\log d\) and absolute value at most
\(|\log d|\). The stipulated logarithmic integrability therefore
permits dominated passage to the full integral against \(d\eta^2\).
The constant term vanishes by the zero mass. The elementary Gaussian
Fourier integral in each of the two real coordinates gives
\[
 -I(\eta)=\frac12\int_0^\infty\frac{dt}{t}
       \iint e^{-t|z-w|^2}\,d\eta(z)d\eta(w),
\]
\[
 \iint e^{-t|z-w|^2}\,d\eta(z)d\eta(w)
 =\frac1{4\pi t}\int_{\mathbb R^2}
       e^{-|\xi|^2/(4t)}
       \left|\int e^{i\xi\cdot z}\,d\eta(z)\right|^2d\xi\geq0.
 \tag{QLG16}
\]
For completeness the scalar Fourier identity is obtained by
differentiating \(\int e^{-u^2/(4t)}e^{isu}du\) with respect to
\(s\), integrating by parts in \(u\), and using its value
\(\sqrt{4\pi t}\) at zero. The product of the two scalar identities
has the displayed factor \((4\pi t)^{-1}\). All finite Gaussian
integrations against the measures are absolutely convergent. This
proves QLG15 without requiring an equilibrium theorem on the plane.

Let \(c_{x,\varepsilon}\) be probability arc measure on the circle
\(z=x+\varepsilon e^{i\theta}\), where the real centre \(x\) is
arbitrary and \(\varepsilon>0\). Its exact logarithmic potential is
\[
 \int\log|z-w|\,dc_{x,\varepsilon}(z)
           =\log\max\{\varepsilon,|x-w|\}.
 \tag{QLG17}
\]
Indeed, if \(|x-w|>\varepsilon\), expand
\(\log|1-re^{i\theta}|=-\sum_{j\geq1}r^j\cos(j\theta)/j\)
with \(r=\varepsilon/|x-w|<1\); all nonconstant terms average to
zero. If \(|x-w|<\varepsilon\), use the same expansion with
\(r=|x-w|/\varepsilon\). At equality the assertion is the elementary
integral \(\int_0^{2\pi}\log|1-e^{i\theta}|d\theta=0\).
To verify this last integral directly, write
\(|1-e^{i\theta}|=2\sin(\theta/2)\), and put
\(A=\int_0^{\pi/2}\log(\sin u)du\). Reflection gives the same
integral for \(\log\cos u\); the double-angle identity then gives
\(2A=-(\pi/2)\log2+A\), whence \(A=-(\pi/2)\log2\).
The endpoint logarithms are integrable, for example from
\(2u/\pi\leq\sin u\leq u\) on \([0,\pi/2]\).

For finite centres \(x,y\), the logarithmic kernel is absolutely
integrable against \(c_{x,\varepsilon}\otimes c_{y,\varepsilon}\).
Its positive part is bounded by a finite constant depending on the
centres and \(\varepsilon\). Applying QLG17 first to this bounded
positive part minus the negative part, Tonelli's theorem and the finite
value \(\log\max\{\varepsilon,|x-w|\}\) show that the negative
part also has finite double integral. This argument includes identical
and tangent circles. The same argument proves absolute integrability
of the circle--\(\mu\) cross kernel, using compactness of the support
of \(\mu\); the \(\mu\)--\(\mu\) term was proved above.
In particular QLG15 applies to a finite mixture of these circles minus
\(\mu\).

Two consequences of QLG17, retaining both self and cross terms, are
\[
 I(c_{x,\varepsilon})=\log\varepsilon,\qquad
 I(c_{x,\varepsilon},c_{y,\varepsilon})
 \geq\log\max\{\varepsilon,|x-y|\}\geq\log|x-y|.
 \tag{QLG18}
\]
The first equality applies QLG17 at \(w\) on its own circle.
For the second inequality, the first circle average is
\(\log\max\{\varepsilon,|x-w|\}\geq\log|x-w|\);
averaging the right side on the second circle and applying QLG17 again
gives its asserted value. When \(x=y\) the last lower bound is
understood in the extended sense.

\paragraph{QLG4. A finite upper estimate on the original particle domain.}
For fixed particles \(x_1,\ldots,x_n>0\), put
\(\nu_\varepsilon=n^{-1}\sum_i c_{x_i,\varepsilon}\).
By QLG18 and QLG15, respectively,
\[
 2\sum_{i<j}\log|x_i-x_j|
       \leq n^2I(\nu_\varepsilon)-n\log\varepsilon,
 \qquad
 I(\nu_\varepsilon)\leq2I(\nu_\varepsilon,\mu)-I(\mu).
 \tag{QLG19}
\]
There is no potential evaluated at a complex particle in these
identities. The exact comparison back to the original real potential is
\[
 0\leq\int P(z)\,dc_{x,\varepsilon}(z)-P(x)
 =\int_{|t-x|<\varepsilon}
               \log\frac{\varepsilon}{|t-x|}\rho(t)\,dt
 \leq2M_*\varepsilon.
 \tag{QLG20}
\]
The last constant is the explicit integral
\(M_*\int_{-\varepsilon}^{\varepsilon}
\log(\varepsilon/|s|)ds=2M_*\varepsilon\).
Combining QLG19--20 and then substituting QLG4--5 proves, for every
original particle configuration,
\[
 2\sum_{i<j}\log|x_i-x_j|-n\sum_iV_{\alpha,\beta}(x_i)
 \leq n^2F-n\sum_iD(x_i)-n\log\varepsilon
                                      +4M_*n^2\varepsilon.
 \tag{QLG21}
\]
The coefficient \(4M_*\) comes from twice the logarithmic potential
and the bound \(2M_*\varepsilon\); the coefficient of the self energy
is \(-n\log\varepsilon\).
Take \(\varepsilon=1/n\). Since \(D\geq0\) and \(n\geq1\),
\(\int e^{-nD}d\lambda\leq T_\lambda\). Integrating QLG21 gives
the finite positive upper bound
\[
 \log J_n^\lambda
 \leq n^2F+4M_*n+n\log T_\lambda+n\log n-\log(n!).
 \tag{QLG22}
\]
In particular the original integrals in QLG7 are finite. The elementary
inequality
\(\log(n!)\geq\int_1^n\log t\,dt=n\log n-n+1\)
also proves
\[
 \log J_n^\lambda
 \leq n^2F+n(1+4M_*+\log T_\lambda)-1.
 \tag{QLG23}
\]
The inequality for \(n!\) follows by monotonicity of \(\log t\)
on each interval \([j-1,j]\), for \(j=2,\ldots,n\), and is exact
at \(n=1\).

\paragraph{QLG5. Quantile lower estimate and the exact factorial cancellation.}
Because \(\rho\) is continuous and positive on \((a,b)\), there are
unique endpoints
\[
 a=t_0<t_1<\cdots<t_n=b,\quad
 \int_{t_{j-1}}^{t_j}\rho(x)\,dx=\frac1n,\quad
 \Delta_j=t_j-t_{j-1}>0,
 \quad p_j(x)=n\rho(x)\mathbf1_{(t_{j-1},t_j)}(x).
 \tag{QLG24}
\]
The endpoints have measure zero. Each \(p_j(x)dx\) is a probability
measure with bounded density, on a compact interval of positive length.
Every pair logarithmic integral below is therefore absolutely
integrable, including self integrals and integrals over touching
intervals, by the same bounded-density estimate used after QLG11.
Also \(\int p_j|\log(p_j/w_\lambda)|dx<\infty\), since
\(p_j\leq nM_*\) and \(\log w_\lambda\) is bounded on the support.

The product cell \(\prod_{j=1}^n(t_{j-1},t_j)\) and its \(n!\)
coordinate permutations are disjoint. Symmetry of the integrand and the
product integration measure in QLG7 implies that the Heine factor
\(1/n!\) cancels these \(n!\) equal integrals exactly. Restricting
to their union therefore gives
\[
 J_n^\lambda\geq
 \int_{\prod_j(t_{j-1},t_j)}
 \exp\!\left(2\sum_{i<j}\log|x_i-x_j|-n\sum_iV(x_i)\right)
                       \prod_j w_\lambda(x_j)\,dx_j.
 \tag{QLG25}
\]
On this cell use the probability density \(\prod_jp_j(x_j)\).
The logarithm of the integrand relative to this probability is
integrable in absolute value, by the preceding estimates and boundedness
of \(V\) on \([a,b]\). The elementary convex inequality
\(e^u\geq e^v(1+u-v)\), integrated with
\(v\) equal to the expectation of this logarithm, proves Jensen's
lower bound directly. Since \(\sum_jp_j=n\rho\), its pair and
potential contributions are exactly
\[
 \sum_{i\ne j}I(p_i,p_j)-n\sum_j\int Vp_j\,dx
   =n^2F-\sum_j I(p_j).
\]
Here \(I(p_i,p_j)=\iint\log|x-y|p_i(x)p_j(y)dxdy\) and
\(I(p_j)=I(p_j,p_j)\). Its entropy contribution is exactly
\[
 \sum_j\int p_j\log\frac{p_j}{w_\lambda}\,dx
    =n\log n+nH_\lambda(\mu).
 \tag{QLG26}
\]
Moreover \(|x-y|\leq\Delta_j\) on the \(j\)-th cell, whence
\(I(p_j)\leq\log\Delta_j\). Concavity of the logarithm, proved
for example by its negative second derivative, and
\(\sum_j\Delta_j=b-a\) now give
\[
 \sum_jI(p_j)\leq\sum_j\log\Delta_j
       \leq n\log\frac{b-a}{n}.
 \tag{QLG27}
\]
Thus the \(n\log n\) in QLG26 cancels that in QLG27, and the complete
lower bound is
\[
 \boxed{\log J_n^\lambda(\alpha,\beta)
       \geq n^2F(\alpha,\beta)
                -n\{\log(b-a)+H_\lambda(\mu)\}.}
 \tag{QLG28}
\]
This also proves strict positivity of QLG7. No lower bound on the
density at an endpoint or on a quantile interval length has been used.

\paragraph{QLG6. The uniform order-\(n\) theorem.}
For every \((\alpha,\beta)\in\mathcal K\), every integer \(n\geq1\),
and each of the two measures QLG6, QLG22 and QLG28 are the explicit
two-sided finite estimate
\[
 \boxed{
 -n\{\log(b-a)+H_\lambda(\mu)\}
 \leq\log J_n^\lambda-n^2F
 \leq 4M_*n+n\log T_\lambda+n\log n-\log(n!).}
 \tag{QLG29}
\]
For entirely parameter-uniform constants set
\[
 L_0=\log(B_*M_*),\qquad
 L_\omega=\log(B_*M_*)+\log2+\tfrac12\log B_*+\sqrt{B_*},
\]
\[
 C_0=\max\{0,L_0,1+4M_*+\log^+T_*\},\qquad
 C_\omega=\max\{0,L_\omega,1+4M_*\}.
 \tag{QLG30}
\]
These constants are finite by QLG8--14. Since \(b-a\leq B_*\),
QLG11 bounds the lower side of QLG29 by \(-nL_\lambda\); QLG12--14
and QLG23 bound its upper side by \(nC_\lambda\). Consequently
\[
 \boxed{\left|\log J_n^{\lambda_0}(\alpha,\beta)
                 -n^2F(\alpha,\beta)\right|\leq C_0n,
 \qquad
 \left|\log J_n^{\lambda_\omega}(\alpha,\beta)
                 -n^2F(\alpha,\beta)\right|\leq C_\omega n.}
 \tag{QLG31}
\]
For the integral \(Z_n^\lambda\) without the Heine factor, the exact
relation \(\log Z_n^\lambda=\log J_n^\lambda+\log(n!)\) gives
\[
 \left|\log Z_n^\lambda-n^2F-\log(n!)\right|\leq C_\lambda n,
 \qquad
 \left|\log Z_n^\lambda-n^2F\right|
                         \leq n\log n+C_\lambda n.
 \tag{QLG32}
\]
The theorem therefore retains the distinction between the original
Heine determinant integral and the integral with labelled particles.
The gain to order \(n\) in QLG31 follows from the proved circle
comparison and the quantile-cell cancellation, with the same exact
potential, positive-half-line particle domain, and equilibrium of
EIQ1--33 throughout.

% END CURRENT QLG

\clearpage
% BEGIN CURRENT OFV
% Uses complete proved providers OSP1--35, EIQ1--33, and QLG1--32.
% No leading law is differentiated without its finite partition remainder.
\section{The leading first variation of the actual outer-strip quotient}
\label{sec:OFV}

Retain every original object and both original source orders of
OSP1--35 and CTV1--22. In addition to the explicit finite conditions
OSP23, take $k\ge257$. This is an eventual threshold on the same
fixed original parameters. Put
\[
\vartheta=\pi/2,\qquad
\ell_s=(s-1)/4,\qquad
\beta_s=\frac{(2\pi)^{s/2}}{\sqrt2\,\Gamma(s/2)}
\quad(s=1,k).
\tag{OFV1}
\]
The symbol $\vartheta$ is the fixed Gamma tail rate; it is not either
original source centre.

We use the equilibrium of EIQ1--33 on its original positive half-line:
\[
\begin{gathered}
V_{\alpha,\beta}(x)=\beta\sqrt{x}-\alpha\log x,\quad
\mu_{\alpha,\beta}(dx)=\rho_{\alpha,\beta}(x)\,dx,\\
F(\alpha,\beta)=
\iint\log|x-y|\,d\mu_{\alpha,\beta}(x)d\mu_{\alpha,\beta}(y)
-\int V_{\alpha,\beta}\,d\mu_{\alpha,\beta}.
\end{gathered}
\tag{OFV2}
\]
Here $\rho$, its two endpoints, the Euler equality and the entire
exterior inequality are exactly the explicit proved EIQ formulas.
There is no new equilibrium assumption. EIQ29--30 proves smoothness
and
$F_\alpha(\alpha,\beta)=\int\log x\,d\mu_{\alpha,\beta}(x)$.
QLG31, with its complete circle and quantile proofs, gives
\[
\left|\log J_n(\alpha,\beta)-n^2F(\alpha,\beta)\right|\le C_0n,
\quad
J_n=\frac1{n!}\int_{(0,\infty)^n}
\prod_{i<j}(x_i-x_j)^2
\prod_i e^{-nV_{\alpha,\beta}(x_i)}\,dx_i
\tag{OFV3}
\]
for every $n\ge1$ and every
$(\alpha,\beta)\in[3/2,5/2]\times[\pi/2,3\pi/2]$.
The constant $C_0$ is the explicit finite constant QLG30.
We use precisely this uniform finite error, including the factor
$1/n!$; the conclusion below is not a derivative of an earlier
unquantified $q^2$ limit.

\subsection{Exact original-order recurrence and a small metric error}

For $a\ge0$, define
\[
\mathcal D_a(Q;s)
=\det\left[\int_{\mathbb R}y^{2Q+i+j}\,dm_s(y)\right]_{0\le i,j<a},
\qquad
\mathcal A_a(Q;\vartheta)
=\det\left[\int_{\mathbb R}y^{2Q+i+j}e^{-\vartheta|y|}\,dy
                                      \right]_{0\le i,j<a},
\tag{OFV4}
\]
with empty determinants one. The $Q$ used here is an integer;
all the instances below have $Q\ge q/2$.
The original recurrence OSP12 gives exactly
\[
|y|^{2Q}dm_s(y)=
\beta_s |y|^{2Q}
\prod_{j=0}^{\ell_s-1}(y^2+t_j^2)\,d\sigma(y),
\qquad t_j=2j+\tfrac12.
\tag{OFV5}
\]
For each actual source centre $c_0\in\{k/2,k/2-4\}$, the
coordinate polynomial of this multiplier is
$f_{\ell_s,c_0}(S)=\prod_j(S-c_0-t_j)$.
Its exact pullback satisfies
$f_{\ell_s,c_0}(c_0+iy)=\prod_j(iy-t_j)
=i^{\ell_s}\prod_j(y+it_j)$.
Thus multiplication by
$y^Q\prod_{j=0}^{\ell_s-1}(y+it_j)$, together with the
retained phase $i^{\ell_s}$, is the exact coefficient map
into the original one-factor Gamma source. Its squared modulus
is the displayed density multiplier, with no change to any $t_j$
or to $\beta_s$.

Set $T=\epsilon q$, $\epsilon=2^{-32}$ as in OSP23, and
\[
S_{\ell_s}=\sum_{j=0}^{\ell_s-1}t_j^2
=\frac{\ell_s(16\ell_s^2-12\ell_s-1)}{12},\qquad
r_{\ell_s}=\frac{S_{\ell_s}}{T^2}+\log(1+K_q),
\tag{OFV6}
\]
where $K_q$ is exactly OSP24. At $\ell_s=0$ we may and do take
$r_0=0$ instead, since the source identity is exact.
For $Q\in\{q',q\}$, $0\le a\le2q+1$,
\[
0\le
\log\mathcal D_a(Q;s)-a\log\beta_s
-\log\mathcal D_a(Q+\ell_s;1)
\le ar_{\ell_s}.
\tag{OFV7}
\]

Here is a proof of the full-form comparison underlying OFV7.
For every real $y$,
$|y|^{2Q}\prod_j(y^2+t_j^2)\ge |y|^{2(Q+\ell_s)}$.
On $|y|\ge T$, the logarithm of the ratio is at most
$\sum_jt_j^2/T^2=S_{\ell_s}/T^2$, by $\log(1+u)\le u$.
For $|y|<T$, its numerator is at most
$(T^2+(k/2)^2)^{Q+\ell_s}$.
OSP15 on the complete coefficient space of degree at most $a-1$,
OSP13 with source $s=1$, and mass $M_1=\sqrt{2\pi}$ therefore
bound its inner quadratic form by $K_q$ times the reference
quadratic form on $[q,2q]$. This is exactly the calculation
OSP22--25 with exponent $Q+\ell_s\ge q/2$, degree $a-1\le2q$,
and root bound $k/2\le R\le T/2$.
No upper bound on that exponent is needed in OSP25, since its
base $(5/4)\epsilon^2$ is less than one.
The full upper ratio is consequently at most
$e^{S_{\ell_s}/T^2}+K_q\le e^{r_{\ell_s}}$.
The lower full ratio is at least one by the pointwise lower
inequality. The positive-Gram eigenvalue argument OSP28 now proves
OFV7 in the original coefficient frames.
Since $S_{\ell_s}\le\ell_s k^2/4$ and $q\ge k^2$,
$r_{\ell_s}=O(1/k)$, with its stated fixed cutoff constant.

\subsection{A uniform finite Gamma-to-exponential determinant estimate}

Put $h=1/q$, and define the actual finite constants
\[
L_h=C_L(\sin h)^{3/2},\qquad
U_h=\frac{\Gamma(1/4)^2}{2\pi}(\sin h)^{-1/2}.
\tag{OFV8}
\]
Then the original one-factor Gamma density satisfies on the entire
real line
\[
L_he^{-(\vartheta+h)|y|}
\le\frac{d\sigma}{dy}(y)
\le U_he^{-(\vartheta-h)|y|}.
\tag{OFV9}
\]
Indeed apply OSP7 to $\Gamma(1/4+iy/2)$ with
$\theta=\pi/2-h$ for the upper bound.
For the lower bound apply it to $\Gamma(3/4-iy/2)$ in the exact
reflection product proved before OSP11. The upper bound on that
squared factor is
$\Gamma(3/4)^2(\sin h)^{-3/2}e^{-(\vartheta-h)|y|}$.
After division, the identity and
$\cosh(\pi y)\le e^{\pi|y|}$ give exactly the lower side of OFV9.
In particular both inequalities hold near zero as well as in
the tails.

The exact change of variables $y=\vartheta^{-1}u$ in OFV4 gives
\[
\mathcal A_a(Q;\vartheta)
=\vartheta^{-a(2Q+a)}\mathcal A_a(Q;1).
\tag{OFV10}
\]
The exponent follows from $a$ moment Jacobians, $2Q$ powers in
each, and the row and column powers
$\sum_{i=0}^{a-1}i+\sum_{j=0}^{a-1}j=a(a-1)$.
Applying OFV9 as full quadratic-form inequalities, and then using
OFV10, proves the separate finite bounds
\[
\begin{split}
a\log L_h-a(2Q+a)\log(1+h/\vartheta)
&\le\log\mathcal D_a(Q;1)-\log\mathcal A_a(Q;\vartheta)\\
&\le a\log U_h-a(2Q+a)\log(1-h/\vartheta).
\end{split}\tag{OFV11}
\]
Thus the absolute difference is at most
\[
B_a(Q)=a\max\{|\log L_h|,|\log U_h|\}
       +a(2Q+a)\frac{h}{\vartheta-h}.
\tag{OFV12}
\]
Here $-\log(1-u)\le u/(1-u)$ and $\log(1+u)\le u$
justify the second summand. Since
$2h/\pi\le\sin h\le h$ for $0<h\le\pi/2$,
the first summand is $O(a\log q)$.
The sine inequalities follow from concavity and the endpoint chord,
and by integrating $\cos t\le1$.
For all the present $a\le2q+1$, $Q\le q+\ell_s\le2q$,
OFV12 is $O(q\log q)$ uniformly in $a,Q$.
Combining OFV7 and OFV11 gives the error-controlled identity
\[
\log\mathcal D_a(Q;s)
=a\log\beta_s+\log\mathcal A_a(Q+\ell_s;\vartheta)
  +R_{a,Q,s},\qquad
|R_{a,Q,s}|\le ar_{\ell_s}+B_a(Q+\ell_s).
\tag{OFV13}
\]
This is a comparison on complete original polynomial coefficient
spaces, with the order shift proved by the literal multiplier
OFV5; it is not an assigned replacement of the original source order.

\subsection{Direct low-dimension estimates, outside the compact equilibrium range}

Define
\[
g(Q)=2Q\log(2Q/\vartheta)-2Q.
\tag{OFV14}
\]
For $Q\ge q/2$ and $1\le a\le\Delta+1$ we prove
\[
\log\mathcal A_a(Q;\vartheta)
=a g(Q)+O\bigl(a^2\log(Q+a)+a\log(Q+a)\bigr),
\tag{OFV15}
\]
with constants independent of $k,Q,a$ in these ranges.
There is no use of OFV3 here; its compact parameter rectangle
does not contain these low-dimension ratios.

For the lower bound take the full interval
$[y_0,y_0+1]$, $y_0=2Q/\vartheta$. On it
$y^{2Q}e^{-\vartheta y}\ge e^{g(Q)-\vartheta}$.
Consequently its moment Gram is bounded below by
$e^{g(Q)-\vartheta}$ times the unweighted monomial Gram on that interval.
Translation has a triangular coefficient matrix with diagonal one,
so the latter Gram determinant equals the one on $[0,1]$.
The Rodrigues norms proved in OSP15 show that it is exactly
\[
\prod_{j=0}^{a-1}
\frac1{(2j+1)\binom{2j}{j}^{\,2}}.
\]
Indeed $\binom{2j}{j}$ is the leading coefficient of
$L_j(2x-1)$, whose squared norm is $1/(2j+1)$.
Using $\binom{2j}{j}\le4^j$ therefore proves
\[
\log\mathcal A_a(Q;\vartheta)
\ge a g(Q)-a\vartheta-a\log(2a)-2a(a-1)\log2.
\tag{OFV16}
\]

For the upper bound, Gram--Schmidt shows that a positive Gram
determinant is at most the product of the squared norms of its
displayed vectors: each orthogonal component has norm no greater
than its original vector. The exact monomial moments give
\[
\log\mathcal A_a(Q;\vartheta)
\le\sum_{j=0}^{a-1}
\log\left(\frac{2(2(Q+j))!}{\vartheta^{2(Q+j)+1}}\right).
\tag{OFV17}
\]
For every positive integer $m$, integrating the increasing function
$\log x$ between successive integers proves
$m\log m-m+1\le\log(m!)\le m\log m-m+1+\log m$.
Each summand in OFV17 is consequently at most
$g(Q+j)+|\log(2/\vartheta)+1|+\log(2(Q+j))$.
Furthermore $g'(u)=2\log(2u/\vartheta)$, so
$g(Q+j)-g(Q)\le2j\log(2(Q+a)/\vartheta)$ in the present range.
This proves the explicit upper bound
\[
\begin{split}
\log\mathcal A_a(Q;\vartheta)\le{}&
a g(Q)+a\{|\log(2/\vartheta)+1|+\log(2(Q+a))\}\\
&+a(a-1)\log(2(Q+a)/\vartheta).
\end{split}\tag{OFV18}
\]
OFV16 and OFV18 prove OFV15 with its complete lower and upper
errors. They also apply at $a=1$.

\subsection{High-dimension partitions and their controlled first variation}

For an integer $a$ write $n_e=\lceil a/2\rceil$ and
$n_o=\lfloor a/2\rfloor$. The two exact changes of variables
$y=\sqrt{x}$ and $y=-\sqrt{x}$ give, with their two positive
Jacobians, the parity identity
\[
\mathcal A_a(Q;\vartheta)
=Z_{n_e}(Q-\tfrac12;\vartheta)\,
 Z_{n_o}(Q+\tfrac12;\vartheta),
\tag{OFV19}
\]
where
\[
Z_n(A;\vartheta)
=\det\left[\int_0^\infty x^{A+i+j}e^{-\vartheta\sqrt{x}}\,dx
                                      \right]_{0\le i,j<n}.
\]
The two branches add to the density $x^{-1/2}e^{-\vartheta\sqrt{x}}$;
thus no factor two is missing from OFV19. The odd block has the
additional factor $x$. Expanding the two Vandermonde determinants
and integrating each finite permutation sum proves the exact
Heine identity for this moment determinant:
\[
Z_n(A;\vartheta)=\frac1{n!}\int_{(0,\infty)^n}
\prod_{i<j}(x_i-x_j)^2\prod_i x_i^Ae^{-\vartheta\sqrt{x_i}}\,dx_i.
\]
Changing variables $x_i=n^2u_i$ now gives
\[
Z_n(A;\vartheta)=n^{2n(A+n)}J_n(A/n,\vartheta).
\tag{OFV20}
\]
There are $2n(A+1)$ powers from the weights and Jacobians
and $2n(n-1)$ from the squared Vandermonde, proving the exponent.

For the high instances we use exactly
\[
(a,Q)\in
\{(q,q+\ell_s),(q+1,q+\ell_s),
  (q+\Delta,q-\Delta+\ell_s),
  (q+\Delta+1,q-\Delta+\ell_s)\}.
\tag{OFV21}
\]
Since $k\ge257$, $\Delta/q\le16/k<1/16$ and
$\ell_s/q\le1/(4k)<1/64$.
For both parity blocks, elementary rounding therefore gives
\[
\frac{15q}{32}\le n\le\frac{9q}{16},\qquad
\frac{29q}{32}\le A\le\frac{33q}{32}.
\]
The first inequalities use $q\ge32$ to absorb the possible
rounding by $1/2$ and the extra endpoint $1$.
Hence $29/18\le A/n\le11/5$, contained strictly in $[3/2,5/2]$.
Every partition in OFV20 is thus in the proved parameter rectangle
OFV3, with $\beta=\vartheta=\pi/2$.

Define the smooth explicit function
\[
\mathcal H(a,Q)=
a(2Q+a)\log(a/2)+\frac{a^2}{2}F(2Q/a,\vartheta)
\quad(a>0,\ Q>0).
\tag{OFV22}
\]
Equations OFV3 and OFV19--20 give uniformly over OFV21
\[
\log\mathcal A_a(Q;\vartheta)=\mathcal H(a,Q)+O(q\log q).
\tag{OFV23}
\]
To verify the rounding error rather than omit it, set
$G(n,A)=2n(A+n)\log n+n^2F(A/n,\vartheta)$.
Each actual pair $(n,A)$ differs from $(a/2,Q)$ by at most $1/2$
in each coordinate. On all connecting segments $n$ is bounded
above and below by positive constant multiples of $q$, while
$A/n$ remains in $[3/2,5/2]$.
Its derivatives are
\[
G_A=2n\log n+nF_\alpha,\qquad
G_n=2(A+2n)\log n+2(A+n)+2nF-AF_\alpha .
\]
They are $O(q\log q)$ there, since the exact EIQ functions
$F,F_\alpha$ are bounded on that compact interval. The mean value
formula bounds each rounding error by $O(q\log q)$.
The two QLG errors sum to at most $C_0a=O(q)$, and
$2G(a/2,Q)=\mathcal H(a,Q)$, proving OFV23.

The derivative needed for the first variation is computed without
altering the direction of the two original degree changes:
\[
(\partial_a-\partial_Q)\mathcal H(a,Q)
=2Q\log(a/2)+(2Q+a)+aF(2Q/a,\vartheta)
 -(Q+a)F_\alpha(2Q/a,\vartheta).
\tag{OFV24}
\]
EIQ29--30 proves that $F_\alpha$ is smooth; in particular
$F,F_\alpha,F_{\alpha\alpha}$ are bounded on the compact interval
just used. Differentiating the right side of OFV24 shows that
its two first partial derivatives are $O(\log q)$ on the segment
$a=q+u$, $Q=q+\ell_s-u$, $0\le u\le\Delta$ and on its
connections to $(q,q)$. Every term follows by the product rule;
the derivatives of $2Q/a$ are $2/a$ and $-2Q/a^2$, both $O(1/q)$.
The mean value formula and then integration in $u$ consequently give
\[
\begin{split}
&\mathcal H(q+\Delta,q-\Delta+\ell_s)
             -\mathcal H(q,q+\ell_s)\\
&\quad=
\Delta q\{2\log(q/2)+3+F(2,\vartheta)-2F_\alpha(2,\vartheta)\}
+O\bigl(\Delta(\Delta+\ell_s)\log q\bigr).
\end{split}\tag{OFV25}
\]
This remainder is $O(q\log q)$ because
$\Delta=O(k)$, $\ell_s=O(k)$ and $q=(k+1)^2$.
The additional endpoint $a+1$ in OFV21 changes $\mathcal H$
by $O(q\log q)$, directly from its displayed derivatives.
Thus all rounding and endpoint changes have been quantified.

\subsection{All seven determinants and the resulting leading constant}

Replace the seven moment determinants in OSP32 by the exact
exponential determinants in OFV13, with their respective
$Q+\ell_s$. Their $a\log\beta_s$ terms cancel because their
signed dimensions are exactly
\[
\Delta+(\Delta+1)-(q+\Delta)-(q+\Delta+1)-1+q+(q+1)=0.
\tag{OFV26}
\]
The sum of their absolute finite errors in OFV13 is
$O(q\log q)$, since the sum of their dimensions is
$4(q+\Delta+1)$, $r_{\ell_s}=O(1/k)$, and OFV12 applies
uniformly with $Q+\ell_s\le2q$.
The low terms, using OFV15--18, are
\[
\begin{split}
&\log\mathcal A_\Delta(q-\Delta+\ell_s;\vartheta)
+\log\mathcal A_{\Delta+1}(q-\Delta+\ell_s;\vartheta)
-\log\mathcal A_1(q+\ell_s;\vartheta)\\
&\qquad=2\Delta g(q)+O(q\log q)
=4\Delta q\{\log(2q/\vartheta)-1\}+O(q\log q).
\end{split}\tag{OFV27}
\]
Indeed replacing each $g(q-\Delta+\ell_s)$ or $g(q+\ell_s)$
by $g(q)$ costs at most
$O(((2\Delta+1)(\Delta+\ell_s)+\ell_s)\log q)=O(q\log q)$,
using $g'(u)=2\log(2u/\vartheta)$ and the actual degree ranges.
The individual low determinant errors are
$O((\Delta+1)^2\log q)=O(q\log q)$.

The remaining four high terms are, by OFV23--25,
\[
\begin{split}
&\log\mathcal A_q(q+\ell_s;\vartheta)
+\log\mathcal A_{q+1}(q+\ell_s;\vartheta)\\
&\quad-\log\mathcal A_{q+\Delta}(q-\Delta+\ell_s;\vartheta)
-\log\mathcal A_{q+\Delta+1}(q-\Delta+\ell_s;\vartheta)\\
&=-2\Delta q
\{2\log(q/2)+3+F(2,\vartheta)-2F_\alpha(2,\vartheta)\}
+O(q\log q).
\end{split}\tag{OFV28}
\]
Combining OFV27 and OFV28 gives the complete constant
\[
\mathfrak C_{\partial}
=4\log(4/\vartheta)-10-2F(2,\vartheta)+4F_\alpha(2,\vartheta),
\qquad \vartheta=\pi/2,
\tag{OFV29}
\]
with every logarithmic scale term canceled by direct algebra.
OSP33 now proves the actual original-root asymptotic
\[
\boxed{\displaystyle
\mathcal F_{\partial,k}^{(s)}
=\Delta q\,\mathfrak C_{\partial}
   +O_{\delta,\gamma}(q\log q),\qquad
\frac{\mathcal F_{\partial,k}^{(s)}}{kq}
\longrightarrow16\,\mathfrak C_{\partial},
\qquad s=1,k.}
\tag{OFV30}
\]
All constants in the error depend only on the fixed original
parameters, the fixed cutoff, and the proved compact equilibrium
constants. The two original source orders have been treated
through their exact recurrence maps and full mass factors.
Since $\Delta/k=16-48/k$ and $\log q/k\to0$, the displayed
limit follows from its quantified remainder.

An equivalent expression uses the logarithmic energy of precisely
the original EIQ limiting potential, $V(x)=\pi\sqrt{x}-2\log x$:
\[
I_{2,\pi}=\iint\log|x-y|\,d\mu_{2,\pi}(x)d\mu_{2,\pi}(y),
\qquad
\mathfrak C_{\partial}=4\log(4/\pi)+2-2I_{2,\pi}.
\tag{OFV31}
\]
To prove it, EIQ31 gives
$\vartheta\int\sqrt{x}\,d\mu_{2,\vartheta}=6$, so
$F(2,\vartheta)=I_{2,\vartheta}-6+2F_\alpha(2,\vartheta)$.
Substituting in OFV29 yields
$4\log(4/\vartheta)+2-2I_{2,\vartheta}$.
The proved EIQ32 dilation $x\mapsto4x$ carries $\mu_{2,\pi}$
to $\mu_{2,\pi/2}$. Its exact logarithmic energy change is
$I_{2,\pi/2}=I_{2,\pi}+\log4$, proving OFV31.
This energy is a finite explicit integral of the fully specified
EIQ density; its logarithmic singularity is integrable because
that density is bounded and compactly supported.

Finally, the exact quotient metrics decrease when the source degree
increases: each smaller-degree lift remains an allowed lift in the
larger source with exactly the same norm. Hence each low-degree
quotient Gram in CTV13 dominates the corresponding high-degree
Gram on the identical target frame. Positive conjugation shows
their determinant ratios are at least one. It follows that
$\mathcal F_{\partial,k}^{(s)}\ge0$, and OFV30 therefore also proves
$\mathfrak C_{\partial}\ge0$.
The already proved finite tail bounds CTV18--20, with
$0\le v\le80$, further give
\[
\mathcal T_k^{(s)}
=\Delta q\,\mathfrak C_{\partial}
+O_{\delta,\gamma}(q\log q),\qquad s=1,k,
\tag{OFV32}
\]
uniformly over the possible actual first-moment indices.
No leading value of the separate original observation quotient
is assigned here.

% END CURRENT OFV

\clearpage
% BEGIN CURRENT FMC
% FMC: a fixed original module presentation and controlled coefficient frames.
\section{Fixed coefficient control of the original residual module}
\label{sec:FMC}

Retain the complete actual simple-quartet period matrix
$\mathsf H=F^{-1}\Pi U\mathcal E_{\rm prim}$, the original substitution
$t=(z_0w_0,z_0w_1,z_1w_0,z_1w_1)^{\mathsf T}$, and
$Q(x)=Q_0(\mathsf H^{-1}x)$, where
$Q_0(t)=t_{++}t_{--}-t_{+-}t_{-+}$.
The original action is $g f(x)=f(D^{-1}x)$ with
$D=\operatorname{diag}(\zeta,\zeta^2,\zeta^3,\zeta^4)$ and
$\zeta=e^{2\pi i/5}$. At the fixed actual period on the proved
five-member orbit locus retain the entire polynomials
$Q_a=g^aQ$ and $\mathcal A=\prod_{a=1}^4Q_a$.
Their coefficients and all phases are unchanged.
The degree-$k$ calculations below use $k=4l+1\ge9$ and
\[
 q=(k+1)^2,\quad q'=(k-7)^2,\quad
 \Delta=16k-48,\quad j=8k-32,\quad r=8k-16.
 \tag{FMC1}
\]
All constants introduced below are determined by this one fixed
actual coefficient datum. They are independent of $k$ and of the
original Gamma cutoff. A fixed presentation will supply the residual
frame in place of dividing new selected minors at each degree.

\subsection{A fixed free module containing the entire invariant image}

Let $P=\mathbb C[x_1,x_2,x_3,x_4]$ and introduce the specified
graded ring injection
\[
 T=\mathbb C[y_1,y_2,y_3,y_4]\longrightarrow P,
 \qquad y_i\longmapsto x_i^5,\qquad \deg y_i=5.
 \tag{FMC2}
\]
It is injective because distinct monomials have distinct exponent
vectors after multiplying each exponent by five. Put
$\mathcal E=\{a\in\mathbb Z^4:0\le a_i\le4\}$ and let $e_a$
be a free $T$-basis element of degree $|a|$. The exact isomorphism is
\[
 \mathscr P:=\bigoplus_{a\in\mathcal E}T e_a
       \xrightarrow{\ \iota\ }P,
 \qquad y^\nu e_a\longmapsto x^{5\nu+a}.
 \tag{FMC3}
\]
Every nonnegative exponent vector has a unique coordinatewise quotient
and remainder on division by five. This proves surjectivity, injectivity,
preservation of degree, and the explicit inverse. In particular this
is a rank-$625$ free module, with coefficient $\ell^1$ norm preserved
exactly by $\iota$. The ring and monomial coefficients in $P$ have
not been rescaled.

Let $\mathcal E_0=\{a\in\mathcal E:
a_1+2a_2+3a_3+4a_4\equiv0\pmod5\}$. The full original invariant
algebra, regarded here as a $T$-module, is
\[
 \mathscr I=P^{\langle g\rangle}
       =\iota\left(\bigoplus_{a\in\mathcal E_0}T e_a\right),
 \qquad |\mathcal E_0|=125.
 \tag{FMC4}
\]
Indeed the monomial $x^{5\nu+a}$ has character
$\zeta^{-(a_1+2a_2+3a_3+4a_4)}$. The diagonal action fixes a
polynomial precisely when each nonzero monomial has character one.
For each of the $5^3$ choices of the first three residues there is
exactly one $a_4$ since $4$ is invertible modulo five. This proves
the count and the displayed equality. These pure fifth powers are
four of the fourteen original generators OCF25; this is a new exact
module presentation of the same entire invariant algebra.

Define two homogeneous $T$-submodules of $\mathscr P$ by
\[
 \begin{split}
 \mathscr H&=\iota^{-1}(QP+\mathcal A P),\\
 \mathscr S&=\mathscr H+
                       \bigoplus_{a\in\mathcal E_0}T e_a.
 \end{split}\tag{FMC5}
\]
They have the fixed finite generating columns
$\iota^{-1}(Qx^a)$ and $\iota^{-1}(\mathcal A x^a)$ for
$a\in\mathcal E$, followed for $\mathscr S$ by the $125$ columns
$e_a$, $a\in\mathcal E_0$. The first two lists have respectively
homogeneous degrees $|a|+2$ and $|a|+8$, all at most $24$.
To obtain every matrix entry explicitly, expand either of the full
original polynomials as $f(x)=\sum_b f_bx^b$ and use
\[
 \iota^{-1}(f x^a)=
       \sum_b f_b\,y^{\lfloor(a+b)/5\rfloor}
                         e_{(a+b)\bmod5}.
 \tag{FMC6}
\]
Both quotient and remainder are coordinatewise. Formula FMC3 verifies
this identity term by term. Since $P$ is free on the $x^a$, these
columns generate precisely the two modules in FMC5; no later degree
is needed to define a new relation matrix.

The full coefficient substitution and conductor quotient give the exact
graded maps
\[
 \mathscr P/\mathscr H\xrightarrow{\ \iota,\,\psi\ }
          \mathcal B/(A\mathcal B)=:\mathcal L,
 \qquad
 \mathscr S/\mathscr H\xrightarrow{\ \sim\ }\overline{\mathscr J},
 \quad A=\psi_8(\mathcal A).
 \tag{FMC7}
\]
Here $\psi f=f(\mathsf Ht)$, and $\overline{\mathscr J}$ is the
image of the original invariant algebra in $\mathcal L$, exactly
OCF16. OCF3 proves $\ker\psi=QP$ and its surjectivity in every
degree. If $\psi f=A b$, lift $b$ to a polynomial $h$ using that
surjectivity; then $f-\mathcal A h\in QP$. Conversely every
$QP+\mathcal AP$ element maps to zero. This proves the first
isomorphism. FMC4--5 show that its restriction has precisely the
second image and kernel. Thus the quotient in FMC7 is the complete
original conductor quotient, and its indicated submodule is the
actual residual certificate image of ROQ4, with no new invariant
premise.

\subsection{Finite reduction data with only fixed actual denominators}

The next construction specifies fixed data for both modules in FMC5.
Order their module monomials $y^\nu e_a$ first lexicographically by
$\nu_1,\nu_2,\nu_3,\nu_4$, and then by a fixed index
$\kappa_a\in\{0,\ldots,624\}$ listing $a$ lexicographically.
Multiplication by a $y$-monomial preserves this order. It is a
well-order: in a decreasing infinite sequence the first coordinate
can decrease only finitely often; after it is fixed the same applies
successively to the other three coordinates and then the finite index.
For a nonzero column $f$, retain its full leading coefficient
$\lambda_f\ne0$ and its leading module monomial $m_f$.
No leading coefficient is deleted from the stored column.

Start with the finite nonzero homogeneous generators of FMC5--6.
Division of a column repeatedly cancels its largest divisible module
monomial $c y^\nu e_a$ using the first available column $f$ for which
$m_f=y^\mu e_a$ divides it: subtract
$(c/\lambda_f)y^{\nu-\mu}f$ and record that exact multiple.
A largest term divisible by none of the current leading monomials is
put into the remainder and removed from the current dividend.
Each operation strictly decreases the largest term still being
processed, so this division terminates. Remainders have no monomial
divisible by the stored leading monomials.

For every pair $f,h$ whose leading positions agree, put
$m=\operatorname{lcm}(m_f,m_h)$ in that position and form the full
column
\[
 \operatorname{Sp}(f,h)=
      \frac{m/m_f}{\lambda_f}f
                  -\frac{m/m_h}{\lambda_h}h.
 \tag{FMC8}
\]
Divide it as just described and append its nonzero remainder to the
list. Reconsider all new pairs. Track every appended column as a
$T$-linear combination of the original generators throughout these
operations. Every column and recorded certificate remains homogeneous:
the two terms in FMC8 have the same weighted degree, because their
leading position agrees and every input column is homogeneous in the
shifted grading of FMC3.

This algorithm stops after finitely many appended columns. Here is the
needed finiteness argument. Every infinite sequence in $\mathbb N^4$
has two terms in their original order with all coordinates weakly
increasing. To prove it, an infinite sequence of nonnegative integers
has either an infinitely repeated value or an increasing infinite
subsequence; use this successively for all four coordinates. For
module monomials first take a subsequence with one common position.
Consequently every infinite sequence of module monomials has an earlier
one dividing a later one. Each appended remainder's leading monomial
is divisible by none of the previous leading monomials. An infinite
sequence of appendages would contradict the proved property.

At termination every element of the generated module has leading
monomial divisible by one of the stored leading monomials. For
completeness, write an element as a finite sum of monomial multiples
of the stored columns. Among all such representations choose one for
which the largest leading module monomial of a summand is smallest,
and among these choose one with the fewest summands at that monomial.
These choices exist by the well-order and the positive integers.
If cancellation of those largest terms occurs, two summands at the
same module position can be combined using a multiple of FMC8,
with its exact two leading coefficients. Explicitly, if their
common leading monomial is $M$ and their leading coefficients are
$a,b$, their sum is
\[
 a\frac{M}{\operatorname{lcm}(m_f,m_h)}\operatorname{Sp}(f,h)
       +(a+b)\frac{M/m_h}{\lambda_h}h.
\]
This follows by substituting FMC8; it retains exactly both original
summands and every coefficient. The recorded division of
this S-column, together with its appended remainder when one was
needed, expresses it in stored columns with all
leading multiples strictly below its canceled common monomial.
Replacing the two summands by this expression either decreases the
largest monomial or decreases the number attaining it; combining
their coefficients if necessary gives the latter alternative.
Both contradict the chosen representation. Thus the largest
remaining monomial is an actual leading monomial of the element,
and it is divisible by one of the stored ones, as claimed.
This also proves that division of a module element has zero remainder.

Write $\mathcal G_H$ and $\mathcal G_S$ for the two resulting fixed
finite lists for FMC5. The preceding argument proves their full
reduction property. A column supported outside the respective leading
module is called a standard remainder. Such a remainder represents
each quotient class uniquely: two representatives of one class differ
by a module element, and a nonzero standard difference would violate
the proved leading-monomial property. Denote these exact linear
remainder maps by $\operatorname{NF}_H$ and $\operatorname{NF}_S$.
Their linearity follows from uniqueness and linearity of quotienting.

All data constructed here depend only on the actual matrices in FMC6.
Every operation is an addition, multiplication, or division by one of
the nonzero actual leading coefficients encountered in this finite
calculation. Thus the complete list of denominators is finite and
independent of $k$. Its precise coefficient chart records each tested
coefficient's vanishing or nonvanishing and every selected leading
monomial. For every actual period exactly its corresponding decisions
give the above terminating calculation. This does not posit a
numerical test for equality of unevaluated periods, nor does it assert
a lower bound for a later sequence of minors. The exact stored
coefficients and their nonzero leading entries define the constants
that follow.

\subsection{A fixed weight bounds every homogeneous reduction path}

The finite reductions admit a linear-in-degree depth bound. Let
$D_0\ge1$ be the largest of one and all absolute differences between
corresponding $y$-exponents in a leading and a nonleading term of
any column in $\mathcal G_H\cup\mathcal G_S$. Empty tail lists
contribute zero. Set
\[
 B_0=D_0+626,\qquad
 \mathfrak w(y^\nu e_a)
   =B_0^4\nu_1+B_0^3\nu_2+B_0^2\nu_3+B_0\nu_4+\kappa_a,
 \qquad L_k=B_0^4\lfloor k/5\rfloor+624.
 \tag{FMC9}
\]
Every nonleading term of either fixed list has weight at least one
less than its leading term. If the $y$-exponents agree, this follows
from the position order. Otherwise let their first differing
coordinate have weight $B_0^m$, $1\le m\le4$. Its exponent drops by
at least one, while the later differences have absolute value at
most $D_0$, and the position difference has absolute value at most
$624$. The total weight drop is at least
\[
 B_0^m-D_0\sum_{h=1}^{m-1}B_0^h-624
 \;\ge\;\frac{625B_0^m}{D_0+625}-624\;>\;1.
 \tag{FMC10}
\]
The middle inequality follows by bounding the geometric sum by
$B_0^m/(B_0-1)$; the last follows from $B_0^m\ge B_0=D_0+626$.
Multiplication by a monomial adds the same weight to the leading
and all its replacement terms, so the same drop holds in every
division step. A homogeneous degree-$k$ monomial has
$5|\nu|+|a|=k$, hence weight between zero and $L_k$. Every branch
of its reduction tree consequently has length at most $L_k$.

For each of the two lists define the complete coefficient constant
\[
 C_H=\max\left\{1,\max_{g\in\mathcal G_H}
                 \sum_{m\ne m_g}|[m]g/\lambda_g|\right\},
 \qquad
 C_S=\max\left\{1,\max_{g\in\mathcal G_S}
                 \sum_{m\ne m_g}|[m]g/\lambda_g|\right\}.
 \tag{FMC11}
\]
The maximum over an empty list is zero. They are finite actual
numbers because the lists are finite and their displayed leading
coefficients are nonzero. For every homogeneous degree-$k$ column,
\[
 \|\operatorname{NF}_H f\|_1\le C_H^{L_k}\|f\|_1,
 \qquad
 \|\operatorname{NF}_S f\|_1\le C_S^{L_k}\|f\|_1.
 \tag{FMC12}
\]
To prove these inequalities start with one monomial of coefficient
one. A reduction replaces its leading term by the complete tail with
total absolute coefficient at most the indicated constant. Along
each branch the nonnegative integer weight drops at least one.
Induction on that weight bounds the total absolute coefficient of
all final leaves by $C^{\mathfrak w(m)}$: a standard leaf contributes
one, and an internal node contributes at most $C$ times the bound at
weight one lower. Cancellation between branches can only decrease the
coefficient $\ell^1$ norm. Linearity, the triangle inequality, and
$\mathfrak w(m)\le L_k$ prove FMC12 for every input.

There is equally explicit control of the recorded original-generator
certificates. For each stored $g\in\mathcal G_S$ let
$c_g$ be its complete certificate in the fixed original generating
list of $\mathscr S$ in FMC5--6. Define
\[
 D_S=\max\left\{1,\max_{g\in\mathcal G_S}
                             \|c_g\|_1/|\lambda_g|\right\}.
 \tag{FMC13}
\]
At a reduction node with input coefficient $c$, the certificate
contribution has norm at most $D_S|c|$. At depth $h$ the total
absolute coefficient entering all nodes is at most $C_S^h$ times
the input norm, by the same tree proof. Summing through depth $L_k$
therefore proves that a certificate of
$f-\operatorname{NF}_S f$ in the original generators has norm at most
$D_S(L_k+1)C_S^{L_k}\|f\|_1$. All certificate entries retain the
original homogeneous grading. Their denominators are among the same
finite fixed reduction data; no degree-$k$ linear system is inverted.

\subsection{The two exact maps back to the original conductor strip}

Let $\mathsf N_k$ and $\mathsf S_k$ be the original OCF strip
remainder and its coordinate insertion, with leading conductor
coefficient $a_*=a_{r_*s_*}$. Define
\[
 A_0=\max\left\{1,
                \sum_{(a,b)\ne(r_*,s_*)}|a_{ab}/a_*|\right\}.
 \tag{FMC14}
\]
Their exact coefficient bounds are
$\|\mathsf N_k\|_1\le A_0^{10k}$ and
$\|\mathsf S_k\|_1=1$. To prove the first, give the biform monomial
$e_{ab}^{(k)}$ the weight $9a+b\in[0,10k]$. Every other nonzero
term of the full $A$ has smaller weight than its lexicographically
largest term: a decrease in the first index outweighs at most eight
increases in the second index, and a decrease in only the second
index also drops the weight. Every translated OCF6 subtraction
therefore has weight drop at least one. Its complete tail has
coefficient norm at most $A_0$. The reduction-tree proof of FMC12
now applies with depth $10k$. The insertion bound is immediate
from its unchanged coordinate coefficients.

Put
\[
 h_+=\max\{1,\max_i\sum_a|\mathsf H_{ia}|\},\qquad
 h_-=\max\{1,\max_a\sum_i|(\mathsf H^{-1})_{ai}|\},
 \quad h_* =\max\{h_+,h_-\}.
 \tag{FMC15}
\]
Substitution $x=\mathsf Ht$ on degree-$k$ polynomial coefficients
has $\ell^1$ norm at most $h_+^k$. Expanding each product of linear
forms and using the multinomial theorem proves this for each
monomial; the triangle inequality proves it for their sum. The Segre
substitution adds coefficients on identical biform monomials and so
has norm at most one. Conversely choose the following exact lift
of $e_{ab}^{(k)}$: with $h=\min(a,b)$ take
\[
 t_{++}^{h}t_{+-}^{a-h}t_{-+}^{b-h}t_{--}^{k-a-b+h},
 \qquad t=\mathsf H^{-1}x.
 \tag{FMC16}
\]
All exponents are nonnegative since $0\le a,b\le k$, and their
sum is $k$. Direct substitution gives the required biform monomial.
This lift, extended linearly and denoted $\mathscr L_k$, has norm
at most $h_-^k$ by the same complete multinomial estimate.

Let $\mathscr M_{H,k}$ be the set of degree-$k$ standard module
monomials for $\mathcal G_H$. Its coefficient space is the exact
degree-$k$ quotient model of FMC7. Define
\[
 \begin{split}
 U_k f&=\mathsf N_k\,\psi\,\iota f,\qquad
          f\in\mathbb C^{\mathscr M_{H,k}},\\
 V_k w&=\operatorname{NF}_H\,\iota^{-1}
                        \mathscr L_k\mathsf S_k w,
          \qquad w\in W_k.
 \end{split}\tag{FMC17}
\]
These are inverse maps. For the first composition, $\operatorname{NF}_H$
changes a lift only by $\mathscr H$, which maps to zero in the
original strip by FMC7, and $\psi\mathscr L_k=I$ together with
$\mathsf N_k\mathsf S_k=I$ proves $U_kV_k=I$.
For the opposite composition, $\mathscr L_k\mathsf S_k
\mathsf N_k\psi\iota f-\iota f$ maps to zero in the strip, so
it belongs to $QP+\mathcal AP$ by FMC7. Its free-module difference
belongs to $\mathscr H$ and is killed by the unique remainder.
The original $f$ is already standard. Hence $V_kU_k=I$.
The bounds already proved give
\[
 \|U_k\|_1\le h_+^k A_0^{10k},\qquad
 \|V_k\|_1\le h_-^k C_H^{L_k}.
 \tag{FMC18}
\]
These bounds use only the displayed original coefficients and the
two fixed homogeneous relation lists.

\subsection{An identity pivot block for the entire residual image}

Since $\mathscr H\subset\mathscr S$, every leading monomial of
$\mathscr H$ is a leading monomial of an element of $\mathscr S$.
The proved reduction property for $\mathcal G_S$ then implies
inclusion of their leading modules. Consequently
\[
 \mathscr M_{S,k}\subset\mathscr M_{H,k},\qquad
 \mathscr U_k=\mathscr M_{H,k}\setminus\mathscr M_{S,k},\qquad
 |\mathscr M_{H,k}|=\Delta,\quad
 |\mathscr M_{S,k}|=r,\quad |\mathscr U_k|=j.
 \tag{FMC19}
\]
The first count is FMC7 and OCF8. The quotient by $\mathscr S$
is exactly the original conductor quotient modulo its invariant
image, which has dimension $r$ by OCF22--24 and OKM3--5.
Subtracting proves the last count. These are counts of actual
standard monomials for the fixed relation lists, not independent
generic ranks.

Order $\mathscr M_{H,k}$ by listing $\mathscr U_k$ first and
$\mathscr M_{S,k}$ second, with the specified monomial order within
each list. For every $u\in\mathscr U_k$ put
$z_u=u-\operatorname{NF}_S(u)$. Both terms are $H$-standard,
and their difference belongs to $\mathscr S$. The $j$ resulting
columns form a basis of $\mathscr S_k/\mathscr H_k$ in this
quotient model. Indeed their coordinates on $\mathscr U_k$ are
the distinct unit vectors. For a standard column in $\mathscr S$,
the remainder under $\operatorname{NF}_S$ is zero. Its coordinates
on $\mathscr M_{S,k}$ must therefore be the negatives of the
remainders of its $\mathscr U_k$ terms. It is exactly the stated
linear combination of the $z_u$. This proves spanning and
independence, including all tails of every column.

Write $B_k$ for the matrix with columns $\operatorname{NF}_S(u)$
in the second standard set. The full exact graph frame and inverse are
\[
 P_k=\begin{pmatrix}I_j&0\\-B_k&I_r\end{pmatrix},\qquad
 P_k^{-1}=\begin{pmatrix}I_j&0\\B_k&I_r\end{pmatrix},\qquad
 \det P_k=1,\quad
 \|P_k\|_1,\|P_k^{-1}\|_1\le1+C_S^{L_k}.
 \tag{FMC20}
\]
The two matrix identities follow by multiplication and block
triangularity, and the norm bound follows from FMC12 applied to
each entire column. All entries of $B_k$ come from fixed reductions;
no minor is selected or divided at degree $k$.

Each $z_u$ has an actual invariant certificate. The recorded reduction
certificate of $u-\operatorname{NF}_S(u)$ from FMC13 expresses its
polynomial representative as
\[
 \iota z_u=Qh_u+\mathcal A a_u+f_u,
 \qquad f_u\in\mathscr I_k,
 \qquad \|f_u\|_{1,x}
                  \le D_S(L_k+1)C_S^{L_k}.
 \tag{FMC21}
\]
To obtain the three displayed pieces, collect the certificate terms
corresponding respectively to $Qx^a$, $\mathcal Ax^a$, and the
$125$ invariant basis elements $x^a$. Every coefficient is a
polynomial in the $y_i=x_i^5$. The last terms therefore belong
to $\mathscr I$, and homogeneity gives exactly degree $k$.
FMC3 preserves coefficient norm, and discarding the first two
certificate lists when bounding only the last cannot increase
their total $\ell^1$ norm. FMC13 proves the bound. Thus $U_kz_u$
is exactly $\mathsf N_k\psi f_u$, an original invariant restriction,
with the entire original periods in its coefficients.

Define the full original strip frame, its residual columns and their
explicit left inverse by
\[
 T_k=U_kP_k,\qquad
 Z_k^{\rm fix}=T_kE_j,\qquad
 L_{Z,k}=E_j^{\mathsf T}P_k^{-1}V_k,\qquad
 L_{Z,k}Z_k^{\rm fix}=I_j.
 \tag{FMC22}
\]
Here $E_j$ inserts the first $j$ coordinates into the new ordered
standard basis, while all output strip coordinates retain their
original OCF order. FMC17 and FMC20 prove the inverse formula
$T_k^{-1}=P_k^{-1}V_k$ and the last identity. By FMC21 these
$Z_k^{\rm fix}$ columns span precisely the original $\overline J_k$.
They provide all actual residual coefficients with a left inverse
obtained from fixed relations.

For the original coefficient Euclidean norms set the explicit number
\[
 M_k=2\sqrt{\Delta}\,h_*^k A_0^{10k}(C_HC_S)^{L_k}\ge1.
 \tag{FMC23}
\]
Then
\[
 \|T_k\|_2,\|T_k^{-1}\|_2,
       \|Z_k^{\rm fix}\|_2,\|L_{Z,k}\|_2\le M_k,
 \quad
 M_k^{-1}\|v\|_2\le\|Z_k^{\rm fix}v\|_2\le M_k\|v\|_2.
 \tag{FMC24}
\]
For the first two bounds, FMC18--20 give respectively the column-sum
bounds $h_+^kA_0^{10k}(1+C_S^{L_k})$ and
$(1+C_S^{L_k})h_-^kC_H^{L_k}$. For a matrix with $\Delta$
input coordinates, its Euclidean norm is at most $\sqrt\Delta$
times its coefficient column-sum norm: use
$\|Av\|_2\le\|Av\|_1\le\|A\|_1\|v\|_1
\le\sqrt\Delta\|A\|_1\|v\|_2$.
All constants are at least one, proving the two stated bounds.
The coordinate insertion and restriction have Euclidean norm one,
so FMC22 proves the remaining operator bounds. Its left-inverse
identity then proves the lower bound on $Z_k^{\rm fix}$.

For every $0\le t\le\Delta$, singular values now give
\[
 \log^+\|\bigwedge^t T_k\|_2,
 \log^+\|\bigwedge^t T_k^{-1}\|_2
       \le t\log M_k\le\Delta\log M_k
       =O_{\rm actual}(k^2)=o(kq).
 \tag{FMC25}
\]
Indeed every singular value of both square matrices is at most
$M_k$, and exterior singular values are their products. This last
fact follows by applying a singular-value decomposition: unitary
factors preserve exterior norms, and the diagonal factor multiplies
each coordinate wedge by the corresponding product. Explicitly
\[
 \log M_k=\log2+\tfrac12\log\Delta+k\log h_*
              +10k\log A_0+L_k\log(C_HC_S).
 \tag{FMC26}
\]
All data except $k,\Delta,L_k$ are fixed actual numbers,
$L_k\le B_0^4 k/5+624$, $\Delta=16k-48$, and
$kq=k(k+1)^2$. These formulas prove the last two statements of
FMC25 directly. The same bounds hold for transpose inverse maps,
whose singular values are unchanged. For $0\le t\le j$, FMC24
also puts every singular value of $Z_k^{\rm fix}$ between
$M_k^{-1}$ and $M_k$, so the same exterior bounds apply to its
map onto its image and its inverse there.

\subsection{Exact relation to every earlier certificate frame}

Let $Z_k^{\rm old}$ be the full actual OCF/ROQ certificate frame,
with its selected invertible $j\times j$ row block
$(Z_k^{\rm old})_I$. Both frames have exactly the same column
space, as proved in FMC21--22. Their exact transition and inverse are
\[
 A_k=((Z_k^{\rm old})_I)^{-1}(Z_k^{\rm fix})_I,
 \qquad A_k^{-1}=L_{Z,k}Z_k^{\rm old},\qquad
 Z_k^{\rm fix}=Z_k^{\rm old}A_k.
 \tag{FMC27}
\]
Each new column has a unique expansion in the old independent columns;
restriction to the old pivot rows gives the first formula. Applying
$L_{Z,k}$ gives the second and both inverse identities. Thus this
is the exact coefficient morphism to the earlier frame, including
its actual pivot determinant. No bound for that old determinant
is needed or asserted by FMC24--26. The new frame and its proved
bounds use only the fixed data of FMC5--23.

The two residual maps on $V_A$ therefore obey
$\Xi_k^{\rm fix}=A_k^{\mathsf T}\Xi_k^{\rm old}$ in their
respective coordinates, since the maps in the original strip dual
are the ordinary transposes of their certificate matrices.
For every original source and cutoff their minimum quotient metrics
satisfy the exact congruence and determinant formulas
\[
 Q_{\Xi,N}^{\rm fix}
        =\overline{A_k}^{-1}Q_{\Xi,N}^{\rm old}A_k^{-\mathsf T},
 \qquad
 \log\det Q_{\Xi,N}^{\rm fix}
        =\log\det Q_{\Xi,N}^{\rm old}-2\log|\det A_k|.
 \tag{FMC28}
\]
The constraints on a lift are transformed by precisely
$A_k^{\mathsf T}$; substituting the inverse coordinates in the
same attained minimum proves the congruence, and determinants give
the second formula. The scalar is independent of $N$ and of source
order. Its coefficient is exactly $1+1-1-1=0$ in the original
four-endpoint return. Thus this return agrees identically with the
entire earlier ROQ/RCA return, even if an earlier selected pivot is
small at some degrees.

The complete inherited period Gram is transported by the same exact
map at each endpoint. Let $L_{\rm old}$ be the primary right
section ROQ6 and let $L_{\rm fix}$ be any primary right section
of $\Xi_k^{\rm fix}$ inside $V_A$, for example the first $j$
columns of FMC30 below. Both $L_{\rm fix}$ and
$L_{\rm old}A_k^{-\mathsf T}$ have residual image $I_j$ under
$\Xi_k^{\rm fix}$, so their difference has columns in the
original $K=\ker\Lambda$. Consequently
\[
 \Lambda L_{\rm fix}=\Lambda L_{\rm old}A_k^{-\mathsf T},\qquad
 P_{\rm per}^{\rm fix}
     =\overline{A_k}^{-1}P_{\rm per}^{\rm old}A_k^{-\mathsf T},
 \qquad
 \frac{\det Q_{\Xi,N}^{\rm fix}}{\det P_{\rm per}^{\rm fix}}
     =\frac{\det Q_{\Xi,N}^{\rm old}}{\det P_{\rm per}^{\rm old}}.
\]
The first equality preserves the literal boundary vectors and their
full cyclic coefficient map. Taking their Gram with the unchanged
inherited metric ROQ9 proves the second equality. Taking both
determinants together with FMC28 proves the last equality pointwise,
before taking any endpoint difference.

\subsection{The unchanged Gamma source in the fixed degree model}

Retain $c=k/2$, the original nodes
$\beta_{ab}=c+(2a-k)\delta+i(2b-k)\gamma$, and
$\omega_{ab}=(\beta_{ab}-c)/i$. For each original source
$s\in\{1,k\}$ retain its complete recurrence and masses
\[
 \begin{gathered}
 p_0=1,\ p_1=y,\quad
 p_{n+1}=yp_n-n(s/2+n-1)p_{n-1},\\
 h_n=(2\pi)^{s/2}n!(s/2)_n,\qquad
 (U_N)_{(a,b),n}=p_n(\omega_{ab})/\sqrt{h_n},\qquad
 G_N=(U_NU_N^*)^{-1}.
 \end{gathered}\tag{FMC29}
\]
Here $N$ remains one of $q-1,q,2q-1,2q$. These are exactly the
original primary quotient metrics of CEP21--22 and ROQ10, with
every source mass and coordinate phase present.

The complete primary frame of $V_A$, in residual coordinates first
and common-kernel coordinates second, is
\[
 J_k=\mathsf N_k^{\mathsf T}T_k^{-\mathsf T},\qquad
 \Xi_k^{\rm fix}J_k=E_j^{\mathsf T},\qquad
 K=\operatorname{im}(J_kE_r).
 \tag{FMC30}
\]
Here $E_r$ inserts the last $r$ coordinates. For a primary vector
$x\in V_A$, ROQ3 gives its strip dual coordinate
$z=\mathsf S_k^{\mathsf T}x$. The new residual map is
$(Z_k^{\rm fix})^{\mathsf T}z=E_j^{\mathsf T}T_k^{\mathsf T}z$.
Substitution of $z=T_k^{-\mathsf T}y$ proves the second identity,
and its kernel is the last coordinate block, proving the third.
Thus this full frame has precisely the original common kernel;
it changes no observation. Its left inverse on $V_A$ is
$T_k^{\mathsf T}\mathsf S_k^{\mathsf T}$.
The original coefficient bounds additionally give
\[
 \|J_k\|_2\le\sqrt q A_0^{10k}M_k=:\widehat M_k,
 \qquad
 \|T_k^{\mathsf T}\mathsf S_k^{\mathsf T}\|_2\le M_k.
 \tag{FMC31}
\]
Indeed FMC14 bounds the norm of $\mathsf N_k$ by its column-sum
norm times $\sqrt q$, transpose preserves Euclidean norm, and
FMC24 applies to $T_k^{-\mathsf T}$. The insertion transpose is
a coordinate restriction of norm one. Both logarithms in FMC31
are $O_{\rm actual}(k+\log k)$ with their full displayed constants.

The exact source Gram in this specified degree model, and its actual
residual quotient, are
\[
 \begin{gathered}
 \mathcal H_N=J_k^*G_NJ_k
       =\overline{T_k}^{-1}
          \overline{\mathsf N_k}G_N\mathsf N_k^{\mathsf T}
                                               T_k^{-\mathsf T},\\
 \mathcal H_N=
       \begin{pmatrix}H_{11,N}&H_{12,N}\\H_{21,N}&H_{22,N}\end{pmatrix},
 \qquad
 Q_{\Xi,N}^{\rm fix}=H_{11,N}-H_{12,N}H_{22,N}^{-1}H_{21,N}.
 \end{gathered}\tag{FMC32}
\]
The first formula is direct substitution of the actual primary
frame. Since $G_N>0$ and $J_k$ is injective, $\mathcal H_N>0$;
therefore $H_{22,N}>0$. For prescribed first coordinate $y$, the
second coordinate minimizing the full quadratic form is
$-H_{22,N}^{-1}H_{21,N}y$. Completing the square, or multiplying
out this substitution, gives exactly the displayed Schur
complement. This is the original minimum on $\Xi_k^{\rm fix}x=y$
inside $V_A$, so it equals FMC28 and the full ROQ metric.
All four blocks, including every off-diagonal entry, are retained.
The original minimum source lift of the full frame is explicitly
$U_N^*G_NJ_k$, whose Gram is $\mathcal H_N$ by $U_NU_N^*=G_N^{-1}$.
It contains the entire original polynomial recurrence and remainder
at each stated cutoff.

There is a useful quantitative determinant comparison which does not
replace that Gamma form by a coefficient form. In strip dual coordinates
write $H_{V,N}=\overline{\mathsf N_k}G_N\mathsf N_k^{\mathsf T}$ and
$Z=Z_k^{\rm fix}$. Define the actual coefficient quotient form and
the relative determinant by
\[
 Q_{\rm coeff}=(Z^{\mathsf T}\overline Z)^{-1},\qquad
 \Phi_{\Xi,N}=\log\det Q_{\Xi,N}^{\rm fix}
                                  -\log\det Q_{\rm coeff}.
 \tag{FMC33}
\]
The first is the minimum quotient of the original coefficient
Euclidean norm in the fixed OCF strip, with the identical constraint
$Z^{\mathsf T}z=y$. The elementary minimum proof of ROQ11 gives
both this formula and
$Q_{\Xi,N}^{\rm fix}=(Z^{\mathsf T}H_{V,N}^{-1}\overline Z)^{-1}$.
Hence the full, directly calculable expression is
\[
 \begin{gathered}
 \Phi_{\Xi,N}=
       \log\det(Z^{\mathsf T}\overline Z)
       -\log\det(Z^{\mathsf T}H_{V,N}^{-1}\overline Z),\\
 |\log\det Q_{\rm coeff}|\le2j\log M_k
                       =O_{\rm actual}(k^2)=o(kq),\\
 F_{\Xi,k}^{(s)}
       =\Phi_{\Xi,q-1}+\Phi_{\Xi,q}
                              -\Phi_{\Xi,2q-1}-\Phi_{\Xi,2q}.
 \end{gathered}\tag{FMC34}
\]
FMC24 puts the singular values of $Z$ between $M_k^{-1}$ and
$M_k$. The positive eigenvalues of $Z^{\mathsf T}\overline Z$
are their squares, proving the determinant bound by summing $j$
logarithms. The first identity follows by the two exact inverse
Gram formulas. The final identity follows by canceling the fixed
$\log\det Q_{\rm coeff}$ with the four original signs and then
using FMC28. Thus all coefficient denominators and frame-volume
effects in this specified construction are controlled below the
$kq$ scale, while the original Gamma compression remains fully
visible in its exact determinant.

For the kernel frame $I_k=J_kE_r$, FMC31 likewise gives the concrete
coefficient-Gram bounds
\[
 M_k^{-2r}\le\det(I_k^*I_k)\le\widehat M_k^{2r},\qquad
 |\log\det(I_k^*I_k)|
             \le2r\log\widehat M_k=O_{\rm actual}(k^2)=o(kq).
 \tag{FMC35}
\]
Its left inverse on its image is
$E_r^{\mathsf T}T_k^{\mathsf T}\mathsf S_k^{\mathsf T}$, with norm
at most $M_k$; FMC31 bounds its forward norm. The singular-value
argument just given proves FMC35. Define the exact primary-coordinate
projection and determinant by
\[
 P_{K,\rm prim}=I_k(I_k^*I_k)^{-1}I_k^*,\qquad
 D_{K,N}^{\rm prim}=[z^r]\det(I_q+zP_{K,\rm prim}G_NP_{K,\rm prim}).
\]
Multiplication proves that this projection is Hermitian, idempotent,
and has exactly the original primary-coordinate $K$ as its image.
For the determinant proof take $U=I_k(I_k^*I_k)^{-1/2}$, so
$U^*U=I_r$, and extend its columns to an orthonormal basis of
$\mathbb C^q$. In that basis the projected metric has block
$U^*G_NU$ and a zero block. Its indicated determinant coefficient
is $\det(U^*G_NU)$, and substitution of the full Gram square root
therefore proves
\[
 \det(I_k^*G_NI_k)=\det(I_k^*I_k)D_{K,N}^{\rm prim}.
\]
The precise map to AKS12--14's original polynomial-coefficient
projector is the ordered Vandermonde: put
$I_{k,\rm poly}=V_\beta^{-1}I_k$ and
$G_{N,\rm poly}=V_\beta^*G_NV_\beta$.
Their compressed Gram is identically $I_k^*G_NI_k$ by
multiplication. Applying the determinant proof just given in
these polynomial coordinates proves the exact bridge
\[
 D_{K,N}^{\rm AKS}
    =\frac{\det(I_k^*I_k)}
            {\det(I_{k,\rm poly}^*I_{k,\rm poly})}
                         D_{K,N}^{\rm prim}.
\]
Here the polynomial projector has image the actual coefficient
kernel and thus agrees with AKS12's projector. The displayed
positive factor is independent of $N$ and source order and
cancels exactly with the four endpoint signs. FMC35 bounds its
primary-coordinate numerator; the polynomial denominator is
retained as the full displayed Vandermonde Gram. No bound for
that different denominator follows from FMC35 alone.

FMC25 is a coefficient-space estimate for specified exact maps.
Its relation to the original Gamma norms is the complete congruence
and Schur calculation FMC29--34, not an assertion that those maps
are Gamma isometries. In particular no upper or lower leading
value for the remaining Gamma compression is inferred from their
coefficient condition numbers. The new result supplies a fixed
degree presentation, actual invariant certificates, bounded forward
and inverse exterior maps, and the exact original determinant on
that model, with every fixed denominator specified and every
source order and endpoint retained.

\subsection{The exact map to the original polynomial degree flag}

The homogeneous degree $k$ in FMC3 is the original tensor degree.
Its precise map to the original polynomial-in-$S$ model is the
ordered evaluation isomorphism, not an identification of the two
degree indices. Retain
$\chi_k(S)=\prod_{a,b=0}^k(S-\beta_{ab})$ and the original
Vandermonde $(V_\beta)_{(a,b),d}=\beta_{ab}^d$, $0\le d<q$.
For a fixed-model primal coordinate column $y\in\mathbb C^\Delta$
the exact polynomial representative is
\[
 p_y(S)=\sum_{a,b=0}^k(J_ky)_{ab}
          \frac{\chi_k(S)}{(S-\beta_{ab})\chi_k'(\beta_{ab})},
 \qquad [p_y]_{\rm coeff}=V_\beta^{-1}J_ky.
 \tag{FMC36}
\]
All denominators are nonzero because the original nodes are
distinct. Evaluation at each node proves the displayed inverse
and hence the equality of the two formulas. It includes the
complete roots, their order, and their actual centre. Thus the
source degree-$N$ metric in FMC32 is equivalently
$(V_\beta^{-1}J_k)^*(V_\beta^*G_NV_\beta)
(V_\beta^{-1}J_k)$, identically $J_k^*G_NJ_k$ by multiplication.
This proves the full typed relation to the original coefficient
Gamma model.

There is an exact further map to the degree-adapted conductor
basis AKJ17. Retain all actual conductor moments $\mu_d$, their
first nonzero index $v\le80$, the literal $B_q(\partial_S)$ of
AKJ14, and $\mathcal J_v S^d=d!S^{d+v}/(d+v)!$. The ordered
polynomial columns of $P_{A,k}$ are
\[
 \begin{gathered}
 1,S,\ldots,S^{v-1},\\
 P_h=\frac{(q'+h+v)!}{(q'+h)!}
       B_q(\partial_S)^{-1}\mathcal J_v(\chi_{k'}S^h),
               \qquad 0\le h\le\Delta-v-1.
 \end{gathered}\tag{FMC37}
\]
Their distinct monic degrees are $0,\ldots,v-1$ and
$q'+v,\ldots,q-1$. AKJ14--17 proves directly that they form a
basis of this same $V_A$, using the complete conductor derivative
and its finite nilpotent inverse. Define $B_{A,k}(p)$ on $V_A$
by successively subtracting the coefficient of the highest
remaining degree times its matching column in FMC37, down
through the high list and then through the low monomials.
Each subtraction uses leading coefficient one. Spanning and the
distinct degree pivots show that it ends at zero and records the
unique coordinate column of $p$; this also proves its linearity
and both inverse identities on $V_A$.
Consequently the precise transition is
\[
 \begin{gathered}
 R_{A,k}=\mathsf S_k^{\mathsf T}V_\beta P_{A,k},\qquad
 R_{A,k}^{-1}z=B_{A,k}\bigl(V_\beta^{-1}\mathsf N_k^{\mathsf T}z\bigr),\\
 V_\beta^{-1}J_k=P_{A,k}R_{A,k}^{-1}T_k^{-\mathsf T},\qquad
 \Xi_{A,k}^{\rm fix}=E_j^{\mathsf T}T_k^{\mathsf T}R_{A,k}.
 \end{gathered}\tag{FMC38}
\]
The columns $V_\beta P_{A,k}$ lie in $V_A$. Therefore ROQ3 gives
$\mathsf N_k^{\mathsf T}R_{A,k}=V_\beta P_{A,k}$, which proves
the first inverse formula by applying the unique degree
coordinates just constructed. Both sides are invertible maps
between the displayed $\Delta$-dimensional coefficient spaces.
Substituting FMC30 gives the third equality, and evaluating
the new residual map on these same primary columns gives the
last. Thus its degree-prefix ranks give exactly
$\dim(K\cap\mathbb C[S]_{\le d})$ by subtracting that prefix
rank from the number of columns of FMC37 of degree at most $d$,
as in AKJ18. Every original factorial, finite inverse, primary
denominator and degree pivot is visible in these maps. The
subleading coefficient-frame estimate proved in FMC25 pertains
to the fixed strip module maps; FMC36--38 retain the separate
polynomial interpolation and its full Gamma metric exactly.

% END CURRENT FMC

\clearpage
% BEGIN CURRENT OVR
\section{The evaluated outer contribution in the original mixed return}
\label{sec:OVR}

Retain the actual simple quartet, original period on the proved five-orbit
locus, full amplitude unit, branch, and original source orders $s=1,k$ of
OAR and OFV. All arithmetic statements retain the original domain
$|u|\ge R_*$ of AMT; the actual period is fixed as $k=4l+1$ grows.
Here $q=(k+1)^2$, $\Delta=16k-48$, and every finite use of OFV retains
its stated threshold, including OSP23. Write
\[
 C_\partial=\mathfrak C_\partial
 =4\log(4/\pi)+2-2I_{2,\pi},\qquad
 I_{2,\pi}=\iint\log|x-y|\,d\mu_{2,\pi}(x)d\mu_{2,\pi}(y).
 \tag{OVR1}
\]
The measure here is the full EIQ density with its uniquely specified
endpoints. OFV proves this exact expression and $C_\partial\ge0$.
No value of the original residual observation is inserted into it.

\subsection{The exact remainder and the sourcewise leading identity}
Define the actual outer remainder and retain OAR3's full conductor and
monic-tail remainder by
\[
 b_{k,s}=\mathcal F_{\partial,k}^{(s)}-\Delta q C_\partial,
 \qquad
 R_{k,s}=b_{k,s}+\mathfrak e_{k,s}.
 \tag{OVR2}
\]
Thus $b_{k,s}=O_{\delta,\gamma}(q\log q)$ by OFV30. The second
summand is precisely the signed finite expression OAR3, with
$-e_{k,s}^-\le\mathfrak e_{k,s}\le e_{k,s}^+$ and
$e_{k,s}^\pm=o(kq)$. In particular $R_{k,s}=o(kq)$, because
$\log q/k\to0$. Its two exact finite error endpoints remain
$b_{k,s}-e_{k,s}^-$ and $b_{k,s}+e_{k,s}^+$; their signs are
not replaced by a symmetry assumption. Substitution into OAR3 proves
\[
 \begin{gathered}
 F_K^{(s)}+\Phi_{k,s}=\Delta q C_\partial+R_{k,s},\\
 \frac{F_K^{(s)}}{kq}+\frac{\Phi_{k,s}}{kq}
       =16C_\partial-\frac{48C_\partial}{k}
                         +\frac{R_{k,s}}{kq}
       \longrightarrow16C_\partial,\qquad s=1,k.
 \end{gathered}\tag{OVR3}
\]
Here $\Phi_{k,s}$ is exactly ROQ19's positive sum of constrained
covariance increments, built from the original period certificate
matrix $D$ and the complete conductor cross covariance. Its coefficient
maps, minimum sections, inherited period metric and graph denominator
are those already proved in ROQ; OVR3 changes none of them.

\subsection{Transfer between the two original Gamma sources}
Let $d_k=F_K^{(k)}-F_K^{(1)}$ denote the actual difference.
AMT44--46 proves $d_k=o(kq)$, with its original finite two-sided
source-comparison bounds. Subtracting the two exact lines OVR3 gives
\[
 \Phi_{k,k}-\Phi_{k,1}=R_{k,k}-R_{k,1}-d_k,
 \qquad
 \frac{\Phi_{k,k}-\Phi_{k,1}}{kq}\longrightarrow0.
 \tag{OVR4}
\]
This proves the transfer using the actual two residual observations,
whose individual metrics retain their respective original sources.
There is no asserted equality between those metrics. The equality in
OVR4 is the exact scalar consequence of their common original-kernel
restriction maps and the proved determinant identities.

\subsection{Arithmetic kernel and its correlated boundary}
Put $d_k^{\rm ar}=F_K^{\rm ar}-F_K^{(1)}$. This is the actual
quantity bounded in OAR8; AMT38 and AMT46 prove
$d_k^{\rm ar}=o(kq)$. Keeping its full value first gives
\[
 \begin{split}
 F_K^{\rm ar}
    &=\Delta q C_\partial-\Phi_{k,1}+R_{k,1}+d_k^{\rm ar},\\
 F_B^{\rm ar}
    &=\mathcal B_k^{\rm ar}-\Delta q C_\partial
          +\Phi_{k,1}-R_{k,1}-d_k^{\rm ar}.
 \end{split}\tag{OVR5}
\]
The first equality adds the literal definition of $d_k^{\rm ar}$
to OVR3 at $s=1$. The second subtracts the first from
$F_K^{\rm ar}+F_B^{\rm ar}=\mathcal B_k^{\rm ar}$.
In particular the same arithmetic total and the same two finite
errors occur with opposite signs in the boundary. OAR8 supplies
the unchanged correlated finite bounds on $d_k^{\rm ar}$.
Dividing these exact equalities gives
\[
 \frac{F_K^{\rm ar}+\Phi_{k,1}}{kq}\longrightarrow16C_\partial,
 \qquad
 \frac{\mathcal B_k^{\rm ar}-F_B^{\rm ar}+\Phi_{k,1}}{kq}
                  \longrightarrow16C_\partial.
 \tag{OVR6}
\]
Thus evaluation of the residual sum in either original Gamma order
passes through to the same original arithmetic kernel and boundary.

\subsection{The signed return retains its original Gamma correction}
The exact signed quantity of OAR11 is
$\mathfrak S_k^{\rm mix}=F_K^{\rm ar}+F_B^{\rm ar}-F_B^{(k)}$.
Applying OVR3 at $s=k$ to that identity yields
\[
 \mathfrak S_k^{\rm mix}+\Phi_{k,k}
     =\Delta q C_\partial+\mathcal H_k+\delta_k^\sigma+R_{k,k}.
 \tag{OVR7}
\]
All finite terms are the original ones. AMT48--49, with the retained
total correction estimates used there, proves
$(\mathcal H_k+\delta_k^\sigma)/(kq)\to-\mathcal C_\Gamma/4$,
where
\[
 \mathcal C_\Gamma
   =2\int\log x\,d\mu_{2,\pi}(x)-4(2\log2-1)>0.
\]
The equilibrium and source maps are exactly the EIQ/GEL providers
used in OFV and AMT. Hence OVR7 and OVR4 prove both limits
\[
 \frac{\mathfrak S_k^{\rm mix}+\Phi_{k,k}}{kq}
       \longrightarrow16C_\partial-\frac{\mathcal C_\Gamma}{4},
 \qquad
 \frac{\mathfrak S_k^{\rm mix}+\Phi_{k,1}}{kq}
       \longrightarrow16C_\partial-\frac{\mathcal C_\Gamma}{4}.
 \tag{OVR8}
\]
The two constants are explicit integrals of the same proved density,
respectively its full logarithmic pair energy and its logarithmic
moment. No identity between these distinct integrals is presumed.

\subsection{Bounds on the remaining residual coefficient}
Each $F_K^{(s)}$ is nonnegative: the original quotient metrics decrease
with increasing source degree, and restriction to the identical kernel
preserves the two Gram inequalities at the paired endpoints. Their
determinant ratios are consequently at least one, as proved in AMT15.
ROQ19 gives $\Phi_{k,s}\ge0$. Thus OVR3 bounds both nonnegative
sequences after division by $kq$. Let
$\kappa_*=32\log(25\log3/(4\pi))$ be the retained AKS asymptotic
kernel ceiling on this actual locus, applied in both original orders.
The actual residual cap OAR5, proved by ROQ21, is
$\Phi_{k,s}\le j(\mathcal L_{s,q}+\mathcal L_{s,q+1})-\lambda_s$,
where $j=8k-32$ and $\lambda_s\ge0$ is the original first increment.
AKS31 gives
$\limsup j(\mathcal L_{s,q}+\mathcal L_{s,q+1})/(kq)\le\kappa_*$
in both original source orders: its same explicit logarithmic
comparison constants apply, and $j/k\to8$ just as $r/k\to8$.
Thus both the original kernel and residual sequences have the same
proved upper ceiling. Combining their two ceilings with OVR3 gives
\[
 \begin{gathered}
 \max\{0,16C_\partial-\kappa_*\}
       \le\liminf\frac{F_K^{(s)}}{kq}
       \le\limsup\frac{F_K^{(s)}}{kq}
       \le\min\{\kappa_*,16C_\partial\},\\
 \max\{0,16C_\partial-\kappa_*\}
       \le\liminf\frac{\Phi_{k,s}}{kq}
       \le\limsup\frac{\Phi_{k,s}}{kq}
       \le\min\{\kappa_*,16C_\partial\}.
 \end{gathered}\tag{OVR9}
\]
For each line, subtract the other bounded nonnegative sequence from
the convergent sum in OVR3. More explicitly, that sum differs from
$16C_\partial$ by a sequence tending to zero; taking its upper and
lower eventual bounds proves
$\liminf(\Phi/(kq))=16C_\partial-\limsup(F_K/(kq))$ and
$\limsup(\Phi/(kq))=16C_\partial-\liminf(F_K/(kq))$.
The same argument with the two sequences interchanged proves the
first line, including its lower bound. The two first positive
increments also survive in the finite common upper bound. In the
original AKS notation let $d_1$ be the first section-corrected kernel
increment of AKS29, in the same source order $s$. AKS28 and OAR5 give
\[
 \begin{split}
 \Delta q C_\partial+R_{k,s}
   &=F_K^{(s)}+\Phi_{k,s}\\
   &\le(r+j)(\mathcal L_{s,q}+\mathcal L_{s,q+1})-d_1-\lambda_s\\
   &=\Delta(\mathcal L_{s,q}+\mathcal L_{s,q+1})-d_1-\lambda_s.
 \end{split}\tag{OVR10}
\]
The last equality is the exact rank identity
$(8k-16)+(8k-32)=16k-48$; both subtractions retain their actual
minimum-section definitions. Adding the eventual upper bounds
already proved and using OVR3 also proves the explicit consequence
\[
 0\le16C_\partial\le2\kappa_*,\qquad
 0\le C_\partial\le\frac{\kappa_*}{8}.
 \tag{OVR11}
\]
Indeed for every positive $\varepsilon$, both sequences are eventually
at most $\kappa_*+\varepsilon$, so their limiting sum is at most
$2\kappa_*+2\varepsilon$; let $\varepsilon$ decrease to zero.
Their nonnegativity proves the lower bound independently.
OVR4 and AMT46 carry the same limiting bounds to the two original
sources and the arithmetic kernel. The outstanding coefficient is
therefore the explicit original positive residual sum; no convergence
or numerical value for that sequence has been assumed.

% END CURRENT OVR


\clearpage
% BEGIN CURRENT OOQ
% Original outer quotient operator control; full root polynomial retained.
\section{All directions of the original outer quotient at the leading scale}
\label{sec:OOQ}

Retain the complete original simple quartet, the source orders
$s=1,k$, $k=4l+1$, $0<\delta<1/2$, $\gamma>2$, and the
polynomials and masses of CTV1--22 and OSP1--35. Put
\[
 q=(k+1)^2,\quad q'=(k-7)^2,\quad d=q-q'=16k-48,\quad
 n=q/2,\quad \vartheta=\pi/2,\quad R=k\sqrt{\delta^2+\gamma^2}.
 \tag{OOQ1}
\]
The integer $q$ is even. The entire outer polynomial is
$P_{\partial,k}=P_k/P_{k-8}$, monic of degree $d$, with every
original root retained. Each root has modulus at most $R$.
Explicitly $P_k(y)=i^{-q}\chi_k(c+iy)$ and
$P_{k-8}(y)=i^{-q'}\chi_{k-8}(c'+iy)$, with the unchanged
original centres $c=k/2$ and $c'=c-4$. These are the exact
CTV3--4 phase and coefficient maps.
Its target coordinates throughout are the coefficients of
$1,y,\ldots,y^{d-1}$ in its unique remainder. For
$L=d-1,d,q+d-1,q+d$, let $Q_{\partial,L}^{(s)}$ be the
minimum quotient Gram of this remainder map from polynomials of
degree at most $L$, with the complete source
\[
 d\zeta_{k',s}(y)=|P_{k-8}(y)|^2dm_s(y),\qquad
 dm_s(y)=\frac{(2\pi)^{s/2}2^{s/2}}{4\pi\Gamma(s/2)}
          |\Gamma(s/4+iy/2)|^2dy .
 \tag{OOQ2}
\]
These are exactly the CTV10--13 Grams, not merely determinants
equal to those Grams. We prove operator estimates in this same
original remainder frame.

\subsection{Uniform equilibrium constants and a polynomial evaluation bound}

Use the complete EIQ equilibrium for
$V_{\alpha,\beta}(x)=\beta\sqrt x-\alpha\log x$ on $(0,\infty)$.
In this proof take the compact rectangle
$\mathcal K=[3/2,5/2]\times[\pi/4,\pi]$. EIQ is proved on the
entire positive quadrant, so it applies to this rectangle.
Write its exact density as $\rho_{\alpha,\beta}$, its endpoints
as $a,b$, its logarithmic potential as $P(z)$, and retain
\[
 V(x)-2P(x)=\ell+D(x),\quad D(x)\ge0,\quad
 D|_{[a,b]}=0,\quad
 P_0=P(0)=\int\log t\,\rho(t)\,dt,\quad W=P_0+\ell/2.
 \tag{OOQ3}
\]
Here and below a parameter on these symbols is understood whenever
two parameter pairs occur in the same expression.

The following fixed constants are all defined by these exact
functions on $\mathcal K$:
\[
\begin{gathered}
 a_*=\min_{\mathcal K}a>0,\quad B_*=\max_{\mathcal K}b<\infty,
 \quad b_*=\min_{\mathcal K}(b-a)>0,\quad
 M_*=\max_{\mathcal K,x}\rho(x)<\infty,\\
 L_*=\max_{\mathcal K,\ x\in[a,b]}|V'(x)|,\quad
 C_{\rm BM}=\sqrt{32}\,e^{3L_*/4},\quad
 C_{\rm quant}=B_*+2M_*,\quad C_0=B_*/a_*,\\
 C_{\rm tail}=\max\{0,\sup_{\mathcal K,\ x>0}
                    [\log(1+\sqrt x)-D(x)]\},\quad
 T_*=\max\{1,\sup_{\mathcal K}\int_0^\infty e^{-D(x)}dx\}.
\end{gathered}\tag{OOQ4}
\]
These constants are finite. Positivity, continuity and boundedness
of the endpoints and density are EIQ4, EIQ10--11 and EIQ29;
the positive gap follows by compactness. The derivative of $V$
is bounded on $[a_*,B_*]$. Uniformly over this rectangle,
$P(x)\le\log(1+x)+\log(1+B_*)$ and $P(x)\ge-2M_*$, as in
the direct logarithmic estimates in QLG9--13. The Euler constant
is bounded by evaluating OOQ3 on the compact support.
Consequently
\[
 D(x)\ge\frac{\pi}{4}\sqrt x-\frac52\log x
                         -2\log(1+x)-C\qquad(x\ge1)
\]
for a fixed finite constant $C$. This proves the assertions at
infinity in the last two definitions of OOQ4. At zero,
$D(x)\to+\infty$ uniformly because $\alpha\ge3/2$ and $P$
is bounded there. For integrability one can use precisely
the finite majorant
$e^C(1+x)^2(1+x^{5/2})e^{-\pi\sqrt x/4}$.
Thus neither of the two supremum definitions assumes a tail estimate.

The function $W$ and its first derivatives are bounded on
$\mathcal K$. Indeed EIQ30 gives $F_\alpha=P_0$, EIQ31 gives
$\beta\int\sqrt x\,d\mu=2(\alpha+1)$, and hence
\[
 W(\alpha,\beta)
 =-F(\alpha,\beta)-(\alpha+1)+(1+\alpha/2)F_\alpha(\alpha,\beta).
 \tag{OOQ5}
\]
The smoothness asserted here follows from EIQ29--30. Fix
\[
 C_W=1+\max_{\mathcal K}|W|
                   +(4+\pi)\max_{\mathcal K}(|W_\alpha|+|W_\beta|).
 \tag{OOQ6}
\]

For an integer $m\ge\max\{1,2/b_*\}$ and a polynomial $p$ of
degree at most $m$, set
$\|p\|_{m,V}^2=\int_0^\infty|p(x)|^2e^{-mV(x)}dx$.
For all $z$ outside $[a,b]$ we have
\[
 |p(z)|\le C_{\rm BM}m^2
             \exp\{m(P(z)+\ell/2)\}\,\|p\|_{m,V}.
 \tag{OOQ7}
\]
Here is a complete proof. Let
$M=\max_{[a,b]}|p(x)|e^{-mV(x)/2}$, attained at $x_0$.
Choose a one-sided interval starting at $x_0$, contained in
$[a,b]$, of length $1/m$; one side has that length because
$b-a\ge2/m$. On that interval,
$|p(x)|\le e^{L_*/2}|p(x_0)|$.

The elementary derivative bound
$\|p'\|_{\infty,J}\le4m^3|J|^{-1}\|p\|_{\infty,J}$
holds also for complex coefficients. To prove it, expand on
the affine copy of $[-1,1]$ in the polynomials $T_j$ defined
by $T_j(\cos\theta)=\cos(j\theta)$.
Orthogonality of cosines gives coefficients of modulus at most
$2\|p\|_\infty$ for $j\ge1$. The identity
$T_j'(\cos\theta)=j\sin(j\theta)/\sin\theta$ and the sum of
$j$ unit complex numbers give $|T_j'|\le j^2$, including
the limiting endpoints. Summing $\sum_{j=1}^m j^2\le m^3$
and applying the affine derivative factor $2/|J|$ proves
the bound.
It follows that on a segment of length
$[8e^{L_*/2}m^4]^{-1}$ starting at $x_0$,
$|p|\ge|p(x_0)|/2$. The weight on that segment is at least
$e^{-L_*}e^{-mV(x_0)}$. Integration gives
$\|p\|_{m,V}^2\ge M^2/(32e^{3L_*/2}m^4)$.

It remains to extend the weighted supremum to complex $z$.
The function $\log|p(z)|-mP(z)$ is subharmonic off $[a,b]$.
On this interval its boundary upper values are at most
$\log M+m\ell/2$ by OOQ3. It extends subharmonically over
infinity: in the local coordinate $w=1/z$ its logarithmic
term is $(m-\deg p)\log|w|$, plus a harmonic function.
The maximum principle on the complement of the interval in
the sphere therefore bounds it by the same boundary constant.
One may first remove small neighborhoods of polynomial zeros
and of the interval and then pass to the limit; $P$ is
continuous because its density is bounded. Combining this
bound with the preceding norm bound proves OOQ7.
For $|z|\le a_*/2$ its useful uniform consequence is
\[
 |p(z)|\le C_{\rm BM}m^2
   e^{mW+2m|z|/a_*}\|p\|_{m,V},
 \tag{OOQ8}
\]
since $|\log|1-z/t||\le2|z|/a_*$ on the support.

\subsection{A matching polynomial lift from actual equilibrium quantiles}

For each $n$ divide the actual equilibrium measure into $n$
intervals of mass $1/n$ and take the midpoint of each interval
as $\xi_i$. Retain the polynomial
\[
 r_n(x)=\prod_{i=1}^n(1-x/\xi_i),\qquad r_n(0)=1.
 \tag{OOQ9}
\]
All its roots lie in $[a_*,B_*]$. For every real $x>0$,
\[
 \log|r_n(x)|\le n(P(x)-P_0)+C_{\rm quant}\sqrt n+C_0.
 \tag{OOQ10}
\]
To prove it, use $f_\eta(t)=\log\max\{|x-t|,\eta\}$, a
$1/\eta$-Lipschitz function. For a quantile interval of length
$\Delta_i$, the difference between its point value and the
average against its probability density is at most
$\Delta_i/\eta$. Summing gives at most $B_*/\eta$.
Also
\[
0\le\int[f_\eta(t)-\log|x-t|]\rho(t)dt
\le M_*\int_{-\eta}^{\eta}\log(\eta/|s|)ds=2M_*\eta .
\]
Choose $\eta=n^{-1/2}$. Finally $\log t$ is
$1/a_*$-Lipschitz on the support, so
$|\sum_i\log\xi_i-nP_0|\le B_*/a_*$. These three facts
prove OOQ10, including at a polynomial zero.
For complex $|z|\le a_*/2$, the same literal product gives
\[
 |r_n(z)|^{-1}\le e^{2n|z|/a_*}.
 \tag{OOQ11}
\]
This uses $|1-z/\xi_i|\ge\exp(-2|z|/a_*)$ separately
for every original quantile factor.

\subsection{The entire analytic remainder map}

For the present $n=q/2$ put
\[
 \rho=n^{-1/8},\qquad
 \widehat P(z)=n^{-d}P_{\partial,k}(nz),\qquad r_*=R/n.
 \tag{OOQ12}
\]
Impose the following explicit eventual conditions, together
with the original OSP23 domain and $k\ge257$:
\[
 n\ge2d+1,\quad n\ge2/b_*,
 \quad r_*\le\rho/2,\quad
 \rho\le\min\{1/2,\sqrt{a_*/2}\}.
 \tag{OOQ13}
\]
They hold for all sufficiently large original $k$ at every fixed
$\delta,\gamma$: $n\asymp k^2$, $d=O(k)$ and $r_*=O(1/k)$.
These inequalities specify the finite range used here; no
assertion at excluded small degrees is required.

For $f$ holomorphic on the closed disk of radius $\rho$ define
its complete remainder by
\[
\operatorname{rem}_{\widehat P}f(z)=
\frac1{2\pi i}\int_{|\zeta|=\rho}
 f(\zeta)\,
 \frac{\widehat P(\zeta)-\widehat P(z)}
      {(\zeta-z)\widehat P(\zeta)}\,d\zeta .
 \tag{OOQ14}
\]
The numerator divided by $\zeta-z$ is a polynomial in $z$
of degree at most $d-1$. At every root of $\widehat P$,
Cauchy's formula gives the value of $f$; all these original
roots are distinct, so this is exactly the unique interpolation
remainder. In particular it is the usual polynomial remainder
when $f$ is polynomial. This proves the exact map and its
agreement with the original target.

Write $\|\cdot\|_1$ for the full monomial coefficient norm.
The root bound in OOQ13 proves
\[
\|\operatorname{rem}_{\widehat P}f\|_1
 \le C_{\rm rem}\max_{|\zeta|=\rho}|f(\zeta)|,\qquad
C_{\rm rem}=\rho d(1+r_*)^d(2/\rho)^d .
 \tag{OOQ15}
\]
Indeed $|\widehat P(\zeta)|\ge(\rho-r_*)^d\ge(\rho/2)^d$.
If $\widehat P(z)=\sum_{a=0}^d p_a z^a$, then
\[
\frac{\widehat P(\zeta)-\widehat P(z)}{\zeta-z}
 =\sum_{a=1}^d p_a\sum_{h=0}^{a-1}\zeta^{a-1-h}z^h,
\]
whose coefficient norm is at most $d\sum_a|p_a|$
because $\rho\le1$. The sum of absolute coefficients of
the full monic product is at most $(1+r_*)^d$.
The contour length supplies the factor $\rho$, proving OOQ15.
In particular $\log C_{\rm rem}=O(d\log n+\log d)$.
No outer root has been moved to zero in this map.

\subsection{High source degrees: a full operator estimate}

First use the precise exponential source
$|y|^{2Q}e^{-\beta|y|}dy$, where
$Q=q'+\ell_s$, $\ell_s=(s-1)/4$, and
$\beta\in\{\vartheta-1/q,\vartheta+1/q\}$.
Let $Q_{L}^{\exp}(Q,\beta)$ be its quotient Gram by the same
entire $P_{\partial,k}$, in the original $y$ remainder frame.
Set
\[
\alpha_0=(Q-\tfrac12)/n,\quad
c_{\rm high}(Q,\beta)=(2Q+1)\log n-2nW(\alpha_0,\beta).
\tag{OOQ16}
\]
The two high source degrees are $L=q+d-1,q+d$.
For their two parity degrees put
$m_e=\lfloor L/2\rfloor$, $m_o=\lfloor(L-1)/2\rfloor$.
The positive-half-line parameters are exactly
\[
(\alpha_e,\beta_e)=((Q-\tfrac12)/m_e,n\beta/m_e),\qquad
(\alpha_o,\beta_o)=((Q+\tfrac12)/m_o,n\beta/m_o).
\tag{OOQ17}
\]
All these pairs and $(\alpha_0,\beta)$ lie in $\mathcal K$
for the stated $k\ge257$. For example
$d/q<1/16$, $\ell_s/q<1/64$ and the rounding errors are at
most one; these give $3/2<\alpha_e,\alpha_o<5/2$.
The rate lies between $\pi/4$ and $\pi$ because
$m_e,m_o=n+O(d)$ and $1/q<\vartheta/4$.
More directly the actual lower bounds
$m_e,m_o\ge n$ and upper bounds $m_e,m_o\le n+d$
together with $d/q<1/16$ verify the rate bounds.

For a polynomial $p(nz)=p_e(z^2)+z p_o(z^2)$, parity and the
two substitutions $z=\pm\sqrt x$ give exactly
\[
n^{-(2Q+1)}\|p\|_{\exp}^2
 =\int_0^\infty |p_e(x)|^2x^{Q-1/2}e^{-n\beta\sqrt x}dx
 +\int_0^\infty |p_o(x)|^2x^{Q+1/2}e^{-n\beta\sqrt x}dx .
\tag{OOQ18}
\]
Both measures on the right are precisely the OOQ7 measures at
OOQ17. There is no missing factor two: each branch contributes
the positive Jacobian $1/(2\sqrt x)$.
The smooth bound OOQ6 gives
\[
|m_eW(\alpha_e,\beta_e)-nW(\alpha_0,\beta)|,\ 
|m_oW(\alpha_o,\beta_o)-nW(\alpha_0,\beta)|
 \le C_W(d+2).
\tag{OOQ19}
\]
For verification, $|m-n|\le d+1$,
$m|\alpha-\alpha_0|\le1+(5/2)|m-n|$ and
$m|\beta_m-\beta|\le\pi|m-n|$.
The mean-value formula on the rectangle and OOQ6 prove OOQ19.

Apply OOQ8 at $x=\zeta^2$ and $|\zeta|=\rho$, then use
OOQ18 and Cauchy--Schwarz on its two norms. If $v$ is the
coefficient vector of the remainder of $p(nz)$ modulo
$\widehat P$, OOQ15 gives
\[
\|v\|_2\le C_{\rm rem}\sqrt2 C_{\rm BM}(n+d)^2
 e^{\,nW+C_W(d+2)+2(n+d)\rho^2/a_*}
 n^{-(2Q+1)/2}\|p\|_{\exp}.
\tag{OOQ20}
\]
This holds for every lift, and proves the lower bound for the
quotient metric in these scaled remainder coordinates.

For the upper bound use the actual equilibrium at
$(\alpha_0,\beta)$ to form $r_n$ in OOQ9.
For the given remainder polynomial $v(z)$ put
\[
a_v(z)=\operatorname{rem}_{\widehat P}
                \left(\frac{v(z)}{r_n(z^2)}\right),\qquad
p_v(nz)=r_n(z^2)a_v(z).
\tag{OOQ21}
\]
The denominator has no zero in the disk because all its
$x$-roots are at least $a_*>\rho^2$.
The polynomial $p_v$ has degree at most $q+d-1\le L$.
At every original outer root its value is exactly $v$;
thus it is an actual lift of the prescribed remainder.
Equations OOQ11 and OOQ15 give
\[
\|a_v\|_1\le C_{\rm rem}\sqrt d\,
                  e^{2n\rho^2/a_*}\|v\|_2 .
\tag{OOQ22}
\]
For each real $x>0$ both
$|a_v(\sqrt x)|$ and $|a_v(-\sqrt x)|$ are at most
$\|a_v\|_1(1+\sqrt x)^{d-1}$.
By OOQ4,
$(1+\sqrt x)^{2(d-1)}
\le e^{2C_{\rm tail}(d-1)}e^{2(d-1)D(x)}$.
Combine this with OOQ10 and
$2nP-nV=-n\ell-nD$.
Since $n-2(d-1)\ge1$, integration using $T_*$ proves
\[
\begin{split}
\|p_v\|_{\exp}^2\le{}&
 e^{c_{\rm high}(Q,\beta)}
 \|a_v\|_1^2
 \exp\{2C_{\rm quant}\sqrt n+2C_0+
                   2C_{\rm tail}(d-1)\}T_* .
\end{split}\tag{OOQ23}
\]
All source scalings in this expression come from the explicit
$y=nz$ substitution; in the original lower-source variable it
is precisely $S'=c'+inz$, with the same original centre and phase.

Define the finite nonnegative errors
\[
\begin{split}
E^-_{\rm high}={}&
2\log\!\left(C_{\rm rem}\sqrt2 C_{\rm BM}(n+d)^2\right)
 +2C_W(d+2)+4(n+d)\rho^2/a_*,\\
E^+_{\rm high}={}&
2\log C_{\rm rem}+\log d+4n\rho^2/a_*
 +2C_{\rm quant}\sqrt n+2C_0+
                   2C_{\rm tail}(d-1)+\log T_*,\\
E_{\rm high}={}&
\max\{0,E^-_{\rm high},E^+_{\rm high}\}+2(d-1)\log n .
\end{split}\tag{OOQ24}
\]
The original and scaled remainder vectors are related by the
exact diagonal map
$D_n=\operatorname{diag}(1,n,\ldots,n^{d-1})$.
Hence the original Gram is $D_n^*Q_{\rm scaled}D_n$.
Since $I\preceq D_n^*D_n\preceq n^{2(d-1)}I$,
OOQ20--24 prove, at both high endpoints,
\[
 e^{c_{\rm high}(Q,\beta)-E_{\rm high}}I_d
 \preceq Q_L^{\exp}(Q,\beta)
 \preceq e^{c_{\rm high}(Q,\beta)+E_{\rm high}}I_d .
\tag{OOQ25}
\]
The errors obey
$E_{\rm high}=O(n^{3/4}+d\log n)$.
This is a full quadratic-form bound for the entire original
outer denominator, not an inference from a determinant estimate.

\subsection{Low source degrees, including the extra ideal column}

Put $g(Q,\beta)=2Q\log(2Q/\beta)-2Q$ and $y_0=2Q/\beta$.
For polynomials of degree at most $d$, their full exponential
source Gram is bounded below by
\[
\frac{e^{g(Q,\beta)-\beta}}
 {(d+1)^2(2d+1)64^d(1+y_0)^{2d}}\,I_{d+1}.
\tag{OOQ26}
\]
To prove it, the weight on $[y_0,y_0+1]$ is at least
$e^{g(Q,\beta)-\beta}$. On $[0,1]$ the orthonormal shifted
Legendre polynomial of degree $j$ is
$\sqrt{2j+1}L_j(2x-1)$. Its coefficient sum is at most
$\sqrt{2j+1}\,8^j$: Rodrigues' formula gives coefficients
$(-1)^{j-r}\binom jr\binom{j+r}r$, and
$\binom{j+r}r\le4^j$.
Thus the coefficient matrix of these orthonormal polynomials
has Euclidean norm at most
$\sqrt{(d+1)(2d+1)}\,8^d$.
The translation back by $-y_0$ has coefficient norm at most
$(1+y_0)^d$ and Euclidean norm at most
$\sqrt{d+1}(1+y_0)^d$.
These two invertible maps prove OOQ26 by direct norm comparison.
The Legendre orthogonality and full Rodrigues norm are proved
in OSP15 and are also obtained by its stated integrations by parts.

The largest eigenvalue of the same full Gram is at most its trace,
which is exactly
\[
\sum_{j=0}^d\frac{2(2Q+2j)!}{\beta^{2Q+2j+1}} .
\tag{OOQ27}
\]
This follows by the two real half-line integrals and repeated
integration by parts. At $L=d-1$ the remainder map is the
identity. At $L=d$ it has matrix $[I_d,-p_{\rm low}]$,
where $p_{\rm low}$ contains every lower coefficient of
$P_{\partial,k}$. Its squared operator norm is at most
$(d+1)(1+R)^{2d}$, because that complete monic polynomial
has coefficient sum at most $(1+R)^d$.
The minimum Euclidean norm on its fibre is therefore at least
the target norm divided by this operator norm; the canonical
degree-less-than-$d$ lift gives the upper bound.

For both low endpoints define
\[
\begin{split}
E_{\rm low}(Q,\beta)=\max\bigg\{0,\ 
&\log\!\left(\sum_{j=0}^d
       \frac{2(2Q+2j)!}{\beta^{2Q+2j+1}}\right)-g(Q,\beta),\\
&\beta+3\log(d+1)+\log(2d+1)+d\log64
       +2d\log(1+y_0)+2d\log(1+R)\bigg\}.
\end{split}\tag{OOQ28}
\]
Equations OOQ26--27 and the exact remainder matrices prove
\[
e^{g(Q,\beta)-E_{\rm low}}I_d
\preceq Q_L^{\exp}(Q,\beta)
\preceq e^{g(Q,\beta)+E_{\rm low}}I_d,\qquad L=d-1,d.
\tag{OOQ29}
\]
In the present range this error is $O_{\delta,\gamma}(d\log q)$.
For the upper side, integrating $\log t$ between consecutive
integers gives
$m\log m-m+1\le\log(m!)\le m\log m-m+1+\log m$.
Consequently each logarithm in OOQ27 differs above from
$g(Q,\beta)$ by $O(d\log(Q+d)+\log(Q+d))$,
uniformly over $\beta\in[\pi/4,\pi]$. The logarithm of the
number of summands adds only $\log(d+1)$.
The displayed lower error has the same order since
$Q\asymp q$, $d=O(k)$ and $R=O_{\delta,\gamma}(k)$.

\subsection{Transport back to both complete original Gamma sources}

Retain the exact source recurrence factors of OFV1--13:
\[
\ell_s=(s-1)/4,\quad
\beta_s=\frac{(2\pi)^{s/2}}{\sqrt2\,\Gamma(s/2)},\quad
Q_s=q'+\ell_s,\quad h=1/q,
\]
\[
L_h=C_L(\sin h)^{3/2},\qquad
U_h=\frac{\Gamma(1/4)^2}{2\pi}(\sin h)^{-1/2},
\qquad C_L=\frac{\pi}{\Gamma(3/4)^2}.
\tag{OOQ30}
\]
Use $\mathcal E_k$ from OSP24 and $r_{\ell_s}$ from OFV6,
with $r_0=0$. For every polynomial of all the degrees used
here, the complete source comparisons give the quadratic-form
inequalities
\[
\begin{split}
\beta_s L_h e^{-\mathcal E_k}
 \int |p(y)|^2|y|^{2Q_s}e^{-(\vartheta+h)|y|}dy
&\le\int |p(y)|^2d\zeta_{k',s}(y)\\
&\le
\beta_s U_h e^{\mathcal E_k+r_{\ell_s}}
 \int |p(y)|^2|y|^{2Q_s}e^{-(\vartheta-h)|y|}dy .
\end{split}\tag{OOQ31}
\]
OSP20--26 supplies the first root-polynomial comparison on the
entire coefficient space, retaining every coefficient of
$P_{k-8}$. OFV5--7 supplies the exact original-order density
multiplier, including every factor $t_j$ and $\beta_s$.
OFV8--9 supplies both pointwise bounds on the one-factor Gamma
density, including near zero. Composing these three proved
inequalities gives exactly OOQ31. The degree bound required
there is satisfied because $L\le q+d\le2q$.
Every minimum here is taken on the same original fibre
$\operatorname{rem}_{P_{\partial,k}}p=v$; therefore OOQ31
passes to the quotient Grams by taking its attained minima.
No multiplication map is presumed to preserve the kernel.

Define the two common scalar centres
\[
\begin{split}
c_{\rm lo}(k)&=g(q,\vartheta),\\
c_{\rm hi}(k)&=2q\log(q/2)-qW(2,\vartheta).
\end{split}\tag{OOQ32}
\]
For an entirely finite common error, take
\[
\begin{split}
\varepsilon_k={}&\mathcal E_k+
 \max_{s=1,k}r_{\ell_s}+\max\{|\log L_h|,|\log U_h|\}\\
&+\max_{\substack{s\in\{1,k\}\\ b\in\{\vartheta-h,\vartheta+h\}}}
\left\{\begin{aligned}
 &E_{\rm low}(Q_s,b)+E_{\rm high}
   +|g(Q_s,b)-c_{\rm lo}(k)|\\
 &\quad+|c_{\rm high}(Q_s,b)-c_{\rm hi}(k)|
 \end{aligned}\right\}.
\end{split}\tag{OOQ33}
\]
All constants in this definition were specified above from
the original parameters and the exact equilibrium.
The resulting operator conclusion is
\[
\begin{array}{ll}
\beta_s e^{c_{\rm lo}-\varepsilon_k}I_d
\preceq Q_{\partial,L}^{(s)}
\preceq\beta_s e^{c_{\rm lo}+\varepsilon_k}I_d,
 &L=d-1,d,\\[3pt]
\beta_s e^{c_{\rm hi}-\varepsilon_k}I_d
\preceq Q_{\partial,L}^{(s)}
\preceq\beta_s e^{c_{\rm hi}+\varepsilon_k}I_d,
 &L=q+d-1,q+d .
\end{array}\tag{OOQ34}
\]
These are inequalities on every full original remainder
coefficient vector.

Moreover
\[
\varepsilon_k=O_{\delta,\gamma}(q^{3/4}+k\log q)=o(q).
\tag{OOQ35}
\]
Indeed OOQ24 gives the first term $q^{3/4}$ because
$n\rho^2=n^{3/4}$; all remainder and diagonal-coordinate
costs are $O(d\log n)$. OOQ28 gives the low error.
The original OSP constant $\mathcal E_k$ is bounded at fixed
$\delta,\gamma$ on its explicit domain, and
$r_{\ell_s}=O(1/k)$ by OFV6--7.
The sine bounds in OFV12 give logarithmic cost $O(\log q)$.
Finally $Q_s-q=-d+\ell_s=O(k)$, the rates differ from
$\vartheta$ by $1/q$, and $W$ and its derivatives are bounded.
The mean-value formula in the displayed scalar definitions
therefore bounds both scalar discrepancies in OOQ33 by
$O(k\log q+\log q)$.
Every original mass remains in the common factor $\beta_s$.

\subsection{Relative operators and all fixed subquotients of this actual target}

The exact leading scalar difference is
\[
c_{\rm lo}(k)-c_{\rm hi}(k)
 =q\{2\log(4/\vartheta)-2+W(2,\vartheta)\}
 =\frac{q\mathfrak C_\partial}{2}.
\tag{OOQ36}
\]
For the second equality use OOQ5 at $\alpha=2$ and precisely
OFV29,
$\mathfrak C_\partial=4\log(4/\vartheta)-10
-2F(2,\vartheta)+4F_\alpha(2,\vartheta)$.
Equivalently it is the full OFV31 energy expression
$4\log(4/\pi)+2-2I_{2,\pi}$.
Thus this constant was obtained from the operator calculation
with the same original source scale.

Let
\[
\begin{split}
\mathcal R_0&=(Q_{\partial,q+d-1}^{(s)})^{-1/2}
 Q_{\partial,d-1}^{(s)}
 (Q_{\partial,q+d-1}^{(s)})^{-1/2},\\
\mathcal R_1&=(Q_{\partial,q+d}^{(s)})^{-1/2}
 Q_{\partial,d}^{(s)}
 (Q_{\partial,q+d}^{(s)})^{-1/2}.
\end{split}\tag{OOQ37}
\]
The positive square roots are the unique positive ones.
For every eigenvalue of either of these actual relative
operators, OOQ34 proves
\[
\left|\log\lambda-\frac{q\mathfrak C_\partial}{2}\right|
\le2\varepsilon_k=o(q).
\tag{OOQ38}
\]
This follows by positive congruence of the two scalar
inequalities; it does not infer eigenvalues from their product.
The positive operator
$\mathcal R_0^{1/2}\mathcal R_1\mathcal R_0^{1/2}$ has
all eigenvalue logarithms within $4\varepsilon_k$ of
$q\mathfrak C_\partial$, and its determinant is exactly the
original four-endpoint determinant ratio. No commutation
of the two relative operators is used or asserted.

For precision, the same result controls every specified
fixed restriction or quotient of this outer target, even
when its coefficient map depends on the original period
and on $k$. For an injective map $I:\mathbb C^t\to\mathbb C^d$,
use the actual restricted Grams $I^*Q_{\partial,L}^{(s)}I$.
For a surjective map $A:\mathbb C^d\to\mathbb C^t$,
use their attained minimum quotient Grams
\[
Q_{A,L}^{(s)}
 =(A(Q_{\partial,L}^{(s)})^{-1}A^*)^{-1}.
\tag{OOQ39}
\]
The restriction coefficient reference is $I^*I$ and the
quotient coefficient reference is $(AA^*)^{-1}$.
Congruence, or minimization on the identical fibres,
transports OOQ34 with exactly those reference forms.
Taking determinants and adding the four original signs
cancels the complete fixed reference determinant and
the factor $\beta_s^t$. Hence in both cases the exact
four-endpoint return $F_t^{(s)}$ satisfies
\[
|F_t^{(s)}-tq\mathfrak C_\partial|
\le4t\varepsilon_k.
\tag{OOQ40}
\]
For $t=O(k)$ this error is $o(kq)$, without any assumed
conditioning of $I$ or $A$: each is the same actual map
at all four endpoints, so its complete coefficient form
has been retained and cancels exactly.

The operator statements OOQ34 and OOQ38 concern the complete
CTV outer-divisor quotient, whose ideal is
$P_{\partial,k}\mathbb C[y]_{\le L-d}$. The current original
residual has the different, explicitly supplied CEP source
graph, containing its complete $\mathcal U_{A,L}$ columns,
the low-$v$ components, and all original $K$ lifts.
OOQ39--40 do not identify that graph with this outer ideal.
The present calculation supplies the full original outer
operator and its exact subquotient transport; the comparison
with the actual residual must use the complete CEP map and
its denominator. No value of the separate original
$\Phi_{k,s}$ is assigned by omitting that map.

% END CURRENT OOQ

\clearpage
% BEGIN CURRENT GIP
\section{The complete residual map on the original conductor graph}
\label{sec:GIP}
Retain every original coefficient and metric of CEP1--49 and FMC1--38.
In particular $m=1$, $k=4l+1\ge9$, $q=(k+1)^2$, $q'=(k-7)^2$,
$\Delta=q-q'$, $j=8k-32$, $r=8k-16$, and the original first
nonzero conductor moment is $0\le v\le80$. For each source
$s\in\{1,k\}$ and endpoint $N\in\{q-1,q,2q-1,2q\}$ put
$t=N-q+1$ and $d=N-v$. All coefficient maps below retain the actual
period matrix and the entire conductor $A=\sum_{a,b=0}^8a_{ab}e_{ab}^{(8)}$.

\subsection{The full graph denominator and its metric}
Use original orthonormal coordinates on both source polynomial spaces,
and define the coefficient space and injective multiplication map
\[
 X_N=\mathbb C^v\oplus\mathbb C[S']_{\le N-v-q'},\qquad
 D_N=\operatorname{diag}(I_v,M_{\chi_{k'}}),\qquad
 \dim X_N=t+\Delta.
 \tag{GIP1}
\]
The second summand of $X_N$ uses the increasing ordinary monomials
in $S'$; $M_{\chi_{k'}}$ includes their exact conversion to the
original lower orthonormal source coordinates. The first summand
is the original upper source's first $v$ orthonormal coordinates.
Thus $D_N$ identifies $X_N$ with $\mathbb C^v\oplus I'_N$,
where $I'_N=\chi_{k'}\mathbb C[S']_{\le N-v-q'}$.
Let $S_N$ be the complete invertible source map CEP28, and let
$U_N$ be the full upper source remainder in original primary
coordinates CEP21. The specified map to the original conductor
kernel and its residual observation are
\[
 \begin{gathered}
 E_N^{\rm graph}=U_NS_N^{-1}D_N:X_N\longrightarrow V_A,\\
 B_N^{\rm graph}
   =E_j^{\mathsf T}T_k^{\mathsf T}\mathsf S_k^{\mathsf T}
                      E_N^{\rm graph}:X_N\longrightarrow\mathbb C^j.
 \end{gathered}\tag{GIP2}
\]
Here $T_k,\mathsf S_k,E_j$ are the complete FMC and OCF maps,
not new sections. CEP31 proves that $E_N^{\rm graph}$ is onto
$V_A$, since $S_N^{-1}D_NX_N=U_N^{-1}(V_A)$.
FMC30 proves that the second map is exactly
$\Xi_k^{\rm fix}E_N^{\rm graph}$ and is onto.

For the $t$ ordinary monomials $h$ of degree at most $N-q$ let
$Z_{\chi,N}h$ be the original upper orthonormal coefficients of
$\chi_kh$. Define, with the complete products and shifts of CEP41--42,
\[
 \begin{gathered}
 \mathcal Z_{A,N}h=
   \begin{pmatrix}P_{<v}Z_{\chi,N}h\\\mathcal U_{A,N-q}h\end{pmatrix},\\
 (\mathcal U_{A,N-q}h)(S')
       =\sum_{a,b=0}^8a_{ab}d_{ab}(S')h(S'+b_{ab}),\qquad
 D_N\mathcal Z_{A,N}=S_NZ_{\chi,N}.
 \end{gathered}\tag{GIP3}
\]
At $t=0$ this is the empty column matrix. CEP42 proves the equality
including the nonzero leading coefficient
$\operatorname{lc}(h)\binom{q+\deg h}{v}\mu_v$ for nonzero $h$. Consequently
$\mathcal Z_{A,N}$ has rank $t$, and
$\ker E_N^{\rm graph}=\operatorname{im}\mathcal Z_{A,N}$:
a vector maps to zero precisely when its $S_N^{-1}D_N$ image
is a full source relation $\chi_kh$.

The entire common kernel is also retained. Put $I_k=J_kE_r$ as
in FMC30, and use its exact original minimum source lift
$Y_{K,N}=U_N^*G_NI_k$. Its columns belong to
$U_N^{-1}(V_A)$, and hence to the domain of the inverse of $D_N$
on its indicated image after application of $S_N$. Define
\[
 W_{K,N}=D_N^{-1}S_NY_{K,N},\qquad
 \ker B_N^{\rm graph}
    =\operatorname{im}[\mathcal Z_{A,N},W_{K,N}],\qquad
 \operatorname{rank}[\mathcal Z_{A,N},W_{K,N}]=t+r.
 \tag{GIP4}
\]
Indeed $E_N^{\rm graph}W_{K,N}=I_k$ because $U_NU_N^*=G_N^{-1}$.
A relation among the two displayed sets, after applying
$E_N^{\rm graph}$, is therefore a relation among the independent
columns of $I_k$, so its $W_{K,N}$ coefficients vanish. The graph
columns are independent as proved above. Conversely a vector killed
by $B_N^{\rm graph}$ has image in $K$ and subtracting the
corresponding $W_{K,N}$ combination leaves a graph relation.
This proves both equality and rank, including every original
low-degree graph component and every common-kernel column.

There are two precisely specified positive metrics on $X_N$:
\[
 H_N^{\rm actual}=D_N^*S_N^{-*}S_N^{-1}D_N,\qquad
 H_N^{\rm graph}=D_N^*D_N
       =\operatorname{diag}(I_v,M_{\chi_{k'}}^*M_{\chi_{k'}}).
 \tag{GIP5}
\]
The first is the exact pullback of the unchanged original source
norm, and the second is the pullback of the original direct-sum
target norm of CEP28. Positivity follows from injectivity of $D_N$
and invertibility of $S_N$. In particular the second block, after
the literal substitution $S'=c'+iy$, integrates
$|f(c'+iy)|^2|P_{k'}(y)|^2dm_s(y)$: the factor $i^{q'}$
in $\chi_{k'}(c'+iy)=i^{q'}P_{k'}(y)$ has its complete modulus one.
The coordinate substitution in $f$ is retained in this statement.

Both residual quotient metrics consequently have full finite formulas
\[
 \begin{gathered}
 Q_{\Xi,N}^{\rm fix}
   =[B_N^{\rm graph}(H_N^{\rm actual})^{-1}(B_N^{\rm graph})^*]^{-1},\\
 \widetilde Q_{\Xi,N}
   =[B_N^{\rm graph}(H_N^{\rm graph})^{-1}(B_N^{\rm graph})^*]^{-1},\qquad
 |\log\det\widetilde Q_{\Xi,N}-\log\det Q_{\Xi,N}^{\rm fix}|
       \le2E_N^{\rm CEP}.
 \end{gathered}\tag{GIP6}
\]
Here $E_N^{\rm CEP}$ is the entire constant CEP35. For either
positive $H$ and onto $B$, the vector
$H^{-1}B^*(BH^{-1}B^*)^{-1}y$ maps to $y$ and is $H$-orthogonal
to $\ker B$. Adding any other lift adds a perpendicular kernel
vector and proves the inverse-Gram formula. The first minimum
is the original minimum over all upper source lifts of a primary
vector with residual $y$, by GIP2 and GIP5; this proves its equality
with FMC32. For the comparison, the actual source map $S_N$
induces the identity of these two quotient coordinate targets.
Minimal lifts followed by orthogonal projection off the entire
kernel GIP4 have exterior norm at most $\|\bigwedge^jS_N\|$;
the inverse has bound $\|\bigwedge^jS_N^{-1}\|$. Orthogonal
projection is an exterior contraction because its singular values
are zero or one. Squaring the two volume inequalities and using
CEP34--35 proves the final bound. Thus all four determinant errors
sum to at most $2\sum_NE_N^{\rm CEP}=o(kq)$ for the fixed
actual period, without identifying any two different relation spaces.

\subsection{The exact natural map to the CTV outer divisor}
Let $\mathcal I_k=\{4,\ldots,k-4\}^2$ and
$\mathcal O_k=\{0,\ldots,k\}^2\setminus\mathcal I_k$.
Let $E_{\rm in},E_{\rm out}$ insert these coordinates into the
original primary order. In the centred $y$ coordinate retain the
full outer Vandermonde
$(V_{\rm out})_{(a,b),h}=\omega_{ab}^h$ for $0\le h<\Delta$
and the diagonal matrix
$D_{\rm in,out}=\operatorname{diag}(P_{k'}(\omega_{ab}))_{\mathcal O_k}$.
Both are invertible: all the original roots are distinct and
no outer root is an inner root. Define
\[
 \begin{gathered}
 \rho_A=V_{\rm out}^{-1}D_{\rm in,out}^{-1}E_{\rm out}^{\mathsf T}|_{V_A}
      :V_A\longrightarrow\mathbb C[y]/(P_{\partial,k}),\\
 C_{\rm in}=CE_{\rm in},\qquad C_{\rm out}=CE_{\rm out},\qquad
 (C_{\rm in})_{(a,b),(i,j)}=a_{i+4-a,j+4-b}\quad(0\le a,b,i,j\le k'),
 \end{gathered}\tag{GIP7}
\]
where an out-of-range coefficient of $A$ is zero. This is the
literal map which first takes outer values and then divides by the
inner root factor from CTV10; on the graph its matrix is
$\rho_AE_N^{\rm graph}$. A source relation $\chi_kh$ is killed
because its outer evaluations vanish. Thus GIP7 is independent
of every source lift and retains exactly the CTV remainder frame.

There is an exact sequence with an entirely explicit obstruction:
\[
 \begin{gathered}
 0\longrightarrow\ker C_{\rm in}
   \xrightarrow{E_{\rm in}}V_A
   \xrightarrow{\rho_A}\mathbb C^\Delta
   \xrightarrow{\partial_A}\operatorname{coker}C_{\rm in}
   \longrightarrow0,\\
 \partial_A f=[C_{\rm out}D_{\rm in,out}V_{\rm out}f],\qquad
 \dim\ker\rho_A=\dim\operatorname{coker}\rho_A
                           =\dim\ker C_{\rm in}.
 \end{gathered}\tag{GIP8}
\]
To prove it, write $x=E_{\rm in}z+E_{\rm out}w$.
Membership in $V_A$ says exactly $C_{\rm in}z+C_{\rm out}w=0$,
while $\rho_Ax=f$ says $w=D_{\rm in,out}V_{\rm out}f$.
These two equations prove exactness at the first three terms.
Surjectivity of $C$ proves
$\operatorname{im}C_{\rm in}+\operatorname{im}C_{\rm out}
=\mathbb C^{q'}$, and the outer diagonal and Vandermonde are
invertible, so $\partial_A$ is onto. Since $C_{\rm in}$ is
square its kernel and cokernel have equal dimension, proving the
dimension assertion. Every possible defect is therefore the
nullity of this specific actual finite convolution section.
The scalar first-moment index $v$ does not by itself evaluate
that matrix or its nullity.

The original residual observation is retained even when this defect
is nonzero. Put $W_A=\Xi_k^{\rm fix}E_{\rm in}\ker C_{\rm in}$.
Then the exact descended residual map and its kernel are
\[
 \overline\Xi_A:\ker\partial_A\longrightarrow\mathbb C^j/W_A,
 \qquad \rho_Ax\longmapsto[\Xi_k^{\rm fix}x],\qquad
 \ker\overline\Xi_A=\rho_A(K).
 \tag{GIP9}
\]
Two lifts differ by $E_{\rm in}\ker C_{\rm in}$ by GIP8,
so the displayed rule is well defined. It is onto because
$\Xi_k^{\rm fix}$ is onto. If its value is zero, subtract an
element of $E_{\rm in}\ker C_{\rm in}$ with the same residual;
the remaining lift lies in $K$ and has unchanged outer image.
This proves the kernel equality; its converse follows by direct
application of $\Xi_k^{\rm fix}$. Equations GIP3--9 thus specify
the full denominator, retained metrics, and exact obstruction to
using the natural outer map for the same residual classes.

% END CURRENT GIP

\clearpage
% BEGIN CURRENT MCE
\section{The original mixed return and its same-class action control}
\label{sec:MCE}

This calculation retains the original arithmetic source, its actual
cyclic inclusion, and the complete period observation. It computes the
action control on those objects and the exact degree-step information
carried by the signed mixed return. The original full-multiplicity
control calculation is stated with its full multiplicity; the subsequent
evaluated outer-return substitution uses precisely the proved $m=1$
domain.

\subsection{The original arithmetic class and all three source metrics}
Fix the stipulated quartet $\rho=1/2+\delta+i\gamma$ and its
reflected and conjugate centres, with $0<\delta<1/2$, $\gamma>2$,
and full common multiplicity $m\ge1$. Retain
\[
 \begin{gathered}
 h(s)=\prod_{\lambda\in\{\rho,\bar\rho,1-\rho,1-\bar\rho\}}(s-\lambda)^m,
 \quad g=2\xi,\quad v_h=g/h,\quad
 w_h(y)=|v_h(1/2+iy)|^2/(2\pi),\\
 k=4l+1\ge9,\quad c=k/2,\quad e=1+k(m-1),\quad q=e(k+1)^2,\\
 \beta_{ab}=c+(2a-k)\delta+i(2b-k)\gamma,\qquad
 \chi(S)=\prod_{a,b=0}^k(S-\beta_{ab})^e,
 \quad E=\mathbb C[S]/(\chi),\quad A=M_S.
 \end{gathered}\tag{MCE1}
\]
The three source measures, all in the literal coordinate $S=c+iy$,
are the arithmetic measure $w_h^{*k}(y)\,dy$ and the original
Gamma measures $dm_s(y)$ of CTV1 for $s=1,k$. Their masses are
respectively $\mu_h^k$ and $(2\pi)^{s/2}$. For a source
$a\in\{\mathrm{ar},1,k\}$ let $p_n^a(S)$ be its monic orthogonal
polynomial, $\omega_n^a$ its full squared norm, and $b_n^a=[p_n^a]_\chi$
in the fixed original coefficient frame of $E$. For the Gamma sources
these are explicitly
\[
 \begin{gathered}
 p_n^s(S)=i^n\,\mathfrak p_n^s((S-c)/i),\qquad
 \mathfrak p_0^s=1,\quad\mathfrak p_1^s=y,\\
 \mathfrak p_{n+1}^s=y\mathfrak p_n^s-n(s/2+n-1)\mathfrak p_{n-1}^s,
 \qquad \omega_n^s=(2\pi)^{s/2}n!(s/2)_n .
 \end{gathered}\tag{MCE2}
\]
Thus the $i^n$ converting FMC29's real-coordinate polynomial to
the monic polynomial in $S$ is retained. In primary coordinates
for $m=1$, its entries are exactly
$i^n\mathfrak p_n^s(\omega_{ab})$. For higher multiplicity the
entries are the complete derivatives at each centre, with the
factorials of the original Hermite coordinate map; none is replaced
by a value-only coordinate.

For every $N\ge q-1$ define the complete covariance and minimum Gram
\[
 R_N^a=\sum_{n=0}^N b_n^a(b_n^a)^*/\omega_n^a,
 \qquad G_N^a=(R_N^a)^{-1},\qquad V_N^a=\det G_N^a.
 \tag{MCE3}
\]
The first $q$ columns are the monic triangular remainder frame and
are independent. Positivity of all norms therefore proves $R_N^a>0$.
In the orthogonal source coefficient frame the minimum section is
$\operatorname{diag}(\omega_0^a,\ldots,\omega_N^a)^{-1}
[b_0^a,\ldots,b_N^a]^*G_N^a$.
Multiplication verifies that it is a section and is orthogonal to
the full relation kernel. The squared norm of every other section
value adds the squared norm of a perpendicular relation. This proves
MCE3's minimum interpretation with all masses intact.

The original arithmetic inclusion is
\[
 \eta:E\hookrightarrow E_h^{\otimes k},\qquad
 \eta[P]=\upsilon_h^{\otimes k}P(A_\Sigma)1,
 \quad\upsilon_h=j_h(v_h),\quad A_\Sigma=\sum_{i=1}^kM_{s_i},
 \qquad \eta A=A_\Sigma\eta.
 \tag{MCE4}
\]
Its injectivity follows from the exact annihilator $\chi$ and the
invertibility of the complete Taylor unit $\upsilon_h^{\otimes k}$.
For clarity the local nilpotent sum has exact order $e$: its
$(e-1)$-st power has coefficient
$[k(m-1)]!/[(m-1)!]^k$ at the product of the top local powers,
and every higher monomial has an exponent at least $m$. All pairs
$(a,b)$ occur among tensor choices and their centres are distinct.
The least common multiple of these local annihilators is exactly
MCE1. This proves the asserted annihilator and injectivity.
The commutation identity follows by polynomial multiplication and
commutation with the retained unit. The supplied source identity
$J^{(k)}\mathcal V_{h,k}=\eta\operatorname{rem}_\chi$ therefore
carries each calculation below to this same arithmetic class.
On an original present support label the map is
$(\lambda,x)\mapsto(\lambda,\eta x)$, and external absence retains
its original image; no coefficient calculation changes that label.

\subsection{The exact rank-two control from the original recurrence}
For this subsection fix any one source and suppress its superscript.
Positivity of the original measure gives real monic orthogonal
polynomials $Q_n(y)=i^{-n}p_n(c+iy)$. Pairing $yQ_n$ with lower
degrees proves
\[
 Sp_n=p_{n+1}+(c+i\alpha_n)p_n
             -\frac{\omega_n}{\omega_{n-1}}p_{n-1},\qquad
 \alpha_n\in\mathbb R,
 \tag{MCE5}
\]
with the lower term absent at $n=0$. Indeed degrees below $n-1$
are orthogonal by self-adjoint multiplication by the real variable;
the coefficient at $Q_{n+1}$ is one by monicity, and pairing with
$Q_{n-1}$ gives $\omega_n/\omega_{n-1}$. For the Gamma measures
$\alpha_n=0$ and MCE5 is precisely MCE2 with its phases.

Reducing MCE5 modulo $\chi$ and inserting it into MCE3 gives
\[
 \begin{aligned}
 AR_N+R_NA^*-kR_N
       &=\frac{b_{N+1}b_N^*+b_Nb_{N+1}^*}{\omega_N},\\
 D_N:=A^*G_N+G_NA-kG_N
       &=\frac{G_Nb_{N+1}b_N^*G_N+G_Nb_Nb_{N+1}^*G_N}{\omega_N}.
 \end{aligned}\tag{MCE6}
\]
For the first equality the diagonal coefficient is
$c+i\alpha_n+c-i\alpha_n-k=0$. At every internal adjacent
pair the remaining coefficient is
$\omega_n^{-1}-(\omega_{n+1}/\omega_n)\omega_{n+1}^{-1}=0$.
Only the two displayed top columns survive. Multiplication on both
sides by $G_N$ proves the second equality, including its signs.

The dagger-stable polynomial obeys
$\overline\chi(k-S)=(-1)^q\chi(S)$, so
$\widehat\chi(y)=i^{-q}\chi(c+iy)$ is real monic. The exact
substitution $[P(S)]\mapsto[P(c+iy)]$ transports $G_N$ by its
full inverse congruence; in that real remainder frame the covariance
and Gram have real entries. Each column has phase $i^n$.
Consequently
\[
 \begin{gathered}
 b_N^*G_Nb_{N+1}=i\omega_N\phi_N,\quad\phi_N\in\mathbb R,\\
 \mathcal H_N=G_N^{-1/2}D_NG_N^{-1/2}
        =\frac{yx^*+xy^*}{\omega_N},\\
 x=G_N^{1/2}b_N,\quad y=G_N^{1/2}b_{N+1}.
 \end{gathered}\tag{MCE7}
\]
Thus $\mathcal H_N$ is Hermitian of rank at most two and trace
zero. Its square has trace
$2(\|x\|^2\|y\|^2-|x^*y|^2)/\omega_N^2$, because
$x^*y=i\omega_N\phi_N$. Its two possibly nonzero eigenvalues
are therefore $\epsilon_N,-\epsilon_N$, where
\[
 \epsilon_N^2
   =\frac{\|x\|^2\|y\|^2-|x^*y|^2}{\omega_N^2}
   =\frac{\|x\wedge y\|^2}{\omega_N^2}.
 \tag{MCE8}
\]
The last equality is the determinant of the two-vector Gram,
proved by expanding in an orthonormal basis. This defines the
literal control allowance, including the case when the rank is zero.

\subsection{The same nonzero arithmetic exterior class}
Let $E_+$ be the full sum of the primary blocks with $a>k/2$,
and put $t_+=e(k+1)^2/2$. Its dimension is $t_+$ because
$k$ is odd. A completely specified primary basis is
$e_\beta(S-\beta)^d$, $0\le d<e$, over those centres, where
\[
 \chi_\beta=\chi/(S-\beta)^e,\qquad
 e_\beta=\chi_\beta\sum_{u=0}^{e-1}
       \frac{(1/\chi_\beta)^{(u)}(\beta)}{u!}(S-\beta)^u
                                  \pmod\chi.
 \tag{MCE9}
\]
The Taylor product is one modulo $(S-\beta)^e$ and vanishes at
the other complete primary factors. This proves the primary
idempotent property and the stated basis. Multiplication by $S$
has diagonal $\beta$ and the exact next-power nilpotent on each
block. Thus $E_+$ is $A$-invariant, and its ordered top wedge
$\xi_+$ is a nonzero eigenvector for the additive exterior action
$A^{[t_+]}$, with eigenvalue $\sum_{a>k/2,b}e\beta_{ab}$.
The map $\bigwedge^{t_+}\eta$ is injective: extend an independent
set in its domain to a basis and use injectivity of $\eta$ on
the resulting independent images. Hence
$\bigwedge^{t_+}\eta(\xi_+)\ne0$ is the same surviving arithmetic
exterior class. MCE4 intertwines the additive exterior actions
term by term, retaining all its unit and nilpotent coefficients.

On $\bigwedge^t(E,G_N)$ the exact relative control is the additive
exterior action of $\mathcal H_N$. In an orthonormal basis this
follows by applying $A$ to each wedge factor and pairing the
result, so its eigenvalues are sums of $t$ distinct eigenvalues
of $\mathcal H_N$. For $1\le t\le q-1$ they lie in
$[-\epsilon_N,\epsilon_N]$, since the only nonzero entries in
the original list are $\epsilon_N,-\epsilon_N$. Applying this
bound to $\xi_+$, or taking the trace of the control on its
invariant primary space, proves
\[
 \mathcal L_{h,k}
  :=2\delta e(k+1)\left\lfloor\frac{(k+1)^2}{4}\right\rfloor
  =2\Re\operatorname{Tr}(A|_{E_+})-kt_+
  \le\epsilon_N,
 \qquad \mathcal L_{h,k}/q\ge\delta k/2>0.
 \tag{MCE10}
\]
The sum of the positive integers $2a-k$ is the displayed floor,
by summing the positive odd or even progression; for the present
odd $k$ it is $(k+1)^2/4$. Multiplication by the $k+1$ imaginary
centres and full length $e$ gives the equality. The last bound
uses $\lfloor(k+1)^2/4\rfloor\ge k(k+1)/4$. The exterior rank
here is $t_+$, and this exact inclusion and action determine its
relation to the arithmetic class; it is not assigned the dimension
of the unrelated period kernel.

\subsection{Complete kernel, boundary and mixed pieces of the control}
Fix the original observation $\Lambda:E\twoheadrightarrow B$
in any one fixed coordinate frame of its actual image, and its
actual kernel inclusion $I:\mathbb C^{r_K}\hookrightarrow E$.
For $m=1$ on the proved period locus, $r_K=8k-16$; in the present
algebraic formulas it remains the actual rank for the stipulated
multiplicity. Define at each source and cutoff
\[
 \begin{gathered}
 H_N=I^*G_NI,\quad Q_N=(\Lambda R_N\Lambda^*)^{-1},\quad
 L_N=R_N\Lambda^*Q_N,\quad
 K_N^{\rm coord}=H_N^{-1}I^*G_N,\\
 \Lambda L_N=I_B,\quad I^*G_NL_N=0,\quad
 I K_N^{\rm coord}+L_N\Lambda=I_E,\\
 F_N=[I,L_N],\quad F_N^*G_NF_N=\operatorname{diag}(H_N,Q_N).
 \end{gathered}\tag{MCE11}
\]
The first section identity follows by multiplication. The second
follows from $\Lambda I=0$. The two summands in the third are
the $G_N$-orthogonal projections onto $K$ and its complement:
the first fixes $I\mathbb C^{r_K}$; subtracting the second from
any vector lies in $K$ and has its stated kernel coordinate.
This proves the third identity and the full Gram equality.

The complete original action in this frame is
\[
 F_N^{-1}AF_N=
 \begin{pmatrix}
 K_N^{\rm coord}AI&K_N^{\rm coord}AL_N\\
 \Lambda AI&\Lambda AL_N
 \end{pmatrix}=:\begin{pmatrix}A_{KK}&A_{KB}\\A_{BK}&A_{BB}\end{pmatrix}.
 \tag{MCE12}
\]
In particular $A_{BK}=\Lambda AI$ is the actual action defect;
no invariant-kernel premise is used. The exact boundary relation is
$\Lambda A-A_{BB}\Lambda=A_{BK}K_N^{\rm coord}$, obtained by
inserting the third identity in MCE11. It is the metric-section
version of OPG22; replacing its fixed section by $L_N$ changes
both displayed blocks by precisely this identity.
The three complete control blocks are
\[
 \begin{split}
 I^*D_NI&=A_{KK}^*H_N+H_NA_{KK}-kH_N,\\
 L_N^*D_NL_N&=A_{BB}^*Q_N+Q_NA_{BB}-kQ_N,\\
 I^*D_NL_N&=A_{BK}^*Q_N+H_NA_{KB}.
 \end{split}\tag{MCE13}
\]
Multiplication of MCE12 by the block Gram MCE11 proves each entry.
For an actual eigenvector $z=Iu+L_Nb$ of eigenvalue $\lambda$,
the two action equations are
$A_{KK}u+A_{KB}b=\lambda u$ and
$A_{BK}u+A_{BB}b=\lambda b$. Consequently its literal control
pairing is
\[
 (2\Re\lambda-k)(u^*H_Nu+b^*Q_Nb)
 =u^*I^*D_NIu+2\Re(u^*I^*D_NL_Nb)+b^*L_N^*D_NL_Nb.
 \tag{MCE14}
\]
This follows either from the two action equations and MCE13 or
directly from $z^*D_Nz$; it preserves the entire mixed pairing.
For generalized primary vectors the full block matrix MCE12,
rather than only this eigenvector equality, remains the action.

The same splitting gives a quantitative exterior identity for the
actual allowance itself. Put
\[
 \begin{gathered}
 x_K=H_N^{-1/2}I^*G_Nb_N,\quad x_B=Q_N^{1/2}\Lambda b_N,\\
 y_K=H_N^{-1/2}I^*G_Nb_{N+1},\quad y_B=Q_N^{1/2}\Lambda b_{N+1},\\
 a_K=\|x_K\|^2,\ a_B=\|x_B\|^2,\ d_K=\|y_K\|^2,\ d_B=\|y_B\|^2,\\
 z_K=x_K^*y_K,\quad z_B=x_B^*y_B.
 \end{gathered}\tag{MCE15}
\]
The map $G_N^{1/2}F_N\operatorname{diag}(H_N,Q_N)^{-1/2}$
is unitary by MCE11. Its inverse sends MCE7's $x,y$ to these
respective ordered pairs. Therefore
\[
 \begin{gathered}
 z_K+z_B=i\omega_N\phi_N,\\
 \omega_N^2\epsilon_N^2
   =\|x_K\wedge y_K\|^2+\|x_B\wedge y_B\|^2
                     +\|x_K\otimes y_B-y_K\otimes x_B\|^2\\
   =(a_Kd_K-|z_K|^2)+(a_Bd_B-|z_B|^2)
                       +a_Kd_B+a_Bd_K-2\Re(z_K\overline{z_B}).
 \end{gathered}\tag{MCE16}
\]
To prove the second line expand $x\wedge y$ in the orthogonal
direct sum $\bigwedge^2K\oplus(K\otimes B)\oplus\bigwedge^2B$.
Its three components are exactly those displayed, including the
minus sign in the mixed tensor. Orthogonality adds their squared
norms. Expansion of each two-vector Gram and of the mixed tensor
inner product proves the last line. Each squared norm in the
second line is nonnegative. The individual $z_K,z_B$ need not be
purely imaginary; their actual sum is the pairing in MCE7.

For completeness the same surviving exterior class also has a full
mixed coordinate map. If $I_+$ is the matrix of MCE9's actual
positive-primary basis, put
$X_+=K_N^{\rm coord}I_+$ and $Y_+=\Lambda I_+$. Then
\[
 I_+=IX_++L_NY_+,\quad
 G_{+,N}=X_+^*H_NX_++Y_+^*Q_NY_+>0,
 \quad
 \mathcal L_{h,k}
   =\operatorname{Tr}\bigl(G_{+,N}^{-1}I_+^*D_NI_+\bigr).
 \tag{MCE17}
\]
The first two equalities are MCE11. Invariance of $E_+$ says
$AI_+=I_+A_+$ with the full primary matrix of MCE9. Thus
$I_+^*D_NI_+=A_+^*G_{+,N}+G_{+,N}A_+-kG_{+,N}$;
taking the trace after multiplication by $G_{+,N}^{-1}$ proves
the last identity by MCE10. Its expansion using MCE13 retains
$X_+^*I^*D_NL_NY_+$ and its adjoint, as well as both diagonal
blocks. The induced map on $\xi_+$ is the literal wedge expansion
of the columns $IX_++L_NY_+$, with the usual permutation signs
when kernel factors are moved before boundary factors. This is
an invertible mixed coordinate description of the same vector
before applying the injective exterior arithmetic map MCE4.

\subsection{Degree-step volumes determine the squared component sizes}
Still fixing one original source, put for $N\ge q-1$
\[
 \kappa_N=\log\frac{\det H_N}{\det H_{N+1}},\qquad
 \beta_N^B=\log\frac{\det Q_N}{\det Q_{N+1}},\qquad
 \tau_N=\log(V_N/V_{N+1}).
 \tag{MCE18}
\]
These are nonnegative. Inserting the next exact covariance column
and multiplying verifies
$G_{N+1}=G_N-G_Nb_{N+1}b_{N+1}^*G_N/
(\omega_{N+1}+b_{N+1}^*G_Nb_{N+1})$.
The determinant of $I+uv^*$ is $1+v^*u$, proved by choosing
a basis with first vector $u$ or by two-block elimination.
Applying this identity first to the boundary covariance and then
to the restricted kernel metric proves
\[
 \begin{gathered}
 e^{\beta_N^B}=1+d_B/\omega_{N+1},\qquad
 e^{\kappa_N}=1+d_K/(\omega_{N+1}+d_B),\qquad
 \tau_N=\kappa_N+\beta_N^B,\\
 d_B/\omega_{N+1}=e^{\beta_N^B}-1,\qquad
 d_K/\omega_{N+1}=e^{\beta_N^B}(e^{\kappa_N}-1).
 \end{gathered}\tag{MCE19}
\]
All denominators are strictly positive. For the boundary step,
\[
 Q_{N+1}^{-1}=Q_N^{-1}
       +(\Lambda b_{N+1})(\Lambda b_{N+1})^*/\omega_{N+1}.
\]
For the kernel step the determinant ratio from the displayed
$G_{N+1}$ is
$1-d_K/(\omega_{N+1}+d_K+d_B)$, whose reciprocal is the stated
formula. Multiplying the two reciprocal ratios gives the complete
covariance ratio and proves the third equality.

For $N\ge q$ the preceding total and boundary covariances are
both positive definite. Subtracting their last column instead of
adding the next gives
\[
 a_B/\omega_N=1-e^{-\beta_{N-1}^B},\qquad
 a_K/\omega_N=e^{-\beta_{N-1}^B}(1-e^{-\kappa_{N-1}}).
 \tag{MCE20}
\]
Indeed the determinant ratio of the previous boundary covariance
to the current one is $1-a_B/\omega_N=e^{-\beta_{N-1}^B}$.
The identical total formula gives
$1-(a_K+a_B)/\omega_N=e^{-\tau_{N-1}}$.
Subtracting and using MCE19 proves MCE20. At the first degree
$N=q-1$ the total pairing is exactly $a_K+a_B=\omega_{q-1}$,
because the first $q$ monic remainder columns are an orthogonal
quotient basis. Its two components are kept as MCE15; no inverse
at $q-2$ is taken or assumed. Thus MCE16, MCE19 and MCE20 express
the actual mixed exterior allowance using the complete consecutive
kernel and boundary returns and their actual cross-pairings.

\subsection{The signed return gives an exact original control-radius budget}
Restore source superscripts. Define the actual signed one-degree term
\[
 \psi_N=\tau_N^{\rm ar}-(\beta_N^B)^{(k)},\qquad
 w_{q-1}=w_{2q-1}=1,\quad w_N=2\ (q\le N\le2q-2),
 \qquad\sum_Nw_N=2q.
 \tag{MCE21}
\]
Telescoping MCE18--19 gives, with every original endpoint,
\[
 \Psi_k:=F_K^{\rm ar}+F_B^{\rm ar}-F_B^{(k)}
       =\sum_{N=q-1}^{2q-1}w_N\psi_N,
 \qquad \tau_N^{\rm ar}=\psi_N+(\beta_N^B)^{(k)}.
 \tag{MCE22}
\]
To see the weights, the first window from $q-1$ to $2q-1$
includes degrees $q-1,\ldots,2q-2$, while the second window
from $q$ to $2q$ includes $q,\ldots,2q-1$. Their intersection
is counted twice and their two extreme degrees once.

For each individual source, MCE7--8 and MCE19--20 give the literal
radius formulas
\[
 \begin{split}
 (\epsilon_{q-1})^2+(\phi_{q-1})^2
    &=\frac{\omega_q}{\omega_{q-1}}(e^{\tau_{q-1}}-1),\\
 \epsilon_N^2+\phi_N^2
    &=\frac{\omega_{N+1}}{\omega_N}
       (1-e^{-\tau_{N-1}})(e^{\tau_N}-1),\qquad N\ge q.
 \end{split}\tag{MCE23}
\]
The first uses exactly $\|x\|^2=\omega_{q-1}$ and the next-column
formula for $\|y\|^2$; the second uses both total formulas already
proved. Subtracting $\phi_N^2=|x^*y|^2/\omega_N^2$ recovers the
actual $\epsilon_N^2$. In particular substitution of
$\tau_N^{\rm ar}=\psi_N+(\beta_N^B)^{(k)}$ in MCE23 is a
pointwise, same-arithmetic-class map from the signed degree data to
the original control. It retains the complete reference boundary
increment and arithmetic phase rather than assigning them zero.

Put $\delta_j^{\rm ar}=V_j^{\rm ar}/V_{j-1}^{\rm ar}$ for $j\ge q$.
The positive lower bound MCE10 and MCE23 imply
$0<\delta_j^{\rm ar}<1$ at every index used next. Define the full
arithmetic monic-norm window and contraction penalties
\[
 \begin{gathered}
 \mathcal W_k^{\rm ar}
     =\log\frac{\omega_{2q-1}^{\rm ar}\omega_{2q}^{\rm ar}}
                       {\omega_{q-1}^{\rm ar}\omega_q^{\rm ar}},\\
 \mathcal P_-=-\log(1-\delta_{2q-1}^{\rm ar})
                  -2\sum_{j=q}^{2q-2}\log(1-\delta_j^{\rm ar}),\\
 \mathcal P_+=-\log(1-\delta_q^{\rm ar})-\log(1-\delta_{2q}^{\rm ar})
                  -2\sum_{j=q+1}^{2q-1}\log(1-\delta_j^{\rm ar}).
 \end{gathered}\tag{MCE24}
\]
They are all finite, and the two penalties are nonnegative.
Multiplying MCE23 across the first window gives
\[
 \prod_{N=q-1}^{2q-2}(\epsilon_N^2+\phi_N^2)
 =\frac{\omega_{2q-1}V_{q-1}}{\omega_{q-1}V_{2q-1}}
   (1-\delta_{2q-1})\prod_{j=q}^{2q-2}(1-\delta_j)^2.
\]
Across the second it gives
\[
 \prod_{N=q}^{2q-1}(\epsilon_N^2+\phi_N^2)
 =\frac{\omega_{2q}V_q}{\omega_qV_{2q}}
   (1-\delta_q)(1-\delta_{2q})
                      \prod_{j=q+1}^{2q-1}(1-\delta_j)^2.
\]
For these two displays all quantities are arithmetic. To verify them,
write $e^{\tau_N}-1=(1-\delta_{N+1})/\delta_{N+1}$.
The monic norm ratios telescope, the reciprocal $\delta$ factors
give the indicated quotient-volume ratio, every internal
$(1-\delta_j)$ occurs twice, and each stated extreme occurs once.
Taking logarithms and using MCE22 proves the exact bridge
\[
 \boxed{
 \sum_{N=q-1}^{2q-1}w_N
       \log\bigl((\epsilon_N^{\rm ar})^2+(\phi_N^{\rm ar})^2\bigr)
 =\Psi_k+F_B^{(k)}+\mathcal W_k^{\rm ar}-\mathcal P_--\mathcal P_+ .}
 \tag{MCE25}
\]
It is an equality on the unchanged arithmetic packet. Every
allowance on its left is precisely MCE6--10's control allowance.

MCE10 now supplies a concrete necessary inequality for those same
data, retaining its full phase contribution:
\[
 \begin{gathered}
 \Psi_k+F_B^{(k)}+\mathcal W_k^{\rm ar}-\mathcal P_--\mathcal P_+
 \ \ge\ 4q\log\mathcal L_{h,k}
       +\sum_Nw_N\log\left(1+
                   \frac{(\phi_N^{\rm ar})^2}{\mathcal L_{h,k}^2}\right),\\
 \min_{q-1\le N\le2q-1}\epsilon_N^{\rm ar}
 \ \le\ \exp\left(
   \frac{\Psi_k+F_B^{(k)}+\mathcal W_k^{\rm ar}-\mathcal P_--\mathcal P_+}
              {4q}\right).
 \end{gathered}\tag{MCE26}
\]
For the first line apply
$\epsilon_N^{\rm ar}\ge\mathcal L_{h,k}$ to each positive
radius, write
$\log(\mathcal L_{h,k}^2+\phi^2)=2\log\mathcal L_{h,k}
+\log(1+\phi^2/\mathcal L_{h,k}^2)$, and use
$\sum w_N=2q$. For the second, the minimum of the allowances
is at most the weighted geometric mean of their larger radii;
take half of the logarithmic mean in MCE25. This proves both
inequalities without assuming any desired upper control estimate.

Finally specialize only the evaluated return substitution to the
proved $m=1$ domain of OAR/OVR/OFV. Retain its exact value
$\Psi_k=\mathcal F_{\partial,k}^{(k)}-\Phi_{k,k}
+\mathcal H_k+\delta_k^\sigma+\mathfrak e_{k,k}$,
where $-e_{k,k}^-\le\mathfrak e_{k,k}\le e_{k,k}^+$ with the
entire OAR2 constants. Then MCE25 becomes
\[
 \begin{split}
 \sum_Nw_N\log\bigl((\epsilon_N^{\rm ar})^2+(\phi_N^{\rm ar})^2\bigr)
 ={}&\mathcal F_{\partial,k}^{(k)}-\Phi_{k,k}
       +\mathcal H_k+\delta_k^\sigma+\mathfrak e_{k,k}\\
    &+F_B^{(k)}+\mathcal W_k^{\rm ar}-\mathcal P_--\mathcal P_+ .
 \end{split}\tag{MCE27}
\]
Substituting the already evaluated OFV outer value
$\mathcal F_{\partial,k}^{(k)}
=\Delta q\mathcal C_{\partial}+\mathfrak r_k$
retains exactly the same equality, with its proved remainder
$\mathfrak r_k=O(q\log q)$ and the actual residual $\Phi_{k,k}$.
In particular its original arithmetic necessary inequality is
\[
 \begin{split}
 \Phi_{k,k}\le{}&\mathcal F_{\partial,k}^{(k)}
       +\mathcal H_k+\delta_k^\sigma+e_{k,k}^+
       +F_B^{(k)}+\mathcal W_k^{\rm ar}-\mathcal P_--\mathcal P_+\\
       &-4q\log\mathcal L_{h,k}
       -\sum_Nw_N\log\left(1+
                    \frac{(\phi_N^{\rm ar})^2}{\mathcal L_{h,k}^2}\right).
 \end{split}\tag{MCE28}
\]
This is MCE26 composed with the exact receiver and its upper error
endpoint; no new positivity premise is introduced. It shows exactly
which additional original boundary, norm, contraction and phase
quantities accompany the mixed-return estimate in the control budget.
The full pointwise control can instead be calculated by MCE15--16
and MCE19--23. In either form, the actual arithmetic action, its
nonzero exterior class, all kernel--boundary defects and the
original four endpoints remain explicit. No sign of $\Psi_k$ alone
has been substituted for a bound on $\epsilon_N^{\rm ar}$.

% END CURRENT MCE

\clearpage
% BEGIN CURRENT MCP
\section{The next marked pairing as an original action and source derivative}
\label{sec:MCP}
Retain MCE1--28 with the full original observation and all source
orders. Give the same source its existing real tilt
$dm_{a,\theta}(y)=e^{\theta y}dm_a(y)$ for real
$|\theta|<\pi/2$. In the arithmetic case $dm_a=w_h^{*k}(y)dy$.
All polynomial moments and their derivatives exist on compact
subintervals by the original exponential-moment bound. At
$\theta=0$ this is exactly the current source. The polynomial
algebra, $A=M_S$, the arithmetic inclusion, and $\Lambda$ are
independent of $\theta$.

In fixed monomial source coordinates let $\mathcal R_N(\theta)$
be the original minimum section of the quotient and let its source
Gram be $\mathsf H_N(\theta)$. The identities
$J_N\mathcal R_N=I$ and orthogonality to $\ker J_N$ imply
\[
 \partial_\theta G_N
   =\mathcal R_N^*(\partial_\theta\mathsf H_N)\mathcal R_N
   =\langle\mathcal R_N,\,y\mathcal R_N\rangle_{a,\theta}.
 \tag{MCP1}
\]
Indeed differentiating $J_N\mathcal R_N=I$ gives
$J_N\mathcal R_N'=0$. The two section-derivative terms in the
derivative of $\mathcal R_N^*\mathsf H_N\mathcal R_N$ therefore
vanish by that orthogonality. Differentiating the integral defining
$\mathsf H_N$ multiplies the integrand by the unchanged variable
$y$, proving the second equality with full source mass.

Here is its exact endpoint evaluation. In the fixed real-coordinate
quotient set $T=M_y$, $Q_n=i^{-n}p_n(c+iy)$ and $d_n=[Q_n]$.
The leading coefficient of the minimum lift of $x$ is
$d_N^*G_Nx/\omega_N$. This follows by expanding that lift in the
orthogonal $Q_n$: its $Q_N$ coefficient is its inner product
with $Q_N$ divided by $\omega_N$, and the section identity gives
that inner product $d_N^*G_Nx$. The polynomial column
\[
 y\mathcal R_N-\mathcal R_NT
       -(Q_{N+1}-\mathcal R_Nd_{N+1})d_N^*G_N/\omega_N
\]
has degree at most $N$ because its highest coefficient cancels,
and its remainder is zero. Pairing it with $\mathcal R_N$
kills it by the same orthogonality. The term $Q_{N+1}$ is
orthogonal to every degree-at-most-$N$ minimum lift. Thus
$G_N'=G_NT-G_Nd_{N+1}d_N^*G_N/\omega_N$ in this specified
real-coordinate frame. Returning through the unchanged substitution
$S=c+iy$, with $b_n=i^nd_n$ and $T=(A-cI)/i$, proves
\[
 G_N'=-iG_N(A-cI)
                   +iG_Nb_{N+1}b_N^*G_N/\omega_N.
 \tag{MCP2}
\]
Every $p_n,b_n,\omega_n$ in this formula is evaluated at the same
$\theta$. No derivative of these columns was inserted: the formula
comes from the exact section minimum and its polynomial remainder.

Use MCE11's $H_N,Q_N,L_N$ at that same parameter, and keep the
original kernel inclusion $I$ fixed. Then
$H_N'=I^*G_N'I$ and $Q_N'=L_N^*G_N'L_N$.
For the latter, differentiating $\Lambda L_N=I_B$ puts $L_N'$
in $K$, while $I^*G_NL_N=0$ kills both differentiated-section
terms. Taking the two traces in MCP2 therefore proves the exact
marked-pairing identities
\[
 \begin{gathered}
 \frac{z_K}{\omega_N}
   =\operatorname{Tr}A_{KK}-cr_K
                    -i\partial_\theta\log\det H_N,\\
 \frac{z_B}{\omega_N}
   =\operatorname{Tr}A_{BB}-c\dim B
                    -i\partial_\theta\log\det Q_N.
 \end{gathered}\tag{MCP3}
\]
For example the first trace of $G_N(A-cI)$ compressed by $I$
is $\operatorname{Tr}A_{KK}-cr_K$. The trace of its marked
rank-one term is precisely $z_K$ of MCE15, by cyclicity of trace.
The second calculation uses $L_N^*G_N=Q_N\Lambda$ and gives
exactly $z_B$. Rearranging
$\partial_\theta\log\det H_N=-i(\operatorname{Tr}A_{KK}-cr_K)
+iz_K/\omega_N$ proves the first line and fixes its sign;
the second line follows in the same manner. Adding them recovers
the full phase MCE7 and EP.9 because
$\operatorname{Tr}A_{KK}+\operatorname{Tr}A_{BB}=\operatorname{Tr}A$.

In particular the original kernel--boundary action is measured by
the real parts of these pairings:
\[
 \begin{gathered}
 2\Re(z_K/\omega_N)=2\Re\operatorname{Tr}A_{KK}-kr_K,\qquad
 2\Re(z_B/\omega_N)=2\Re\operatorname{Tr}A_{BB}-k\dim B,\\
 \frac{\|x_K\otimes y_B-y_K\otimes x_B\|^2}{\omega_N^2}
  =\frac{a_Kd_B+a_Bd_K}{\omega_N^2}
       -2\Re\left[
       (\alpha_K-i\partial_\theta\log\det H_N)
       \overline{(\alpha_B-i\partial_\theta\log\det Q_N)}\right],\\
 \alpha_K=\operatorname{Tr}A_{KK}-cr_K,\qquad
 \alpha_B=\operatorname{Tr}A_{BB}-c\dim B.
 \end{gathered}\tag{MCP4}
\]
The derivatives are real because the two Gram determinants are
strictly positive for real $\theta$. This proves the first line.
Substituting MCP3 into MCE16's exact mixed tensor norm proves the
second, with all phases and off-diagonal action effects retained.
Its two size terms are the explicit degree-return quantities in
MCE19--20. Thus the next original pointwise calculation is the
marked pair $z_K,z_B$, equivalently these actual action traces
and source derivatives. An outer quotient operator estimate must
be transported through the full GIP graph and actual marked
columns before it can bound these quantities. No derivative
estimate is inferred by differentiating an uncontrolled remainder
in a scalar asymptotic law.

% END CURRENT MCP

\clearpage
% BEGIN CURRENT RFX
\section{The exact original reflection and the retained mixed phase}
\label{sec:RFX}
Retain the full original quartet, its multiplicity $m$, the complete
unit $v_h=2\xi/h$, every tensor factor and all source masses of
MCE1--4. The argument concerns the original arithmetic source and
both original Gamma orders $s=1,k$. It also specifies the exact
reflection map on the real tilted family of MCP1--4. No observation
invariance is assumed.

\subsection{The entire arithmetic unit and the source involution}
The functional equation of the original completed function is
$\xi(1-s)=\xi(s)$. Reflection permutes the four factors of the
full quartet polynomial; each factor contributes its minus sign
with its full multiplicity, giving
\[
 h(1-s)=(-1)^{4m}h(s)=h(s),\qquad v_h(1-s)=v_h(s).
 \tag{RFX1}
\]
The second identity first holds off the finite zero set of $h$.
Both sides are the entire quotient already supplied by the original
arithmetic construction, so continuity at every removed centre proves
the same identity there, including all Taylor coefficients. At
$s=1/2+iy$ it implies
\[
 w_h(-y)=w_h(y),\qquad w_h^{*k}(-y)=w_h^{*k}(y).
 \tag{RFX2}
\]
For the second equality, in the convolution integral substitute
$y_i\mapsto-y_i$ in every integration variable and use the first
equality in every factor. The integrals exist by the original
finite-mass estimates; Tonelli's theorem applies to these nonnegative
integrands. Each Gamma density in OOQ2 is also even, since
$\overline{\Gamma(s/4+iy/2)}=\Gamma(s/4-iy/2)$ for the retained
positive real source order. The full source constants remain equal
on the two sides, rather than being divided out.

Let $c=k/2$ and let $\mathcal P_N$ consist of polynomials in the
literal variable $S$. Define the complex-linear involution
\[
 (\mathscr S_Np)(S)=p(k-S),\qquad
 (\mathscr S_N)_{ab}=(-1)^a\binom ba k^{b-a}
 \quad(0\le a\le b\le N),
 \tag{RFX3}
\]
and set every other matrix entry to zero. The displayed matrix is
the binomial expansion in the original increasing monomial frame;
substitution twice proves $\mathscr S_N^2=I$. For each original
source $a$ and its real tilt $dm_{a,\theta}=e^{\theta y}dm_a$, direct
substitution $y\mapsto-y$ proves
\[
 \|\mathscr S_Np\|_{a,-\theta}^2=\|p\|_{a,\theta}^2,
 \qquad |\theta|<\pi/2.
 \tag{RFX4}
\]
The full exponential-moment domain is the original one used in MCP;
the polynomial factors and their derivatives remain integrable on
its compact subintervals. Equation RFX4 therefore states an exact
isometry between these two specified source spaces, including both
full tilted source measures.

\subsection{The quotient, primary jets and arithmetic class}
The primary centres obey
$k-\beta_{ab}=\beta_{k-a,k-b}$. The degree
$q=e(k+1)^2$ is even for the original odd $k$. Consequently
\[
 \chi(k-S)=(-1)^q\chi(S)=\chi(S),\qquad
 \mathscr S_N(\chi f)=\chi\,\mathscr S_{N-q}f.
 \tag{RFX5}
\]
The relation space is thus carried onto itself at every cutoff.
There is an induced complex-linear involution $\Sigma$ on
$E=\mathbb C[S]/(\chi)$, with the literal identities
\[
 \operatorname{rem}_\chi\mathscr S_N
     =\Sigma\operatorname{rem}_\chi,\qquad
 \Sigma^2=I_E,\qquad \Sigma A\Sigma=kI_E-A.
 \tag{RFX6}
\]
The first follows from RFX5, and the third follows by applying
reflection to $S p(S)$. In the degree-less-than-$q$ frame, $\Sigma$
is the matrix in RFX3 with $N=q-1$. In original primary derivative
coordinates its complete formula is
\[
 \bigl(\Sigma[P]\bigr)^{(d)}(\beta_{ab})
       =(-1)^dP^{(d)}(\beta_{k-a,k-b}),\qquad 0\le d<e.
 \tag{RFX7}
\]
Dividing both sides by the same retained $d!$ gives the formula in
divided-derivative coordinates; no factorial or top nilpotent layer
is omitted. The matrix of this permutation with jet signs squares
to the identity and is precisely the conjugate of RFX3 by the full
original Hermite coordinate map.

The arithmetic inclusion also has an exact reflection map. On
$E_h=\mathbb C[s]/(h)$ put $\sigma_h[f(s)]=[f(1-s)]$.
RFX1 makes this a well-defined involution and fixes the full Taylor
unit $\upsilon_h=j_h(v_h)$. Let
$\sigma_h^{\otimes k}$ act on every tensor factor. Then
\[
 \sigma_h^{\otimes k} A_\Sigma\sigma_h^{\otimes k}
      =kI-A_\Sigma,\qquad
 \sigma_h^{\otimes k}\eta=\eta\Sigma.
 \tag{RFX8}
\]
The first identity follows by summing $M_{1-s_i}=I-M_{s_i}$.
For the second, apply that identity to
$\eta[P]=\upsilon_h^{\otimes k}P(A_\Sigma)1$; the unit and $1$
are fixed, giving exactly $\upsilon_h^{\otimes k}P(k-A_\Sigma)1$.
Thus the same injective arithmetic map is used. On exterior powers
the relation follows on every decomposable wedge and then by
linearity. In particular the positive-primary exterior class is
carried to the negative-primary class with all signs and nilpotent
coordinates specified by RFX7, and its image remains nonzero by
injectivity of $\bigwedge\eta$. Explicitly order the positive primary
blocks lexicographically and their local powers by $d=0,\ldots,e-1$,
and use the same convention on the negative side. With
$B=(k+1)^2/2$ blocks, the resulting top-wedge sign is exactly
\[
 (-1)^{\,e^2B(B-1)/2+B e(e-1)/2}.
\]
Reflection reverses the order of the $B$ blocks, contributing
$e^2$ transpositions per reversed pair, and contributes
$\sum_{d=0}^{e-1}d=e(e-1)/2$ within-block sign exponents in each
block by RFX7. This proves the displayed sign while retaining
the same full unit in the arithmetic image.
A present support label is retained
under these maps; the original external absent element keeps its
specified absent image. No amplitude cancellation changes a label.

\subsection{The full metric and the total endpoint pairing}
Give the quotient at cutoff $N\ge q-1$ its original minimum norm,
and denote its Gram by $G_N^a(\theta)$. Isometry RFX4 and relation
identity RFX5 take the whole affine fibre of a class onto the fibre
of its reflected class. Taking both attained minima therefore proves
\[
 \Sigma^*G_N^a(-\theta)\Sigma=G_N^a(\theta),\qquad
 R_N^a(-\theta)=\Sigma R_N^a(\theta)\Sigma^*.
 \tag{RFX9}
\]
The inverse formula follows by taking inverses and using
$\Sigma^{-1}=\Sigma$. The exact minimum sections satisfy
$\mathcal R_N^a(-\theta)\Sigma=\mathscr S_N\mathcal R_N^a(\theta)$:
the right side is a lift in the required reflected fibre and has the
same minimum norm, so uniqueness identifies it with the left side.

At $\theta=0$, the original real monic orthogonal polynomial
$Q_n(y)$ has parity $n$. Indeed $(-1)^nQ_n(-y)$ is monic of degree
$n$; its pairing with any lower polynomial vanishes by evenness of
the original measure and the substitution $y\mapsto-y$. Uniqueness
of the monic orthogonal polynomial proves the claim. In the original
$S$ coordinate $p_n(S)=i^nQ_n((S-c)/i)$, so its remainder obeys
\[
 \Sigma b_n=(-1)^n b_n,
 \qquad b_N^*G_Nb_{N+1}=0,
 \qquad \phi_N=0\quad(N\ge q-1).
 \tag{RFX10}
\]
For the middle equality use RFX9 at zero to equate the pairing with
$(\Sigma b_N)^*G_N(\Sigma b_{N+1})$. This is its negative by the
first equality; over $\mathbb C$ the pairing must vanish. The last
identity is the exact definition in MCE7, with $\omega_N>0$.
All three original sources and every original multiplicity are
included. More generally their tilted monic polynomials obey
$\mathscr S_np_n^a(\theta)=(-1)^np_n^a(-\theta)$ with unchanged
monic norms, by the same substitution and uniqueness. The vanishing
claim RFX10 is made at the original source $\theta=0$.

Since the determinant of the involution has squared modulus one,
RFX9 also gives
\[
 \det G_N^a(-\theta)=\det G_N^a(\theta),\qquad
 \left.\partial_\theta\log\det G_N^a(\theta)\right|_{\theta=0}=0.
 \tag{RFX11}
\]
Differentiability follows from the original exponential-moment
bound and finite positive Gram inversion, as proved in MCP1.
This is a proved derivative identity for the exact minimum metric;
it does not differentiate an asymptotic error estimate.

\subsection{The original mixed blocks remain present}
Retain the original observation $\Lambda$, kernel inclusion $I$, and
metric section $L_N$ of MCE11. There is no assumption that $\Sigma$
preserves $K=\ker\Lambda$. At the original source RFX10 implies,
for the complete component pairings of MCE15,
\[
 z_B=-z_K,
 \quad \omega_N^2\epsilon_N^2
   =(a_Kd_K-|z_K|^2)+(a_Bd_B-|z_K|^2)
                  +(a_Kd_B+a_Bd_K+2|z_K|^2).
 \tag{RFX12}
\]
The first identity is $z_K+z_B=b_N^*G_Nb_{N+1}=0$.
Inserting it into each of the three separate squared norms of MCE16
gives exactly the three displayed terms. Each term is the squared
norm of its stated exterior component; the mixed one contains
$+2|z_K|^2$. Thus RFX12 retains the full individual pairings even
when they are nonzero. Adding the three exact terms yields
$\omega_N^2\epsilon_N^2=(a_K+a_B)(d_K+d_B)$, consistently with
MCE8 and RFX10.

Here is the exact relationship to the reflected observation, without
requiring it to coincide with the original one. Define
$\Lambda^\Sigma=\Lambda\Sigma$ and $I^\Sigma=\Sigma I$ with the
same boundary and kernel coordinate spaces. Then
$\ker\Lambda^\Sigma=\Sigma K$. At $\theta=0$, RFX9 and MCE11 give
\[
 H_N^\Sigma=H_N,\quad Q_N^\Sigma=Q_N,\quad L_N^\Sigma=\Sigma L_N,
 \quad (K_N^{\rm coord})^\Sigma=K_N^{\rm coord}\Sigma.
 \tag{RFX13}
\]
For the section equality, use
$R_N\Sigma^*=\Sigma R_N$, a consequence of RFX9, in the exact
formula $L_N^\Sigma=R_N(\Lambda\Sigma)^*Q_N$.
For the kernel coordinate equality use
$\Sigma^*G_N=G_N\Sigma$. The two diagonal action blocks and both
leakage blocks consequently satisfy
\[
 A_{KK}^\Sigma=kI_K-A_{KK},\quad
 A_{BB}^\Sigma=kI_B-A_{BB},\quad
 A_{KB}^\Sigma=-A_{KB},\quad A_{BK}^\Sigma=-A_{BK}.
 \tag{RFX14}
\]
Insert RFX13 and $\Sigma A\Sigma=kI-A$ into the four exact
definitions MCE12. The off-diagonal constant terms vanish because
$K_N^{\rm coord}L_N=0$ and $\Lambda I=0$. This proves all four
relations while retaining the actual action defect $\Lambda AI$.
The reflected component vectors are $(-1)^N(x_K,x_B)$ and
$(-1)^{N+1}(y_K,y_B)$, by RFX10 and RFX13. Their squared sizes
are unchanged and both marked pairings change sign, an exact
description of their relation to the original mixed components.

In the original observation, put
$g_K=\partial_\theta\log\det H_N|_0$ and
$g_B=\partial_\theta\log\det Q_N|_0$. The complete metric splitting
has endpoint frame $[I,L_N(\theta)]$. Its determinant is independent
of $\theta$, because $L_N(\theta)-L_N(0)$ lies in $K$ and hence
changes that frame by a triangular matrix with diagonal identity.
Taking determinants in MCE11 and differentiating gives
$g_K+g_B=0$ by RFX11. Also the trace of $A$ is $cq$, by summing
all actual primary centres with multiplicity $e$. MCP3 therefore
gives the additional exact formulas
\[
 \alpha_B=-\alpha_K,\qquad g_B=-g_K,\qquad
 \frac{z_K}{\omega_N}=\alpha_K-ig_K,
 \quad\frac{z_B}{\omega_N}=-\alpha_K+ig_K,
 \quad\alpha_K=\operatorname{Tr}A_{KK}-cr_K.
 \tag{RFX15}
\]
They identify the original marked term to be controlled. Neither
$\alpha_K$ nor $g_K$ is assumed to vanish. In particular the mixed
term in RFX12 is
$a_Kd_B+a_Bd_K+2\omega_N^2|\alpha_K-ig_K|^2$.

\subsection{Propagation through the same arithmetic control budget}
Every allowance in the original arithmetic source is strictly
positive by MCE10. Substitution of the proved RFX10 into MCE25
therefore gives the exact full-multiplicity equality
\[
 2\sum_{N=q-1}^{2q-1}w_N\log\epsilon_N^{\rm ar}
 =\Psi_k+F_B^{(k)}+\mathcal W_k^{\rm ar}
                          -\mathcal P_--\mathcal P_+.
 \tag{RFX16}
\]
Every endpoint weight, monic-norm window and contraction penalty
is exactly MCE21--24. With $\sum_Nw_N=2q$, this proves
\[
 \begin{gathered}
 \mathcal L_{h,k}\le\min_N\epsilon_N^{\rm ar}
 \le\exp\!\left\{
 \frac{\Psi_k+F_B^{(k)}+\mathcal W_k^{\rm ar}
                   -\mathcal P_--\mathcal P_+}{4q}\right\},\\
 4q\log\mathcal L_{h,k}\le
 \Psi_k+F_B^{(k)}+\mathcal W_k^{\rm ar}-\mathcal P_--\mathcal P_+.
 \end{gathered}
 \tag{RFX17}
\]
The exponential is now the exact weighted geometric mean of the
allowances. A positive weighted geometric mean is at least the
minimum, and each allowance is at least $\mathcal L_{h,k}$ by the
same surviving arithmetic exterior class; these facts prove both
inequalities. No new arithmetic upper bound is assumed.

On precisely the $m=1$ evaluated-return domain of OAR/OVR/OFV,
retain the actual signed error and substitute MCE27 to obtain
\[
 \begin{aligned}
 2\sum_Nw_N\log\epsilon_N^{\rm ar}
  ={}&\mathcal F_{\partial,k}^{(k)}-\Phi_{k,k}
       +\mathcal H_k+\delta_k^\sigma+\mathfrak e_{k,k}\\
       &+F_B^{(k)}+\mathcal W_k^{\rm ar}-\mathcal P_--\mathcal P_+,\\
 \Phi_{k,k}\le{}&\mathcal F_{\partial,k}^{(k)}
       +\mathcal H_k+\delta_k^\sigma+e_{k,k}^+\\
       &+F_B^{(k)}+\mathcal W_k^{\rm ar}-\mathcal P_--\mathcal P_+
       -4q\log\mathcal L_{h,k}.
 \end{aligned}\tag{RFX18}
\]
The second line follows from the first, RFX17, and the original
upper error endpoint $\mathfrak e_{k,k}\le e_{k,k}^+$.
The complete evaluated value
$\mathcal F_{\partial,k}^{(k)}=\Delta q\mathfrak C_\partial
+O_{\delta,\gamma}(q\log q)$ may be inserted with its proved
error and original domain. The arithmetic boundary, norm window,
contraction penalties and residual remain explicit. Equations
RFX12--15 retain the individual mixed blocks required for further
control; the total phase has been evaluated through a proved source,
quotient and arithmetic involution, with the full data preserved.

% END CURRENT RFX

\clearpage
% BEGIN CURRENT OMR
\section{The outer operator, the full residual graph, and original action control}
\label{sec:OMR}

This receiving calculation uses the complete OOQ1--40, GIP1--9 and
MCE1--28 proofs, with their original source maps and constants.
The preceding MRI and BRI equations remain valid. Here we give the
actual substitutions into those equations and into the control of
the same arithmetic class.

\subsection{Original objects, domains and the same-target operator interval}

For the original quartet retain $0<\delta<1/2$, $\gamma>2$,
$m\ge1$, $k=4l+1\ge9$, $e=1+k(m-1)$, $c=k/2$ and
\[
 \begin{gathered}
 q=e(k+1)^2,\\
 \chi(S)=\prod_{a,b=0}^k
 [S-c-(2a-k)\delta-i(2b-k)\gamma]^e,\\
 E=\mathbb C[S]/(\chi),\qquad A=M_S:E\longrightarrow E.
 \end{gathered}\tag{OMR1}
\]
The original source measures are $w_h^{*k}(y)\,dy$ and the
two full Gamma measures of orders $s=1,k$, evaluated at the
unchanged coordinate $S=c+iy$. Their masses remain $\mu_h^k$
and $(2\pi)^{s/2}$. All action identities below retain this full
multiplicity. The outer estimates and the residual graph statements
are applied only on their proved simple-quartet domain $m=1$,
with $q=(k+1)^2$, $q'=(k-7)^2$, $\Delta=q-q'=16k-48$,
$r=8k-16$ and $j=8k-32$. They retain the actual period and the
full explicit finite domains of OSP23, OOQ13 and OAR1--3.
No eventual bound is assigned to an input outside its stated domain.

Write $\varepsilon_k^{\rm out}$ for exactly OOQ33's finite
$\varepsilon_k$, with all OOQ4--33 constants, the full root
polynomials and both source orders. It is distinct from the
action radius $\epsilon_N^{\rm ar}$ below and from
$E_N^{\rm CEP}$, the conductor exterior error. Put
\[
 \begin{gathered}
 \mathfrak C_\partial=4\log(4/\pi)+2-2I_{2,\pi},\\
 I_{2,\pi}=\iint\log|x-y|\,d\mu_{2,\pi}(x)d\mu_{2,\pi}(y),\\
 D_k^{\rm out}=4\Delta\varepsilon_k^{\rm out},\qquad
 \mathcal O_k^-=\Delta q\mathfrak C_\partial-D_k^{\rm out},\quad
 \mathcal O_k^+=\Delta q\mathfrak C_\partial+D_k^{\rm out}.
 \end{gathered}\tag{OMR2}
\]
The measure $\mu_{2,\pi}$ is the complete EIQ equilibrium with
its original mass one, endpoints and density, as retained in OFV
and OVR. This is the same constant in MRI4r--u and BRI6h--i.

For clarity write $Q_{{\rm lo},0}^{(s)}=Q_{\partial,\Delta-1}^{(s)}$,
$Q_{{\rm lo},1}^{(s)}=Q_{\partial,\Delta}^{(s)}$,
$Q_{{\rm hi},0}^{(s)}=Q_{\partial,q+\Delta-1}^{(s)}$ and
$Q_{{\rm hi},1}^{(s)}=Q_{\partial,q+\Delta}^{(s)}$.
Every one acts on the original coefficient remainder frame
$1,y,\ldots,y^{\Delta-1}$. OOQ34 gives, for $a=0,1$,
\[
 \begin{gathered}
 \beta_s e^{c_{\rm lo}-\varepsilon_k^{\rm out}}I_\Delta
 \preceq Q_{{\rm lo},a}^{(s)}
 \preceq\beta_s e^{c_{\rm lo}+\varepsilon_k^{\rm out}}I_\Delta,\\
 \beta_s e^{c_{\rm hi}-\varepsilon_k^{\rm out}}I_\Delta
 \preceq Q_{{\rm hi},a}^{(s)}
 \preceq\beta_s e^{c_{\rm hi}+\varepsilon_k^{\rm out}}I_\Delta,\\
 \beta_s=\frac{(2\pi)^{s/2}}{\sqrt2\,\Gamma(s/2)},\qquad
 c_{\rm lo}-c_{\rm hi}=q\mathfrak C_\partial/2.
 \end{gathered}\tag{OMR3}
\]
All original masses are retained in $\beta_s$. Inverting the high
inequalities and conjugating the low inequalities by their positive
square roots proves that every eigenvalue of
$\mathscr R_a=(Q_{{\rm hi},a}^{(s)})^{-1/2}
Q_{{\rm lo},a}^{(s)}(Q_{{\rm hi},a}^{(s)})^{-1/2}$ lies between
$e^{q\mathfrak C_\partial/2-2\varepsilon_k^{\rm out}}$ and
$e^{q\mathfrak C_\partial/2+2\varepsilon_k^{\rm out}}$.
Conjugating the bounds on $\mathscr R_1$ by $\mathscr R_0^{1/2}$
then proves the corresponding bounds on
$\mathscr R_0^{1/2}\mathscr R_1\mathscr R_0^{1/2}$, without any
commutation assumption. Its determinant is exactly the product
of the two original low/high determinant ratios. Therefore
\[
 \mathcal O_k^-\le\mathcal F_{\partial,k}^{(s)}
                      \le\mathcal O_k^+,\qquad s=1,k.
 \tag{OMR4}
\]
The error is $O_{\delta,\gamma}(\Delta q^{3/4}+\Delta k\log q)
=o(kq)$ by OOQ35. The earlier OFV estimate with error $O(q\log q)$
also remains valid; OMR4 supplies a full operator origin for an
additional interval, and does not assert that this scalar interval
is tighter than that earlier bound at every finite input.

More generally let $I:\mathbb C^t\hookrightarrow\mathbb C^\Delta$
be a fixed injection, or let $C_0:\mathbb C^\Delta\twoheadrightarrow
\mathbb C^t$ be a fixed surjection. Here ``fixed'' means the same
actual map at all four endpoints; dependence on $k$ and the period
is allowed. The restriction and quotient forms are exactly
\[
\begin{gathered}
 Q_I=I^*QI,\quad Q_I^{\rm ref}=I^*I,\\
 Q_{C_0}=(C_0Q^{-1}C_0^*)^{-1},\quad
 Q_{C_0}^{\rm ref}=(C_0C_0^*)^{-1},\\
 |\mathcal F_t^{(s)}-tq\mathfrak C_\partial|
          \le4t\varepsilon_k^{\rm out}.
\end{gathered}\tag{OMR5}
\]
Restriction of a positive-form inequality preserves it. For a
quotient, inversion reverses the order, multiplication by the onto
map preserves it, and the second inversion reverses it again.
Thus the same scalar bounds hold relative to the displayed exact
reference form. In the four determinants its logarithm has
coefficient $1+1-1-1=0$; the complete frame cost cancels.
The $t$ eigenvalue bounds now prove the last inequality.

\subsection{The full graph denominator and its exact outer morphism}

We now give the actual map relating OMR5's target to the original
residual. This retains the full conductor $A_{\rm cond}$ of CEP,
which is a different named map from $A=M_S$ in OMR1. For each
original $N=q-1,q,2q-1,2q$ let $t=N-q+1$, retain the first
nonzero conductor moment index $0\le v\le80$, and define
\[
 \begin{gathered}
 X_N=\mathbb C^v\oplus\mathbb C[S']_{\le N-v-q'},\quad
 D_N=\operatorname{diag}(I_v,M_{\chi_{k-8}}),\quad
 \dim X_N=t+\Delta,\\
 E_N^{\rm graph}=U_NS_N^{-1}D_N:X_N\twoheadrightarrow V_A,\\
 B_N^{\rm graph}=\Xi_k^{\rm fix}E_N^{\rm graph}
       =E_j^{\mathsf T}T_k^{\mathsf T}\mathsf S_k^{\mathsf T}
                       U_NS_N^{-1}D_N:X_N\twoheadrightarrow\mathbb C^j.
 \end{gathered}\tag{OMR6}
\]
These are the full GIP1--2 maps: $U_N$ is the original source
remainder in primary coordinates and $S_N$ is CEP28's complete
source map. The multiplication block of $D_N$ retains its exact
conversion from ordinary lower monomials to orthonormal source
coordinates. Let $Z_{\chi,N}h$ be the original upper coefficients
of $\chi_kh$. With every shift and root product of CEP41--42,
\[
 \begin{gathered}
 \mathcal Z_{A,N}h=
   \begin{pmatrix}P_{<v}Z_{\chi,N}h\\
     \displaystyle\sum_{a,b=0}^8a_{ab}d_{ab}(S')h(S'+b_{ab})
   \end{pmatrix},\qquad
 D_N\mathcal Z_{A,N}=S_NZ_{\chi,N},\\
 I_k=J_kE_r,\quad Y_{K,N}=U_N^*G_NI_k,\quad
 W_{K,N}=D_N^{-1}S_NY_{K,N},\\
 \ker E_N^{\rm graph}=\operatorname{im}\mathcal Z_{A,N},\quad
 \ker B_N^{\rm graph}
       =\operatorname{im}[\mathcal Z_{A,N},W_{K,N}],\quad
 \operatorname{rank}[\mathcal Z_{A,N},W_{K,N}]=t+r.
 \end{gathered}\tag{OMR7}
\]
The graph columns are independent because for every nonzero $h$
their complete leading coefficient is
$\binom{q+\deg h}{v}\mu_v\operatorname{lc}(h)\ne0$;
for the specified monomial columns $\operatorname{lc}(h)=1$.
Applying $E_N^{\rm graph}$ to the last $r$ columns gives $I_k$,
since $U_NU_N^*=G_N^{-1}$. Thus a relation among both sets first
has zero coefficients in $W_{K,N}$ and then in $\mathcal Z_{A,N}$.
Conversely, subtracting the corresponding $W_{K,N}$ lift from any
element of $\ker B_N^{\rm graph}$ leaves a source relation.
This proves the complete kernel formula, including the $t=0$ case
where the source-relation matrix has no columns.

The two actual positive metrics and their residual quotient are
\[
 \begin{gathered}
 H_N^{\rm actual}=D_N^*S_N^{-*}S_N^{-1}D_N,\quad
 H_N^{\rm graph}=D_N^*D_N,\\
 Q_{\Xi,N}^{\rm fix}
   =[B_N^{\rm graph}(H_N^{\rm actual})^{-1}
                              (B_N^{\rm graph})^*]^{-1},\\
 \widetilde Q_{\Xi,N}
   =[B_N^{\rm graph}(H_N^{\rm graph})^{-1}
                              (B_N^{\rm graph})^*]^{-1},\\
 |\log\det\widetilde Q_{\Xi,N}-\log\det Q_{\Xi,N}^{\rm fix}|
                           \le2E_N^{\rm CEP}.
 \end{gathered}\tag{OMR8}
\]
For any onto $B$ and positive $H$, the vector
$H^{-1}B^*(BH^{-1}B^*)^{-1}y$ is the unique lift of $y$
orthogonal to $\ker B$; every other lift adds an orthogonal
kernel vector. Its squared norm proves the displayed quotient
formula. The first metric is the exact pullback of the original
upper source norm, so this is precisely MRI4v's Schur complement.
The second metric retains the full low-$v$ summand and the source
$|P_{k-8}(y)|^2dm_s(y)$ after the unchanged lower substitution.
The full $j$th exterior bounds of $S_N$ and $S_N^{-1}$, followed
by orthogonal projection off OMR7's entire kernel, give the last
inequality as in GIP6. Their four endpoint costs total at most
$2\sum_NE_N^{\rm CEP}=o(kq)$ for the fixed actual period.

Let $E_{\rm in},E_{\rm out}$ insert the original inner square
$\{4,\ldots,k-4\}^2$ and its complement in primary order.
Retain the full outer Vandermonde $V_{\rm out}$ and the diagonal
$D_{\rm in,out}=\operatorname{diag}
(P_{k-8}(\omega_{ab}))_{\rm out}$. The exact natural morphism is
\[
 \begin{gathered}
 \rho_A=V_{\rm out}^{-1}D_{\rm in,out}^{-1}
                         E_{\rm out}^{\mathsf T}|_{V_A},\\
 C_{\rm in}=CE_{\rm in},\quad C_{\rm out}=CE_{\rm out},\quad
 (C_{\rm in})_{(a,b),(i,j)}=a_{i+4-a,j+4-b},\\
 0\longrightarrow\ker C_{\rm in}
   \xrightarrow{E_{\rm in}}V_A
   \xrightarrow{\rho_A}\mathbb C^\Delta
   \xrightarrow{\partial_A}\operatorname{coker}C_{\rm in}
   \longrightarrow0,\\
 \partial_A f=[C_{\rm out}D_{\rm in,out}V_{\rm out}f],\quad
 \dim\ker\rho_A=\dim\operatorname{coker}\rho_A
                              =\dim\ker C_{\rm in}.
 \end{gathered}\tag{OMR9}
\]
Indeed, writing a primary vector as $E_{\rm in}z+E_{\rm out}w$,
its conductor equation is $C_{\rm in}z+C_{\rm out}w=0$ and its
outer image $f$ specifies $w=D_{\rm in,out}V_{\rm out}f$.
These two equations prove exactness; the surjectivity of $C$
proves surjectivity onto the cokernel. The inner block is square,
so its kernel and cokernel dimensions agree. On the original
graph the map is exactly $\rho_AE_N^{\rm graph}$.

The residual descends with its full possible defect as
\[
 W_A=\Xi_k^{\rm fix}E_{\rm in}\ker C_{\rm in},\qquad
 \overline\Xi_A:\ker\partial_A\twoheadrightarrow\mathbb C^j/W_A,
 \quad \rho_Ax\longmapsto[\Xi_k^{\rm fix}x],\qquad
 \ker\overline\Xi_A=\rho_A(K).
 \tag{OMR10}
\]
Two lifts differ by $E_{\rm in}\ker C_{\rm in}$, which proves
well-definedness. Surjectivity follows from that of $\Xi_k^{\rm fix}$.
If a class has zero residual modulo $W_A$, subtract a vector in
$E_{\rm in}\ker C_{\rm in}$ with the same residual. Its remaining
lift lies in $K$ and has the same outer image, proving the kernel
identity. Equations OMR6--10 are the exact typed bridge used here.
They retain the entire denominator and metrics. Consequently the
fixed-subquotient bound OMR5 is applied only to a specified map
of its own full outer target; this receiving calculation assigns
no leading coefficient to $\Phi_{k,s}$ by replacing the actual
graph denominator with an outer polynomial ideal.

\subsection{Backward propagation through the original kernel and signed return}

Keep the actual nonnegative residual sum $\Phi_{k,s}$ in MRI4l,
its first increment $\lambda_s$, and its full earlier cap $V_s$.
Write $[L_s^0,U_s^0]$ and $[L_{\rm ar}^0,U_{\rm ar}^0]$ for
exactly MRI4q's existing intervals. On the common proved domain,
the full operator estimate gives the additional same-scalar intervals
\[
 \begin{gathered}
 L_s^{\rm op}=\max\{L_s^0,
       \mathcal O_k^--\Phi_{k,s}-e_{k,s}^-\},\\
 U_s^{\rm op}=\min\{U_s^0,
       \mathcal O_k^+-\Phi_{k,s}+e_{k,s}^+\},\\
 L_s^{\rm op}\le F_K^{(s)}\le U_s^{\rm op},\qquad s=1,k.
 \end{gathered}\tag{OMR11}
\]
To prove this, substitute OMR4 into the exact MRI4m equality
$F_K^{(s)}=\mathcal F_{\partial,k}^{(s)}-\Phi_{k,s}
+\mathfrak e_{k,s}$, with
$-e_{k,s}^-\le\mathfrak e_{k,s}\le e_{k,s}^+$.
Intersecting with the previously proved interval gives OMR11.
The intersection is nonempty because the same original scalar
belongs to both intervals. All earlier branches remain available;
outside the new domain one uses their original endpoints.

Define the unchanged correlated arithmetic allowances by
\[
 \begin{gathered}
 a_{K,k}^-=\max\{-C_K^-,\delta_k^\sigma-C_B^+\},\quad
 a_{K,k}^+=\min\{C_K^+,\delta_k^\sigma+C_B^-\},\\
 a_{K,k}^-\le d_k^{\rm ar}=F_K^{\rm ar}-F_K^{(1)}
                                      \le a_{K,k}^+,\\
 L_{\rm ar}^{\rm op}=\max\{L_{\rm ar}^0,L_1^{\rm op}+a_{K,k}^-\},\\
 U_{\rm ar}^{\rm op}=\min\{U_{\rm ar}^0,U_1^{\rm op}+a_{K,k}^+\},\\
 F_K^{\rm ar}\in[L_{\rm ar}^{\rm op},U_{\rm ar}^{\rm op}],\qquad
 F_B^{\rm ar}\in[\mathcal B_k^{\rm ar}-U_{\rm ar}^{\rm op},
                 \mathcal B_k^{\rm ar}-L_{\rm ar}^{\rm op}].
 \end{gathered}\tag{OMR12}
\]
The first bounds are exactly OAR8/MRI4o, with their original
constants. Adding them to the source-$1$ inequalities proves the
kernel interval. Subtraction from the same arithmetic total proves
the boundary interval and its reversed endpoints. Their errors
remain correlated through this equality.

The earlier signed scalar is exactly MCE22's $\Psi_k$:
\[
 \begin{gathered}
 \Psi_k=\mathfrak S_k^{\rm mix}
      =F_K^{\rm ar}+F_B^{\rm ar}-F_B^{(k)}
      =F_K^{(k)}+\mathcal H_k+\delta_k^\sigma,\\
 L_k^{\rm op}+\mathcal H_k+\delta_k^\sigma
 \le\Psi_k\le
 U_k^{\rm op}+\mathcal H_k+\delta_k^\sigma.
 \end{gathered}\tag{OMR13}
\]
The first equality uses the original total identity BRI6d, including
$\mathcal H_k=-T_k$. The second line adds this exact signed correction
to OMR11 at the original source order $s=k$. Every earlier finite
bound for this same $\Psi_k$ may also be intersected here. No
arithmetic allowance or Gamma sign is changed.

\subsection{The original action, its mixed components, and its positive class}

For any of the actual sources let $p_N$ be the monic polynomial
in $S$, $\omega_N$ its full squared source norm, and $b_N=[p_N]_\chi$.
In the original coefficient coordinates put
\[
 \begin{gathered}
 R_N=\sum_{n=0}^N b_nb_n^*/\omega_n,\quad G_N=R_N^{-1},\quad
 \mathscr D_N=A^*G_N+G_NA-kG_N,\\
 x_N=G_N^{1/2}b_N,\quad y_N=G_N^{1/2}b_{N+1},\quad
 x_N^*y_N=i\omega_N\phi_N,\quad \phi_N\in\mathbb R,\\
 G_N^{-1/2}\mathscr D_NG_N^{-1/2}
       =(y_Nx_N^*+x_Ny_N^*)/\omega_N,\\
 \operatorname{spec}(G_N^{-1/2}\mathscr D_NG_N^{-1/2})
       =\{\epsilon_N,-\epsilon_N,0,\ldots,0\},\quad
 \epsilon_N^2=\frac{\|x_N\wedge y_N\|^2}{\omega_N^2}.
 \end{gathered}\tag{OMR14}
\]
The monic recurrence in MCE5 telescopes in $AR_N+R_NA^*-kR_N$
to $(b_{N+1}b_N^*+b_Nb_{N+1}^*)/\omega_N$.
Multiplication by $G_N$ on both sides gives the control identity.
Its trace is $2\Re\operatorname{Tr}A-kq=0$, which proves that
$x_N^*y_N$ is purely imaginary. The square of this rank-at-most-two
Hermitian operator has half trace
$(\|x_N\|^2\|y_N\|^2-|x_N^*y_N|^2)/\omega_N^2$.
This proves the two opposite eigenvalues and the exterior formula.

For the original observation $\Lambda:E\twoheadrightarrow B$
and inclusion $I:K\hookrightarrow E$, retain
\[
 \begin{gathered}
 H_N=I^*G_NI,\quad Q_N=(\Lambda R_N\Lambda^*)^{-1},\quad
 L_N=R_N\Lambda^*Q_N,\quad K_N^{\rm coord}=H_N^{-1}I^*G_N,\\
 \Lambda L_N=I_B,\quad I^*G_NL_N=0,\quad
 I K_N^{\rm coord}+L_N\Lambda=I_E,\\
 A_{KK}=K_N^{\rm coord}AI,\quad A_{KB}=K_N^{\rm coord}AL_N,\quad
 A_{BK}=\Lambda AI,\quad A_{BB}=\Lambda AL_N,\\
 \Lambda A-A_{BB}\Lambda=A_{BK}K_N^{\rm coord},\\
 I^*\mathscr D_NL_N=A_{BK}^*Q_N+H_NA_{KB}.
 \end{gathered}\tag{OMR15}
\]
The first three identities follow by multiplication and the
minimum-lift orthogonality. Applying $A$ in this direct-sum frame
gives the action blocks and the displayed observation defect.
Multiplication of $A^*G_N+G_NA-kG_N$ by $I^*$ and $L_N$
gives the final identity, since $I^*G_NL_N=0$. This retains the
actual action on the observation; it does not impose invariance
of its kernel.

Define the marked components exactly by
$x_K=H_N^{-1/2}I^*G_Nb_N$, $x_B=Q_N^{1/2}\Lambda b_N$,
and the same formulas with $b_{N+1}$ for $y_K,y_B$. Put
$a_K=\|x_K\|^2$, $a_B=\|x_B\|^2$,
$d_K=\|y_K\|^2$, $d_B=\|y_B\|^2$,
$z_K=x_K^*y_K$ and $z_B=x_B^*y_B$. Then
\[
 \begin{gathered}
 z_K+z_B=i\omega_N\phi_N,\\
 \omega_N^2\epsilon_N^2
 =\|x_K\wedge y_K\|^2+\|x_B\wedge y_B\|^2
          +\|x_K\otimes y_B-y_K\otimes x_B\|^2\\
 =(a_Kd_K-|z_K|^2)+(a_Bd_B-|z_B|^2)
       +a_Kd_B+a_Bd_K-2\Re(z_K\overline{z_B}).
 \end{gathered}\tag{OMR16}
\]
The frame $[I,L_N]$ has Gram $\operatorname{diag}(H_N,Q_N)$,
so the marked components give an isometry to $K\oplus B$.
The exterior square decomposes orthogonally into
$\bigwedge^2K\oplus(K\otimes B)\oplus\bigwedge^2B$.
Expanding the middle norm gives the last three terms, including
its complete complex cross-pairing. This proves every equality.

The full original positive primary part $E_+$ consists of all
jets with $a>k/2$, has dimension $q/2$, and its top exterior
class is nonzero under the original injection
$\eta:E\hookrightarrow E_h^{\otimes k}$ because $\eta$ is
injective and $\eta A=A_{\rm sum}\eta$. Its action excess is
\[
 \mathcal L_{h,k}=2\Re\operatorname{Tr}(A|_{E_+})-kq/2
     =2\delta e(k+1)\left\lfloor\frac{(k+1)^2}{4}\right\rfloor
     \le\epsilon_N,\qquad N\ge q-1.
 \tag{OMR17}
\]
The nilpotent primary terms have zero trace. Summing the real
parts over $a>k/2$, all $b$ and all $e$ jet directions gives the
displayed coefficient. The induced additive Hermitian action on
any exterior degree strictly between zero and $q$ has largest
eigenvalue $\epsilon_N$, since OMR14's only nonzero eigenvalues
are $\epsilon_N,-\epsilon_N$. Evaluating it on the positive
class therefore proves the inequality. In the frame OMR15 its
representative has both components
$X_+=K_N^{\rm coord}I_+$ and $Y_+=\Lambda I_+$, with Gram
$X_+^*H_NX_++Y_+^*Q_NY_+$; its expectation includes the cross
block in OMR15. Thus this is the control of that same original
class, with no component omitted.

\subsection{The exact arithmetic radius budget and its backward bound}

Define $\kappa_N=\log(\det H_N/\det H_{N+1})$,
$\beta_N^B=\log(\det Q_N/\det Q_{N+1})$, and
$\tau_N=\log(V_N/V_{N+1})=\kappa_N+\beta_N^B$,
where $V_N=\det G_N$. The original covariance rank-one update
and its restriction give
\[
 \begin{gathered}
 e^{\beta_N^B}=1+d_B/\omega_{N+1},\qquad
 e^{\kappa_N}=1+d_K/(\omega_{N+1}+d_B),\\
 \frac{a_B}{\omega_N}=1-e^{-\beta_{N-1}^B},\quad
 \frac{a_K}{\omega_N}=e^{-\beta_{N-1}^B}(1-e^{-\kappa_{N-1}})
                    \quad(N\ge q),\\
 \epsilon_{q-1}^2+\phi_{q-1}^2
   =\frac{\omega_q}{\omega_{q-1}}(e^{\tau_{q-1}}-1),\\
 \epsilon_N^2+\phi_N^2
   =\frac{\omega_{N+1}}{\omega_N}
          (1-e^{-\tau_{N-1}})(e^{\tau_N}-1)\quad(N\ge q).
 \end{gathered}\tag{OMR18}
\]
The first two formulas are the determinant lemma applied to the
boundary covariance and then to the restricted Sherman--Morrison
update for $G_N$. Applying the determinant lemma to removal of
the last column proves the next two formulas. At $N=q-1$ the
first $q$ monic columns form the complete quotient basis, so
$a_K+a_B=\omega_{q-1}$; no inverse at degree $q-2$ is used.
Multiplying the total $a$ and $d$ formulas and using OMR14 proves
the two radius identities.

Set $w_{q-1}=w_{2q-1}=1$ and $w_N=2$ for
$q\le N\le2q-2$, so that $\sum_Nw_N=2q$.
For the arithmetic source alone, put
\[
 \begin{gathered}
 \delta_j^{\rm ar}=V_j^{\rm ar}/V_{j-1}^{\rm ar}\quad(j\ge q),\\
 \mathcal W_k^{\rm ar}
   =\log\frac{\omega_{2q-1}^{\rm ar}\omega_{2q}^{\rm ar}}
                    {\omega_{q-1}^{\rm ar}\omega_q^{\rm ar}},\\
 \mathcal P_-=-\log(1-\delta_{2q-1}^{\rm ar})
                 -2\sum_{j=q}^{2q-2}\log(1-\delta_j^{\rm ar}),\\
 \mathcal P_+=-\log(1-\delta_q^{\rm ar})-\log(1-\delta_{2q}^{\rm ar})
                 -2\sum_{j=q+1}^{2q-1}\log(1-\delta_j^{\rm ar}),\\
 C_k^{\rm budget}=F_B^{(k)}+\mathcal W_k^{\rm ar}-\mathcal P_--\mathcal P_+,\\
 J_k^{\rm action}=\sum_{N=q-1}^{2q-1}w_N
                  \log((\epsilon_N^{\rm ar})^2+(\phi_N^{\rm ar})^2)
                  =\Psi_k+C_k^{\rm budget}.
 \end{gathered}\tag{OMR19}
\]
Here $0<\delta_j^{\rm ar}<1$: positivity gives the first
inequality and OMR17--18 forces the strict contraction at every
index in the formula. All logarithms are therefore finite.
To prove the budget identity, in each of the two windows
$N=q-1,\ldots,2q-2$ and $N=q,\ldots,2q-1$ substitute
$e^{\tau_N}-1=(1-\delta_{N+1})/\delta_{N+1}$ into OMR18.
The monic norm ratios telescope to $\mathcal W_k^{\rm ar}$;
the reciprocal $\delta$ factors give $F_K^{\rm ar}+F_B^{\rm ar}$.
Each interior contraction occurs twice and the extreme ones once,
giving exactly $-\mathcal P_--\mathcal P_+$. Finally subtract
and restore the same reference $F_B^{(k)}$ using OMR13.
This proves the equality with every original endpoint.

The phases and the lower control allowance give the exact receiving
inequalities
\[
 \begin{gathered}
 Z_k^{\rm phase}=\sum_{N=q-1}^{2q-1}w_N
       \log\left(1+\frac{(\phi_N^{\rm ar})^2}{\mathcal L_{h,k}^2}\right),\\
 L_J=L_k^{\rm op}+\mathcal H_k+\delta_k^\sigma+C_k^{\rm budget},\quad
 U_J=U_k^{\rm op}+\mathcal H_k+\delta_k^\sigma+C_k^{\rm budget},\\
 4q\log\mathcal L_{h,k}+Z_k^{\rm phase}
       \le J_k^{\rm action}\in[L_J,U_J],\\
 \mathcal L_{h,k}\le\min_{q-1\le N\le2q-1}\epsilon_N^{\rm ar}
       \le\exp\left(\frac{J_k^{\rm action}}{4q}\right)
       \le\exp\left(\frac{U_J}{4q}\right).
 \end{gathered}\tag{OMR20}
\]
Indeed $\epsilon_N^{\rm ar}\ge\mathcal L_{h,k}>0$ term by term.
Thus each logarithm in OMR19 is at least
$2\log\mathcal L_{h,k}+
\log(1+(\phi_N^{\rm ar})^2/\mathcal L_{h,k}^2)$.
The weights sum to $2q$, giving the first bound.
Also $\log((\epsilon_N^{\rm ar})^2+(\phi_N^{\rm ar})^2)
\ge2\log\min_M\epsilon_M^{\rm ar}$; summing proves the
middle upper bound. OMR13 and the exact budget equality prove
$L_J\le J_k^{\rm action}\le U_J$ and the final bound.

Finally the same exact signed substitution as BRI6f gives
\[
 \begin{gathered}
 J_k^{\rm action}
 =\mathcal F_{\partial,k}^{(k)}-\Phi_{k,k}
       +\mathcal H_k+\delta_k^\sigma+\mathfrak e_{k,k}
       +C_k^{\rm budget},\\
 \Phi_{k,k}\le
 \mathcal F_{\partial,k}^{(k)}+\mathcal H_k+\delta_k^\sigma
       +e_{k,k}^++C_k^{\rm budget}
       -4q\log\mathcal L_{h,k}-Z_k^{\rm phase},\\
 \Phi_{k,k}\le
 \mathcal O_k^++\mathcal H_k+\delta_k^\sigma
       +e_{k,k}^++C_k^{\rm budget}
       -4q\log\mathcal L_{h,k}-Z_k^{\rm phase}.
 \end{gathered}\tag{OMR21}
\]
The first equality preserves the complete original OAR signed
error. Move $\Phi_{k,k}$ to the other side, use OMR20's lower
bound for $J_k^{\rm action}$, and then use
$\mathfrak e_{k,k}\le e_{k,k}^+$ to prove the second line.
OMR4 proves the third. These caps may be intersected with the
previous $\lambda_k\le\Phi_{k,k}\le V_k$ and all its earlier
proved bounds, on their common domain.

Equations OMR19--21 retain the complete reference boundary, the
arithmetic monic-norm window, both literal contraction penalties
and every marked phase. The operator interval propagates to the
original scalar and then to the control of the same surviving
arithmetic class by these exact maps. They provide no sign-based
exclusion from $\Psi_k$ alone and assign no unevaluated leading
value to $\Phi_{k,s}$.

\subsection{The exact marked derivative and original reflection specialization}

Use the original real source tilt $dm_{a,\theta}(y)=e^{\theta y}dm_a(y)$,
with $|\theta|<\pi/2$, leaving $E,A,\Lambda,I$ and the original
arithmetic inclusion fixed. Every expression in this paragraph is
evaluated at the same real $\theta$. The full exponential-moment
bound permits differentiation of every finite Gram integral on
compact subintervals. If $\mathcal R_N$ is the minimum source
section in fixed monomial source coordinates, differentiation of
$J_N\mathcal R_N=I$ puts $\mathcal R_N'$ in $\ker J_N$.
Its pairings with the minimum section vanish. Thus differentiating
the minimum Gram inserts exactly $y$ in its source integral.
The leading coefficient of a minimum lift of $x$ is
$d_N^*G_Nx/\omega_N$ in the real monic frame
$Q_n^{\rm mon}=i^{-n}p_n(c+iy)$, $d_n=[Q_n^{\rm mon}]$.
The polynomial
$y\mathcal R_N-\mathcal R_NM_y
 -(Q_{N+1}^{\rm mon}-\mathcal R_Nd_{N+1})d_N^*G_N/\omega_N$
has zero remainder and degree at most $N$. Its pairing with
$\mathcal R_N$ vanishes, and $Q_{N+1}^{\rm mon}$ is orthogonal to that
section. Returning by the exact substitution $S=c+iy$ gives
\[
 \begin{gathered}
 \partial_\theta G_N
   =-iG_N(A-cI)+iG_Nb_{N+1}b_N^*G_N/\omega_N,\\
 \frac{z_K}{\omega_N}
   =\operatorname{Tr}A_{KK}-c\dim K
               -i\partial_\theta\log\det H_N,\\
 \frac{z_B}{\omega_N}
   =\operatorname{Tr}A_{BB}-c\dim B
               -i\partial_\theta\log\det Q_N.
 \end{gathered}\tag{OMR22}
\]
The first equality is the full MCP1--2 calculation. For the second,
multiply it by $I^*$ and $I$, then take its $H_N^{-1}$ trace.
The first term is $-i(\operatorname{Tr}A_{KK}-c\dim K)$ and
the rank-one term has trace $iz_K/\omega_N$ by cyclicity.
Rearrangement proves the stated sign. For the third use
$Q_N'=L_N^*G_N'L_N$, whose section-derivative terms vanish by
$I^*G_NL_N=0$, and $L_N^*G_N=Q_N\Lambda$.
This proves the boundary formula with its complete marked pairing.
No derivative estimate follows from differentiating an asymptotic
remainder; OMR22 concerns the exact finite source family.

Here the original, untilted source has a proved additional
involution. The completed function and full quartet obey
$\xi(1-s)=\xi(s)$ and $h(1-s)=(-1)^{4m}h(s)=h(s)$, so their
entire quotient satisfies $v_h(1-s)=v_h(s)$, including every
removed zero and its Taylor coefficients. Consequently $w_h$ is
even and substitution in every convolution factor makes $w_h^{*k}$
even. Both original Gamma densities are also even, with their full
masses. Define in the original coefficient frame
\[
 \begin{gathered}
 (\mathscr S_Np)(S)=p(k-S),\qquad
 (\mathscr S_N)_{ab}=(-1)^a\binom ba k^{b-a}\quad(0\le a\le b\le N),\\
 \chi(k-S)=\chi(S),\quad
 \operatorname{rem}_\chi\mathscr S_N=\Sigma\operatorname{rem}_\chi,\quad
 \Sigma^2=I,\quad \Sigma A\Sigma=kI-A,\\
 (\Sigma[P])^{(d)}(\beta_{ab})
           =(-1)^dP^{(d)}(\beta_{k-a,k-b})\quad(0\le d<e),\\
 \sigma_h^{\otimes k}\eta=\eta\Sigma,\qquad
 \sigma_h[f(s)]=[f(1-s)].
 \end{gathered}\tag{OMR23}
\]
The first formula is the literal binomial expansion. Reflection
permutes the centres by $(a,b)\mapsto(k-a,k-b)$ and $q$ is even,
so it preserves the full relation space. This proves the quotient
map; differentiation proves every jet sign. The local involution
$\sigma_h$ fixes the full Taylor unit, and conjugates the tensor
sum action to $kI-A_{\rm sum}$. Applying it to
$\eta[P]=\upsilon_h^{\otimes k}P(A_{\rm sum})1$ proves the last
identity. Thus the original arithmetic class is carried by the
same injective map, with all tensor factors and jets retained.

Changing variables $y\mapsto-y$ carries the whole source fibre
at tilt $\theta$ isometrically onto the reflected fibre at
$-\theta$. The unique minimum sections and their Grams therefore
obey
\[
 \begin{gathered}
 \Sigma^*G_N(-\theta)\Sigma=G_N(\theta),\qquad
 \mathcal R_N(-\theta)\Sigma=\mathscr S_N\mathcal R_N(\theta),\\
 \Sigma b_n(0)=(-1)^n b_n(0),\qquad
 b_N(0)^*G_N(0)b_{N+1}(0)=0,\qquad\phi_N(0)=0,\\
 \left.\partial_\theta\log\det G_N(\theta)\right|_0=0.
 \end{gathered}\tag{OMR24}
\]
Indeed for an even measure the polynomial $(-1)^nQ_n^{\rm mon}(-y)$ is
monic and has the same orthogonality as $Q_n^{\rm mon}(y)$, so uniqueness
proves its parity. The quotient isometry then takes the adjacent
pairing to its negative; over $\mathbb C$ it is zero. This gives
$\phi_N(0)=0$ by its original definition. Finally
$|\det\Sigma|^2=1$ makes $\det G_N$ an even differentiable
function of $\theta$, proving the derivative equality. These
statements hold for every original multiplicity and all three
actual sources.

The total phase equality preserves the individual mixed terms.
At the original source let
$g_K=\partial_\theta\log\det H_N|_0$,
$g_B=\partial_\theta\log\det Q_N|_0$, and
$\alpha_K=\operatorname{Tr}A_{KK}-c\dim K$.
The exact consequences are
\[
 \begin{gathered}
 z_B=-z_K,\qquad g_B=-g_K,\qquad
 z_K/\omega_N=\alpha_K-ig_K,\\
 \|x_K\otimes y_B-y_K\otimes x_B\|^2
           =a_Kd_B+a_Bd_K+2|z_K|^2,\\
 \Lambda^\Sigma=\Lambda\Sigma,\quad I^\Sigma=\Sigma I,\quad
 H_N^\Sigma=H_N,\quad Q_N^\Sigma=Q_N,\quad L_N^\Sigma=\Sigma L_N,\\
 A_{KK}^\Sigma=kI_K-A_{KK},\quad
 A_{BB}^\Sigma=kI_B-A_{BB},\quad
 A_{KB}^\Sigma=-A_{KB},\quad A_{BK}^\Sigma=-A_{BK}.
 \end{gathered}\tag{OMR25}
\]
The first identity follows from OMR16 and OMR24. For the second,
$[I,L_N(\theta)]$ has constant determinant: its second block
changes only by columns in $K$, so the frame change is triangular
with diagonal identity. The exact splitting determinant therefore
gives $g_K+g_B=\partial_\theta\log\det G_N|_0=0$.
The third identity is OMR22 and the mixed norm follows by
substitution in OMR16. For the reflected observation the metric
identities follow from OMR24 and the exact minimum section.
Substituting $\Sigma A\Sigma=kI-A$ into the four OMR15 blocks
proves their signs; both constant off-diagonal terms vanish
because $K_N^{\rm coord}L_N=0$ and $\Lambda I=0$.
These are the exact maps to the reflected observation even when
its kernel differs from $K$. Neither $z_K$, $\alpha_K$ nor $g_K$
has been set to zero.

In particular the untilted arithmetic budget in OMR19 becomes
\[
 \begin{gathered}
 J_k^{\rm action}=2\sum_{N=q-1}^{2q-1}w_N\log\epsilon_N^{\rm ar}
              =\Psi_k+C_k^{\rm budget},\qquad Z_k^{\rm phase}=0,\\
 \exp\left(\frac{J_k^{\rm action}}{4q}\right)
          =\prod_{N=q-1}^{2q-1}(\epsilon_N^{\rm ar})^{w_N/(2q)},\\
 4q\log\mathcal L_{h,k}\le J_k^{\rm action}\in[L_J,U_J],\qquad
 \mathcal L_{h,k}\le\min_N\epsilon_N^{\rm ar}
          \le\prod_N(\epsilon_N^{\rm ar})^{w_N/(2q)}
          \le e^{U_J/(4q)},\\
 \Phi_{k,k}\le
 \mathcal F_{\partial,k}^{(k)}+\mathcal H_k+\delta_k^\sigma
       +e_{k,k}^++C_k^{\rm budget}-4q\log\mathcal L_{h,k},\\
 \Phi_{k,k}\le
 \mathcal O_k^++\mathcal H_k+\delta_k^\sigma
       +e_{k,k}^++C_k^{\rm budget}-4q\log\mathcal L_{h,k}.
 \end{gathered}\tag{OMR26}
\]
All allowances are positive by OMR17. Therefore the logarithms
exist, substitution of the proved zero total phase gives the
first line, and exponentiation gives the exact weighted geometric
mean in the second. The third line follows from the minimum and
the same positive-class bound. The last two inequalities are the
proved OMR21 substitutions at $\theta=0$, retaining every other
term and its sign. The $[L_J,U_J]$ and outer substitutions keep
their original simple-quartet domain; the budget and geometric
mean identities retain full multiplicity. The generic MCE and
MCP tilted-family formulas remain in full, and OMR25 retains the
actual mixed phase still to be controlled.

% END CURRENT OMR

\clearpage
% BEGIN CURRENT PCL
\section{The literal original conductor limit and an invertible pivot restriction}
\label{sec:PCL}
Retain all original hypotheses and maps of PQO, RPC, OCF and CEP.
In particular $0<\delta<1/2$, $\gamma>2$, the original unit values
are $(a,\bar a,\bar a,a)$ with $a\ne0$, and the same logarithm
branch is used throughout. Write $\beta=2(\gamma^2-\delta^2)$,
$\eta=(\delta^2+\gamma^2)^2$, and
$\mathfrak a=1/160+\beta/24+\eta/2$. No source order is changed.

\subsection{An entire coefficient matrix in the inverse original period}
Let $P_0$ be RPC3's centered Fourier period matrix with its factor
$e^{\mathfrak a/u}$ removed by the stated exact identity, and retain
$G_u=\operatorname{diag}(u^{a/5}g_a)_{a=1}^4$ with RPC8's full
Gamma constants. Introduce the coefficient matrix $\mathcal R$ by
the following identity, rather than replacing the original period:
\[
 P= e^{\mathfrak a/u}P_0
   = e^{\mathfrak a/u}G_u\mathcal R(u^{-1}).
 \tag{PCL1}
\]
The row index below is $1\le a\le4$, and the column index is
$0\le r\le3$. The complete series is
\[
 \begin{split}
 \mathcal R_{a r}(z)&=\sum_{n=0}^{\infty}R_{a r,n}z^n,\\
 R_{a r,n}&=
 \sum_{\substack{p,h\ge0\\2p+4h=5n+r+1-a}}
 \frac{\beta^p\eta^h}{3^p p!h!}
      (-1)^{L}5^L(a/5)_L,
 \qquad L=p+h-n.
 \end{split}\tag{PCL2}
\]
Every sum at a fixed $n$ is finite. For each contributing pair,
$r+3p+h+1=a+5L$ and $L\ge0$. To prove the latter assertion,
the left side is positive and $1\le a\le4$, so a negative multiple
of five cannot be their difference. The exponent $n$ is nonnegative:
its numerator $2p+4h+a-r-1$ is an integer multiple of five and is
at least $-3$. Thus a negative value of $n$ is impossible.

For a direct derivation, divide every term of the absolutely
convergent original series PQO9 by $u^{a/5}g_a$. The quotient of
its Gamma factors is exactly
\[
 \frac{5^{(a+5L)/5-1}e^{\pi i(a+5L)/5}
                  \Gamma(a/5+L)}{g_a}
        =(-1)^L5^L(a/5)_L.
 \tag{PCL3}
\]
Its remaining power of $u$ is $u^{L-p-h}=u^{-n}$.
These identities prove PCL2 including its signs, factorials and
all unit-independent period coefficients. Absolute convergence of
PQO9 permits grouping these terms. For each prescribed complex
$z\ne0$, put $u=1/z$ on any of its original branches. The resulting
sum of absolute values is finite. Consequently PCL2 converges
absolutely for every complex $z$, and defines an entire matrix.
This also proves that this matrix is independent of the branch;
all branch factors remain in $G_u$ in PCL1.

The complete constant matrix in PCL2 is
\[
 J_\infty=
 \begin{pmatrix}
 1&0&-\beta/3&0\\
 0&1&0&-2\beta/3\\
 0&0&1&0\\
 0&0&0&1
 \end{pmatrix},\qquad
 J_\infty^{-1}=
 \begin{pmatrix}
 1&0&\beta/3&0\\
 0&1&0&2\beta/3\\
 0&0&1&0\\
 0&0&0&1
 \end{pmatrix}.
 \tag{PCL4}
\]
Indeed $n=0$ forces $2p+4h=r+1-a\le3$, so either $p=h=0$
or $p=1,h=0,r+1-a=2$. PCL3 gives respectively the diagonal
entries and the two negative entries displayed. In particular,
the original row-scaled period limit includes both cubic-phase
corrections. It is not a diagonal coefficient map.

RPC's determinant proof holds on the complete coefficient family
and gives $\det P_0=\det G_u$. Alternatively divide PQO11 by
the product of the four unchanged Gamma ray factors. Therefore
\[
       \det\mathcal R(z)=1,\qquad
       \mathcal R(z)^{-1}=\operatorname{adj}\mathcal R(z)
                       \quad(z\in\mathbb C).
 \tag{PCL5}
\]
The equality first holds for $z\ne0$ by PCL1 and extends to zero
by continuity. Both the matrix and its inverse are entire. For
any fixed $R>0$, the finite, completely specified constant
\[
 C_R=\sum_{n\ge1}\|R_n\|_1 R^{n-1}<\infty
 \tag{PCL6}
\]
gives $\|\mathcal R(z)-J_\infty\|_1\le C_R|z|$ on $|z|\le R$.
Here $\|\cdot\|_1$ is the induced column-sum matrix norm and
$R_n=(R_{a r,n})$. Convergence follows from the entire power
series in finite dimension, also at any larger radius. Thus
this is a bound in the actual coefficients, with no uncomputed
nonzero minor used as a premise.

\subsection{All eighty-one actual conductor coefficients and their limit}
Let $V$ be the original centered evaluation Vandermonde in PQO12,
$U=\operatorname{diag}(a,\bar a,\bar a,a)$, and
$D=\operatorname{diag}(\zeta,\zeta^2,\zeta^3,\zeta^4)$.
The complete period matrix from primary to Fourier coordinates is
\[
 \mathsf H(u)=e^{\mathfrak a/u}G_u\mathcal R(u^{-1})V^{-1}U.
 \tag{PCL7}
\]
For $1\le j\le4$ define the entire matrix and biform
\[
 \begin{split}
 L_j(z)&=U^{-1}V\mathcal R(z)^{-1}D^{-j}
                              \mathcal R(z)V^{-1}U,\\
 t(Z,W)&=(ZW,Z,W,1)^{\mathsf T},\\
 A_z(Z,W)&=\prod_{j=1}^4 Q_0(L_j(z)t(Z,W))
                  =\sum_{r,s=0}^8 a_{rs}(z)Z^rW^s.
 \end{split}\tag{PCL8}
\]
These are precisely all original OCF4 coefficients at $z=u^{-1}$.
To prove this, compute
\[
 \mathsf H(u)^{-1}D^{-j}\mathsf H(u)=L_j(u^{-1}):
 \tag{PCL9}
\]
the scalar $e^{\mathfrak a/u}$ cancels with its inverse, and the
two diagonal matrices $G_u,D^{-j}$ commute. OCF's definition
$Q_j(x)=Q_0(\mathsf H^{-1}D^{-j}x)$ followed by $x=\mathsf Ht$
then proves PCL8. The homogeneous version uses
$t=(z_0w_0,z_0w_1,z_1w_0,z_1w_1)$ and retains the bidegree
$(8,8)$. Thus no unit or period entry has been omitted by the
affine notation.

For an explicit finite formula for the individual coefficients,
write $L_{j,b e}$ for a matrix entry, and associate the four pairs
$(\rho_e,\sigma_e)=(1,1),(1,0),(0,1),(0,0)$ to its columns.
Put
\[
 \begin{split}
 c_{j,rs}(z)&=
 \sum_{\substack{1\le e,f\le4\\
       \rho_e+\rho_f=r,\ \sigma_e+\sigma_f=s}}
 (L_{j,1e}L_{j,4f}-L_{j,2e}L_{j,3f}),\\
 a_{rs}(z)&=
 \sum_{\substack{0\le r_j,s_j\le2\\
             \sum_jr_j=r,\ \sum_js_j=s}}
                       \prod_{j=1}^4c_{j,r_js_j}(z).
 \end{split}\tag{PCL10}
\]
Expanding the two products defining $Q_0$ proves the first line;
expanding its four original orbit factors proves the second.
Both lines are finite and retain all ordered-pair multiplicities.
Equations PCL2, PCL5 and PCL10 are a convergent formula for every
actual coefficient, with no freely assigned amplitude.

Their exact leading values are obtained by the explicit matrices
\[
 L_j(0)=U^{-1}VJ_\infty^{-1}D^{-j}J_\infty V^{-1}U,
 \qquad a_{rs}(u^{-1})=a_{rs}(0)+O_{\rm actual}(|u|^{-1}).
 \tag{PCL11}
\]
The equality for $L_j(0)$ follows from PCL4, and the bound follows
from the entire coefficient series. More explicitly, if
$a_{rs}(z)=\sum_{n\ge0}a_{rs,n}z^n$, the finite constant
$\sum_{r,s}\sum_{n\ge1}|a_{rs,n}|R^{n-1}$ bounds the sum of
all eighty-one errors divided by $|z|$ on $|z|\le R$.
The coefficient functions are single-valued entire functions of
$u^{-1}$ even though the original period itself retains its
logarithm branch. This assertion follows from the exact conjugation
PCL9, not from a change of the original branch.

At any specified $k$, the complete central convolution section is
therefore the explicitly evaluated entire matrix
\[
 C_{\rm in}(z)_{(a,b),(i,j)}=a_{i+4-a,j+4-b}(z),
 \qquad 0\le a,b,i,j\le k-8,
 \tag{PCL12}
\]
with all out-of-range coefficients zero. PCL10--12 give its full
constant matrix and all successive coefficients. They do not
discard off-diagonal coefficients or replace this matrix by its
central coefficient $a_{44}$. In particular the calculation of
the limit has not asserted a sign or diagonal-dominance inequality
for this complete original matrix.

\subsection{An exact positivity slice for the literal positive period ray}
For $u>0$ on any of its original logarithm lifts, PQO21 gives
$P_0=E_{\rm br}R$ with $R$ real invertible and $E_{\rm br}$ diagonal
unitary commuting with $D$. Its exact centered quadratic coefficient
matrix is $S$ of PQO13; all its entries are purely imaginary.
Consequently, for real $y\in\mathbb R^4$,
\[
 \overline{Q_j(E_{\rm br}y)}=-Q_{5-j}(E_{\rm br}y)
              \quad(1\le j\le4).
 \tag{PCL13}
\]
Indeed the original quadratic matrix in these coordinates is
$e^{-2\mathfrak a/u}R^{-\mathsf T}SR^{-1}$, which is purely
imaginary. Complex conjugation changes its sign and changes
$D^{-j}$ into $D^j$; this is precisely the displayed formula.

Let $\mathsf P_{\rm conj}$ be the permutation matrix exchanging
primary indices $1,2$ and $3,4$. For real $Z$ and $|W|=1$, direct substitution
gives $\mathsf P_{\rm conj}\overline{t(Z,W)}=W^{-1}t(Z,W)$.
The original roots and unit satisfy
$\bar V=\mathsf P_{\rm conj}V$ and
$\mathsf P_{\rm conj}\bar U=U\mathsf P_{\rm conj}$. Thus
\[
 \overline{V^{-1}Ut(Z,W)}=W^{-1}V^{-1}Ut(Z,W).
 \tag{PCL14}
\]
Choose either square root $w$ of $W$. Then
$p=w^{-1}V^{-1}Ut(Z,W)$ is real by PCL14, and
$\mathsf Ht=wE_{\rm br}y$ for the real vector
$y=e^{\mathfrak a/u}Rp$. Homogeneity and PCL13 prove the exact
nonnegative slice identity
\[
 \boxed{\quad
 W^{-4}A_{u^{-1}}(Z,W)
   =|Q_1(E_{\rm br}y)Q_2(E_{\rm br}y)|^2\ge0
       \quad(Z\in\mathbb R,\ |W|=1).\quad}
 \tag{PCL15}
\]
The choice of sign of $w$ has no effect because the quadratic
forms have even degree. In particular, comparing the Laurent
coefficients on the unit circle, and then the polynomial
coefficients in real $Z$, gives the exact symmetry
\[
             a_{rs}(u^{-1})=\overline{a_{r,8-s}(u^{-1})}.
 \tag{PCL16}
\]
PCL11 extends this identity and PCL15 to the coefficient limit
at $u^{-1}=0$ by continuity. The slice in PCL15 has one real
coordinate and one unit-circle coordinate; it has not changed
either coordinate into a second unit circle.

\subsection{A complete invertible restriction supplied by the original pivot}
At each actual $u$ in the proved five-orbit locus, use precisely
OCF5's lexicographically largest nonzero coefficient $(r_*,s_*)$,
$a_*=a_{r_*s_*}\ne0$, and its original pivot rectangle
$\mathcal P_k=\{(r_*+i,s_*+j):0\le i,j\le k-8\}$.
Let $\mathcal E_k$ be its complement and $S_k$ the original
insertion of its coordinates. No new coefficient has been selected.
The full original maps satisfy
\[
 C N_k^{\mathsf T}=0,\qquad
 S_k^{\mathsf T}N_k^{\mathsf T}=I_\Delta,\qquad
 \operatorname{im}N_k^{\mathsf T}=V_A.
 \tag{PCL17}
\]
The first is the transpose of $N_kC^{\mathsf T}=0$ in OCF7;
the second is the transpose of $N_kS_k=I$. Dimension $\Delta$
then proves the last equality, since $C$ has full row rank.
It follows that restriction to the actual strip is the exact
isomorphism
\[
 S_k^{\mathsf T}:V_A\longrightarrow\mathbb C^{\mathcal E_k},
        \qquad (S_k^{\mathsf T}|_{V_A})^{-1}=N_k^{\mathsf T}.
 \tag{PCL18}
\]
Equivalently its square pivot block has nonzero diagonal $a_*$
and is triangular in the elimination order; OCF's complete
division supplies its inverse. This is an alternative restriction
to the central one in PCL12, with the exact changed set of roots.

Let
$P_{\mathcal P,k}(y)=\prod_{(a,b)\in\mathcal P_k}(y-\omega_{ab})$
and $P_{\mathcal E,k}(y)=\prod_{(a,b)\in\mathcal E_k}(y-\omega_{ab})$,
where $\omega_{ab}=(2b-k)\gamma-i(2a-k)\delta$. Thus
$P_k=P_{\mathcal P,k}P_{\mathcal E,k}$ exactly. Let
$V_{\mathcal E}$ evaluate $1,y,\ldots,y^{\Delta-1}$ at those
actual outer roots and set
$D_{\mathcal P,\mathcal E}=\operatorname{diag}
(P_{\mathcal P,k}(\omega_{ab}))_{\mathcal E_k}$. These two matrices
are invertible by distinctness of the original roots. The full
polynomial-target restriction and its exact inverse are
\[
 \begin{split}
 \rho_{\mathcal E}
 &=V_{\mathcal E}^{-1}D_{\mathcal P,\mathcal E}^{-1}
                         S_k^{\mathsf T}|_{V_A},\\
 \rho_{\mathcal E}^{-1}
 &=N_k^{\mathsf T}D_{\mathcal P,\mathcal E}V_{\mathcal E}.
 \end{split}\tag{PCL19}
\]
Multiplying these matrices and using PCL17 proves both inverse
identities, with all inner products of root differences retained.
In particular the original residual becomes the following
specified onto map, with its exact complete kernel:
\[
 \Xi_{\mathcal E}=\Xi_k^{\rm fix}
       N_k^{\mathsf T}D_{\mathcal P,\mathcal E}V_{\mathcal E},
 \qquad\ker\Xi_{\mathcal E}=\rho_{\mathcal E}(K).
 \tag{PCL20}
\]
This follows by composition with the proved isomorphism, not by
deleting any common-kernel vector.

Finally the original metric and its entire source graph are carried
by this isomorphism as follows. Write
$J_{\mathcal E}=N_k^{\mathsf T}D_{\mathcal P,\mathcal E}V_{\mathcal E}$.
At every original endpoint the actual conductor metric in these
outer coordinates is exactly
\[
 H_{\mathcal E,N}=J_{\mathcal E}^*G_NJ_{\mathcal E},\qquad
 Q_{\Xi,N}^{\rm fix}
     =(\Xi_{\mathcal E}H_{\mathcal E,N}^{-1}
                      \Xi_{\mathcal E}^*)^{-1}.
 \tag{PCL21}
\]
The first formula is the pullback along $\rho_{\mathcal E}^{-1}$;
the minimum-section proof of GIP6 gives the second. On the full
GIP graph its target map is
$\rho_{\mathcal E}E_N^{\rm graph}$, with the same complete
$t$ ideal columns, and composing with $\Xi_{\mathcal E}$ gives
exactly $B_N^{\rm graph}$. Its residual kernel is therefore all
$t+r$ columns in GIP4, including the low-$v$ coordinates.
The coefficient-space bound for $N_k$ does not remove the displayed
root-value and Gamma matrices from PCL21. These formulas retain
the full metric transport needed for the next comparison.

\subsection{The exact inverse-period reflection of the central section}
Let $E_{\rm par}=\operatorname{diag}(1,-1,1,-1)$ in the centered
coefficient order, and let $\mathsf P_{\rm opp}$ exchange primary
indices $1,4$ and $2,3$. Equation PCL2 implies
\[
 \mathcal R(-z)=E_{\rm par}\mathcal R(z)E_{\rm par},\qquad
 V E_{\rm par}=\mathsf P_{\rm opp}V,\qquad
 U\mathsf P_{\rm opp}=\mathsf P_{\rm opp}U.
 \tag{PCL22}
\]
For the first equality a contributing coefficient satisfies
$5n+r+1-a=2p+4h$, so $n\equiv a-r-1\pmod2$. The signs from
the two parity matrices on entry $(a,r)$ are
$(-1)^{a-1+r}=(-1)^n$. Summing the entire series proves the
identity. The second equality evaluates each monomial at the
opposite original root. The third uses the full original unit
values $(a,\bar a,\bar a,a)$, which agree at opposite roots.

The two diagonal matrices $E_{\rm par}$ and $D$ commute. Substitution
into PCL8 and use of PCL22 therefore give the exact maps
\[
 L_j(-z)=\mathsf P_{\rm opp}L_j(z)\mathsf P_{\rm opp},\qquad
 A_{-z}(Z,W)=Z^8W^8 A_z(Z^{-1},W^{-1}).
 \tag{PCL23}
\]
For the second identity, $Q_0(\mathsf P_{\rm opp}t)=Q_0(t)$
by its two explicit products. Also
$\mathsf P_{\rm opp}t(Z,W)=ZWt(Z^{-1},W^{-1})$ when $ZW\ne0$.
Each of the four quadratic factors therefore contributes $(ZW)^2$.
This proves the identity on that open set; both sides are
polynomials after multiplication by $Z^8W^8$, so it holds
identically including both coordinate axes.

Comparing every coefficient in PCL23 proves
\[
 a_{rs}(-z)=a_{8-r,8-s}(z),\qquad
 C_{\rm in}(-z)=C_{\rm in}(z)^{\mathsf T},\qquad
 \det C_{\rm in}(-z)=\det C_{\rm in}(z).
 \tag{PCL24}
\]
Indeed the $(a,b),(i,j)$ entry on the left of the middle identity
is $a_{4+a-i,4+b-j}(z)$, precisely the transposed entry of
PCL12. Out-of-range coefficients are zero on both sides.
The determinant equality follows by the permutation expansion
of a determinant and its transpose. Thus the exact central
determinant is an even entire function of $u^{-1}$, with every
coefficient supplied by PCL2 and PCL10--12. This retains its
actual constant matrix and does not assign that determinant
a value by deleting any of its entries.

% END CURRENT PCL

\clearpage
% BEGIN CURRENT ROG
% Complete original relation operator; no source ideal or low component omitted.
\section{The original relation operator and its finite-rank graph transport}
\label{sec:ROG}

Retain the complete original data and maps of CEP1--49 and GIP1--9.
In particular the period and its logarithm branch are fixed on the
actual five-orbit locus, $k=4l+1\ge9$, and
\[
 q=(k+1)^2,\quad q'=(k-7)^2,\quad \Delta=q-q',\quad
 c=k/2,\quad c'=c-4,\quad 0\le v\le80.
 \tag{ROG1}
\]
All $81$ original coefficients $a_{ab}$ and their actual shifts
$b_{ab}=4+(2a-8)\delta+i(2b-8)\gamma$ are retained. Write
$\mu_j=\sum a_{ab}b_{ab}^j$, with $\mu_0=\cdots=\mu_{v-1}=0$
and $\mu_v\ne0$. For an original endpoint
$N\in\{q-1,q,2q-1,2q\}$ put
\[
 L=N-q,\qquad \mathfrak m=\Delta-v,\qquad
 M=L+\mathfrak m,\qquad d=N-v=q'+M.
 \tag{ROG2}
\]
Thus $L=-1,0,q-1,q$ at the four endpoints, respectively. The
polynomial spaces of degree at most $-1$ are zero. Since
$\Delta\ge96>v$, every displayed $M$ is nonnegative.

\subsection{The complete positive-shift relation operator}

Let $\chi=\chi_k$ and $\chi'=\chi_{k-8}$ denote the original
monic polynomials in $S$ and $S'$, respectively; the prime here
denotes the lower-grid polynomial and never a derivative.
For $0\le a,b\le8$ define the actual monic product
\[
 d_{ab}(S')=
 \prod_{(i,j)\notin\{a,\ldots,a+k-8\}\times\{b,\ldots,b+k-8\}}
       (S'+b_{ab}-\beta_{ij}),\qquad 0\le i,j\le k.
 \tag{ROG3}
\]
The root identity $\beta'_{ij}+b_{ab}=\beta_{i+a,j+b}$ gives
$\chi(S'+b_{ab})=\chi'(S')d_{ab}(S')$ by matching individual
factors, with the same leading coefficient one. Consequently
\[
 \begin{split}
 (\mathcal U_Lh)(S')
   &=\sum_{a,b=0}^8a_{ab}d_{ab}(S')h(S'+b_{ab}),\\
 \chi'\mathcal U_Lh&=\mathcal T_A(\chi h),\qquad
 \mathcal T_Ap=\sum_{a,b}a_{ab}p(S'+b_{ab}).
 \end{split}\tag{ROG4}
\]
These are polynomial equalities with positive shifts. They include
all original period coefficients and contain no conjugate or
binomial multiplier from the biform pairing: multiplication of
the original biform monomials adds their two indices with
coefficient one, and their complex-linear primary pairing evaluates
$p$ at the sum of their roots, exactly as in CEP3--4.

For $0\le j\le L$ write
\[
 u_j=\mathcal U_L(S^j),\qquad
 \lambda_j=\binom{q+j}{v}\mu_v.
 \tag{ROG5}
\]
Every $u_j$ has degree $\mathfrak m+j$ and leading coefficient
$\lambda_j$. Indeed Taylor expansion gives
$\mathcal T_Ap=\sum_{t\ge v}\mu_t p^{(t)}/t!$, a finite sum
on every polynomial. Applying it to the monic degree-$q+j$
polynomial $\chi S^j$ gives leading coefficient
$\binom{q+j}{v}\mu_v$ in degree $q+j-v$. Division by the
monic degree-$q'$ polynomial gives the asserted degree and
coefficient. In particular $\mathcal U_L$ is injective.

\subsection{An explicit finite inverse with coefficient control}

The following constants are expressed in the original coordinates:
\[
 \begin{gathered}
 R=|c|+k\sqrt{\delta^2+\gamma^2},\qquad
 R'=|c'|+(k-8)\sqrt{\delta^2+\gamma^2},\\
 B=\max_{a,b}|b_{ab}|,\qquad A_{\rm abs}=\sum_{a,b}|a_{ab}|,
 \qquad C_v=\frac{A_{\rm abs}B^v}{|\mu_v|}\ge1.
 \end{gathered}\tag{ROG6}
\]
All roots of $\chi$ and $\chi'$ have modulus at most $R$ and
$R'$, respectively. $B\ge4$ and $C_v\ge1$ follow directly
from the shift list and $|\mu_v|\le A_{\rm abs}B^v$.
For $L\ge0$ put
\[
 C_L=qR+2qB+(q'+L)R',\qquad
 \tau_L=\min\left\{1,\frac{\log(1+(2C_v)^{-1})}{C_L}\right\}>0.
 \tag{ROG7}
\]
For every $0\le h\le j\le L$, the complete coefficient estimate is
\[
 \left|[S'^{\mathfrak m+j-h}]\frac{u_j(S')}{\lambda_j}\right|
 \le C_v\frac{C_L^h}{h!}.
 \tag{ROG8}
\]
Here the coefficient for $h=0$ is exactly one, which will be
retained separately in the inverse.

To prove ROG8, write the coefficient of $S^{q-a}$ in $\chi$
as $\chi_a$. Its absolute value is bounded by
$\binom qa R^a\le(qR)^a/a!$, by expanding the product of
its actual roots. For $n=q+j$, the term of derivative order
$v+b$ applied to $\chi_aS^{n-a}$ contributes
$\chi_a\mu_{v+b}\binom{n-a}{v+b}$; it is zero when the
derivative order exceeds the degree. Its ratio to
$\mu_v\binom nv$ is bounded by
\[
 C_v\frac{(qR)^a}{a!}\frac{(nB)^b}{b!},
 \tag{ROG9}
\]
because
$|\mu_{v+b}/\mu_v|\le C_vB^b$ and
$\binom{n-a}{v+b}/\binom nv\le n^b/b!$.
For the latter assertion first remove $a$ using monotonicity
of a binomial coefficient in its top entry; then the ratio is
$\prod_{r=0}^{b-1}(n-v-r)/(v+1+r)\le n^b/b!$.

In the formal expansion at infinity,
\[
 \frac1{\chi'(S')}=S'^{-q'}\sum_{t\ge0}H_tS'^{-t},\qquad
 |H_t|\le\binom{q'+t-1}{t}(R')^t
       \le\frac{((q'+h)R')^t}{t!}\quad(t\le h).
 \tag{ROG10}
\]
Expansion of each $(1-\beta'_{ab}/S')^{-1}$ proves the
complete homogeneous coefficient formula and this estimate.
To obtain the specified coefficient of $u_j$, sum ROG9--10
over $a+b+t=h$. Since $n\le2q$ and $h\le L$, the
multinomial identity gives exactly ROG8. Formal division
agrees with ordinary polynomial division here because ROG4
proved exact divisibility before the expansion was made.

Let $H_L$ be the square matrix consisting of the coefficients
of $u_0,\ldots,u_L$ in degrees $\mathfrak m,\ldots,\mathfrak m+L$. Then
\[
 H_L=\widehat H_L\operatorname{diag}(\lambda_0,\ldots,\lambda_L),
 \qquad \det\widehat H_L=1,
 \qquad \det H_L=\mu_v^{L+1}\prod_{j=0}^L\binom{q+j}{v}.
 \tag{ROG11}
\]
The matrix $\widehat H_L$ is upper triangular, with diagonal
one and offset-$h$ entries bounded by ROG8 for $h\ge1$.
For $g$ in the actual image of $\mathcal U_L$, its unique
preimage is therefore the finite formula
\[
 \operatorname{coef}h=
 \operatorname{diag}(\lambda_0^{-1},\ldots,\lambda_L^{-1})
 \sum_{r=0}^L(I-\widehat H_L)^r
          ([S'^{\mathfrak m}]g,\ldots,[S'^{\mathfrak m+L}]g)^T.
 \tag{ROG12}
\]
The finite sum is its inverse because $I-\widehat H_L$ is
strictly upper triangular and its $(L+1)$st power is zero.
Thus this formula retains every coefficient of the original
operator, and every complex factor $\lambda_j$.

In increasing monomial coefficient Euclidean norms it satisfies
\[
 \|\operatorname{coef}h\|_2\le I_L\|\operatorname{coef}g\|_2,
 \quad
 I_L=\frac{1+2\sum_{h=1}^L\tau_L^{-h}}
                {|\mu_v|\binom qv}.
 \tag{ROG13}
\]
To check the bound, products of strictly upper triangular
entries at a total offset $h$ are majorized by the coefficient
of $z^h$ in
\[
 \sum_{r\ge0}\{C_v(e^{C_Lz}-1)\}^r
       =\frac1{1-C_v(e^{C_Lz}-1)}.
 \tag{ROG14}
\]
The absolute value of the expression inside braces is at most
$1/2$ on $|z|=\tau_L$, so Cauchy's coefficient formula bounds
these coefficients by $2\tau_L^{-h}$. The actual constant
diagonal is one. Every absolute row and column sum of the
inverse is consequently at most the numerator in ROG13.
The Euclidean operator norm is at most the square root of
the product of these row and column bounds, by applying
Cauchy--Schwarz in each row and summing columns. Finally
$|\lambda_j|\ge|\mu_v|\binom qv$ proves ROG13.

For fixed original period and quartet, $L\le q$ gives
\[
 \log^+ I_L=O_{\rm actual}(q\log(q+1)).
 \tag{ROG15}
\]
Indeed $R,R'=O(k)$, $B,C_v,\mu_v,v$ are fixed,
$C_L=O_{\rm actual}(qk)$, and the numerator in ROG13 is
at most $(2L+1)\max\{1,\tau_L^{-L}\}$. No lower bound
for a moment varying with the period is used.

\subsection{Complete original-source bounds for the low graph map}

For source $s\in\{1,k\}$ retain
\[
 b=s/2,\quad M_s=(2\pi)^{s/2},\quad h_n=M_sn!(b)_n,
 \qquad \|p\|_{s,c_0}^2=\int|p(c_0+iy)|^2dm_s(y).
 \tag{ROG16}
\]
The density $dm_s$ is exactly the original Gamma density in
GIP5 and OOQ2. We now supply elementary finite coefficient
bounds, so no implicit conditioning of a polynomial frame is
needed in the graph estimate.
For $D\ge0$ and real centre $c_0$ define
\[
 \begin{split}
 A^+_{D,s,c_0}
  &=\sqrt{M_s(D+1)}
      \max\{1,|c_0|+2\sqrt{D(b+D)}\}^{D},\\
 A^-_{D,s,c_0}
  &=\frac{\sqrt{D+1}}{\sqrt{M_s\min\{1,b\}}}
      \max\{1,(1+|c_0|)(1+D(b+D))\}^{D}.
 \end{split}\tag{ROG17}
\]
Then for every polynomial of degree at most $D$,
\[
 \|p\|_{s,c_0}\le A^+_{D,s,c_0}\|\operatorname{coef}p\|_2,
 \qquad
 \|\operatorname{coef}p\|_2\le A^-_{D,s,c_0}\|p\|_{s,c_0}.
 \tag{ROG18}
\]
For the first inequality, the original monic Gamma recurrence
gives in its orthonormal basis
\[
 y\varphi_n=\sqrt{(n+1)(b+n)}\varphi_{n+1}
             +\sqrt{n(b+n-1)}\varphi_{n-1}.
\]
On the degrees needed to multiply a polynomial of degree at
most $D-1$, the upper and lower shift norms are at most
$\sqrt{D(b+D)}$. Iterating the resulting bound for
$S=c_0+iy$, starting with $\|1\|=\sqrt{M_s}$, gives
$\|S^a\|\le\sqrt{M_s}(|c_0|+2\sqrt{D(b+D)})^a$
for $a\le D$. Cauchy--Schwarz on the coefficient sum proves
the first formula in ROG18.

For the second, the coefficient absolute sum of the monic
$p_n(y)$ is at most $(1+D(b+D))^n$ for $n\le D$.
The recurrence $p_{n+1}=yp_n-n(b+n-1)p_{n-1}$ proves this
by induction: with $A=D(b+D)$, the required scalar inequality
is $(1+A)+A\le(1+A)^2$. Substitution of $(S-c_0)/i$
multiplies the degree-$n$ coefficient bound by at most
$(1+|c_0|)^n$, retaining the phase $i^{-n}$ in the polynomial
itself. Also $h_n\ge M_s\min\{1,b\}$: it holds at
$n=0,1$, and $h_n/h_{n-1}=n(b+n-1)\ge1$ for $n\ge2$.
The Euclidean norm of the coefficient column of each
$p_n((S-c_0)/i)/\sqrt{h_n}$ is therefore bounded by the
corresponding term in ROG17. Expanding in that orthonormal
basis and applying Cauchy--Schwarz proves the second formula.

Define the explicit divisor constant
\[
 D_{M,\chi'}=\sum_{a=0}^{M}\binom{q'+a-1}{a}(R')^a.
 \tag{ROG19}
\]
For $f=\chi'g$ of degree at most $d=q'+M$, ROG10 gives
\[
 \|\operatorname{coef}g\|_2
       \le D_{M,\chi'}\|\operatorname{coef}f\|_2.
 \tag{ROG20}
\]
Indeed the high coefficients of the quotient are obtained by
convolution with the first $M+1$ coefficients of $1/\chi'$.
Every absolute row and column sum is bounded by ROG19;
the same row and column argument proves the norm assertion.
This bound is for the exact divisible input and retains the
whole divisor. It does not assert a contractive division map.

Equip $\mathbb C[S']_{\le M}$ with the weighted original norm
\[
 \|g\|_{\chi',s,c'}=\|\chi'g\|_{s,c'},\qquad
 G_M=\big(\langle\chi'S'^a,\chi'S'^b\rangle_{s,c'}\big)_{a,b=0}^M.
 \tag{ROG21}
\]
The Gram is positive definite because multiplication by the
nonzero polynomial is injective. Let $\mathbf U_L$ be the
complete coefficient matrix of ROG4. For $L\ge0$ the exact
weighted orthogonal projection onto its image is
\[
 \Pi_{U,L}=\mathbf U_L(\mathbf U_L^*G_M\mathbf U_L)^{-1}
                         \mathbf U_L^*G_M.
 \tag{ROG22}
\]
Its Gram inverse exists by ROG5. Multiplication verifies
idempotence, weighted self-adjointness, and its image, so it
has weighted norm one. At $L=-1$ put $\Pi_{U,L}=0$.

Let $Z_{\chi,N}h$ be the original upper orthonormal
coefficient column of $\chi h$, exactly as in GIP3. Define
the complete low graph map on this weighted space by
\[
 Q_{N,s}g=P_{<v}Z_{\chi,N}\,\mathcal U_L^{-1}\Pi_{U,L}g
      :\mathbb C[S']_{\le M}\longrightarrow\mathbb C^v.
 \tag{ROG23}
\]
The inverse is used only on its actual image, where ROG12
proved it uniquely. For $L=-1$ this map is zero. For $L\ge0$
its norm is at most the following entirely finite constant:
\[
 H_{N,s}=A^+_{N,s,c}(1+R)^q I_L D_{M,\chi'} A^-_{d,s,c'}.
 \tag{ROG24}
\]
To verify it, begin with $\Pi_{U,L}g$, whose weighted norm
does not exceed $\|g\|_{\chi',s,c'}$. Apply ROG18 to
$\chi'\Pi_{U,L}g$, then ROG20 and ROG13 to recover the
ordinary coefficients of its unique $h$. Multiplication by
$\chi$ has coefficient Euclidean norm at most $(1+R)^q$:
the sum of the absolute coefficients of its full monic product
is at most that number, and coefficient convolution has row
and column sums bounded by this sum. The first bound in ROG18
now gives the original upper source norm of $\chi h$.
Finally $P_{<v}$ is an orthogonal coordinate projection and
has norm at most one. These are precisely the factors in ROG24.

At every original endpoint with $L\ge0$,
\[
 \log(1+H_{N,s})=O_{\rm actual}(q\log(q+1)),
                    \qquad s=1,k.
 \tag{ROG25}
\]
For detail, $N,d\le2q$, $M\le2q$, $b\le k/2$ and
$|c|,|c'|,R,R'=O_{\delta,\gamma}(k)$. Both logarithms
of the powered factors in ROG17 are $O(q\log(q+1))$.
The factors $\sqrt{M_s}$ and $1/\sqrt{M_s}$ in ROG24
cancel at this specified composition; the individual source
norms ROG16--18 retain their full mass. The factor
$1/\sqrt{\min\{1,b\}}$ is at most $\sqrt2$.
ROG19 is bounded by
$(M+1)\max\{1,(q'+M)R'\}^{M}$, which has the same
logarithmic order. ROG15 and $q\log(1+R)=O(q\log(q+1))$
control the other two factors. These estimates prove ROG25.

\subsection{The exact graph shear and its complementary singular pairs}

Use exactly the GIP source space and product metric
\[
 X_N=\mathbb C^v\oplus\mathbb C[S']_{\le M},\qquad
 \|(x,g)\|^2=\|x\|_2^2+\|g\|_{\chi',s,c'}^2.
 \tag{ROG26}
\]
The full relation graph and its image under the specified
invertible map are
\[
 \begin{gathered}
 \mathcal R_{N,s}(x,g)=(x-Q_{N,s}g,g),\qquad
 \mathcal R_{N,s}^{-1}(x,g)=(x+Q_{N,s}g,g),\qquad
 \det\mathcal R_{N,s}=1,\\
 \mathcal R_{N,s}\operatorname{im}
       \begin{pmatrix}P_{<v}Z_{\chi,N}\\\mathbf U_L\end{pmatrix}
       =\{0\}\oplus\operatorname{im}\mathbf U_L .
 \end{gathered}\tag{ROG27}
\]
The last equality follows by applying ROG23 to $\mathcal U_Lh$;
the actual low-degree column is subtracted with its full value.
The inverse and determinant follow from the displayed triangular
blocks. Thus ROG27 is an exact morphism between the full
original graph and the indicated relation image.

Put $r_Q=\operatorname{rank}Q_{N,s}\le v$ and let its nonzero
singular values in the specified metrics be $\sigma_a$.
Singular value decompositions on its domain and codomain put
the shear into $r_Q$ orthogonal two-dimensional blocks
$\left(\begin{smallmatrix}1&-\sigma_a\\0&1\end{smallmatrix}\right)$
and identity blocks. The two singular values of such a block
are exactly
\[
 t_a=\frac{\sqrt{\sigma_a^2+4}+\sigma_a}{2},\qquad t_a^{-1}.
 \tag{ROG28}
\]
Indeed its positive Gram has determinant one and trace
$2+\sigma_a^2$; solving its quadratic characteristic
polynomial gives these positive square roots. In particular
$1\le t_a\le1+\sigma_a\le1+H_{N,s}$.
For every exterior degree $e$, including zero and full degree,
both the forward and inverse shear therefore satisfy
\[
 \log\|\bigwedge^e\mathcal R_{N,s}^{\pm1}\|
       \le\min\{e,v\}\log(1+H_{N,s}).
 \tag{ROG29}
\]
There are at most $v$ singular values above one; any remaining
factors are at most one. This is the complementary-singular-value
estimate with the actual rank, rather than multiplying the
scalar inverse bound by the full source dimension. At $v=0$
the map is the identity and this error is exactly zero.

\subsection{The unchanged residual target and its four endpoint metrics}

Retain the complete onto residual map $B_N^{\rm graph}$ and
its common-kernel lift columns $W_{K,N}$ from GIP2--4. Define
\[
 B_{N,s}^{\rm str}=B_N^{\rm graph}\mathcal R_{N,s}^{-1},\qquad
 \ker B_{N,s}^{\rm str}
   =\operatorname{im}
      \left[\begin{pmatrix}0\\\mathbf U_L\end{pmatrix},
                   \mathcal R_{N,s}W_{K,N}\right].
 \tag{ROG30}
\]
The kernel equality is the image under the invertible shear
of the full GIP4 kernel. It retains all $r=8k-16$ original
common-kernel columns and all $L+1$ relation columns; at
$L=-1$ the latter list is empty. Let $H_N^{\rm graph}$ be
the product Gram of ROG26, which is exactly GIP5, and put
\[
 Q_{N,s}^{\rm str}
    =[B_{N,s}^{\rm str}(H_N^{\rm graph})^{-1}
                      (B_{N,s}^{\rm str})^*]^{-1}.
 \tag{ROG31}
\]
This is the attained minimum metric on the same original
residual coordinate space $\mathbb C^j$, $j=8k-32$.
For an onto matrix $B$ and positive $H$, its minimum lift
of $y$ is $H^{-1}B^*(BH^{-1}B^*)^{-1}y$; subtraction of
this lift from any other lift gives an $H$-orthogonal kernel
vector. This proves the inverse-Gram assertion directly.

The shear induces the identity of the residual targets with
these two quotient norms. A minimum lift is sent to its
shear image and then projected off the entire target kernel.
Orthogonal projection contracts every exterior power because
its singular values are zero or one. Apply this observation
also to the inverse shear and then square the two volume
bounds in ROG29. It proves
\[
 |\log\det Q_{N,s}^{\rm str}
        -\log\det\widetilde Q_{\Xi,N}|
       \le2\min\{j,v\}\log(1+H_{N,s}).
 \tag{ROG32}
\]
Here $\widetilde Q_{\Xi,N}$ is precisely GIP6 with its
original complete graph metric. At $L=-1$ set $H_{N,s}=0$;
the two metrics are then identical.
Combining with the already proved GIP6 source-map estimate
gives the fully specified original residual comparison
\[
 |\log\det Q_{N,s}^{\rm str}
       -\log\det Q_{\Xi,N}^{\rm fix}|
 \le2E_N^{\rm CEP}+2v\log(1+H_{N,s}).
 \tag{ROG33}
\]
Thus, with the unchanged signs $(1,1,-1,-1)$ and original
endpoints,
\[
 \left|\sum_N\sigma_N\log\det Q_{N,s}^{\rm str}
            -\sum_N\sigma_N\log\det Q_{\Xi,N}^{\rm fix}\right|
 \le2\sum_NE_N^{\rm CEP}
       +2v\sum_N\log(1+H_{N,s})=o(kq).
 \tag{ROG34}
\]
CEP36 gives the first estimate, while ROG25 and fixed
$v\le80$ give $O_{\rm actual}(q\log(q+1))=o(kq)$ for
the second. This proves the transport for both complete
original sources without removing a low component by assumption.

The resulting relation quotient is, with its product quotient norm,
\[
 \mathbb C^v\oplus
   \left(\mathbb C[S']_{\le L+\Delta-v}/
                    \mathcal U_L\mathbb C[S]_{\le L}\right).
 \tag{ROG35}
\]
The residual map on it is exactly the descended ROG30 map.
In particular the second denominator remains the full operator
ROG4; no equality with multiplication by the outer divisor is
used in ROG1--35. Its complete leading coefficient and exact
inverse are ROG5 and ROG12. The natural map to that outer
divisor and its exact kernel and cokernel remain GIP7--9;
their finite-convolution obstruction is carried unchanged by
ROG27 and ROG30. The new finite-rank transport proves that
the full low-$v$ graph itself has only the explicit subleading
cost ROG34 in the original residual return.

% END CURRENT ROG

\clearpage
% BEGIN CURRENT RDS
% Complete relation-annihilator realization in the original finite symbol.
\section{The exact differential realization of the original relation quotient}
\label{sec:RDS}

Retain the original $\chi_k(S)$, $\chi_{k-8}(S')$, all $81$
coefficients and shifts of CEP41--42 and ROG1--6. Write
$\chi=\chi_k$, $\chi'=\chi_{k-8}$, where the prime denotes
the lower-grid polynomial, and put $\mathfrak m=q-q'-v$.
The cutoff-independent polynomial map is
\[
 \mathcal U:\mathbb C[S]\longrightarrow\mathbb C[S'],\qquad
 \mathcal Uh=\sum_{a,b=0}^8a_{ab}d_{ab}(S')h(S'+b_{ab}).
 \tag{RDS1}
\]
It restricts to every original $\mathcal U_L$. For all $j\ge0$,
$u_j=\mathcal U(S^j)$ has degree $\mathfrak m+j$ and leading
coefficient $\lambda_j=\binom{q+j}{v}\mu_v\ne0$ by exactly
the Taylor and monic-division proof in ROG4--5, which holds at
every finite polynomial degree.

\subsection{A finite-dimensional quotient and its analytic dual}

Repeated subtraction of the unique highest-degree multiple
$u_j/\lambda_j$ gives the direct vector-space decomposition
\[
 \mathbb C[S']=\mathbb C[S']_{<\mathfrak m}
                         \oplus\mathcal U\mathbb C[S].
 \tag{RDS2}
\]
At each step the degree strictly decreases. No nonzero sum
$\mathcal Uh$ can have degree below $\mathfrak m$, by RDS1,
which proves directness and uniqueness. Thus the resulting
remainder map is linear and independent of the polynomial
cutoff. Its restrictions identify every CEP43 quotient with
the same vector space of dimension $\mathfrak m$.

For a complex-linear functional $\ell$ on polynomials use its
exponential generating series
\[
 f_\ell(z)=\sum_{n\ge0}\ell(S'^n)\frac{z^n}{n!}.
 \tag{RDS3}
\]
Initially this is a formal series. The subspace annihilating
$\operatorname{im}\mathcal U$ has dimension $\mathfrak m$:
its first $\mathfrak m$ values are arbitrary, and the equations
$\ell(u_j)=0$ uniquely determine the value in degree
$\mathfrak m+j$ from earlier ones, since its coefficient is
the nonzero $\lambda_j$. We now realize this whole formal
space as convergent series using the original symbol.

Put $D=\partial_z$ and define
\[
 E_A(z)=\sum_{a,b=0}^8a_{ab}e^{b_{ab}z},\qquad
 \mathscr L_A=\sum_{a,b=0}^8a_{ab}e^{b_{ab}z}d_{ab}(D).
 \tag{RDS4}
\]
The coefficients multiplying powers of $D$ are on their
displayed left side. The precise adjoint and intertwining
identities, in this formal-series pairing, are
\[
 \ell(\mathcal U e^{Sz})=\mathscr L_Af_\ell(z),\qquad
 \mathscr L_A\chi'(D)=\chi(D)M_{E_A}.
 \tag{RDS5}
\]
For the first, multiplication of $S'^r$ inside the functional
is differentiation $D^r$, while the positive shift contributes
$e^{b_{ab}z}$. Thus every term is exactly
$a_{ab}e^{b_{ab}z}d_{ab}(D)f_\ell$.
For the second, the full polynomial product in CEP41 gives
$d_{ab}(D)\chi'(D)=\chi(D+b_{ab})$. Leibniz' rule gives
$\chi(D)(e^{b_{ab}z}F)=e^{b_{ab}z}\chi(D+b_{ab})F$
for every formal $F$. Summing these identities proves RDS5,
including the operator order and all coefficients.

Define the two original exponential spaces and a jet kernel by
\[
 \begin{gathered}
 \mathscr E_k=\operatorname{span}\{e^{\beta_{ij}z}:0\le i,j\le k\},
 \qquad
 \mathscr E_{k-8}=\operatorname{span}\{e^{\beta'_{ab}z}:0\le a,b\le k-8\},\\
 \mathscr N_v=\{N\in\mathscr E_k:N^{(a)}(0)=0\ (0\le a<v)\}.
 \end{gathered}\tag{RDS6}
\]
The $q$ upper and $q'$ lower roots are distinct. Their full
Vandermonde derivative matrices at zero prove independence of
the respective exponentials. The first $v$ upper derivative
rows have rank $v$, since any chosen $v$ distinct root columns
give a nonzero Vandermonde determinant. Hence
$\dim\mathscr N_v=q-v$, also when $v=0$.

The original sum identity gives the injective multiplication map
\[
 M_{E_A}:\mathscr E_{k-8}\hookrightarrow\mathscr N_v,
 \quad
 E_A(z)e^{\beta'_{ab}z}
       =\sum_{r,s=0}^8a_{rs}e^{\beta_{a+r,b+s}z}.
 \tag{RDS7}
\]
Its numerical coefficient matrix is exactly $C^{\mathsf T}$
from CEP4, with the original order and no conjugation.
The product vanishes to order at least $v$ at zero because
$E_A$ has that exact zero order. Multiplication by this
nonzero analytic function is injective: if its product with
an analytic function is zero, that function vanishes on a
punctured neighbourhood with nonzero $E_A$ and hence identically.

There is the exact sequence and specified quotient isomorphism
\[
 \begin{gathered}
 0\longrightarrow\mathscr E_{k-8}
    \xrightarrow{M_{E_A}}\mathscr N_v
    \xrightarrow{\mathfrak F_A}
                 (\mathbb C[S']/\mathcal U\mathbb C[S])^\vee
    \longrightarrow0,\\
 f_{\mathfrak F_A(N)}(z)=\chi'(D)\left(\frac{N(z)}{E_A(z)}\right).
 \end{gathered}\tag{RDS8}
\]
The ratio is analytic at zero because its numerator vanishes
to order at least $v$, the exact order of its denominator.
It defines a functional by the Taylor derivatives in RDS3.
RDS5 shows that this functional annihilates every $u_j$,
since $\chi(D)N=0$ for the original upper exponential space.
If its image is zero, $F=N/E_A$ solves $\chi'(D)F=0$ near
zero. Every analytic solution of this constant-coefficient
equation is in $\mathscr E_{k-8}$: its first $q'$ derivatives
determine all later derivatives by the monic equation, and
the $q'$ independent displayed exponentials realize arbitrary
initial values by the invertible Vandermonde. This proves
$N\in E_A\mathscr E_{k-8}$. The reverse kernel inclusion
follows by applying $\chi'(D)$ to those exponentials.
The image dimension is consequently
$q-v-q'=\mathfrak m$, equal to the full annihilator dimension
proved after RDS3, so the map is onto. This also proves that
every formal annihilating series converges near zero.

The transpose of the injection in RDS7 is the exact primary
conductor $C$; this specifies the direction of the dual maps.
Thus RDS8 relates the original polynomial relation quotient
to the full original exponential module by an explicit
quotient and differential map, without a new polynomial ideal.

\subsection{Every possible singularity and its original multiplicity}

The nonzero entire function $E_A$ has isolated zeros. For
$N\in\mathscr N_v$, the formula in RDS8 extends its series
meromorphically to the whole $z$ plane. Every possible pole
lies at a zero of this same fixed original finite symbol;
zero itself is removable. If $\zeta\ne0$ is a zero of
order $h$, put
\[
 e_\zeta(w)=w^{-h}E_A(\zeta+w),\qquad
 c_{-\nu}= [w^{h-\nu}]\frac{N(\zeta+w)}{e_\zeta(w)}
                  \quad(1\le\nu\le h).
 \tag{RDS9}
\]
Here $e_\zeta$ is analytic and nonzero at zero, so these
are its actual finite Taylor coefficients. Write
$\chi'(X)=\sum_{r=0}^{q'}\chi'_rX^r$, with the coefficients
in increasing order in this display. The complete principal
part of the function in RDS8 at $\zeta$ is
\[
 \sum_{\nu=1}^h\sum_{r=0}^{q'}
       \chi'_r(-1)^r(\nu)_r c_{-\nu}
                        (z-\zeta)^{-\nu-r}.
 \tag{RDS10}
\]
Laurent expansion of $N/E_A$ gives its principal part
$\sum c_{-\nu}(z-\zeta)^{-\nu}$; differentiating a term
$r$ times gives $(-1)^r(\nu)_r(z-\zeta)^{-\nu-r}$.
Its holomorphic part remains holomorphic under differentiation,
which proves the entire formula. Pole cancellations are retained
by the actual coefficients, and the pole order is at most
$h+q'$. More precisely, let $t$ be the vanishing order of
$N$ at $\zeta$. If $t\ge h$, the quotient and its image
under $\chi'(D)$ are holomorphic. If $t<h$, the leading
Laurent coefficient of $N/E_A$ has order $-(h-t)$ and is
nonzero. The unique highest-order term after applying the
monic $\chi'(D)$ is its $q'$th derivative: it has order
$-(h-t+q')$ and leading coefficient multiplied by
$(-1)^{q'}(h-t)_{q'}\ne0$. Every lower derivative and
every subsequent Laurent term has strictly smaller pole
order. Thus the exact pole order is $h-t+q'$ in this case.
The identically zero numerator gives the zero function.
No simplicity of a symbol zero is used.

\subsection{An explicit original-data analytic disk and source coefficients}

Retain $B,A_{\rm abs},C_v,R$ of ROG6 and set
\[
 R_0=\frac1B\min\left\{1,\frac{v+1}{2eC_v}\right\}>0.
 \tag{RDS11}
\]
The exact entire quotient
$B_\infty(z)=v!E_A(z)/(\mu_vz^v)$, with value one at zero,
obeys $|B_\infty(z)-1|\le1/2$ on $|z|\le R_0$.
Indeed its remaining coefficient sum is at most
$C_vB|z|e^{B|z|}/(v+1)$ by expanding each original
exponential and using
$v!/(v+1+j)!\le1/((v+1)j!)$. Definition RDS11 gives
the asserted bound. This proves that this disk has no
nonzero zero of $E_A$.

Write $N=\sum\nu_{ij}e^{\beta_{ij}z}$. Its first $v$
vanishing Taylor coefficients and the elementary remainder
bound for the exponential give
\[
 |N(z)|\le\|\nu\|_1\frac{(R|z|)^v}{v!}e^{R|z|},
 \qquad
 \left|\frac{N(z)}{E_A(z)}\right|
      \le\frac{2R^v}{|\mu_v|}e^{RR_0}\|\nu\|_1
                      \quad(|z|\le R_0).
 \tag{RDS12}
\]
The removable value at zero follows by continuity. Cauchy's
derivative formula on this disk gives, for $|z|\le R_0/2$,
the explicit estimate
\[
 |f_{\mathfrak F_A(N)}(z)|\le C_{k,A}\|\nu\|_1,
 \quad
 C_{k,A}=\frac{2R^v}{|\mu_v|}\,q'!
       \left(\frac2{R_0}\right)^{q'}
                    \exp(RR_0+R'R_0/2).
 \tag{RDS13}
\]
For verification, the derivative of order $r$ of $N/E_A$
is bounded by $r!(2/R_0)^r$ times the bound in RDS12.
The coefficient of $X^{q'-j}$ in $\chi'(X)$ has modulus
at most $\binom{q'}j(R')^j$. Summing the derivative bounds
therefore gives
$q'!(2/R_0)^{q'}\sum_{j=0}^{q'}(R'R_0/2)^j/j!$,
which proves RDS13. Its logarithm is
$O_{\rm actual}(q\log(q+1))$ at the fixed original period.

Finally, for any polynomial functional the following identity
is formal. For $\ell=\mathfrak F_A(N)$ its right side is
analytic near zero by RDS8 and RDS13, so the identity also
proves convergence of its left side. For either unchanged
source order $s=1,k$ it is
\[
 \sum_{n\ge0}\ell\left(p_n\left(\frac{S'-c'}i\right)\right)
                     \frac{z^n}{n!}
  =(1+z^2)^{-s/4}e^{ic'\arctan z}f_\ell(-i\arctan z).
 \tag{RDS14}
\]
It follows by applying the complex-linear functional to the
original Gamma generating series; the centre and phase are
the literal equality $(S'-c')/i=-iS'+ic'$. Let
$\rho=\tanh(R_0/4)$. On $|z|=\rho$,
$|\arctan z|\le\operatorname{arctanh}\rho=R_0/4$,
so RDS13 and Cauchy's coefficient formula yield
\[
 \left|\ell\left(\frac{p_n((S'-c')/i)}{\sqrt{h_n}}\right)\right|
 \le\sqrt{\frac{n!}{M_s(s/2)_n}}\,\rho^{-n}
       (1-\rho^2)^{-s/4}e^{|c'|R_0/4}C_{k,A}\|\nu\|_1.
 \tag{RDS15}
\]
The square root uses the full original
$h_n=M_sn!(s/2)_n$ and $M_s=(2\pi)^{s/2}$.
This identity gives the source coefficient realization of the
actual relation-annihilating functionals. The weighted relation
metric itself remains ROG21, including multiplication by
the full $\chi'$. Equations RDS8--15 provide its exact
finite-symbol analytic model and quantitative disk; the
rank-$O(k)$ residual determinant is not assigned a value by
multiplying this scalar coefficient bound by its rank.

% END CURRENT RDS

\clearpage
% BEGIN CURRENT FEX
% FEX: new written continuation, 15 September 2026.
\section{The fixed-period problem and the result proved here}
\label{sec:FEX-scope}
The current handoff already proves the absolute $O(kq)$ kernel estimate,
the conductor frame, its residual observation, and the finite inverse.
The new calculation here removes the $O(kq)$ \emph{exterior-inverse
allowance} left by AIE29. It replaces that allowance by an explicit
$O_{\rm actual}(q\log q)=o(kq)$ error and transfers the actual
four-endpoint determinant to the smaller conductor coefficient space.
It does not identify this transfer error with the original kernel
volume, and it does not evaluate the latter's leading coefficient.

Retain an actual hypothetical simple quartet
\[
 \rho=\tfrac12+\delta+i\gamma,\qquad0<\delta<\tfrac12,\quad\gamma>2,
 \qquad h(z)=\prod_{\alpha\in\{\rho,\bar\rho,1-\rho,1-\bar\rho\}}(z-\alpha),
 \quad v_h=2\xi/h.
\]
The original full unit and period observation are
\[
 U=M_{j_hv_h},\qquad
 \jmath\Lambda_k=P_k\Pi^{\otimes k}U^{\otimes k}\iota_k,
 \qquad P_k=\frac15\sum_{b=0}^4(C^{-b})^{\otimes k}.
 \tag{FEX1}
\]
Fix one original period and logarithm branch on the proved five-orbit
locus, for example any fixed point with $|u|\ge R_*$. Here $R_*$ is the
unchanged RPC/AMT radius. Its formulas and the proof of the rank are in
the supplied complete source; no new rank premise is imposed.
Throughout the principal application,
\[
\begin{gathered}
 k=4l+1\ge9,\quad q=(k+1)^2,\quad c=k/2,\quad k'=k-8,\quad c'=c-4,\\
 q'=(k-7)^2,\quad a=q-q'=16k-48,\quad r=8k-16,\quad j_k=8k-32,\\
 \beta_{ab}^{(k)}=c+(2a-k)\delta+i(2b-k)\gamma,\quad
 \chi=\prod_{a,b=0}^k(S-\beta_{ab}^{(k)}),\quad \psi=\chi_{k'}.
\end{gathered}\tag{FEX2}
\]
The letter $a$ without indices denotes the degree difference, not an
amplitude. Let $K=\ker\Lambda_k\subset E=\mathbb C[S]/(\chi)$, using its
fixed original coefficient inclusion. The actual ranks satisfy
$a=r+j_k$. Every class below has its unique representative of degree
less than $q$ whenever a polynomial representative is required.

Retain the complete original conductor biform and its coefficients
\[
A=\sum_{r_0,s_0=0}^8a_{r_0s_0}X_1^{r_0}X_0^{8-r_0}Y_1^{s_0}Y_0^{8-s_0}.
\]
Write $\nu=(r_0,s_0)$ for an index in this sum. Define
\[
\begin{gathered}
 b_\nu=4+(2r_0-8)\delta+i(2s_0-8)\gamma,\quad
 \mu_j=\sum_\nu a_\nu b_\nu^j,\quad
 v=\min\{j:\mu_j\ne0\}\le80,\\
 (\mathcal T_Ap)(S')=\sum_\nu a_\nu p(S'+b_\nu),\qquad
 \mathcal C_A=\frac{v!}{\mu_v}\mathcal T_A,\\
 B_q(X)=\sum_{n=0}^{q-1-v}
       \frac{v!\mu_{v+n}}{(v+n)!\mu_v}X^n.
\end{gathered}\tag{FEX3}
\]
The coefficients are those of the actual $\mathsf H=F^{-1}\Pi U
\mathcal E_{\rm prim}$, fixed as $k$ changes. The original distinct-node
Vandermonde proves the stated finite $v$ and $\mu_v\ne0$. In particular,
$B_q(0)=1$. The scalar in $\mathcal C_A$ is displayed explicitly; it is
not a renormalization of an arithmetic source or observation.

Only the two original Gamma sources are used:
\[
\begin{gathered}
 s\in\{1,k\},\quad b=s/2,\quad M_s=(2\pi)^{s/2},\quad
 dm_{0,s}(y)=\frac{M_s2^b}{4\pi\Gamma(b)}
                 |\Gamma(b/2+iy/2)|^2\,dy,\\
 h_n=M_s n!(b)_n,\quad w_n=\frac{n!}{(b)_n},\qquad
 \|p\|_{s,c_0}^2=\int |p(c_0+iy)|^2\,dm_{0,s}(y).
\end{gathered}\tag{FEX4}
\]
In particular $s=k'$ is never inserted as a source order. The polynomial
$\psi$ in FEX2 will enter as a literal multiplication factor in these
unchanged measures.

The main conclusion will be the finite interval FEX35 and the consequence
\[
 \boxed{\quad F_K^{(s)}=\widehat F_k^{(s)}
          +O_{\mathsf H,\delta,\gamma}(q\log q),\qquad s=1,k.\quad}
 \tag{FEX5}
\]
Here $\widehat F$ is defined in FEX29--30 using the actual residual
observation and all original transformed relation columns. The error
is proved, not postulated. At the end, the unchanged arithmetic
comparison transfers the same leading-limit question to $F_K^{\rm ar}$.

\section{A full-flag determinant bound}
\label{sec:FEX-flag}
Let $T:(V,g)\to(V',g')$ be an isomorphism of $n$-dimensional complex
Hilbert spaces. Write $\|T\|\le M$ and let $J>0$ be its squared whole-space
Jacobian. For $R\subset W\subset V$, set $t=\dim W-\dim R$. Give $W/R$
and $TW/TR$ their actual Hilbert quotient norms. For any coefficient
basis of $W/R$, let $D_Q$ be the ratio of its target and source Gram
determinants under the induced map. Then
\[
 \boxed{\qquad J M^{-2(n-t)}\le D_Q\le M^{2t}.\qquad}
 \tag{FEX6}
\]
All empty factors have determinant one. The upper bound follows by
applying $T$ to each lift and taking the infimum over the same fibre.
For the lower bound, use the flag $0\subset R\subset W\subset V$.
Choose any basis adapted to it, on both sides using its transported
basis. Successive block elimination of the complete positive Gram gives
\[
 J=D_R D_Q D_{V/W}.
 \tag{FEX7}
\]
Each block elimination uses its full Schur complement. The other two
induced maps also have norm at most $M$, by restriction or minimization,
and their dimensions total $n-t$. Thus their squared determinant product
is at most $M^{2(n-t)}$, which proves FEX6. This proves the assertion
without asserting orthogonality of a fixed coefficient section.

An equivalent inverse form is useful. If $T$ is an endomorphism of a
single Hilbert space and $|\det T|=1$, then, for every full-rank $n$ by
$t$ coefficient matrix $X$,
\[
 \boxed{\quad
 M^{-2t}\le
 \frac{\det(X^*T^{-*}T^{-1}X)}{\det(X^*X)}
 \le M^{2(n-t)}.
 \quad}\tag{FEX8}
\]
For example, this follows by taking $W=T^{-1}\operatorname{im}X$, $R=0$
in FEX6 and inverting the determinant ratio. It also follows by taking
the product of the largest $t$ inverse singular values and cancelling
against the full product. The statement is a bound: $\det T=1$ does
\emph{not} say that a restricted Gram determinant is unchanged.
For $T=\left(\begin{smallmatrix}1&4\\0&1\end{smallmatrix}\right)$,
restriction to the second coordinate has squared volume $17$.

\section{Polynomial forward control replaces the exponential exterior loss}
\label{sec:FEX-forward}
For completeness, the derivative bound used here follows directly from
the original polynomial norms. The recurrence and generating identity are
\[
\begin{gathered}
 p_0=1,\quad p_1=y,\quad
 p_{n+1}=yp_n-n(b+n-1)p_{n-1},\\
 \sum_{n\ge0}\frac{p_n(y)}{n!}z^n
       =(1+z^2)^{-b/2}e^{y\arctan z},\qquad \|p_n\|^2=h_n.
\end{gathered}\tag{FEX9}
\]
They are the supplied original Gamma identities. Differentiation of the
finite coefficient at zero, followed by the full normalization, gives
\[
 \partial_y\varphi_n=
 \sum_{\substack{1\le j\le n\\j\ \mathrm{odd}}}
 \frac{(-1)^{(j-1)/2}}{j}\sqrt{\frac{w_n}{w_{n-j}}}\,
 \varphi_{n-j},\qquad \varphi_n=p_n/\sqrt{h_n}.
 \tag{FEX10}
\]
The mass $M_s$ cancels from this ratio, not from the underlying norms.
Let $H_D=\sum_{j=1}^D1/j$ and $H_0=0$. For $n>m\ge1$, $b\ge1/2$ gives
$w_n/w_m\le\sqrt{n/m}$ and $w_n\le2\sqrt n$. Indeed square each factor
$j/(j-1/2)$ and use $j(j-1)\le(j-1/2)^2$; the product telescopes.
Apply the weighted Schur estimate to the absolute matrix in FEX10 with
weights $w_n^{-1/2}$. Its row sum is at most $H_D$; its column sum at $n$
is at most $H_n+4$. The excess for $1\le m<n$ is bounded by
\[
 \frac{\sqrt{n/m}-1}{n-m}
 =\frac1{\sqrt m(\sqrt n+\sqrt m)}\le\frac1{\sqrt{nm}},
\]
whose sum is at most two; the $m=0$ term is at most $2/\sqrt n\le2$.
Enlarging the odd-index sums is legitimate because all terms are
nonnegative. Cauchy--Schwarz in each weighted row and then summing the
column estimate proves
\[
 \|\partial_y\|_{\mathcal P_D}\le\sqrt{H_D(H_D+4)}\le H_D+2,
 \qquad \|f(y+z)\|\le e^{|z|(H_D+2)}\|f\|.
 \tag{FEX11}
\]
At $D=0$ the derivative is zero. Translation is the exact finite Taylor
exponential. The coordinate substitution $S=c_0+iy$ gives
$\partial_S=-i\partial_y$ with the same norm.

Retain every actual coefficient in
\[
 A_*^{\rm abs}=\sum_\nu|a_\nu|,\quad B_*^{\rm sh}=\max_\nu|b_\nu|,
 \quad C_v=\frac{A_*^{\rm abs}(B_*^{\rm sh})^v}{|\mu_v|}\ge1,
 \qquad P_D=C_v e^{B_*^{\rm sh}(H_D+2)}.
 \tag{FEX12}
\]
Since $|\mu_{v+n}|\le A_*^{\rm abs}(B_*^{\rm sh})^{v+n}$ and
$v!/(v+n)!\le1/n!$, the triangle inequality in FEX10--11 gives
\[
 \|B_q(\partial_S)\|_{s,c_0;\mathcal P_D}\le P_D.
 \tag{FEX13}
\]
This bound holds for the literal fixed polynomial $B_q$ on any explicitly
specified finite degree space. On $D\le q-1$ it is the original AIE
operator, whose matrix is triangular with diagonal one. FEX8 therefore
proves its new exterior estimate
\[
\boxed{\quad
 P_D^{-2t}\le
 \frac{\det(X^*(B_q^{-1})^*B_q^{-1}X)}{\det(X^*X)}
 \le \min\{N_{D,s}^{2t},P_D^{2(D+1-t)}\}.
 \quad}\tag{FEX14}
\]
Here $X$ is any original $t$-plane in $\mathcal P_D$ and $N_{D,s}$ is the
unchanged AIE10 inverse bound. The new bound requires neither a new
zero-free disk nor a lower bound on an unevaluated restricted minor.
For $D=q-1$, $t=r$, its upper logarithm is at most
\[
 2(q-r)\{\log C_v+B_*^{\rm sh}(H_{q-1}+2)\}
       =O_{\mathsf H,\delta,\gamma}(q\log q)=o(kq).
 \tag{FEX15}
\]
Thus the factor $rD$ in the old exterior allowance is unnecessary. The
forward norm, together with the full Jacobian, controls the product;
raising the worst inverse singular value to the $r$th power loses that
information.

There is a direct high-source extension of this statement, distinguished
from changing AIE's cutoff. For a fixed $k$ use the same polynomial
$B_q(\partial_S)$ on each $\mathcal P_N$. Map the exact flag
\[
 \chi\mathcal P_{N-q}\ \subset\ I_K\mathbb C^r+\chi\mathcal P_{N-q}
\]
to its literal image under $B_q$. Its relation space is
$B_q(\partial_S)(\chi\mathcal P_{N-q})$, not the original ideal by
assumption. Its whole Jacobian is one, and FEX13 is valid on this newly
specified domain. If $\widetilde F_K$ is the signed kernel determinant
using the native source norm on these transported flags, FEX6 gives
\[
 -2(2q+1)\log P_{2q}\le\widetilde F_K-F_K
              \le2(4q+1)\log P_{2q}.
 \tag{FEX16}
\]
Indeed at endpoint $N$ its logarithmic Gram change belongs to
$[-2(N+1-r)\log P_N,2r\log P_N]$. Insert the four original signs and
use $P_N\le P_{2q}$. This is an exact extension and a signed estimate;
no commutation with the original remainder map has been assumed.

\section{The literal conductor at every source endpoint}
\label{sec:FEX-conductor}
The previous high-source extension is useful, but the actual conductor
provides a more informative reduction. On $\mathcal P_N$, $N\ge q-1$,
its exact Taylor identity is
\[
\begin{gathered}
 \mathcal C_A=\partial_S^v B^{[N]}(\partial_S),\qquad
 B^{[N]}(X)=\sum_{n=0}^{N-v}
       \frac{v!\mu_{v+n}}{(v+n)!\mu_v}X^n,\\
 B^{[q-1]}=B_q,\qquad
 L:=\ker(\mathcal C_A|_{\mathcal P_N})=\mathcal P_{v-1},\quad
 \mathcal C_A:\mathcal P_N/L\xrightarrow{\sim}\mathcal P_{N-v}.
\end{gathered}\tag{FEX17}
\]
For $v=0$, $L=0$. These assertions follow either from the leading
nonzero derivative or from the commuting invertible triangular factor
$B^{[N]}$. Thus the operator on the high endpoints is defined by the
original translations themselves. It is not obtained by silently
changing the cutoff in the definition of $B_q$.

The domain has its original quotient Gamma norm at centre $c$; the
codomain has its native Gamma norm at centre $c'$. Define
\[
 M_N=\max\left\{1,\frac{v!}{|\mu_v|}
       \sum_\nu|a_\nu|e^{|b_\nu-4|(H_N+2)}\right\},\qquad m_N=\log M_N.
 \tag{FEX18}
\]
The difference of centres is essential: $S'+b_\nu=c+iy+(b_\nu-4)$.
FEX11 proves $\|\mathcal C_Ap\|_{s,c'}\le M_N\|p\|_{s,c}$.
Taking the infimum over the same $L$-coset gives the same bound for the
induced isomorphism in FEX17.

Its squared whole-space Jacobian is explicitly
\[
 \boxed{\quad
 \mathscr J_{N,s}
  =\prod_{d=v}^{N}\frac{w_d}{w_{d-v}}
  =\frac{\prod_{d=N-v+1}^{N}w_d}{\prod_{d=0}^{v-1}w_d}.
 \quad}\tag{FEX19}
\]
Products are one when $v=0$. To prove the formula, take the orthonormal
quotient basis represented by $\varphi_v((S-c)/i),\ldots,
\varphi_N((S-c)/i)$, with the normalization in FEX10, and the corresponding
native basis at $c'$. The triangular diagonal of the map is
$i^{-v}(d)_v\sqrt{h_{d-v}/h_d}$, where $(d)_v=d!/(d-v)!$ is a falling
factorial in this formula. Its squared modulus is $w_d/w_{d-v}$.
The source mass is present in both $h$ factors before cancellation;
translation of the centre affects only strictly lower coefficients.
Telescoping proves the second equality, including overlapping product
ranges. For the unnormalized $\mathcal T_A$ the squared Jacobian is
$|\mu_v/v!|^{2(N+1-v)}\mathscr J_{N,s}$.

Put $\zeta_{N,s}=\log\mathscr J_{N,s}$. The following finite bounds
will suffice:
\[
\begin{array}{ll}
 0\le\zeta_{N,1}\le v\log(2\sqrt N),&s=1,\\[2pt]
 0\le-\zeta_{N,s}\le v(b-1)H_N,&s=k\ge9.
\end{array}\tag{FEX20}
\]
For the first use monotonicity of $w_n$ and $1\le w_n\le2\sqrt n$.
For the second, $w_n$ decreases and
$-\log w_n=\sum_{j=1}^n\log(1+(b-1)/j)\le(b-1)H_n$; then use the second
product formula in FEX19. These proofs also show directly that the
Jacobian error at the tensor source is $O(vk\log q)$, rather than an
uncontrolled determinant of dimension $q$.

The old explicit AIE inverse is still compatible with this estimate.
For example, let $R_A$ be that right inverse for the unnormalized
$\mathcal T_A:\mathcal P_D\to\mathcal P_{D-v}$, and let $R_{\min}$ be
its actual minimum-norm right inverse. Their difference takes values in
$L$ and is orthogonal to $R_{\min}$. On a $t$-plane $X$ in the target,
\[
 \det((R_AX)^*R_AX)
 \le \det((R_{\min}X)^*R_{\min}X)
       (1+M_{\rm raw}^2 C_{A,D,s}^2)^{\min(t,v)}.
 \tag{FEX21}
\]
Here $\|\mathcal T_A\|\le M_{\rm raw}$ and $\|R_A\|\le C_{A,D,s}$
are the original finite forward and AIE22 bounds. Indeed the excess
Gram has rank at most $v$, and $R_{\min}^*R_{\min}\succeq
M_{\rm raw}^{-2}I$. The operator $R_A-R_{\min}$ is an orthogonal
projection of $R_A$, so its norm is at most $C_{A,D,s}$. Apply the
matrix determinant lemma, or congruence with the minimum Gram, to obtain
FEX21. FEX6 bounds the minimum inverse determinant by
$J_{\rm raw}^{-1}M_{\rm raw}^{2(D+1-v-t)}\det(X^*X)$.
Only the fixed rank $v$, not $t$, multiplies the exponential allowance
for the raw low-degree correction. The main reduction below avoids that
raw correction by retaining its exact two quotient metrics instead.

\section{The residual observation and the smaller coefficient space}
\label{sec:FEX-residual}
AKJ12--25 identifies the original conductor preimage as
\[
 V_A=\{p\in\mathcal P_{q-1}:\psi\mid\mathcal T_Ap\}.
 \tag{FEX22}
\]
Its complete degree-adapted basis is $1,S,\ldots,S^{v-1}$ followed by
\[
 P_h=\frac{(q'+h+v)!}{(q'+h)!}
            B_q(\partial_S)^{-1}\mathcal J_v(\psi(S)S^h),
 \quad 0\le h<a-v,
 \qquad \mathcal J_vS^d=\frac{d!}{(d+v)!}S^{d+v}.
 \tag{FEX23}
\]
The variable in this inverse coefficient expression is identified with
$S'$ after applying the translation map. Each $P_h$ is monic of degree
$q'+v+h$. Its normalized conductor image is exactly
\[
 \mathcal C_AP_h=d_h\psi(S')S'^h,
 \qquad d_h=\frac{(q'+h+v)!}{(q'+h)!}.
 \tag{FEX24}
\]
The full original residual matrix in this basis is
$\Xi=(\Xi_{\rm lo},\Xi_{\rm hi})$, of size $j_k$ by $a$ and rank $j_k$.
Its kernel maps isomorphically onto the original $K$. These are the
supplied actual coefficient maps, not independent matrix parameters.
Put
\[
 r_{\rm lo}=\operatorname{rank}\Xi_{\rm lo},\quad
 t_0=v-r_{\rm lo}=\dim(K\cap L),\quad t=r-t_0.
 \tag{FEX25}
\]
Choose a quotient coordinate map $Q_{\rm lo}$ with kernel exactly
$\operatorname{im}\Xi_{\rm lo}$, and a full-column-rank basis $W_k$ of
\[
 \ker(Q_{\rm lo}\Xi_{\rm hi})\subset\mathbb C^{a-v}.
 \tag{FEX26}
\]
The residual map has rank $j_k-r_{\rm lo}$ because the full matrix has
rank $j_k$. Hence the space in FEX26 has dimension $t$. These choices
are made using actual nonzero minors and then kept fixed as the source
endpoint changes. No generic minor or generic rank is substituted.
Define the polynomial column matrix
\[
 \boxed{\qquad Z_k=\operatorname{diag}(d_h)W_k
             \quad\hbox{in the frame }(1,S',\ldots,S'^{a-v-1}).\qquad}
 \tag{FEX27}
\]
Then $\mathcal C_AK=\psi\operatorname{im}Z_k$. In fact, for a high
coordinate $z$ in FEX26 the equation
$\Xi_{\rm lo}\ell+\Xi_{\rm hi}z=0$ has a solution. Its different
solutions differ by $K\cap L$; FEX24 kills that ambiguity and gives
precisely $\psi Z_k$. Thus this is the actual image, including the
original unit and period dependence in $\Xi$ and the conductor
coefficients. The inverse finite-symbol factor no longer appears in the
columns $Z_k$.

For every $j\ge0$ define the exact transformed relation polynomial
\[
 \boxed{\quad
 R_j^A(S')=\frac{\mathcal C_A(\chi S^j)(S')}{\psi(S')}
 =\frac{v!}{\mu_v}\sum_\nu a_\nu
       \frac{\chi(S'+b_\nu)}{\psi(S')}(S'+b_\nu)^j.
 \quad}\tag{FEX28}
\]
Every displayed quotient is a polynomial. Indeed, if $S'$ is the
$(a',b')$ inner-grid root, then $S'+b_{r_0s_0}$ is the original root
with indices $(a'+r_0,b'+s_0)$. All these indices are between $0$ and $k$.
The degree of $R_j^A$ is $a-v+j$ and its leading coefficient is
$(q+j)_v$; this follows from the first nonzero derivative in FEX17.
It is nonzero, and these distinct degrees prove linear independence of
the relation columns. In particular the high-source terms have not
been replaced by an unjustified continuation of the low-degree symbol.

At endpoint $N$ put $m_N^{\rm pol}=N-v-q'$ and $d_N^{\rm rel}=N+1-q$.
Form the complete polynomial Gram $H_{N,s}^{\psi}$ of
$(1,S',\ldots,S'^{m_N^{\rm pol}})$ in the measure
$|\psi(c'+iy)|^2dm_{0,s}(y)$ at $S'=c'+iy$. Zero-extend $Z_k$ to this
frame and let $R_N^A$ have columns $R_0^A,\ldots,R_{d_N^{\rm rel}-1}^A$.
Define
\[
\begin{split}
 \widehat H_{K,N}^{(s)}={}&Z_k^*H_{N,s}^{\psi}Z_k\\
 &-Z_k^*H_{N,s}^{\psi}R_N^A
    ((R_N^A)^*H_{N,s}^{\psi}R_N^A)^{-1}
                (R_N^A)^*H_{N,s}^{\psi}Z_k,\\
 \widehat F_k^{(s)}={}&
 \sum_{i=0}^3\epsilon_i\log\det\widehat H_{K,N_i}^{(s)},\\
 & (N_0,N_1,N_2,N_3)=(q-1,q,2q-1,2q),\\
 & (\epsilon_0,\epsilon_1,\epsilon_2,\epsilon_3)=(1,1,-1,-1).
\end{split}\tag{FEX29}
\]
The empty relation inverse is absent at $N=q-1$. The four relation
ranks are still exactly $0,1,q,q+1$. They have not been truncated to the
smaller kernel dimension. Positivity follows because a nonzero linear
combination of $Z_k$ has degree at most $a-v-1$, whereas every nonzero
combination of the $R_j^A$ has degree at least $a-v$. The complete Gram
is positive in the original positive Gamma density. Minimization over
an increasing relation span proves
\[
 \widehat H_{K,q-1}^{(s)}\succeq\widehat H_{K,q}^{(s)}
 \succeq\cdots,\qquad \widehat F_k^{(s)}\ge0.
 \tag{FEX30}
\]
A change of the fixed basis $W_k$ gives the same four-endpoint value;
its four identical Gram factors cancel with the original signs.

\section{A finite error for the actual compressed determinant}
\label{sec:FEX-error}
Let $K^+=K+L\subset E$. This is only an auxiliary space for an exact
metric calculation, not an enlargement of the observation kernel.
All quotient and restriction metrics below are induced from the same
original $G_N^{(s)}$ and their indicated fixed coefficient frames.
Let $F_{K^+/L}$ and $F_{K^+/K}$ denote their four-endpoint volumes.
The two block Gram factorizations of $K^+$ give
\[
 \boxed{\qquad F_K=F_{K^+/L}+F_L-F_{K^+/K}.\qquad}
 \tag{FEX31}
\]
Before summing there are fixed basis Jacobians; they cancel only after
all four signs. Both factorizations use the actual Schur complements.
The two subquotient dimensions are $\dim L=v$ and
$\dim(K^+/K)=v-t_0$. The existing full AKS covariance inequality passes
by restriction and minimization to every fixed subquotient. Thus, with
$\mathcal L_s=\mathcal L_{s,q}+\mathcal L_{s,q+1}$ from AKS,
\[
 0\le F_L\le v\mathcal L_s,\qquad
 0\le F_{K^+/K}\le(v-t_0)\mathcal L_s.
 \tag{FEX32}
\]
No angle condition between $K$ and $L$ is used.

For a fully explicit existing allowance one may use the AKS Newton
bound in place of a sharper available $\mathcal L_s$. Write
\[
\begin{gathered}
 R=k\sqrt{\delta^2+\gamma^2},\quad
 t_q=\pi/2-1/q,\quad a_q=2\log3/t_q,\quad
 s_q=(a_q^q-1)/(a_q-1),\\
 b_{s,q}=[(9/25)\sin(1/q)]^{-s/2}
          \exp(2R(\log3+t_q/2)),\\
 \mathcal L_{s,j}=\log\left(1+b_{s,q}s_q^2
                  \sum_{d=q}^{q+j-1}\frac{d!}{(s/2)_d}(25/16)^d\right).
\end{gathered}\tag{FEX33}
\]
These are the original complete coefficient bounds with both masses
retained before norm-ratio cancellation. They satisfy
$\mathcal L_s=O_{\delta,\gamma}(q)$ for $s=1,k$ along this family.
Indeed $\log s_q=O(q)$, the last geometric-weighted sum has logarithm
$O(q+\log q)$ because $d!/(s/2)_d\le2\sqrt d$, and
$\log b_{s,q}=O(k\log q+R)=o(q)+O(q)$.

Apply FEX6 to the source $\mathcal P_N/L$ and its conductor image.
The relation subspace is $(\chi\mathcal P_{N-q}+L)/L$. The larger
flag term is generated also by $K$, giving quotient dimension $t$.
By FEX28, its image is the span of $\psi R_N^A$ and $\psi Z_k$.
Multiplication by $\psi$ is an exact isometry from the indicated
weighted coefficient space onto this polynomial image. The source
quotient metric is exactly the metric of $K^+/L$, since successive
minimization first by $\chi\mathcal P_{N-q}$ and then by $L$ is
minimization over their sum. Thus, in the fixed transported frames, put
\[
\begin{gathered}
 e_N=\log\det\widehat H_{K,N}^{(s)}-
                    \log\det H_{K^+/L,N}^{(s)},\qquad n_N=N+1-v,\\
 e^-_{N,s}=\zeta_{N,s}-2(n_N-t)m_N,\quad e^+_{N,s}=2tm_N,
 \quad e^-_{N,s}\le e_N\le e^+_{N,s},\\
 D_s^-=e^-_{q-1,s}+e^-_{q,s}-e^+_{2q-1,s}-e^+_{2q,s},\qquad
 D_s^+=e^+_{q-1,s}+e^+_{q,s}-e^-_{2q-1,s}-e^-_{2q,s}.
\end{gathered}\tag{FEX34}
\]
The signs on $e_{N,s}^{\pm}$ denote new finite scalar bounds; they do not
redefine an original source scalar or covariance. All endpoint quantities are computed from FEX18--20 and the
actual finite ranks in FEX25. Original frame factors have already been
accounted for by transported quotient frames. Combining the four
inequalities and FEX31--32 proves the central finite result:
\[
 \boxed{\quad
 \widehat F_k^{(s)}-D_s^+-(v-t_0)\mathcal L_s
 \ \le F_K^{(s)}\le\
 \widehat F_k^{(s)}-D_s^-+v\mathcal L_s.
 \quad}\tag{FEX35}
\]
This is a bound for the actual kernel determinant, not for a different
ambient determinant. Every residual coefficient and relation cross
term remains in FEX29.

For a simpler explicit error let
\[
 J_{\max,s}=\max_{i=0}^3|\zeta_{N_i,s}|,\quad
 m_{\max}=m_{2q},\quad
 \mathcal E_{k,s}=v\mathcal L_s+2J_{\max,s}+2(4q+1)m_{\max}.
 \tag{FEX36}
\]
Then
\[
 \boxed{|F_K^{(s)}-\widehat F_k^{(s)}|\le\mathcal E_{k,s}
              =O_{\mathsf H,\delta,\gamma}(q\log q)=o(kq).}
 \tag{FEX37}
\]
To check the finite factors, the lower signed bound for $\sum\epsilon_i e_{N_i}$
is at least $-2J_{\max,s}-2(2q+1-2v)m_{\max}$, and the upper bound is at
most $2J_{\max,s}+2(4q+1-2v)m_{\max}$. Enlarge to FEX36 and use
FEX32. Also
\[
 m_{\max}\le
 \log\max\{1,v!A_*^{\rm abs}/|\mu_v|\}
 +8\sqrt{\delta^2+\gamma^2}(H_{2q}+2).
 \tag{FEX38}
\]
Together with FEX20 and FEX33 this proves the order claim. All constants
are fixed at the actual fixed period. Dividing by $kq$ uses
$\log q/k\to0$. A uniform claim over periods approaching a locus where
$\mu_v$ vanishes is not made: the displayed inverse coefficient must
remain in the finite bound. Uniformity holds on a fixed coefficient
stratum with this coefficient ratio bounded and the same $v$.

It follows in particular that $F_K^{(s)}/(kq)$ and
$\widehat F_k^{(s)}/(kq)$ have exactly the same limit points. This does
\emph{not} assert that either sequence has a unique limit. The old
uniform ceiling for $F_K$ remains true, while FEX35--37 constitute a
new, fixed-period transfer at strictly smaller error scale.

\section{The original arithmetic and signed receiving identities}
\label{sec:FEX-receivers}
Let $[L_s^{\rm old},U_s^{\rm old}]$ denote the current MRI4e interval,
including its first compressed increment and the minimum of its existing
finite upper bounds. The following is an add-only improvement at the
same four endpoints:
\[
\begin{split}
 L_s^{\rm new}&=\max\{L_s^{\rm old},0,
            \widehat F_k^{(s)}-D_s^+-(v-t_0)\mathcal L_s\},\\
 U_s^{\rm new}&=\min\{U_s^{\rm old},
            \widehat F_k^{(s)}-D_s^-+v\mathcal L_s\},\\
 L_s^{\rm new}&\le F_K^{(s)}\le U_s^{\rm new},\qquad s=1,k.
\end{split}\tag{FEX39}
\]
The intersection need not improve every finite instance. The comparison
of scales, however, is unconditional at each fixed actual period in
the stated domain.

The full original Gamma multiplier is
\[
 f_l(S)=\prod_{j=0}^{l-1}(S-c-2j-\tfrac12),\qquad
 dm_{0,k}=\beta_k|f_l(c+iy)|^2d\sigma,
 \quad \beta_k=\frac{(2\pi)^{k/2}}{\sqrt2\Gamma(k/2)}.
 \tag{FEX40}
\]
Its transported full-jet kernel, conditional covariance and mixed blocks
are AMT23--30, unchanged. The existing difference interval is
$-r\Omega_k^{\rm low}\le F_K^{(1)}-F_K^{(k)}\le r\Omega_k^{\rm high}$,
where
\[
\begin{gathered}
 \Pi_l=\prod_{j=0}^{l-1}(2j+\tfrac12)^2,\quad
 U_{l,N}^{\rm poly}=\prod_{j=0}^{l-1}
                  [4(N+j+1)^2+(2j+\tfrac12)^2],\\
 \Omega_k^{\rm low}=\log\frac{U_{l,q-1}^{\rm poly}U_{l,q}^{\rm poly}}{\Pi_l^2},
 \qquad
 \Omega_k^{\rm high}=\log\frac{U_{l,2q-1}^{\rm poly}U_{l,2q}^{\rm poly}}{\Pi_l^2}.
\end{gathered}\tag{FEX41}
\]
Consequently one may further intersect the tensor-source interval with
$[L_1^{\rm new}-r\Omega_k^{\rm high},
 U_1^{\rm new}+r\Omega_k^{\rm low}]$. The fixed-source interval has the
corresponding reverse intersection. No independent error assignment is
substituted for their common total.

For the arithmetic measure retain exactly
\[
 w_h(y)=|v_h(\tfrac12+iy)|^2/(2\pi),\qquad
 dm_{h,k}(y)=w_h^{*k}(y)\,dy,\qquad
 \delta_k^\sigma=\mathcal B_k^{\rm ar}-\mathcal B_k^\sigma.
 \tag{FEX42}
\]
Let $H_N^{\rm ar},H_N^\sigma$ be the full coefficient-source Grams.
AMT6 retains $(a_k/D_N)H_N^\sigma\preceq H_N^{\rm ar}\preceq A_kH_N^\sigma$.
For $m=1$ its original constants are $p=7/2$, $B_h=50$, $M=25$,
$\alpha=\pi/2$, $D_N=[1+4(N+M)^2]^M$, and
\[
 a_k=\frac{c_h\vartheta_h^{k-3}e^{-\alpha(k-3)}(k-2)^{-B_h}}
                  {C_\Gamma2^{B_h/2}},\qquad
 A_k=\frac{C_h^{\rm tilt}}{c_\Gamma}k^{p+1}M_\alpha^{k-1}.
 \tag{FEX43}
\]
Here $\vartheta_h=\int_{-1}^1w_h$ and $M_\alpha=\int e^{\alpha y}w_h(y)dy$
are the original integrals, and the named positive envelope constants
are the unchanged TW/BT constants. In particular none is inferred from
a diagnostic packet.

Keep the AMT exact deficits
\[
\begin{gathered}
 L_N^{\rm ar}=\log(A_kD_N/a_k),\quad
 d_{E,N}=q\log A_k-\log\frac{\det G_N^{\rm ar}}{\det G_N^\sigma},\\
 c_{X,N}=\min\{r_XL_N^{\rm ar},d_{E,N}\},\quad
 C_X^-=c_{X,q-1}+c_{X,q},\quad C_X^+=c_{X,2q-1}+c_{X,2q},\quad X=K,B.
\end{gathered}\tag{FEX44}
\]
The source-evaluated caps replacing $d_{E,N}$ by $\Xi_{k,N}$ remain
available. The complete nested positive deficits and their exact common
sum are those of AMT13--18. Using either version of the caps, FEX39
therefore gives precisely the new MRI4f interval
\[
\begin{split}
 L_{\rm ar}^{\rm new}&=\max\{0,L_1^{\rm new}-C_K^-,
                         L_1^{\rm new}+\delta_k^\sigma-C_B^+\},\\
 U_{\rm ar}^{\rm new}&=\min\{\mathcal B_k^{\rm ar},U_1^{\rm new}+C_K^+,
                         U_1^{\rm new}+\delta_k^\sigma+C_B^-\},\\
 L_{\rm ar}^{\rm new}&\le F_K^{\rm ar}\le U_{\rm ar}^{\rm new}.
\end{split}\tag{FEX45}
\]
Subtracting the appropriate kernel endpoints from the unchanged total
$\mathcal B_k^{(s)}$ or $\mathcal B_k^{\rm ar}$ gives the exact boundary
interval. This keeps correlation with the actual signed arithmetic
value. Its QRI/QGQ finite intervals, including the simple-multiplicity
nonzero radius, are not altered by this calculation.

The BRI6d identity remains, with its original orientation,
\[
\begin{gathered}
 F_K^{\rm ar}+F_B^{\rm ar}-F_B^{(k)}
       =F_K^{(k)}+\mathcal H_k+\delta_k^\sigma,\qquad
 \mathcal H_k=\mathcal B_k^\sigma-\mathcal B_k^0=-\mathfrak T_k,\\
 L_k^{\rm new}+\mathcal H_k+\delta_k^\sigma
 \le F_K^{\rm ar}+F_B^{\rm ar}-F_B^{(k)}
 \le U_k^{\rm new}+\mathcal H_k+\delta_k^\sigma.
\end{gathered}\tag{FEX46}
\]
The already supplied source differences are $o(kq)$, and the complete
Gamma return has $\mathcal H_k/(kq)\to-\mathcal C_\Gamma/4$ with the
unchanged positive equilibrium constant. Consequently the new reduction
has the exact leading-scale receiving form
\[
 \boxed{\quad
 \frac{F_K^{\rm ar}+F_B^{\rm ar}-F_B^{(k)}}{kq}
  =\frac{\widehat F_k^{(1)}}{kq}-\frac{\mathcal C_\Gamma}{4}+o(1).
 \quad}\tag{FEX47}
\]
In particular $F_K^{\rm ar}/(kq)$, $F_K^{(1)}/(kq)$,
$F_K^{(k)}/(kq)$, and both new $\widehat F$ ratios have the same limit
points. The leading value of FEX29 is still needed to decide the sign
in FEX47. Neither the old ceiling nor the new $o(kq)$ transfer assigns
that value.

\section{Actions, representatives, and the retained mixed-support maps}
\label{sec:FEX-maps}
The conductor is a coefficient calculation, not a replacement for the
arithmetic action or the original observation. Its precise action defect
on polynomials is
\[
 \mathcal C_A(Sp)-S'\mathcal C_Ap
       =\frac{v!}{\mu_v}\sum_\nu a_\nu b_\nu p(S'+b_\nu).
 \tag{FEX48}
\]
For the representative of degree less than $q$, write
$\tau(p)=[S^{q-1}]p$. The original action is
$A_\chi[p]=[Sp-\chi\tau(p)]$. Its transported polynomial defect is the
right side of FEX48 minus $\mathcal C_A(\chi\tau(p))$. This formula is
in the full polynomial target; divisibility by $\psi$ is not asserted
for an action image outside the actual conductor preimage. In
particular $K$ is not assumed invariant, and the original OKA9
observation-action defect is retained.

Nor is the auxiliary $L$ an invariant subspace of that action when
$v>0$. On its top monomial $S^{v-1}$, multiplication followed by
$\mathcal C_A$ gives the nonzero constant $v!$, whereas all its lower
monomials give zero. This is the explicit rank-one defect. No quotient
Euler action on $K^+/L$ was needed in the metric proof.

One can retain all coefficients while using the calculation. Let
$R_{C,N}$ be any of the explicitly specified finite right inverses of
$\mathcal C_A$ on the full space $\mathcal P_{N-v}$, for example
$(B^{[N]}(\partial_S))^{-1}\mathcal J_v$. Then
\[
 p\longmapsto(p-R_{C,N}\mathcal C_Ap,\mathcal C_Ap)
       \in L\oplus\mathcal P_{N-v},\qquad
 (\ell,g)\longmapsto\ell+R_{C,N}g
 \tag{FEX49}
\]
are inverse coefficient maps. This explicitly retains the low
component rather than treating it as zero. The original remainder map,
full unit, and literal period observation may be composed with the
inverse in FEX49; they are then the same original maps on the same
class, not newly assigned maps on a smaller fibre. The volume identity
FEX31 restores both low-component quotient metrics exactly.

If an original minimum section is changed by $\chi Q$, its original
cochain primitive is still the AKP primitive. Namely divide
$\chi(s_1+\cdots+s_k)$ successively by the original $h(s_j)$ to obtain
$\sum_jh(s_j)Q_j(\mathbf s)$. Then use the original tensor differential
and the primitive
\[
 \sum_{j=1}^k(-1)^{j-1}Q_j(\mathbf D)Q(D_1+\cdots+D_k)
 \bigl(F_h^{\otimes(j-1)}\otimes\phi_*\otimes
                     F_h^{\otimes(k-j)}\bigr).
 \tag{FEX50}
\]
Its differential is the literal difference of the original theta
representatives, because $h(D)F_h=\Theta\phi_*$. Minus-only legs use the
unchanged lift $(0,-\mathcal F_{\rm Fourier}\phi)$ and transported
$1-D$ action. This statement applies to an actual relation difference;
no element of $K$ or $L$ is newly declared a theta boundary.
All support-face insertions and signed coefficient permutations commute
with the reconstructed linear maps componentwise. A supported zero
stays at its original label and is not identified with the external
$\tau$. The additional proper-source Taylor kernels from the mixed-primary
lane remain, not being replaced by a conductor kernel.

The linear-algebra and Gamma-polynomial lemmas apply at any finite
multiplicity with their actual dimensions. The explicit OCF conductor,
its four additional orbit factors, and the rank $8k-16$ used for
FEX22--47 belong only to the stated simple-quartet case. No such rank or
frame is assigned to the full-multiplicity OMG case. Its nilpotents,
row-residue annihilator, and observed-iterate kernel remain the supplied
ones. This continuation makes no Frobenius identification or arithmetic
purity assertion.

\section{Scope of verification and the next exact quantity}
\label{sec:FEX-status}
The analytic assertions proved here are FEX6--21 and the exact transfer
FEX22--47 from the explicitly named supplied inputs. Finite exact
fixtures test the flag identity, its inverse determinant bounds, the
source Jacobian, and the full four-endpoint reduction with nonzero cross
terms. They do not evaluate actual xi zeros or actual period coefficients.
The accompanying execution record distinguishes those computations
from the written proofs and from inherited formal certificates.

The next unresolved quantity is the fixed-period determinant FEX29,
with the actual $Z_k$, $R_j^A$, and weighted original Gamma Gram. Its
leading value is \emph{not} established by this paper. It is now separated
from the finite-symbol inverse with an explicit error smaller than $kq$.
In particular, repeating an operator inverse bound of exponential size
and raising it to $r$ is no longer the remaining task. The complete
original high relation ranks, their mixed Gram blocks, and the period
position of $\operatorname{im}Z_k$ are the data that still need estimation.
No RH contradiction or disproof follows from the result proved here.

\subsection*{Source dependencies}
The fixed objects and current ranks are the supplied OCS/OCF/OKM/OKA,
AKJ12--25, and RPC/AMT statements in \texttt{14\_CURRENT\_KERNEL\_AND\_MULTIPLICITY\_PROOFS.tex}.
The retained inverse is AIE1--29; the original covariance bound is
AKS20--32. The signed source transfer is AMT13--49 and the unchanged
GEL/QGT/QRI/QGQ provider in \texttt{02\_FULL\_GAMMA\_PROOFS(1).tex}.
All complete current source files and their exact hashes accompany this
continuation. PR32's residue and finite observation certificates are
inherited inputs only; no new Lean execution is claimed here.

% END CURRENT FEX

\clearpage
% BEGIN CURRENT FGC
% Exact comparison for September 15 intake. No current or sealed source edited.
\section{The fixed-column conductor quotient and the full GIP graph}
\label{sec:FGC-intake}

Fix exactly the original simple-quartet and fixed-period data of FEX1--4,
CEP1--6 and GIP1--3. In particular $k=4l+1\ge9$,
$q=(k+1)^2$, $q'=(k-7)^2$, $\Delta=q-q'$, $0\le v\le80$,
$\chi=\chi_k$, $\psi=\chi_{k-8}$, $c=k/2$, $c'=c-4$,
and $s\in\{1,k\}$. Retain the exact nonzero scalar
\[
 \alpha_A=\frac{v!}{\mu_v},\qquad
 \mathcal C_A=\alpha_A\mathcal T_A.
 \tag{FGC1}
\]
Every polynomial norm is the original Gamma norm at its stated centre,
with the full mass $(2\pi)^{s/2}$. The scalar in FGC1 multiplies the
specified coefficient map and does not alter the measure.

For each $N\in\{q-1,q,2q-1,2q\}$ put
$t_N=N+1-q$ and $b_N=N-v-q'$. Use the exact GIP coefficient space,
now with a different name to keep $t_N$ distinct from a kernel dimension:
\[
 \mathscr X_N=\mathbb C^v\oplus\mathbb C[S']_{\le b_N},\qquad
 L_Nh=P_{<v}Z_{\chi,N}h,\qquad
 U_N^Ah=\frac{\mathcal T_A(\chi h)}{\psi}
 \quad(h\in\mathbb C[S]_{\le N-q}).
 \tag{FGC2}
\]
Here $P_{<v}$ gives the first $v$ original upper orthonormal coordinates,
$Z_{\chi,N}h$ is the complete upper coordinate column of $\chi h$,
and $U_N^A$ is exactly GIP3, equivalently CEP41--42. Thus the relation
subspace of GIP is $\{(L_Nh,U_N^Ah):\deg h\le N-q\}$.
The top coordinate is retained in every formula below.

Define the explicit coefficient isomorphism and the transformed graph
\[
 D_{\alpha,N}(\ell,z)=(\ell,\alpha_Az),\qquad
 D_{\alpha,N}^{-1}(\ell,z)=(\ell,\alpha_A^{-1}z),\qquad
 \mathscr G_N^\alpha=
 \{(L_Nh,\alpha_AU_N^Ah):\deg h\le N-q\}.
 \tag{FGC3}
\]
The complete GIP graph quotient is carried by $D_{\alpha,N}$ to
$\mathscr X_N/\mathscr G_N^\alpha$. Its map from the original
conductor kernel $V_A$ sends the class of a polynomial $p$ to
\[
 [p]\longmapsto
 \left[\left(P_{<v}p,\frac{\mathcal C_Ap}{\psi}\right)\right].
 \tag{FGC4}
\]
The second coordinate exists exactly for $[p]\in V_A$, by CEP5--6.
Changing $p$ by $\chi h$ adds the displayed member of
$\mathscr G_N^\alpha$ and therefore does not change its class.
The inverse is the GIP inverse $E_N^{\rm graph}$ composed with
$D_{\alpha,N}^{-1}$, so FGC4 is an isomorphism. This proves the
relation to the full earlier graph at each of its actual cutoffs.

\subsection{The precise projection and its complete kernel}

Define the lower coefficient quotient
\[
 \mathscr Y_N=
 \mathbb C[S']_{\le b_N}/\operatorname{im}(\alpha_AU_N^A),
 \qquad
 \pi_N:\mathscr X_N/\mathscr G_N^\alpha\longrightarrow\mathscr Y_N,
 \quad[(\ell,z)]\longmapsto[z].
 \tag{FGC5}
\]
The map is well defined because its relation images are exactly
$\operatorname{im}(\alpha_AU_N^A)$, and it is onto because every
lower column has the lift $(0,z)$. If its value is zero, write
$z=\alpha_AU_N^Ah$ and subtract the full graph vector
$(L_Nh,\alpha_AU_N^Ah)$. The resulting representative is
$(\ell-L_Nh,0)$. A vector $(\ell,0)$ belongs to the graph only
if $U_N^Ah=0$. This implies $h=0$: for nonzero $h$, the degree
of $\mathcal T_A(\chi h)$ is $q+\deg h-v$, with nonzero leading
coefficient $\operatorname{lc}(h)\mu_v\binom{q+\deg h}{v}$, by the first nonzero
moment identity. Therefore $\ell=0$. We have proved the full
exact sequence, including the original low coordinate and every
relation component,
\[
 0\longrightarrow\mathbb C^v
  \xrightarrow{\ell\mapsto[(\ell,0)]}
 \mathscr X_N/\mathscr G_N^\alpha
  \xrightarrow{\pi_N}\mathscr Y_N\longrightarrow0.
 \tag{FGC6}
\]
Under FGC4 the first term is exactly the original subspace
$L=\mathbb C[S]_{<v}\subset E_k$, because the conductor kills
exactly those polynomials and its first coordinate is their original
orthonormal coordinate. Thus FGC6 is, on the same original primary
space, the quotient sequence
\[
 0\longrightarrow L\longrightarrow V_A
       \longrightarrow V_A/L\longrightarrow0,
 \qquad V_A/L\xrightarrow{\ \simeq\ }\mathscr Y_N.
 \tag{FGC7}
\]
In particular its target has dimension $\Delta-v$ at each original
endpoint. The inverse to the last map is completely specified: choose
a representative $z$, solve $\mathcal C_Ap=\psi z$ by the finite
right inverse of FEX17, and take $[p]$ modulo $L$ and the source
relations. A different representative changes $p$ by an original
$\chi h$ plus an element of $L$, so the same target class results.

Now retain $K=\ker\Lambda_k\subset V_A$, put
$t_0=\dim(K\cap L)$, and set $K^+=K+L$. Restricting FGC7 gives
\[
 0\longrightarrow K\cap L\longrightarrow K
       \longrightarrow\pi_N(K)\longrightarrow0,
 \qquad \pi_N(K)\simeq K^+/L,
 \quad \dim\pi_N(K)=r-t_0.
 \tag{FGC8}
\]
The second isomorphism sends $k+L$ to the same image as $k$.
It is onto, and a zero image says $k\in L$, which proves injectivity
on $K^+/L$. No preservation of an Euler or multiplication action is
asserted by these coefficient quotient maps.

\subsection{The actual fixed columns and the Schur metric}

Use the original AKJ basis and residual matrix
$\Xi=(\Xi_{\rm lo},\Xi_{\rm hi})$ from FEX23--26. Let $W_k$
be its specified basis of $\ker(Q_{\rm lo}\Xi_{\rm hi})$ and
$d_h=(q'+h+v)!/(q'+h)!$. The first nonzero moment identity gives
$\mathcal C_AP_h=d_h\psi S'^h$ for every original high AKJ column.
The low AKJ columns are precisely $L$ and are killed by $\mathcal C_A$.
Thus the exact image in FGC8 has the fixed, endpoint-independent
representative matrix
\[
 Z_k=\operatorname{diag}(d_h)W_k,
 \qquad \pi_N(K)=
   (\operatorname{im}Z_k+\operatorname{im}(\alpha_AU_N^A))
                         /\operatorname{im}(\alpha_AU_N^A).
 \tag{FGC9}
\]
Solvability of $\Xi_{\rm lo}\ell+\Xi_{\rm hi}z=0$ is exactly
the equation $Q_{\rm lo}\Xi_{\rm hi}z=0$, proving both inclusions
in this identity. The graph's lower relation columns are
\[
 R_j^A=\alpha_AU_N^A(S^j)=
 \frac{\mathcal C_A(\chi S^j)}{\psi},\qquad0\le j<t_N.
 \tag{FGC10}
\]
This is FEX28, and also the lower block of the complete original
relation graph in FGC3. The nonzero leading coefficient is
$(q+j)_v$ and its degree is $\Delta-v+j$. The columns of $Z_k$
have degree at most $\Delta-v-1$. Consequently their quotient
classes are independent at every cutoff, including the two high
relation ranks $q$ and $q+1$.

Give the lower polynomial space its exact norm
$\int|z(c'+iy)|^2|\psi(c'+iy)|^2dm_{0,s}(y)$, whose full Gram
in the increasing monomials is $H_{N,s}^\psi$. Then the quotient
metric of FGC5 on FGC9 is exactly
\[
 \begin{split}
 \widehat H_{K,N}^{(s)}
 &=Z_k^*H_{N,s}^\psi Z_k
 -Z_k^*H_{N,s}^\psi R_N^A
   ((R_N^A)^*H_{N,s}^\psi R_N^A)^{-1}
                  (R_N^A)^*H_{N,s}^\psi Z_k.
 \end{split}
 \tag{FGC11}
\]
Indeed minimizing $(Z_kx+R_N^Ah)^*H_{N,s}^\psi(Z_kx+R_N^Ah)$
over $h$ gives the unique minimizer
$h=-((R_N^A)^*H_{N,s}^\psi R_N^A)^{-1}(R_N^A)^*H_{N,s}^\psi Z_kx$.
The relation Gram is positive by the independence in FGC10, and
substitution proves FGC11 including both mixed factors. At $t_N=0$
the minimization has no variable and the second term is absent.
This proves the FEX29 matrix directly from the full GIP relation
graph by the explicit quotient FGC5. It has not identified the
transformed relation span with a chosen polynomial ideal.

The original exact metric in these normalized graph coordinates is
$D_{\alpha,N}^{-*}H_N^{\rm actual}D_{\alpha,N}^{-1}$, with
$H_N^{\rm actual}$ from GIP5. It is generally different from the
native lower metric in FGC11. On the whole normalized graph space
the native direct-sum metric is $\operatorname{diag}(I_v,H_{N,s}^\psi)$.
Taking its quotient first by $\mathscr G_N^\alpha$ and then by the
low image is the same minimization as taking the quotient by their
sum. In that sum the low coordinate is unrestricted, so its minimum
is zero, and the second-coordinate minimization is precisely FGC11.
Thus the native-metric projection also matches the stated Schur matrix;
the comparison to the original metric still uses FEX34--35 or CEP's
full source-map estimates. The scalar $\alpha_A$ has remained in the
coordinate map and both metrics throughout.

\subsection{Both original low metric terms and their signed identity}

For any fixed positive original endpoint metric $G_N^{(s)}$ on $E_k$,
restrict it to $K^+$. In fixed coefficient frames the two complete
Schur factorizations give
\[
 \det H_{K^+,N}=J_1\det H_{L,N}\det H_{K^+/L,N}
              =J_2\det H_{K,N}\det H_{K^+/K,N},
 \tag{FGC12}
\]
where, if $A_1,A_2$ are the matrices of the two adapted frames in
the fixed reference frame of $K^+$, then
$J_1=|\det A_1|^{-2}$ and $J_2=|\det A_2|^{-2}$.
They are independent of the endpoint. To prove
each equality, complete the chosen subspace frame by fixed quotient
lifts, write the full positive Gram in that adapted frame, and
eliminate its off-diagonal block. The quotient block is exactly
the minimum Gram because solving its first-block linear equation
is the same minimization. Its determinant times the first block's
determinant is the determinant in the adapted frame. This proves
the formula including its fixed Jacobian.

Taking the four original signs $(1,1,-1,-1)$ cancels these two
frame constants, and gives exactly
\[
 F_K^{(s)}=F_{K^+/L}^{(s)}+F_L^{(s)}-F_{K^+/K}^{(s)},\qquad
 \dim L=v,\quad\dim(K^+/K)=v-t_0.
 \tag{FGC13}
\]
The original AKS endpoint covariance comparison passes to a restriction
by evaluating the quadratic forms on that subspace, and passes to a
quotient by taking the infimum on the same fibre. Its scalar two-sided
bounds therefore give FEX32 for these actual subquotients. In particular
both nonnegative low-volume terms in FGC13 remain; a projection to
FGC11 alone does not remove them. FEX34--37 additionally compares
$F_{K^+/L}$ with $\widehat F_k$ by the original conductor's complete
Jacobian, proving the stated $O_{\rm actual}(q\log q)$ error.

Equations FGC3--13 exhibit exactly how the new fixed-column presentation
fits the existing CEP/GIP construction. The existing RUI5--15 and
CEP29--39 already prove the $o(kq)$ source/exterior transport. FEX adds
an explicit realization on the lower fixed columns $Z_k$, with the
full original transformed relation span and the two original low
metric corrections restored. These equations do not turn that span
into the CTV outer ideal or remove the actual residual observation.

% END CURRENT FGC

\clearpage
% BEGIN CURRENT ECL
% Independent proof in the original EIQ/OFV normalization.
\section{The exact closed equilibrium coefficient}
\label{sec:closed-equilibrium-intake}

This section retains the potential, measure, and energy of EIQ1--33:
\begin{equation}\tag{ECL1}
\begin{gathered}
V_{\alpha,\beta}(x)=\beta\sqrt{x}-\alpha\log x,\qquad
\alpha>0,\quad\beta>0,\quad x>0,\\
\mu_{\alpha,\beta}(dx)=\rho_{\alpha,\beta}(x)\,dx,\qquad
F(\alpha,\beta)=
\iint\log|x-y|\,d\mu_{\alpha,\beta}(x)d\mu_{\alpha,\beta}(y)
-\int V_{\alpha,\beta}\,d\mu_{\alpha,\beta},\\
\mathcal L(\alpha,\beta)=\int\log x\,d\mu_{\alpha,\beta}(x).
\end{gathered}
\end{equation}
The preceding complete EIQ proof supplies the unique minimizing
probability measure, its positive density and exact support
$[a,b]=[u^2,v^2]$, and its Euler equality. In particular, no assertion
in this section changes the minimization problem or chooses another
measure. The endpoints are exactly
\begin{equation}\tag{ECL2}
uK(k)=\frac{\alpha\pi}{\beta},\qquad
vE(k)=\frac{(\alpha+2)\pi}{\beta},\qquad
k^2=1-\frac{u^2}{v^2},\qquad 0<u<v.
\end{equation}
Here $K$ and $E$ have the modulus convention of EIQ2, not the
elliptic-parameter convention. Define the original endpoint quantities
\begin{equation}\tag{ECL3}
w=\frac{u+v}{2},\qquad z=uv,\qquad
c_{\mathrm{cap}}=\frac{v^2-u^2}{4},\qquad
d\nu(x)=\frac{dx}{\pi\sqrt{(b-x)(x-a)}}\quad(a<x<b).
\end{equation}
The letter $z$ in ECL3 is a positive scalar. A Cauchy-transform variable
below is denoted by $\lambda$. The measure $\nu$ has mass one by
$x=(a+b)/2+(b-a)\cos\theta/2$, $0\leq\theta\leq\pi$.

\subsection{Arcsine identities with their original constants}

Put $m=(a+b)/2$, $d=(b-a)/2$, and $\eta=(v-u)/(v+u)$.
Then $0<\eta<1$ and
\[
x=m+d\cos\theta
 =w^2(1+2\eta\cos\theta+\eta^2),\qquad
z=w^2(1-\eta^2).
\]
The absolutely uniformly convergent series
\[
\begin{split}
\log x&=2\log w+
  2\sum_{n\geq1}\frac{(-1)^{n+1}\eta^n}{n}\cos(n\theta),\\
\frac1x&=\frac1z\left(1+
  2\sum_{n\geq1}(-\eta)^n\cos(n\theta)\right)
\end{split}
\]
follow from the geometric series and the integral of that series.
Termwise integration, and orthogonality of the cosines, prove
\begin{equation}\tag{ECL4}
\begin{split}
\int\log x\,d\nu(x)&=2\log w,\qquad
\int\frac{d\nu(x)}x=\frac1z,\\
\int\frac{\log x}{x}\,d\nu(x)
 &=\frac{2\log w+2\log(1-\eta^2)}z
 =\frac{2\log z-2\log w}z.
\end{split}
\end{equation}
In the product of the two Fourier series, the contribution of degree
$n\geq1$ to the constant coefficient is
$-2\eta^{2n}/(nz)$, giving the displayed logarithm without an omitted
constant.

The logarithmic potential of $\nu$ on its entire support is
\begin{equation}\tag{ECL5}
\int_a^b\log|x-t|\,d\nu(x)=\log c_{\mathrm{cap}},
\qquad a\leq t\leq b.
\end{equation}
Here is a direct proof, including the endpoint values. Write
$t=m+d\cos\phi$ and use
\[
|m+d\cos\theta-t|
=\frac d2
 |e^{i\theta}-e^{i\phi}|\,
 |e^{i\theta}-e^{-i\phi}|.
\]
The integrand depending on $\cos\theta$ has the same mean on
$[0,\pi]$ as on $[0,2\pi]$. For $0<r<1$, the power series of
$\log|1-re^{i\theta}|$ has zero circle mean. This mean tends to
the mean at $r=1$: near its only singular angle,
$|1-re^{i\theta}|^2=(1-r)^2+2r(1-\cos\theta)$,
so for $r\geq1/2$ the absolute value of its negative logarithm
is bounded by $C+|\log|\theta||$, and its positive logarithm
is bounded by $\log2$. These are integrable bounds. A rotation
proves that each of the two circle means above is zero, for
every $\phi$, including $\phi=0,\pi$. Hence their sum is
$\log(d/2)=\log c_{\mathrm{cap}}$, proving ECL5.

All logarithmic double integrals used below are absolutely integrable.
In fact the absolute logarithm for two arcsine variables is bounded by
a constant plus
\[
\left|\log\left|\sin\frac{\theta-\phi}{2}\right|\right|
+\left|\log\left|\sin\frac{\theta+\phi}{2}\right|\right|.
\]
Each term is integrable on $[0,\pi]^2$; near a zero, use
$\sin t\geq 2t/\pi$ for $0\leq t\leq\pi/2$, and integrate
$|\log t|$. Thus the same assertion holds for any two signed
measures whose absolute values are bounded by constant multiples
of $\nu$. The EIQ density also has this domination, because
$\rho(x)/[1/(\pi\sqrt{(b-x)(x-a)})]$ is bounded on $[a,b]$.

Let $\ell$ have the exact EIQ17 convention
\[
V_{\alpha,\beta}(x)
-2\int\log|x-t|\,d\mu_{\alpha,\beta}(t)=\ell
\quad(a\leq x\leq b).
\]
Integrating this equality against $\nu$, using ECL4--5 and
EIQ7, gives
\begin{equation}\tag{ECL6}
\boxed{\displaystyle
\ell=2(\alpha+2)-2\alpha\log w-2\log c_{\mathrm{cap}}.}
\end{equation}
Indeed $\int\sqrt{x}\,d\nu(x)=2vE(k)/\pi$, whose product with
$\beta$ is $2(\alpha+2)$ by ECL2. Fubini is justified by the
absolute integrability just proved and the unit mass of $\mu$.
Smoothness of $\ell$ now follows from the endpoint smoothness in EIQ6.

\subsection{Differentiation with the moving endpoints retained}

All derivatives in this subsection are with respect to $\alpha$ at
fixed $\beta$. Put
\[
p_\alpha(x)=\frac{h_{\alpha,\beta}(x)}{2\pi},\qquad
D(x)=(b-x)(x-a).
\]
In this subsection a dot denotes a parameter derivative, so that
$\dot p(x)$ differentiates the function $p_\alpha(x)$ at fixed $x$.
The function $h_{\alpha,\beta}$ is exactly EIQ10; the subscript on
$p_\alpha$ labels the parameter and does not denote differentiation.
EIQ10 and EIQ29 prove that $p_\alpha$ and its parameter and spatial
derivatives are bounded on a fixed positive compact neighbourhood
of the supports when parameters range in a sufficiently small
compact neighbourhood.

The derivative of the probability measure, interpreted by integration
against every continuously differentiable function on this neighbourhood,
is the finite signed measure
\begin{equation}\tag{ECL7}
\begin{split}
\sigma(dx)&=\left[
\dot p(x)\sqrt{D(x)}
+\frac{p_\alpha(x)}{2\sqrt{D(x)}}
 \{\dot b(x-a)-\dot a(b-x)\}\right]\mathbf1_{(a,b)}(x)\,dx,\\
\frac{\partial}{\partial\alpha}\int f\,d\mu_{\alpha,\beta}
&=\int f\,d\sigma,\qquad \sigma([a,b])=0.
\end{split}
\end{equation}
This identity includes the moving support. To prove it without
postulating endpoint terms, set $\delta=b-a$ and $x=a+\delta y$.
The measure in the fixed variable $y$ is
\[
\delta^2 p_\alpha(a+\delta y)\sqrt{y(1-y)}\,dy.
\]
Its differentiated integral against $f$ is
\[
\int_0^1\sqrt{y(1-y)}
\left[
2\delta\dot\delta\,p_\alpha f
+\delta^2(\dot p+p_\alpha'\dot x)f
+\delta^2p_\alpha f'\dot x
\right]dy,\qquad \dot x=\dot a+\dot\delta y.
\]
Replace $\delta f'$ in the last summand by $d f(x)/dy$ and
integrate by parts. Its boundary term is
$[\delta p_\alpha(x)\dot x\sqrt{y(1-y)}f(x)]_0^1=0$.
The remaining terms combine, on returning to $x$, to
$\dot p\sqrt D+p_\alpha\dot D/(2\sqrt D)$, where
$\dot D=\dot b(x-a)-\dot a(b-x)$, exactly ECL7.
This integration by parts is justified first on
$[\varepsilon,1-\varepsilon]$ and then by passage to the limit:
the boundary expression tends to zero and the differentiated
square-root factor is bounded by a constant times
$[y(1-y)]^{-1/2}$, an integrable function. The differentiated
fixed-interval integral is valid by the uniform bounds above.
Taking $f=1$ proves the zero total mass in ECL7.
The displayed density also proves $|\sigma|\leq C\nu$.
In particular no endpoint atom has appeared or been discarded.

For fixed $x\in(a,b)$, differentiating the Euler equality gives
\begin{equation}\tag{ECL8}
-\log x-2\int\log|x-t|\,d\sigma(t)=\dot\ell.
\end{equation}
For precision, choose a smooth function of $t$ equal to one near
$x$, with support in a common interior subinterval for nearby
parameters. On this subinterval the density and its parameter
derivative are uniformly bounded, so differentiation of its
logarithmic integral follows by domination by
$C|\log|x-t||$. The complementary integrand is a continuously
differentiable test function, so ECL7 applies. This proves ECL8
despite the logarithmic singularity and the moving endpoints.

We next identify $\sigma$ exactly:
\begin{equation}\tag{ECL9}
\boxed{\displaystyle
\frac{\partial}{\partial\alpha}\mu_{\alpha,\beta}(dx)
=\sigma(dx)=\frac12\left(1-\frac z x\right)d\nu(x).}
\end{equation}
The derivative has the test-function meaning proved in ECL7,
which in particular applies to the logarithmic and square-root
moments because $a>0$.

Here is an explicit uniqueness argument establishing ECL9.
Denote its proposed right side by $\sigma_0$.
By ECL4 it has zero total mass, and $|\sigma_0|\leq C\nu$.
Retain the branch
$R(\lambda)=\sqrt{(\lambda-a)(\lambda-b)}$ with
$R(\lambda)/\lambda\to1$ at infinity, exactly as in EIQ12.
The arcsine transform proved there and partial fractions yield
\begin{equation}\tag{ECL10}
\begin{split}
\int\frac{d\sigma_0(t)}{\lambda-t}
&=\frac1{2R(\lambda)}
-\frac z{2\lambda}
 \left(\frac1z+\frac1{R(\lambda)}\right)\\
&=\frac{1-z/\lambda}{2R(\lambda)}-\frac1{2\lambda}.
\end{split}
\end{equation}
There is no pole at zero: $R(0)=-uv=-z$, so the numerator
$(\lambda-z)/R(\lambda)-1$ vanishes there. The integral
itself is analytic in a neighbourhood of zero, as is also
immediate from $t\geq a>0$.
For $a<x<b$, the real principal value of the arcsine transform
is zero, because $R(x+i0)$ is purely imaginary. This limiting
principal value also follows directly by subtracting the local
value of the smooth interior density in its Cauchy integral.
Consequently ECL10 gives
\[
\operatorname{PV}\int\frac{d\sigma_0(t)}{x-t}=-\frac1{2x}.
\]
Differentiating the logarithmic potential on each interior
subinterval is valid by that same local subtraction; the endpoint
parts are integrable and are a positive distance away.
It follows that
$-\log x-2\int\log|x-t|\,d\sigma_0(t)$ is constant on $(a,b)$.
Subtract this identity from ECL8. For
$\sigma_1=\sigma-\sigma_0$, the logarithmic potential
$\int\log|x-t|\,d\sigma_1(t)$ is constant on $(a,b)$.
Since $\sigma_1$ has zero mass and no endpoint atoms, and
$|\sigma_1|\leq C\nu$, absolute integrability proved after ECL5 gives
\[
-\iint\log|x-t|\,d\sigma_1(x)d\sigma_1(t)=0.
\]
The exact zero-mass energy argument EIQ20--22, including its
Fourier-uniqueness proof, forces $\sigma_1=0$. This establishes
the claimed derivative as a measure, with every endpoint
contribution in ECL7 retained.

\subsection{The moment and energy formulas}

Integrate ECL8 against $\nu$. By ECL5 and $\sigma([a,b])=0$,
the iterated logarithmic term is zero. Therefore
\begin{equation}\tag{ECL11}
\dot\ell=-2\log w.
\end{equation}
All exchanges again have the absolute integrability established
after ECL5; the equality ECL8 holds $\nu$-almost everywhere.
Using ECL9 with $f(x)=\log x$, and then ECL4, proves
\begin{equation}\tag{ECL12}
\boxed{\displaystyle
\mathcal L_\alpha
=\frac12\int\left(1-\frac z x\right)\log x\,d\nu(x)
=2\log w-\log z=\log(w^2/z).}
\end{equation}

Write $I=\iint\log|x-t|\,d\mu(x)d\mu(t)$, with the original
unscaled logarithmic kernel. The EIQ31 dilation identity is
$\beta\int\sqrt{x}\,d\mu=2(\alpha+1)$.
Integration of the Euler equality against $\mu$ gives
$\ell=\int V\,d\mu-2I$. Combining this with
$F=I-\int V\,d\mu$ yields the exact identity
\begin{equation}\tag{ECL13}
\ell=-2F-2(\alpha+1)+\alpha\mathcal L.
\end{equation}
EIQ30 proves $F_\alpha=\mathcal L$. Differentiating ECL13 at
fixed $\beta$, with the differentiability just proved, gives
$\dot\ell=-\mathcal L-2+\alpha\mathcal L_\alpha$.
Substituting ECL11--12 now proves, without an integration
constant or a limiting parameter choice,
\begin{equation}\tag{ECL14}
\boxed{\displaystyle
\mathcal L=2(\alpha+1)\log w-\alpha\log z-2,\qquad
F=-\frac{\ell}{2}-(\alpha+1)+\frac{\alpha\mathcal L}{2}.}
\end{equation}
Equations ECL6, ECL12, and ECL14 therefore evaluate the same
Robin constant, logarithmic moment, and variational energy as
the original EIQ integrals.

\subsection{The original outer coefficient and the two scale conventions}

The coefficient of OFV29 is defined by its full original expression
\begin{equation}\tag{ECL15}
\mathfrak C_{\partial}
=4\log(4/\vartheta)-10
-2F(2,\vartheta)+4\mathcal L(2,\vartheta),
\qquad \vartheta=\pi/2.
\end{equation}
To express this same number at the original EIQ limiting potential
$V_{2,\pi}(x)=\pi\sqrt{x}-2\log x$, retain the actual dilation
$x\mapsto(\pi/\beta)^2x$ of EIQ32. Its effect on the interaction,
potential integral, and moment is, respectively,
\[
\begin{split}
I(\alpha,\beta)&=I(\alpha,\pi)+2\log(\pi/\beta),\\
\int V_{\alpha,\beta}\,d\mu_{\alpha,\beta}
&=\int V_{\alpha,\pi}\,d\mu_{\alpha,\pi}
   -2\alpha\log(\pi/\beta),\\
\mathcal L(\alpha,\beta)&=\mathcal L(\alpha,\pi)
   +2\log(\pi/\beta).
\end{split}
\]
Thus $F(\alpha,\beta)=F(\alpha,\pi)
+2(\alpha+1)\log(\pi/\beta)$, with its stated sign.
At $\alpha=2$ the added coefficient in ECL15 is
$(4-12+8)\log(\pi/\beta)=0$. It follows exactly that
\begin{equation}\tag{ECL16}
c_{\partial}:=\mathfrak C_{\partial}
=4\log(4/\pi)-10-2F(2,\pi)+4\mathcal L(2,\pi).
\end{equation}
For the rest of this subsection $u,v,w,z,c_{\mathrm{cap}}$ are
precisely their ECL2--3 values at $(\alpha,\beta)=(2,\pi)$.
They therefore satisfy $uK(k)=2$ and $vE(k)=4$.
Substitution of the established formulas, with their constants,
gives
\begin{equation}\tag{ECL17}
\boxed{\displaystyle
c_{\partial}
=4\log(4/\pi)+8\log w-4\log z-2\log c_{\mathrm{cap}}.}
\end{equation}
For example, ECL6 gives
$\ell=8-4\log w-2\log c_{\mathrm{cap}}$, ECL14 gives
$\mathcal L=6\log w-2\log z-2$, and ECL13 gives
$F=-\ell/2-3+\mathcal L$. In ECL16 these yield
$4\log(4/\pi)-4+\ell+2\mathcal L$, in which the constant
terms are exactly $-4+8-4=0$, proving ECL17.
Equivalently $F(2,\pi)=I(2,\pi)-6+2\mathcal L(2,\pi)$
in ECL16 gives
$c_{\partial}=4\log(4/\pi)+2-2I(2,\pi)$,
which is precisely the OFV31 energy convention.
This section evaluates this coefficient. It does not assign it
as the leading coefficient of a separate observation residual
or a separate kernel determinant.

\subsection{An integer-only outward certificate}

The accompanying \texttt{certify\_closed\_equilibrium.py} proves
\begin{equation}\tag{ECL18}
\boxed{21.66851180520<16c_{\partial}<21.66851180521.}
\end{equation}
The receipt stores every interval endpoint as an integer divided by
$S=10^{70}$, including the original endpoints and all quantities in
ECL17. The script uses no floating-point inputs or numerical quadrature.
The following argument proves the correctness of its finite algorithm.

Each stored interval $[A,B]$ denotes $[A/S,B/S]$. Addition and
subtraction are exact. For multiplication, the smallest and largest
of the four endpoint products are divided by $S$ and rounded
respectively down and up. For division by an interval excluding zero,
the four ratios of endpoint integers, multiplied by $S$, are
rounded outwards. Their extrema enclose every possible product or
quotient because these operations attain their extrema at corners
of their rectangular domains. The rounding uses integer floor
division and $\lceil n/d\rceil=-\lfloor-n/d\rfloor$ after making
$d>0$. These formulas apply to negative numerators as well.
Consequently induction proves containment after every arithmetic
operation, without a probabilistic or floating-point error estimate.

Let $h=u/v$ in this paragraph only; this is a scalar ratio and
does not replace the density function $h(x)$ of EIQ10. Put
$m_e=1-h^2$, so $k=\sqrt{m_e}$. Define
$b_n=\binom{2n}{n}/4^n$ and $t_n=b_n^2m_e^n$.
Expansion of the binomial series, followed by its uniformly
convergent integration for $m_e<1$, proves
\begin{equation}\tag{ECL19}
\frac{2K(\sqrt{m_e})}{\pi}=\sum_{n=0}^{\infty}t_n,\qquad
\frac{2E(\sqrt{m_e})}{\pi}
=1-\sum_{n=1}^{\infty}\frac{t_n}{2n-1},\qquad
\frac{t_n}{t_{n-1}}=m_e\frac{(2n-1)^2}{4n^2}.
\end{equation}
Indeed integration by parts gives
$\int_0^{\pi/2}\sin^{2n}\theta\,d\theta=(\pi/2)b_n$;
the two binomial coefficients are $b_n$ and
$-b_n/(2n-1)$ for $n\geq1$. Since
$t_{n+1}/t_n\leq m_e$, the positive omitted tails after
indices $0,\ldots,N-1$ have the rigorous bounds
\begin{equation}\tag{ECL20}
0\leq\sum_{n=N}^{\infty}t_n\leq\frac{t_N}{1-m_e},
\qquad
0\leq\sum_{n=N}^{\infty}\frac{t_n}{2n-1}
\leq\frac{t_N}{(2N-1)(1-m_e)}.
\end{equation}
The script takes $N=3200$, applies the recurrence with outward
rounding at each step, and adds these tails with their proper
signs. They remain valid for interval values of $m_e$ by the
interval arithmetic just proved.

The two exact terminating decimals used for the root bracket are
\begin{equation}\tag{ECL21}
\begin{split}
h_-&=0.160101670957718355884567453115,\\
h_+&=0.160101670957718355884567453116.
\end{split}
\end{equation}
Let $K_b=2K/\pi$, $E_b=2E/\pi$ as in ECL19.
The receipt gives strictly negative upper bound for
$2h_-K_b(1-h_-^2)-E_b(1-h_-^2)$ and strictly positive lower
bound for the corresponding expression at $h_+$.
Here the arguments of $K_b,E_b$ denote $m_e$ as defined by
their series, not the modulus argument of EIQ2.
The common factor $\pi/2>0$ in the original equation cancels
exactly. EIQ6 proves that $hK(\sqrt{1-h^2})/
E(\sqrt{1-h^2})$ increases strictly from zero to one as
$h$ increases from zero to one. Thus ECL21 encloses the
unique root of the original endpoint equation
$hK(\sqrt{1-h^2})=E(\sqrt{1-h^2})/2$.

The value of $\pi$ is enclosed independently by
\begin{equation}\tag{ECL22}
\pi=16\arctan(1/5)-4\arctan(1/239).
\end{equation}
To prove the identity, if $\theta=\arctan(1/5)$, then
$\tan(2\theta)=5/12$ and $\tan(4\theta)=120/119$.
The subtraction formula gives
$\tan(4\theta-\arctan(1/239))=1$.
This angle lies in $(0,\pi/2)$:
$4\theta<4/5<1<\pi/2$, while
$4\theta\geq4(1/5)/(1+1/25)>1/239$.
The inequality $\pi/2>1$ follows from
$\arctan1=\int_0^1(1+t^2)^{-1}dt>1/2$.
Consequently the angle is $\pi/4$, proving ECL22.
Each arctangent is computed from its alternating series with
128 terms. The first omitted term bounds the absolute remainder;
its sign is positive for this even number of terms. Outward
arithmetic and that signed tail therefore enclose the exact $\pi$.

For every positive logarithm argument $x$ put $t=(x-1)/(x+1)$.
The geometric series, integrated on the segment from zero to $t$,
gives
\begin{equation}\tag{ECL23}
\log x=2\sum_{n=0}^{M-1}\frac{t^{2n+1}}{2n+1}
+R_M,\qquad
|R_M|\leq
\frac{2|t|^{2M+1}}{(2M+1)(1-|t|^2)}.
\end{equation}
The script takes $M=256$. For an interval argument it uses
the maximum absolute endpoint of the enclosing $t$ interval
in this tail bound and adds a symmetric outward tail. Thus
ECL23 encloses the logarithm throughout the argument interval.

The enclosed original endpoint values are computed in order by
\[
\begin{gathered}
K=(\pi/2)K_b,\qquad E=(\pi/2)E_b,\qquad v=4/E,\qquad u=hv,\\
w=(u+v)/2,\qquad z=uv,\qquad c_{\mathrm{cap}}=(v^2-u^2)/4.
\end{gathered}
\]
The script checks $0<u<v$, positive capacity, strict endpoint
signs, and that every division interval excludes zero. It applies
ECL23 to the four arguments $4/\pi,w,z,c_{\mathrm{cap}}$
and evaluates the full expression ECL17 with its original
coefficients. Integer comparison of the final bounds proves
the stronger enclosure
\begin{equation}\tag{ECL24}
\begin{split}
21.668511805204674126143836401790042003
&<16c_{\partial}\\
&<21.668511805204674126143836403323660355.
\end{split}
\end{equation}
The decimal endpoints in ECL24 are rounded outwards from the
stored integer endpoints. Their comparison with the exact
terminating decimals in ECL18 proves ECL18.

The same operations also enclose the endpoint expression
\begin{equation}\tag{ECL25}
\frac18\log(w^2/z)-\frac14\log2
\ \in\ (-0.080453400358,\,-0.080453400357).
\end{equation}
This is an evaluation of this explicit expression. If it is used
as an earlier approximation centre, its previously specified
error radius is not changed by this calculation.
The ordinary and optimized Python executions produce identical
JSON receipts. All mathematical checks use explicit exceptions,
so disabling Python assertions does not disable any certificate
condition. The receipt includes the script hash, integer endpoint
signs, elliptic-tail bounds, every resulting interval, and the
strict target comparisons.

% END CURRENT ECL

\clearpage
% BEGIN CURRENT EQR
% EQR backward propagation. Frozen ECL1--25 is unchanged.
% Full original provider: flat_web_stage/built/02_FULL_GAMMA_PROOFS.tex
% SHA256 80c8e48b04cf032311a93dfc40c01969c9999412eab113068fc26d6eb2b8c7dd
% Complete QGQ1--9 section: original provider lines 5128--5342.
% QGT9 and QGT12 are in the same full provider.
\section{The closed coefficient in the original QGQ return centre}
\label{sec:EQR}

Retain the complete original QGQ1--9 theorem and its full provider
\texttt{02\_FULL\_GAMMA\_PROOFS.tex}, section
\emph{The finite original \(m=1\) return bound at the \(q\) scale}.
Retain its actual Gamma functional $W_k$, finite centre
$\mathcal V^{\rm c}_{q,l}$, constants $\mathcal C_\Gamma,A,M,C$,
and its original packet
\[
m=1,\qquad k=4l+1,\qquad q=(4l+2)^2,\qquad n=q/2,\qquad l\geq1.
\]
The support and endpoint values $u,v,w,z$ in this section are
exactly those of ECL2--3 at $(\alpha,\beta)=(2,\pi)$.
In particular $uK(k_{\rm ell})=2$, $vE(k_{\rm ell})=4$, and
$k_{\rm ell}^2=1-u^2/v^2$; the subscript distinguishes this
elliptic modulus from the original integer packet degree $k$.
No endpoint or original degree is changed by this notation.

QGQ4 defines, for the original EIQ logarithmic moment,
$\mathcal L_1=\partial_\alpha\mathcal L(2,\pi)$ and
$\mathcal C_\Gamma^{[q,1]}=\mathcal L_1/8-(\log2)/4$.
The derivative proved in ECL12 is a derivative of precisely
this same moment, at fixed original $\beta$. Therefore the
following is an exact identification of the original centre:
\begin{equation}\tag{EQR1}
\boxed{\begin{gathered}
\mathcal L_1=\log(w^2/z),\qquad
\mathcal C_\Gamma^{[q,1]}
=\frac18\log(w^2/z)-\frac14\log2,\\
-0.080453400358
<\mathcal C_\Gamma^{[q,1]}
<-0.080453400357.
\end{gathered}}
\end{equation}
The first two equalities follow by substitution into the literal
QGQ4 definition; there is no new centre or change of energy
convention. The last line is consequently the strict outward
integer enclosure ECL25 for that exact original centre. Its
proof and stored integer endpoints are those of ECL19--25.

The original finite-centre theorem QGQ8 now gives
\begin{equation}\tag{EQR2}
\boxed{\displaystyle
\frac{\mathcal V^{\rm c}_{q,l}-lq\mathcal C_\Gamma}{q}
\longrightarrow
\frac18\log(w^2/z)-\frac14\log2
\qquad(l\longrightarrow\infty).}
\end{equation}
To check the connection to the original finite calculation
explicitly, QGQ7 retains its first term
$(2\mathcal L_1-4\log2)l^2$ and all its other original terms.
The complete termwise estimates following QGQ7 bound their
sum by $O(l)$; in particular its actual remainder
$\zeta_{q,l}$ is $O(l)$ by QGQ6. On the literal packet,
$l^2/q\to1/16$ and $l/q\to0$. Thus dividing that same identity
by $q$ gives
$(2\mathcal L_1-4\log2)/16=\mathcal L_1/8-(\log2)/4$.
Substitution of EQR1 proves EQR2 with the original finite centre,
source, low endpoint, and rank correction intact.

The original error radius remains the following nonzero number:
\begin{equation}\tag{EQR3}
\boxed{\displaystyle
\limsup_{l\to\infty}
\left|
\frac{W_k-lq\mathcal C_\Gamma}{q}
-\left(\frac18\log(w^2/z)-\frac14\log2\right)
\right|
\leq\frac A{\sqrt2},
\qquad A=M/2+2C>0.}
\end{equation}
Here $M$ and $C$ retain their complete original QGT12 and QGT9
definitions. To verify the transfer, retain the finite QGT22
inequality used in QGQ9:
\[
-p^-_{q,l}-\mathcal U_{q,l}-E_{q,l}
\leq W_k-\mathcal V^{\rm c}_{q,l}
\leq p^+_{q,l}+\mathcal U_{q,l}+T_{q,l}+E_{q,l}.
\]
QGQ1--2 proves that each of
$p^-_{q,l}/q,p^+_{q,l}/q,T_{q,l}/q,E_{q,l}/q$ tends to zero,
while $\mathcal U_{q,l}/q\to A/\sqrt2$. Hence the difference
$(W_k-\mathcal V^{\rm c}_{q,l})/q$ is bounded below by
$-A/\sqrt2-o(1)$ and above by $A/\sqrt2+o(1)$.
Adding the convergent expression in EQR2, and using the triangle
inequality, proves EQR3. Both unequal finite errors remain in the
displayed original inequality; their passage to the same limiting
radius does not replace that finite interval by a new one.

The radius is strictly positive from the original definitions
themselves. QGT9 gives $C\geq0$, and QGT12 gives
\[
M=\max_{(\alpha,\beta)\in\mathcal K}
|\partial_\alpha\mathcal L(\alpha,\beta)|
\geq\partial_\alpha\mathcal L(2,\pi)=\log(w^2/z)>0,
\]
because $(2,\pi)\in\mathcal K$ and
$w^2/z=1+(v-u)^2/(4uv)>1$ by $0<u<v$.
Consequently $A=M/2+2C>0$. This proves that the retained
nonzero radius has not disappeared upon evaluating its centre.
The only convergence asserted in EQR2 concerns the original
finite centre. Neither EQR1 nor EQR3 asserts a limit for the
actual quantity $(W_k-lq\mathcal C_\Gamma)/q$.

% END CURRENT EQR

\clearpage
% BEGIN CURRENT RWC
% RWC: independent reconstruction of the original full-window residual update.
\section{Strict full-window growth of the original residual covariance}
\label{sec:RWC}

This calculation starts from the complete ROQ1--27 construction above.
The incoming first-variation handoff states a full-window update and
trace--determinant bounds. The following proof derives those statements
from the original columns, and adds exact exterior-power formulas for
every intermediate coefficient. No unavailable FVR proof is used.
Retain $k=4l+1\ge9$, $q=(k+1)^2$, $q'=(k-7)^2$, $j=8k-32$,
the fixed actual period in the five-orbit locus, and $s\in\{1,k\}$.
The actual matrices $C:\mathbb C^q\to\mathbb C^{q'}$ and
$D:\mathbb C^q\to\mathbb C^j$ are ROQ1 and ROQ4, with all their
coefficients. In particular $C$ is onto and $D|_{\ker C}$ is onto.
The original roots, monic Gamma polynomials and their squared norms are
\begin{equation}\tag{RWC1}
\begin{gathered}
\omega_{ab}=(2b-k)\gamma-i(2a-k)\delta,\quad 0\le a,b\le k,
\qquad \delta>0,\quad\gamma>0,\\
p_0=1,\quad p_1=y,\quad
p_{n+1}=yp_n-n(s/2+n-1)p_{n-1},\\
h_n=(2\pi)^{s/2}n!(s/2)_n,\qquad
u_n=(p_n(\omega_{ab})/\sqrt{h_n})_{ab},\qquad
R_N=\sum_{n=0}^N u_nu_n^*.
\end{gathered}
\end{equation}
The positive original measure $dm_s$ has the orthogonality and norms
in RWC1, as proved in the Gamma provider. The order of the $q$ roots
is the original ordered primary order; stars below are conjugate
transposes. The first $q$ monic polynomials have invertible evaluation
matrix, since these $q$ roots are distinct. Consequently $R_N>0$
for $N\ge q-1$.

\subsection{Every consecutive full block is invertible}
The complete real polynomial of the original root set is
\begin{equation}\tag{RWC2}
\begin{split}
P_k(y)&=\prod_{a=0}^k\prod_{b=0}^k(y-\omega_{ab})\\
&=\prod_{a=0}^{(k-1)/2}\prod_{b=0}^k
\bigl((y-(2b-k)\gamma)^2+(k-2a)^2\delta^2\bigr)>0
\quad(y\in\mathbb R).
\end{split}
\end{equation}
Indeed the factor indexed by $k-a,b$ is the conjugate of the
factor indexed by $a,b$, and no factor in the second product can
vanish on the real line. This proves the identity with its monic
constant and all factors retained.

For any integer $N\ge q-1$ put
\begin{equation}\tag{RWC3}
U_N^+=[u_{N+1},\ldots,u_{N+q}]\in\operatorname{Mat}_{q\times q}(\mathbb C).
\qquad \det U_N^+\ne0.
\end{equation}
To prove the assertion suppose $U_N^+c=0$ and form
$f(y)=\sum_{a=1}^q c_a p_{N+a}(y)/\sqrt{h_{N+a}}$.
It vanishes at all $q$ distinct roots, so polynomial division gives
$f=P_k r$ with $\deg r\le N$. Orthogonality of every $p_{N+a}$
to every polynomial of degree at most $N$ gives
\[
0=\int_{\mathbb R}f(y)\overline{r(y)}\,dm_s(y)
 =\int_{\mathbb R}P_k(y)|r(y)|^2\,dm_s(y).
\]
All integrals exist, because the measure has every polynomial moment.
Its strictly positive density and RWC2 force $r=0$. The distinct
degrees of the monic $p_{N+a}$ then force every $c_a=0$.
This argument applies at both $N=q-1$ and $N=q$, and also at
every larger integer cutoff; it makes no genericity choice of $D$.

\subsection{The simultaneous conditioned update}
Use the complete original matrices
\begin{equation}\tag{RWC4}
\begin{gathered}
A_N=CR_NC^*,\quad
W_N=R_N-R_NC^*A_N^{-1}CR_N,\quad
\Omega_N=DW_ND^*>0,\quad Q_{\Xi,N}=\Omega_N^{-1},\\
\Pi_N=I_q-R_NC^*A_N^{-1}C,\quad
X_N=CU_N^+,\quad Z_N=D\Pi_N U_N^+,\\
B_N=I_q+X_N^*A_N^{-1}X_N>0,\qquad
V_N=Z_NB_N^{-1}Z_N^*.
\end{gathered}
\end{equation}
These use $A_N$ solely for this covariance block, rather than for
the conductor biform. Since $A_N>0$, every displayed inverse exists.
The equalities $C\Pi_N=0$ and $\Pi_Nx=x$ for $Cx=0$ show that
$\Pi_N$ is the exact projection onto $\ker C$ with complementary
image $R_NC^*\mathbb C^{q'}$. It is generally not the Euclidean
orthogonal projection. RWC3 and surjectivity of $D|_{\ker C}$ give
\begin{equation}\tag{RWC5}
\operatorname{im}(\Pi_N U_N^+)=\ker C,\qquad
\operatorname{rank}Z_N=j,\qquad V_N>0.
\end{equation}
For the last assertion, $Z_N^*$ is injective and $B_N^{-1}>0$,
so $z^*V_Nz>0$ for each nonzero $z\in\mathbb C^j$.

There are exact identities
\begin{equation}\tag{RWC6}
\begin{gathered}
W_{N+q}=W_N+(\Pi_NU_N^+)B_N^{-1}(\Pi_NU_N^+)^*,\\
\Omega_{N+q}=\Omega_N+V_N,\qquad
Q_{\Xi,N+q}=(\Omega_N+V_N)^{-1}.
\end{gathered}
\end{equation}
Here is the complete matrix calculation. Temporarily write
$R=R_N$, $U=U_N^+$, $A=A_N$, $X=CU$, $B=B_N$ and
$T=RC^*A^{-1}X$. Direct multiplication proves
\[
(A+XX^*)^{-1}
 =A^{-1}-A^{-1}XB^{-1}X^*A^{-1}.
\]
Indeed $(A+XX^*)A^{-1}X=XB$, so the subtracted term cancels
$XX^*A^{-1}$. Since $CR_{N+q}C^*=A+XX^*$, the new subtracted
covariance product equals
\[
RC^*A^{-1}CR+TU^*+UT^*+U(B-I_q)U^*
 -[T+U(B-I_q)]B^{-1}[T+U(B-I_q)]^*.
\]
Subtract this expression from $R+UU^*$ and then subtract $W_N$.
Writing $T+U(B-I_q)=UB-(U-T)$, expansion leaves exactly
$(U-T)B^{-1}(U-T)^*$. But $U-T=\Pi_NU$, proving the first
identity in RWC6. Multiplication by $D,D^*$ proves the second,
and ROQ16 proves the third. In particular the full $q\times q$
matrix $B_N$ and all its off-diagonal entries remain in the update.

\subsection{All exterior coefficients and finite bounds}
Let
\begin{equation}\tag{RWC7}
Y_N=\Omega_N^{-1/2}Z_NB_N^{-1/2},\quad
T_N=Y_NY_N^*>0,\quad
t_N=\operatorname{Tr}T_N,\quad
c_N=\det T_N=\frac{\det V_N}{\det\Omega_N}>0.
\end{equation}
The square roots are the positive Hermitian square roots; no
source mass, coordinate phase or rank has been changed.
If $\lambda_1,\ldots,\lambda_j>0$ are the eigenvalues of $T_N$,
then RWC6 gives
\begin{equation}\tag{RWC8}
L_N:=\log\frac{\det\Omega_{N+q}}{\det\Omega_N}
 =\sum_{a=1}^j\log(1+\lambda_a)
 =\log\left(\sum_{d=0}^j e_d(T_N)\right),
\quad
e_d(T_N)=\|\bigwedge^d Y_N\|_{\mathrm{HS}}^2.
\end{equation}
The exterior norm is the Hilbert--Schmidt norm in the standard
orthonormal wedge bases, and $e_0=1$. To verify the last equality,
the singular values of $Y_N$ are $\sqrt{\lambda_a}$, and its
$d$th exterior power has squared singular values
$\prod_{a\in I}\lambda_a$ for the $d$-element subsets $I$.
Their sum is the elementary symmetric polynomial $e_d$.
Equivalently Cauchy--Binet writes it as the sum of squared moduli
of all $d\times d$ minors of the literal $Y_N$. This proves
every coefficient in the determinant expansion, including all
intermediate exterior powers; the sole determinant $c_N$ is its
last coefficient.

The following finite inequalities now hold:
\begin{equation}\tag{RWC9}
\max\{\log(1+t_N),\ j\log(1+c_N^{1/j})\}
\le L_N\le j\log(1+t_N/j).
\end{equation}
For the first lower bound, expand $\prod(1+\lambda_a)$: the
constant term is one, its degree-one term is $t_N$, and all other
terms are positive. For the second, the arithmetic--geometric
mean inequality on the $\binom jd$ positive $d$-fold products
gives
$e_d\ge\binom jd c_N^{d/j}$, since each eigenvalue occurs in
$\binom{j-1}{d-1}$ products. Summing this inequality for
$0\le d\le j$ yields $(1+c_N^{1/j})^j$.
Finally the concavity of $\log(1+x)$ on $x>0$ gives the upper
bound by averaging the $j$ positive eigenvalues. These arguments
also prove useful finite refinements: retaining any chosen initial
set $e_0,\ldots,e_d$ in RWC8 gives its exact positive partial-sum
lower bound; no rank fraction is substituted for an eigenvalue.

\subsection{Return to the original four endpoints}
Use precisely the ordered cutoffs $(q-1,q,2q-1,2q)$, with signs
$(1,1,-1,-1)$. ROQ14 and ROQ16 give
\begin{equation}\tag{RWC10}
\Phi_{k,s}=F_{\Xi,k}^{(s)}=L_{q-1}+L_q>0.
\end{equation}
Define the actual finite bounds
\begin{equation}\tag{RWC11}
\begin{gathered}
\ell_{k,s}=\sum_{N\in\{q-1,q\}}
\max\{\log(1+t_N),j\log(1+c_N^{1/j})\}>0,\\
u_{k,s}=\sum_{N\in\{q-1,q\}}j\log(1+t_N/j),
\qquad \ell_{k,s}\le\Phi_{k,s}\le u_{k,s}.
\end{gathered}
\end{equation}
These are complete expressions in the actual $C,D,u_n$ and their
Gram matrices. Together with ROQ21 they give the stronger endpoints
$\Phi^-_{k,s}=\max(\ell_{k,s},\lambda_s)$ and
$\Phi^+_{k,s}=\min(u_{k,s},
j(\mathcal L_{s,q}+\mathcal L_{s,q+1})-\lambda_s)$.
Here $\lambda_s=\log(1+\eta_{q-1})$ and the full original
$\mathcal L_{s,t}$ are exactly ROQ17--21; all masses are retained.

For the already proved original seven-moment receiving equality,
retain the finite domain of PMR1--3 and PMR7:
$k\ge29$ and $k\sqrt{\delta^2+\gamma^2}\le2^{-33}(k+1)^2$.
Write $\mathcal F_{0,k}^{(s)}$ for the full seven-moment outer
expression, and set $e_{k,s}^{0,\pm}=e_{k,s}^{\pm}+c_k^\circ$
with the exact OAR error endpoints and PMR2's complete root-to-moment
error $c_k^\circ$. Substitution, without changing either endpoint, proves
\begin{equation}\tag{RWC12}
\begin{split}
\mathcal F_{0,k}^{(s)}-e_{k,s}^{0,-}-\Phi^+_{k,s}
&\le F_K^{(s)}\\
&\le\mathcal F_{0,k}^{(s)}+e_{k,s}^{0,+}-\Phi^-_{k,s}.
\end{split}
\end{equation}
Outside that PMR domain the original OAR outer-return formula gives
the same interval with $\mathcal F_{\partial,k}^{(s)}$ in place
of $\mathcal F_{0,k}^{(s)}$ and $e_{k,s}^{\pm}$ in place of
$e_{k,s}^{0,\pm}$, throughout OAR's stated domain. RWC1--11
remain valid at every $k\ge9$ specified at the start of this section.
Every stronger previously proved kernel bound can be intersected
with these endpoints. For the arithmetic statement also retain
the original AMT period domain $|u|\ge R_*$, as in OVR's opening
paragraph. In the complete OAR/OVR signed return
$\Psi_k=F_K^{\rm ar}+F_B^{\rm ar}-F_B^{(k)}$, the same literal
substitution gives
\begin{equation}\tag{RWC13}
\begin{split}
\mathcal F_{0,k}^{(k)}+\mathcal H_k+\delta_k^\sigma
 -e_{k,k}^{0,-}-\Phi^+_{k,k}
&\le\Psi_k\\
&\le\mathcal F_{0,k}^{(k)}+\mathcal H_k+\delta_k^\sigma
 +e_{k,k}^{0,+}-\Phi^-_{k,k}.
\end{split}
\end{equation}
The original correlated arithmetic deficit $\delta_k^\sigma$,
boundary return $\mathcal H_k$, and each unequal source error
are unchanged. Thus every computed exterior coefficient in RWC8
has a proved map into both the kernel and signed-return estimates.
The strict finite positivity in RWC10 alone supplies no rate for
$\Phi_{k,s}/(kq)$. The quantitative next calculation is the growth
of these same conditioned matrices or their exterior coefficients
in the fixed actual period, not a replacement by an arbitrary
rank-$j$ positive matrix.

The new lower endpoint strictly strengthens the original first-increment
endpoint, through an exact comparison on the same covariance. ROQ18
telescopes to
\begin{equation}\tag{RWC14}
V_{q-1}=\sum_{a=q-1}^{2q-2}\frac{z_az_a^*}{a_a},\qquad
t_{q-1}=\sum_{a=q-1}^{2q-2}
\frac{z_a^*\Omega_{q-1}^{-1}z_a}{a_a}
\ge\eta_{q-1},\qquad \ell_{k,s}>\lambda_s.
\end{equation}
Here $z_a,a_a,\eta_a$ are precisely ROQ17 at their original
individual cutoffs; the sum retains every later mixed covariance.
Each summand is nonnegative, and the first is exactly
$\eta_{q-1}$. Thus the first-window trace lower bound in RWC11
is at least $\log(1+\eta_{q-1})=\lambda_s$. The second-window
lower bound is strictly positive by RWC5 and RWC7. Adding them
proves the strict inequality. Consequently the original maximum
in $\Phi^-_{k,s}=\max(\ell_{k,s},\lambda_s)$ equals
$\ell_{k,s}$ by this proved comparison; both constituent bounds
remain recorded, with their original definitions and domains.

% END CURRENT RWC

\clearpage
% BEGIN CURRENT OER
\section{The fixed-period determinant and closed coefficients in the original receiver}
\label{sec:OER}

Retain the complete original objects, source orders, observations and
arithmetic maps of OMR1--26. The PCL, ROG, FEX, FGC, ECL, EQR and RDS
proofs are supplied in full. This section composes their exact maps
with the current receiver. The conductor and fixed-period determinant
applications keep $m=1$, $k=4l+1\ge9$, $q=(k+1)^2$,
$q'=(k-7)^2$, $\Delta=16k-48$, $r=8k-16$, $j=8k-32$,
$c=k/2$, $c'=c-4$, and the fixed actual five-orbit period and
its original logarithm branch. Each earlier eventual comparison
retains its explicit finite domain.

\subsection{The existing outer coefficient has an exact endpoint expression}

To distinguish the endpoint values from the conductor moment index
$v$, write $(u_E,v_E,w_E,z_E,c_E)$ for exactly ECL2--3 at
$(\alpha,\beta)=(2,\pi)$. This is the literal dictionary
$u_E=u$, $v_E=v$, $w_E=(u+v)/2$, $z_E=uv$,
$c_E=(v^2-u^2)/4$ in that provider. With $\kappa_E$ denoting its
elliptic modulus, the original endpoint equations and coefficient are
\[
\begin{gathered}
 u_EK(\kappa_E)=2,\qquad v_EE(\kappa_E)=4,\qquad
 \kappa_E^2=1-u_E^2/v_E^2,\qquad 0<u_E<v_E,\\
 \mathfrak C_\partial
 =4\log(4/\pi)+8\log w_E-4\log z_E-2\log c_E,\\
 21.66851180520<16\mathfrak C_\partial<21.66851180521 .
\end{gathered}\tag{OER1}
\]
Here $K,E$ are the original complete elliptic integrals with modulus
convention. ECL6 and ECL14 give
$\ell=8-4\log w_E-2\log c_E$,
$\mathcal L=6\log w_E-2\log z_E-2$ and
$F=-\ell/2-3+\mathcal L$. Substitution in the exact OFV/ECL16
expression $4\log(4/\pi)-10-2F+4\mathcal L$ gives OER1:
the constant terms are exactly $-4+8-4=0$.
Also $F=I_{2,\pi}-6+2\mathcal L$, so this number is exactly
$4\log(4/\pi)+2-2I_{2,\pi}$ in OMR2, MRI4r and BRI6h.
The integer outward certificate ECL19--24 proves its displayed
enclosure; all stored endpoints and its complete code are retained.

Consequently the existing exact relations and their errors become
\[
\begin{gathered}
 \mathcal O_k^\pm=\Delta q
  [4\log(4/\pi)+8\log w_E-4\log z_E-2\log c_E]
                         \ \pm4\Delta\varepsilon_k^{\rm out},\\
 F_K^{(s)}+\Phi_{k,s}=\Delta q\mathfrak C_\partial+R_{k,s},
 \quad b_{k,s}-e_{k,s}^-\le R_{k,s}\le b_{k,s}+e_{k,s}^+,\\
 \Psi_k+\Phi_{k,k}
  =\Delta q\mathfrak C_\partial+\mathcal H_k+\delta_k^\sigma+R_{k,k}.
\end{gathered}\tag{OER2}
\]
These are substitutions into OMR2, MRI4s and BRI6h. The complete
OOQ error, the earlier $O(q\log q)$ OFV error, and both signed OAR
errors retain their original finite definitions. Put
$c_-=21.66851180520/16$, $c_+=21.66851180521/16$.
The original outer value is also enclosed by
$[\Delta qc_- -4\Delta\varepsilon_k^{\rm out},
  \Delta qc_+ +4\Delta\varepsilon_k^{\rm out}]$.
This follows from $\Delta q>0$ and OER1 and retains the exact
coefficient in all identities.

\subsection{The actual period coefficients and the pivot morphism}

PCL1--12 gives the full entire coefficient matrix
$\mathcal R(z)=\sum_{n\ge0}R_nz^n$, with every coefficient
specified by the finite sums in PCL2, and the exact identities
\[
\begin{gathered}
 \mathsf H(u)=e^{\mathfrak a/u}G_u\mathcal R(u^{-1})V^{-1}U,
 \qquad \det\mathcal R(z)=1,\\
 \mathcal R(0)=
 \begin{pmatrix}1&0&-\beta/3&0\\0&1&0&-2\beta/3\\
                 0&0&1&0\\0&0&0&1\end{pmatrix},\\
 L_a(z)=U^{-1}V\mathcal R(z)^{-1}D^{-a}
                         \mathcal R(z)V^{-1}U,\qquad
 A_z(Z,W)=\prod_{a=1}^4Q_0(L_a(z)(ZW,Z,W,1)^{\mathsf T}),\\
 (C_{\rm in}(z))_{(a,b),(i,j)}=a_{i+4-a,j+4-b}(z),
 \qquad C_{\rm in}(-z)=C_{\rm in}(z)^{\mathsf T}.
\end{gathered}\tag{OER3}
\]
All $81$ coefficients are those of this displayed product, with
PCL10's full sums; $U,V,G_u,\mathfrak a,\beta,D$ retain their
original definitions. Conjugating $D^{-a}$ by $\mathsf H$
cancels its scalar factors; $G_u$ commutes with $D^{-a}$,
leaving exactly $L_a$. PCL2's $n=0$ terms give the constant
matrix, including both nonzero off-diagonal entries.
PCL22--24 proves the last identity coefficient by coefficient.
This is the actual matrix family in GIP8 and OMR9; no entry
has been removed in taking its coefficient limit.

For the original OCF leading pivot let $\mathcal P_k$ be its
actual pivot rectangle and $\mathcal E_k$ its complement.
With the original insertion $S_k$ and full reduction $N_k$,
\[
\begin{gathered}
 S_k^{\mathsf T}N_k^{\mathsf T}=I_\Delta,\qquad
 \operatorname{im}N_k^{\mathsf T}=V_A,\\
 J_{\mathcal E}=N_k^{\mathsf T}
                D_{\mathcal P,\mathcal E}V_{\mathcal E},
 \quad \rho_{\mathcal E}=V_{\mathcal E}^{-1}
       D_{\mathcal P,\mathcal E}^{-1}S_k^{\mathsf T}|_{V_A},
 \quad \rho_{\mathcal E}^{-1}=J_{\mathcal E},\\
 \Xi_{\mathcal E}=\Xi_k^{\rm fix}J_{\mathcal E},\quad
 \ker\Xi_{\mathcal E}=\rho_{\mathcal E}(K),\\
 H_{\mathcal E,N}=J_{\mathcal E}^*G_NJ_{\mathcal E},\qquad
 Q_{\Xi,N}^{\rm fix}
   =(\Xi_{\mathcal E}H_{\mathcal E,N}^{-1}\Xi_{\mathcal E}^*)^{-1}.
\end{gathered}\tag{OER4}
\]
The diagonal $D_{\mathcal P,\mathcal E}$ consists of the full
root products $P_{\mathcal P,k}(\omega_{ab})$, and
$V_{\mathcal E}$ is the evaluation Vandermonde in that actual
complementary root set. Distinct roots make both invertible.
Multiplication proves both inverse identities; composition gives
the kernel. Pullback of the unchanged Gamma form and minimum
over the same residual fibre prove the last line. All root factors
and the full metric in PCL19--21 remain present. No central-grid
OOQ bound is applied to the changed root set without that transport.

\subsection{The full low graph has its own finite-rank transport}

At each original endpoint put $L_N=N-q$ and
$M_N^{\rm pol}=L_N+\Delta-v$. Let $\mathbf U_{L_N}$ be the
complete ROG4 positive-shift relation matrix and
$\Pi_{U,L_N}$ its weighted orthogonal projection in
$\|g\|_{\chi',s,c'}=\|\chi'g\|_{s,c'}$, $\chi'=\chi_{k-8}$.
The finite inverse ROG12 gives the exact maps
\[
\begin{gathered}
 Q_{N,s}^{\rm low}g
 =P_{<v}Z_{\chi,N}\mathcal U_{L_N}^{-1}\Pi_{U,L_N}g,\qquad
 \operatorname{rank}Q_{N,s}^{\rm low}\le v,\qquad
 \|Q_{N,s}^{\rm low}\|\le H_{N,s}^{\rm ROG},\\
 \mathcal R_{N,s}(x,g)=(x-Q_{N,s}^{\rm low}g,g),\quad
 B_{N,s}^{\rm str}=B_N^{\rm graph}\mathcal R_{N,s}^{-1},\\
 \ker B_{N,s}^{\rm str}
 =\operatorname{im}\left[
   \begin{pmatrix}0\\\mathbf U_{L_N}\end{pmatrix},
                       \mathcal R_{N,s}W_{K,N}\right].
\end{gathered}\tag{OER5}
\]
$H_{N,s}^{\rm ROG}$ is exactly ROG24's product
$A^+_{N,s,c}(1+R)^q I_{L_N}D_{M_N^{\rm pol},\chi'}
A^-_{N-v,s,c'}$, with all ROG6--24 constants.
At $L_N=-1$ the map and its bound are zero.
The shear sends
$(P_{<v}Z_{\chi,N}h,\mathcal U_{L_N}h)$ to
$(0,\mathcal U_{L_N}h)$; applying it to the full GIP kernel
proves the last line, retaining every common-kernel column.

For the exact product source Gram define $Q_{N,s}^{\rm str}$
by ROG31 and
$\Phi_{k,s}^{\rm str}=\sum_N\sigma_N\log\det Q_{N,s}^{\rm str}$,
with the original signs $(1,1,-1,-1)$. Then
\[
\begin{gathered}
 E_s^{\rm str}
 =2\sum_NE_N^{\rm CEP}
       +2v\sum_N\log(1+H_{N,s}^{\rm ROG}),\\
 |\Phi_{k,s}^{\rm str}-\Phi_{k,s}|
       \le E_s^{\rm str}=o(kq).
\end{gathered}\tag{OER6}
\]
Each singular value $\sigma$ of the low map gives two shear
singular values $(\sqrt{\sigma^2+4}+\sigma)/2$ and its reciprocal.
At most $v$ exceed one. Thus every exterior forward and inverse
logarithmic cost is at most $v\log(1+H_{N,s}^{\rm ROG})$.
Minimum lifts and projection off the entire kernel contract
exterior norms. Squaring gives the determinant bound; adding
the full GIP source-map cost and the four signs proves OER6.
ROG25 gives $\log(1+H_{N,s}^{\rm ROG})
=O_{\rm actual}(q\log(q+1))$; fixed $v\le80$ makes it $o(kq)$.
The full ROQ/FMC32 minimum-form identity identifies the actual
four-endpoint residual with $\Phi_{k,s}$.

\subsection{The fixed columns arise from an exact quotient of that graph}

Put $\alpha_A=v!/\mu_v$, $\mathcal C_A=\alpha_A\mathcal T_A$,
$L=\mathbb C[S]_{<v}\subset E$, $K^+=K+L$ and
$t_0=\dim(K\cap L)$. The explicit coordinate isomorphism
$(\ell,z)\mapsto(\ell,\alpha_Az)$ sends the graph relations to
$(P_{<v}Z_{\chi,N}h,\alpha_A\mathcal U_{L_N}h)$.
Projection to the second coordinate gives the exact sequences
\[
\begin{gathered}
 0\longrightarrow L\longrightarrow V_A
   \longrightarrow V_A/L
   \simeq\mathbb C[S']_{\le N-v-q'}/
                 \alpha_A\mathcal U_{L_N}\mathbb C[S]_{\le N-q}
   \longrightarrow0,\\
 0\longrightarrow K\cap L\longrightarrow K
       \longrightarrow K^+/L\longrightarrow0,\qquad
 \dim(K^+/L)=t=r-t_0 .
\end{gathered}\tag{OER7}
\]
For a zero lower class, subtract its full graph vector from
$(\ell,\alpha_A\mathcal U h)$, leaving
$(\ell-P_{<v}Z_{\chi,N}h,0)$. Such a low vector is a relation
only when $\mathcal U h=0$, which forces $h=0$ by the nonzero
leading coefficient
$\operatorname{lc}(h)\mu_v\binom{q+\deg h}{v}$.
This proves the first kernel and identifies it with the original
$L$. Restriction to $K$ proves the second sequence, as in FGC3--8.

In the complete AKJ basis let $\Xi=(\Xi_{\rm lo},\Xi_{\rm hi})$
and let $Q_{\rm lo}$ have kernel $\operatorname{im}\Xi_{\rm lo}$.
Choose the specified fixed full basis $W_k$ of
$\ker(Q_{\rm lo}\Xi_{\rm hi})$ using actual nonzero minors.
The original low and high columns give
\[
\begin{gathered}
 Z_k=\operatorname{diag}(d_h)W_k,\qquad
 d_h=\frac{(q'+h+v)!}{(q'+h)!},\qquad
 \mathcal C_AP_h=d_h\psi(S')S'^h,\quad\psi=\chi_{k-8},\\
 R_a^A(S')=\frac{\mathcal C_A(\chi S^a)}{\psi}
      =\alpha_A\sum_\nu a_\nu
             \frac{\chi(S'+b_\nu)}{\psi(S')}(S'+b_\nu)^a .
\end{gathered}\tag{OER8}
\]
Solvability of $\Xi_{\rm lo}\ell+\Xi_{\rm hi}z=0$ is exactly
the condition defining $W_k$, so $Z_k$ represents $K^+/L$.
Each relation polynomial has degree $\Delta-v+a$ and nonzero
leading coefficient $(q+a)_v$, while the $Z_k$ columns have
degree below $\Delta-v$. Their quotient classes are independent.

Let $H_{N,s}^{\psi}$ be the full monomial Gram in the original
measure $|\psi(c'+iy)|^2dm_{0,s}(y)$, and $R_N^A$ the first
$N+1-q$ relation columns of OER8. Their four ranks are exactly
$0,1,q,q+1$. The complete minimum matrix is
\[
\begin{split}
 \widehat H_{K,N}^{(s)}
 ={}&Z_k^*H_{N,s}^{\psi}Z_k
 -Z_k^*H_{N,s}^{\psi}R_N^A
   ((R_N^A)^*H_{N,s}^{\psi}R_N^A)^{-1}
       (R_N^A)^*H_{N,s}^{\psi}Z_k,\\
 \widehat F_k^{(s)}
 ={}&\sum_N\sigma_N\log\det\widehat H_{K,N}^{(s)} .
\end{split}\tag{OER9}
\]
Minimizing the norm of $Z_kx+R_N^Ah$ gives
$h=-((R_N^A)^*H_{N,s}^{\psi}R_N^A)^{-1}
(R_N^A)^*H_{N,s}^{\psi}Z_kx$. Substitution proves the full
Schur formula including both cross factors. At the first endpoint
there is no relation variable. Independence proves positivity.
This is exactly FEX29/FGC11 on the native lower metric.
The original graph metric in these coordinates is
$D_{\alpha,N}^{-*}H_N^{\rm actual}D_{\alpha,N}^{-1}$;
the next bounds retain its change to that native metric.

\subsection{Both low terms and all one-sided errors enter the current interval}

Set $e_{N,s}^{\rm FEX}=\log\det\widehat H_{K,N}^{(s)}
 -\log\det H_{K^+/L,N}^{(s)}$. Keep exactly FEX18--20/34:
\[
\begin{gathered}
 n_N=N+1-v,\quad
 m_N=\log\max\left\{1,\frac{v!}{|\mu_v|}
           \sum_\nu|a_\nu|e^{|b_\nu-4|(H_N+2)}\right\},
 \quad H_N=\sum_{a=1}^{N}\frac1a,\\
 \zeta_{N,s}=\log\prod_{d=v}^{N}\frac{w_d}{w_{d-v}},
 \qquad w_d=\frac{d!}{(s/2)_d},\\
 e_{N,s}^-=\zeta_{N,s}-2(n_N-t)m_N,\qquad
 e_{N,s}^+=2tm_N,\qquad
 e_{N,s}^-\le e_{N,s}^{\rm FEX}\le e_{N,s}^+,\\
 D_s^-=e_{q-1,s}^-+e_{q,s}^--e_{2q-1,s}^+-e_{2q,s}^+,\qquad
 D_s^+=e_{q-1,s}^++e_{q,s}^+-e_{2q-1,s}^--e_{2q,s}^- .
\end{gathered}\tag{OER10}
\]
These $e^\pm$ are the FEX flag bounds, distinct from OAR's
$e_{k,s}^\pm$. Here $H_N$ is the harmonic number, distinct
from the kernel Gram. The full flag identity gives
$\mathscr J M^{-2(n_N-t)}\le D_Q\le M^{2t}$,
with the exact squared Jacobian $\mathscr J=e^{\zeta_{N,s}}$
and norm bound $M=e^{m_N}$. This proves the endpoint bounds
and their four signed sums, keeping both original Gamma masses
until their exact Jacobian cancellation.

The two full Schur factorizations of the same $K^+$ give
\[
\begin{gathered}
 F_K^{(s)}=\widehat F_k^{(s)}
       -\sum_N\sigma_N e_{N,s}^{\rm FEX}
       +F_L^{(s)}-F_{K^+/K}^{(s)},\\
 0\le F_L^{(s)}\le v\mathcal L_s,\quad
 0\le F_{K^+/K}^{(s)}\le(v-t_0)\mathcal L_s,\quad
 \mathcal L_s=\mathcal L_{s,q}+\mathcal L_{s,q+1},\\
 f_s^-=\widehat F_k^{(s)}-D_s^+-(v-t_0)\mathcal L_s
 \le F_K^{(s)}\le
 f_s^+=\widehat F_k^{(s)}-D_s^-+v\mathcal L_s .
\end{gathered}\tag{OER11}
\]
In fixed adapted frames the determinant of $K^+$ is both its
$L$ determinant times that of $K^+/L$, and its $K$ determinant
times that of $K^+/K$, including fixed Jacobians. Their four
endpoint constants cancel because the signs sum to zero.
This proves the first identity. The original AKS bound passes
to both subquotients by restriction and minimization, proving
the second line. OER10 then proves the last line.
If the two actual low volumes are evaluated, the stronger
interval retains their exact difference in place of the caps.
Both low terms remain part of the identity in either case.

All existing OMR11--12 intervals remain. The new same-scalar
intersection is
\[
\begin{gathered}
 L_s^{\rm ext}=\max\{L_s^{\rm op},f_s^-,
   \mathcal F_{\partial,k}^{(s)}-\Phi_{k,s}^{\rm str}
                          -E_s^{\rm str}-e_{k,s}^-\},\\
 U_s^{\rm ext}=\min\{U_s^{\rm op},f_s^+,
   \mathcal F_{\partial,k}^{(s)}-\Phi_{k,s}^{\rm str}
                          +E_s^{\rm str}+e_{k,s}^+\},\\
 L_{\rm ar}^{\rm ext}
   =\max\{L_{\rm ar}^{\rm op},L_1^{\rm ext}+a_{K,k}^-\},
 \qquad
 U_{\rm ar}^{\rm ext}
   =\min\{U_{\rm ar}^{\rm op},U_1^{\rm ext}+a_{K,k}^+\},\\
 F_K^{\rm ar}\in[L_{\rm ar}^{\rm ext},U_{\rm ar}^{\rm ext}],
 \qquad
 F_B^{\rm ar}\in[\mathcal B_k^{\rm ar}-U_{\rm ar}^{\rm ext},
                 \mathcal B_k^{\rm ar}-L_{\rm ar}^{\rm ext}].
\end{gathered}\tag{OER12}
\]
The first lines combine OER6, the original MRI4m equality,
and OER11. Every interval bounds the same $F_K^{(s)}$.
The last lines add the unchanged correlated arithmetic allowances
of OMR12 and subtract from the same total. All prior finite
branches remain available outside any added domain. FEX40--41
also permits the existing reverse source intersections with
its full asymmetric $\Omega$ bounds.

Define $L_J^{\rm ext}=L_k^{\rm ext}+\mathcal H_k
+\delta_k^\sigma+C_k^{\rm budget}$ and
$U_J^{\rm ext}=U_k^{\rm ext}+\mathcal H_k
+\delta_k^\sigma+C_k^{\rm budget}$. Then
\[
\begin{gathered}
 L_k^{\rm ext}+\mathcal H_k+\delta_k^\sigma
 \le\Psi_k\le U_k^{\rm ext}+\mathcal H_k+\delta_k^\sigma,\\
 J_k^{\rm action}=2\sum_Nw_N\log\epsilon_N^{\rm ar}
   \in[L_J^{\rm ext},U_J^{\rm ext}],\qquad
 4q\log\mathcal L_{h,k}\le J_k^{\rm action},\\
 \mathcal L_{h,k}\le\min_N\epsilon_N^{\rm ar}
 \le\prod_N(\epsilon_N^{\rm ar})^{w_N/(2q)}
 \le\exp(U_J^{\rm ext}/(4q)),\\
 \frac{\Psi_k}{kq}
 =\frac{\widehat F_k^{(1)}}{kq}-\frac{\mathcal C_\Gamma}{4}+o(1).
\end{gathered}\tag{OER13}
\]
The first line is the exact OMR13 equality with $\mathcal H_k=-T_k$.
The next lines use the same arithmetic budget, exact original
reflection, and same positive class in OMR19--26.
FEX36--38 proves $|F_K^{(s)}-\widehat F_k^{(s)}|
=O_{\rm actual}(q\log q)=o(kq)$; the original source difference
and arithmetic correction are $o(kq)$, and the proved Gamma
return divided by $kq$ tends to $-\mathcal C_\Gamma/4$.
Substitution proves the last line without asserting a limit
for the unevaluated fixed-column determinant.

\subsection{The relation quotient and the earlier finite Gamma centre}

With every original coefficient, put
$E_A(z)=\sum a_{ab}e^{b_{ab}z}$, $D=\partial_z$ and
$\mathscr L_A=\sum a_{ab}e^{b_{ab}z}d_{ab}(D)$.
The exact analytic dual of the denominator above is
\[
\begin{gathered}
 \mathscr L_A\chi'(D)=\chi(D)M_{E_A},\\
 0\longrightarrow\mathscr E_{k-8}
   \xrightarrow{M_{E_A}}\mathscr N_v
   \xrightarrow{\mathfrak F_A}
      (\mathbb C[S']/\mathcal U\mathbb C[S])^\vee
   \longrightarrow0,\\
 f_{\mathfrak F_A(N)}(z)=\chi'(D)(N(z)/E_A(z)).
\end{gathered}\tag{OER14}
\]
$\mathscr N_v$ is the original upper exponential space with its
first $v$ Taylor values zero; the operator order and all spaces
are RDS6--8. RDS5 proves the intertwining by the literal
positive shifts and Leibniz' rule. Its kernel consists precisely
of $E_A\mathscr E_{k-8}$, since the monic constant-coefficient
equation $\chi'(D)F=0$ has exactly that lower exponential
solution space; the full Vandermonde proves its initial-value
uniqueness. Dimension $q-v-q'=\Delta-v$ then proves surjectivity.
Because $\alpha_A\ne0$, the relation subspace
$\alpha_A\mathcal U\mathbb C[S]$ in OER7 is the same subspace.
RDS's finite-degree direct decomposition identifies its dual
at every original cutoff. The full source coefficient identity,
original phases and masses, analytic disk and pole multiplicities
remain RDS8--15; the Schur metric in OER9 is retained.

Finally ECL12 applies to precisely QGQ4's original derivative
$\mathcal L_1=\partial_\alpha\mathcal L(2,\pi)$. Hence
\[
\begin{gathered}
 c_q^\Gamma=\mathcal C_\Gamma^{[q,1]}
   =\frac18\log(w_E^2/z_E)-\frac14\log2,\qquad
 -0.080453400358<c_q^\Gamma<-0.080453400357,\\
 \frac{\mathcal V^{\rm c}_{q,l}-lq\mathcal C_\Gamma}{q}
                  \longrightarrow c_q^\Gamma,\qquad
 \limsup_{l\to\infty}
 \left|\frac{W_k-lq\mathcal C_\Gamma}{q}-c_q^\Gamma\right|
             \le A/\sqrt2,\quad A=M/2+2C>0 .
\end{gathered}\tag{OER15}
\]
This uses exactly the QGQ domain $m=1$, $k=4l+1$,
$q=(4l+2)^2$, $l\ge1$. EQR1--3 proves the connection:
ECL12 gives $\mathcal L_1=\log(w_E^2/z_E)$; the complete
QGQ7 centre has correction $(2\mathcal L_1-4\log2)l^2$
and remaining $O(l)$ terms. Since $l^2/q\to1/16$,
division gives the stated centre limit. The full unequal
finite errors remain
\[
 -p^-_{q,l}-\mathcal U_{q,l}-E_{q,l}
 \le W_k-\mathcal V^{\rm c}_{q,l}
 \le p^+_{q,l}+\mathcal U_{q,l}+T_{q,l}+E_{q,l}.
 \tag{OER16}
\]
All terms except $\mathcal U_{q,l}$ become zero after division
by $q$, while $\mathcal U_{q,l}/q\to A/\sqrt2$.
Adding the finite-centre limit proves OER15 with the same
radius. The original definitions give $C\ge0$ and
$M\ge\log(w_E^2/z_E)>0$, since
$w_E^2/z_E=1+(v_E-u_E)^2/(4u_Ev_E)>1$.
The radius remains strictly positive. The evaluation does not
assert a $q$-scale limit for the actual $W_k$ or change its
finite interval.

\subsection{Full-window residual bounds in these same receivers}
Retain the literal RWC4 matrices for each source $s=1,k$ and
the two low cutoffs $N=q-1,q$. Their full-window positive update,
including the full $q\times q$ conditioning matrix, is
\[
\begin{gathered}
 U_N^+=[u_{N+1},\ldots,u_{N+q}],\quad
 B_N=I_q+(CU_N^+)^*(CR_NC^*)^{-1}(CU_N^+),\\
 Z_N=D[I_q-R_NC^*(CR_NC^*)^{-1}C]U_N^+,\quad
 V_N=Z_NB_N^{-1}Z_N^*>0,\quad
 \Omega_{N+q}=\Omega_N+V_N .
\end{gathered}\tag{OER17}
\]
The original full conductor rows and residual rows here are precisely
ROQ1 and ROQ4; $D$ in this subsection denotes those rows.
RWC2--6 proves the rank and update by the positive original quartet
polynomial, orthogonality of the complete consecutive source block,
and the displayed full covariance expansion. With
$Y_N=\Omega_N^{-1/2}Z_NB_N^{-1/2}$, the exact four-endpoint
residual and its finite bounds are
\[
\begin{gathered}
 \Phi_{k,s}=\sum_{N\in\{q-1,q\}}
 \log\left(\sum_{d=0}^j\|\bigwedge^dY_N\|_{\rm HS}^2\right),\\
 \lambda_s<\ell_{k,s}=\Phi^-_{k,s}\le\Phi_{k,s}
 \le\Phi^+_{k,s},
\end{gathered}\tag{OER18}
\]
where the complete trace, determinant and earlier finite-cap
definitions of $\ell,\Phi^+$ are RWC7--11. RWC8 proves the
exterior equality by all singular products; RWC14 proves the
strict improvement over the original first increment $\lambda_s$.
Thus every exterior coefficient supplies a specified finite
subtraction in the same kernel equation.

On the exact PMR range $k\ge29$,
$k\sqrt{\delta^2+\gamma^2}\le2^{-33}(k+1)^2$, put
\[
\begin{gathered}
 L_s^{\rm win}=\max\{L_s^{\rm ext},
   \mathcal F_{0,k}^{(s)}-e_{k,s}^{0,-}-\Phi^+_{k,s}\},\\
 U_s^{\rm win}=\min\{U_s^{\rm ext},
   \mathcal F_{0,k}^{(s)}+e_{k,s}^{0,+}-\Phi^-_{k,s}\},
 \qquad e_{k,s}^{0,\pm}=e_{k,s}^{\pm}+c_k^\circ,\\
 L_s^{\rm win}\le F_K^{(s)}\le U_s^{\rm win}.
\end{gathered}\tag{OER19}
\]
RWC12 proves that each added interval contains the actual kernel
volume, so taking its intersection with OER12 proves OER19.
Outside that PMR range use the original OAR expression
$\mathcal F_{\partial,k}^{(s)}$ and its original $e_{k,s}^{\pm}$
in precisely the same construction on OAR's domain.
For the arithmetic statements retain $|u|\ge R_*$ as in AMT.
Substitute $L_s^{\rm win},U_s^{\rm win}$ for the source bounds
in the last two lines of OER12, and for source $k$ in OER13.
This gives respectively the same correlated arithmetic interval,
its exact boundary complement, and
\[
\begin{gathered}
 L_k^{\rm win}+\mathcal H_k+\delta_k^\sigma
 \le\Psi_k\le U_k^{\rm win}+\mathcal H_k+\delta_k^\sigma,\\
 J_k^{\rm action}\le U_k^{\rm win}+\mathcal H_k
                  +\delta_k^\sigma+C_k^{\rm budget}.
\end{gathered}\tag{OER20}
\]
The proof is addition of the unchanged scalars to the actual
kernel bounds and then the exact OMR26 identity. The full
budget, its same-class lower bound, both low corrections and
every unequal finite error remain present. No estimate of
the individual $kq$-scale residual coefficient is inferred
from strict finite positivity alone.

% END CURRENT OER


\clearpage
% BEGIN CURRENT RCF
% Exact ideal hull and factor sequence of the full original relation image.
\section{The complete common factor of the original polynomial relations}
\label{sec:RCF}

Retain the full original $\mathcal U$, $d_{ab}$, roots, coefficients,
and degree $\mathfrak m=\Delta-v$ of RDS1--2 and ROG1--6.
Let $\mathcal S_A=\{(a,b):a_{ab}\ne0\}$ be the actual nonempty
coefficient support, and define its four exact integer extrema
\[
 a_-=\min_{\mathcal S_A}a,\quad a_+=\max_{\mathcal S_A}a,\quad
 b_-=\min_{\mathcal S_A}b,\quad b_+=\max_{\mathcal S_A}b,\qquad
 d_a=a_+-a_-,\quad d_b=b_+-b_-.
 \tag{RCF1}
\]
These are cases determined by the original coefficients; none is
assigned a generic support. Write $k'=k-8$ and extend only the
root-index formula to integer pairs by
\[
 \widetilde\beta_{ij}=c'+(2i-k')\delta+i(2j-k')\gamma.
 \tag{RCF2}
\]
The $i$ multiplying the imaginary term is the imaginary unit.
The positive real numbers $\delta,\gamma$ make this indexed map
injective on all integer pairs by comparing real and imaginary parts.

Each original shifted outer product $d_{ab}$ has precisely the
simple roots indexed by
\[
 ([-a,k-a]\times[-b,k-b])\setminus([0,k']\times[0,k']),
 \tag{RCF3}
\]
where all intervals contain their integer points. Indeed
$\beta_{ij}-b_{ab}=\widetilde\beta_{i-a,j-b}$, and the
rectangle removed in CEP41 is exactly the inner rectangle in
these translated indices. This verifies the whole root set,
not just its cardinality.

Every rectangle in RCF3 contains that same inner rectangle.
Their intersection is
$[-a_-,k-a_+]\times[-b_-,k-b_+]$. Consequently the exact
monic greatest common divisor of the active products is
\[
 \begin{split}
 G_A(S')&=
 \prod_{(i,j)\in
   ([-a_-,k-a_+]\times[-b_-,k-b_+])\setminus[0,k']^2}
          (S'-\widetilde\beta_{ij}),\\
 g_A:=\deg G_A&=(k+1-d_a)(k+1-d_b)-q'.
 \end{split}\tag{RCF4}
\]
Distinctness of the roots proves that the intersection product
is their monic greatest common divisor, including every root
with its multiplicity one. Empty products give $G_A=1$.
The rectangle is nonempty and contains the inner square because
$0\le a_-\le a_+\le8$ and $0\le b_-\le b_+\le8$.

This product is also the monic generator of the polynomial
ideal generated by the complete linear relation image:
\[
 (\mathcal U\mathbb C[S])=G_A\mathbb C[S'].
 \tag{RCF5}
\]
Here parentheses mean the ideal generated by that image, whose
underlying linear subspace is still the full original image.
To prove the equality, list the $t=|\mathcal S_A|$ active pairs
as $\alpha_1,\ldots,\alpha_t$ and form the matrix
\[
 V(S')_{n,j}=(S'+b_{\alpha_j})^n,
       \quad 0\le n<t,\quad 1\le j\le t.
 \tag{RCF6}
\]
Its determinant, in the displayed order, is the nonzero
constant $\prod_{i<j}(b_{\alpha_j}-b_{\alpha_i})$.
The shift nodes are distinct by the same real and imaginary
comparison as RCF2. Its adjugate divided by this constant is
therefore an inverse matrix over $\mathbb C[S']$.
The vector of the first $t$ exact relations satisfies
\[
 (\mathcal U1,\mathcal US,\ldots,\mathcal US^{t-1})^T
       =V(S')(a_{\alpha_j}d_{\alpha_j}(S'))_{j=1}^t.
 \tag{RCF7}
\]
Thus these two finite lists generate the same polynomial
ideal. All nonzero $a_{\alpha_j}$ are invertible constants,
so it is the ideal generated by the active $d_{\alpha_j}$.
The Euclidean algorithm for polynomials makes this the ideal
of their monic greatest common divisor RCF4. For completeness,
each polynomial remainder step $f=qg+r$ leaves the two
generated ideals equal; decreasing degree terminates and
identifies the final nonzero remainder, divided by its leading
coefficient, with the monic common divisor. The other relations
are polynomial combinations of the same active $d_{\alpha_j}$,
by their full formula RDS1. This proves RCF5 for the whole image.

In particular the actual linear image is a polynomial ideal
if and only if
\[
                         g_A=\mathfrak m.
 \tag{RCF8}
\]
If it is an ideal, it equals its ideal hull RCF5, whose
codimension is $g_A$. RDS2 gives its actual codimension
$\mathfrak m$, proving necessity. Conversely divide each
active $d_{ab}$ by the full $G_A$, and define the exact
polynomial operator
\[
 \mathcal V_Ah=\sum_{(a,b)\in\mathcal S_A}
         a_{ab}\frac{d_{ab}(S')}{G_A(S')}h(S'+b_{ab}),
       \qquad \mathcal U=M_{G_A}\mathcal V_A.
 \tag{RCF9}
\]
Its degree on a nonzero degree-$j$ input is exactly
$j+\mathfrak m-g_A$ and its leading coefficient is still
$\binom{q+j}{v}\mu_v$ times the input leading coefficient.
This follows from monic exact division of RDS1; also
$g_A\le\mathfrak m$ because $G_A$ divides the nonzero
$\mathcal U1$ of degree $\mathfrak m$. If equality holds,
$\mathcal V_A$ preserves degree with nonzero diagonal
coefficients in every finite monomial frame. Repeated
highest-degree elimination therefore proves that it is
onto all polynomials. RCF9 then proves sufficiency in RCF8.

For every actual support, including the unequal-degree cases,
the precise finite-dimensional factor sequence is
\[
 \begin{gathered}
 0\longrightarrow
 \mathbb C[S']/\mathcal V_A\mathbb C[S]
 \xrightarrow{M_{G_A}}
 \mathbb C[S']/\mathcal U\mathbb C[S]
 \xrightarrow{\operatorname{rem}_{G_A}}
 \mathbb C[S']/(G_A)\longrightarrow0,\\
 \dim\bigl(\mathbb C[S']/\mathcal V_A\mathbb C[S]\bigr)
    =\mathfrak m-g_A
    =(d_a+d_b)(k+1)-d_ad_b-v.
 \end{gathered}\tag{RCF10}
\]
The first map is well defined by RCF9. If $G_Af$ belongs
to $\mathcal U\mathbb C[S]=G_A\mathcal V_A\mathbb C[S]$,
cancellation of the nonzero polynomial proves that $f$
belongs to the first denominator, so that map is injective.
The last map is onto by ordinary monic remainder, and its
kernel consists exactly of the classes represented by
multiples of $G_A$, the image of the first map. The same
degree-elimination proof as RDS2, applied to RCF9, gives its
first dimension; expansion of RCF4 gives the last expression.
If $G_A=1$, the rightmost quotient is zero and these maps
retain the entire original relation quotient on the left.

The source metric transported by the first map is exactly
\[
 \|G_Af\|_{\chi',s,c'}^2
      =\int|f(c'+iy)|^2|\chi'(c'+iy)G_A(c'+iy)|^2dm_s(y).
 \tag{RCF11}
\]
This is direct substitution into the original weighted norm
ROG21. It retains every factor and the original source order
and mass. Thus RCF9--11 are a specified map to the polynomial
ideal supplied by the actual relations, together with the
full remaining quotient. They do not identify its metric
with the CTV outer-divisor metric: the original natural map
to that divisor remains GIP7--9, while RDS8 gives the exact
analytic realization of the complete intermediate quotient.
The original residual map and its common-kernel subspace
remain the transported ROG30 map on that same intermediate
space.

\subsection{Every original finite degree window and its exact metric sequence}

For the original $N=q-1,q,2q-1,2q$ put
$L=N-q$, $M=L+\mathfrak m$ as in ROG2, and set
$r_A=\mathfrak m-g_A\ge0$. At all four endpoints the exact
finite version of RCF10 is
\[
 \begin{gathered}
 0\longrightarrow
 \mathbb C[S']_{\le M-g_A}/\mathcal V_A\mathbb C[S]_{\le L}
 \xrightarrow{M_{G_A}}
 \mathbb C[S']_{\le M}/\mathcal U\mathbb C[S]_{\le L}
 \xrightarrow{\operatorname{rem}_{G_A}}
 \mathbb C[S']_{<g_A}\longrightarrow0,\\
 \dim(\text{first space})=r_A,\quad
 \dim(\text{middle space})=\mathfrak m,\quad
 \dim(\text{last space})=g_A.
 \end{gathered}\tag{RCF12}
\]
Here the first numerator is zero when $M-g_A=-1$, and every
denominator is zero at $L=-1$. The degree identity in RCF9
puts its image in precisely the first numerator. The maps
and cancellations used in RCF10 preserve these degree bounds,
proving exactness. Moreover $M\ge\mathfrak m-1\ge g_A-1$,
so all last-space remainder classes have representatives in
the displayed middle numerator. The first numerator dimension
is $L+r_A+1$, while its denominator dimension is $L+1$
when $L\ge0$; at $L=-1$ it is $r_A$ with zero denominator.
These counts prove the dimensions, including the empty cases.

Use the quotient norm of ROG21 on the middle space, the
quotient norm of RCF11 on the first space, and the attained
minimum quotient norm of the same ROG21 polynomial source
under $\operatorname{rem}_{G_A}$ on the last space. These
are exactly a Hilbert restriction and its quotient. Indeed,
every lift of a middle class $[G_Af]$ differs from it by
$\mathcal Uh=G_A\mathcal V_Ah$; thus all its lifts remain
multiples of $G_A$, and its minimum is precisely the first
space's stated norm. Quotienting further by that first space
allows every change by a multiple of $G_A$ within the finite
window, exactly the last remainder norm. No polynomial
division map was presumed to be orthogonal in this argument.

Let $H_{U,N}^{(s)}$, $H_{V,N}^{(s)}$, and $Q_{G,N}^{(s)}$
be these three Grams in their canonical remainder coefficient
frames $1,S',\ldots$ of dimensions $\mathfrak m,r_A,g_A$.
Those remainder frames exist by the degree elimination proofs
RDS2 and RCF9. Then the exact determinant identity is
\[
 \det H_{U,N}^{(s)}
             =\det H_{V,N}^{(s)}\det Q_{G,N}^{(s)}.
 \tag{RCF13}
\]
To prove it, use in the middle quotient the fixed frame
$[G_A,G_AS',\ldots,G_AS'^{r_A-1},1,S',\ldots,S'^{g_A-1}]$.
Reordering it by increasing leading degree gives a triangular
monic coefficient matrix. Its original order therefore has
determinant $(-1)^{g_Ar_A}$ and modulus one, retaining every
lower coefficient of $G_A$ in the matrix. The first Gram
block is $H_{V,N}^{(s)}$; its Schur complement is the
attained quotient Gram $Q_{G,N}^{(s)}$ by the minimum-lift
calculation. Block Gaussian elimination proves RCF13,
including zero-dimensional determinants equal to one.

The four-return identity thus has the unchanged endpoint signs:
\[
 F_U^{(s)}=F_V^{(s)}+F_G^{(s)},\qquad
 F_G^{(s)}=\sum_{N=q-1,q,2q-1,2q}\sigma_N\log\det Q_{G,N}^{(s)},
 \quad\sigma=(1,1,-1,-1).
 \tag{RCF14}
\]
The other two returns are the same sums of their Grams in
RCF13. The full source orders are $s=1,k$. Thus RCF14
splits an actual original relation quotient metric; the
remaining factor has the explicit weighted source RCF11
and the complete operator RCF9.

There is also a fixed original primary map, before any metric
choice or shear:
\[
 \rho_G:V_A\twoheadrightarrow\mathbb C[S']_{<g_A},\qquad
 [p]\longmapsto\operatorname{rem}_{G_A}
                   \left(\frac{\mathcal T_Ap}{\chi'}\right).
 \tag{RCF15}
\]
Membership in $V_A$ makes the quotient polynomial by CEP5.
A change of representative by $\chi h$ changes it by
$\mathcal Uh$, divisible by $G_A$, proving independence.
The map is onto: for a prescribed polynomial $f$ of degree
below $g_A$, the source surjectivity CEP6 gives a preimage
of $\chi'f$ of degree at most $q'+g_A-1+v\le q-1$.
This preimage belongs to $V_A$ and has the stated remainder.
At $g_A=0$ the target is zero and surjectivity has its usual
zero-dimensional meaning. The map is consequently the same
one at all original endpoints and in both original source norms.

% END CURRENT RCF

\clearpage
% BEGIN CURRENT RCM
% Actual common-factor quotient at the original four source windows.
\section{Original Gamma control on the actual common-factor quotient}
\label{sec:RCM}

Retain RCF1--15, all fixed original period coefficients and their
actual support, and the complete original sources $s=1,k$.
Write $g=g_A$, $\mathfrak m=\Delta-v$, $M=N-q+\mathfrak m$.
If $g=0$, the common-factor target and every quotient of it
are zero, their determinants are one, and all conclusions below
hold with zero return. In the proof of its positive-dimensional
operator inequalities suppose $g\ge1$; this includes every
remaining actual support case.

\subsection{The exact original divisor, windows, phases and source}

Use the unchanged real source coordinate $S'=c'+iy$ and define
\[
 G_y(y)=i^{-g}G_A(c'+iy),\qquad
 \chi'(c'+iy)=i^{q'}P_{k-8}(y),\qquad
 R_G=(k+8)\sqrt{\delta^2+\gamma^2}.
 \tag{RCM1}
\]
These are monic polynomials with exactly the displayed phases.
By RCF2--4 the roots of $G_y$ are
$(2j-k')\gamma-i(2i-k')\delta$ in its stated integer rectangle.
Every $|2i-k'|,|2j-k'|$ there is at most $k+8$; hence every
root has modulus at most $R_G$. They are distinct. No root or
coefficient is set to a surrogate value in RCM1.

The target coefficients of $1,S',\ldots,S'^{g-1}$ and of
$1,y,\ldots,y^{g-1}$ are related by the fixed invertible matrix
\[
 T_{a b}=\begin{cases}\binom ba(c')^{b-a}i^a,&a\le b,\\0,&a>b,
              \end{cases}\qquad 0\le a,b<g,
 \quad \det T=i^{g(g-1)/2}.
 \tag{RCM2}
\]
This follows by expanding $(c'+iy)^b$. In particular $|\det T|=1$,
while the full matrix, including its centre, is retained in every
operator congruence. Let $Q_{G,M}^{y,(s)}$ be the minimum remainder
Gram by $G_y$ on polynomials of degree at most $M$, with source
\[
 d\zeta_{k',s}(y)=|P_{k-8}(y)|^2dm_s(y).
 \tag{RCM3}
\]
The complete phases in RCM1 have modulus one. Substitution proves
$Q_{G,N}^{(s)}=T^*Q_{G,M}^{y,(s)}T$ for the RCF12--14 Gram.
The four actual $M$ values are
\[
 \mathfrak m-1,\quad\mathfrak m,\quad
 q+\mathfrak m-1,\quad q+\mathfrak m.
 \tag{RCM4}
\]
They include all additional columns beyond the divisor's degree:
$g\le\mathfrak m\le\Delta$ and $\mathfrak m\ge16$.

\subsection{Complete finite bounds at the low source windows}

For the comparison measure $|y|^{2Q}e^{-\eta|y|}dy$ use the
full positive source rate $\eta>0$ and put
\[
 \mathfrak g(Q,\eta)=2Q\log(2Q/\eta)-2Q,\qquad y_0=2Q/\eta.
 \tag{RCM5}
\]
OOQ26--27 proves the following bounds on the entire polynomial
coefficient space of degree at most $\Delta$:
\[
 \frac{e^{\mathfrak g(Q,\eta)-\eta}}
 { (\Delta+1)^2(2\Delta+1)64^\Delta(1+y_0)^{2\Delta}}I
 \preceq H_{\Delta}^{\exp}
 \preceq
 \left(\sum_{j=0}^{\Delta}\frac{2(2Q+2j)!}{\eta^{2Q+2j+1}}\right)I.
 \tag{RCM6}
\]
Here $Q$ will be the original integer $q'+(s-1)/4$. These
estimates retain the positive-interval integration bound and the
full two-half-line moment trace from OOQ, respectively.

For every monic degree-$g$ polynomial whose roots have modulus
at most $R_G$, the ordinary remainder map on degree at most
$\Delta$ has coefficient operator norm at most
\[
 C_{\rm lo}=1+(1+R_G)^\Delta(\Delta+1)
                         \max\{1,\Delta R_G\}^{\Delta}.
 \tag{RCM7}
\]
To prove this for the actual $G_y$, its coefficient absolute sum
is at most $(1+R_G)^g$. Polynomial division uses the first
$\Delta-g+1$ complete homogeneous coefficients of its inverse
at infinity. Their sum is at most
$\sum_{a=0}^{\Delta-g}\binom{g+a-1}{a}R_G^a
 \le(\Delta+1)\max\{1,\Delta R_G\}^{\Delta}$.
The coefficient convolution argument in ROG19--20 bounds the
quotient operator by this sum. The exact remainder is the input
minus its quotient multiplied by $G_y$, which proves RCM7.
This proof does not truncate the actual root polynomial.

At either low window in RCM4, taking the minimum over all
lifts of a remainder $z$ gives a lower norm bound from RCM6
and $\|p\|_{\rm coef}\ge\|z\|_2/C_{\rm lo}$.
The canonical degree-less-than-$g$ representative is an allowed
lift and gives the upper bound from RCM6. Thus define
\[
 \begin{split}
 E_{\rm lo}^G(Q,\eta)=\max\bigg\{0,\ &
 \log\left(\sum_{j=0}^{\Delta}
          \frac{2(2Q+2j)!}{\eta^{2Q+2j+1}}\right)-\mathfrak g(Q,\eta),\\
 &\eta+2\log(\Delta+1)+\log(2\Delta+1)+\Delta\log64
       +2\Delta\log(1+y_0)+2\log C_{\rm lo}\bigg\}.
 \end{split}\tag{RCM8}
\]
The full low remainder Gram obeys
\[
 e^{\mathfrak g(Q,\eta)-E_{\rm lo}^G}I_g
 \preceq Q_{G,M}^{\exp}(Q,\eta)
 \preceq e^{\mathfrak g(Q,\eta)+E_{\rm lo}^G}I_g.
 \tag{RCM9}
\]
The extra low columns $g,\ldots,M$ have all been included
in the lower minimization. The factorial bounds proved after
OOQ29 and $\log C_{\rm lo}=O_{\delta,\gamma}(\Delta\log q)$
give $E_{\rm lo}^G=O_{\delta,\gamma}(k\log q)$ in the
source and rate ranges specified below.

\subsection{The high windows with the actual analytic remainder and lift}

Retain the equilibrium constants of OOQ3--6 on its entire
specified compact parameter rectangle. Set
\[
 n=q/2,\quad\rho=n^{-1/8},\quad r_*=R_G/n,\quad
 C_{\rm rem}=\rho\Delta(1+r_*)^\Delta(2/\rho)^\Delta.
 \tag{RCM10}
\]
Use the original OSP23 source-comparison range, $k\ge257$,
and the explicit eventual conditions
\[
 n\ge2\Delta+1,\quad n\ge2/b_*,\quad r_*\le\rho/2,
 \quad\rho\le\min\{1/2,\sqrt{a_*/2}\}.
 \tag{RCM11}
\]
They hold for all sufficiently large original $k$ at the fixed
quartet, since $R_G=O(k)$ and $\Delta=O(k)$ while $n\asymp k^2$.

Apply the exact Cauchy remainder OOQ14 to
$\widehat G(z)=n^{-g}G_y(nz)$. Its $g$ roots are distinct
and lie inside $|z|\le r_*$. Its coefficient norm bound is
$\rho g(1+r_*)^g(2/\rho)^g$, at most $C_{\rm rem}$ because
$1\le g\le\Delta$ and $2(1+r_*)/\rho>1$. The Cauchy
formula interpolates the actual roots and is exactly the
remainder by $\widehat G$, as proved in OOQ14--15.

For either high $M$ in RCM4, the parity degrees
$m_e=\lfloor M/2\rfloor$, $m_o=\lfloor(M-1)/2\rfloor$
belong to $[n,n+\Delta]$. The exact half-line parameters are
$((Q-1/2)/m_e,n\eta/m_e)$ and
$((Q+1/2)/m_o,n\eta/m_o)$. They lie in OOQ's compact
rectangle for the same bounds as OOQ17, since
$16\le\mathfrak m\le\Delta$. The two-branch Jacobian,
weighted polynomial evaluation bound, and smooth variation
estimate are therefore precisely OOQ18--20, with the upper
degree bound $n+\Delta$. These steps give the same lower
quotient bound with $C_{\rm rem}$ from RCM10.

For completeness the upper lift is the literal polynomial
\[
 a_{\varphi}=\operatorname{rem}_{\widehat G}
                   \left(\frac{\varphi}{r_n(\cdot^2)}\right),
 \qquad p_{\varphi}(nx)=r_n(x^2)a_{\varphi}(x),
 \tag{RCM12}
\]
where $\varphi$ denotes the prescribed
degree-less-than-$g$ remainder polynomial and $r_n$ is the
unchanged quantile polynomial OOQ9. Its denominator is nonzero
on the integration disk because $\rho^2<a_*$. The lift has
degree at most $q+g-1\le q+\mathfrak m-1\le M$ and agrees
at every actual root of $\widehat G$ with the prescribed
remainder. Thus it is an allowed original high-window lift.
Its coefficient and integral estimates are OOQ22--23 with
$g\le\Delta$ and the tail power bounded by $2(\Delta-1)$.

To retain all finite errors, put
\[
 \begin{split}
 E^-={}&2\log(C_{\rm rem}\sqrt2 C_{\rm BM}(n+\Delta)^2)
           +2C_W(\Delta+2)+4(n+\Delta)\rho^2/a_*,\\
 E^+={}&2\log C_{\rm rem}+\log\Delta+4n\rho^2/a_*
       +2C_{\rm quant}\sqrt n+2C_0\\
       &{}+2C_{\rm tail}(\Delta-1)+\log T_*,\\
 E_{\rm hi}^G={}&\max\{0,E^-,E^+\}+2(\Delta-1)\log n,\\
 c_{\rm hi}^{\exp}(Q,\eta)
       ={}&(2Q+1)\log n-2nW((Q-1/2)/n,\eta).
 \end{split}\tag{RCM13}
\]
The exact original-versus-scaled remainder matrix is
$\operatorname{diag}(1,n,\ldots,n^{g-1})$; its squared
norm is at most $n^{2(\Delta-1)}$, the last term in RCM13.
The preceding bounds prove at both high windows
\[
 e^{c_{\rm hi}^{\exp}(Q,\eta)-E_{\rm hi}^G}I_g
 \preceq Q_{G,M}^{\exp}(Q,\eta)
 \preceq e^{c_{\rm hi}^{\exp}(Q,\eta)+E_{\rm hi}^G}I_g,
 \quad E_{\rm hi}^G=O_{\delta,\gamma}(q^{3/4}+k\log q).
 \tag{RCM14}
\]
Every imported estimate here has the same actual source,
parity and centre scalings as OOQ. The only changes in its
remainder proof are the specified divisor $G_y$, its proved
root bound, and the degree allowances verified above.

\subsection{Both full Gamma sources and the four-return coefficient}

Retain the exact source constants of OOQ30--31:
\[
 \vartheta=\pi/2,\quad h=1/q,\quad\ell_s=(s-1)/4,\quad Q_s=q'+\ell_s,
 \quad\beta_s=\frac{(2\pi)^{s/2}}{\sqrt2\Gamma(s/2)},
\]
\[
 L_h=\frac{\pi}{\Gamma(3/4)^2}(\sin h)^{3/2},\quad
 U_h=\frac{\Gamma(1/4)^2}{2\pi}(\sin h)^{-1/2},
 \tag{RCM15}
\]
and the complete root/source errors $\mathcal E_k,r_{\ell_s}$
from OSP24 and OFV6, with $r_0=0$. The full quadratic-form
comparison OOQ31 applies to every polynomial used here,
since $M\le q+\Delta\le2q$ on RCM11's range. Taking minima
over the same actual $G_y$ fibres preserves those two
inequalities, including all source masses.

Define the two unchanged leading centres
\[
 c_{\rm lo}=\mathfrak g(q,\vartheta),\qquad
 c_{\rm hi}=2q\log(q/2)-qW(2,\vartheta),
 \tag{RCM16}
\]
and a single finite nonnegative error
\[
 \begin{split}
 \varepsilon_k^G={}&\mathcal E_k+\max_{s=1,k}r_{\ell_s}
                  +\max\{|\log L_h|,|\log U_h|\}\\
 &+\max_{\substack{s\in\{1,k\}\\\eta\in\{\vartheta-h,\vartheta+h\}}}
 \bigl\{E_{\rm lo}^G(Q_s,\eta)+E_{\rm hi}^G
       +|\mathfrak g(Q_s,\eta)-c_{\rm lo}|
       +|c_{\rm hi}^{\exp}(Q_s,\eta)-c_{\rm hi}|\bigr\}.
 \end{split}\tag{RCM17}
\]
All its constants are supplied above or by the complete
original source and equilibrium proofs cited there. The
resulting operator inequalities, in the original $S'$ frame,
are
\[
 \beta_se^{c_\bullet-\varepsilon_k^G}T^*T
 \preceq Q_{G,N}^{(s)}
 \preceq\beta_se^{c_\bullet+\varepsilon_k^G}T^*T,
 \tag{RCM18}
\]
where $c_\bullet=c_{\rm lo}$ at $N=q-1,q$ and
$c_\bullet=c_{\rm hi}$ at $N=2q-1,2q$.
Moreover
\[
 \varepsilon_k^G=O_{\delta,\gamma}(q^{3/4}+k\log q)=o(q).
 \tag{RCM19}
\]
Indeed RCM8 and RCM14 control the new errors. The source
errors have the orders proved in OOQ35. Also
$Q_s-q=-\Delta+\ell_s=O(k)$, the rates differ by $1/q$,
and $W$ with its first derivatives is bounded on the same
compact rectangle. The mean-value estimates in OOQ33--35
give $O(k\log q)$ for both scalar discrepancies in RCM17.
This proves RCM19 without a support-dependent lower bound.

OOQ36 gives the exact difference
$c_{\rm lo}-c_{\rm hi}=q\mathfrak C_\partial/2$ with its
unchanged equilibrium constant. Taking four determinants
in RCM18, and retaining the signs and fixed frame of RCF14,
therefore proves
\[
 |F_G^{(s)}-gq\mathfrak C_\partial|
               \le4g\varepsilon_k^G=o(kq).
 \tag{RCM20}
\]
The source factor $\beta_s^g$ and $\det(T^*T)$ cancel with
coefficients $1+1-1-1$. All individual Grams retain them.
The exact relation return consequently obeys the proved
partial evaluation
\[
 F_U^{(s)}=F_V^{(s)}+gq\mathfrak C_\partial+\epsilon_G^{(s)},
       \qquad |\epsilon_G^{(s)}|\le4g\varepsilon_k^G,
 \tag{RCM21}
\]
where the remaining $F_V$ is the complete smaller operator
and weighted source in RCF9--14, of rank $\mathfrak m-g$.

\subsection{The resulting original residual quotient and its retained complement}

Let $\Xi=\Xi_k^{\rm fix}:V_A\twoheadrightarrow B_{\Xi}$
be the same original residual observation, with
$B_{\Xi}=\mathbb C^j$, $j=8k-32$, and kernel $K$.
Use the fixed onto map $\rho_G$ proved in RCF15, and put
\[
 V_0=\ker\rho_G,\quad W=\Xi(V_0),\quad
 t_G=\dim(B_{\Xi}/W)=g-\dim\rho_G(K).
 \tag{RCM22}
\]
There is the precise onto map
\[
 A_G:\mathbb C[S']_{<g}\twoheadrightarrow B_{\Xi}/W,
       \quad \rho_Gx\longmapsto[\Xi x],\qquad
 \ker A_G=\rho_G(K).
 \tag{RCM23}
\]
Two lifts differ by $V_0$, so the map is well defined.
It is onto because $\Xi$ is onto. If $\Xi x\in W$,
subtract an element of $V_0$ with that same image; the
remaining element belongs to $K$ and has the same $\rho_G$
value. This proves the kernel equality and RCM22's dimension.
All these maps are independent of source order and cutoff.

In the graph source $X_N$ of GIP1, the composite target map
is exactly $(x,f)\mapsto A_G\operatorname{rem}_{G_A}f$.
It ignores the low coordinate and kills every full graph
relation, since $\mathcal Uh$ is divisible by $G_A$.
Minimizing its product norm first over the free low coordinate
and then over polynomial remainder fibres proves that its
quotient Gram is exactly
\[
 Q_{A_G,G,N}^{(s)}
       =[A_G(Q_{G,N}^{(s)})^{-1}A_G^*]^{-1}.
 \tag{RCM24}
\]
Thus no equality of different relation subspaces enters this
identification; only the proved composite map and its full
kernel are used.

Let $\overline Q_{\Xi,N}^{(s)}$ be the original metric on
$B_{\Xi}/W$ induced by the original residual quotient.
The actual source map $S_N$ of CEP28 induces the identity
of this target with RCM24. Its minimum-lift exterior argument
CEP34--38, at the actual rank $t_G$, gives
\[
 |\log\det\overline Q_{\Xi,N}^{(s)}
                  -\log\det Q_{A_G,G,N}^{(s)}|
                          \le2E_N^{\rm CEP}.
 \tag{RCM25}
\]
The source preimage is precisely $\mathcal W_N$ of CEP31.
On its $S_N$ image it is $\mathbb C^v\oplus\chi'\mathbb C[S']_{\le M}$;
the composite map is the one just minimized in RCM24.
Orthogonal projection off the entire composite kernel is
an exterior contraction, for both $S_N$ and its inverse.
This verifies all spaces in the bound and retains the
original source map determinant from CEP29.

Taking the same fixed quotient of RCM18 replaces $T^*T$
by the full fixed reference
$[A_G(T^*T)^{-1}A_G^*]^{-1}$, through minimum on identical
fibres. Its determinant and the original mass cancel only
in the four-sign return. Consequently
\[
 \left|\sum_N\sigma_N\log\det\overline Q_{\Xi,N}^{(s)}
                  -t_Gq\mathfrak C_\partial\right|
      \le4t_G\varepsilon_k^G+2\sum_NE_N^{\rm CEP}=o(kq).
 \tag{RCM26}
\]
No coefficient conditioning of the fixed quotient $A_G$ has
been assumed; its complete reference form is retained and
cancels in the displayed return.

Finally choose any full fixed original frame $I_W$ for $W$
and put
\[
 F_W^{(s)}=\sum_N\sigma_N\log\det(I_W^*Q_{\Xi,N}^{\rm fix}I_W).
\]
Block elimination in a fixed frame of $B_{\Xi}$ extending
$I_W$ gives its restriction determinant times the quotient
determinant. The full fixed frame factor cancels in the
four signs, exactly as in CEP26. Hence the original residual
return has the established decomposition
\[
 F_{\Xi,k}^{(s)}=F_W^{(s)}+t_Gq\mathfrak C_\partial+\epsilon_\Xi^{(s)},
 \quad
 |\epsilon_\Xi^{(s)}|\le4t_G\varepsilon_k^G+2\sum_NE_N^{\rm CEP}.
 \tag{RCM27}
\]
The remaining rank is $\dim W=j-t_G$. Its original metric
and embedding have been specified, not replaced by a scalar
or by a source norm. Also $F_W^{(s)}\ge0$: the original
source polynomial spaces are nested, with the same remainder
and residual target, so minimum quotient metrics decrease
as $N$ increases. Restriction to the same $W$ preserves this
positive-form order, and the two low-versus-high determinant
comparisons are nonnegative. Therefore RCM27 also gives the
actual lower estimate
$F_{\Xi,k}^{(s)}\ge t_Gq\mathfrak C_\partial
 -4t_G\varepsilon_k^G-2\sum_NE_N^{\rm CEP}$.
The value of $t_G$ uses the exact active support and original
map $\rho_G|_K$; no generic nonvanishing or zero-gcd case has
been substituted for that calculation.

% END CURRENT RCM

\clearpage
% BEGIN CURRENT DNE
% Complete original Gamma dual norm; every extension and quotient frame retained.
\section{The original weighted relation metric as a minimum dual extension}
\label{sec:DNE}

Retain the original simple quartet, fixed five-orbit period and
logarithm branch, all coefficients and shifts of ROG and RDS, and
the same two source orders $s\in\{1,k\}$. Put
\[
 \begin{gathered}
 q=(k+1)^2,\quad q'=(k-7)^2,\quad
 \Delta=q-q',\quad \mathfrak m=\Delta-v,\quad c'=k/2-4,\\
 N\in\{q-1,q,2q-1,2q\},\quad L=N-q,\quad
 d=N-v,\quad M=N-v-q'=L+\mathfrak m .
 \end{gathered}\tag{DNE1}
\]
Here $0\le v\le80$, $\mu_v\ne0$, and $\chi'$ is the complete
lower-grid polynomial $\chi_{k-8}$, never a derivative.
All spaces of degree at most $-1$ are zero. The source norm is
the original norm at $S'=c'+iy$, with mass $M_s=(2\pi)^{s/2}$.

\subsection{A fixed original exponential quotient frame}

We specify a finite frame using only the actual original matrices.
Let $\mathbb V$ be the upper-root Vandermonde with rows
$a=0,\ldots,q-1$ and entries
$\mathbb V_{a,(i,j)}=\beta_{ij}^{a}$.
Let $P_v:\mathbb C^{q-v}\to\mathbb C^q$ insert zero in the
first $v$ positions, and let $C$ be the exact primary conductor
of CEP4, with $q'$ rows and $q$ columns. Define
\[
 W_v=\mathbb V^{-1}P_v,\qquad
 J_v=P_v^{\mathsf T}\mathbb V C^{\mathsf T},\qquad
 W_vJ_v=C^{\mathsf T}.
 \tag{DNE2}
\]
The Vandermonde is invertible by the original distinct roots.
The first $v$ rows of $\mathbb V C^{\mathsf T}$ vanish because
the corresponding exponential columns are
$E_Ae^{\beta'_{ab}z}$, whose zero order is at least $v$.
This proves the last equality. Consequently $W_v$ is an
isomorphism onto the coefficient vectors of $\mathscr N_v$,
and $J_v$ has rank $q'$, since $C^{\mathsf T}$ is injective.

Choose the first $q'$ row positions $\mathcal I$ in a fixed
ordering of all subsets for which $\det (J_v)_{\mathcal I}\ne0$,
and let $\mathcal J$ be their complement. Such a subset exists
by the proved rank. Let $E_{\mathcal J}$ insert these
$\mathfrak m$ complementary coordinates and put
\[
 Z_{\exp}=W_vE_{\mathcal J},\qquad
 \mathcal N_a(z)=
   \sum_{i,j}(Z_{\exp})_{(i,j),a}e^{\beta_{ij}z},
 \qquad F_a(z)=\frac{\mathcal N_a(z)}{E_A(z)}
       \quad(1\le a\le\mathfrak m).
 \tag{DNE3}
\]
The columns of $[J_v,E_{\mathcal J}]$ are independent:
restriction of a relation to $\mathcal I$ first kills its
$J_v$ coefficients, and then its other coefficients vanish.
Thus the $\mathcal N_a$ give a basis of the exact quotient
$\mathscr N_v/E_A\mathscr E_{k-8}$. Every root value and
actual minor defining this frame is retained. The frame is
fixed as the endpoint $N$ and source order $s$ change.

For an analytic germ $F$ at zero define the complex-linear
functional $\Lambda_F$ by
$\Lambda_F(S'^n)=F^{(n)}(0)$. Let $\ell_a$ be the
functional from RDS8. Then
\[
 \ell_a(g)=\Lambda_{F_a}(\chi'g),\qquad
 f_{\ell_a}=\chi'(D)F_a,\qquad
 A_{\ell,ba}=\ell_a(S'^b)\quad(0\le b<\mathfrak m).
 \tag{DNE4}
\]
Indeed multiplication by $\chi'$ in the functional is precisely
the constant-coefficient operator $\chi'(D)$ on its generating
series. RDS8 proves that these functionals form the entire dual
of the relation quotient. RDS2 identifies that quotient with
the increasing monomials below degree $\mathfrak m$, so
$A_\ell$ is invertible. This supplies its existence without
assuming a nonzero unevaluated minor.

\subsection{The exact original source extension problem}

On $\mathbb C[S']_{\le M}$ use the original weighted norm
$\|g\|_{\chi',s,c'}=\|\chi'g\|_{s,c'}$. Let $G_M$ be its full
monomial Gram, and let $\mathbf U_L$ be the complete ROG4
relation matrix. The quotient
\[
 Y_N=\mathbb C[S']_{\le M}/\operatorname{im}\mathcal U_L
 \simeq\mathbb C[S']_{<\mathfrak m}
\]
has Gram, in its fixed low-monomial representative frame,
\[
 H_{Y,N}^{(s)}
 =E_{<\mathfrak m}^*G_ME_{<\mathfrak m}
 -E_{<\mathfrak m}^*G_M\mathbf U_L
  (\mathbf U_L^*G_M\mathbf U_L)^{-1}
             \mathbf U_L^*G_ME_{<\mathfrak m}.
 \tag{DNE5}
\]
The second term is absent when $L=-1$. Distinct leading
degrees of the columns of $\mathcal U_L$ show that the low
representatives and all relation columns form a basis.
Positivity of $G_M$ and minimization over the full relation
columns therefore prove DNE5 and its strict positivity.

For $x\in\mathbb C^{\mathfrak m}$ write
$\mathcal N_x=\sum_a x_a\mathcal N_a$ and
$F_x=\mathcal N_x/E_A$. Its functional $\ell_x$ on $Y_N$
has the exact norm
\[
 \|\ell_x\|_{Y_N^\vee}^{\,2}
 =\min_{\eta\in\mathbb C^{q'}}
   \left\|\Lambda_{F_x}
          +\sum_{a,b}\eta_{ab}\operatorname{ev}_{\beta'_{ab}}
   \right\|_{(\mathbb C[S']_{\le d},\,\|\cdot\|_{s,c'})^\vee}^{\,2}.
 \tag{DNE6}
\]
Here the dual consists of complex-linear functionals.
To prove the equality, multiply by $\chi'$ to identify the
weighted polynomial space isometrically with
$I_N=\chi'\mathbb C[S']_{\le M}\subset\mathbb C[S']_{\le d}$.
The restriction of $\Lambda_{F_x}$ to $I_N$ is exactly
$g\mapsto\ell_x(g)$ by DNE4, and it annihilates
$\chi'\operatorname{im}\mathcal U_L$ by RDS5.

A functional annihilating a subspace has the same norm on
the Hilbert quotient: the minimum lift of each quotient
class has the same functional value, and every other lift
has larger norm. A functional on a Hilbert subspace has a
unique extension of minimum norm to the ambient space,
namely its composition with the orthogonal projection onto
that subspace. Every other extension adds a functional on
the perpendicular complement; their squared dual norms add.
These two facts reduce its quotient dual norm to the
minimum over all ambient extensions of its restriction.

The difference of two extensions annihilates $I_N$.
That annihilator has dimension
$(d+1)-(M+1)=q'$. Evaluation at each original lower root
annihilates $I_N$, and these $q'$ evaluation functionals are
independent on $\mathbb C[S']_{\le d}$, by the first $q'$
monomials and the distinct-root Vandermonde. They therefore
constitute its whole annihilator. This proves DNE6 and
retains every permitted extension.

Let
\[
 \phi_{n,s,c'}(S')
 =\frac{p_n((S'-c')/i)}{\sqrt{h_n}},\qquad
 h_n=M_sn!(s/2)_n,\qquad 0\le n\le d
\]
be the unchanged original orthonormal source frame. Form the
complete finite matrices
\[
 \begin{gathered}
 T_{N,s;na}=\Lambda_{F_a}(\phi_{n,s,c'}),\qquad
 V_{N,s;n,(i,j)}=\phi_{n,s,c'}(\beta'_{ij}),\\
 \mathcal D_N^{(s)}
 =T_{N,s}^*
  \left[I-V_{N,s}(V_{N,s}^*V_{N,s})^{-1}V_{N,s}^*\right]T_{N,s}.
 \end{gathered}
 \tag{DNE7}
\]
The matrix $V_{N,s}$ has rank $q'$ by the same Vandermonde
argument, since the displayed orthonormal polynomials span
all degrees at most $d$. In an orthonormal basis the squared
norm of a complex-linear functional is the sum of the
squared moduli of its values. Thus the minimum in DNE6 is
$\min_\eta\|T_{N,s}x+V_{N,s}\eta\|_2^2$.
The minimizing vector is
$-(V_{N,s}^*V_{N,s})^{-1}V_{N,s}^*T_{N,s}x$.
Substitution proves DNE7, including both mixed factors,
and gives $\|\ell_x\|^2=x^*\mathcal D_N^{(s)}x$.
Positivity follows because the $\ell_a$ are an independent
dual frame of $Y_N$.

\subsection{Complex-linear duality and the actual fixed kernel columns}

For a positive Gram $H$ and a complex-linear functional
$y\mapsto a^{\mathsf T}y$, its squared dual norm is
$a^{\mathsf T}H^{-1}\bar a
 =a^*\overline{H^{-1}}a$.
Indeed its Riesz vector is $H^{-1}\bar a$, and evaluation
of its squared norm gives the first expression; the second
is its equal real conjugate. Therefore the precise metric
identity and determinant consequence of DNE4--7 are
\[
 \mathcal D_N^{(s)}
   =A_\ell^*\overline{(H_{Y,N}^{(s)})^{-1}}A_\ell,\qquad
 \det\mathcal D_N^{(s)}
   =\frac{|\det A_\ell|^2}{\det H_{Y,N}^{(s)}}.
 \tag{DNE8}
\]
The conjugation records the ordinary complex-linear dual
coordinates; it is not discarded by changing the source frame.
In the four-endpoint sum, the fixed factor $|\det A_\ell|^2$
cancels with signs $(1,1,-1,-1)$.

Retain the actual fixed matrix $Z_k$ of FEX27 and FGC9,
with $t=r-t_0$ independent columns in the low-monomial
frame of dimension $\mathfrak m$. Multiplying every
relation column by the same nonzero scalar
$\alpha_A=v!/\mu_v$ does not change its span or its
Schur complement. Consequently its actual native
FEX29 matrix and the exact dual quotient map are
\[
 \widehat H_{K,N}^{(s)}=Z_k^*H_{Y,N}^{(s)}Z_k,\qquad
 B_K=Z_k^{\mathsf T}A_\ell,\qquad
 Q_{K^\vee,N}^{(s)}
 =\left[B_K(\mathcal D_N^{(s)})^{-1}B_K^*\right]^{-1}
 =\overline{(\widehat H_{K,N}^{(s)})^{-1}}.
 \tag{DNE9}
\]
The map $B_K$ evaluates each relation-annihilating functional
on the actual columns $Z_k$. It is onto, because $A_\ell$
is invertible and $Z_k$ has rank $t$. The first inverse-Gram
formula is the minimum dual metric on that same evaluation
target. Substitution of DNE8 yields
\[
 B_K(\mathcal D_N^{(s)})^{-1}B_K^*
 =Z_k^{\mathsf T}\overline{H_{Y,N}^{(s)}}\bar Z_k
 =\overline{\widehat H_{K,N}^{(s)}},
\]
which proves the final equality with every transpose in
its original direction. Empty rank has determinant one.
In particular the full fixed-column return is exactly
\[
 \widehat F_k^{(s)}
 =-\sum_N\sigma_N\log\det Q_{K^\vee,N}^{(s)}.
 \tag{DNE10}
\]
Its relation to the original $F_K^{(s)}$ retains both low
quotient volumes and the full source-map error of FGC12--13
and FEX34--37; DNE10 removes none of those terms.

\subsection{A finite source bound before the high-order derivative}

Use the actual $R,R',B,C_v$ and $R_0$ of ROG6 and RDS11,
and put
\[
 \rho=\tanh(R_0/4),\qquad
 C_{F,k,A}=\frac{2R^v}{|\mu_v|}e^{RR_0}.
 \tag{DNE11}
\]
For any $\mathcal N=\sum\nu_{ij}e^{\beta_{ij}z}
\in\mathscr N_v$, RDS12 proves
$|\mathcal N/E_A|\le C_{F,k,A}\|\nu\|_1$
on the complete disk $|z|\le R_0$.
Apply the original Gamma generating identity directly to
the functional $\Lambda_{\mathcal N/E_A}$:
\[
 \sum_{n\ge0}\Lambda_{\mathcal N/E_A}
          \left(p_n((S'-c')/i)\right)\frac{z^n}{n!}
 =(1+z^2)^{-s/4}e^{ic'\arctan z}
                       \frac{\mathcal N(-i\arctan z)}
                            {E_A(-i\arctan z)}.
 \tag{DNE12}
\]
The quotient at zero is removable. On $|z|=\rho$,
$|\arctan z|\le R_0/4$.
Cauchy's coefficient estimate therefore gives the exact bound
\[
 \begin{gathered}
 \|\Lambda_{\mathcal N/E_A}\|_
             {(\mathbb C[S']_{\le d},\|\cdot\|_{s,c'})^\vee}
       \le \mathcal B_{d,s,k,A}\|\nu\|_1,\\
 \mathcal B_{d,s,k,A}^{\,2}
 =\frac{C_{F,k,A}^2}{M_s}
    (1-\rho^2)^{-s/2}e^{|c'|R_0/2}
       \sum_{n=0}^d\frac{n!}{(s/2)_n}\rho^{-2n}.
 \end{gathered}\tag{DNE13}
\]
Indeed the bound for the $n$th orthonormal coefficient is
the square root of its displayed summand times the common
factor and $\|\nu\|_1$; summing squared coefficients gives
the dual norm. The full Gamma mass remains explicitly in
this formula. No application of $\chi'(D)$ and no
$q'!$ factor is needed in this exact extension norm.

For a fully specified comparison reference, use Euclidean
norm in the original upper exponential coefficient order
$\nu_{ij}$, with its actual quotient by
$\operatorname{im}C^{\mathsf T}$. Its positive Gram in the
fixed DNE3 frame is
\[
 G_{\exp}
 =Z_{\exp}^*
  \left[I-C^{\mathsf T}
     (\bar C C^{\mathsf T})^{-1}\bar C\right]Z_{\exp},
 \qquad
 0\prec\mathcal D_N^{(s)}
       \preceq q\,\mathcal B_{d,s,k,A}^{\,2}G_{\exp}.
 \tag{DNE14}
\]
The reference is explicitly specified here for comparison;
the actual weighted Gamma quotient remains DNE7.
The inverse in $G_{\exp}$ exists by the full row rank of
$C$. Its expression is the orthogonal-projection formula,
and is positive on this quotient frame by DNE3.
Every representative $\nu+C^{\mathsf T}\eta$ remains in
$\mathscr N_v$ and changes $\mathcal N/E_A$ by precisely a
lower exponential, by RDS7. Combining DNE6 and DNE13,
and then using $\|w\|_1\le\sqrt q\|w\|_2$, gives
\[
 x^*\mathcal D_N^{(s)}x
 \le\mathcal B_{d,s,k,A}^{\,2}
      \inf_\eta\|Z_{\exp}x+C^{\mathsf T}\eta\|_1^2
 \le q\,\mathcal B_{d,s,k,A}^{\,2}x^*G_{\exp}x.
\]
This proves DNE14 without a bound on an omitted frame inverse.

At fixed actual period and quartet the finite constant satisfies
\[
 \log\left(1+q\mathcal B_{d,s,k,A}^{\,2}\right)
       =O_{\rm actual}(q),\qquad s=1,k,\quad d=N-v.
 \tag{DNE15}
\]
For detail, $R_0,\rho,v,\mu_v$ are fixed, $R=O(k)$,
$|c'|=O(k)$ and $s\le k$. Thus
$\log C_{F,k,A}=O_{\rm actual}(k+\log(k+1))$.
The original Gamma ratios obey
$n!/(s/2)_n\le2\sqrt n$ for $n\ge1$ and are one at $n=0$,
as proved by the product estimate in FEX10--11. Hence the
sum in DNE13 is at most
$(d+1)\max\{1,2\sqrt d\}\rho^{-2d}$.
Since $d\le2q$, its logarithm is $O_{\rm actual}(q)$.
The remaining factors have logarithms $O_{\rm actual}(k)$,
with the original positive mass retained before estimation.
These bounds prove DNE15.

Restriction of DNE14 to any specified plane, and minimization
under any specified onto quotient map, preserve its scalar
upper bound. In particular it applies to the actual map $B_K$
with the full reference quotient
$(B_KG_{\exp}^{-1}B_K^*)^{-1}$.
This is a finite norm comparison of the original metric
with a completely retained reference form. The reference
contains the full original Vandermonde and conductor data;
neither its determinant nor the original return in DNE10
has been assigned a leading value by this upper bound.

% END CURRENT DNE

\clearpage
% BEGIN CURRENT PEC
% Original symbol poles, complete Gamma dual metric and its four-endpoint return.
\section{Subleading positive exterior expansion of the original dual metric}
\label{sec:PEC}

Fix the original simple quartet, its full unit, a fixed actual five-orbit
period and logarithm branch. Retain all coefficients of the degree-$(8,8)$
conductor $A$, all shifts $b_{ab}$ and every original root. The coefficient
formula PCL8 depends on this fixed period and quartet and is independent
of $k$. Thus the entire symbol
\[
 E_A(w)=\sum_{a,b=0}^8a_{ab}e^{b_{ab}w},\qquad
 v=\operatorname{ord}_0E_A,\quad \mu_v=E_A^{(v)}(0)\ne0
 \tag{PEC1}
\]
is the same function throughout the sequence $k=4l+1\ge9$.
Keep $q=(k+1)^2$, $q'=(k-7)^2$, $\Delta=q-q'=16k-48$,
$\mathfrak m=\Delta-v$, $c'=k/2-4$, both source orders $s=1,k$,
and the original endpoints $N=q-1,q,2q-1,2q$. Write $d=N-v$.
The letter $w$ in this section is the analytic symbol variable.
It does not replace a primary root, polynomial coordinate, period,
or the absolute base $\tau$.

Use the complete DNE frame $Z_{\exp}$, numerators $\mathcal N_a$,
dual Grams $\mathcal D_N^{(s)}$ and coefficient reference $G_{\exp}$.
Their literal definitions, including the complex-linear dual, are
\[
\begin{gathered}
 F_a=\mathcal N_a/E_A,\qquad
 T_{N,s;na}=\Lambda_{F_a}(\phi_{n,s,c'}),\qquad
 V_{N,s;n,(i,j)}=\phi_{n,s,c'}(\beta'_{ij}),\\
 \mathcal D_N^{(s)}=T_{N,s}^*
 [I-V_{N,s}(V_{N,s}^*V_{N,s})^{-1}V_{N,s}^*]T_{N,s},\\
 G_{\exp}=Z_{\exp}^*
 [I-C^{\mathsf T}(\bar C C^{\mathsf T})^{-1}\bar C]Z_{\exp}>0.
\end{gathered}\tag{PEC2}
\]
Here $C$ is the entire original conductor, and the exact Gamma
orthonormal frame has $h_n=M_sn!(s/2)_n$, $M_s=(2\pi)^{s/2}$.
DNE6--8 proves these metrics, retaining the full lower evaluation
space and both conditioning factors. In particular $G_{\exp}$ is
the minimum Euclidean norm of an original exponential coefficient
vector modulo $\operatorname{im}C^{\mathsf T}$; it is not substituted
for $\mathcal D_N^{(s)}$.

\subsection{The actual symbol-zero jet subspaces}
On $|z|<1$ use the branch of $\arctan z$ vanishing at zero and set
\[
 w(z)=-i\arctan z,\qquad
 \mathcal E_A(z)=z^{-v}E_A(w(z)),\qquad
 \mathcal E_A(0)=\frac{\mu_v(-i)^v}{v!}\ne0.
 \tag{PEC3}
\]
The power series for $\arctan$ is holomorphic on this disk and gives
$|\arctan z|\le\operatorname{arctanh}|z|$. Its derivative is
$(1+z^2)^{-1}$, so $w'$ is nonzero there. Moreover $w$ is injective:
$e^{-2w(z)}=(1+iz)/(1-iz)$, and equality of two $w$ values implies
equality under this injective fractional linear function. The displayed
removable value in PEC3 follows by substituting the first nonzero
Taylor term of PEC1. Thus $\mathcal E_A$ is a nonzero holomorphic
function on the entire open unit disk.

Call $0<r<1$ admissible when $E_A(w(z))\ne0$ for $|z|=r$.
For every such $r$, let $\mathcal Z_r$ be the finite set of nonzero
zeros $\omega$ of $E_A$ in $w(\{|z|<r\})$, and write $h_\omega$
for their exact orders. Define
\[
 H_r=\sum_{\omega\in\mathcal Z_r}h_\omega,
 \qquad m_r=\min_{|z|=r}|E_A(w(z))|>0.
 \tag{PEC4}
\]
Finiteness follows from isolation of the zeros of $\mathcal E_A$
in a compact subdisk. Positivity of $m_r$ follows from continuity
and admissibility on the compact circle. Every compact subdisk contains
only finitely many excluded radii. Consequently there are admissible
radii arbitrarily close to one. These are properties of the specified
symbol, with no assumed zero-free unit disk. Its zero count can also
be written exactly as
\[
 H_r=\frac1{2\pi i}\int_{|z|=r}
             \frac{\mathcal E_A'(z)}{\mathcal E_A(z)}\,dz .
 \tag{PEC5}
\]
Indeed factor out the finitely many interior zeros with their
multiplicities. The logarithmic derivative of the remaining nonvanishing
factor integrates to zero on the disk, while each factor $(z-z_0)^h$
contributes $h$. The nonzero derivative of $w$ preserves the orders
$h_\omega$. This proves PEC5 in the original symbol variables.

In the DNE coefficient frame, form the full jet matrix
\[
 (J_{r,k})_{(\omega,b),a}
   =\frac{\mathcal N_a^{(b)}(\omega)}{b!},
 \quad 0\le b<h_\omega,\quad
 c_{r,k}=\operatorname{rank}J_{r,k},\quad
 \mathcal H_{r,k}=\ker J_{r,k}\subset\mathbb C^{\mathfrak m}.
 \tag{PEC6}
\]
Its entries are the complete sums
$\sum_{i,j}(Z_{\exp})_{(i,j),a}\beta_{ij}^{b}
e^{\beta_{ij}\omega}/b!$. Thus every original root, factorial and
coefficient is retained. Its codimension is $c_{r,k}\le
\min\{\mathfrak m,H_r\}$. Addition of an original coefficient vector
$C^{\mathsf T}\eta$ changes its numerator by $E_A$ times the
corresponding lower exponential polynomial, by RDS7. That change has
zero jets in PEC6. Therefore the cancellation subspace is a specified
subspace of the original quotient
$\mathscr N_v/(E_A\mathscr E_{k-8})$, independent of the choice of
representative and of both $N$ and $s$. In particular the minimum
Euclidean representative defining $G_{\exp}$ stays in this subspace.

For $x\in\mathcal H_{r,k}$, all interior poles of
$\mathcal N_x/E_A$ cancel, including zero by membership in $\mathscr N_v$.
Admissibility and discreteness give a neighborhood of the closed $r$
disk with no additional boundary pole. Thus
$\mathcal N_x(w(z))/E_A(w(z))$ is holomorphic on that neighborhood.
This is an exact pole cancellation imposed by the actual numerator
jets, not a removal of poles from the original symbol.

\subsection{A complete finite Gamma bound on each cancellation subspace}
Use the original root bound $R=|k/2|+k\sqrt{\delta^2+\gamma^2}$
and put $a_r=\operatorname{arctanh}r$. For an original exponential
coefficient vector $\nu$ of a canceled numerator, the exact generating
function DNE12 satisfies, on $|z|=r$,
\[
\left|(1+z^2)^{-s/4}e^{ic'\arctan z}
          \frac{\sum\nu_{ij}e^{\beta_{ij}w(z)}}{E_A(w(z))}\right|
 \le \frac{(1-r^2)^{-s/4}}{m_r}
             e^{(R+|c'|)a_r}\|\nu\|_1 .
 \tag{PEC7}
\]
This follows term by term from $|w(z)|\le a_r$ and the root bound,
with the original exponential coefficient order. The holomorphic
branch of $(1+z^2)^{-s/4}$ is the one taking value one at zero;
its modulus is $|1+z^2|^{-s/4}$. No source mass has yet canceled.
Cauchy's coefficient integral and $h_n=M_sn!(s/2)_n$ therefore give
\[
 \|\Lambda_{\mathcal N/E_A}\|_{(\mathcal P_d,\|\cdot\|_{s,c'})^\vee}^2
 \le B_{r,N,s}^{2}\|\nu\|_1^2,\qquad
 B_{r,N,s}^{2}=
 \frac{e^{2(R+|c'|)a_r}(1-r^2)^{-s/2}}{M_sm_r^2}
 \sum_{n=0}^{d}\frac{n!}{(s/2)_n}r^{-2n}.
 \tag{PEC8}
\]
The squared dual norm is the sum of the squared values on the complete
orthonormal frame of degrees zero through $d$. The exponential generating
series gives an $n!$ before division by $\sqrt{h_n}$, yielding precisely
the ratio $n!/(s/2)_n$ above. This verifies every term of PEC8.

Every allowed addition $E_A\mathscr E_{k-8}$ preserves cancellation.
Minimize over these representatives as in DNE6 and use
$\|\nu\|_1^2\le q\|\nu\|_2^2$. This proves the quadratic-form bound
\[
 x^*\mathcal D_N^{(s)}x\le C_{r,N,s}x^*G_{\exp}x
 \quad(x\in\mathcal H_{r,k}),\qquad C_{r,N,s}=qB_{r,N,s}^{2}.
 \tag{PEC9}
\]
It uses the entire lower conditioning space, not a chosen extension.
For each fixed admissible radius, $E_A,m_r,a_r$ are independent of $k$.
The product estimate $n!/(s/2)_n\le2\sqrt n$ for $n\ge1$
and its value one at $n=0$ give the explicit bound
\[
 C_{r,N,s}\le
 \frac{q\,e^{2(R+|c'|)a_r}(1-r^2)^{-s/2}}{M_sm_r^2}
 (d+1)\max\{1,2\sqrt d\}\,r^{-2d}.
 \tag{PEC10}
\]
Here $R+|c'|=O_{\delta,\gamma}(k)$, $s\le k$ and $d\le2q$.
Since $M_s>1$, taking the positive part of the logarithm gives
\[
 \log^+ C_{r,N,s}\le 2d\log(1/r)
                    +O_{{\rm actual},r}(k+\log q).
 \tag{PEC11}
\]
The finite expression PEC10 supplies the constants in this assertion;
no asymptotic assumption on an unevaluated numerator is used.

\subsection{The whole positive logarithmic spectrum}
Let $\lambda_1\ge\cdots\ge\lambda_{\mathfrak m}>0$ be the
eigenvalues of the specified relative operator
\[
 T_N^{(s)}=G_{\exp}^{-1/2}\mathcal D_N^{(s)}G_{\exp}^{-1/2}.
 \tag{PEC12}
\]
The positive square root is unique by diagonalizing the positive
Hermitian matrix. This congruence is the exact comparison of the
two retained forms. DNE14 gives $\lambda_a\le C_{0,N,s}$ for all $a$,
where $C_{0,N,s}=q\mathcal B_{d,s,k,A}^{2}$ is its complete explicit
constant, and DNE15 gives $\log^+C_{0,N,s}=O_{\rm actual}(q)$.
PEC9 bounds the Rayleigh quotient on a subspace of codimension $c_{r,k}$.
Hence at most $c_{r,k}$ eigenvalues can exceed $C_{r,N,s}$: otherwise
the span of $c_{r,k}+1$ corresponding orthogonal eigenvectors would
intersect that subspace nontrivially, contradicting the bound there.
Consequently
\[
\begin{gathered}
 \mathcal P_{N,s}:=\sum_{a=1}^{\mathfrak m}\log^+\lambda_a
       \le E_{r,N,s},\\
 E_{r,N,s}=c_{r,k}\log^+C_{0,N,s}
        +(\mathfrak m-c_{r,k})\log^+C_{r,N,s},\qquad
 \varepsilon_{N,s}^{\vee}=\inf_{r\ \mathrm{admissible}}E_{r,N,s}.
\end{gathered}\tag{PEC13}
\]
All terms are finite and nonnegative, and the admissible set is nonempty.
Thus $\mathcal P_{N,s}\le\varepsilon_{N,s}^{\vee}$ even when the
infimum is not attained. For every fixed admissible $r$, $c_{r,k}\le H_r$
is bounded independently of $k$. Since $\mathfrak m/k\to16$ and
$d/q\le2$, PEC11 and DNE15 imply
\[
 0\le\limsup_{k\to\infty}
 \frac{\varepsilon_{N,s}^{\vee}}{kq}\le64\log(1/r).
 \tag{PEC14}
\]
This holds simultaneously for the two sources and four indicated
endpoint sequences, which form a finite list. Choose admissible radii
arbitrarily close to one only after taking this bound for each fixed
radius. The right side tends to zero. We have proved
\[
 \boxed{\quad \mathcal P_{N,s}\le\varepsilon_{N,s}^{\vee}=o(kq).
                  \quad}\tag{PEC15}
\]
Only the original source orders $1,k$ are used in the two limiting
steps; the actual period and operator remain fixed. No additional
source order or presumed spectral gap is inserted. Although $H_r$ can increase as $r$ tends to one,
it is finite before each $k$ limit, which is exactly what PEC14 uses.

This also controls every exterior degree of the comparison map.
The singular values squared of the identity from $(\mathbb C^{\mathfrak m},
G_{\exp})$ to $(\mathbb C^{\mathfrak m},\mathcal D_N^{(s)})$ are
the $\lambda_a$. In orthonormal wedge frames, each squared singular
value is a product over a subset of these numbers. Therefore
\[
 \log\|\bigwedge^t\operatorname{id}\|
     \le\tfrac12\mathcal P_{N,s}
     \le\tfrac12\varepsilon_{N,s}^{\vee}
       \quad(0\le t\le\mathfrak m).
 \tag{PEC16}
\]
This is an all-degree bound in the exact two metrics. It does not
bound the inverse exterior map; the small relative eigenvalues are
retained below.

\subsection{Every fixed restriction and quotient keeps this bound}
For an injective matrix $J$ and an onto matrix $B$, retain respectively
the full paired metrics
\[
\begin{gathered}
 D_J=J^*\mathcal D_N^{(s)}J,\quad G_J=J^*G_{\exp}J,\\
 D_B=[B(\mathcal D_N^{(s)})^{-1}B^*]^{-1},\quad
 G_B=[BG_{\exp}^{-1}B^*]^{-1}.
\end{gathered}\tag{PEC17}
\]
Their relative eigenvalues $\mu_1\ge\cdots\ge\mu_t>0$ satisfy
$\mu_a\le\lambda_a$ for $1\le a\le t$. For a restriction,
write $U=G_{\exp}^{1/2}J G_J^{-1/2}$, so $U^*U=I$ and its
relative operator is $U^*T_N^{(s)}U$. The variational description
of its $a$th largest eigenvalue gives the asserted upper bound.
Explicitly the preimage of any subspace of codimension $a-1$ has
codimension at most $a-1$, so the usual minimum of maximal Rayleigh
quotients on such subspaces cannot increase under restriction.
This variational formula itself follows by expanding vectors in an
orthonormal eigenbasis: the complement of the first $a-1$ vectors
gives the upper bound, and dimension intersection with their first
$a$ vectors gives the lower bound.

For the quotient put $V=(BG_{\exp}^{-1}B^*)^{-1/2}BG_{\exp}^{-1/2}$.
Then $VV^*=I$ and its relative operator is
$(V(T_N^{(s)})^{-1}V^*)^{-1}$. The $a$th smallest eigenvalue of
the compression of $(T_N^{(s)})^{-1}$ is at least its $a$th smallest
full eigenvalue, by the same variational argument applied to smallest
eigenvalues. Taking reciprocals proves $\mu_a\le\lambda_a$.
It follows that for either metric pair in PEC17,
\[
 \sum_{a=1}^t\log^+\mu_a\le\mathcal P_{N,s}
                         \le\varepsilon_{N,s}^{\vee}.
 \tag{PEC18}
\]
No bound for a coefficient inverse of $B$ or $J$ occurs: their full
reference metrics in PEC17 remain part of the assertion.

\subsection{The exact negative spectral return on the actual kernel}
Retain $Z_k$, $A_\ell$ and the actual onto matrix
$B_K=Z_k^{\mathsf T}A_\ell$ of DNE9. Put
\[
 D_{K,N}=[B_K(\mathcal D_N^{(s)})^{-1}B_K^*]^{-1}
       =\overline{(\widehat H_{K,N}^{(s)})^{-1}},\qquad
 G_K=[B_KG_{\exp}^{-1}B_K^*]^{-1}.
 \tag{PEC19}
\]
These are exactly the native original kernel dual and its fixed
reference, with no changed relation or factorial column. For this
pair define, using all its positive relative eigenvalues $\mu_{a,N}$,
\[
 p_{K,N}=\sum_a\log^+\mu_{a,N},\quad
 n_{K,N}=\sum_a\log^+(\mu_{a,N}^{-1}),\quad
 A_{K,k,s}=n_{K,q-1}+n_{K,q}-n_{K,2q-1}-n_{K,2q}.
 \tag{PEC20}
\]
The original weighted polynomial spaces are nested and the complete
relation maps restrict to the earlier cutoffs. Minimizing a primal
norm over their increasing lift sets decreases $H_{Y,N}^{(s)}$.
By DNE8 this increases $\mathcal D_N^{(s)}$. Taking minima over the
same $B_K$ fibres preserves this increase for $D_{K,N}$. Relative
eigenvalues in its fixed $G_K$ metric therefore increase by the
variational formula. Hence $p_{K,N}$ increases, $n_{K,N}$ decreases,
and $A_{K,k,s}\ge0$.

Since $\log\mu=\log^+\mu-\log^+(1/\mu)$, the complete determinant
formula DNE10 yields the exact identity and finite bound
\[
\begin{gathered}
 \widehat F_k^{(s)}=A_{K,k,s}+P_{K,k,s},\\
 P_{K,k,s}=p_{K,2q-1}+p_{K,2q}-p_{K,q-1}-p_{K,q},\qquad
 0\le P_{K,k,s}\le
 E_{k,s}^{\vee}:=\varepsilon_{2q-1,s}^{\vee}
                       +\varepsilon_{2q,s}^{\vee}=o(kq).
\end{gathered}\tag{PEC21}
\]
The reference determinant cancels with the four original signs only
after PEC19's complete dual identity has been used. Formula PEC18
bounds each of the two high positive parts, proving the last inequality.
Thus the entire leading native kernel return is carried by the specified
decrease of its small relative eigenvalue deficits; none has been set
to zero or inferred from the rank.

Retain the complete unequal $D_s^-,D_s^+$ and original $\mathcal L_s$
from FEX32--35, with $t_0=\dim(K\cap L)$ and $L=\mathbb C[S]_{<v}$.
Substituting PEC21 into that proved interval gives
\[
\begin{split}
 A_{K,k,s}-D_s^+-(v-t_0)\mathcal L_s
 &\le F_K^{(s)}\\
 &\le A_{K,k,s}+E_{k,s}^{\vee}-D_s^-+v\mathcal L_s.
\end{split}\tag{PEC22}
\]
Both original low corrections and both unequal flag errors are retained.
FEX37 separately gives the useful exact symmetric allowance
\[
 |F_K^{(s)}-A_{K,k,s}|
   \le\mathcal E_{k,s}^{\rm FEX}+E_{k,s}^{\vee}=o(kq),
 \tag{PEC23}
\]
where $\mathcal E_{k,s}^{\rm FEX}$ is the full finite FEX36 expression.
This is a proved original-kernel metric reduction, not an evaluated
individual leading coefficient.

\subsection{The complementary original residual and its low subspace}
Use the fixed onto map $P_Y:\mathbb C^{\mathfrak m}\to\mathbb C^{j_0}$
with kernel $\operatorname{im}Z_k$, where
$j_0=\mathfrak m-(r-t_0)=j-v+t_0$ and $j=8k-32$.
For an explicit choice, let $\mathcal I$ be the first set of $r-t_0$
row indices in the fixed ordering for which $E_{\mathcal I}^{\mathsf T}Z_k$
is invertible, and let $\mathcal J$ be the complementary indices. Define
\[
 P_Y=E_{\mathcal J}^{\mathsf T}
       -E_{\mathcal J}^{\mathsf T}Z_k
        (E_{\mathcal I}^{\mathsf T}Z_k)^{-1}E_{\mathcal I}^{\mathsf T}.
\]
Independence of the columns, proved in FGC10, ensures such a minor exists.
Multiplication gives $P_YZ_k=0$ and $P_YE_{\mathcal J}=I$.
Its rank is $j_0$ and its kernel has dimension $r-t_0$, so that kernel
is exactly $\operatorname{im}Z_k$. Empty rank uses the identity quotient.
Set
\[
 J_0=A_\ell^{-1}P_Y^{\mathsf T},\quad
 D_{0,N}=J_0^*\mathcal D_N^{(s)}J_0,\quad G_0=J_0^*G_{\exp}J_0,
 \quad Q_{0,N}=[P_Y(H_{Y,N}^{(s)})^{-1}P_Y^*]^{-1}.
 \tag{PEC24}
\]
Then $B_KJ_0=0$ and $J_0$ is a basis of $\ker B_K$ by dimensions.
Direct substitution in DNE8 proves the complete dual identity
\[
 D_{0,N}=\bar P_Y\,\overline{(H_{Y,N}^{(s)})^{-1}}P_Y^{\mathsf T}
        =\overline{Q_{0,N}^{-1}}.
 \tag{PEC25}
\]
Thus PEC24 is the annihilator of the actual kernel columns, with its
full metric, not a complementary coordinate plane chosen for convenience.

Define $p_{0,N},n_{0,N}$ from the relative eigenvalues of $(D_{0,N},G_0)$
exactly as in PEC20, and $A_{0,k,s}$, $P_{0,k,s}$ by the same four
endpoint differences. Restriction in PEC18 and the preceding monotonicity
proof give
\[
 F_{0,k,s}:=\sum_N\sigma_N\log\det Q_{0,N}
     =A_{0,k,s}+P_{0,k,s},\quad
 A_{0,k,s}\ge0,\quad0\le P_{0,k,s}\le E_{k,s}^{\vee}.
 \tag{PEC26}
\]
The scalar $F_{0,k,s}$ here denotes this displayed quotient return only;
it is not the seven-moment quantity $\mathcal F_{0,k}^{(s)}$.

To identify the original target exactly, put $W_0=\Xi(L)\subset B_\Xi$.
The map $L\to W_0$ has kernel $K\cap L$, so $\dim W_0=v-t_0$.
There is the onto map, in the original FGC coordinates,
\[
 V_A\longrightarrow Y_N/\operatorname{im}Z_k,
 \quad[p]\longmapsto
 \left[\frac{\alpha_A\mathcal T_Ap}{\chi'}\right],
 \quad\alpha_A=v!/\mu_v,
 \qquad Y_N/\operatorname{im}Z_k\simeq B_\Xi/W_0.
 \tag{PEC27}
\]
FGC4--10 proves independence of representatives and surjectivity.
Its kernel is $K+L$: a zero image subtracts a member of $K$ with
that image in $Y_N$, leaving an element of $L$ by FGC6. Conversely
both $K$ and $L$ map to zero. The induced isomorphism with $B_\Xi/W_0$
sends the class of $p$ to the class of $\Xi[p]$. This verifies both
directions of the target identification.

At each endpoint the source map is the original $S_N$ of CEP28
followed by $D_{\alpha,N}(x,f)=(x,\alpha_Af)$ on GIP's full graph.
The complete quotient map PEC27 kills its entire relation graph and
kernel columns. In the product source norm, minimization first over
the unrestricted low coordinate and then over the polynomial and
$Z_k$ relations is exactly $Q_{0,N}$. The exterior norms of
$D_{\alpha,N}^{\pm1}$ in these same product metrics are at most
$\exp(j_0|\log|\alpha_A||)$ in degree $j_0$.
The minimum-lift and projection argument of GIP6 applied to this
composite source map therefore proves, for the original induced
metric $\overline Q_{\Xi,N}$ on $B_\Xi/W_0$ in the transported frame,
\[
 |\log\det\overline Q_{\Xi,N}-\log\det Q_{0,N}|
 \le2E_N^{\rm CEP}+2j_0|\log|\alpha_A||.
 \tag{PEC28}
\]
Orthogonal projection off the entire kernel is an exterior contraction,
and the inverse composite gives the other direction. This checks the
full source, target, metric and rank in the comparison.

Choose a fixed original frame $I_{W_0}$ of $W_0$ and retain its exact
restriction return
$F_{W_0}^{(s)}=\sum_N\sigma_N\log\det(I_{W_0}^*Q_{\Xi,N}^{\rm fix}I_{W_0})$.
Completing that frame by fixed quotient lifts gives a block Gram whose
Schur complement is $\overline Q_{\Xi,N}$. Its fixed frame determinant
cancels with the four signs. Consequently the original residual is
\[
\begin{gathered}
 \Phi_{k,s}=F_{W_0}^{(s)}+A_{0,k,s}+P_{0,k,s}+\delta_{0,k,s},\\
 |\delta_{0,k,s}|\le E_{k,s}^{\rm tr}
 :=2\sum_NE_N^{\rm CEP}+8j_0|\log|\alpha_A||=o(kq),\\
 0\le F_{W_0}^{(s)}\le(v-t_0)\mathcal L_s=O_{\rm actual}(q).
\end{gathered}\tag{PEC29}
\]
The last inequality follows by restricting and minimizing the original
AKS covariance comparison on this fixed subquotient of the primary
space, exactly as in FEX32. Its metrics decrease on increasing source
spaces, proving nonnegativity. The full finite $\mathcal L_s$ is FEX33,
whose proved order is $O(q)$, and $v\le80$. Thus no original low
contribution has been dropped. In particular the finite and asymptotic
consequences are
\[
\begin{gathered}
 A_{0,k,s}-E_{k,s}^{\rm tr}\le\Phi_{k,s}
 \le A_{0,k,s}+E_{k,s}^{\vee}+E_{k,s}^{\rm tr}
                          +(v-t_0)\mathcal L_s,\\
 \Phi_{k,s}=A_{0,k,s}+o(kq).
\end{gathered}\tag{PEC30}
\]

\subsection{Backward and forward return to the same arithmetic class}
The complete OVR/ECL identity $F_K^{(s)}+\Phi_{k,s}
=\Delta q\mathfrak C_\partial+o(kq)$ and PEC23,PEC30 give
\[
\begin{gathered}
 A_{K,k,s}+A_{0,k,s}=\Delta q\mathfrak C_\partial+o(kq),\\
 \frac{A_{K,k,s}+A_{0,k,s}}{kq}\longrightarrow16\mathfrak C_\partial,\\
 21.66851180520<16\mathfrak C_\partial<21.66851180521.
\end{gathered}
 \tag{PEC31}
\]
Both deficits are on the exact DNE dual quotient and its actual
kernel annihilator, with the full $G_{\exp}$ references. The equality
evaluates their sum and retains their individual distributions.
AMT's already proved source-order agreement and PEC23,PEC30 also
give $(A_{K,k,1}-A_{K,k,k})/(kq)\to0$ and the same statement for $A_0$.

For finite original kernel bounds intersect PEC22 with OER19. Retain
the full AMT arithmetic intervals and take their exact boundary
complements. On the original arithmetic period domain $|u|\ge R_*$,
the signed identity $\Psi_k=F_K^{(k)}+\mathcal H_k+\delta_k^\sigma$
then proves, with unchanged signs,
\[
\begin{split}
 A_{K,k,k}-D_k^+-(v-t_0)\mathcal L_k
       +\mathcal H_k+\delta_k^\sigma
 &\le\Psi_k\\
 &\le A_{K,k,k}+E_{k,k}^{\vee}-D_k^-+v\mathcal L_k
       +\mathcal H_k+\delta_k^\sigma .
\end{split}\tag{PEC32}
\]
Here $D_k^\pm$ means the original FEX bounds at source $s=k$,
and $\mathcal L_k$ is the same source's FEX33 allowance. The original
same-class action inequality OER20 uses precisely the right endpoint
of PEC32 plus the complete $C_k^{\rm budget}$. This addition preserves
the original monic window, contraction penalties, mixed pairings,
action leakage and positive-primary exterior class. Where the seven
moments are used, retain PMR's finite range and $e^{0,\pm}=e^\pm+c_k^\circ$;
outside that range retain the original OAR outer expression.

The calculations PEC1--32 prove a subleading positive exterior budget
and return it to the original kernel and residual, whose leading
negative spectral deficits are explicitly specified. Calculating
these actual small eigenvalues and the resulting same-class exclusion
bound is the remaining analytic step in these formulas. No assertion
about their individual leading values or an RH conclusion follows
merely from the evaluated sum PEC31.

% END CURRENT PEC

\clearpage
% BEGIN CURRENT ACR
\section{The actual common-factor map on the original kernel}
\label{sec:ACR}

Fix the original simple quartet, its complete unit and period data,
and the original logarithm branch on the proved five-orbit domain.
Use $k=4l+1\geq9$, $k'=k-8$, and retain
\[
 \begin{gathered}
 q=(k+1)^2,\quad q'=(k-7)^2,\quad
 \Delta=q-q'=16k-48,\\ r=8k-16,\quad j=8k-32,
 \qquad c=k/2,\quad c'=k'/2.
 \end{gathered}
 \tag{ACR1}
\]
The original primary points, also at every integer index in the
second formula, are
\[
 \begin{aligned}
 \beta_{ab}&=c+(2a-k)\delta+i(2b-k)\gamma,
                         &&0\leq a,b\leq k,\\
 \tau_{ij}&=c'+(2i-k')\delta+i(2j-k')\gamma,
                         &&(i,j)\in\mathbb Z^2,\\
 b_{ab}&=4+(2a-8)\delta+i(2b-8)\gamma,
                         &&0\leq a,b\leq8.
 \end{aligned}\tag{ACR2}
\]
In the imaginary terms $i$ is the imaginary unit; subscripts are
integer indices. The functions of the two integer indices are
injective, by comparison of real and imaginary parts and
$\delta>0$, $\gamma>2$. In particular
$\tau_{ij}+b_{ab}=\beta_{i+a,j+b}$ whenever the latter indices
are in their displayed range.

Write $\chi=\prod_{a,b=0}^k(S-\beta_{ab})$ and
$\chi'=\prod_{i,j=0}^{k'}(S'-\tau_{ij})$. The prime names
the original lower-grid polynomial, and is not differentiation.
All dual pairings until a Hermitian Gram is explicitly mentioned
are complex-linear.

\subsection{Support and its precise extended coefficient rectangle}

Let the actual PCL8--11 biform be
\[
 A=\sum_{a,b=0}^8a_{ab}
       z_0^a z_1^{8-a}w_0^b w_1^{8-b}\ne0,
 \qquad
 \mathcal S=\{(a,b):a_{ab}\ne0\}.
 \tag{ACR3}
\]
Here, literally, its factors are
$Q_0(L_h(u^{-1})t)$, $1\leq h\leq4$, where
$t=(z_0w_0,z_0w_1,z_1w_0,z_1w_1)^T$ and
$L_h=U^{-1}V\mathcal R^{-1}D^{-h}\mathcal R V^{-1}U$.
Thus every coefficient is the complete finite sum of products
in PCL10, with the entire original matrix series PCL2.
No coefficient is assigned in this calculation. Put
\[
 a_- =\min_{\mathcal S}a,\quad a_+=\max_{\mathcal S}a,
 \quad b_-=\min_{\mathcal S}b,\quad b_+=\max_{\mathcal S}b,
 \quad d_a=a_+-a_-,\quad d_b=b_+-b_-.
 \tag{ACR4}
\]
The following finite index sets have their literal integer intervals:
\[
 \begin{gathered}
 I_{\rm all}=[-a_-,k-a_+]\times[-b_-,k-b_+],\qquad
 I_{\rm in}=[0,k']^2,\qquad I_G=I_{\rm all}\setminus I_{\rm in},\\
 G_A(S')=\prod_{(i,j)\in I_G}(S'-\tau_{ij}),\qquad
 g=|I_G|=(k+1-d_a)(k+1-d_b)-q'.
 \end{gathered}\tag{ACR5}
\]
Every interval is nonempty and the first rectangle contains
$I_{\rm in}$, since all active indices are in $[0,8]^2$.
The intersection of the shifted full-grid rectangles
$[-a,k-a]\times[-b,k-b]$, over the actual active pairs,
is exactly $I_{\rm all}$. Each full shifted polynomial
$\chi(S'+b_{ab})/\chi'(S')$ has the simple roots in its
rectangle with $I_{\rm in}$ removed. Distinctness in ACR2
therefore proves that $G_A$ is precisely their monic greatest
common divisor, with the degree in ACR5. This reproves the
RCF common factor from the actual support. At $g=0$ its
polynomial is the empty product one.

Use the original ordered biform coefficient space
$\mathcal B_k$, with basis
$e_{ab}^{(k)}=z_0^a z_1^{k-a}w_0^b w_1^{k-b}$.
For each $(i,j)\in I_{\rm all}$ define the complete column
\[
 f_{ij}=\sum_{(a,b)\in\mathcal S}a_{ab}e_{i+a,j+b}^{(k)}.
 \tag{ACR6}
\]
All indices in this expression lie in $[0,k]^2$: the lower
bound follows from $i\geq-a_-$, $j\geq-b_-$; the upper
bound follows from $i\leq k-a_+$, $j\leq k-b_+$.
Let $C_{\rm in}$, $C_G$ and $C_{\rm all}$ have these columns
for $I_{\rm in}$, $I_G$ and $I_{\rm all}$, in fixed
lexicographic orders. These coefficient matrices use every
active coefficient, each once, and no conjugate or binomial
multiplier.

The full boundary factor and its remaining biform are
\[
 \begin{gathered}
 M_\partial=z_0^{a_-}z_1^{8-a_+}w_0^{b_-}w_1^{8-b_+},
 \qquad A=M_\partial A_{\rm red},\\
 \deg A_{\rm red}=(d_a,d_b),\qquad
 A_{\rm red}=\sum_{a,b}a_{ab}
       z_0^{a-a_-}z_1^{a_+-a}w_0^{b-b_-}w_1^{b_+-b}.
 \end{gathered}\tag{ACR7}
\]
The formula retains the full scalar and every coefficient of
$A$; only the displayed coordinate monomial is factored.
For bidegrees $(p,t)$ let $\mathcal B_{p,t}$ be the complete
biform space of those bidegrees and set it to zero if either
degree is negative. Multiplication by a nonzero biform is
injective, since the ambient polynomial ring is an integral
domain. ACR6 and ACR7 then give the exact coefficient images
\[
 \begin{gathered}
 H_k:=\operatorname{im}C_{\rm all}
       =A_{\rm red}\mathcal B_{k-d_a,k-d_b},\qquad
 \operatorname{im}C_{\rm in}=A\mathcal B_{k-8,k-8},\\
 \operatorname{rank}C_{\rm all}=q'+g,
 \qquad \operatorname{rank}C_{\rm in}=q',\qquad
 H_k=\operatorname{im}C_{\rm in}\oplus\operatorname{im}C_G.
 \end{gathered}\tag{ACR8}
\]
Indeed $f_{ij}$ is $A_{\rm red}$ times the monomial of
bidegree $(k-d_a,k-d_b)$ whose two indices are
$(i+a_-,j+b_-)$. As $(i,j)$ ranges over $I_{\rm all}$,
these are all its basis monomials. Their products are
independent by injectivity, proving every image and rank
in ACR8. For the inner subset the same multiplication is
exactly $A$ times its original degree-$(k-8,k-8)$ basis.

\subsection{The residual map evaluated on every original primary coordinate}

Let $E=\mathbb C[S]/(\chi)$ and let
$\mathcal V_\beta:E\to\mathbb C^{[0,k]^2}$ be the original
primary evaluation map. Define the original conductor
annihilator $V_A$ and the original kernel $K$ by
\[
 V_A=\{[p]:C_{\rm in}^T\mathcal V_\beta[p]=0\},
 \qquad K=\ker\Lambda\subseteq V_A.
 \tag{ACR9}
\]
The inclusion follows from the full original trace identity
$A\mathcal B_{k-8,k-8}\subseteq J_k$, where $J_k$ is the
image of the original invariant algebra under $x=\mathsf Ht$;
the complex-linear annihilator of $J_k$ is $K$ by OQC8.
Equivalently, direct evaluation proves
\[
 \mathcal T_Ap(S')=\sum_{a,b}a_{ab}p(S'+b_{ab}),\qquad
 (C_{\rm in}^T\mathcal V_\beta p)_{ij}
                 =(\mathcal T_Ap)(\tau_{ij}),\qquad
 [p]\in V_A\ \Longleftrightarrow\ \chi'\mid\mathcal T_Ap.
 \tag{ACR10}
\]
The roots of $\chi'$ are distinct, so vanishing at them is
equivalent to its polynomial divisibility. Replacing $p$ by
$p+\chi h$ preserves this property, because
$\chi(S'+b_{ab})=\chi'(S')d_{ab}(S')$ for the original
monic shifted factors. Thus ACR9 and ACR10 agree on the
original quotient, with all original positive shifts.

For every $\ell=(i,j)\in I_G$ one has $\chi'(\tau_\ell)\ne0$.
The full common-factor residual has the exact evaluation
\[
 \boxed{\quad
 (\rho_G[p])(\tau_\ell)
   =\frac{1}{\chi'(\tau_\ell)}
        \sum_{a,b}a_{ab}p(\beta_{i+a,j+b})
   =\frac{f_\ell^T\mathcal V_\beta[p]}{\chi'(\tau_\ell)}.
 \quad}\tag{ACR11}
\]
Here $\rho_G[p]=\operatorname{rem}_{G_A}
 (\mathcal T_Ap/\chi')$ is in the original remainder
coefficient frame. Divisibility ACR10 gives a polynomial
quotient. Evaluation at a simple root of $G_A$ commutes
with its monic remainder and gives the first expression.
The point identity ACR2 proves the second. A representative
change adds $\sum a_{ab}d_{ab}h(S'+b_{ab})$, divisible by
the full $G_A$, because it divides every active $d_{ab}$.
This proves the stated map on $V_A$, independently of
representatives, and re-establishes the exact RCF15 map.

Let $W_G=(\tau_\ell^d)_{\ell\in I_G,\,0\leq d<g}$
be the original ordered Vandermonde and
$D_G=\operatorname{diag}(\chi'(\tau_\ell)^{-1})$.
For $g>0$, both are invertible, since all their indicated
differences and diagonal entries are nonzero. The same
notation has its unique empty meaning at $g=0$.
Let $\pi:\mathcal B_k\twoheadrightarrow\mathbb C^r$
be any exact OCF coefficient quotient with $\ker\pi=J_k$.
OKM proves that the original kernel frame is
$I_K=\mathcal V_\beta^{-1}\pi^T$. Therefore
\[
 \boxed{\quad
 \rho_G I_K=W_G^{-1}D_G C_G^T\pi^T,
 \qquad B_G:=\pi C_G\in\operatorname{Mat}_{r\times g}(\mathbb C),
 \qquad t_G=g-\operatorname{rank}B_G.
 \quad}\tag{ACR12}
\]
The first equality follows by applying ACR11 on every
kernel column and then the specified inverse evaluation
map. The matrix $C_G^T\pi^T=B_G^T$ uses ordinary
transpose. Multiplication by the two invertible target
maps does not change rank, proving the last equality
with the definition in RCM22. Every entry of the actual
matrix needed for this rank is explicitly
\[
 (B_G)_{h,(i,j)}=\sum_{(a,b)\in\mathcal S}
                       a_{ab}\pi_{h,(i+a,j+b)}.
 \tag{ACR13}
\]
Thus this is a matrix with at most $81$ terms per entry,
$r=8k-16$ rows and $g\leq16k-48$ columns, computed from
the fixed original period coefficients and original kernel
quotient. It requires no full symmetric-power matrix.

For an exact all-case scalar prescription put
$\Delta_0(B_G)=1$, and let $\Delta_h(B_G)$ be the sum
of the squared absolute values of all its $h$ by $h$ minors.
Then
\[
 t_G=g-\max\{h:\Delta_h(B_G)>0\}.
 \tag{ACR14}
\]
A nonzero minor gives that many independent columns;
conversely Gaussian elimination on independent columns
gives that many pivot rows and a nonzero minor. This
proves the equality with rank. A sum of nonnegative real
terms vanishes exactly when all those minors vanish.
All $a_{ab}$ and the full OCF pivot denominators remain
in those minors. Formula ACR14 neither assigns generic
rank nor supplies a numerical zero test for an unevaluated
analytic coefficient.

\subsection{The exact retained residual quotient in biform coordinates}

The independent columns of $C_G$ stay independent modulo
$\operatorname{im}C_{\rm in}$ by ACR8. Hence the induced
map $\mathbb C^g\to V_A^\vee$ represented by $C_G$ is
injective. Its transpose is onto, and the invertible factors
in ACR11 prove that $\rho_G:V_A\to\mathbb C[S']_{<g}$
is onto. This proves its surjectivity directly from the
original coefficient map.

The annihilator calculation also gives
\[
 \begin{gathered}
 V_0:=\ker\rho_G=H_k^\perp\subseteq V_A,\qquad
 t_G=\dim\frac{J_k\cap H_k}{A\mathcal B_{k-8,k-8}},\\
 J_k^{\rm red}:=\operatorname{im}
       \left(J_k\longrightarrow\mathcal B_k/H_k\right),
 \qquad t_G=j-\dim J_k^{\rm red}.
 \end{gathered}\tag{ACR15}
\]
The first statement follows from the simultaneous vanishing
of the inner and exterior rows in ACR11 and ACR8. The
map $B_G$ in ACR12 has kernel precisely those exterior
linear combinations belonging to $J_k$. The direct sum in
ACR8 and the inclusion $A\mathcal B_{k-8,k-8}\subseteq J_k$
identify this kernel with the quotient in the first line.
Finally $\dim J_k=q-r$, while the inner space has
dimension $q'$. Applying rank-nullity to $J_k\to J_k^{\rm red}$
gives $(q-r)-\dim J_k^{\rm red}-q'=j-\dim J_k^{\rm red}$,
proving the second line.

This is the same invariant image as FMC7: its map is induced
by the identical original substitution $x=\mathsf Ht$,
and the inner quotient is the same $A\mathcal B$ quotient.
Factoring the explicit boundary monomial in ACR7 provides
the additional quotient map to $\mathcal B/H$, with its
kernel computed in ACR15; it does not replace that module
by a newly chosen invariant space.

For the original residual observation
$\Xi:V_A\twoheadrightarrow B_\Xi$ with kernel $K$, put
$W=\Xi(V_0)$. Define
\[
 A_G:\mathbb C[S']_{<g}\twoheadrightarrow B_\Xi/W,
 \qquad A_G(\rho_Gx)=[\Xi x].
 \tag{ACR16}
\]
Two lifts differ by $V_0$, so the map is well defined;
it is onto since both maps from $V_A$ are onto. Its kernel
is exactly $\rho_GK$: if $\Xi x\in W$, choose $y\in V_0$
with $\Xi y=\Xi x$; then $x-y\in K$ and has the same
$\rho_G$ value. The reverse containment is direct.
Consequently $\dim W=\dim J_k^{\rm red}$, while
$\dim(B_\Xi/W)=t_G$. In particular
\[
 \begin{array}{ll}
 d_a=d_b=8:&g=0,\quad t_G=0,\quad V_0=V_A,\quad W=B_\Xi;\\
 d_a=d_b=0:&g=\Delta,\quad H_k=\mathcal B_k,
              \quad t_G=j,\quad W=0.
 \end{array}\tag{ACR17}
\]
The first row follows from ACR5 and the zero target of
$\rho_G$; its full residual still has rank $j$. In the
second row $A_{\rm red}$ is a nonzero constant, so the
quotient defining $J_k^{\rm red}$ is zero and ACR15
gives the result. This second row classifies that exact
coefficient case; it does not assert that a given original
period attains it. Every intermediate case retains the
actual matrices ACR12--15 and the full original quotient
ACR16.

\subsection{A smaller invariant recursion evaluating every exceptional rank}

Here is a second exact computation of the remaining rank
which never uses the full symmetric power. In every diagonal
degree $d\geq0$ set
\[
 H_d=A_{\rm red}\mathcal B_{d-d_a,d-d_b},\qquad
 \mathcal W_d=\mathcal B_{d,d}/H_d,
 \qquad
 n_d=(d+1)^2-\max(d-d_a+1,0)\max(d-d_b+1,0).
 \tag{ACR18}
\]
Multiplication by $A_{\rm red}$ is injective, so this is
the exact dimension of the quotient. Also
$H_d\mathcal B_{e,e}\subseteq H_{d+e}$; therefore the
direct sum of these quotients is an actual graded algebra.

Choose the lexicographically largest nonzero coefficient
$a_*$ of the complete $A_{\rm red}$, at the biform
indices $(r_*,s_*)$. In the full degree-$d$ grid eliminate
the pivot rectangle
\[
 \{(r_*+i,s_*+h):0\leq i\leq d-d_a,
                             \ 0\leq h\leq d-d_b\}
 \tag{ACR19}
\]
in descending lexicographic order, subtracting at each
pivot its current coefficient divided by $a_*$ times
the corresponding monomial multiple of $A_{\rm red}$.
When a degree bound is negative the rectangle is empty.
Every other term in a subtracted polynomial is strictly
smaller in the additive lexicographic order. Earlier zero
pivots therefore stay zero. The result is a linear map
$N_d^{\rm red}$ to the complementary coefficient strip,
with its literal zero-insertion section $S_d^{\rm red}$.
Summing the subtractions gives a decomposition into the
strip and $H_d$. Their intersection is zero: the leading
term of every nonzero $A_{\rm red}b$ lies in the pivot
rectangle. Thus
\[
 \ker N_d^{\rm red}=H_d,\qquad
 N_d^{\rm red}S_d^{\rm red}=I_{n_d}.
 \tag{ACR20}
\]
All entries are finite rational expressions in the original
$a_{ab}$ with powers of the selected actual $a_*$ in their
denominators. Selecting the first nonzero coefficient covers
all support and leading-coefficient cases. If $A_{\rm red}$
is constant, this construction eliminates the full grid,
as it must.

Use the following fourteen actual invariant monomials:
\[
 \begin{array}{c|l}
 e&\mathcal G_e\\\hline
 2&x_1x_4,\ x_2x_3\\
 3&x_1^2x_3,\ x_1x_2^2,\ x_2x_4^2,\ x_3^2x_4\\
 4&x_1^3x_2,\ x_1x_3^3,\ x_2^3x_4,\ x_3x_4^3\\
 5&x_1^5,\ x_2^5,\ x_3^5,\ x_4^5.
 \end{array}\tag{ACR21}
\]
These generate every invariant polynomial, as follows
directly from the original weights $1,2,3,4$ modulo five.
An indecomposable invariant monomial has degree at most
five: among six partial sums of any five factor weights,
two agree and supply an invariant divisor. An opposite
pair of weights gives one of the displayed quadratics;
an indecomposable containing it is that quadratic. In the
absence of opposite pairs the support has at most two
variables. One variable gives a fifth power. The remaining
supports are $\{1,2\},\{1,3\},\{2,4\},\{3,4\}$;
the positive exponent solutions of total degree at most
five are respectively $(3,1),(1,2)$ for the first support
and $(2,1),(1,3)$ for the second, with the last two obtained
by multiplying weights by $2^{-1}$ and $3^{-1}$ modulo
five. These are precisely the displayed cubics and quartics.
Repeated factorization into indecomposables terminates by
decreasing total degree, proving the generator assertion.

For $f\in\mathcal G_e$, retain its complete biform
$\widetilde f=f(\mathsf Ht)$, whose coefficients are the
OCF13 finite multinomial expressions in the actual period
and unit matrix. Define
\[
 m_{f,d}^{\rm red}=N_d^{\rm red}M_{\widetilde f}S_{d-e}^{\rm red},
 \qquad
 J_0^{\rm red}=\operatorname{span}\{N_0^{\rm red}1\},\qquad
 J_d^{\rm red}=\sum_{\substack{f\in\mathcal G_e\\e\leq d}}
                     m_{f,d}^{\rm red}J_{d-e}^{\rm red}.
 \tag{ACR22}
\]
Negative degrees give zero spaces. Multiplying two
representatives differing by $H_{d-e}$ changes their
product by $H_d$, so ACR20 proves representative independence
of the displayed multiplication. Every invariant monomial
of positive degree factors by ACR21 into one generator
times an earlier invariant monomial. This proves both
inclusions in the recursion and hence its equality with
the original invariant image in every degree. The base
case includes the zero class of one when the quotient
algebra is zero.

To evaluate its rank, concatenate the indicated multiplication
matrices applied to bases of the preceding invariant images,
and choose independent columns by their exact minors. Its
rank is the maximum size of a nonzero minor, with rank zero
for the empty image. The selected columns supply the next
basis with their full original coefficient expressions.
This proves the algorithm by induction, without assigning
any unknown rank. All exceptional minors are covered by
choosing the first nonzero one in a fixed finite order.
For $d\geq8$, $n_d=(d_a+d_b)(d+1)-d_ad_b\leq16d+16$;
the candidate matrix has at most
$\sum_{e=2}^5|\mathcal G_e|n_{d-e}$ columns. Its row and
column dimensions are therefore bounded linearly in $d$.
Temporary full biform columns have $(d+1)^2$ entries, as
in OCF; no claim that all storage objects have linear
size is made. The exact output evaluates the residual
integer by the proved identity
\[
             \boxed{\ t_G=8k-32-\dim J_k^{\rm red}.\ }
 \tag{ACR23}
\]
Thus all coefficient cases on the original five-orbit
domain have explicit finite formulas, with the unchanged
original residual quotient ACR16. The separate literal
period support calculation can determine where its
first case ACR17 holds; its use then retains the full
complementary residual, rather than inferring that this
residual has vanished.

% END CURRENT ACR

\clearpage
% BEGIN CURRENT LBRE
\section{Exact exclusion of all four boundary rulings at the original period limit}
\label{sec:LBRE}
Retain the original hypotheses $0<\delta<1/2$, $\gamma>2$ and
$a\ne0$. In particular, no phase of the original amplitude is selected.
Set
\[
 x=\delta+i\gamma,\qquad y=\delta-i\gamma,\qquad
 \beta=2(\gamma^2-\delta^2)=-(x^2+y^2),\qquad
 t_{xy}=x/y,\quad r=a/\bar a,\quad R=r^2.
 \tag{LBRE1}
\]
All these ratios are defined, $|t_{xy}|=|r|=1$, and
$x^2-y^2=4i\delta\gamma\ne0$. Thus $t_{xy}^2\ne1$.
The root order is exactly $(x,y,-y,-x)$, the full unit is
$U=\operatorname{diag}(a,\bar a,\bar a,a)$, and $V$ evaluates
$1,w,w^2,w^3$ in this root order. For $j\in\{1,2,3,4\}$ put
$\rho=\zeta^{-j}$, where $\zeta=e^{2\pi i/5}$. Retain PCL4's
complete matrix
\[
 J_\infty=
 \begin{pmatrix}1&0&-\beta/3&0\\0&1&0&-2\beta/3\\
 0&0&1&0\\0&0&0&1\end{pmatrix},\qquad
 L=U^{-1}VJ_\infty^{-1}
 \operatorname{diag}(\rho,\rho^2,\rho^3,\rho^4)
 J_\infty V^{-1}U.
 \tag{LBRE2}
\]
The original form is $Q_0(t)=t_1t_4-t_2t_3$.

\subsection{The exact even and odd evaluation maps}
Define the following finite coefficients:
\[
 \begin{gathered}
 A=\frac{2t_{xy}^2-1}{3(t_{xy}^2-1)},\quad
 B=\frac{2-t_{xy}^2}{3(t_{xy}^2-1)},\quad h=\rho^3-\rho,\\
 e=\rho+hA,\quad f=-hA,\quad g=hB,\quad d=\rho-hB.
 \end{gathered}
 \tag{LBRE3}
\]
Here $A-B=1$. The notation $d$ in this section denotes precisely
this scalar and does not denote the original root displacement
$\delta$.
For any primary vector $t\in\mathbb C^4$ set
\[
 v_+(t)=\frac12\binom{t_1+t_4}{t_2+t_3},\qquad
 v_-(t)=\frac12\binom{t_1-t_4}{t_2-t_3},\qquad
 J_2=\operatorname{diag}(1,-1).
 \tag{LBRE4}
\]
These are explicit coordinate maps on the original primary space;
their inverse sends $(v_+,v_-)$ to
$(v_{+,1}+v_{-,1},v_{+,2}+v_{-,2},
v_{+,2}-v_{-,2},v_{+,1}-v_{-,1})$.
Consequently
\[
 Q_0(t)=v_+(t)^{\mathsf T}J_2v_+(t)
        -v_-(t)^{\mathsf T}J_2v_-(t).
 \tag{LBRE5}
\]
The complete matrices induced by $L$ in these coordinates are
\[
 v_+(Lt)=E v_+(t),\qquad v_-(Lt)=O v_-(t),\qquad
 E=\begin{pmatrix}e&f/r\\rg&d\end{pmatrix},\qquad
 O=\rho\begin{pmatrix}d&t_{xy} g/r\\rf/t_{xy}&e\end{pmatrix}.
 \tag{LBRE6}
\]
To prove the formulas, write an original polynomial as
$p_0+p_1w+p_2w^2+p_3w^3$. On its even coefficients the middle
conjugation in LBRE2 is
\[
 \begin{pmatrix}\rho&\frac{\beta}{3}(\rho^3-\rho)\\
 0&\rho^3\end{pmatrix},
 \tag{LBRE7}
\]
while on its odd coefficients it is
\[
 \begin{pmatrix}\rho^2&\frac{2\beta}{3}(\rho^4-\rho^2)\\
 0&\rho^4\end{pmatrix}.
 \tag{LBRE8}
\]
The even evaluation matrix is
$\left(\begin{smallmatrix}1&x^2\\1&y^2\end{smallmatrix}\right)$;
the odd evaluation matrix is its left product with
$\operatorname{diag}(x,y)$. Direct evaluation of LBRE7 gives
$\left(\begin{smallmatrix}e&f\\g&d\end{smallmatrix}\right)$,
because
\[
 \frac{x^2+\beta/3}{x^2-y^2}=A,\quad
 \frac{y^2+\beta/3}{x^2-y^2}=B.
 \tag{LBRE9}
\]
For LBRE8 the two corresponding fractions are $-B,-A$, since
$(x^2+2\beta/3)/(x^2-y^2)=-B$ and
$(y^2+2\beta/3)/(x^2-y^2)=-A$. Also
$\rho^4-\rho^2=\rho h$. Thus its evaluated matrix is
$\rho\left(\begin{smallmatrix}d&t_{xy} g\\f/t_{xy}&e\end{smallmatrix}\right)$.
Finally conjugation by the full unit acts on both parity spaces
by $\operatorname{diag}(a,\bar a)$, proving LBRE6 including every
amplitude factor.

\subsection{The four complete quadratic restrictions}
Define
\[
 \begin{aligned}
 n_1&=e^2-\rho^2d^2,& c_1&=\rho^2f^2/t_{xy}^2-g^2,\\
 n_2&=f^2-\rho^2t_{xy}^2g^2,& c_2&=\rho^2e^2-d^2,\\
 \mathcal A&=n_1+Rc_1,&
 \mathcal C&=(n_2+Rc_2)/R,\\
 m_E&=ef/r-rgd,&
 m_O&=\rho^2(dt_{xy} g/r-rfe/t_{xy}),\\
 \mathcal B_-&=m_E-m_O,&\mathcal B_+&=m_E+m_O.
 \end{aligned}\tag{LBRE10}
\]
For the affine Segre vector $t(Z,W)=(ZW,Z,W,1)^{\mathsf T}$,
write $q_j(Z,W)=Q_0(Lt(Z,W))$. The restrictions below also give
the exact homogeneous restrictions at infinity:
\[
 \begin{array}{ll}
 Z=0:&Q_0(L(0,0,W,1)^{\mathsf T})
   =\tfrac14(\mathcal A+2\mathcal B_-W+\mathcal CW^2),\\[2pt]
 Z=\infty:&Q_0(L(W,1,0,0)^{\mathsf T})
   =\tfrac14(\mathcal AW^2+2\mathcal B_-W+\mathcal C),\\[2pt]
 W=0:&Q_0(L(0,Z,0,1)^{\mathsf T})
   =\tfrac14(\mathcal A+2\mathcal B_+Z+\mathcal CZ^2),\\[2pt]
 W=\infty:&Q_0(L(Z,0,1,0)^{\mathsf T})
   =\tfrac14(\mathcal AZ^2+2\mathcal B_+Z+\mathcal C).
 \end{array}\tag{LBRE11}
\]
For example the two parity vectors for $(W,1,0,0)$ both equal
$(W,1)^{\mathsf T}/2$, so LBRE5--6 give the second row directly.
For $(0,0,W,1)$ they are $(1,W)^{\mathsf T}/2$ and its negative,
giving the first row. For $(Z,0,1,0)$ they are
$(Z,1)^{\mathsf T}/2$ and $(Z,-1)^{\mathsf T}/2$, giving the
last row. For $(0,Z,0,1)$ they are $(1,Z)^{\mathsf T}/2$ and
$(-1,Z)^{\mathsf T}/2$, giving the third row. These substitutions
also prove every mixed-coefficient sign and the factor $1/4$.
In particular any identically zero boundary restriction would
require $\mathcal A=\mathcal C=0$.

\subsection{A nonzero exact eliminant for every original amplitude}
Put
\[
 q=t_{xy}^2+t_{xy}^{-2}\in[-2,2],\qquad
 c=\rho^2+\rho^3\in[-2,2],\qquad
 F(q,c)=5(c-1)q^2+(20c+25)q-61c-173.
 \tag{LBRE12}
\]
The first interval follows from $|t_{xy}|=1$. For the second,
$\rho^3=\overline{\rho^2}$ because $\rho^5=1$.
Direct algebra in LBRE3 and LBRE10 gives the exact identity
\[
 n_1c_2-c_1n_2
   =\frac{t_{xy}^2}{81(t_{xy}^2-1)^2}F(q,c).
 \tag{LBRE13}
\]
Here is a finite expansion verifying the identity without any
numerical specialization. Before using the fifth-root relation,
the left side has denominator $81t_{xy}^2(t_{xy}^2-1)^2$ and
numerator
\[
 \begin{split}
 -\rho^6(\rho^2-1)^2\bigl[
 &(t_{xy}^8+1)(\rho^4-2\rho^2+1)\\
 &+(t_{xy}^6+t_{xy}^2)(13\rho^4+19\rho^2+13)\\
 &-t_{xy}^4(57\rho^4+120\rho^2+57)\bigr].
 \end{split}\tag{LBRE14}
\]
Use $1+\rho+\rho^2+\rho^3+\rho^4=0$ to replace this numerator by
\[
 5(c-1)(t_{xy}^8+1)+(20c+25)(t_{xy}^6+t_{xy}^2)
                         -(51c+183)t_{xy}^4.
 \tag{LBRE15}
\]
The identities
$(t_{xy}^8+1)/t_{xy}^4=q^2-2$ and
$(t_{xy}^6+t_{xy}^2)/t_{xy}^4=q$ now give LBRE13.
All denominators used here are nonzero by LBRE1.

For $-2\le q\le2$, the polynomial $5q^2+20q-61$ has maximum
$-1$: its derivative is $10q+20\ge0$ on that interval and its
value at $2$ is $-1$. Since $c\ge-2$, it follows that
\[
 \begin{split}
 F(q,c)
 &=c(5q^2+20q-61)-5q^2+25q-173\\
 &\le -2(5q^2+20q-61)-5q^2+25q-173\\
 &=-15(q+\tfrac12)^2-\tfrac{189}{4}<0.
 \end{split}\tag{LBRE16}
\]
Thus LBRE13 is nonzero. If $\mathcal A=\mathcal C=0$, then
$n_1+Rc_1=n_2+Rc_2=0$ since $R\ne0$. Multiplying the first
equation by $c_2$, the second by $c_1$, and subtracting would give
$n_1c_2-c_1n_2=0$, a contradiction. Consequently
\[
 (\mathcal A,\mathcal C)\ne(0,0)
 \quad\text{for every allowed }\delta,\gamma,a
 \text{ and every }j\in\{1,2,3,4\}.
 \tag{LBRE17}
\]
There is also a fully explicit coefficient margin. Indeed
$t_{xy}-t_{xy}^{-1}=4i\delta\gamma/(\delta^2+\gamma^2)$, and therefore
\[
 \begin{split}
 \mathcal D_j:=n_1c_2-c_1n_2
  &=-\frac{(\delta^2+\gamma^2)^2}{1296\delta^2\gamma^2}F(q,c)
    \ge\frac{7(\delta^2+\gamma^2)^2}{192\delta^2\gamma^2}>0,\\
 \max(|\mathcal A|,|\mathcal C|)
  &\ge\frac{\mathcal D_j}{|c_1|+|c_2|}>0.
 \end{split}\tag{LBRE18}
\]
The last denominator is positive because simultaneous vanishing of
$c_1,c_2$ would contradict $\mathcal D_j>0$. For its proof use
$\mathcal D_j=\mathcal Ac_2-R\mathcal Cc_1$ and $|R|=1$, then
apply the triangle inequality. Thus on every boundary at least one
of its two endpoint coefficients has modulus at least
$\mathcal D_j/[4(|c_1|+|c_2|)]$, with the exact scalar $1/4$ from
LBRE11 retained.
Each polynomial in LBRE11 is therefore nonzero. In homogeneous
Segre coordinates, none of $z_0,z_1,w_0,w_1$ divides the actual
limiting biform $Q_0(Lt)$. This conclusion includes every amplitude
phase and all eigenvalue-collision phases; its proof never excludes
such a phase. It also gives the corresponding statement for the
full limiting conductor product: restriction to any one boundary
is the product of its four nonzero quadratic restrictions, which
is nonzero because the polynomial ring over $\mathbb C$ has no
zero divisors. Hence none of the four homogeneous boundary
coordinates divides $A_0$ of PCL8.

Finally the biform coefficients are entire in the original inverse
period $z$ by PCL8--11. For each of the finitely many pairs consisting
of an orbit index and a boundary, LBRE11 and LBRE17 supply a
nonzero coefficient at $z=0$. That exact coefficient function is
continuous, so it stays nonzero on a disk about $0$. The minimum
of these finitely many positive radii is positive. Therefore all
four boundary exclusions hold simultaneously on some neighborhood
of the original inverse-period limit, retaining the original
period and unit throughout. No explicit size for this last
neighborhood is asserted here.

% END CURRENT LBRE

\clearpage
% BEGIN CURRENT ASR
\section{Actual boundary multiplicities and a finite full-support radius}
\label{sec:ASR}

Retain the original quartet, unit, period matrices, logarithm branch and
cyclic orientation of PCL1--11, PQO1--25 and RPC1--26. In particular,
$r_1=\delta+i\gamma$, $r_2=\delta-i\gamma$, $r_3=-r_2$,
$r_4=-r_1$, $0<\delta<1/2$, $\gamma>2$ and
$U=\operatorname{diag}(a,\bar a,\bar a,a)$ with the full actual
$a\ne0$. Put $z=u^{-1}$ and retain
\[
\begin{gathered}
 L_j(z)=U^{-1}V\mathcal R(z)^{-1}D^{-j}\mathcal R(z)V^{-1}U,
 \\
 q_j(z;Z,W)=Q_0(L_j(z)(ZW,Z,W,1)^{\mathsf T}),
 \quad A_z=\prod_{j=1}^4q_j(z;Z,W).
\end{gathered}
 \tag{ASR1}
\]
These are exact identities in the original period map by PCL9.
Every entry of $L_j$ is entire, since $\det\mathcal R=1$ and
$\mathcal R^{-1}=\operatorname{adj}\mathcal R$. In particular all
coefficients below are the entire functions evaluated by PCL2 and
PCL10, with their original signs and ordered-pair multiplicities.

\subsection{Every finite-period boundary multiplicity}
Let $S_0$ be the symmetric matrix with entries
$(S_0)_{14}=(S_0)_{41}=1/2$,
$(S_0)_{23}=(S_0)_{32}=-1/2$ and all other entries zero, and set
$S_j(z)=L_j(z)^{\mathsf T}S_0L_j(z)$. Thus
\[
\begin{split}
 q_j={}&(S_j)_{11}Z^2W^2+2(S_j)_{12}Z^2W+(S_j)_{22}Z^2
       +2(S_j)_{13}ZW^2\\
 &+2((S_j)_{14}+(S_j)_{23})ZW+2(S_j)_{24}Z
       +(S_j)_{33}W^2+2(S_j)_{34}W+(S_j)_{44}.
\end{split}\tag{ASR2}
\]
This follows by substituting all four coordinates into $t^{\mathsf T}S_jt$;
each off-diagonal term occurs twice. Introduce the following six
three-entry vectors, with the displayed order of their coefficients:
\[
\begin{gathered}
 E_{Z,0,j}=((S_j)_{33},2(S_j)_{34},(S_j)_{44}),\\
 E_{Z,\infty,j}=((S_j)_{11},2(S_j)_{12},(S_j)_{22}),\\
 E_{W,0,j}=((S_j)_{22},2(S_j)_{24},(S_j)_{44}),\\
 E_{W,\infty,j}=((S_j)_{11},2(S_j)_{13},(S_j)_{33}),\\
 M_{Z,j}=(2(S_j)_{13},2((S_j)_{14}+(S_j)_{23}),2(S_j)_{24}),\\
 M_{W,j}=(2(S_j)_{12},2((S_j)_{14}+(S_j)_{23}),2(S_j)_{34}).
\end{gathered}\tag{ASR3}
\]
On the original five-orbit locus, for $X=Z,W$ and $b=0,\infty$
define the following integer by
literal vanishing of these actual coefficient vectors:
\[
 \nu_{X,b,j}=
 \begin{cases}
 0,&E_{X,b,j}\ne0,\\
 1,&E_{X,b,j}=0,\ M_{X,j}\ne0,\\
 2,&E_{X,b,j}=0,\ M_{X,j}=0.
 \end{cases}
 \tag{ASR4}
\]
On the original five-orbit locus, $q_j$ is nonzero. To verify this
without a geometric identification, ASR2 is zero exactly when its
nine displayed coefficients vanish. These equations say precisely
$S_j=cS_0$ for a scalar $c$, so $Q_0(L_jt)=cQ_0(t)$ in the four
independent primary variables. Conjugation by the original period
matrix identifies this with projective equality of orbit quadrics,
which is excluded for $1\le j\le4$ on that locus. Since $S_j$
is invertible, the scalar in such an equality could not be zero.

The homogeneous form of $q_j$ has bidegree $(2,2)$ with
$t=(z_0w_0,z_0w_1,z_1w_0,z_1w_1)$. For $X=Z=z_0/z_1$,
its three successive powers of $z_0$ have coefficient vectors
$E_{Z,0,j},M_{Z,j},E_{Z,\infty,j}$. This proves that ASR4 is
exactly its order along $z_0=0$. Reversing the powers proves
the assertion for $z_1=0$. The two identical calculations with
$W=w_0/w_1$ prove the remaining two orders. In the last case
of ASR4 the opposite vector is nonzero, since $q_j\ne0$.

Write $a_-,a_+$ for the minimum and maximum $Z$ exponents in
$A_z$, and $b_-,b_+$ for its minimum and maximum $W$ exponents.
The exact actual support extrema are therefore
\[
\boxed{\begin{aligned}
 a_- =\sum_{j=1}^4\nu_{Z,0,j},\qquad
 &a_+ =8-\sum_{j=1}^4\nu_{Z,\infty,j},\\
 b_- =\sum_{j=1}^4\nu_{W,0,j},\qquad
 &b_+ =8-\sum_{j=1}^4\nu_{W,\infty,j}.
\end{aligned}}
 \tag{ASR5}
\]
Indeed the lowest nonzero $Z$ coefficient of a product is the
product of its lowest coefficients in the integral domain
$\mathbb C[W]$, and the same is true of the highest coefficients.
Thus neither extreme can cancel. The argument over $\mathbb C[Z]$
gives both $W$ extremes. Equivalently the complete boundary divisor
and the remaining bihomogeneous factor are specified by
\[
 M_{\partial}=z_0^{a_-}z_1^{8-a_+}w_0^{b_-}w_1^{8-b_+},
 \qquad A=M_{\partial}A_{\mathrm{red}},\qquad
 (d_a,d_b)=(a_+-a_-,b_+-b_-).
 \tag{ASR6}
\]
This retains the order at each of the four original boundary rulings.

\subsection{The exact symplectic map underlying the coefficient tests}
Put $B=V^{-1}U$ and retain RPC5's alternating matrix $J$.
The original alternating form on primary coordinates is
\[
\begin{gathered}
 \Omega_{\mathrm{prim}}=B^{\mathsf T}JB,\\
 (\Omega_{\mathrm{prim}})_{14}=\frac{a^2}{8i\delta\gamma r_1},
 \qquad
 (\Omega_{\mathrm{prim}})_{23}=-\frac{\bar a^2}{8i\delta\gamma r_2},
\end{gathered}
 \tag{ASR7}
\]
with the negative transposed entries and all others zero. In fact
RPC4 evaluates this form as
$\sum_i a_i^2t_iv_{5-i}/H'(r_i)$, where
$(a_1,a_2,a_3,a_4)=(a,\bar a,\bar a,a)$,
$H'(r_1)=8i\delta\gamma r_1$ and
$H'(r_2)=-8i\delta\gamma r_2$. This proves every entry and its sign.
Using $P_0=G_u\mathcal R(z)$ in RPC9 gives
$\mathcal R(z)^{-\mathsf T}J\mathcal R(z)^{-1}=J|_{\beta=0}$.
The diagonal $D$ preserves this last alternating matrix by its paired
weights. Conjugating the equality by $\mathcal R(z)B$ proves
\[
 L_j(z)^{\mathsf T}\Omega_{\mathrm{prim}}L_j(z)
       =\Omega_{\mathrm{prim}}.
 \tag{ASR8}
\]

Let $s=r_1/a^2$, retain its literal complex conjugate
$\bar s=r_2/\bar a^2$, and define
$T=\operatorname{diag}(s,\bar s,-\bar s,-s)$.
The entries of ASR7 prove the exact matrix identity
\[
 S_0=-4i\delta\gamma\,\Omega_{\mathrm{prim}}T.
 \tag{ASR9}
\]
Each original boundary plane is one of
$E_{12},E_{34},E_{13},E_{24}$, where $E_{ab}$ is spanned by
the two indicated original coordinate vectors. Each is Lagrangian
for ASR7: its restriction vanishes, its dimension is two, and the
ambient alternating matrix is invertible. Equation ASR8 shows that
its image $F=L_jE_{ab}$ is also Lagrangian. Polarization of ASR9
now proves the exact equivalence
\[
 Q_0|_{L_jE_{ab}}=0
 \quad\Longleftrightarrow\quad
 T L_jE_{ab}\subset L_jE_{ab}
 \quad\Longleftrightarrow\quad
 \operatorname{rank}[L_jE_{ab},TL_jE_{ab}]=2.
 \tag{ASR10}
\]
Here the last matrix has the two columns of $L_j$ indexed by $a,b$
followed by their images under $T$. For the first equivalence,
polarization says $\Omega_{\mathrm{prim}}(v,Tw)=0$ for all $v,w\in F$;
since $F^{\perp_{\Omega}}=F$, this is exactly $TF\subset F$.
For the last equivalence the first two columns are independent.
Thus all its three-by-three minors vanish precisely in that case.
These are the same boundary-vanishing tests as the corresponding
three entries of ASR3, with the exact map between the two tests proved.
No eigenvalue is cancelled when $s$ is real or purely imaginary.

\subsection{An explicit radius from the actual limiting corner calculation}
The complete limiting calculation LBRE1--18 proves, for every
$1\le j\le4$ and every original nonzero $a$, that
\[
\begin{gathered}
 d_{j,1}(0)=d_{j,4}(0),\qquad d_{j,2}(0)=d_{j,3}(0),\\
 m_j:=\max\{|d_{j,1}(0)|,|d_{j,2}(0)|\}>0,
 \quad d_{j,i}(z):=(S_j(z))_{ii}.
\end{gathered}
 \tag{ASR11}
\]
Indeed LBRE11 gives
$d_{j,1}(0)=d_{j,4}(0)=\mathcal A/4$ and
$d_{j,2}(0)=d_{j,3}(0)=\mathcal C/4$ with the complete original
amplitude in LBRE10. LBRE13--18 prove the positive lower bound
\[
 m_j\ge\frac{\mathcal D_j}{4(|c_1|+|c_2|)}>0,
 \qquad \mathcal D_j\ge
       \frac{7(\delta^2+\gamma^2)^2}{192\delta^2\gamma^2},
\]
where $c_1,c_2,\mathcal D_j$ here are exactly LBRE10 and LBRE18's
finite coefficients for this orbit index. That proof evaluates the
literal $J_\infty$ in PCL4 and the original $r_1/r_2$ and
$\zeta^{-j}$. The two equalities also follow directly from
PCL22--23 at $z=0$, because the opposite-root permutation fixes
$Q_0$ and conjugates $L_j(0)$ to itself.

Expand each of these actual entire diagonal entries as
$d_{j,i}(z)=\sum_{n\ge0}d_{j,i,n}z^n$ by PCL2 and PCL10, and
define the explicitly specified finite constant
\[
 C_j=\sum_{n\ge1}\max_{1\le i\le4}|d_{j,i,n}|,
 \qquad
 R_{\mathrm{supp}}=
 \max\left\{R_*,1,\frac{2C_1}{m_1},\frac{2C_2}{m_2},
                       \frac{2C_3}{m_3},\frac{2C_4}{m_4}\right\}.
 \tag{ASR12}
\]
Each series is finite in value: the four entire Taylor series converge
absolutely at $z=1$, and the maximum is bounded by their sum of absolute
values. Every denominator is strictly positive by ASR11. If $C_j=0$,
its displayed quotient is zero; no division by $C_j$ is performed.
For $|z|\le1$, summing the same series gives
\[
 |d_{j,i}(z)-d_{j,i}(0)|\le C_j|z|\qquad(1\le i\le4).
 \tag{ASR13}
\]
Choose $i_j=1$ when $|d_{j,1}(0)|\ge|d_{j,2}(0)|$, and $i_j=2$
otherwise. This deterministic choice selects a coefficient already
present in the actual period calculation. For $|u|\ge R_{\mathrm{supp}}$,
ASR11--13 give
\[
 |d_{j,i_j}(u^{-1})|\ge\frac{m_j}{2}>0,
 \qquad |d_{j,5-i_j}(u^{-1})|\ge\frac{m_j}{2}>0.
 \tag{ASR14}
\]
Indeed both initial values are equal with modulus $m_j$ and their
errors are at most $C_j/|u|\le m_j/2$. Each of the four boundary
planes in ASR10 contains exactly one of these two opposite coordinate
vectors. Thus its boundary restriction has a nonzero diagonal
coefficient, so all four vectors $E_{X,b,j}$ in ASR3 are nonzero.
This proves, at every such original $u$ on every logarithm branch,
\[
\boxed{\quad
 a_-=b_-=0,\qquad a_+=b_+=8,\qquad
 (d_a,d_b)=(8,8),\qquad M_{\partial}=1,\quad A_{\mathrm{red}}=A.
 \quad}\tag{ASR15}
\]
Branch independence follows from the exact original-period conjugation
PCL9, which leaves $L_j$ single-valued in $u^{-1}$. The radius includes
$R_*$, so all these points are on the proved original five-orbit locus.
The estimate has introduced no new source order and no freely specified
amplitude or period coefficient.

At arbitrary nonzero original periods, each boundary failure is the
simultaneous zero set of its three actual entire entries in ASR3.
ASR11 supplies a nonzero value at $z=0$ for at least one of those
entries. Such an entry has isolated zeros: at any zero its convergent
Taylor series has a first nonzero term, and factoring that power
leaves a nonvanishing factor in a sufficiently small disk. Thus each
boundary restriction-vanishing set is discrete and closed in the finite $z$ plane,
and its finite union over the sixteen boundary/index pairs is also
discrete and closed. This boundary restriction-vanishing set in the nonzero $u$ plane
is contained in $|u|<R_{\mathrm{supp}}$ and has no accumulation
point other than possibly $u=0$. ASR3--5 retain every multiplicity
on its intersection with the original five-orbit locus. At an orbit-collapse
period, ASR2's nine coefficients detect an identically zero factor;
no finite ruling multiplicity is assigned to such a factor.

\subsection{Consequence for the actual common factor and residual}
Retain the common-factor and grid maps of ACR1--23,
with original $k=4\ell+1\ge9$ and both unchanged source orders $1,k$.
Its exact dimension formula and its actual invariant-image quotient are
\[
\begin{gathered}
 g_A=(k+1-d_a)(k+1-d_b)-(k-7)^2,\\
 t_G=\dim\frac{J_k\cap A_{\mathrm{red}}
                       \mathcal B_{k-d_a,k-d_b}}
                      {A\mathcal B_{k-8,k-8}}.
\end{gathered}
 \tag{ASR16}
\]
Here $\mathcal B_{p,q}$ is the space of bihomogeneous polynomials of
bidegree $(p,q)$ in the two original pairs $(z_0,z_1),(w_0,w_1)$,
and $J_k$ is the original invariant image. Its exact map from the
affine coefficient space used in ASR1--5 is
\[
 f(Z,W)\longmapsto z_1^p w_1^q f(z_0/z_1,w_0/w_1),
 \qquad F\longmapsto F(Z,1,W,1).
\]
The respective bounded degrees ensure that the first expression is
a polynomial. Expansion in its monomials proves that these two maps
are inverse and that multiplication is preserved with the displayed
bidegrees. ACR5 and ACR15 prove the two formulas in ASR16;
ACR8--9 prove that the denominator
is a subspace of that numerator. At every $|u|\ge R_{\mathrm{supp}}$,
ASR15 first makes the two products in $g_A$ equal. It then makes the
numerator of the displayed quotient a subspace of its denominator;
the opposite inclusion was already proved by the exact receiving map.
Consequently
\[
 g_A=t_G=0\qquad(|u|\ge R_{\mathrm{supp}},\quad k=4\ell+1\ge9).
 \tag{ASR17}
\]
This removes precisely the common boundary factor calculated by that
map. The original remaining residual subspace is still $W=B_\Xi$,
and the receiving calculation's full residual rank remains
\[
 j=\bigl((k+1)^2-(k-7)^2\bigr)-(8k-16)=8k-32.
\]
Here the first parenthesis is the original conductor quotient dimension
from its injective multiplication by the nonzero biform $A$, and
$8k-16$ is the exact original invariant-image kernel dimension of
OQX and OCF on this five-orbit locus. Neither that factor nor those residual directions have
been set to zero. For periods below this explicit support radius that
remain on the original five-orbit locus, ASR3--6 give all actual boundary
losses without assuming full widths.

\subsection{Only finitely many support exceptions on the original exterior domain}
Define the exact coefficient set and its corresponding original periods by
\[
\begin{gathered}
 \mathcal Z_{\partial}^{*}=
 \left\{z:0<|z|\le R_*^{-1},\quad
 E_{X,b,j}(z)=0\text{ for some }X\in\{Z,W\},\
 b\in\{0,\infty\},\ 1\le j\le4\right\},\\
 \mathcal U_{\partial}^{*}=\{z^{-1}:z\in\mathcal Z_{\partial}^{*}\}.
\end{gathered}
\]
For each of the sixteen boundary/index pairs, ASR11 selects one
diagonal coefficient among the three entries of $E_{X,b,j}$ that is
nonzero at $z=0$. This is an entire function with isolated zeros by
the local Taylor argument given above. It has finitely many zeros
on the compact disk $|z|\le R_*^{-1}$: an infinite set would have an
accumulation point in that disk, contradicting the proved isolation.
The common zero set of the three entries is a subset of this finite
set. Taking the union of the sixteen finite sets proves finiteness
of $\mathcal Z_{\partial}^{*}$ and hence of $\mathcal U_{\partial}^{*}$.
None of these zeros is zero, by the selected nonzero value there.
ASR14 excludes every inverse period with $|u|\ge R_{\mathrm{supp}}$.
On the whole domain $|u|\ge R_*$, the factors are nonzero by the
original five-orbit theorem, so ASR4--5 show that this is precisely
the set where at least one of the four support extrema loses its
full value. Thus
\[
\boxed{\begin{gathered}
 \mathcal U_{\partial}^{*}\text{ is finite},\qquad
 \mathcal U_{\partial}^{*}\subseteq
       \{u:R_*\le|u|<R_{\mathrm{supp}}\},\\
 |u|\ge R_*,\ u\notin\mathcal U_{\partial}^{*}
 \quad\Longrightarrow\quad
 (a_-,a_+,b_-,b_+)=(0,8,0,8),\\
 g_A=0,\qquad t_G=0,\qquad W=B_\Xi,\qquad j=8k-32
 \quad(k=4\ell+1\ge9).
\end{gathered}}\tag{ASR18}
\]
At the finite exceptional periods every boundary multiplicity is
given by ASR3--5, and the unchanged matrices ACR12--15 evaluate the
remaining common-factor map and its residual quotient. No statement
that this finite set is empty is needed for any displayed conclusion.

% END CURRENT ASR

\clearpage
% BEGIN CURRENT PBR
% BEGIN PBR COMPLETE RECEIVING PROOF
\section{The two exact dual subquotients in the original receiver}
\label{sec:PBR}

Fix the original simple quartet, full unit, actual five-orbit period
and logarithm branch. All coefficients of the actual conductor $A$
and all original roots are retained. In this section
\[
\begin{gathered}
 k=4l+1\ge9,\quad q=(k+1)^2,\quad q'=(k-7)^2,\\
 \Delta=q-q'=16k-48,\quad r_K=8k-16,\quad j=8k-32,\\
 E_A(w)=\sum_{a,b=0}^8a_{ab}e^{b_{ab}w},\quad
 v=\operatorname{ord}_0E_A\le80,\quad
 \mu_v=E_A^{(v)}(0)\ne0,\quad \alpha_A=v!/\mu_v,\\
 \mathfrak m=\Delta-v,\quad c'=k/2-4,\quad
 s\in\{1,k\},\quad
 \mathcal N=\{q-1,q,2q-1,2q\},\\
 (\sigma_{q-1},\sigma_q,\sigma_{2q-1},\sigma_{2q})=(1,1,-1,-1).
\end{gathered}\tag{PBR1}
\]
The original relation quotient has Gram $H_{Y,N}^{(s)}$ in the
fixed FGC frame, and $Z_k$ is its actual FGC10 kernel-column matrix.
Here $K=\ker\Xi$ in the original primary space $V_A$,
$L=\C[S]_{<v}\subset V_A$, $K^+=K+L$, $t_0=\dim(K\cap L)$,
and $t=r_K-t_0$ is the number of columns of $Z_k$.
The symbol $r_K$ denotes this rank parameter, whereas $\rho$ below
denotes an analytic radius. No independent source order is introduced.

Retain the complete DNE coefficient frame $Z_{\exp}$, the functions
$F_a=\mathcal N_a/E_A$, and
$\ell_a(g)=\Lambda_{F_a}(\chi'g)$.
Put $(A_\ell)_{ba}=\ell_a(S'^b)$ for $0\le b<\mathfrak m$.
The matrix $A_\ell$ is invertible by DNE4. In the original Gamma
orthonormal frame $\phi_{n,s,c'}$, $0\le n\le N-v$, whose squared
monic norms are $h_n=M_sn!(s/2)_n$, $M_s=(2\pi)^{s/2}$, define
\[
\begin{gathered}
 T_{na}=\Lambda_{F_a}(\phi_{n,s,c'}),\qquad
 V_{n,(i,j)}=\phi_{n,s,c'}(\beta'_{ij}),\\
 \mathcal D_N^{(s)}
 =T^*[I-V(V^*V)^{-1}V^*]T
 =A_\ell^*\overline{(H_{Y,N}^{(s)})^{-1}}A_\ell,\\
 G_{\exp}=Z_{\exp}^*
 [I-C^{\mathsf T}(\bar C C^{\mathsf T})^{-1}\bar C]Z_{\exp}>0,\\
 B_K=Z_k^{\mathsf T}A_\ell,\qquad
 D_{K,N}=[B_K(\mathcal D_N^{(s)})^{-1}B_K^*]^{-1}
       =\overline{(\widehat H_{K,N}^{(s)})^{-1}},\qquad
 G_K=[B_KG_{\exp}^{-1}B_K^*]^{-1}.
\end{gathered}\tag{PBR2}
\]
The onto map $B_K$ is the exact restriction of the complex-linear
dual to the actual kernel columns. The conjugation in PBR2 follows
by substituting the displayed $A_\ell$ identity and
$B_K=Z_k^{\mathsf T}A_\ell$; it gives
$B_K\mathcal D_N^{-1}B_K^*
=Z_k^{\mathsf T}\overline{H_{Y,N}}\bar Z_k
=\overline{Z_k^*H_{Y,N}Z_k}$ before inversion.
The first equality for $\mathcal D_N$ minimizes the full extension
over every lower evaluation functional. Its bracket is the
orthogonal projection off their entire column span. Thus the full
conditional covariance and the coefficient reference are both
specified; neither is replaced by the other.

For the complementary quotient choose exactly the first invertible
$t$-row minor of $Z_k$, with row set $\mathcal I$ and complementary
set $\mathcal J$, and put
\[
\begin{gathered}
 j_0=\mathfrak m-t=j-v+t_0,\qquad
 P_Y=E_{\mathcal J}^{\mathsf T}
 -E_{\mathcal J}^{\mathsf T}Z_k
 (E_{\mathcal I}^{\mathsf T}Z_k)^{-1}E_{\mathcal I}^{\mathsf T},\\
 J_0=A_\ell^{-1}P_Y^{\mathsf T},\qquad
 D_{0,N}=J_0^*\mathcal D_N^{(s)}J_0,\quad G_0=J_0^*G_{\exp}J_0,\\
 Q_{0,N}=[P_Y(H_{Y,N}^{(s)})^{-1}P_Y^*]^{-1},\qquad
 D_{0,N}=\overline{Q_{0,N}^{-1}}.
\end{gathered}\tag{PBR3}
\]
Empty rank uses the identity quotient. Multiplication shows
$P_YZ_k=0$ and $P_YE_{\mathcal J}=I$; hence
$\ker P_Y=\operatorname{im}Z_k$. Also $B_KJ_0=0$ and
$J_0$ has $j_0$ independent columns, proving
$\operatorname{im}J_0=\ker B_K$. Substitution in PBR2 gives
$D_{0,N}=\bar P_Y\overline{H_{Y,N}^{-1}}P_Y^{\mathsf T}$,
which proves the last identity in PBR3. These are the actual
kernel quotient and its full dual annihilator.

\subsection{Finite positive-spectrum allowance and every small eigenvalue}

For clarity the finite allowance used below is the complete PEC13
quantity, with no hidden replacement by an asymptotic zero.
For $0<\rho<1$ with $E_A(-i\arctan z)\ne0$ on $|z|=\rho$, put
$m_\rho=\min_{|z|=\rho}|E_A(-i\arctan z)|>0$.
The actual nonzero symbol zeros inside this image have orders
$h_\omega$. Form the full jet matrix
$(J_{\rho,k})_{(\omega,b),a}=\mathcal N_a^{(b)}(\omega)/b!$,
$0\le b<h_\omega$, and put $c_{\rho,k}=\operatorname{rank}J_{\rho,k}$.
With $R=|k/2|+k\sqrt{\delta^2+\gamma^2}$ and $d=N-v$, retain
\[
\begin{gathered}
 C_{\rho,N,s}=
 \frac{q\,e^{2(R+|c'|)\operatorname{arctanh}\rho}
 (1-\rho^2)^{-s/2}}{M_sm_\rho^2}
 \sum_{n=0}^{d}\frac{n!}{(s/2)_n}\rho^{-2n},\\
 C_{0,N,s}=q\mathcal B_{d,s,k,A}^{\,2},\qquad
 \varepsilon_{N,s}^{\vee}
 =\inf_{\rho\ {\rm admissible}}
 \{c_{\rho,k}\log^+C_{0,N,s}
       +(\mathfrak m-c_{\rho,k})\log^+C_{\rho,N,s}\},\\
 E_{k,s}^{\vee}=\varepsilon_{2q-1,s}^{\vee}
                       +\varepsilon_{2q,s}^{\vee}=o(kq).
\end{gathered}\tag{PBR4}
\]
Here $\mathcal B_{d,s,k,A}$ is the full finite expression DNE11--13,
including its chosen zero-free radius, Taylor coefficient and Gamma
mass. The full DNE proof is included as a provider, not assumed.
Pole cancellation on $\ker J_{\rho,k}$ and Cauchy's coefficient
integral give $C_{\rho,N,s}$ exactly: the $n$th squared coefficient
is weighted by $n!/(M_s(s/2)_n)$. Minimization over the entire
conductor coefficient fibre preserves these cancellations. The
full DNE bound gives $C_{0,N,s}$ on all vectors. At most
$c_{\rho,k}$ eigenvalues can exceed $C_{\rho,N,s}$, by intersecting
the span of $c_{\rho,k}+1$ such eigenvectors with the cancellation
subspace. This proves the positive logarithmic sum bound PBR4.
PEC17--18 proves the same bound for the fixed quotient and restriction
in PBR2--3, by the variational principle applied respectively to the
inverse compression and the compression. The exact reference metrics
$G_K,G_0$ are retained in both applications.

For completeness, at each fixed admissible radius the symbol has
only finitely many zeros in the compact disk, so $c_{\rho,k}$ is
bounded independently of $k$. The full finite constants above and
DNE15 give
$\limsup\varepsilon_{N,s}^{\vee}/(kq)\le64\log(1/\rho)$.
Admissible radii approach one because the zeros are isolated.
Taking that limit after the fixed-radius bound proves the displayed
$o(kq)$ for the original two source orders. This bounds every
forward exterior degree by $\varepsilon_{N,s}^{\vee}/2$ in logarithmic
norm, by multiplying subsets of the relative eigenvalues.
For the inverse map its squared singular values are the reciprocals.
They are retained in the following exact formulas.

For $X=K,0$, let $\mu_{X,a,N}$ be every eigenvalue of
$G_X^{-1/2}D_{X,N}G_X^{-1/2}$ and define
\[
\begin{gathered}
 p_{X,N}=\sum_a\log^+\mu_{X,a,N},\qquad
 n_{X,N}=\sum_a\log^+(1/\mu_{X,a,N}),\\
 A_{X,k,s}=n_{X,q-1}+n_{X,q}-n_{X,2q-1}-n_{X,2q},\\
 P_{X,k,s}=p_{X,2q-1}+p_{X,2q}-p_{X,q-1}-p_{X,q},\\
 \widehat F_k^{(s)}=A_{K,k,s}+P_{K,k,s},\qquad
 \sum_N\sigma_N\log\det Q_{0,N}=A_{0,k,s}+P_{0,k,s},\\
 A_{X,k,s}\ge0,\qquad 0\le P_{X,k,s}\le E_{k,s}^{\vee}.
\end{gathered}\tag{PBR5}
\]
Increasing $N$ enlarges the original polynomial lift space without
changing its norm on previous polynomials or its relation map.
Minimum primal lift norms therefore decrease. Inversion in PBR2
increases $\mathcal D_N$, and restriction or minimization under
the same onto $B_K$ preserves that increase. The variational
principle then increases each relative eigenvalue in the fixed
reference metric. Thus $p$ increases and $n$ decreases, proving
all nonnegativities in PBR5. The two high positive sums are bounded
by PBR4. Finally $\log\mu=\log^+\mu-\log^+(1/\mu)$ and
$D_{K,N}=\overline{\widehat H_{K,N}^{-1}}$ give the first return
identity with the four stated signs; PBR3 gives the other.
The fixed reference determinants cancel only in these four
signed sums. In particular the small eigenvalue contributions
$n_{X,N}$ have not been bounded away or discarded.

\subsection{The original metric on the low image}

Let $G_N^{(s)}$ be the original primary quotient metric and let
$Q_{\Xi,N}^{\rm fix}$ be its induced metric on $B_\Xi=\Xi(V_A)$
in the original fixed image frame. Set $W_0=\Xi(L)$.
There is the exact isometric isomorphism
\[
\begin{gathered}
 \iota_N:(K+L)/K\longrightarrow W_0,\qquad
 [p]\longmapsto\Xi p,\\
 \|\iota_N[p]\|_{Q_{\Xi,N}^{\rm fix}}^2
 =\min_{a\in K}G_N^{(s)}(p+a,p+a)
 =\|[p]\|_{(K+L)/K,N}^{\,2},\\
 \dim W_0=v-t_0,\qquad
 F_{W_0}^{(s)}
 :=\sum_N\sigma_N\log\det
       (I_{W_0}^*Q_{\Xi,N}^{\rm fix}I_{W_0})
 =F_{K^+/K}^{(s)}.
\end{gathered}\tag{PBR6}
\]
Indeed $\Xi p\in\Xi(L)$ if and only if $p\in K+L$.
The kernel of the displayed map is exactly $K$ and it is onto.
The full fibre of $\Xi p$ in $V_A$ is $p+K$, already contained
in $K+L$. Hence minimization in the full image norm is exactly
the quotient minimization on $K+L$, proving the norm identity
at each cutoff. In transported frames the Grams agree exactly.
For any other fixed frames their congruence has the same fixed
Jacobian at all four endpoints, which cancels with $\sum\sigma_N=0$.
This proves the last identity without identifying $W_0$ with
the entire residual image.

In the FGC coordinates the onto map
$[p]\mapsto[\alpha_A\mathcal T_Ap/\chi']$ from $V_A$ to
$Y_N/\operatorname{im}Z_k$ has kernel $K+L$.
The corresponding quotient map sends this class to $[\Xi p]$ in
$B_\Xi/W_0$. The kernel assertion follows directly: subtract the
kernel polynomial giving the same FGC image, then the remaining
zero image under $\mathcal T_A$ is in $L$; the converse follows
from the two kernels. This proves the exact target isomorphism
in PEC27.

The source map comparing this quotient to the native Gamma one
is the original CEP map $S_N$, followed by
$D_{\alpha,N}(x,f)=(x,\alpha_A f)$ on the full GIP graph.
Retaining every relation column and every $Z_k$ column, minimization
first in the low coordinate and then in these polynomial fibres
gives precisely $Q_{0,N}$ in PBR3. Let $\bar Q_{\Xi,N}$ be the
induced original metric on $B_\Xi/W_0$ in the transported frame.
PEC28's full exterior comparison yields
\[
\begin{gathered}
 |\log\det\bar Q_{\Xi,N}-\log\det Q_{0,N}|
 \le2E_N^{\rm CEP}+2j_0|\log|\alpha_A||,\\
 E_{k,s}^{\rm tr}=2\sum_NE_N^{\rm CEP}
                     +8j_0|\log|\alpha_A||=o(kq),\\
 \Phi_{k,s}=F_{W_0}^{(s)}+A_{0,k,s}+P_{0,k,s}
                      +\delta_{0,k,s},\qquad
 |\delta_{0,k,s}|\le E_{k,s}^{\rm tr}.
\end{gathered}\tag{PBR7}
\]
To see the error signs explicitly, put
$d_N=\log\det\bar Q_{\Xi,N}-\log\det Q_{0,N}$.
In a fixed frame adapted to $W_0$, the full Gram determinant is
its restriction determinant times the Schur complement determinant,
including its fixed frame Jacobian. Therefore
$\delta_{0,k,s}=\sum_N\sigma_Nd_N$; the triangle inequality gives
the stated bound. The projection used to take each quotient is
an exterior contraction; the inverse source map gives the reverse
inequality. This proves that the factors $\alpha_A$, all four
original CEP errors, and the full graph denominator occur in
this very metric comparison.

\subsection{Both low terms and the full one-sided kernel errors}

Keep the original FEX endpoint constants explicitly. Write
$h_N^{\rm harm}=\sum_{a=1}^N1/a$ and $w_d=d!/(s/2)_d$.
The symbol $t$ here remains the kernel-column count from PBR1.
\[
\begin{gathered}
 n_N=N+1-v,\qquad
 m_N=\log\max\left\{1,\frac{v!}{|\mu_v|}
       \sum_\nu|a_\nu|e^{|b_\nu-4|(h_N^{\rm harm}+2)}\right\},\\
 \zeta_{N,s}=\log\prod_{d=v}^N\frac{w_d}{w_{d-v}},\qquad
 e_{N,s}^-=\zeta_{N,s}-2(n_N-t)m_N,\quad e_{N,s}^+=2tm_N,\\
 e_{N,s}^{\rm FEX}
   =\log\det\widehat H_{K,N}^{(s)}
          -\log\det H_{K^+/L,N}^{(s)},\qquad
 e_{N,s}^-\le e_{N,s}^{\rm FEX}\le e_{N,s}^+,\\
 D_s^-=e_{q-1,s}^-+e_{q,s}^--e_{2q-1,s}^+-e_{2q,s}^+,\quad
 D_s^+=e_{q-1,s}^++e_{q,s}^+-e_{2q-1,s}^--e_{2q,s}^-,\\
 \mathcal E_s^{\rm flag}=\sum_N\sigma_N e_{N,s}^{\rm FEX}
       \in[D_s^-,D_s^+].
\end{gathered}\tag{PBR8}
\]
These are the original FEX18--20/34 expressions, with both source
masses retained in the underlying norms. The endpoint inequalities
come from the full flag determinant bounds, with squared Jacobian
$e^{\zeta_{N,s}}$ and both operator bounds $e^{m_N}$.
The signs of the high endpoint bounds are reversed when forming
the four sum, giving exactly $D_s^\pm$ above.

Retain also the exact finite AKS/FEX allowance
\[
\begin{gathered}
 R_{\rm cen}=k\sqrt{\delta^2+\gamma^2},\qquad
 \tau_q=\pi/2-1/q,\quad a_q=2\log3/\tau_q,\quad
 s_q=(a_q^q-1)/(a_q-1),\\
 b_{s,q}=[(9/25)\sin(1/q)]^{-s/2}
                      e^{2R_{\rm cen}(\log3+\tau_q/2)},\\
 \mathcal L_{s,h}=\log\left(1+b_{s,q}s_q^2
          \sum_{d=q}^{q+h-1}\frac{d!}{(s/2)_d}(25/16)^d\right),
 \qquad \mathcal L_s=\mathcal L_{s,q}+\mathcal L_{s,q+1},\\
 0\le F_L^{(s)}\le v\mathcal L_s,\qquad
 0\le F_{W_0}^{(s)}\le(v-t_0)\mathcal L_s,\qquad
 \lambda_s^{\rm low}=F_L^{(s)}-F_{W_0}^{(s)},\\
 F_K^{(s)}
 =A_{K,k,s}+P_{K,k,s}-\mathcal E_s^{\rm flag}
                                      +\lambda_s^{\rm low}.
\end{gathered}\tag{PBR9}
\]
The low metric estimates are the original restriction and quotient
consequences of AKS, as proved in FEX32--33. Applying the two complete
Schur factorizations to $K^+$ gives OER11, and substituting PBR5
and the isometry PBR6 gives the last identity. The quantities
$F_L$ and $F_{W_0}$ remain their separate original volumes.
In particular they are not individually removed by using their
difference. Combining PBR7 and PBR9 gives the additional exact
identity
\[
 F_K^{(s)}+\Phi_{k,s}
 =A_{K,k,s}+A_{0,k,s}+P_{K,k,s}+P_{0,k,s}
       -\mathcal E_s^{\rm flag}+F_L^{(s)}+\delta_{0,k,s}.
 \tag{PBR10}
\]
The $W_0$ terms cancel in this sum because they are precisely the
same original metric by PBR6. They remain in each summand PBR7/9.

\subsection{The complete current finite interval}

First retain the exact low volumes in the new finite kernel bounds,
and separately display the previously proved coarse bounds:
\[
\begin{gathered}
 \kappa_s^-=A_{K,k,s}-D_s^++\lambda_s^{\rm low},\qquad
 \kappa_s^+=A_{K,k,s}+E_{k,s}^{\vee}-D_s^-+\lambda_s^{\rm low},
 \qquad \kappa_s^-\le F_K^{(s)}\le\kappa_s^+,\\
 A_{K,k,s}-D_s^+-(v-t_0)\mathcal L_s\le F_K^{(s)}
 \le A_{K,k,s}+E_{k,s}^{\vee}-D_s^-+v\mathcal L_s,\\
 \varphi_s^-=F_{W_0}^{(s)}+A_{0,k,s}-E_{k,s}^{\rm tr},\qquad
 \varphi_s^+=F_{W_0}^{(s)}+A_{0,k,s}+E_{k,s}^{\vee}+E_{k,s}^{\rm tr},
 \qquad \varphi_s^-\le\Phi_{k,s}\le\varphi_s^+,\\
 A_{0,k,s}-E_{k,s}^{\rm tr}\le\Phi_{k,s}
 \le A_{0,k,s}+E_{k,s}^{\vee}+E_{k,s}^{\rm tr}
                                      +(v-t_0)\mathcal L_s.
\end{gathered}\tag{PBR11}
\]
PBR5/8/9 proves the first line by taking the least and greatest
allowed signed flag errors; PBR7 proves the third line. The
original low caps then prove the second and fourth lines, which
are exactly PEC22 and PEC30. Thus using the exact low difference
has a specified morphism to each coarser bound.

Let $\Phi^-_{k,s},\Phi^+_{k,s}$ be the full-window RWC bounds of
OER17--18, and retain their entire conditional matrices
\[
\begin{gathered}
 U_N^+=[u_{N+1},\ldots,u_{N+q}],\quad
 B_N=I_q+(CU_N^+)^*(CR_NC^*)^{-1}(CU_N^+),\\
 Z_N=D[I_q-R_NC^*(CR_NC^*)^{-1}C]U_N^+,\quad
 \Omega_{N+q}=\Omega_N+Z_NB_N^{-1}Z_N^*,\\
 \Phi_{k,s}=\sum_{N=q-1,q}
 \log\left(\sum_{a=0}^j
 \left\|\bigwedge^a(\Omega_N^{-1/2}Z_NB_N^{-1/2})\right\|_{\rm HS}^2\right).
\end{gathered}\tag{PBR12}
\]
Here $D$ denotes the original residual rows, as in OER17.
All $q$ columns, the complete conditioning inverse, and every
exterior coefficient remain. The proof of the covariance identity
is the full Schur update RWC4--6; applying the determinant lemma
and expanding all singular-value products gives the last line.
Its finite bounds therefore apply to the same $\Phi$ as PBR7.
Define
$\Phi_{k,s}^{-,\vee}=\max\{\Phi^-_{k,s},\varphi_s^-\}$ and
$\Phi_{k,s}^{+,\vee}=\min\{\Phi^+_{k,s},\varphi_s^+\}$.

Retain the original OAR expression $\mathcal F_{\partial,k}^{(s)}$
and its full errors
$-e_{k,s}^-\le\mathfrak e_{k,s}\le e_{k,s}^+$ in
$F_K^{(s)}=\mathcal F_{\partial,k}^{(s)}-\Phi_{k,s}
+\mathfrak e_{k,s}$. These errors are distinct from PBR8.
With all earlier OER19 branches still included, the new receiver is
\[
\begin{gathered}
 L_s^\vee=\max\{L_s^{\rm win},\kappa_s^-,
  \mathcal F_{\partial,k}^{(s)}-\Phi_{k,s}^{+,\vee}-e_{k,s}^-\},\\
 U_s^\vee=\min\{U_s^{\rm win},\kappa_s^+,
  \mathcal F_{\partial,k}^{(s)}-\Phi_{k,s}^{-,\vee}+e_{k,s}^+\},
 \qquad L_s^\vee\le F_K^{(s)}\le U_s^\vee .
\end{gathered}\tag{PBR13}
\]
Each lower entry is at most the actual scalar and each upper entry
is at least that scalar, proving the intersection and its nonempty
order without assuming an unproved relation between the entries.
On PMR's exact domain
$k\ge29$ and $k\sqrt{\delta^2+\gamma^2}\le2^{-33}(k+1)^2$,
also include in this same maximum and minimum respectively
$\mathcal F_{0,k}^{(s)}-\Phi_{k,s}^{+,\vee}-e_{k,s}^{0,-}$
and $\mathcal F_{0,k}^{(s)}-\Phi_{k,s}^{-,\vee}+e_{k,s}^{0,+}$,
where $e_{k,s}^{0,\pm}=e_{k,s}^{\pm}+c_k^\circ$.
Outside that domain these two entries are omitted; the original
OAR expression and errors remain. The seven-moment
$\mathcal F_{0,k}^{(s)}$ is distinct from the quotient return
$\sum\sigma_N\log\det Q_{0,N}$ in PBR5.

For the arithmetic map retain $|u|\ge R_*$ and the same correlated
AMT allowances, with all their original finite constants:
\[
\begin{gathered}
 a_{K,k}^-=\max\{-C_K^-,\delta_k^\sigma-C_B^+\},\qquad
 a_{K,k}^+=\min\{C_K^+,\delta_k^\sigma+C_B^-\},\\
 L_{\rm ar}^{\rm win}
 =\max\{L_{\rm ar}^{\rm op},L_1^{\rm win}+a_{K,k}^-\},\qquad
 U_{\rm ar}^{\rm win}
 =\min\{U_{\rm ar}^{\rm op},U_1^{\rm win}+a_{K,k}^+\},\\
 L_{\rm ar}^{\vee}
 =\max\{L_{\rm ar}^{\rm win},L_1^\vee+a_{K,k}^-\},\qquad
 U_{\rm ar}^{\vee}
 =\min\{U_{\rm ar}^{\rm win},U_1^\vee+a_{K,k}^+\},\\
 F_K^{\rm ar}\in[L_{\rm ar}^{\vee},U_{\rm ar}^{\vee}],\qquad
 F_B^{\rm ar}\in[\mathcal B_k^{\rm ar}-U_{\rm ar}^{\vee},
                  \mathcal B_k^{\rm ar}-L_{\rm ar}^{\vee}].
\end{gathered}\tag{PBR14}
\]
AMT gives the identical scalar
$F_K^{\rm ar}-F_K^{(1)}$ in $[a_{K,k}^-,a_{K,k}^+]$.
Adding this interval to PBR13 proves the new entries. The earlier
entries remain; subtracting from the same total
$F_K^{\rm ar}+F_B^{\rm ar}=\mathcal B_k^{\rm ar}$ reverses the
two endpoints and proves the final line. No allowance has changed.
Previously available reverse source intersections remain valid.

\subsection{The signed return and the complete same-class action budget}

Keep $\mathcal H_k=-T_k$, the full arithmetic correction
$\delta_k^\sigma$, and exactly the original budget
\[
\begin{gathered}
 C_k^{\rm budget}
 =F_B^{(k)}+\mathcal W_k^{\rm ar}-\mathcal P_--\mathcal P_+,\qquad
 \Psi_k=F_K^{(k)}+\mathcal H_k+\delta_k^\sigma,\\
 L_k^\vee+\mathcal H_k+\delta_k^\sigma
 \le\Psi_k\le U_k^\vee+\mathcal H_k+\delta_k^\sigma,\\
 J_k^{\rm action}
 =2\sum_{N=q-1}^{2q-1}w_N\log\epsilon_N^{\rm ar}
 =\Psi_k+C_k^{\rm budget},\\
 L_J^\vee=L_k^\vee+\mathcal H_k+\delta_k^\sigma+C_k^{\rm budget},
 \quad
 U_J^\vee=U_k^\vee+\mathcal H_k+\delta_k^\sigma+C_k^{\rm budget},\\
 4q\log\mathcal L_{h,k}\le J_k^{\rm action}
       \in[L_J^\vee,U_J^\vee],\qquad
 \mathcal L_{h,k}\le\min_N\epsilon_N^{\rm ar}
 \le\prod_N(\epsilon_N^{\rm ar})^{w_N/(2q)}
 \le e^{U_J^\vee/(4q)} .
\end{gathered}\tag{PBR15}
\]
The weights are $w_{q-1}=w_{2q-1}=1$ and $w_N=2$ in between,
so $\sum w_N=2q$. The allowance $\epsilon_N^{\rm ar}$ is the
original rank-two action radius and
$\mathcal L_{h,k}=2\delta e(k+1)\lfloor(k+1)^2/4\rfloor$,
with $e=1$ in this section. MCE's action identities still hold
for every original multiplicity. RFX gives $\phi_N(0)=0$ in
the original untilted source, yielding the displayed exact
geometric mean. It retains $z_B=-z_K$, and hence the mixed
term $a_Kd_B+a_Bd_K+2|z_K|^2$. The full generic real-tilt MCE
identity with $\log(\epsilon_N^2+\phi_N^2)$ and the exact MCP
derivative remain unchanged. Addition to PBR13 proves the
signed bounds; addition of the complete budget proves the
action bounds. Applying the same-class lower radius at each
degree proves the geometric-mean inequalities. The whole
reference boundary, monic window, two contraction penalties
and action leakage are retained.

In particular the explicit coarse PEC32 entries become
\[
\begin{split}
 A_{K,k,k}-D_k^+-(v-t_0)\mathcal L_k
          +\mathcal H_k+\delta_k^\sigma
 &\le\Psi_k\\
 &\le A_{K,k,k}+E_{k,k}^{\vee}-D_k^-+v\mathcal L_k
          +\mathcal H_k+\delta_k^\sigma .
\end{split}\tag{PBR16}
\]
The exact-low entries of PBR13 can strengthen these entries;
every older bound is still in the intersection. This is an
actual control interval in the original signed direction. It
does not remove $C_k^{\rm budget}$ from the action inequality.

\subsection{Actual support, the common-factor map, and the retained residual}

Let $a_\pm,b_\pm$ be the actual extrema of the nonzero coefficients
of $A$, and let $d_a=a_+-a_-$, $d_b=b_+-b_-$.
The complete RCF/ACR maps give
\[
\begin{gathered}
 \mathcal I_G=
 ([-a_-,k-a_+]\times[-b_-,k-b_+])\setminus[0,k-8]^2,\quad
 g_A=(k+1-d_a)(k+1-d_b)-q',\\
 \rho_G:V_A\twoheadrightarrow\C[S']_{<g_A},\qquad
 V_0=\ker\rho_G,\quad W=\Xi(V_0),\\
 \rho_G I_K=W_G^{-1}D_GC_G^{\mathsf T}\pi^{\mathsf T},
 \qquad B_G=\pi C_G,\quad
 t_G=\dim(B_\Xi/W)=g_A-\operatorname{rank}B_G .
\end{gathered}\tag{PBR17}
\]
All columns of $C_G$ are the actual shifted coefficient vectors
$f_{ij}=\sum_{a,b}a_{ab}e_{i+a,j+b}$ for $(i,j)\in\mathcal I_G$.
The root Vandermonde $W_G$, evaluated-polynomial diagonal
$D_G=\operatorname{diag}(\chi'(\tau_\ell)^{-1})$, and original
exterior selector $\pi$ are the full ACR11--13 matrices.
Here $\chi'$ is the original lower-grid polynomial; the prime
does not denote differentiation.
The exact morphism from $V_A$ to the last quotient sends $p$ to
its class modulo $V_0+K$. The first isomorphism theorem identifies
it with $\C[S']_{<g_A}/\rho_G(K)$; composing with $\Xi$ identifies
it with $B_\Xi/W$. The displayed matrix for $\rho_G I_K$ has the
same rank as $B_G$, proving the last formula with all actual
coefficients. This retains the original linear relation quotient;
RCF's ideal hull does not replace its denominator. The full
common-factor metric comparison is RCM23--27 and its actual
coefficient-rank specialization is ACR18--23.

LBRE proves the exclusion of all four boundary rulings for the
actual limiting two-plane, with its full coefficient margin.
ASR uses the actual entire boundary coefficient functions to
produce the explicit radius
$R_{\rm supp}=\max\{R_*,1,2C_1/m_1,\ldots,2C_4/m_4\}$,
where $m_j>0$ is the larger of the two limiting coordinate
moduli and $C_j$ is the full convergent absolute tail sum.
For $|u|\ge R_{\rm supp}$ each selected pair has a nonzero
coordinate by the bound $C_j/|u|\le m_j/2$.
On the remaining compact reciprocal-period disk, each chosen
entire coordinate is nonzero at zero and consequently has
finitely many zeros. Taking the actual common zeros of each
three-coordinate boundary vector gives exactly ASR18's finite
exception set $\mathcal U_{\partial}^*$ in
$\{R_*\le|u|<R_{\rm supp}\}$. Thus, on the original arithmetic
period domain outside this finite set,
\[
\begin{gathered}
 (a_-,a_+,b_-,b_+)=(0,8,0,8),\quad
 d_a=d_b=8,\quad g_A=t_G=0,\quad V_0=V_A,\quad W=B_\Xi,\\
 B_\Xi/W_0\simeq Y_N/\operatorname{im}Z_k,\qquad
 \dim W_0=v-t_0,\quad\dim(B_\Xi/W_0)=j-v+t_0 .
\end{gathered}\tag{PBR18}
\]
The first line follows by substitution into PBR17: the target
of $\rho_G$ is zero, so its kernel is all of $V_A$.
The second line is the separate exact map PBR6--7. In particular
the full residual image $W=B_\Xi$ is still equipped with
$Q_{\Xi,N}^{\rm fix}$, its low subspace and its complementary
dual deficit. Neither $t_0$ nor either metric return is forced
to vanish by $g_A=t_G=0$. On the finite exceptional set the
literal ACR rank matrices and RCM metric comparison remain;
PBR2--16 did not assume full support and remain on their
original hypotheses.

Finally the complete OVR/ECL outer identity and the proved
finite errors imply
\[
\begin{gathered}
 A_{K,k,s}+A_{0,k,s}=\Delta q\mathfrak C_\partial+o(kq),\qquad
 \frac{A_{K,k,s}+A_{0,k,s}}{kq}\longrightarrow16\mathfrak C_\partial,\\
 21.66851180520<16\mathfrak C_\partial<21.66851180521,\qquad
 \frac{\Psi_k}{kq}
   =\frac{A_{K,k,1}}{kq}-\frac{\mathcal C_\Gamma}{4}+o(1).
\end{gathered}\tag{PBR19}
\]
Indeed PBR4, PBR7--9 and the complete FEX36 allowance give
$F_K^{(s)}=A_{K,k,s}+o(kq)$ and
$\Phi_{k,s}=A_{0,k,s}+o(kq)$.
Substitute these equalities in the original evaluated outer sum.
The second line uses the certified ECL coefficient in the same
original EIQ variables. For the last equality use the original
AMT source-order agreement and the proved signed Gamma return
$\mathcal H_k/(kq)\to-\mathcal C_\Gamma/4$, while the original
arithmetic correction is $o(kq)$. All these are the accepted
complete proof providers, preserved in the current full source.
This determines the sum of the two specified spectral deficits
and the exact signed receiver in terms of the kernel one; it
does not assign an individual leading value to either deficit.
OER15--16's actual $q$-scale Gamma interval, including its
strictly positive radius and both unequal finite endpoint
errors, is unchanged. The next analytic quantity in this
receiver is the actual deficit distribution in PBR5 and its
complete action comparison PBR15.
% END PBR COMPLETE RECEIVING PROOF
% END CURRENT PBR

\clearpage
% BEGIN CURRENT LJR
% Original low endpoint: every root, moment, Gamma mass and quotient retained.
\section{Exact low-jet reconstruction and its original quotient spectrum}
\label{sec:LJR}

Retain the simple-quartet fixed-period hypotheses of RDS, DNE and PEC.
The original conductor and all its coefficients remain fixed as
$k=4l+1\ge9$ varies. Set
\[
 \begin{gathered}
 q=(k+1)^2,\quad q'=(k-7)^2,\quad \Delta=q-q',\quad N=q-1,\\
 d=q-v-1,\quad n=d+1=q-v,\quad \mathfrak m=q-q'-v,
 \end{gathered}
 \tag{LJR1}
\]
with $0\le v\le80$, $\mu_v\ne0$, and the same centers $c=k/2$,
$c'=c-4$. The original source order is $s\in\{1,k\}$,
$b=s/2$, $M_s=(2\pi)^{s/2}$ and $h_a=M_sa!(b)_a$.
The roots are
$\beta_{ij}=c+(2i-k)\delta+i(2j-k)\gamma$ and the analogous
lower roots $\beta'_{ij}$, with $0<\delta<1/2$, $\gamma>2$.
Write $\sigma=(\delta^2+\gamma^2)^{1/2}$.

\subsection{The exact finite reconstruction and its inverse}

Let $\phi_a(S')=p_a((S'-c')/i)/\sqrt{h_a}$ be the complete
original orthonormal source frame through degree $d$. Define
$B_s$ and the original upper Vandermonde by
\[
 S'^a=\sum_{r=0}^d(B_s)_{ra}\phi_r(S'),\qquad
 \mathbb V_{j,(a,b)}=\beta_{ab}^{\,j}
       \quad(0\le j<q,\ 0\le a,b\le k).
 \tag{LJR2}
\]
These matrices are not assigned new source norms. The coefficient
matrix $B_s$ is upper triangular, with diagonal
$i^a\sqrt{h_a}$, since the leading coefficient of $\phi_a$
is $i^{-a}/\sqrt{h_a}$. In particular
$|\det B_s|^2=\prod_{a=0}^dh_a$.
For an explicit formula, let $\mathsf J_b$ be the real symmetric
$(d+1)$-by-$(d+1)$ Jacobi matrix with adjacent entries
$\sqrt{(r+1)(b+r)}$. Then
\[
 (B_s)_{ra}=\sqrt{M_s}\,
                 [(c'I+i\mathsf J_b)^a]_{r0}.
 \tag{LJR3}
\]
The original three-term recurrence proves this formula by repeated
multiplication by $S'=c'+iy$. A path of length at most $d$
starting at degree zero never needs a degree above $d$, so the
finite Jacobi truncation is exact in all the displayed columns.

For a vector $w=(w_0,\ldots,w_d)^{\mathsf T}$ define the
complex-linear functional $\Lambda_w$ by
$\Lambda_w(\phi_a)=w_a$. Its first $n$ ordinary moments are
$B_s^{\mathsf T}w$. Define the full moment-convolution matrix
\[
 (J_E)_{ja}=\binom ja\mu_{j-a}
       \quad(0\le j<q,\ 0\le a\le d),\qquad
 \mathsf R_s=\mathbb V^{-1}J_E B_s^{\mathsf T}.
 \tag{LJR4}
\]
Terms with $a>j$ are zero. The first $v$ rows of $J_E$ vanish,
and its remaining square block $\widetilde J_E$ is lower triangular
with diagonal $\binom{v+a}{a}\mu_v$, $0\le a\le d$.
Therefore $\mathsf R_s$ is an isomorphism from $\mathbb C^n$
onto the coefficient space of the original $\mathscr N_v$.

To verify the map exactly, let $F$ be any formal series with
$F^{(a)}(0)=(B_s^{\mathsf T}w)_a$ for $a\le d$.
Leibniz' rule gives
\[
 (E_AF)^{(j)}(0)
 =\sum_{a=0}^{d}\binom ja\mu_{j-a}F^{(a)}(0)
 \qquad(0\le j<q).
 \tag{LJR5}
\]
No unknown derivative of $F$ enters, because $\mu_{j-a}=0$
when $a>j-v$. The upper Vandermonde then reconstructs the
unique upper exponential polynomial $\mathcal N$ with precisely
these first $q$ jets. This proves LJR4 as the actual map of
original polynomial functionals to original exponential coefficients.

Let $\kappa_r$ denote the coefficients of the finite inverse
series of
$B_\infty(z)=v!E_A(z)/(\mu_vz^v)$, so
$B_\infty(z)^{-1}=\sum_{r\ge0}\kappa_rz^r$ near zero.
For $\nu\in\mathscr N_v$ write $u=\mathbb V\nu$. The full inverse is
\[
 \begin{gathered}
 f_r=\frac{v!}{\mu_v}r!
       \sum_{a=0}^{r}\kappa_{r-a}\frac{u_{v+a}}{(v+a)!}
          \quad(0\le r\le d),\\
 \mathsf R_s^{-1}\nu=B_s^{-\mathsf T}(f_0,\ldots,f_d)^{\mathsf T}.
 \end{gathered}\tag{LJR6}
\]
Indeed multiplication of exponential Taylor coefficients by
$E_A=(\mu_v/v!)z^vB_\infty$ gives precisely this triangular
inverse. The constants $\kappa_0,\ldots,\kappa_d$ are determined
by the finite recurrence
$\kappa_0=1$,
$\kappa_r=-\sum_{a=1}^r[v!\mu_{v+a}/((v+a)!\mu_v)]\kappa_{r-a}$.
Thus every moment, factorial and inverse coefficient is supplied.

\subsection{Both exact quotient kernels and their metrics}

Let $V_{\rm lo,s}$ be the $n$-by-$q'$ matrix
$(V_{\rm lo,s})_{r,(a,b)}=\phi_r(\beta'_{ab})$, and let
$C$ be the entire original primary conductor. Then
\[
 \mathsf R_s V_{\rm lo,s}=C^{\mathsf T}.
 \tag{LJR7}
\]
For each column, its functional is evaluation at the corresponding
lower root, so $F(z)=e^{\beta'_{ab}z}$. The original root-sum
identity gives
$E_A(z)e^{\beta'_{ab}z}=\sum_{r,t}a_{rt}
e^{\beta_{a+r,b+t}z}$. Its coefficient column is exactly
$C^{\mathsf T}$, and LJR5 reconstructs it, proving LJR7.
Both matrices have rank $q'$, by the lower Vandermonde and
the original injectivity of the conductor transpose.

Consequently LJR4 induces the exact isomorphism
\[
 \overline{\mathsf R}_s:
 \mathbb C^n/\operatorname{im}V_{\rm lo,s}
 \longrightarrow
 \mathscr N_v/\operatorname{im}C^{\mathsf T}.
 \tag{LJR8}
\]
Its quotient kernel is zero: if $\mathsf R_sw=C^{\mathsf T}\eta$,
then LJR7 and injectivity give $w=V_{\rm lo,s}\eta$.
Its inverse is induced by LJR6. The first space has the actual
native Gamma dual quotient norm, using the Euclidean norm on the
orthonormal functional-value coordinates $w$. The second has the
unchanged original exponential coefficient quotient norm.

Use precisely DNE's fixed $Z_{\exp}$ on the right and set
$W_s=\mathsf R_s^{-1}Z_{\exp}$ on the left. In these paired frames,
\[
 \begin{gathered}
 D_s=W_s^*[I-V_{\rm lo,s}(V_{\rm lo,s}^*V_{\rm lo,s})^{-1}
                          V_{\rm lo,s}^*]W_s
                =\mathcal D_{q-1}^{(s)},\\
 G=Z_{\exp}^*[I-C^{\mathsf T}(\bar C C^{\mathsf T})^{-1}
                          \bar C]Z_{\exp}=G_{\exp}.
 \end{gathered}\tag{LJR9}
\]
The first equality with DNE follows because LJR6 gives exactly
the first $n$ orthonormal values of
$\Lambda_{\mathcal N_a/E_A}$; the allowed extensions are all
lower evaluations by LJR7. Both projection formulas are attained
minimum norms. In particular the reconstruction is represented by
the identity from $(\mathbb C^{\mathfrak m},D_s)$ to
$(\mathbb C^{\mathfrak m},G)$, not by an asserted isometry of
these two different forms.

Let $\lambda_a$ be the relative eigenvalues of $(D_s,G)$ as
in PEC12. The squared singular values of LJR8, in decreasing
order, are $\lambda_a^{-1}$ with the order reversed. Thus
\[
 \log\det(\overline{\mathsf R}_s^*
                   \overline{\mathsf R}_s)
       =\log\frac{\det G}{\det D_s},\qquad
 \sum_a\log^+\sigma_a(\overline{\mathsf R}_s)^2
       =\sum_a\log^+(\lambda_a^{-1}).
 \tag{LJR10}
\]
The adjoint and singular values here use the two specified
quotient norms.

For an actual onto map $B_0$, or an actual injective map $J_0$,
the identity in these coordinates induces the paired quotients
or restrictions with complete forms
\[
 \begin{gathered}
 D_{B_0}=(B_0D_s^{-1}B_0^*)^{-1},\quad
 G_{B_0}=(B_0G^{-1}B_0^*)^{-1},\\
 D_{J_0}=J_0^*D_sJ_0,\quad G_{J_0}=J_0^*GJ_0.
 \end{gathered}\tag{LJR11}
\]
The exact common kernels and targets are transported by LJR8.
For $B_0=B_K=Z_k^{\mathsf T}A_\ell$, DNE9 gives
$D_{B_K}=\overline{(\widehat H_{K,q-1}^{(s)})^{-1}}$.
For the full annihilator frame $J_0=A_\ell^{-1}P_Y^{\mathsf T}$
from PEC24, $D_{J_0}=\overline{Q_{0,q-1}^{-1}}$.
These identities retain the original low-quotient corrections
of FGC when compared with the original kernel and residual.

\subsection{The complete full-space and quotient Jacobians}

Let $\mathbb V_{<v}$ be the first $v$ rows of $\mathbb V$ and put
$\mathcal J_v=\det(\mathbb V_{<v}\mathbb V_{<v}^*)$, with
$\mathcal J_0=1$. Then the squared full-space Jacobian is
\[
 \det(\mathsf R_s^*\mathsf R_s)
 =\frac{\mathcal J_v}{|\det\mathbb V|^2}
   \left|\frac{\mu_v}{v!}\right|^{2n}
    \prod_{a=0}^{d}\left(\frac{(v+a)!}{a!}\right)^2
    \prod_{a=0}^{d}h_a .
 \tag{LJR12}
\]
Indeed $\mathsf R_s=\mathbb V^{-1}P_v\widetilde J_E
B_s^{\mathsf T}$, where $P_v$ inserts the last $n$ coordinates.
The complementary principal-minor identity for the positive
matrix $\mathbb V\mathbb V^*$ gives
$\det(P_v^*(\mathbb V\mathbb V^*)^{-1}P_v)
=\mathcal J_v/|\det\mathbb V|^2$.
For completeness this identity follows by block Gaussian
elimination: the determinant of a block matrix equals the
determinant of its first principal block times the determinant
of its Schur complement, while the bottom block of its inverse
is the inverse Schur complement. All principal blocks are
positive. The determinants of $\widetilde J_E$ and $B_s$,
computed above, prove LJR12.

The exact restriction Jacobian on the subspace in LJR7 is
$\det(\bar C C^{\mathsf T})/
\det(V_{\rm lo,s}^*V_{\rm lo,s})$.
Completing this subspace to a full source frame, and its image
to the transported target frame, factors the full Gram into
restriction and quotient Schur determinants. It follows that
\[
 \mathfrak J_{k,s}
 :=\frac{\det G}{\det D_s}
 =\det(\mathsf R_s^*\mathsf R_s)
       \frac{\det(V_{\rm lo,s}^*V_{\rm lo,s})}
            {\det(\bar C C^{\mathsf T})}.
 \tag{LJR13}
\]
No orthogonality of a chosen coefficient complement is assumed.

Let $\mathbb V'$ be the lower-root evaluation Vandermonde
of degrees $0,\ldots,q'-1$, and let $H_{\psi,s}$ be the full
$\mathfrak m$-by-$\mathfrak m$ original Gram of
$\chi'S'^a$, $0\le a<\mathfrak m$. Then
\[
 \det(V_{\rm lo,s}^*V_{\rm lo,s})
 =\frac{|\det\mathbb V'|^2\det H_{\psi,s}}
              {\prod_{a=0}^{d}h_a},
 \tag{LJR14}
\]
and therefore
\[
 \boxed{\quad
 \mathfrak J_{k,s}
 =\frac{\mathcal J_v|\det\mathbb V'|^2}
        {|\det\mathbb V|^2\det(\bar C C^{\mathsf T})}
   \left|\frac{\mu_v}{v!}\right|^{2n}
    \prod_{a=0}^{d}\left(\frac{(v+a)!}{a!}\right)^2
        \det H_{\psi,s}.
 \quad}\tag{LJR15}
\]
To prove LJR14, split the native polynomial space into the
low monomials through degree $q'-1$ and the ideal columns
$\chi'S'^a$. Their complete monomial coefficient matrix has
determinant one because $\chi'$ is monic. The full source
Gram determinant is $\prod h_a$, since the original monic
$S'$-polynomials have squared norms $h_a$ and their phases
have modulus one. Schur elimination therefore gives the
low-remainder quotient determinant $\prod h_a/\det H_{\psi,s}$.
Evaluation at the lower roots is precisely $\mathbb V'$
on that remainder. The inverse quotient determinant is thus
$|\det\mathbb V'|^2\det H_{\psi,s}/\prod h_a$.
In an orthonormal frame the inverse evaluation quotient Gram is the conjugate of
$V_{\rm lo,s}^*V_{\rm lo,s}$; their determinants agree because
they are positive real. This proves the identity with its
full original source mass before cancellation in LJR15.

\subsection{Finite reconstruction bounds on the actual grid}

Put $A_{\rm abs}=\sum_{r,t=0}^8|a_{rt}|$ and
$H_D=\sum_{a=1}^D1/a$, with $H_0=0$. For the original roots define
\[
 \begin{gathered}
 \mathfrak d_{ab}=\prod_{(i,j)\ne(a,b)}(\beta_{ab}-\beta_{ij}),\qquad
 L_{ab}(S)=\mathfrak d_{ab}^{-1}
                         \prod_{(i,j)\ne(a,b)}(S-\beta_{ij}),\\
 \Pi_{k,s}=\prod_{a=1}^{q-1}[2\sqrt{a(b+a)}+k\sigma],\\
 M_{k,s}^{\rm rec}=\max\left\{1,\,
 A_{\rm abs}e^{8\sigma(H_{q-1}+2)}\sqrt{M_s}\,
            \Pi_{k,s}\left(\sum_{a,b=0}^k|\mathfrak d_{ab}|^{-2}
                     \right)^{1/2}\right\}.
 \end{gathered}\tag{LJR16}
\]
The derivative product $\mathfrak d_{ab}$ is distinguished from
the complete lower polynomial $\chi'$. It retains its actual
complex sign in $L_{ab}$. Then
\[
 \|\mathsf R_s\|\le M_{k,s}^{\rm rec},\qquad
 \|\overline{\mathsf R}_s\|\le M_{k,s}^{\rm rec}.
 \tag{LJR17}
\]
Here all norms have precisely the source and target metrics of
LJR4 and LJR8.

To prove the first bound, evaluation of the reconstructed
functional on $L_{ab}$ gives
$\nu_{ab}=\Lambda_w(T_A L_{ab})$. Its coefficient row in the
$w$ coordinates therefore has Euclidean norm
$\|T_A L_{ab}\|_{s,c'}$. This is an exact row identity, not
a comparison with an unweighted polynomial coefficient norm.
In the original coordinates,
$T_Ap(c'+iy)=\sum a_{rt}p(c+i y+(b_{rt}-4))$.
Since $b_{rt}-4=(2r-8)\delta+i(2t-8)\gamma$ has modulus
at most $8\sigma$, the exact finite Taylor translation and
the original derivative estimate
$\|\partial_S\|_{\le q-1,s,c}\le H_{q-1}+2$
give the factor $A_{\rm abs}e^{8\sigma(H_{q-1}+2)}$.
The factor $-i$ in $\partial_S=-i\partial_y$ has modulus one
and is retained in the Taylor map.

Multiplication by $y$ from degree at most $a-1$ to degree at
most $a$ is the sum of the two original Jacobi shifts, each
bounded by $\sqrt{a(b+a)}$. Therefore each factor
$S-\beta_{ij}$ has multiplication norm at most
$2\sqrt{a(b+a)}+|c-\beta_{ij}|$ at this step.
Starting from the original constant norm $\sqrt{M_s}$ proves
$\|L_{ab}\|_{s,c}\le\sqrt{M_s}\Pi_{k,s}/|\mathfrak d_{ab}|$.
The sum of the squared row norms is the Hilbert--Schmidt norm
squared of $\mathsf R_s$, proving the first inequality.
For the quotient, choose the attained minimum representative
on the left and project its image orthogonally off
$\operatorname{im}C^{\mathsf T}$. This proves the second
inequality, with exactly the denominator in LJR8.

Every spacing in the denominator admits the exact factorization
\[
 \begin{split}
 |\mathfrak d_{ab}|={}&
 (2\delta)^k a!(k-a)!\,
 (2\gamma)^{k(k+1)}[b!(k-b)!]^{k+1}\\
 &\quad\cdot
 \prod_{i=0}^k\prod_{\substack{0\le j\le k\\j\ne b}}
 \left(1+\frac{\delta^2(a-i)^2}{\gamma^2(b-j)^2}\right)^{1/2}.
 \end{split}\tag{LJR18}
\]
In this formula only, $b$ is the second root index; the source
parameter inside $\Pi_{k,s}$ remains $s/2$.
Separate the $k$ factors in the row $j=b$, and factor
$2\gamma|b-j|$ from every other row to prove LJR18.
The elementary bounds
$a!(k-a)!=k!/\binom ka\ge(k/(2e))^k$ and the same bound
for $b$ give
\[
 |\mathfrak d_{ab}|\ge(k/e)^{q-1}
                   \delta^k\gamma^{k(k+1)},\qquad
 \Pi_{k,s}\le[k^2(3+\sigma)]^{q-1}.
 \tag{LJR19}
\]
The factorial estimate follows from
$\sum_{a=1}^k\log a\ge\int_1^k\log t\,dt$ and
$\binom ka\le2^k$.
For the second inequality use
$2\sqrt{a(a+b)}\le2a+b$,
$b\le k/2$, and
$2(1+1/k)^2+(1/2+\sigma)/k\le3+\sigma$ for $k\ge9$.
Consequently
\[
 \begin{split}
 \log M_{k,s}^{\rm rec}\le
 \max\{0,\ &(q-1)\log k+\log A_{\rm abs}
 +8\sigma(H_{q-1}+2)+\tfrac12\log q+\tfrac12\log M_s\\
 &+(q-1)\log[e(3+\sigma)]-k\log\delta
                              -k(k+1)\log\gamma\}\\
 &=q\log k+O_{\rm actual}(q).
 \end{split}\tag{LJR20}
\]
The last equality denotes an upper bound with that order of
error. It is uniform for both original source orders, since
$\log M_s=(s/2)\log(2\pi)=O(k)$.
The finite definition LJR16, with every actual spacing and
all conductor coefficients, remains the bound used below.

\subsection{A direct estimate of the nested determinant difference}

Let $\mathcal I=\{4,\ldots,k-4\}^2$ and let $\mathcal O$ be
its complement in the upper index grid. Then
$|\mathcal I|=q'$, $|\mathcal O|=\Delta$, and
$\beta_{i+4,j+4}=\beta'_{ij}+4$.
The coefficient map
$(S'+4)^a=\sum_{r\le a}\binom ar4^{a-r}S'^r$
is upper triangular with diagonal one. Hence the Vandermonde
on the central upper roots has exactly the lower Vandermonde
determinant. Let $\mathcal P_{\mathcal O}$ denote the unordered
distinct pairs with at least one index in $\mathcal O$. Thus
\[
 \begin{gathered}
 2M_{\mathcal O}=2|\mathcal P_{\mathcal O}|
       =q(q-1)-q'(q'-1)=\Delta(q+q'-1),\\
 \mathscr D_k:=\log\frac{|\det\mathbb V|^2}{|\det\mathbb V'|^2}
       =2\sum_{\{x,y\}\in\mathcal P_{\mathcal O}}
                                      \log|\beta_x-\beta_y|.
 \end{gathered}\tag{LJR21}
\]
This is cancellation of the identical central-central factors
in the exact Vandermonde products.

For $d(x,y)=\|x-y\|_\infty$ one has
$2\delta d(x,y)\le|\beta_x-\beta_y|\le2\sigma d(x,y)$.
For each fixed $x$, the number of other grid indices within
distance $t$ is at most $4t^2+4t$, for $1\le t\le k$.
The exact identity
$\log(k/d)=\int_d^kdt/t$ consequently gives
$\sum_{y\ne x}\log(k/d(x,y))\le2q-8$.
Summing this nonnegative estimate over the $\Delta$ outer
indices counts every pair in LJR21 at least once. It proves
\[
 \begin{split}
 2M_{\mathcal O}\log k+2M_{\mathcal O}\log(2\delta)
                        -2\Delta(2q-8)
 &\le\mathscr D_k\\
 &\le2M_{\mathcal O}\log k+2M_{\mathcal O}\log(2\sigma).
 \end{split}\tag{LJR22}
\]
In particular
$\mathscr D_k=\Delta(q+q'-1)\log k+O_{\rm actual}(kq)$.
This error is obtained from the affected pairs themselves,
not by subtracting two estimates with errors of order $q^2$.

\subsection{The complete remaining native Gram}

The original Gram $H_{\psi,s}$ in LJR15 is positive definite:
a nonzero polynomial has nonzero modulus off a finite set,
and the original Gamma density is positive on the real line.
The following bounds include its full mass and apply at every
$k\ge9$ under LJR1. Define
\[
 \begin{gathered}
 R_{\psi,k}=\max\{(k-8)\sigma,|c'|\},\quad
 T_*=\max\{1,2d+s/2+R_{\psi,k}\},\\
 Y=2(q+k\sigma+1),\quad h=1/q,\quad
 L_h=\frac{\pi}{\Gamma(3/4)^2}(\sin h)^{3/2},\\
 \beta_s=\frac{M_s}{\sqrt2\,\Gamma(s/2)},\qquad
 A_{k,s}=\beta_s L_h e^{-(\pi/2+h)(Y+1)},\\
 \mathcal H_{\mathfrak m}
 =\prod_{a=0}^{\mathfrak m-1}
             [(2a+1)\binom{2a}{a}^{\,2}]^{-1}.
 \end{gathered}\tag{LJR23}
\]
Then the finite two-sided estimate is
\[
 (A_{k,s}q^{2q'})^{\mathfrak m}\mathcal H_{\mathfrak m}
 \le\det H_{\psi,s}
 \le M_s^{\mathfrak m}
                T_*^{\,2\mathfrak m q'+\mathfrak m(\mathfrak m-1)}.
 \tag{LJR24}
\]

For the upper bound, $\chi'S'^a$ has the $q'$ actual lower
roots and $a$ additional zero roots. All are at distance at
most $R_{\psi,k}$ from $c'$. Applying the same original Jacobi
multiplication bound as in LJR17 through degree $q'+a\le d$
gives its norm at most $\sqrt{M_s}T_*^{q'+a}$.
Hadamard's Gram determinant inequality proves the upper side.

For the lower bound, the exact original Gamma recurrence and
the global density estimate give
\[
 \begin{gathered}
 dm_s(y)=\beta_s
       \prod_{j=0}^{(s-1)/4-1}(y^2+(2j+\tfrac12)^2)\,d\sigma(y),\\
 \frac{d\sigma}{dy}(y)
 \ge L_h e^{-(\pi/2+h)|y|}.
 \end{gathered}\tag{LJR25}
\]
These are the complete identities OSP12 and OFV9. To make
the lower-density constant explicit, Euler's integral rotated
through $\pi/2-h$ gives
\[
 |\Gamma(3/4-iy/2)|^2\le
 \Gamma(3/4)^2(\sin h)^{-3/2}e^{-(\pi/2-h)|y|}.
\]
The exact reflection product with
$|\Gamma(1/4+iy/2)|^2$ is $2\pi^2/\cosh(\pi y)$.
Dividing and using $\cosh(\pi y)\le e^{\pi|y|}$
proves the displayed lower density, including the factor
$1/(2\pi)$ in $d\sigma$.
The recurrence prefactor $\beta_s$ follows by applying
$\Gamma(z+1)=z\Gamma(z)$ exactly $(s-1)/4$ times and
clearing the factors $4^{-1}$ from the squared moduli.
Thus no original source mass has been discarded.

On the full interval $[Y,Y+1]$, every factor of
$\chi'(c'+iy)$ has modulus at least
$y-(k-8)\sigma\ge q$, and the polynomial density
multiplier in LJR25 is at least one because $y\ge1$.
Therefore the weighted density is at least
$A_{k,s}q^{2q'}$. The Gram of $(c'+iy)^a$ on this interval
is obtained from the Gram of $t^a$ on $[0,1]$ by the
translation $y=Y+t$ and coefficients with diagonal $i^a$.
This map has determinant of modulus one.

For completeness, Rodrigues' polynomial
$L_a(t)=(a!)^{-1}(d/dt)^a[t^a(t-1)^a]$
has leading coefficient $\binom{2a}{a}$.
Integration by parts $a$ times proves orthogonality to
every lower power; all boundary terms vanish.
It also gives
$\int_0^1L_a(t)^2dt=1/(2a+1)$, since
$\int_0^1t^a(1-t)^adt=(a!)^2/(2a+1)!$.
Dividing by the square of the leading coefficient gives
the monic norm in LJR23. Triangular orthogonal elimination
therefore proves that the complete interval Gram determinant
is exactly $\mathcal H_{\mathfrak m}$.
The positive-form inequality on that interval proves the
lower side of LJR24.

Both sides imply the useful direct scale
\[
 \log\det H_{\psi,s}
      =2\mathfrak m q'\log q+O_{\rm actual}(kq).
 \tag{LJR26}
\]
Here $\mathfrak m=O(k)$, $d=O(q)$ and
$\log T_*=\log q+O_{\rm actual}(1)$.
Moreover
$-\log\mathcal H_{\mathfrak m}
 \le\mathfrak m\log(2\mathfrak m+1)
                 +2\mathfrak m(\mathfrak m-1)\log2$,
using $\binom{2a}{a}\le4^a$.
Since $2h/\pi\le\sin h\le h$, $\log L_h=O(\log q)$.
Writing $s/2=2\ell_s+1/2$,
$\Gamma(s/2)=\sqrt\pi\prod_{j=0}^{2\ell_s-1}(j+1/2)$
shows $|\log\beta_s|=O(k\log(k+1))$.
Finally $Y=O_{\rm actual}(q)$.
All errors in the logarithms of LJR24 are consequently
$O_{\rm actual}(kq)$; in particular the terms
$k^2\log q$ are smaller than $kq$.

\subsection{The exact remainder and every paired exterior volume}

Let $a_*$ be the actual nonzero pivot coefficient chosen by
the original lexicographic conductor proof. Its triangular
$q'$-by-$q'$ minor has determinant $a_*^{q'}$.
Cauchy--Binet and the row and column sum bounds on the full
convolution matrix give
\[
 |a_*|^{2q'}\le\det(\bar C C^{\mathsf T})
                              \le A_{\rm abs}^{2q'}.
 \tag{LJR27}
\]
The matrix in this determinant is the conjugate of $CC^*$;
their positive real determinants agree.
For the low-jet minor, Hadamard gives
$\mathcal J_v\le q^v[k(1/2+\sigma)]^{v(v-1)}$.
For $k\ge v-1$, choose the $v$ consecutive roots with first
index zero in its Cauchy--Binet sum. Their Vandermonde gives
$\mathcal J_v\ge(2\gamma)^{v(v-1)}
(\prod_{a=1}^{v-1}a!)^2$.
At $v=0$ both bounds mean $\mathcal J_0=1$.
Thus $\log\mathcal J_v=O_{\rm actual}(\log(k+1))$.
Also
$0\le\sum_{a=0}^d\log((v+a)!/a!)
\le vn\log(q+v)$.
The fixed period supplies the actual nonzero $\mu_v,a_*$
and $A_{\rm abs}$ in these statements; none is selected afresh.

Define the following exact, still unevaluated remainder:
\[
 \begin{split}
 \mathscr R_{k,s}:={}&
 \log\mathcal J_v-\mathscr D_k
 -\log\det(\bar C C^{\mathsf T})
 +2n\log|\mu_v/v!|\\
 &+2\sum_{a=0}^d\log\frac{(v+a)!}{a!}
 +\log\det H_{\psi,s}-2\mathfrak m q\log k .
 \end{split}\tag{LJR28}
\]
Every term is an original finite determinant, moment or root
product already specified above. Then
\[
 \log\mathfrak J_{k,s}
       =2\mathfrak m q\log k+\mathscr R_{k,s},\qquad
 |\mathscr R_{k,s}|=O_{\rm actual}(kq).
 \tag{LJR29}
\]
For the bound substitute LJR22, LJR26 and LJR27 into LJR28.
The only large logarithmic expression that remains is
$2\mathfrak m q'\log q-\Delta(q+q'-1)\log k$.
Since $\log q=2\log k+2\log(1+1/k)$ and
$\Delta=\mathfrak m+v$, subtracting $2\mathfrak m q\log k$
from this expression gives exactly
\[
 [(-3\mathfrak m+v)\Delta+\Delta-2vq]\log k
          +4\mathfrak m q'\log(1+1/k).
 \tag{LJR30}
\]
The coefficient of $\log k$ is $O_{\rm actual}(k^2)$;
the last term is $O_{\rm actual}(k^2)$.
Both are $o(kq)$. This proves LJR29 by an explicit
difference calculation.

Consider any fixed paired restriction, quotient, or
restriction followed by quotient of LJR8, with rank $t$.
Its two forms are the restrictions and attained quotients
of $D_s,G$ along the same specified algebraic flag.
Write $\mathfrak J_{t,k,s}$ for its squared Jacobian.
Then the exact finite bounds are
\[
 \log\mathfrak J_{k,s}
       -2(\mathfrak m-t)\log M_{k,s}^{\rm rec}
 \le\log\mathfrak J_{t,k,s}
 \le2t\log M_{k,s}^{\rm rec}.
 \tag{LJR31}
\]
To prove this for a flag $A\subset B$ in the common
$\mathfrak m$-dimensional coordinate space, factor each
full Gram determinant successively into the restriction
on $A$, the quotient $B/A$, and the quotient of the full
space by $B$. The factors are their actual Schur
determinants. Their Jacobians multiply to
$\mathfrak J_{k,s}$.
Every restriction of the map has norm at most $M_{k,s}^{\rm rec}$;
every induced quotient does too, by the minimum-lift and
projection proof of LJR17. The two complementary factors
have total rank $\mathfrak m-t$. Bounding them above and
dividing proves the lower inequality. Bounding the rank-$t$
factor directly proves the upper inequality.

Uniformly over all these actual flags and $0\le t\le\mathfrak m$,
\[
 \log\mathfrak J_{t,k,s}
       =2tq\log k+O_{\rm actual}(kq).
 \tag{LJR32}
\]
For a fully retained finite version put
$u_{k,s}=\log M_{k,s}^{\rm rec}-q\log k$ and
$\mathscr R_{t,k,s}=\log\mathfrak J_{t,k,s}-2tq\log k$.
LJR31 gives
$\mathscr R_{k,s}-2(\mathfrak m-t)u_{k,s}
 \le\mathscr R_{t,k,s}\le2t u_{k,s}$.
LJR20 bounds $u_{k,s}$ above by $O_{\rm actual}(q)$,
which proves both sides of LJR32 using LJR29.
No inverse coefficient norm or conditioning premise for
the flag occurs.

In particular apply LJR31 to the actual $B_K$ quotient and
the actual $J_0$ annihilator restriction in LJR11. If $p_t,n_t$
are their two positive logarithmic spectral parts in the
orientation $(D,G)$, then exactly
\[
 n_{t,q-1}-p_{t,q-1}
 =\log\mathfrak J_{t,k,s}
 =2tq\log k+\mathscr R_{t,k,s},\qquad
 0\le p_{t,q-1}\le\varepsilon^\vee_{q-1,s}=o(kq).
 \tag{LJR33}
\]
The last bound is PEC18 on precisely these forms. Thus
$n_{t,q-1}=2tq\log k+O_{\rm actual}(kq)$.
The exact remaining terms are $\mathscr R_{t,k,s}$ and
$p_{t,q-1}$, with LJR28 and LJR31 giving finite controls.
Their four-endpoint differences have not been assigned
a coefficient: LJR33 proves the low-endpoint leading scale
and supplies the literal maps and remainders to be carried
into that next calculation.

% END CURRENT LJR

\clearpage
% BEGIN CURRENT LRR
% BEGIN LRR COMPLETE RECEIVING PROOF
\section{First-endpoint reconstruction in the actual kernel and annihilator}
\label{sec:LRR}

Retain every original hypothesis, source order, metric, coefficient,
root, mass and sign of PBR1--19. The following applies the complete
LJR1--33 reconstruction at the actual first endpoint $N=q-1$.
The symbol $t_0$ continues to mean $\dim(K\cap L)$.
Define the two ranks without changing that symbol:
\[
 t_K=r_K-t_0,\qquad
 t_R=j-v+t_0,\qquad
 t_K+t_R=\mathfrak m.
 \tag{LRR1}
\]
PBR2 proves that $B_K=Z_k^{\mathsf T}A_\ell$ is onto rank
$t_K$, and PBR3 proves that
$J_0=A_\ell^{-1}P_Y^{\mathsf T}$ is injective rank $t_R$.
Indeed $Z_k$ has the actual independent kernel columns,
$A_\ell$ is invertible, $P_YZ_k=0$, and $P_Y$ is onto with
kernel $\operatorname{im}Z_k$. Thus
$\operatorname{im}J_0=\ker B_K$.

Let $\mathsf R_s=\mathbb V^{-1}J_EB_s^{\mathsf T}$ be the full
LJR4 reconstruction, including the original upper Vandermonde,
moment convolution and original Gamma orthonormal-value matrix.
LJR5 proves the exact first-$q$-jet identity by Leibniz' rule,
and LJR7 gives $\mathsf R_sV_{\rm lo,s}=C^{\mathsf T}$.
It therefore induces the isomorphism
$\overline{\mathsf R}_s:
\mathbb C^{q-v}/\operatorname{im}V_{\rm lo,s}
\to\mathscr N_v/\operatorname{im}C^{\mathsf T}$.
Its inverse is the full finite moment recurrence LJR6.
In the paired frames $\mathsf R_s^{-1}Z_{\exp}$ and $Z_{\exp}$,
this map is the identity from
$(\mathbb C^{\mathfrak m},\mathcal D_{q-1}^{(s)})$ to
$(\mathbb C^{\mathfrak m},G_{\exp})$, because the complete
minimum-extension norms are exactly LJR9/PBR2.

Taking the quotient under the same onto $B_K$ gives precisely
the two forms $(D_{K,q-1},G_K)$ of PBR2; restricting along the
same $J_0$ gives precisely $(D_{0,q-1},G_0)$ of PBR3.
The quotient metric is obtained by minimizing over the full
kernel of $B_K$ in each metric, giving
$(B_KD^{-1}B_K^*)^{-1}$ and $(B_KG^{-1}B_K^*)^{-1}$.
Restriction gives $J_0^*DJ_0$ and $J_0^*GJ_0$ directly.
Use $X=R$ for this residual restriction, with the exact dictionary
$(D_R,G_R,p_R,n_R,A_R,P_R)=(D_0,G_0,p_0,n_0,A_0,P_0)$
from PBR3/5; $t_0$ still denotes only $\dim(K\cap L)$.
Consequently their squared reconstruction Jacobians satisfy
\[
\begin{gathered}
 \mathfrak J_{X,k,s}
   =\frac{\det G_X}{\det D_{X,q-1}},\qquad X=K,R,\\
 \log\mathfrak J_{X,k,s}
   =n_{X,q-1}-p_{X,q-1}
   =2t_Xq\log k+\mathscr R_{X,k,s}.
\end{gathered}\tag{LRR2}
\]
The first equality uses the actual source and target norms of
the reconstruction, in that direction. For the second, a relative
eigenvalue $\mu$ of $(D,G)$ contributes $-\log\mu
=\log^+(1/\mu)-\log^+\mu$ to the reconstruction determinant.
Thus these are exactly the $n,p$ of PBR5, with no inversion
of their orientation.

Here every finite remainder is retained. Put
$u_{k,s}=\log M_{k,s}^{\rm rec}-q\log k$, where
$M_{k,s}^{\rm rec}$ is the complete finite LJR16 constant,
including all actual root spacings and all conductor coefficients.
Keep the exact full-space remainder LJR28:
\[
\begin{gathered}
 \mathscr R_{k,s}
 =\log\mathcal J_v-\mathscr D_k
 -\log\det(\bar C C^{\mathsf T})+2(q-v)\log|\mu_v/v!|\\
 \hspace{8mm}
 +2\sum_{a=0}^{q-v-1}\log\frac{(v+a)!}{a!}
 +\log\det H_{\psi,s}-2\mathfrak m q\log k,\\
 \mathscr R_{k,s}-2(\mathfrak m-t_X)u_{k,s}
 \le\mathscr R_{X,k,s}\le2t_Xu_{k,s},\qquad
 0\le p_{X,q-1}\le\varepsilon_{q-1,s}^{\vee},\\
 n_{X,q-1}
   =2t_Xq\log k+\mathscr R_{X,k,s}+p_{X,q-1}
   =2t_Xq\log k+O_{\rm actual}(kq).
\end{gathered}\tag{LRR3}
\]
The meanings of the finite terms are literal:
$\mathcal J_v=\det(\mathbb V_{<v}\mathbb V_{<v}^*)$ with
$\mathcal J_0=1$,
$\mathscr D_k=\log(|\det\mathbb V|^2/|\det\mathbb V'|^2)$,
and $H_{\psi,s}$ is the full original Gram of
$\chi'S'^a$, $0\le a<\mathfrak m$, retaining its Gamma mass.
The prime on $\chi'$ means the lower-grid polynomial.
LJR12--15 proves the full Jacobian by exact triangular and
Schur factorizations. LJR31 factors its determinant along the
same algebraic flag into the rank-$t_X$ part and its
complements. Every part has forward norm at most
$M_{k,s}^{\rm rec}$ by its actual minimum lift. Bounding the
complementary product above gives the lower inequality for
$\mathscr R_{X,k,s}$; bounding the selected product gives
the upper inequality. PEC18 bounds the positive part on
these very quotient and restriction forms. This proves
every finite assertion in LRR3.

For the final order statement, LJR20 gives
$u_{k,s}\le O_{\rm actual}(q)$ and LJR22--30 gives
$|\mathscr R_{k,s}|=O_{\rm actual}(kq)$ by the affected-root-pair
calculation and the complete Gram bounds. Substitution into
the two finite inequalities gives
$|\mathscr R_{X,k,s}|=O_{\rm actual}(kq)$ uniformly for
$0\le t_X\le\mathfrak m$; no lower bound for an
unretained inverse norm is needed. PEC15 gives
$\varepsilon_{q-1,s}^{\vee}=o(kq)$.
Thus the last equality in LRR3 follows, while its full finite
remainder and positive part remain in the preceding equality.

In the actual four-endpoint receiver this yields exactly
\[
\begin{gathered}
 A_{X,k,s}
 =2t_Xq\log k+\mathscr R_{X,k,s}+p_{X,q-1}
       +n_{X,q}-n_{X,2q-1}-n_{X,2q},\qquad X=K,R,\\
 F_K^{(s)}
 =2t_Kq\log k+\mathscr R_{K,k,s}+p_{K,q-1}
       +n_{K,q}-n_{K,2q-1}-n_{K,2q}
       +P_{K,k,s}-\mathcal E_s^{\rm flag}
       +F_L^{(s)}-F_{W_0}^{(s)},\\
 \Phi_{k,s}
 =F_{W_0}^{(s)}+2t_Rq\log k
       +\mathscr R_{R,k,s}+p_{R,q-1}
       +n_{R,q}-n_{R,2q-1}-n_{R,2q}
       +P_{R,k,s}+\delta_{0,k,s}.
\end{gathered}\tag{LRR4}
\]
The first equality substitutes LRR3 into the definition of
$A_X$ in PBR5. The next two substitute that same identity
into PBR9 and PBR7, respectively. Hence both original low
volumes, all three other endpoint deficits, every positive
spectral part, and all unequal finite flag and transfer errors
remain. PBR13--16 continues to use these exact $A_X$ in its
original source, arithmetic, signed and action intervals.
The first endpoint has been evaluated at its leading scale;
the remaining endpoint terms and the full $kq$-scale remainder
in LRR3 have not been assigned an individual coefficient.
% END LRR COMPLETE RECEIVING PROOF
% END CURRENT LRR

\clearpage
% BEGIN CURRENT HJR
% Literal compatible high jets, original Gamma extension and complete flags.
\section{Compatible high jets and the original positive return}
\label{sec:HJR}

Retain all fixed original quartet, amplitude, period and branch data
of LJR, DNE, PEC and OVR. In particular the arithmetic domain remains
$|u|\ge R_*$, the actual period is fixed while $k=4l+1$ grows, and
the source order is exactly $s=1$ or $s=k$. All asymptotic uses of
OFV retain its proved eventual threshold. Write
\[
 \begin{gathered}
 q=(k+1)^2,\quad q'=(k-7)^2,\quad
 n=q-v,\quad \mathfrak m=q-q'-v,\quad d_0=q-v-1,\\
 q-1\le N\le2q,\qquad D=N-v,\qquad
 e_N=N-q+1,\qquad D+1=n+e_N.
 \end{gathered}\tag{HJR1}
\]
The polynomial degree $D$ is the actual native source degree after
the conductor. The original centers remain $c=k/2$ and $c'=c-4$.
No change of $s$ is made between different endpoints.

\subsection{The exact compatible high-jet graph}

Let $B_{s,D}$ be the original matrix of $S'^a$ in the complete
orthonormal Gamma frame $\phi_0,\ldots,\phi_D$, as in LJR2--3,
with $h_a=M_sa!(s/2)_a$ and $M_s=(2\pi)^{s/2}$. Define
\[
 \begin{gathered}
 (\mathbb V_N)_{j,(a,b)}=\beta_{ab}^{\,j}
                      \quad(0\le j\le N),\\
 (J_{E,N})_{ja}=\binom ja\mu_{j-a}
                 \quad(0\le j\le N,\ 0\le a\le D),\\
 \widetilde J_{E,N}=(J_{E,N})_{\{v,\ldots,N\},\{0,\ldots,D\}},
 \qquad P_{v,N}:\mathbb C^{D+1}\longrightarrow\mathbb C^{N+1}.
 \end{gathered}\tag{HJR2}
\]
Here $P_{v,N}$ inserts $v$ initial zero coordinates, and an
out-of-range term in the convolution is zero. The square
$\widetilde J_{E,N}$ is lower triangular with diagonal
$\mu_v\binom{v+a}{a}$ and is therefore invertible.

On the original coefficient subspace $\mathscr N_v$ put
\[
 \mathsf F_{N,s}
  =B_{s,D}^{-\mathsf T}\widetilde J_{E,N}^{-1}
                P_{v,N}^{\mathsf T}\mathbb V_N,\qquad
 \mathsf E_{N,s}=\mathsf F_{N,s}\mathsf R_s
                  =\begin{pmatrix}I_n\\ H_{N,s}\end{pmatrix}.
 \tag{HJR3}
\]
The right factor $\mathsf R_s$ is the exact low reconstruction
LJR4. Thus $\mathsf F_{N,s}$ takes original upper exponential
coefficients to their original native functional values through
degree $D$; $\mathsf E_{N,s}$ extends the first $n$ such values.
For a literal formula for the inverse convolution, if
$u_j=(\mathbb V_N\nu)_j$, then
\[
 f_a=\frac{v!}{\mu_v}a!
       \sum_{r=0}^{a}\kappa_{a-r}\frac{u_{v+r}}{(v+r)!},
 \qquad
 \mathsf F_{N,s}\nu=B_{s,D}^{-\mathsf T}(f_0,\ldots,f_D)^{\mathsf T}.
 \tag{HJR4}
\]
The coefficients $\kappa_a$ are all the actual inverse-symbol
coefficients of LJR6, with their proved finite recurrence.
The source coefficient transformation in HJR4 contains the
original phase $i^{-a}$ and every $h_a$.

To prove these statements, use the analytic germs
\[
 \mathcal N(z)=\sum\nu_{ab}e^{\beta_{ab}z},\qquad
 E_A(z)=(\mu_v/v!)z^vB_\infty(z).
\]
The order condition
$\nu\in\mathscr N_v$ makes the quotient analytic at zero.
Formal Taylor multiplication gives HJR4 and HJR2, and the
triangular $B_{s,D}^{-\mathsf T}$ changes ordinary moments to
orthonormal functional values. Its first $n$ moments and
orthonormal values depend only on the first $q$ jets $u_j$.
LJR6 therefore proves that the top block of HJR3 is $I_n$.
The bottom block $H_{N,s}$ is the specified remaining $e_N$ rows.

For a fully intrinsic description, write
$\chi(S)=S^q+\sum_{a=0}^{q-1}\chi_aS^a$. Define
$\mathcal U_{N,s}\subset\mathbb C^{D+1}$ by requiring that
$u=J_{E,N}B_{s,D}^{\mathsf T}w$ satisfy
\[
 u_{q+j}+\sum_{a=0}^{q-1}\chi_a u_{a+j}=0,
                      \qquad 0\le j<e_N .
 \tag{HJR5}
\]
Then $\mathcal U_{N,s}=\operatorname{im}\mathsf E_{N,s}$,
and its dimension is $n$. Indeed each original root is
annihilated by $\chi$, so HJR3 satisfies every recurrence.
Conversely the first $q$ entries of $u$, with its first $v$
zeros, determine a unique $\nu$ by the invertible upper
Vandermonde. The monic recurrences determine every later
$u_j$ uniquely and give $\mathbb V_N\nu$. The invertible
lower convolution block then determines $w$ uniquely.
This proves the graph equality, its dimension, and the
independence of its $e_N$ constraints. The free ambient later
functional values have not been identified with upper
exponential coefficients; the exact identification uses HJR5.

Let $V_{N,s}$ denote the $(D+1)$-by-$q'$ lower evaluation matrix
$(V_{N,s})_{a,\beta'}=\phi_a(\beta')$, and let $V=V_{q-1,s}$.
The complete lower-exponential identity gives
\[
 \mathsf E_{N,s}V=V_{N,s},\qquad
 \mathcal U_{N,s}/\operatorname{im}V_{N,s}
       \simeq\mathscr N_v/\operatorname{im}C^{\mathsf T}.
 \tag{HJR6}
\]
For its proof use $\nu=C^{\mathsf T}\eta$ in HJR4:
$\mathcal N/E_A$ is exactly the sum of the lower exponentials
with coefficients $\eta$, at all orders. Since $\mathsf E$
is injective, its induced quotient kernel is exactly
$\operatorname{im}V$. This proves the typed quotient
identification with the original denominators.

\subsection{One lower-evaluation minimum over the complete window}

Keep precisely DNE's fixed coefficient quotient frame $Z_{\exp}$.
Put $T_N=\mathsf F_{N,s}Z_{\exp}$, $T_0=T_{q-1}$, and
\[
 \begin{gathered}
 A=V^*V,\qquad K_0=A^{-1}V^*T_0,\qquad
 Z_0=T_0-VK_0,\qquad Q_0=VA^{-1/2},\\
 D_0=Z_0^*Z_0=\mathcal D_{q-1}^{(s)},\qquad
 V^*Z_0=0,\qquad Q_0^*Q_0=I_{q'}.
 \end{gathered}\tag{HJR7}
\]
The matrices $A$ and its positive square root are the actual
lower-evaluation Gram, retaining the complete root and Gamma
conditioning. No replacement evaluation metric is made.
With $H=H_{N,s}$, define
\[
 P_N=HZ_0,\qquad Q_N=HQ_0,\qquad
 \mathcal C_N=(I_{e_N}+Q_NQ_N^*)^{-1/2}P_N .
 \tag{HJR8}
\]
At $e_N=0$ all these matrices have zero rows. The exact native
metric at every endpoint is
\[
 \boxed{\quad
 \mathcal D_N^{(s)}
   =D_0+P_N^*(I_{e_N}+Q_NQ_N^*)^{-1}P_N
   =D_0+\mathcal C_N^*\mathcal C_N .
 \quad}\tag{HJR9}
\]
In particular its increment is positive and has rank
$\operatorname{rank}P_N\le\min\{\mathfrak m,e_N\}$.
The first endpoint update $N=q$ has rank at most one.
The last two endpoints have $e_{2q-1}=q$ and $e_{2q}=q+1$, respectively.

Here is a proof with the complete minimizer. Subtract the fixed
lower extension $\mathsf E VK_0$ from $T_N=\mathsf E T_0$.
Every remaining allowed extension is
$\mathsf E(Z_0x+Q_0t)$, with one vector $t\in\mathbb C^{q'}$
used for both its low and high rows. Its squared norm is
\[
 x^*D_0x+\|t\|^2+\|P_Nx+Q_Nt\|^2 .
 \tag{HJR10}
\]
The minimum occurs at
$t=-(I+Q_N^*Q_N)^{-1}Q_N^*P_Nx$.
Completing the square leaves
$x^*D_0x+x^*P_N^*[I-Q_N(I+Q_N^*Q_N)^{-1}Q_N^*]P_Nx$.
Multiplication by $I+Q_NQ_N^*$ shows that the bracket is
$(I+Q_NQ_N^*)^{-1}$, proving HJR9. DNE7 proves that these
and only these lower evaluations are the permitted native
extension differences. Thus HJR10 is exactly its full
minimum, without separate minima on different row blocks.

\subsection{The original kernel and annihilator updates}

Write $B=B_K$ and $J=J_0$ for the actual maps in PEC19,24,
so $B$ is onto, $J$ is injective, and
$\operatorname{im}J=\ker B$. Define the full initial forms
and minimum lift by
\[
 \begin{gathered}
 D_{B,0}=(BD_0^{-1}B^*)^{-1},\qquad
 D_{J,0}=J^*D_0J,\\
 L_0=D_0^{-1}B^*D_{B,0},\qquad
 \mathsf W_0=J(D_{J,0})^{-1/2}.
 \end{gathered}\tag{HJR11}
\]
Then $BL_0=I$, $B\mathsf W_0=0$,
$\mathsf W_0^*D_0\mathsf W_0=I$, and $\mathsf W_0^*D_0L_0=0$.
These are literal original kernel and annihilator maps;
the displayed $\mathsf W_0$ is a coefficient matrix for
$\ker B$ with its full $D_0$ metric.
The induced exact updates are
\[
 \begin{gathered}
 D_{J,N}=D_{J,0}+(\mathcal C_NJ)^*(\mathcal C_NJ),\\
 D_{B,N}=D_{B,0}
 +(\mathcal C_NL_0)^*
  [I+\mathcal C_N\mathsf W_0\mathsf W_0^*\mathcal C_N^*]^{-1}
                                      (\mathcal C_NL_0).
 \end{gathered}\tag{HJR12}
\]
The first is restriction of HJR9. For the second, every
$B$-lift of $y$ is $L_0y+\mathsf W_0t$. Its low squared
norm is $y^*D_{B,0}y+\|t\|^2$, and its added squared norm
is $\|\mathcal C_NL_0y+\mathcal C_N\mathsf W_0t\|^2$.
The same completion as HJR10 proves the second formula.
It preserves both the original lower-evaluation denominator
in $\mathcal C_N$ and the original kernel denominator in HJR12.
Empty source, kernel or target ranks use their unique empty
matrices and determinant one.

The exact identities DNE9 and PEC25 are
\[
 D_{B,N}=\overline{(\widehat H_{K,N}^{(s)})^{-1}},
 \qquad D_{J,N}=\overline{Q_{0,N}^{-1}}.
 \tag{HJR13}
\]
Thus HJR12 is a full original native kernel and complementary
residual calculation, rather than a freely chosen coefficient
compression.

\subsection{Exact one-row updates and their allocation}

The original orthonormal polynomial sequence is independent
of its finite cutoff. At any $M$ with $q-1\le M<2q$, let
$d=M-v$, and write $T,V$ for the complete matrices through
degree $d$. Their next rows, at degree $d+1$, are $t,v_+$.
Every entry of $t$ is
$\Lambda_{\mathcal N_a/E_A}(\phi_{d+1})$, computed by
HJR4 and the original monomial matrix $B_{s,d+1}$ through
degree $d+1$; every entry of $v_+$ is
$\phi_{d+1}(\beta')$. Set
\[
 \begin{gathered}
 A_M=V^*V,\quad K_M=A_M^{-1}V^*T,\quad
 r_M=t-v_+K_M,\\
 \lambda_M=1+v_+A_M^{-1}v_+^*>0,\qquad
 h_M=\lambda_M^{-1/2}r_M .
 \end{gathered}\tag{HJR14}
\]
Then the exact update is
\[
 D_{M+1}=D_M+h_M^*h_M,\qquad D_M=\mathcal D_M^{(s)}.
 \tag{HJR15}
\]
To prove this, subtract the attained old minimum from the
new extension problem. The old residual is orthogonal to
all old lower evaluations, so the remaining minimization
for coefficient correction $\eta$ is
$\eta^*A_M\eta+\|r_Mx-v_+\eta\|^2$.
Its minimizer is
$(A_M+v_+^*v_+)^{-1}v_+^*r_Mx$.
The same identity used in HJR10 gives the remaining form
$r_M^*r_M/(1+v_+A_M^{-1}v_+^*)$, proving HJR15.
Every evaluation row has participated in the old $A_M$;
the update does not reset the lower-evaluation coefficient.

At this endpoint put
$D_{B,M}=(BD_M^{-1}B^*)^{-1}$,
$L_M=D_M^{-1}B^*D_{B,M}$ and $D_{J,M}=J^*D_MJ$.
The exact inverse-metric decomposition is
\[
 D_M^{-1}
 =J(D_{J,M})^{-1}J^*
          +L_M(D_{B,M})^{-1}L_M^* .
 \tag{HJR16}
\]
Indeed $[J,L_M]$ is a full coefficient frame whose Gram in
$D_M$ is $\operatorname{diag}(D_{J,M},D_{B,M})$.
The off-diagonal blocks vanish because $BJ=0$.
Inverting that frame Gram proves HJR16.
Define the two actual nonnegative scalar row energies
\[
 x_M=h_MJ(D_{J,M})^{-1}J^*h_M^*,\qquad
 y_M=h_ML_M(D_{B,M})^{-1}L_M^*h_M^* .
 \tag{HJR17}
\]
Both retain $\lambda_M^{-1}$ from HJR14 and every original
period-dependent kernel coordinate. The determinant lemma,
proved by extending a nonzero vector to a basis or by
expanding a rank-one determinant, now gives
\[
 \begin{gathered}
 g_M=\log\frac{\det D_{M+1}}{\det D_M}
                  =\log(1+x_M+y_M),\\
 g_{0,M}=\log\frac{\det D_{J,M+1}}{\det D_{J,M}}
                  =\log(1+x_M),\\
 g_{K,M}=\log\frac{\det D_{B,M+1}}{\det D_{B,M}}
                  =\log\left(1+\frac{y_M}{1+x_M}\right),\qquad
 g_M=g_{0,M}+g_{K,M}.
 \end{gathered}\tag{HJR18}
\]
The restriction update gives the second equality directly.
For the quotient use the one-row instance of HJR12: its
added form is $(h_ML_M)^*(h_ML_M)/(1+x_M)$.
Its determinant relative to $D_{B,M}$ proves the third.
This proves the allocation of each full original row
increment between the actual kernel quotient and its
annihilator, with its complete positive denominator.

\subsection{The four original signs and an exterior bound at scale \(kq\)}

For $X$ equal to the full space, $K$, or $0$, define
\[
 I_{X,N}=\log\frac{\det D_{X,N}}{\det D_{X,q-1}},
 \qquad I_N=I_{{\rm full},N}.
 \tag{HJR19}
\]
All these numbers are nonnegative and increase with $N$.
At each $N$, completing $J$ by any fixed right inverse of
$B$ factors the full Gram into its restriction and quotient.
The fixed frame determinant cancels between endpoints, so
\[
 I_N=I_{K,N}+I_{0,N},\qquad
 I_{X,N}=\sum_{M=q-1}^{N-1}g_{X,M}.
 \tag{HJR20}
\]
In the sum for the full space $g_{{\rm full},M}=g_M$.
The source dual identities HJR13 therefore give exactly
\[
 \begin{gathered}
 \widehat F_k^{(s)}
       =I_{K,2q-1}+I_{K,2q}-I_{K,q},\\
 F_{0,k,s}=I_{0,2q-1}+I_{0,2q}-I_{0,q},\\
 \mathcal T_{k,s}:=\widehat F_k^{(s)}+F_{0,k,s}
       =I_{2q-1}+I_{2q}-I_q.
 \end{gathered}\tag{HJR21}
\]
Equivalently every line is the weighted sum of its original
row increments, with weights
\[
 w_M=
 \begin{cases}
 1,&M=q-1\text{ or }M=2q-1,\\
 2,&q\le M\le2q-2.
 \end{cases}
 \tag{HJR22}
\]
For example $I_{2q-1}$ contains all rows through $2q-2$,
$I_{2q}$ adds the last row, and $I_q$ removes one copy of
the first. This proves all endpoint signs.
Writing $a=I_q$, $b=I_{2q-1}-I_q$ and
$c=I_{2q}-I_{2q-1}$ gives
$\mathcal T=a+2b+c$ with $a,b,c\ge0$. In particular
\[
 0\le I_N\le I_{2q}\le\mathcal T_{k,s}
                         \quad(q-1\le N\le2q).
 \tag{HJR23}
\]

To give its finite original bound, retain
$\delta_{K,k,s}=\widehat F_k^{(s)}-F_K^{(s)}$ with
$|\delta_{K,k,s}|\le\mathcal E_{k,s}^{\rm FEX}$,
and precisely $\delta_{0,k,s}$ and $F_{\Xi(L)}^{(s)}$
from PEC29. The latter is the original restriction return
called $F_{W_0}^{(s)}$ there.
OAR3 and PEC29 give the exact equality
\[
 \mathcal T_{k,s}
 =\mathcal F_{\partial,k}^{(s)}
    +\mathfrak e_{k,s}+\delta_{K,k,s}
            -F_{\Xi(L)}^{(s)}-\delta_{0,k,s}.
 \tag{HJR24}
\]
The original finite signs are consequently
\[
 \begin{split}
 \max\{0,\mathcal F_{\partial,k}^{(s)}
       -e_{k,s}^--\mathcal E_{k,s}^{\rm FEX}
       -E_{k,s}^{\rm tr}-(v-t_0)\mathcal L_s\}
 &\le\mathcal T_{k,s}\\
 &\le U_{k,s}^{\rm jet}
 :=\mathcal F_{\partial,k}^{(s)}
       +e_{k,s}^++\mathcal E_{k,s}^{\rm FEX}+E_{k,s}^{\rm tr}.
 \end{split}\tag{HJR25}
\]
All errors on the right and left are the explicit accepted
FEX, OAR and PEC expressions. The inequality uses
$0\le F_{\Xi(L)}^{(s)}\le(v-t_0)\mathcal L_s$ and
$|\delta_0|\le E^{\rm tr}$ only after retaining their
actual values in HJR24. In particular
\[
 \mathcal T_{k,s}
       =\Delta q\mathfrak C_\partial+o(kq),\qquad
 U_{k,s}^{\rm jet}
       =\Delta q\mathfrak C_\partial+o(kq).
 \tag{HJR26}
\]
This follows from the fully proved OFV/OVR outer first
variation and the stated orders of the finite errors.
The actual period remains fixed; no uniformity over a
moving period is asserted.

Now take any paired restriction, quotient, or subquotient
of the common coefficient space at the two metrics
$(D_0,D_N)$. Let its rank be $t$ and its determinant
increment be $I_{t,N}$. These are attained restrictions
and quotients along the same algebraic flag, retaining
all its complete metrics. Since $D_N\ge D_0$, every
such form increases. Full flag factorization, with
nonnegative log increments on both complementary flags,
proves
\[
 0\le I_{t,N}\le I_N\le U_{k,s}^{\rm jet}=O_{\rm actual}(kq).
 \tag{HJR27}
\]
For the same reason all generalized eigenvalues for
$(D_N,D_0)$ are at least one; the sum of all their
logarithms is $I_N$. Hence every exterior degree of the
identity between these two metrics has norm at most
$\exp(U_{k,s}^{\rm jet}/2)$. The same statement holds
on every specified paired flag. This is a full exterior
bound with no additional rank multiplier and with the
complete lower-evaluation and kernel conditioning retained.

\subsection{Exact cancellation of the low leading scale}

Retain LJR's actual finite remainder $\mathscr R_{t,k,s}$,
including its full-space formula LJR28 and the two-sided
flag bounds LJR31. Keep the fixed $G_{\exp}$ quotient
or restriction associated with this same flag.
Then HJR19 gives the exact identity at every endpoint
\[
 \log\frac{\det G_t}{\det D_{t,N}}
       =2tq\log k+\mathscr R_{t,k,s}-I_{t,N}.
 \tag{HJR28}
\]
HJR27 therefore propagates the same low leading scale
$2tq\log k$ to all four endpoints, with remainder
$O_{\rm actual}(kq)$. It does so by the recurrence-compatible
graph and the full Schur minima, not by assigning the
later free jets the low dimension.

Let $p_{t,N},n_{t,N}$ be PEC's two nonnegative logarithmic
spectral parts for $(D_{t,N},G_t)$. The same identity is
\[
 n_{t,N}
   =2tq\log k+\mathscr R_{t,k,s}-I_{t,N}+p_{t,N},
 \qquad 0\le p_{t,N}\le\varepsilon_{N,s}^{\vee}=o(kq).
 \tag{HJR29}
\]
At the actual kernel quotient and annihilator restriction,
write $X=K,0$ respectively. Substituting HJR29 with the
four original signs cancels the four copies of both
$2tq\log k$ and $\mathscr R_{t,k,s}$ exactly and gives
\[
 \begin{gathered}
 A_{X,k,s}
       =I_{X,2q-1}+I_{X,2q}-I_{X,q}-P_{X,k,s},\\
 P_{X,k,s}=p_{X,2q-1}+p_{X,2q}-p_{X,q-1}-p_{X,q},\qquad
 0\le P_{X,k,s}\le E_{k,s}^{\vee}.
 \end{gathered}\tag{HJR30}
\]
The signs of the $p$ terms follow by substituting their
displayed positions in HJR29. Their monotonicity and
finite bound are PEC21,26 on these same maps.
Together with HJR18 and HJR22 this supplies the explicit
original calculation to carry forward:
\[
 \begin{gathered}
 A_{K,k,s}
   =\sum_{M=q-1}^{2q-1}w_M
          \log\left(1+\frac{y_M}{1+x_M}\right)-P_{K,k,s},\\
 A_{0,k,s}
   =\sum_{M=q-1}^{2q-1}w_M\log(1+x_M)-P_{0,k,s}.
 \end{gathered}\tag{HJR31}
\]
The same original $x_M,y_M$ sum in the full row energy
$\log(1+x_M+y_M)$. Every numerator, denominator and
moment used to compute these scalars is fixed in
HJR2--4,HJR14,HJR16--17. Their weighted total has the
proved leading value HJR26; their separate allocation is
retained by the exact formulas HJR31.
Thus the next-order $kq$ problem is the specified allocation
of these original Gamma row energies, with all low
factorial and Vandermonde terms canceled only through
the exact equality HJR28.

% END CURRENT HJR

\clearpage
% BEGIN CURRENT HCR
% BEGIN HCR COMPLETE RECEIVING PROOF
\section{The compatible high rows in the original receiving interval}
\label{sec:HCR}

Retain the full fixed-period simple-quartet hypotheses and all
original source, coefficient and metric data in PBR and LRR.
For each original source $s=1,k$ the following uses exactly
$B=B_K=Z_k^{\mathsf T}A_\ell$ and
$J=J_0=A_\ell^{-1}P_Y^{\mathsf T}$, with
$\operatorname{im}J=\ker B$ and ranks $t_K,t_R$ of LRR1.
In this section $D_M=\mathcal D_M^{(s)}$ is the full original
dual metric; it is distinct from the original lower-grid polynomial
$\chi'$ and from the residual-row matrix in PBR12.

The complete recurrence-compatible high-jet graph is
$\mathsf E_{N,s}=\mathsf F_{N,s}\mathsf R_s
=\begin{pmatrix}I_n\\ H_{N,s}\end{pmatrix}$,
where $\mathsf R_s$ is the original LJR4 reconstruction and
$\mathsf F_{N,s}$ is the full inverse-convolution and original
Gamma coordinate map HJR3--4. HJR5 proves that its image is
exactly the vectors whose convolved moments satisfy the monic
$\chi$ recurrence. HJR6 proves
$\mathsf E_{N,s}V_{q-1,s}=V_{N,s}$, so the entire lower
evaluation denominator is transported through the whole graph.

At a single original degree $q-1\le M<2q$ use the full rows
$T,V$ through degree $M-v$ and the actual next rows $t,v_+$.
Their entries are the original functional values of
$\mathcal N_a/E_A$ and the original lower-root evaluations
on the degree-$M-v+1$ Gamma orthonormal polynomial.
The exact full conditional update is
\[
\begin{gathered}
 A_M=V^*V,\quad K_M=A_M^{-1}V^*T,\quad
 r_M=t-v_+K_M,\quad \lambda_M=1+v_+A_M^{-1}v_+^*>0,\\
 h_M=\lambda_M^{-1/2}r_M,\qquad
 D_{M+1}=D_M+h_M^*h_M .
\end{gathered}\tag{HCR1}
\]
Indeed subtracting the attained old extension minimum leaves
the quadratic minimization
$\eta^*A_M\eta+\|r_Mx-v_+\eta\|^2$.
Its minimum is
$\|r_Mx\|^2/(1+v_+A_M^{-1}v_+^*)$, by completing the square
with $(A_M+v_+^*v_+)^{-1}$.
The old lower coefficient is consequently retained throughout
the complete row block; it is not minimized separately at
each piece of the graph. This proves HCR1 with the full
conditioning denominator and every original moment.

Put
$D_{B,M}=(BD_M^{-1}B^*)^{-1}$,
$D_{J,M}=J^*D_MJ$, and
$L_M=D_M^{-1}B^*D_{B,M}$.
Then the exact inverse decomposition and two row energies are
\[
\begin{gathered}
 D_M^{-1}=J D_{J,M}^{-1}J^*
                  +L_M D_{B,M}^{-1}L_M^*,\\
 x_M=h_MJ D_{J,M}^{-1}J^*h_M^*\ge0,\qquad
 y_M=h_ML_M D_{B,M}^{-1}L_M^*h_M^*\ge0,\\
 \log\frac{\det D_{J,M+1}}{\det D_{J,M}}=\log(1+x_M),\\
 \log\frac{\det D_{B,M+1}}{\det D_{B,M}}
                 =\log\left(1+\frac{y_M}{1+x_M}\right).
\end{gathered}\tag{HCR2}
\]
The columns $[J,L_M]$ form a full coefficient frame because
$BL_M=I$ and $\operatorname{im}J=\ker B$.
Its $D_M$ Gram is diagonal in these two blocks since
$J^*D_ML_M=J^*B^*D_{B,M}=0$.
Inverting this exact frame Gram proves the first line.
Restriction of HCR1 and the rank-one determinant lemma proves
the third line. For the quotient, every lift of a vector $z$
is $L_Mz+Jt$. Minimize its full updated norm over $t$.
After using the original $D_M$ orthogonal decomposition, the
additional form is
$(h_ML_M)^*(h_ML_M)/(1+x_M)$.
The determinant lemma proves the final line, including the
denominator $1+x_M$. Both $x_M,y_M$ also retain
$\lambda_M^{-1}$ in HCR1.

Use the original weights $w_{q-1}=w_{2q-1}=1$ and
$w_M=2$ for $q\le M\le2q-2$, and define
\[
\begin{gathered}
 \mathcal V_{K,k,s}^{\rm jet}
   =\sum_{M=q-1}^{2q-1}w_M
                      \log\left(1+\frac{y_M}{1+x_M}\right),\qquad
 \mathcal V_{0,k,s}^{\rm jet}
   =\sum_{M=q-1}^{2q-1}w_M\log(1+x_M),\\
 \mathcal V_{K,k,s}^{\rm jet}=A_{K,k,s}+P_{K,k,s},\qquad
 \mathcal V_{0,k,s}^{\rm jet}=A_{0,k,s}+P_{0,k,s},\\
 F_K^{(s)}
   =\mathcal V_{K,k,s}^{\rm jet}-\mathcal E_s^{\rm flag}
                                      +F_L^{(s)}-F_{W_0}^{(s)},\\
 \Phi_{k,s}=F_{W_0}^{(s)}+\mathcal V_{0,k,s}^{\rm jet}
                                      +\delta_{0,k,s}.
\end{gathered}\tag{HCR3}
\]
For each of the two metrics let
$I_{X,N}=\log(\det D_{X,N}/\det D_{X,q-1})$.
Summing the consecutive row identities in HCR2 gives
$I_{X,N}=\sum_{M=q-1}^{N-1}g_{X,M}$.
The expression
$I_{X,2q-1}+I_{X,2q}-I_{X,q}$ therefore has precisely the
weights displayed above, since the first row loses one
copy and the last row occurs once. On the other hand the
full dual identities
$D_{B,N}=\overline{\widehat H_{K,N}^{-1}}$ and
$D_{J,N}=\overline{Q_{0,N}^{-1}}$ identify that weighted
sum with the four original native returns. PBR5 then proves
the second line. Equivalently HJR29--31 cancels the common
low reconstruction scale and remainder in the four signs
and gives $A_X=\mathcal V_X^{\rm jet}-P_X$ exactly.
Substitution into PBR9/7 proves the final two lines:
the $P_X$ terms cancel algebraically, rather than being
estimated or set to zero. Both low terms remain separately,
with their original metrics from PBR6.

The full finite errors are unchanged:
$\mathcal E_s^{\rm flag}\in[D_s^-,D_s^+]$ of PBR8 and
$|\delta_{0,k,s}|\le E_{k,s}^{\rm tr}$ of PBR7, including
$\alpha_A$ and all four CEP errors.
Define the exact row-sum endpoints
\[
\begin{gathered}
 a_s^{\rm row,-}=\mathcal V_{K,k,s}^{\rm jet}-D_s^+
                                      +\lambda_s^{\rm low},\qquad
 a_s^{\rm row,+}=\mathcal V_{K,k,s}^{\rm jet}-D_s^-
                                      +\lambda_s^{\rm low},\\
 b_s^{\rm row,-}=F_{W_0}^{(s)}+\mathcal V_{0,k,s}^{\rm jet}
                                      -E_{k,s}^{\rm tr},\qquad
 b_s^{\rm row,+}=F_{W_0}^{(s)}+\mathcal V_{0,k,s}^{\rm jet}
                                      +E_{k,s}^{\rm tr},\\
 a_s^{\rm row,-}\le F_K^{(s)}\le a_s^{\rm row,+},\qquad
 b_s^{\rm row,-}\le\Phi_{k,s}\le b_s^{\rm row,+},\\
 \Phi_{k,s}^{-,{\rm row}}
    =\max\{\Phi_{k,s}^{-,\vee},b_s^{\rm row,-}\},\qquad
 \Phi_{k,s}^{+,{\rm row}}
    =\min\{\Phi_{k,s}^{+,\vee},b_s^{\rm row,+}\},\\
 L_s^{\rm row}
   =\max\{L_s^\vee,a_s^{\rm row,-},
       \mathcal F_{\partial,k}^{(s)}
             -\Phi_{k,s}^{+,{\rm row}}-e_{k,s}^-\},\\
 U_s^{\rm row}
   =\min\{U_s^\vee,a_s^{\rm row,+},
       \mathcal F_{\partial,k}^{(s)}
             -\Phi_{k,s}^{-,{\rm row}}+e_{k,s}^+\}.
\end{gathered}\tag{HCR4}
\]
Here $\lambda_s^{\rm low}=F_L^{(s)}-F_{W_0}^{(s)}$ is the
unchanged difference of the two separately retained volumes.
The first two interval assertions follow immediately from
the exact HCR3 equalities and the stated finite error bounds.
Subtracting the resulting residual interval in the exact
OAR equality proves the last two lines.
Every entry bounds the same original scalar; hence the
intersection retains all prior branches and is nonempty.
On the exact PMR range $k\ge29$,
$k\sqrt{\delta^2+\gamma^2}\le2^{-33}(k+1)^2$, also include
$\mathcal F_{0,k}^{(s)}-\Phi_{k,s}^{+,{\rm row}}-e_{k,s}^{0,-}$
in the maximum and
$\mathcal F_{0,k}^{(s)}-\Phi_{k,s}^{-,{\rm row}}+e_{k,s}^{0,+}$
in the minimum, where $e^{0,\pm}=e^\pm+c_k^\circ$.
Outside that range the original OAR entries remain.
Because $0\le P_X\le E_{k,s}^{\vee}$, the exact kernel
row-sum interval lies inside PBR11's corresponding
$A_K$ interval, and its residual interval lies inside
PBR11's corresponding $A_0$ interval. This containment
follows by substituting $\mathcal V_X=A_X+P_X$ and
comparing the two endpoints. It uses no distribution
coefficient for $x_M$ or $y_M$.

On the original arithmetic period domain $|u|\ge R_*$,
retain precisely $a_{K,k}^\pm$ and all their finite
constants from PBR14. The complete current arithmetic,
signed and action receivers are
\[
\begin{gathered}
 L_{\rm ar}^{\rm row}
   =\max\{L_{\rm ar}^\vee,L_1^{\rm row}+a_{K,k}^-\},\qquad
 U_{\rm ar}^{\rm row}
   =\min\{U_{\rm ar}^\vee,U_1^{\rm row}+a_{K,k}^+\},\\
 F_K^{\rm ar}\in[L_{\rm ar}^{\rm row},U_{\rm ar}^{\rm row}],
 \quad F_B^{\rm ar}\in
 [\mathcal B_k^{\rm ar}-U_{\rm ar}^{\rm row},
  \mathcal B_k^{\rm ar}-L_{\rm ar}^{\rm row}],\\
 L_k^{\rm row}+\mathcal H_k+\delta_k^\sigma
 \le\Psi_k\le U_k^{\rm row}+\mathcal H_k+\delta_k^\sigma,\\
 L_J^{\rm row}=L_k^{\rm row}+\mathcal H_k+\delta_k^\sigma
                                   +C_k^{\rm budget},\qquad
 U_J^{\rm row}=U_k^{\rm row}+\mathcal H_k+\delta_k^\sigma
                                   +C_k^{\rm budget},\\
 4q\log\mathcal L_{h,k}\le J_k^{\rm action}
       \in[L_J^{\rm row},U_J^{\rm row}],\qquad
 \mathcal L_{h,k}\le\min_M\epsilon_M^{\rm ar}
 \le\prod_M(\epsilon_M^{\rm ar})^{w_M/(2q)}
 \le e^{U_J^{\rm row}/(4q)} .
\end{gathered}\tag{HCR5}
\]
The arithmetic source difference is the identical correlated
scalar used in PBR14; add its exact interval and intersect.
Subtracting from the unchanged total reverses the boundary
endpoints. The signed equality adds $\mathcal H_k=-T_k$ and
the original arithmetic correction. Adding the complete
$C_k^{\rm budget}=F_B^{(k)}+\mathcal W_k^{\rm ar}
-\mathcal P_--\mathcal P_+$ then gives the action interval.
The original untilted reflection and the same-class lower
radius give the weighted geometric mean exactly as in
PBR15. Every mixed pairing, both contraction penalties,
the full monic window and action leakage remain. Under
a real tilt use the unchanged full MCE phase formula.

Finally the exact same-row allocation is
\[
\begin{gathered}
 \mathcal V_{K,k,s}^{\rm jet}+\mathcal V_{0,k,s}^{\rm jet}
   =\sum_{M=q-1}^{2q-1}w_M\log(1+x_M+y_M),\\
 F_K^{(s)}+\Phi_{k,s}
   =\mathcal V_{K,k,s}^{\rm jet}+\mathcal V_{0,k,s}^{\rm jet}
        -\mathcal E_s^{\rm flag}+F_L^{(s)}+\delta_{0,k,s},\\
 \mathcal V_{K,k,s}^{\rm jet}+\mathcal V_{0,k,s}^{\rm jet}
      =\Delta q\mathfrak C_\partial+o(kq).
\end{gathered}\tag{HCR6}
\]
The first equality is
$(1+x_M)(1+y_M/(1+x_M))=1+x_M+y_M$ for every actual row.
The second adds the two HCR3 identities, cancelling the
identical original $W_0$ term only in that sum.
The final equality is HJR24--26 with every finite FEX,
transfer and outer error retained before the fixed-period
limit. The individual allocations are exactly the two
specified sums in HCR3; no individual leading coefficient
is assigned by the evaluated total.
% END HCR COMPLETE RECEIVING PROOF
% END CURRENT HCR

\end{document}
